\documentclass{article}
\usepackage{amsmath,amssymb,amsthm}
\IfFileExists{hyperref.sty}{%
  \IfFileExists{color.sty}{\usepackage{color}\definecolor{paperlinkblue}{rgb}{0.10,0.25,0.42}%
    \usepackage[colorlinks=true,linkcolor=paperlinkblue]{hyperref}}{\usepackage[hidelinks]{hyperref}}%
}{}
\ifdefined\hyperlink
  \newcommand{\paperanchor}[1]{\hypertarget{#1}{}}%
  \newcommand{\paperlink}[2]{\hyperlink{#1}{\begin{NoHyper}#2\end{NoHyper}}}%
\else
  \ifdefined\pdfextension
    \ifnum\outputmode>0
      \def\paperdest{\pdfextension dest}\def\paperstart{\pdfextension startlink}%
      \def\paperend{\pdfextension endlink}\def\papercolor{\pdfextension colorstack}%
    \fi
  \else\ifdefined\pdfstartlink
    \ifnum\pdfoutput>0
      \let\paperdest\pdfdest\let\paperstart\pdfstartlink\let\paperend\pdfendlink\let\papercolor\pdfcolorstack
    \fi
  \fi\fi
  \ifdefined\paperdest
    \newcommand{\paperanchor}[1]{\paperdest name{#1} xyz}%
    \newcommand{\paperlink}[2]{\leavevmode\paperstart attr{/Border [0 0 0]} goto name{#1}%
      \papercolor 0 push {0.10 0.25 0.42 rg 0.10 0.25 0.42 RG}#2\papercolor 0 pop\paperend}%
  \else\ifdefined\XeTeXversion
    \newcommand{\paperanchor}[1]{\special{pdf:dest (#1) [@thispage /XYZ null @ypos null]}}%
    \newcommand{\paperlink}[2]{\leavevmode\special{pdf:bann << /Subtype /Link /Border [0 0 0] /A << /S /GoTo /D (#1) >> >>}%
      \special{color push rgb 0.10 0.25 0.42}#2\special{color pop}\special{pdf:eann}}%
  \else
    \PackageError{lean-to-latex}{No PDF hyperlink backend available}{Use XeLaTeX, pdfLaTeX, LuaLaTeX, or install hyperref.}%
  \fi\fi
\fi
\newcommand{\paperlineanchor}[1]{\raisebox{10pt}[0pt][0pt]{\paperanchor{#1}}}
\DeclareRobustCommand{\paperresultname}[1]{\paperlineanchor{theorem:\arabic{theorem}}#1}
\newtheoremstyle{paperresult}{8pt}{5pt}{\normalfont\raggedright}{}{\bfseries}{.}{\newline}{}
\theoremstyle{paperresult}
\newtheorem{theorem}{\paperresultname{Theorem}}
\newcounter{paperpremise}
\newtheorem{lemma}[theorem]{\paperresultname{Lemma}}
\newtheorem{corollary}[theorem]{\paperresultname{Corollary}}
\newcounter{definingequation}
\newenvironment{proofstep}{\par\smallskip\ifdim\leftskip<3em\relax\advance\leftskip by 1em\relax\fi\noindent\ignorespaces}{\par\smallskip}
\setlength{\textwidth}{6.5in}
\setlength{\oddsidemargin}{0in}
\AddToHook{env/proof/begin}{\raggedright}
\setlength{\emergencystretch}{\linewidth}
\title{Cardinalities of finite sets}
\author{}
\date{}
\begin{document}
\maketitle
\section*{Definitions}
\paragraph{Cardinality.} \par\noindent\refstepcounter{definingequation}\label{defeq:1}\paperlineanchor{defeq:1}\textup{(\paperlink{defeq:1}{D1})} Let \(\alpha\) be a type. Let \(s\) be a finite subset of \(\alpha\). Then \(\left|s\right| = \left|\operatorname{underlying\,multiset}(s)\right|\).

\paragraph{Strong induction.} \par\noindent\refstepcounter{definingequation}\label{defeq:2}\paperlineanchor{defeq:2}\textup{(\paperlink{defeq:2}{D2})} Let \(\alpha\) be a type. Let \(p\) be a function of type \(\mathcal{P}_{\mathrm{fin}}(\alpha) \to \mathsf{Sort}\). Let \(H\) be a function of type \(\Pi_{s\, :\, \mathcal{P}_{\mathrm{fin}}(\alpha)}\, (\Pi_{t\, :\, \mathcal{P}_{\mathrm{fin}}(\alpha)}\, t \subsetneq s \to p(t)) \to p(s)\). Let \(x\) be a finite subset of \(\alpha\). Then \(\operatorname{strong\,Induction}(H,\allowbreak{} x) = H(x,\allowbreak{} t \mapsto h \mapsto \operatorname{strong\,Induction}(H,\allowbreak{} t))\).

\paragraph{Strong induction on.} \par\noindent\refstepcounter{definingequation}\label{defeq:3}\paperlineanchor{defeq:3}\textup{(\paperlink{defeq:3}{D3})} Let \(\alpha\) be a type. Let \(p\) be a function of type \(\mathcal{P}_{\mathrm{fin}}(\alpha) \to \mathsf{Sort}\). Let \(s\) be a finite subset of \(\alpha\). Let \(H\) be a function of type \(\Pi_{s_{1}\, :\, \mathcal{P}_{\mathrm{fin}}(\alpha)}\, (\Pi_{t\, :\, \mathcal{P}_{\mathrm{fin}}(\alpha)}\, t \subsetneq s_{1} \to p(t)) \to p(s_{1})\). Then \(\operatorname{strong\,Induction\,On}\left(s, H\right) = \operatorname{strong\,Induction}\left(H, s\right)\).

\paragraph{Strong downward induction.} \par\noindent\refstepcounter{definingequation}\label{defeq:4}\paperlineanchor{defeq:4}\textup{(\paperlink{defeq:4}{D4})} Let \(\alpha\) be a type. Let \(p\) be a function of type \(\mathcal{P}_{\mathrm{fin}}(\alpha) \to \mathsf{Sort}\). Let \(n\) be a natural number. Let \(H\) be a function of type \(\Pi_{t_{1}\, :\, \mathcal{P}_{\mathrm{fin}}(\alpha)}\, (\Pi_{t_{2}\, :\, \mathcal{P}_{\mathrm{fin}}(\alpha)}\, |t_{2}| \le n \to t_{1} \subsetneq t_{2} \to p(t_{2})) \to |t_{1}| \le n \to p(t_{1})\). Let \(x\) be a finite subset of \(\alpha\). Then \(\operatorname{strong\,Downward\,Induction}(H,\allowbreak{} x) = H(x,\allowbreak{} t \mapsto \mathit{ht} \mapsto h \mapsto  \operatorname{strong\,Downward\,Induction}(H,\allowbreak{} t)|_{\mathit{this}=\operatorname{Finset\mathord{:}card\,lt\,card}(h)})\).

\paragraph{Strong downward induction on.} \par\noindent\refstepcounter{definingequation}\label{defeq:5}\paperlineanchor{defeq:5}\textup{(\paperlink{defeq:5}{D5})} Let \(\alpha\) be a type. Let \(n\) be a natural number. Let \(p\) be a function of type \(\mathcal{P}_{\mathrm{fin}}(\alpha) \to \mathsf{Sort}\). Let \(s\) be a finite subset of \(\alpha\). Let \(H\) be a function of type \(\Pi_{t_{1}\, :\, \mathcal{P}_{\mathrm{fin}}(\alpha)}\, (\Pi_{t_{2}\, :\, \mathcal{P}_{\mathrm{fin}}(\alpha)}\, |t_{2}| \le n \to t_{1} \subsetneq t_{2} \to p(t_{2})) \to |t_{1}| \le n \to p(t_{1})\). Then \(\operatorname{strong\,Downward\,Induction\,On}(s,\allowbreak{} H) = \operatorname{strong\,Downward\,Induction}(H,\allowbreak{} s)\).

\subsection*{Prerequisites}
\begin{description}
\item[\paperlineanchor{premise:436c6173736963616c2e6f725f6966665f6e6f745f696d705f6c656674}\paperlink{premise:436c6173736963616c2e6f725f6966665f6e6f745f696d705f6c656674}{P1.}] \refstepcounter{paperpremise}\label{premise:436c6173736963616c2e6f725f6966665f6e6f745f696d705f6c656674}\textbf{Or if and only if not imp left.} Let \(a\) and \(b\) be propositions. Then \(a \lor b \iff \neg a \implies b\).
\item[\paperlineanchor{premise:446563696461626c652e6e6f745f6e6f74}\paperlink{premise:446563696461626c652e6e6f745f6e6f74}{P2.}] \refstepcounter{paperpremise}\label{premise:446563696461626c652e6e6f745f6e6f74}\textbf{Not not.} Let \(p\) be a proposition. Assume \(p\) is decidable. Then \(\neg \left(\neg p\right) \iff p\).
\item[\paperlineanchor{premise:45712e6765}\paperlink{premise:45712e6765}{P3.}] \refstepcounter{paperpremise}\label{premise:45712e6765}\textbf{Equality implies the reverse weak inequality.} Let \(\alpha\) be a type. Equip \(\alpha\) with a preorder. Let \(a\) and \(b\) be elements of \(\alpha\). Assume \(a = b\). Then \(b \le a\).
\item[\paperlineanchor{premise:45712e73756273657427}\paperlink{premise:45712e73756273657427}{P4.}] \refstepcounter{paperpremise}\label{premise:45712e73756273657427}\textbf{Alias.} Let \(\alpha\) be a type. Equip \(\alpha\) with a subset relation. Let \(a\) and \(b\) be elements of \(\alpha\). Assume \((\mathit{x1} \mapsto \mathit{x2} \mapsto \mathit{x1} \subseteq \mathit{x2})\text{ is reflexive on }\alpha\). Assume \(a = b\). Then \(a \subseteq b\).
\item[\paperlineanchor{premise:45712e73796d6d}\paperlink{premise:45712e73796d6d}{P5.}] \refstepcounter{paperpremise}\label{premise:45712e73796d6d}\textbf{Symm.} Let \(\alpha\) be a sort. Let \(a\) and \(b\) be elements of \(\alpha\). Assume \(a = b\). Then \(b = a\).
\item[\paperlineanchor{premise:45717569762e6170706c795f73796d6d5f6170706c79}\paperlink{premise:45717569762e6170706c795f73796d6d5f6170706c79}{P6.}] \refstepcounter{paperpremise}\label{premise:45717569762e6170706c795f73796d6d5f6170706c79}\textbf{Apply symm apply.} Let \(\alpha\) be a sort. Let \(\beta\) be a sort. Let \(e\) be a bijection from \(\alpha\) to \(\beta\). Let \(x\) be an element of \(\beta\). Then \(e\left(\left(e^{-1}\right)\left(x\right)\right) = x\).
\item[\paperlineanchor{premise:45717569762e73796d6d5f6170706c795f6170706c79}\paperlink{premise:45717569762e73796d6d5f6170706c795f6170706c79}{P7.}] \refstepcounter{paperpremise}\label{premise:45717569762e73796d6d5f6170706c795f6170706c79}\textbf{Symm apply apply.} Let \(\alpha\) be a sort. Let \(\beta\) be a sort. Let \(e\) be a bijection from \(\alpha\) to \(\beta\). Let \(x\) be an element of \(\alpha\). Then \(\left(e^{-1}\right)\left(e\left(x\right)\right) = x\).
\item[\paperlineanchor{premise:46696e7365742e4e6f6e656d7074792e6578697374735f65715f73696e676c65746f6e5f6f725f6e6f6e7472697669616c}\paperlink{premise:46696e7365742e4e6f6e656d7074792e6578697374735f65715f73696e676c65746f6e5f6f725f6e6f6e7472697669616c}{P8.}] \refstepcounter{paperpremise}\label{premise:46696e7365742e4e6f6e656d7074792e6578697374735f65715f73696e676c65746f6e5f6f725f6e6f6e7472697669616c}\textbf{Exists equals singleton or nontrivial.} Let \(\alpha\) be a type. Let \(s\) be a finite subset of \(\alpha\). Assume \(s\ne\varnothing\). Then \(\left(\exists a_{1}\, :\, \alpha,\; s = \{a_{1}\}\right) \lor s\text{ contains two distinct elements}\).
\item[\paperlineanchor{premise:46696e7365742e5375627365742e7265666c}\paperlink{premise:46696e7365742e5375627365742e7265666c}{P9.}] \refstepcounter{paperpremise}\label{premise:46696e7365742e5375627365742e7265666c}\textbf{Refl.} Let \(\alpha\) be a type. Let \(s\) be a finite subset of \(\alpha\). Then \(s \subseteq s\).
\item[\paperlineanchor{premise:46696e7365742e636f655f617474616368}\paperlink{premise:46696e7365742e636f655f617474616368}{P10.}] \refstepcounter{paperpremise}\label{premise:46696e7365742e636f655f617474616368}\textbf{Coercion attach.} Let \(\alpha\) be a type. Let \(s\) be a finite subset of \(\alpha\). Then \(\operatorname{underlying\,set}(\operatorname{Finset\mathord{:}attach}\left(s\right)) = s\).
\item[\paperlineanchor{premise:46696e7365742e636f655f696d616765}\paperlink{premise:46696e7365742e636f655f696d616765}{P11.}] \refstepcounter{paperpremise}\label{premise:46696e7365742e636f655f696d616765}\textbf{Coercion image.} Let \(\alpha\) and \(\beta\) be types. Assume equality on \(\beta\) is decidable. Let \(f\) be a function of type \(\alpha \to \beta\). Let \(s\) be a finite subset of \(\alpha\). Then \(\operatorname{underlying\,set}(f\left[s\right]) = f\left[\operatorname{underlying\,set}(s)\right]\).
\item[\paperlineanchor{premise:46696e7365742e636f655f696e6a656374697665}\paperlink{premise:46696e7365742e636f655f696e6a656374697665}{P12.}] \refstepcounter{paperpremise}\label{premise:46696e7365742e636f655f696e6a656374697665}\textbf{Coercion injective.} Let \(\alpha\) be a type. Then \(\operatorname{underlying\,set}\text{ is injective}\).
\item[\paperlineanchor{premise:46696e7365742e636f655f6d656d}\paperlink{premise:46696e7365742e636f655f6d656d}{P13.}] \refstepcounter{paperpremise}\label{premise:46696e7365742e636f655f6d656d}\textbf{Coercion membership.} Let \(\alpha\) be a type. Let \(s\) be a finite subset of \(\alpha\). Let \(x\) be an element of \(s\). Then \(x \in s\).
\item[\paperlineanchor{premise:46696e7365742e636f655f72616e6765}\paperlink{premise:46696e7365742e636f655f72616e6765}{P14.}] \refstepcounter{paperpremise}\label{premise:46696e7365742e636f655f72616e6765}\textbf{Coercion range.} Let \(n\) be a natural number. Then \(\operatorname{underlying\,set}(\{i\in\mathbb{N}\mid i<n\}) = \operatorname{Iio}\left(n\right)\).
\item[\paperlineanchor{premise:46696e7365742e636f655f736f72745f636f65}\paperlink{premise:46696e7365742e636f655f736f72745f636f65}{P15.}] \refstepcounter{paperpremise}\label{premise:46696e7365742e636f655f736f72745f636f65}\textbf{Coercion sort coercion.} Let \(\alpha\) be a type. Let \(s\) be a finite subset of \(\alpha\). Then \(s = s\).
\item[\paperlineanchor{premise:46696e7365742e636f655f737562736574}\paperlink{premise:46696e7365742e636f655f737562736574}{P16.}] \refstepcounter{paperpremise}\label{premise:46696e7365742e636f655f737562736574}\textbf{Coercion subset.} Let \(\alpha\) be a type. Let \(s_{1}\) and \(s_{2}\) be finite subsets of \(\alpha\). Then \(\operatorname{underlying\,set}(s_{1}) \subseteq \operatorname{underlying\,set}(s_{2}) \iff s_{1} \subseteq s_{2}\).
\item[\paperlineanchor{premise:46696e7365742e636f6e735f65715f696e73657274}\paperlink{premise:46696e7365742e636f6e735f65715f696e73657274}{P17.}] \refstepcounter{paperpremise}\label{premise:46696e7365742e636f6e735f65715f696e73657274}\textbf{Cons equals insert.} Let \(\alpha\) be a type. Assume equality on \(\alpha\) is decidable. Let \(a\) be an element of \(\alpha\). Let \(s\) be a finite subset of \(\alpha\). Let \(h\) be a proof of \(a \notin s\). Then \(\operatorname{Finset\mathord{:}cons}\left(a, s\right) = \{a\} \cup s\).
\item[\paperlineanchor{premise:46696e7365742e636f6e735f696e64756374696f6e5f6f6e}\paperlink{premise:46696e7365742e636f6e735f696e64756374696f6e5f6f6e}{P18.}] \refstepcounter{paperpremise}\label{premise:46696e7365742e636f6e735f696e64756374696f6e5f6f6e}\textbf{Cons induction on.} Let \(\alpha\) be a type. Let \(\mathit{motive}\) be a function of type \(\mathcal{P}_{\mathrm{fin}}(\alpha) \to \mathsf{Prop}\). Let \(s\) be a finite subset of \(\alpha\). Assume \(\mathit{motive}\left(\varnothing\right)\). Assume \(\forall a\, :\, \alpha,\; \forall s_{1}\, :\, \mathcal{P}_{\mathrm{fin}}(\alpha),\; \forall h\, :\, a \notin s_{1},\; \mathit{motive}(s_{1}) \implies \mathit{motive}(\operatorname{Finset\mathord{:}cons}(a,\allowbreak{} s_{1}))\). Then \(\mathit{motive}\left(s\right)\).
\item[\paperlineanchor{premise:46696e7365742e636f6e735f6e655f656d707479}\paperlink{premise:46696e7365742e636f6e735f6e655f656d707479}{P19.}] \refstepcounter{paperpremise}\label{premise:46696e7365742e636f6e735f6e655f656d707479}\textbf{Cons differs from empty.} Let \(\alpha\) be a type. Let \(s\) be a finite subset of \(\alpha\). Let \(a\) be an element of \(\alpha\). Let \(h\) be a proof of \(a \notin s\). Then \(\operatorname{Finset\mathord{:}cons}\left(a, s\right) \ne \varnothing\).
\item[\paperlineanchor{premise:46696e7365742e6469736a6f696e745f66696c7465725f66696c7465725f6e6567}\paperlink{premise:46696e7365742e6469736a6f696e745f66696c7465725f66696c7465725f6e6567}{P20.}] \refstepcounter{paperpremise}\label{premise:46696e7365742e6469736a6f696e745f66696c7465725f66696c7465725f6e6567}\textbf{Disjoint filter filter neg.} Let \(\alpha\) be a type. Let \(s\) and \(t\) be finite subsets of \(\alpha\). Let \(p\) be a function of type \(\alpha \to \mathsf{Prop}\). Assume the predicate \(p\) is decidable. Fix a structure \(\mathit{inst}_{1}\) of type \(\Pi_{x\, :\, \alpha}\, \operatorname{Decidable}\left(\neg p\left(x\right)\right)\). Then \(\operatorname{Disjoint}\left(\{x\in s \mid p(x)\}, \{a\in t \mid \neg p\left(a\right)\}\right)\).
\item[\paperlineanchor{premise:46696e7365742e6469736a6f696e745f6966665f696e7465725f65715f656d707479}\paperlink{premise:46696e7365742e6469736a6f696e745f6966665f696e7465725f65715f656d707479}{P21.}] \refstepcounter{paperpremise}\label{premise:46696e7365742e6469736a6f696e745f6966665f696e7465725f65715f656d707479}\textbf{Disjoint if and only if inter equals empty.} Let \(\alpha\) be a type. Assume equality on \(\alpha\) is decidable. Let \(s\) and \(t\) be finite subsets of \(\alpha\). Then \(\operatorname{Disjoint}\left(s, t\right) \iff s \cap t = \varnothing\).
\item[\paperlineanchor{premise:46696e7365742e6469736a6f696e745f73646966665f696e746572}\paperlink{premise:46696e7365742e6469736a6f696e745f73646966665f696e746572}{P22.}] \refstepcounter{paperpremise}\label{premise:46696e7365742e6469736a6f696e745f73646966665f696e746572}\textbf{Disjoint sdiff inter.} Let \(\alpha\) be a type. Assume equality on \(\alpha\) is decidable. Let \(s\) and \(t\) be finite subsets of \(\alpha\). Then \(\operatorname{Disjoint}\left(s \setminus t, s \cap t\right)\).
\item[\paperlineanchor{premise:46696e7365742e656d7074795f737562736574}\paperlink{premise:46696e7365742e656d7074795f737562736574}{P23.}] \refstepcounter{paperpremise}\label{premise:46696e7365742e656d7074795f737562736574}\textbf{Empty subset.} Let \(\alpha\) be a type. Let \(s\) be a finite subset of \(\alpha\). Then \(\varnothing \subseteq s\).
\item[\paperlineanchor{premise:46696e7365742e65715f656d7074795f6f725f6e6f6e656d707479}\paperlink{premise:46696e7365742e65715f656d7074795f6f725f6e6f6e656d707479}{P24.}] \refstepcounter{paperpremise}\label{premise:46696e7365742e65715f656d7074795f6f725f6e6f6e656d707479}\textbf{Equals empty or nonempty.} Let \(\alpha\) be a type. Let \(s\) be a finite subset of \(\alpha\). Then \(s = \varnothing \lor s\ne\varnothing\).
\item[\paperlineanchor{premise:46696e7365742e65715f6f665f766571}\paperlink{premise:46696e7365742e65715f6f665f766571}{P25.}] \refstepcounter{paperpremise}\label{premise:46696e7365742e65715f6f665f766571}\textbf{Equals of veq.} Let \(\alpha\) be a type. Let \(s\) and \(t\) be finite subsets of \(\alpha\). Assume \(\operatorname{underlying\,multiset}(s) = \operatorname{underlying\,multiset}(t)\). Then \(s = t\).
\item[\paperlineanchor{premise:46696e7365742e65715f73696e676c65746f6e5f6966665f756e697175655f6d656d}\paperlink{premise:46696e7365742e65715f73696e676c65746f6e5f6966665f756e697175655f6d656d}{P26.}] \refstepcounter{paperpremise}\label{premise:46696e7365742e65715f73696e676c65746f6e5f6966665f756e697175655f6d656d}\textbf{Equals singleton if and only if unique membership.} Let \(\alpha\) be a type. Let \(s\) be a finite subset of \(\alpha\). Let \(a\) be an element of \(\alpha\). Then \(s = \{a\} \iff a \in s \land \left(\forall x\, :\, \alpha,\; x \in s \implies x = a\right)\).
\item[\paperlineanchor{premise:46696e7365742e65726173655f65715f6f665f6e6f745f6d656d}\paperlink{premise:46696e7365742e65726173655f65715f6f665f6e6f745f6d656d}{P27.}] \refstepcounter{paperpremise}\label{premise:46696e7365742e65726173655f65715f6f665f6e6f745f6d656d}\textbf{Erase equals of not membership.} Let \(\alpha\) be a type. Assume equality on \(\alpha\) is decidable. Let \(a\) be an element of \(\alpha\). Let \(s\) be a finite subset of \(\alpha\). Assume \(a \notin s\). Then \(s \setminus \{a\} = s\).
\item[\paperlineanchor{premise:46696e7365742e65726173655f7375627365745f6966665f6f665f6d656d}\paperlink{premise:46696e7365742e65726173655f7375627365745f6966665f6f665f6d656d}{P28.}] \refstepcounter{paperpremise}\label{premise:46696e7365742e65726173655f7375627365745f6966665f6f665f6d656d}\textbf{Erase subset if and only if of membership.} Let \(\alpha\) be a type. Assume equality on \(\alpha\) is decidable. Let \(s\) and \(t\) be finite subsets of \(\alpha\). Let \(a\) be an element of \(\alpha\). Assume \(a \in t\). Then \(s \setminus \{a\} \subseteq t \iff s \subseteq t\).
\item[\paperlineanchor{premise:46696e7365742e657874}\paperlink{premise:46696e7365742e657874}{P29.}] \refstepcounter{paperpremise}\label{premise:46696e7365742e657874}\textbf{Ext.} Let \(\alpha\) be a type. Let \(s_{1}\) and \(s_{2}\) be finite subsets of \(\alpha\). Assume \(\forall a\, :\, \alpha,\; a \in s_{1} \iff a \in s_{2}\). Then \(s_{1} = s_{2}\).
\item[\paperlineanchor{premise:46696e7365742e66696265725f6e6f6e656d7074795f6966665f6d656d5f696d616765}\paperlink{premise:46696e7365742e66696265725f6e6f6e656d7074795f6966665f6d656d5f696d616765}{P30.}] \refstepcounter{paperpremise}\label{premise:46696e7365742e66696265725f6e6f6e656d7074795f6966665f6d656d5f696d616765}\textbf{Fiber nonempty if and only if membership image.} Let \(\alpha\) and \(\beta\) be types. Assume equality on \(\beta\) is decidable. Let \(f\) be a function of type \(\alpha \to \beta\). Let \(s\) be a finite subset of \(\alpha\). Let \(y\) be an element of \(\beta\). Then \(\{x\in s \mid f\left(x\right) = y\}\ne\varnothing \iff y \in f\left[s\right]\).
\item[\paperlineanchor{premise:46696e7365742e66696c7465725f65715f656d7074795f696666}\paperlink{premise:46696e7365742e66696c7465725f65715f656d7074795f696666}{P31.}] \refstepcounter{paperpremise}\label{premise:46696e7365742e66696c7465725f65715f656d7074795f696666}\textbf{Filter equals empty if and only if.} Let \(\alpha\) be a type. Let \(p\) be a function of type \(\alpha \to \mathsf{Prop}\). Assume the predicate \(p\) is decidable. Let \(s\) be a finite subset of \(\alpha\). Then \(\{x\in s \mid p(x)\} = \varnothing \iff (\forall x\, :\, \alpha,\; x \in s \implies \neg p(x))\).
\item[\paperlineanchor{premise:46696e7365742e66696c7465725f65715f73656c66}\paperlink{premise:46696e7365742e66696c7465725f65715f73656c66}{P32.}] \refstepcounter{paperpremise}\label{premise:46696e7365742e66696c7465725f65715f73656c66}\textbf{Filter equals self.} Let \(\alpha\) be a type. Let \(p\) be a function of type \(\alpha \to \mathsf{Prop}\). Assume the predicate \(p\) is decidable. Let \(s\) be a finite subset of \(\alpha\). Then \(\{x\in s \mid p(x)\} = s \iff \left(\forall x\, :\, \alpha,\; x \in s \implies p\left(x\right)\right)\).
\item[\paperlineanchor{premise:46696e7365742e66696c7465725f73737562736574}\paperlink{premise:46696e7365742e66696c7465725f73737562736574}{P33.}] \refstepcounter{paperpremise}\label{premise:46696e7365742e66696c7465725f73737562736574}\textbf{Filter ssubset.} Let \(\alpha\) be a type. Let \(p\) be a function of type \(\alpha \to \mathsf{Prop}\). Assume the predicate \(p\) is decidable. Let \(s\) be a finite subset of \(\alpha\). Then \(\{x\in s \mid p(x)\} \subsetneq s \iff (\exists x\, :\, \alpha,\; x \in s \land \neg p(x))\).
\item[\paperlineanchor{premise:46696e7365742e66696c7465725f737562736574}\paperlink{premise:46696e7365742e66696c7465725f737562736574}{P34.}] \refstepcounter{paperpremise}\label{premise:46696e7365742e66696c7465725f737562736574}\textbf{Filter subset.} Let \(\alpha\) be a type. Let \(p\) be a function of type \(\alpha \to \mathsf{Prop}\). Assume the predicate \(p\) is decidable. Let \(s\) be a finite subset of \(\alpha\). Then \(\{x\in s \mid p(x)\} \subseteq s\).
\item[\paperlineanchor{premise:46696e7365742e66696c7465725f756e696f6e5f66696c7465725f6e65675f6571}\paperlink{premise:46696e7365742e66696c7465725f756e696f6e5f66696c7465725f6e65675f6571}{P35.}] \refstepcounter{paperpremise}\label{premise:46696e7365742e66696c7465725f756e696f6e5f66696c7465725f6e65675f6571}\textbf{Filter union filter neg equals.} Let \(\alpha\) be a type. Let \(p\) be a function of type \(\alpha \to \mathsf{Prop}\). Assume the predicate \(p\) is decidable. Assume equality on \(\alpha\) is decidable. Fix a structure \(\mathit{inst}_{2}\) of type \(\Pi_{x\, :\, \alpha}\, \operatorname{Decidable}\left(\neg p\left(x\right)\right)\). Let \(s\) be a finite subset of \(\alpha\). Then \(\{x\in s \mid p(x)\} \cup \{a\in s \mid \neg p\left(a\right)\} = s\).
\item[\paperlineanchor{premise:46696e7365742e696d6167655f7375627365745f696666}\paperlink{premise:46696e7365742e696d6167655f7375627365745f696666}{P36.}] \refstepcounter{paperpremise}\label{premise:46696e7365742e696d6167655f7375627365745f696666}\textbf{Image subset if and only if.} Let \(\alpha\) and \(\beta\) be types. Assume equality on \(\beta\) is decidable. Let \(f\) be a function of type \(\alpha \to \beta\). Let \(s\) be a finite subset of \(\alpha\). Let \(t\) be a finite subset of \(\beta\). Then \(f\left[s\right] \subseteq t \iff \left(\forall x\, :\, \alpha,\; x \in s \implies f\left(x\right) \in t\right)\).
\item[\paperlineanchor{premise:46696e7365742e696d6167655f76616c5f6f665f696e6a4f6e}\paperlink{premise:46696e7365742e696d6167655f76616c5f6f665f696e6a4f6e}{P37.}] \refstepcounter{paperpremise}\label{premise:46696e7365742e696d6167655f76616c5f6f665f696e6a4f6e}\textbf{Image val of injectivity on.} Let \(\alpha\) and \(\beta\) be types. Assume equality on \(\beta\) is decidable. Let \(f\) be a function of type \(\alpha \to \beta\). Let \(s\) be a finite subset of \(\alpha\). Assume \(f\text{ is injective on }\operatorname{underlying\,set}(s)\). Then \(\operatorname{underlying\,multiset}(f[s]) = \operatorname{Multiset\mathord{:}map}(f,\allowbreak{} \operatorname{underlying\,multiset}(s))\).
\item[\paperlineanchor{premise:46696e7365742e696e64756374696f6e5f6f6e}\paperlink{premise:46696e7365742e696e64756374696f6e5f6f6e}{P38.}] \refstepcounter{paperpremise}\label{premise:46696e7365742e696e64756374696f6e5f6f6e}\textbf{Induction on.} Let \(\alpha\) be a type. Let \(\mathit{motive}\) be a function of type \(\mathcal{P}_{\mathrm{fin}}(\alpha) \to \mathsf{Prop}\). Assume equality on \(\alpha\) is decidable. Let \(s\) be a finite subset of \(\alpha\). Assume \(\mathit{motive}\left(\varnothing\right)\). Assume \(\forall a\, :\, \alpha,\; \forall s_{1}\, :\, \mathcal{P}_{\mathrm{fin}}(\alpha),\; a \notin s_{1} \implies \mathit{motive}(s_{1}) \implies \mathit{motive}(\{a\} \cup s_{1})\). Then \(\mathit{motive}\left(s\right)\).
\item[\paperlineanchor{premise:46696e7365742e696e736572745f6571}\paperlink{premise:46696e7365742e696e736572745f6571}{P39.}] \refstepcounter{paperpremise}\label{premise:46696e7365742e696e736572745f6571}\textbf{Insert equals.} Let \(\alpha\) be a type. Assume equality on \(\alpha\) is decidable. Let \(a\) be an element of \(\alpha\). Let \(s\) be a finite subset of \(\alpha\). Then \(\{a\} \cup s = \{a\} \cup s\).
\item[\paperlineanchor{premise:46696e7365742e696e736572745f65715f6f665f6d656d}\paperlink{premise:46696e7365742e696e736572745f65715f6f665f6d656d}{P40.}] \refstepcounter{paperpremise}\label{premise:46696e7365742e696e736572745f65715f6f665f6d656d}\textbf{Insert equals of membership.} Let \(\alpha\) be a type. Assume equality on \(\alpha\) is decidable. Let \(s\) be a finite subset of \(\alpha\). Let \(a\) be an element of \(\alpha\). Assume \(a \in s\). Then \(\{a\} \cup s = s\).
\item[\paperlineanchor{premise:46696e7365742e696e736572745f6572617365}\paperlink{premise:46696e7365742e696e736572745f6572617365}{P41.}] \refstepcounter{paperpremise}\label{premise:46696e7365742e696e736572745f6572617365}\textbf{Insert erase.} Let \(\alpha\) be a type. Assume equality on \(\alpha\) is decidable. Let \(s\) be a finite subset of \(\alpha\). Let \(a\) be an element of \(\alpha\). Assume \(a \in s\). Then \(\{a\} \cup \left(s \setminus \{a\}\right) = s\).
\item[\paperlineanchor{premise:46696e7365742e696e736572745f7375627365745f696666}\paperlink{premise:46696e7365742e696e736572745f7375627365745f696666}{P42.}] \refstepcounter{paperpremise}\label{premise:46696e7365742e696e736572745f7375627365745f696666}\textbf{Insert subset if and only if.} Let \(\alpha\) be a type. Assume equality on \(\alpha\) is decidable. Let \(s\) and \(t\) be finite subsets of \(\alpha\). Let \(a\) be an element of \(\alpha\). Then \(\{a\} \cup s \subseteq t \iff a \in t \land s \subseteq t\).
\item[\paperlineanchor{premise:46696e7365742e696e737449735265666c537562736574}\paperlink{premise:46696e7365742e696e737449735265666c537562736574}{P43.}] \refstepcounter{paperpremise}\label{premise:46696e7365742e696e737449735265666c537562736574}\textbf{Inst is refl subset.} Let \(\alpha\) be a type. Then \((\mathit{x1} \mapsto \mathit{x2} \mapsto \mathit{x1} \subseteq \mathit{x2})\text{ is reflexive on }\mathcal{P}_{\mathrm{fin}}(\alpha)\).
\item[\paperlineanchor{premise:46696e7365742e696e7465725f656d707479}\paperlink{premise:46696e7365742e696e7465725f656d707479}{P44.}] \refstepcounter{paperpremise}\label{premise:46696e7365742e696e7465725f656d707479}\textbf{Inter empty.} Let \(\alpha\) be a type. Assume equality on \(\alpha\) is decidable. Let \(s\) be a finite subset of \(\alpha\). Then \(s \cap \varnothing = \varnothing\).
\item[\paperlineanchor{premise:46696e7365742e696e7465725f696e736572745f6f665f6d656d}\paperlink{premise:46696e7365742e696e7465725f696e736572745f6f665f6d656d}{P45.}] \refstepcounter{paperpremise}\label{premise:46696e7365742e696e7465725f696e736572745f6f665f6d656d}\textbf{Inter insert of membership.} Let \(\alpha\) be a type. Assume equality on \(\alpha\) is decidable. Let \(s_{1}\) and \(s_{2}\) be finite subsets of \(\alpha\). Let \(a\) be an element of \(\alpha\). Assume \(a \in s_{1}\). Then \(s_{1} \cap \left(\{a\} \cup s_{2}\right) = \{a\} \cup \left(s_{1} \cap s_{2}\right)\).
\item[\paperlineanchor{premise:46696e7365742e696e7465725f696e736572745f6f665f6e6f745f6d656d}\paperlink{premise:46696e7365742e696e7465725f696e736572745f6f665f6e6f745f6d656d}{P46.}] \refstepcounter{paperpremise}\label{premise:46696e7365742e696e7465725f696e736572745f6f665f6e6f745f6d656d}\textbf{Inter insert of not membership.} Let \(\alpha\) be a type. Assume equality on \(\alpha\) is decidable. Let \(s_{1}\) and \(s_{2}\) be finite subsets of \(\alpha\). Let \(a\) be an element of \(\alpha\). Assume \(a \notin s_{1}\). Then \(s_{1} \cap \left(\{a\} \cup s_{2}\right) = s_{1} \cap s_{2}\).
\item[\paperlineanchor{premise:46696e7365742e696e7465725f73646966665f73656c66}\paperlink{premise:46696e7365742e696e7465725f73646966665f73656c66}{P47.}] \refstepcounter{paperpremise}\label{premise:46696e7365742e696e7465725f73646966665f73656c66}\textbf{Inter sdiff self.} Let \(\alpha\) be a type. Assume equality on \(\alpha\) is decidable. Let \(s_{1}\) and \(s_{2}\) be finite subsets of \(\alpha\). Then \(s_{1} \cap \left(s_{2} \setminus s_{1}\right) = \varnothing\).
\item[\paperlineanchor{premise:46696e7365742e696e7465725f7375627365745f6c656674}\paperlink{premise:46696e7365742e696e7465725f7375627365745f6c656674}{P48.}] \refstepcounter{paperpremise}\label{premise:46696e7365742e696e7465725f7375627365745f6c656674}\textbf{Inter subset left.} Let \(\alpha\) be a type. Assume equality on \(\alpha\) is decidable. Let \(s_{1}\) and \(s_{2}\) be finite subsets of \(\alpha\). Then \(s_{1} \cap s_{2} \subseteq s_{1}\).
\item[\paperlineanchor{premise:46696e7365742e696e7465725f7375627365745f7269676874}\paperlink{premise:46696e7365742e696e7465725f7375627365745f7269676874}{P49.}] \refstepcounter{paperpremise}\label{premise:46696e7365742e696e7465725f7375627365745f7269676874}\textbf{Inter subset right.} Let \(\alpha\) be a type. Assume equality on \(\alpha\) is decidable. Let \(s_{1}\) and \(s_{2}\) be finite subsets of \(\alpha\). Then \(s_{1} \cap s_{2} \subseteq s_{2}\).
\item[\paperlineanchor{premise:46696e7365742e6d656d5f617474616368}\paperlink{premise:46696e7365742e6d656d5f617474616368}{P50.}] \refstepcounter{paperpremise}\label{premise:46696e7365742e6d656d5f617474616368}\textbf{Membership attach.} Let \(\alpha\) be a type. Let \(s\) be a finite subset of \(\alpha\). Let \(x\) be an element of \(s\). Then \(x \in \operatorname{Finset\mathord{:}attach}\left(s\right)\).
\item[\paperlineanchor{premise:46696e7365742e6d656d5f636f65}\paperlink{premise:46696e7365742e6d656d5f636f65}{P51.}] \refstepcounter{paperpremise}\label{premise:46696e7365742e6d656d5f636f65}\textbf{Membership coercion.} Let \(\alpha\) be a type. Let \(a\) be an element of \(\alpha\). Let \(s\) be a finite subset of \(\alpha\). Then \(a \in \operatorname{underlying\,set}(s) \iff a \in s\).
\item[\paperlineanchor{premise:46696e7365742e6d656d5f696d616765}\paperlink{premise:46696e7365742e6d656d5f696d616765}{P52.}] \refstepcounter{paperpremise}\label{premise:46696e7365742e6d656d5f696d616765}\textbf{Membership image.} Let \(\alpha\) and \(\beta\) be types. Assume equality on \(\beta\) is decidable. Let \(f\) be a function of type \(\alpha \to \beta\). Let \(s\) be a finite subset of \(\alpha\). Let \(b\) be an element of \(\beta\). Then \(b \in f\left[s\right] \iff \left(\exists a\, :\, \alpha,\; a \in s \land f\left(a\right) = b\right)\).
\item[\paperlineanchor{premise:46696e7365742e6d656d5f696d6167655f6f665f6d656d}\paperlink{premise:46696e7365742e6d656d5f696d6167655f6f665f6d656d}{P53.}] \refstepcounter{paperpremise}\label{premise:46696e7365742e6d656d5f696d6167655f6f665f6d656d}\textbf{Membership image of membership.} Let \(\alpha\) and \(\beta\) be types. Assume equality on \(\beta\) is decidable. Let \(s\) be a finite subset of \(\alpha\). Let \(f\) be a function of type \(\alpha \to \beta\). Let \(a\) be an element of \(\alpha\). Assume \(a \in s\). Then \(f\left(a\right) \in f\left[s\right]\).
\item[\paperlineanchor{premise:46696e7365742e6d656d5f696e73657274}\paperlink{premise:46696e7365742e6d656d5f696e73657274}{P54.}] \refstepcounter{paperpremise}\label{premise:46696e7365742e6d656d5f696e73657274}\textbf{Membership insert.} Let \(\alpha\) be a type. Assume equality on \(\alpha\) is decidable. Let \(s\) be a finite subset of \(\alpha\). Let \(a\) and \(b\) be elements of \(\alpha\). Then \(a \in \{b\} \cup s \iff a = b \lor a \in s\).
\item[\paperlineanchor{premise:46696e7365742e6d656d5f696e746572}\paperlink{premise:46696e7365742e6d656d5f696e746572}{P55.}] \refstepcounter{paperpremise}\label{premise:46696e7365742e6d656d5f696e746572}\textbf{Membership inter.} Let \(\alpha\) be a type. Assume equality on \(\alpha\) is decidable. Let \(a\) be an element of \(\alpha\). Let \(s_{1}\) and \(s_{2}\) be finite subsets of \(\alpha\). Then \(a \in s_{1} \cap s_{2} \iff a \in s_{1} \land a \in s_{2}\).
\item[\paperlineanchor{premise:46696e7365742e6d656d5f72616e6765}\paperlink{premise:46696e7365742e6d656d5f72616e6765}{P56.}] \refstepcounter{paperpremise}\label{premise:46696e7365742e6d656d5f72616e6765}\textbf{Membership range.} Let \(n\) and \(m\) be natural numbers. Then \(m \in \{i\in\mathbb{N}\mid i<n\} \iff m < n\).
\item[\paperlineanchor{premise:46696e7365742e6d656d5f7364696666}\paperlink{premise:46696e7365742e6d656d5f7364696666}{P57.}] \refstepcounter{paperpremise}\label{premise:46696e7365742e6d656d5f7364696666}\textbf{Membership sdiff.} Let \(\alpha\) be a type. Assume equality on \(\alpha\) is decidable. Let \(s\) and \(t\) be finite subsets of \(\alpha\). Let \(a\) be an element of \(\alpha\). Then \(a \in s \setminus t \iff a \in s \land a \notin t\).
\item[\paperlineanchor{premise:46696e7365742e6d656d5f73696e676c65746f6e}\paperlink{premise:46696e7365742e6d656d5f73696e676c65746f6e}{P58.}] \refstepcounter{paperpremise}\label{premise:46696e7365742e6d656d5f73696e676c65746f6e}\textbf{Membership singleton.} Let \(\alpha\) be a type. Let \(a\) and \(b\) be elements of \(\alpha\). Then \(b \in \{a\} \iff b = a\).
\item[\paperlineanchor{premise:46696e7365742e6d656d5f73696e676c65746f6e5f73656c66}\paperlink{premise:46696e7365742e6d656d5f73696e676c65746f6e5f73656c66}{P59.}] \refstepcounter{paperpremise}\label{premise:46696e7365742e6d656d5f73696e676c65746f6e5f73656c66}\textbf{Membership singleton self.} Let \(\alpha\) be a type. Let \(a\) be an element of \(\alpha\). Then \(a \in \{a\}\).
\item[\paperlineanchor{premise:46696e7365742e6d656d5f756e696f6e}\paperlink{premise:46696e7365742e6d656d5f756e696f6e}{P60.}] \refstepcounter{paperpremise}\label{premise:46696e7365742e6d656d5f756e696f6e}\textbf{Membership union.} Let \(\alpha\) be a type. Assume equality on \(\alpha\) is decidable. Let \(s\) and \(t\) be finite subsets of \(\alpha\). Let \(a\) be an element of \(\alpha\). Then \(a \in s \cup t \iff a \in s \lor a \in t\).
\item[\paperlineanchor{premise:46696e7365742e6e655f656d7074795f6f665f6d656d}\paperlink{premise:46696e7365742e6e655f656d7074795f6f665f6d656d}{P61.}] \refstepcounter{paperpremise}\label{premise:46696e7365742e6e655f656d7074795f6f665f6d656d}\textbf{Differs from empty of membership.} Let \(\alpha\) be a type. Let \(a\) be an element of \(\alpha\). Let \(s\) be a finite subset of \(\alpha\). Assume \(a \in s\). Then \(s \ne \varnothing\).
\item[\paperlineanchor{premise:46696e7365742e6e6f6e656d7074795f6966665f6e655f656d707479}\paperlink{premise:46696e7365742e6e6f6e656d7074795f6966665f6e655f656d707479}{P62.}] \refstepcounter{paperpremise}\label{premise:46696e7365742e6e6f6e656d7074795f6966665f6e655f656d707479}\textbf{Nonempty if and only if differs from empty.} Let \(\alpha\) be a type. Let \(s\) be a finite subset of \(\alpha\). Then \(s\ne\varnothing \iff s \ne \varnothing\).
\item[\paperlineanchor{premise:46696e7365742e6e6f745f6d656d5f656d707479}\paperlink{premise:46696e7365742e6e6f745f6d656d5f656d707479}{P63.}] \refstepcounter{paperpremise}\label{premise:46696e7365742e6e6f745f6d656d5f656d707479}\textbf{Not membership empty.} Let \(\alpha\) be a type. Let \(a\) be an element of \(\alpha\). Then \(a \notin \varnothing\).
\item[\paperlineanchor{premise:46696e7365742e6e6f745f6d656d5f6572617365}\paperlink{premise:46696e7365742e6e6f745f6d656d5f6572617365}{P64.}] \refstepcounter{paperpremise}\label{premise:46696e7365742e6e6f745f6d656d5f6572617365}\textbf{Not membership erase.} Let \(\alpha\) be a type. Assume equality on \(\alpha\) is decidable. Let \(a\) be an element of \(\alpha\). Let \(s\) be a finite subset of \(\alpha\). Then \(a \notin s \setminus \{a\}\).
\item[\paperlineanchor{premise:46696e7365742e6e6f745f6d656d5f73646966665f6f665f6d656d5f7269676874}\paperlink{premise:46696e7365742e6e6f745f6d656d5f73646966665f6f665f6d656d5f7269676874}{P65.}] \refstepcounter{paperpremise}\label{premise:46696e7365742e6e6f745f6d656d5f73646966665f6f665f6d656d5f7269676874}\textbf{Not membership sdiff of membership right.} Let \(\alpha\) be a type. Assume equality on \(\alpha\) is decidable. Let \(s\) and \(t\) be finite subsets of \(\alpha\). Let \(a\) be an element of \(\alpha\). Assume \(a \in t\). Then \(a \notin s \setminus t\).
\item[\paperlineanchor{premise:46696e7365742e6e6f745f6d656d5f73696e676c65746f6e}\paperlink{premise:46696e7365742e6e6f745f6d656d5f73696e676c65746f6e}{P66.}] \refstepcounter{paperpremise}\label{premise:46696e7365742e6e6f745f6d656d5f73696e676c65746f6e}\textbf{Not membership singleton.} Let \(\alpha\) be a type. Let \(a\) and \(b\) be elements of \(\alpha\). Then \(a \notin \{b\} \iff a \ne b\).
\item[\paperlineanchor{premise:46696e7365742e6e6f745f6e6f6e656d7074795f6966665f65715f656d707479}\paperlink{premise:46696e7365742e6e6f745f6e6f6e656d7074795f6966665f65715f656d707479}{P67.}] \refstepcounter{paperpremise}\label{premise:46696e7365742e6e6f745f6e6f6e656d7074795f6966665f65715f656d707479}\textbf{Not nonempty if and only if equals empty.} Let \(\alpha\) be a type. Let \(s\) be a finite subset of \(\alpha\). Then \(\neg \left(s\ne\varnothing\right) \iff s = \varnothing\).
\item[\paperlineanchor{premise:46696e7365742e73646966665f6469736a6f696e74}\paperlink{premise:46696e7365742e73646966665f6469736a6f696e74}{P68.}] \refstepcounter{paperpremise}\label{premise:46696e7365742e73646966665f6469736a6f696e74}\textbf{Sdiff disjoint.} Let \(\alpha\) be a type. Assume equality on \(\alpha\) is decidable. Let \(s\) and \(t\) be finite subsets of \(\alpha\). Then \(\operatorname{Disjoint}\left(t \setminus s, s\right)\).
\item[\paperlineanchor{premise:46696e7365742e73646966665f65715f656d7074795f6966665f737562736574}\paperlink{premise:46696e7365742e73646966665f65715f656d7074795f6966665f737562736574}{P69.}] \refstepcounter{paperpremise}\label{premise:46696e7365742e73646966665f65715f656d7074795f6966665f737562736574}\textbf{Sdiff equals empty if and only if subset.} Let \(\alpha\) be a type. Assume equality on \(\alpha\) is decidable. Let \(s\) and \(t\) be finite subsets of \(\alpha\). Then \(s \setminus t = \varnothing \iff s \subseteq t\).
\item[\paperlineanchor{premise:46696e7365742e73646966665f696e7465725f73656c665f7269676874}\paperlink{premise:46696e7365742e73646966665f696e7465725f73656c665f7269676874}{P70.}] \refstepcounter{paperpremise}\label{premise:46696e7365742e73646966665f696e7465725f73656c665f7269676874}\textbf{Sdiff inter self right.} Let \(\alpha\) be a type. Assume equality on \(\alpha\) is decidable. Let \(s\) and \(t\) be finite subsets of \(\alpha\). Then \(s \setminus \left(t \cap s\right) = s \setminus t\).
\item[\paperlineanchor{premise:46696e7365742e73646966665f737562736574}\paperlink{premise:46696e7365742e73646966665f737562736574}{P71.}] \refstepcounter{paperpremise}\label{premise:46696e7365742e73646966665f737562736574}\textbf{Sdiff subset.} Let \(\alpha\) be a type. Assume equality on \(\alpha\) is decidable. Let \(s\) and \(t\) be finite subsets of \(\alpha\). Then \(s \setminus t \subseteq s\).
\item[\paperlineanchor{premise:46696e7365742e73646966665f756e696f6e5f696e746572}\paperlink{premise:46696e7365742e73646966665f756e696f6e5f696e746572}{P72.}] \refstepcounter{paperpremise}\label{premise:46696e7365742e73646966665f756e696f6e5f696e746572}\textbf{Sdiff union inter.} Let \(\alpha\) be a type. Assume equality on \(\alpha\) is decidable. Let \(s\) and \(t\) be finite subsets of \(\alpha\). Then \(\left(s \setminus t\right) \cup \left(s \cap t\right) = s\).
\item[\paperlineanchor{premise:46696e7365742e73646966665f756e696f6e5f6f665f737562736574}\paperlink{premise:46696e7365742e73646966665f756e696f6e5f6f665f737562736574}{P73.}] \refstepcounter{paperpremise}\label{premise:46696e7365742e73646966665f756e696f6e5f6f665f737562736574}\textbf{Sdiff union of subset.} Let \(\alpha\) be a type. Assume equality on \(\alpha\) is decidable. Let \(s_{1}\) and \(s_{2}\) be finite subsets of \(\alpha\). Assume \(s_{1} \subseteq s_{2}\). Then \(\left(s_{2} \setminus s_{1}\right) \cup s_{1} = s_{2}\).
\item[\paperlineanchor{premise:46696e7365742e73646966665f756e696f6e5f73656c665f65715f756e696f6e}\paperlink{premise:46696e7365742e73646966665f756e696f6e5f73656c665f65715f756e696f6e}{P74.}] \refstepcounter{paperpremise}\label{premise:46696e7365742e73646966665f756e696f6e5f73656c665f65715f756e696f6e}\textbf{Sdiff union self equals union.} Let \(\alpha\) be a type. Assume equality on \(\alpha\) is decidable. Let \(s\) and \(t\) be finite subsets of \(\alpha\). Then \(\left(s \setminus t\right) \cup t = s \cup t\).
\item[\paperlineanchor{premise:46696e7365742e73696e676c65746f6e5f696e7465725f6f665f6d656d}\paperlink{premise:46696e7365742e73696e676c65746f6e5f696e7465725f6f665f6d656d}{P75.}] \refstepcounter{paperpremise}\label{premise:46696e7365742e73696e676c65746f6e5f696e7465725f6f665f6d656d}\textbf{Singleton inter of membership.} Let \(\alpha\) be a type. Assume equality on \(\alpha\) is decidable. Let \(a\) be an element of \(\alpha\). Let \(s\) be a finite subset of \(\alpha\). Assume \(a \in s\). Then \(\{a\} \cap s = \{a\}\).
\item[\paperlineanchor{premise:46696e7365742e73696e676c65746f6e5f696e7465725f6f665f6e6f745f6d656d}\paperlink{premise:46696e7365742e73696e676c65746f6e5f696e7465725f6f665f6e6f745f6d656d}{P76.}] \refstepcounter{paperpremise}\label{premise:46696e7365742e73696e676c65746f6e5f696e7465725f6f665f6e6f745f6d656d}\textbf{Singleton inter of not membership.} Let \(\alpha\) be a type. Assume equality on \(\alpha\) is decidable. Let \(a\) be an element of \(\alpha\). Let \(s\) be a finite subset of \(\alpha\). Assume \(a \notin s\). Then \(\{a\} \cap s = \varnothing\).
\item[\paperlineanchor{premise:46696e7365742e73696e676c65746f6e5f7375627365745f696666}\paperlink{premise:46696e7365742e73696e676c65746f6e5f7375627365745f696666}{P77.}] \refstepcounter{paperpremise}\label{premise:46696e7365742e73696e676c65746f6e5f7375627365745f696666}\textbf{Singleton subset if and only if.} Let \(\alpha\) be a type. Let \(s\) be a finite subset of \(\alpha\). Let \(a\) be an element of \(\alpha\). Then \(\{a\} \subseteq s \iff a \in s\).
\item[\paperlineanchor{premise:46696e7365742e737375627365745f696e73657274}\paperlink{premise:46696e7365742e737375627365745f696e73657274}{P78.}] \refstepcounter{paperpremise}\label{premise:46696e7365742e737375627365745f696e73657274}\textbf{Ssubset insert.} Let \(\alpha\) be a type. Assume equality on \(\alpha\) is decidable. Let \(s\) be a finite subset of \(\alpha\). Let \(a\) be an element of \(\alpha\). Assume \(a \notin s\). Then \(s \subsetneq \{a\} \cup s\).
\item[\paperlineanchor{premise:46696e7365742e7375627365745f6572617365}\paperlink{premise:46696e7365742e7375627365745f6572617365}{P79.}] \refstepcounter{paperpremise}\label{premise:46696e7365742e7375627365745f6572617365}\textbf{Subset erase.} Let \(\alpha\) be a type. Assume equality on \(\alpha\) is decidable. Let \(a\) be an element of \(\alpha\). Let \(s\) and \(t\) be finite subsets of \(\alpha\). Then \(s \subseteq t \setminus \{a\} \iff s \subseteq t \land a \notin s\).
\item[\paperlineanchor{premise:46696e7365742e7375627365745f696666}\paperlink{premise:46696e7365742e7375627365745f696666}{P80.}] \refstepcounter{paperpremise}\label{premise:46696e7365742e7375627365745f696666}\textbf{Subset if and only if.} Let \(\alpha\) be a type. Let \(s_{1}\) and \(s_{2}\) be finite subsets of \(\alpha\). Then \(s_{1} \subseteq s_{2} \iff \left(\forall x\, :\, \alpha,\; x \in s_{1} \implies x \in s_{2}\right)\).
\item[\paperlineanchor{premise:46696e7365742e7375627365745f696e73657274}\paperlink{premise:46696e7365742e7375627365745f696e73657274}{P81.}] \refstepcounter{paperpremise}\label{premise:46696e7365742e7375627365745f696e73657274}\textbf{Subset insert.} Let \(\alpha\) be a type. Assume equality on \(\alpha\) is decidable. Let \(a\) be an element of \(\alpha\). Let \(s\) be a finite subset of \(\alpha\). Then \(s \subseteq \{a\} \cup s\).
\item[\paperlineanchor{premise:46696e7365742e7375627365745f73696e676c65746f6e5f69666627}\paperlink{premise:46696e7365742e7375627365745f73696e676c65746f6e5f69666627}{P82.}] \refstepcounter{paperpremise}\label{premise:46696e7365742e7375627365745f73696e676c65746f6e5f69666627}\textbf{Subset singleton iff'.} Let \(\alpha\) be a type. Let \(s\) be a finite subset of \(\alpha\). Let \(a\) be an element of \(\alpha\). Then \(s \subseteq \{a\} \iff \left(\forall b\, :\, \alpha,\; b \in s \implies b = a\right)\).
\item[\paperlineanchor{premise:46696e7365742e756e696f6e5f636f6d6d}\paperlink{premise:46696e7365742e756e696f6e5f636f6d6d}{P83.}] \refstepcounter{paperpremise}\label{premise:46696e7365742e756e696f6e5f636f6d6d}\textbf{Union commutation.} Let \(\alpha\) be a type. Assume equality on \(\alpha\) is decidable. Let \(s_{1}\) and \(s_{2}\) be finite subsets of \(\alpha\). Then \(s_{1} \cup s_{2} = s_{2} \cup s_{1}\).
\item[\paperlineanchor{premise:46696e7365742e756e696f6e5f656d707479}\paperlink{premise:46696e7365742e756e696f6e5f656d707479}{P84.}] \refstepcounter{paperpremise}\label{premise:46696e7365742e756e696f6e5f656d707479}\textbf{Union empty.} Let \(\alpha\) be a type. Assume equality on \(\alpha\) is decidable. Let \(s\) be a finite subset of \(\alpha\). Then \(s \cup \varnothing = s\).
\item[\paperlineanchor{premise:46696e7365742e756e696f6e5f696e73657274}\paperlink{premise:46696e7365742e756e696f6e5f696e73657274}{P85.}] \refstepcounter{paperpremise}\label{premise:46696e7365742e756e696f6e5f696e73657274}\textbf{Union insert.} Let \(\alpha\) be a type. Assume equality on \(\alpha\) is decidable. Let \(a\) be an element of \(\alpha\). Let \(s\) and \(t\) be finite subsets of \(\alpha\). Then \(s \cup \left(\{a\} \cup t\right) = \{a\} \cup \left(s \cup t\right)\).
\item[\paperlineanchor{premise:46696e7365742e756e696f6e5f73646966665f73656c665f65715f756e696f6e}\paperlink{premise:46696e7365742e756e696f6e5f73646966665f73656c665f65715f756e696f6e}{P86.}] \refstepcounter{paperpremise}\label{premise:46696e7365742e756e696f6e5f73646966665f73656c665f65715f756e696f6e}\textbf{Union sdiff self equals union.} Let \(\alpha\) be a type. Assume equality on \(\alpha\) is decidable. Let \(s\) and \(t\) be finite subsets of \(\alpha\). Then \(s \cup \left(t \setminus s\right) = s \cup t\).
\item[\paperlineanchor{premise:46696e7365742e756e696f6e5f737562736574}\paperlink{premise:46696e7365742e756e696f6e5f737562736574}{P87.}] \refstepcounter{paperpremise}\label{premise:46696e7365742e756e696f6e5f737562736574}\textbf{Union subset.} Let \(\alpha\) be a type. Assume equality on \(\alpha\) is decidable. Let \(s\), \(t\), and \(u\) be finite subsets of \(\alpha\). Assume \(s \subseteq u\). Assume \(t \subseteq u\). Then \(s \cup t \subseteq u\).
\item[\paperlineanchor{premise:46696e7365742e76616c5f65715f7a65726f}\paperlink{premise:46696e7365742e76616c5f65715f7a65726f}{P88.}] \refstepcounter{paperpremise}\label{premise:46696e7365742e76616c5f65715f7a65726f}\textbf{Val equals zero.} Let \(\alpha\) be a type. Let \(s\) be a finite subset of \(\alpha\). Then \(\operatorname{underlying\,multiset}(s) = 0 \iff s = \varnothing\).
\item[\paperlineanchor{premise:46696e7365742e76616c5f696e6a}\paperlink{premise:46696e7365742e76616c5f696e6a}{P89.}] \refstepcounter{paperpremise}\label{premise:46696e7365742e76616c5f696e6a}\textbf{Val injectivity.} Let \(\alpha\) be a type. Let \(s\) and \(t\) be finite subsets of \(\alpha\). Then \(\operatorname{underlying\,multiset}(s) = \operatorname{underlying\,multiset}(t) \iff s = t\).
\item[\paperlineanchor{premise:46696e7365742e76616c5f6c655f696666}\paperlink{premise:46696e7365742e76616c5f6c655f696666}{P90.}] \refstepcounter{paperpremise}\label{premise:46696e7365742e76616c5f6c655f696666}\textbf{Val at most if and only if.} Let \(\alpha\) be a type. Let \(s_{1}\) and \(s_{2}\) be finite subsets of \(\alpha\). Then \(\operatorname{underlying\,multiset}(s_{1}) \le \operatorname{underlying\,multiset}(s_{2}) \iff s_{1} \subseteq s_{2}\).
\item[\paperlineanchor{premise:46696e7365742e76616c5f6c745f696666}\paperlink{premise:46696e7365742e76616c5f6c745f696666}{P91.}] \refstepcounter{paperpremise}\label{premise:46696e7365742e76616c5f6c745f696666}\textbf{Val less than if and only if.} Let \(\alpha\) be a type. Let \(s_{1}\) and \(s_{2}\) be finite subsets of \(\alpha\). Then \(\operatorname{underlying\,multiset}(s_{1}) < \operatorname{underlying\,multiset}(s_{2}) \iff s_{1} \subsetneq s_{2}\).
\item[\paperlineanchor{premise:46756e6374696f6e2e6e655f696666}\paperlink{premise:46756e6374696f6e2e6e655f696666}{P92.}] \refstepcounter{paperpremise}\label{premise:46756e6374696f6e2e6e655f696666}\textbf{Differs from if and only if.} Let \(\alpha\) be a sort. Let \(\beta\) be a function of type \(\alpha \to \mathsf{Sort}\). Let \(f_{1}\) be a function of type \(\Pi_{a\, :\, \alpha}\, \beta\left(a\right)\). Let \(f_{2}\) be a function of type \(\Pi_{a\, :\, \alpha}\, \beta\left(a\right)\). Then \(f_{1} \ne f_{2} \iff \left(\exists a\, :\, \alpha,\; f_{1}\left(a\right) \ne f_{2}\left(a\right)\right)\).
\item[\paperlineanchor{premise:4966662e6d70}\paperlink{premise:4966662e6d70}{P93.}] \refstepcounter{paperpremise}\label{premise:4966662e6d70}\textbf{Mp.} Let \(a\) and \(b\) be propositions. Assume \(a \iff b\). Assume \(a\). Then \(b\).
\item[\paperlineanchor{premise:4966662e6e65}\paperlink{premise:4966662e6e65}{P94.}] \refstepcounter{paperpremise}\label{premise:4966662e6e65}\textbf{Differs from.} Let \(\alpha\) be a sort. Let \(\beta\) be a sort. Let \(a\) and \(b\) be elements of \(\alpha\). Let \(c\) and \(d\) be elements of \(\beta\). Assume \(a = b \iff c = d\). Then \(a \ne b \iff c \ne d\).
\item[\paperlineanchor{premise:4966662e6f665f6571}\paperlink{premise:4966662e6f665f6571}{P95.}] \refstepcounter{paperpremise}\label{premise:4966662e6f665f6571}\textbf{Of equals.} Let \(a\) and \(b\) be propositions. Assume \(a = b\). Then \(a \iff b\).
\item[\paperlineanchor{premise:4966662e72666c}\paperlink{premise:4966662e72666c}{P96.}] \refstepcounter{paperpremise}\label{premise:4966662e72666c}\textbf{Rfl.} Let \(a\) be a proposition. Then \(a \iff a\).
\item[\paperlineanchor{premise:496e742e6164645f6f6e655f6c655f6f665f6c74}\paperlink{premise:496e742e6164645f6f6e655f6c655f6f665f6c74}{P97.}] \refstepcounter{paperpremise}\label{premise:496e742e6164645f6f6e655f6c655f6f665f6c74}\textbf{Sum one at most of less than.} Let \(a\) and \(b\) be integers. Assume \(a < b\). Then \(a + 1 \le b\).
\item[\paperlineanchor{premise:496e742e6f664e61745f616464}\paperlink{premise:496e742e6f664e61745f616464}{P98.}] \refstepcounter{paperpremise}\label{premise:496e742e6f664e61745f616464}\textbf{Of nat sum.} Let \(n\) and \(m\) be natural numbers. Then \(n + m = n + m\).
\item[\paperlineanchor{premise:496e742e6f664e61745f6e6f6e6e6567}\paperlink{premise:496e742e6f664e61745f6e6f6e6e6567}{P99.}] \refstepcounter{paperpremise}\label{premise:496e742e6f664e61745f6e6f6e6e6567}\textbf{Of nat nonnegative.} Let \(n\) be a natural number. Then \(0 \le n\).
\item[\paperlineanchor{premise:496e742e7375625f65715f7a65726f5f6f665f6571}\paperlink{premise:496e742e7375625f65715f7a65726f5f6f665f6571}{P100.}] \refstepcounter{paperpremise}\label{premise:496e742e7375625f65715f7a65726f5f6f665f6571}\textbf{Difference equals zero of equals.} Let \(a\) and \(b\) be integers. Assume \(a = b\). Then \(a - b = 0\).
\item[\paperlineanchor{premise:496e742e7375625f6e6f6e6e65675f6f665f6c65}\paperlink{premise:496e742e7375625f6e6f6e6e65675f6f665f6c65}{P101.}] \refstepcounter{paperpremise}\label{premise:496e742e7375625f6e6f6e6e65675f6f665f6c65}\textbf{Difference nonnegative of at most.} Let \(a\) and \(b\) be integers. Assume \(b \le a\). Then \(0 \le a - b\).
\item[\paperlineanchor{premise:496e76496d6167652e7766}\paperlink{premise:496e76496d6167652e7766}{P102.}] \refstepcounter{paperpremise}\label{premise:496e76496d6167652e7766}\textbf{Wf.} Let \(\alpha\) be a sort. Let \(\beta\) be a sort. Let \(r\) be a relation from \(\beta\) to \(\beta\). Let \(f\) be a function of type \(\alpha \to \beta\). Assume \(\operatorname{Well\,Founded}\left(r\right)\). Then \(\operatorname{Well\,Founded}\left(\operatorname{Inv\,Image}\left(r, f\right)\right)\).
\item[\paperlineanchor{premise:4973456d7074792e666f72616c6c5f696666}\paperlink{premise:4973456d7074792e666f72616c6c5f696666}{P103.}] \refstepcounter{paperpremise}\label{premise:4973456d7074792e666f72616c6c5f696666}\textbf{Forall if and only if.} Let \(\alpha\) be a sort. Assume \(\alpha\text{ is empty}\). Let \(p\) be a function of type \(\alpha \to \mathsf{Prop}\). Then \(\left(\forall a\, :\, \alpha,\; p\left(a\right)\right) \iff \mathrm{True}\).
\item[\paperlineanchor{premise:4c452e6c652e65715f6966665f6e6f745f6c74}\paperlink{premise:4c452e6c652e65715f6966665f6e6f745f6c74}{P104.}] \refstepcounter{paperpremise}\label{premise:4c452e6c652e65715f6966665f6e6f745f6c74}\textbf{Alias.} Let \(\alpha\) be a type. Equip \(\alpha\) with a partial order. Let \(a\) and \(b\) be elements of \(\alpha\). Assume \(a \le b\). Then \(a = b \iff \neg \left(a < b\right)\).
\item[\paperlineanchor{premise:4c452e6c652e6c655f6966665f6571}\paperlink{premise:4c452e6c652e6c655f6966665f6571}{P105.}] \refstepcounter{paperpremise}\label{premise:4c452e6c652e6c655f6966665f6571}\textbf{At most if and only if equals.} Let \(\alpha\) be a type. Equip \(\alpha\) with a partial order. Let \(a\) and \(b\) be elements of \(\alpha\). Assume \(a \le b\). Then \(b \le a \iff b = a\).
\item[\paperlineanchor{premise:4c452e6c652e7472616e73}\paperlink{premise:4c452e6c652e7472616e73}{P106.}] \refstepcounter{paperpremise}\label{premise:4c452e6c652e7472616e73}\textbf{Alias.} Let \(\alpha\) be a type. Equip \(\alpha\) with a preorder. Let \(a\), \(b\), and \(c\) be elements of \(\alpha\). Assume \(a \le b\). Assume \(b \le c\). Then \(a \le c\).
\item[\paperlineanchor{premise:4c452e6c652e7472616e735f6571}\paperlink{premise:4c452e6c652e7472616e735f6571}{P107.}] \refstepcounter{paperpremise}\label{premise:4c452e6c652e7472616e735f6571}\textbf{Alias.} Let \(\alpha\) be a type. Let \(a\), \(b\), and \(c\) be elements of \(\alpha\). Fix a structure \(\mathit{inst}\) of type \(\operatorname{LE}\left(\alpha\right)\). Assume \(a \le b\). Assume \(b = c\). Then \(a \le c\).
\item[\paperlineanchor{premise:4c542e6c742e6e6f745f6c65}\paperlink{premise:4c542e6c742e6e6f745f6c65}{P108.}] \refstepcounter{paperpremise}\label{premise:4c542e6c742e6e6f745f6c65}\textbf{Alias.} Let \(\alpha\) be a type. Equip \(\alpha\) with a preorder. Let \(a\) and \(b\) be elements of \(\alpha\). Assume \(a < b\). Then \(\neg \left(b \le a\right)\).
\item[\paperlineanchor{premise:4c65616e2e4f6d6567612e436f6e73747261696e742e616464457175616c6974795f736174}\paperlink{premise:4c65616e2e4f6d6567612e436f6e73747261696e742e616464457175616c6974795f736174}{P109.}] \refstepcounter{paperpremise}\label{premise:4c65616e2e4f6d6567612e436f6e73747261696e742e616464457175616c6974795f736174}\textbf{Sum equality sat.} Let \(c\) be an integer. Let \(x\) and \(y\) be elements of \(\operatorname{Coeffs}\). Assume \(c + \operatorname{dot}\left(x, y\right) = 0\). Then \(\operatorname{sat^{\prime}}(\operatorname{Lean\mathord{:}Omega\mathord{:}Constraint\mathord{:}mk}(\operatorname{some}(-c),\allowbreak{} \operatorname{some}(-c)),\allowbreak{} x,\allowbreak{} y) = \mathrm{true}\).
\item[\paperlineanchor{premise:4c65616e2e4f6d6567612e436f6e73747261696e742e616464496e657175616c6974795f736174}\paperlink{premise:4c65616e2e4f6d6567612e436f6e73747261696e742e616464496e657175616c6974795f736174}{P110.}] \refstepcounter{paperpremise}\label{premise:4c65616e2e4f6d6567612e436f6e73747261696e742e616464496e657175616c6974795f736174}\textbf{Sum inequality sat.} Let \(c\) be an integer. Let \(x\) and \(y\) be elements of \(\operatorname{Coeffs}\). Assume \(c + \operatorname{dot}\left(x, y\right) \ge 0\). Then \(\operatorname{sat^{\prime}}(\operatorname{Lean\mathord{:}Omega\mathord{:}Constraint\mathord{:}mk}(\operatorname{some}(-c),\allowbreak{} \operatorname{none}),\allowbreak{} x,\allowbreak{} y) = \mathrm{true}\).
\item[\paperlineanchor{premise:4c65616e2e4f6d6567612e436f6e73747261696e742e636f6d62696e655f73617427}\paperlink{premise:4c65616e2e4f6d6567612e436f6e73747261696e742e636f6d62696e655f73617427}{P111.}] \refstepcounter{paperpremise}\label{premise:4c65616e2e4f6d6567612e436f6e73747261696e742e636f6d62696e655f73617427}\textbf{Combine sat'.} Let \(s\) and \(t\) be elements of \(\operatorname{Constraint}\). Let \(x\) and \(y\) be elements of \(\operatorname{Coeffs}\). Assume \(\operatorname{sat^{\prime}}\left(s, x, y\right) = \mathrm{true}\). Assume \(\operatorname{sat^{\prime}}\left(t, x, y\right) = \mathrm{true}\). Then \(\operatorname{sat^{\prime}}\left(\operatorname{combine}\left(s, t\right), x, y\right) = \mathrm{true}\).
\item[\paperlineanchor{premise:4c65616e2e4f6d6567612e436f6e73747261696e742e6e6f745f736174275f6f665f6973496d706f737369626c65}\paperlink{premise:4c65616e2e4f6d6567612e436f6e73747261696e742e6e6f745f736174275f6f665f6973496d706f737369626c65}{P112.}] \refstepcounter{paperpremise}\label{premise:4c65616e2e4f6d6567612e436f6e73747261696e742e6e6f745f736174275f6f665f6973496d706f737369626c65}\textbf{Not sat' of is impossible.} Let \(c\) be an element of \(\operatorname{Constraint}\). Assume \(\operatorname{is\,Impossible}\left(c\right) = \mathrm{true}\). Let \(x\) and \(y\) be elements of \(\operatorname{Coeffs}\). Then \(\neg \left(\operatorname{sat^{\prime}}\left(c, x, y\right) = \mathrm{true}\right)\).
\item[\paperlineanchor{premise:4c65616e2e4f6d6567612e496e742e6164645f636f6e6772}\paperlink{premise:4c65616e2e4f6d6567612e496e742e6164645f636f6e6772}{P113.}] \refstepcounter{paperpremise}\label{premise:4c65616e2e4f6d6567612e496e742e6164645f636f6e6772}\textbf{Sum congruence.} Let \(a\), \(b\), \(c\), and \(d\) be integers. Assume \(a = b\). Assume \(c = d\). Then \(a + c = b + d\).
\item[\paperlineanchor{premise:4c65616e2e4f6d6567612e496e742e6f664e61745f636f6e6772}\paperlink{premise:4c65616e2e4f6d6567612e496e742e6f664e61745f636f6e6772}{P114.}] \refstepcounter{paperpremise}\label{premise:4c65616e2e4f6d6567612e496e742e6f664e61745f636f6e6772}\textbf{Of nat congruence.} Let \(a\) and \(b\) be natural numbers. Assume \(a = b\). Then \(a = b\).
\item[\paperlineanchor{premise:4c65616e2e4f6d6567612e496e742e6f664e61745f6c655f6f665f6c65}\paperlink{premise:4c65616e2e4f6d6567612e496e742e6f664e61745f6c655f6f665f6c65}{P115.}] \refstepcounter{paperpremise}\label{premise:4c65616e2e4f6d6567612e496e742e6f664e61745f6c655f6f665f6c65}\textbf{Of nat at most of at most.} Let \(x\) and \(y\) be natural numbers. Assume \(x \le y\). Then \(x \le y\).
\item[\paperlineanchor{premise:4c65616e2e4f6d6567612e496e742e6f664e61745f6c745f6f665f6c74}\paperlink{premise:4c65616e2e4f6d6567612e496e742e6f664e61745f6c745f6f665f6c74}{P116.}] \refstepcounter{paperpremise}\label{premise:4c65616e2e4f6d6567612e496e742e6f664e61745f6c745f6f665f6c74}\textbf{Of nat less than of less than.} Let \(x\) and \(y\) be natural numbers. Assume \(x < y\). Then \(x < y\).
\item[\paperlineanchor{premise:4c65616e2e4f6d6567612e496e742e6f664e61745f7375625f646963686f746f6d79}\paperlink{premise:4c65616e2e4f6d6567612e496e742e6f664e61745f7375625f646963686f746f6d79}{P117.}] \refstepcounter{paperpremise}\label{premise:4c65616e2e4f6d6567612e496e742e6f664e61745f7375625f646963686f746f6d79}\textbf{Of nat difference dichotomy.} Let \(a\) and \(b\) be natural numbers. Then \(b \le a \land a - b = a - b \lor a < b \land a - b = 0\).
\item[\paperlineanchor{premise:4c65616e2e4f6d6567612e496e742e7375625f636f6e6772}\paperlink{premise:4c65616e2e4f6d6567612e496e742e7375625f636f6e6772}{P118.}] \refstepcounter{paperpremise}\label{premise:4c65616e2e4f6d6567612e496e742e7375625f636f6e6772}\textbf{Difference congruence.} Let \(a\), \(b\), \(c\), and \(d\) be integers. Assume \(a = b\). Assume \(c = d\). Then \(a - c = b - d\).
\item[\paperlineanchor{premise:4c65616e2e4f6d6567612e4c696e656172436f6d626f2e6164645f6576616c}\paperlink{premise:4c65616e2e4f6d6567612e4c696e656172436f6d626f2e6164645f6576616c}{P119.}] \refstepcounter{paperpremise}\label{premise:4c65616e2e4f6d6567612e4c696e656172436f6d626f2e6164645f6576616c}\textbf{Sum eval.} Let \(l_{1}\) and \(l_{2}\) be elements of \(\operatorname{Linear\,Combo}\). Let \(v\) be an element of \(\operatorname{Coeffs}\). Then \(\operatorname{eval}(l_{1} + l_{2},\allowbreak{} v) = \operatorname{eval}(l_{1},\allowbreak{} v) + \operatorname{eval}(l_{2},\allowbreak{} v)\).
\item[\paperlineanchor{premise:4c65616e2e4f6d6567612e4c696e656172436f6d626f2e636f6f7264696e6174655f6576616c5f30}\paperlink{premise:4c65616e2e4f6d6567612e4c696e656172436f6d626f2e636f6f7264696e6174655f6576616c5f30}{P120.}] \refstepcounter{paperpremise}\label{premise:4c65616e2e4f6d6567612e4c696e656172436f6d626f2e636f6f7264696e6174655f6576616c5f30}\textbf{Coordinate eval 0.} Let \(\mathit{a0}\) be an integer. Let \(t\) be a list of elements of \(\mathbb{Z}\). Then \(\operatorname{eval}(\operatorname{coordinate}(0),\allowbreak{} \operatorname{Lean\mathord{:}Omega\mathord{:}Coeffs\mathord{:}of\,List}(\langle \mathit{a0}\rangle \mathbin{\Vert} t)) = \mathit{a0}\).
\item[\paperlineanchor{premise:4c65616e2e4f6d6567612e4c696e656172436f6d626f2e636f6f7264696e6174655f6576616c5f31}\paperlink{premise:4c65616e2e4f6d6567612e4c696e656172436f6d626f2e636f6f7264696e6174655f6576616c5f31}{P121.}] \refstepcounter{paperpremise}\label{premise:4c65616e2e4f6d6567612e4c696e656172436f6d626f2e636f6f7264696e6174655f6576616c5f31}\textbf{Coordinate eval 1.} Let \(\mathit{a0}\) and \(\mathit{a1}\) be integers. Let \(t\) be a list of elements of \(\mathbb{Z}\). Then \(\operatorname{eval}(\operatorname{coordinate}(1),\allowbreak{} \operatorname{Lean\mathord{:}Omega\mathord{:}Coeffs\mathord{:}of\,List}(\langle \mathit{a0}\rangle \mathbin{\Vert} (\langle \mathit{a1}\rangle \mathbin{\Vert} t))) = \mathit{a1}\).
\item[\paperlineanchor{premise:4c65616e2e4f6d6567612e4c696e656172436f6d626f2e636f6f7264696e6174655f6576616c5f32}\paperlink{premise:4c65616e2e4f6d6567612e4c696e656172436f6d626f2e636f6f7264696e6174655f6576616c5f32}{P122.}] \refstepcounter{paperpremise}\label{premise:4c65616e2e4f6d6567612e4c696e656172436f6d626f2e636f6f7264696e6174655f6576616c5f32}\textbf{Coordinate eval 2.} Let \(\mathit{a0}\), \(\mathit{a1}\), and \(\mathit{a2}\) be integers. Let \(t\) be a list of elements of \(\mathbb{Z}\). Then \(\operatorname{eval}(\operatorname{coordinate}(2),\allowbreak{} \operatorname{Lean\mathord{:}Omega\mathord{:}Coeffs\mathord{:}of\,List}(\langle \mathit{a0}\rangle \mathbin{\Vert} (\langle \mathit{a1}\rangle \mathbin{\Vert} (\langle \mathit{a2}\rangle \mathbin{\Vert} t)))) = \mathit{a2}\).
\item[\paperlineanchor{premise:4c65616e2e4f6d6567612e4c696e656172436f6d626f2e636f6f7264696e6174655f6576616c5f33}\paperlink{premise:4c65616e2e4f6d6567612e4c696e656172436f6d626f2e636f6f7264696e6174655f6576616c5f33}{P123.}] \refstepcounter{paperpremise}\label{premise:4c65616e2e4f6d6567612e4c696e656172436f6d626f2e636f6f7264696e6174655f6576616c5f33}\textbf{Coordinate eval 3.} Let \(\mathit{a0}\), \(\mathit{a1}\), \(\mathit{a2}\), and \(\mathit{a3}\) be integers. Let \(t\) be a list of elements of \(\mathbb{Z}\). Then \(\operatorname{eval}(\operatorname{coordinate}(3),\allowbreak{} \operatorname{Lean\mathord{:}Omega\mathord{:}Coeffs\mathord{:}of\,List}(\langle \mathit{a0}\rangle \mathbin{\Vert} (\langle \mathit{a1}\rangle \mathbin{\Vert} (\langle \mathit{a2}\rangle \mathbin{\Vert} (\langle \mathit{a3}\rangle \mathbin{\Vert} t))))) = \mathit{a3}\).
\item[\paperlineanchor{premise:4c65616e2e4f6d6567612e4c696e656172436f6d626f2e636f6f7264696e6174655f6576616c5f34}\paperlink{premise:4c65616e2e4f6d6567612e4c696e656172436f6d626f2e636f6f7264696e6174655f6576616c5f34}{P124.}] \refstepcounter{paperpremise}\label{premise:4c65616e2e4f6d6567612e4c696e656172436f6d626f2e636f6f7264696e6174655f6576616c5f34}\textbf{Coordinate eval 4.} Let \(\mathit{a0}\), \(\mathit{a1}\), \(\mathit{a2}\), \(\mathit{a3}\), and \(\mathit{a4}\) be integers. Let \(t\) be a list of elements of \(\mathbb{Z}\). Then \(\operatorname{eval}(\operatorname{coordinate}(4),\allowbreak{} \operatorname{Lean\mathord{:}Omega\mathord{:}Coeffs\mathord{:}of\,List}(\langle \mathit{a0}\rangle \mathbin{\Vert} (\langle \mathit{a1}\rangle \mathbin{\Vert} (\langle \mathit{a2}\rangle \mathbin{\Vert} (\langle \mathit{a3}\rangle \mathbin{\Vert} (\langle \mathit{a4}\rangle \mathbin{\Vert} t)))))) = \mathit{a4}\).
\item[\paperlineanchor{premise:4c65616e2e4f6d6567612e4c696e656172436f6d626f2e636f6f7264696e6174655f6576616c5f35}\paperlink{premise:4c65616e2e4f6d6567612e4c696e656172436f6d626f2e636f6f7264696e6174655f6576616c5f35}{P125.}] \refstepcounter{paperpremise}\label{premise:4c65616e2e4f6d6567612e4c696e656172436f6d626f2e636f6f7264696e6174655f6576616c5f35}\textbf{Coordinate eval 5.} Let \(\mathit{a0}\), \(\mathit{a1}\), \(\mathit{a2}\), \(\mathit{a3}\), \(\mathit{a4}\), and \(\mathit{a5}\) be integers. Let \(t\) be a list of elements of \(\mathbb{Z}\). Then \(\operatorname{eval}(\operatorname{coordinate}(5),\allowbreak{} \operatorname{Lean\mathord{:}Omega\mathord{:}Coeffs\mathord{:}of\,List}(\langle \mathit{a0}\rangle \mathbin{\Vert} (\langle \mathit{a1}\rangle \mathbin{\Vert} (\langle \mathit{a2}\rangle \mathbin{\Vert} (\langle \mathit{a3}\rangle \mathbin{\Vert} (\langle \mathit{a4}\rangle \mathbin{\Vert} (\langle \mathit{a5}\rangle \mathbin{\Vert} t))))))) = \mathit{a5}\).
\item[\paperlineanchor{premise:4c65616e2e4f6d6567612e4c696e656172436f6d626f2e7375625f6576616c}\paperlink{premise:4c65616e2e4f6d6567612e4c696e656172436f6d626f2e7375625f6576616c}{P126.}] \refstepcounter{paperpremise}\label{premise:4c65616e2e4f6d6567612e4c696e656172436f6d626f2e7375625f6576616c}\textbf{Difference eval.} Let \(l_{1}\) and \(l_{2}\) be elements of \(\operatorname{Linear\,Combo}\). Let \(v\) be an element of \(\operatorname{Coeffs}\). Then \(\operatorname{eval}(l_{1} - l_{2},\allowbreak{} v) = \operatorname{eval}(l_{1},\allowbreak{} v) - \operatorname{eval}(l_{2},\allowbreak{} v)\).
\item[\paperlineanchor{premise:4c65616e2e4f6d6567612e636f6d626f5f73617427}\paperlink{premise:4c65616e2e4f6d6567612e636f6d626f5f73617427}{P127.}] \refstepcounter{paperpremise}\label{premise:4c65616e2e4f6d6567612e636f6d626f5f73617427}\textbf{Combo sat'.} Let \(s\) and \(t\) be elements of \(\operatorname{Constraint}\). Let \(a\) be an integer. Let \(x\) be an element of \(\operatorname{Coeffs}\). Let \(b\) be an integer. Let \(y\) and \(v\) be elements of \(\operatorname{Coeffs}\). Assume \(\operatorname{sat^{\prime}}\left(s, x, v\right) = \mathrm{true}\). Assume \(\operatorname{sat^{\prime}}\left(t, y, v\right) = \mathrm{true}\). Then \(\operatorname{sat^{\prime}}(\operatorname{Lean\mathord{:}Omega\mathord{:}Constraint\mathord{:}combo}(a,\allowbreak{} s,\allowbreak{} b,\allowbreak{} t),\allowbreak{} \operatorname{Lean\mathord{:}Omega\mathord{:}Coeffs\mathord{:}combo}(a,\allowbreak{} x,\allowbreak{} b,\allowbreak{} y),\allowbreak{} v) = \mathrm{true}\).
\item[\paperlineanchor{premise:4c65616e2e4f6d6567612e746964795f736174}\paperlink{premise:4c65616e2e4f6d6567612e746964795f736174}{P128.}] \refstepcounter{paperpremise}\label{premise:4c65616e2e4f6d6567612e746964795f736174}\textbf{Tidy sat.} Let \(s\) be an element of \(\operatorname{Constraint}\). Let \(x\) and \(v\) be elements of \(\operatorname{Coeffs}\). Assume \(\operatorname{sat^{\prime}}\left(s, x, v\right) = \mathrm{true}\). Then \(\operatorname{sat^{\prime}}(\operatorname{tidy\,Constraint}(s,\allowbreak{} x),\allowbreak{} \operatorname{tidy\,Coeffs}(s,\allowbreak{} x),\allowbreak{} v) = \mathrm{true}\).
\item[\paperlineanchor{premise:4d6174686c69622e5461637469632e436f6e747261706f73652e6d7472}\paperlink{premise:4d6174686c69622e5461637469632e436f6e747261706f73652e6d7472}{P129.}] \refstepcounter{paperpremise}\label{premise:4d6174686c69622e5461637469632e436f6e747261706f73652e6d7472}\textbf{Mtr.} Let \(p\) and \(q\) be propositions. Assume \(\neg q \implies \neg p\). Assume \(p\). Then \(q\).
\item[\paperlineanchor{premise:4d6174686c69622e5461637469632e507573684e65672e6e6f745f616e645f6571}\paperlink{premise:4d6174686c69622e5461637469632e507573684e65672e6e6f745f616e645f6571}{P130.}] \refstepcounter{paperpremise}\label{premise:4d6174686c69622e5461637469632e507573684e65672e6e6f745f616e645f6571}\textbf{Not and equals.} Let \(p\) and \(q\) be propositions. Then \(\left(\neg \left(p \land q\right)\right) = \left(p \implies \neg q\right)\).
\item[\paperlineanchor{premise:4d6174686c69622e5461637469632e507573684e65672e6e6f745f6578697374735f6571}\paperlink{premise:4d6174686c69622e5461637469632e507573684e65672e6e6f745f6578697374735f6571}{P131.}] \refstepcounter{paperpremise}\label{premise:4d6174686c69622e5461637469632e507573684e65672e6e6f745f6578697374735f6571}\textbf{Not exists equals.} Let \(\alpha\) be a sort. Let \(s\) be a function of type \(\alpha \to \mathsf{Prop}\). Then \((\neg (\exists x\, :\, \alpha,\; s(x))) = (\forall x\, :\, \alpha,\; \neg s(x))\).
\item[\paperlineanchor{premise:4d6174686c69622e5461637469632e507573684e65672e6e6f745f6c745f6571}\paperlink{premise:4d6174686c69622e5461637469632e507573684e65672e6e6f745f6c745f6571}{P132.}] \refstepcounter{paperpremise}\label{premise:4d6174686c69622e5461637469632e507573684e65672e6e6f745f6c745f6571}\textbf{Not less than equals.} Let \(\beta\) be a type. Equip \(\beta\) with a linear order. Let \(a\) and \(b\) be elements of \(\beta\). Then \(\left(\neg \left(a < b\right)\right) = \left(b \le a\right)\).
\item[\paperlineanchor{premise:4d6174686c69622e5461637469632e507573684e65672e6e6f745f6e655f6571}\paperlink{premise:4d6174686c69622e5461637469632e507573684e65672e6e6f745f6e655f6571}{P133.}] \refstepcounter{paperpremise}\label{premise:4d6174686c69622e5461637469632e507573684e65672e6e6f745f6e655f6571}\textbf{Not differs from equals.} Let \(\alpha\) be a sort. Let \(x\) and \(y\) be elements of \(\alpha\). Then \(\left(\neg \left(x \ne y\right)\right) = \left(x = y\right)\).
\item[\paperlineanchor{premise:4d6174686c69622e5461637469632e507573684e65672e6e6f745f6f725f6571}\paperlink{premise:4d6174686c69622e5461637469632e507573684e65672e6e6f745f6f725f6571}{P134.}] \refstepcounter{paperpremise}\label{premise:4d6174686c69622e5461637469632e507573684e65672e6e6f745f6f725f6571}\textbf{Not or equals.} Let \(p\) and \(q\) be propositions. Then \(\left(\neg \left(p \lor q\right)\right) = \left(\neg p \land \neg q\right)\).
\item[\paperlineanchor{premise:4d756c74697365742e636172645f616464}\paperlink{premise:4d756c74697365742e636172645f616464}{P135.}] \refstepcounter{paperpremise}\label{premise:4d756c74697365742e636172645f616464}\textbf{Cardinality sum.} Let \(\alpha\) be a type. Let \(s\) and \(t\) be multisets of elements of \(\alpha\). Then \(\left|s + t\right| = \left|s\right| + \left|t\right|\).
\item[\paperlineanchor{premise:4d756c74697365742e636172645f617474616368}\paperlink{premise:4d756c74697365742e636172645f617474616368}{P136.}] \refstepcounter{paperpremise}\label{premise:4d756c74697365742e636172645f617474616368}\textbf{Cardinality attach.} Let \(\alpha\) be a type. Let \(m\) be a multiset of elements of \(\alpha\). Then \(\left|\operatorname{Multiset\mathord{:}attach}\left(m\right)\right| = \left|m\right|\).
\item[\paperlineanchor{premise:4d756c74697365742e636172645f636f6e73}\paperlink{premise:4d756c74697365742e636172645f636f6e73}{P137.}] \refstepcounter{paperpremise}\label{premise:4d756c74697365742e636172645f636f6e73}\textbf{Cardinality cons.} Let \(\alpha\) be a type. Let \(a\) be an element of \(\alpha\). Let \(s\) be a multiset of elements of \(\alpha\). Then \(\left|\operatorname{Multiset\mathord{:}cons}\left(a, s\right)\right| = \left|s\right| + 1\).
\item[\paperlineanchor{premise:4d756c74697365742e636172645f65715f6f6e65}\paperlink{premise:4d756c74697365742e636172645f65715f6f6e65}{P138.}] \refstepcounter{paperpremise}\label{premise:4d756c74697365742e636172645f65715f6f6e65}\textbf{Cardinality equals one.} Let \(\alpha\) be a type. Let \(s\) be a multiset of elements of \(\alpha\). Then \(\left|s\right| = 1 \iff \left(\exists a\, :\, \alpha,\; s = \{a\}\right)\).
\item[\paperlineanchor{premise:4d756c74697365742e636172645f65715f7a65726f}\paperlink{premise:4d756c74697365742e636172645f65715f7a65726f}{P139.}] \refstepcounter{paperpremise}\label{premise:4d756c74697365742e636172645f65715f7a65726f}\textbf{Cardinality equals zero.} Let \(\alpha\) be a type. Let \(s\) be a multiset of elements of \(\alpha\). Then \(\left|s\right| = 0 \iff s = 0\).
\item[\paperlineanchor{premise:4d756c74697365742e636172645f65726173655f6164645f6f6e65}\paperlink{premise:4d756c74697365742e636172645f65726173655f6164645f6f6e65}{P140.}] \refstepcounter{paperpremise}\label{premise:4d756c74697365742e636172645f65726173655f6164645f6f6e65}\textbf{Cardinality erase sum one.} Let \(\alpha\) be a type. Assume equality on \(\alpha\) is decidable. Let \(a\) be an element of \(\alpha\). Let \(s\) be a multiset of elements of \(\alpha\). Assume \(a \in s\). Then \(\left|\operatorname{Multiset\mathord{:}erase}\left(s, a\right)\right| + 1 = \left|s\right|\).
\item[\paperlineanchor{premise:4d756c74697365742e636172645f65726173655f65715f697465}\paperlink{premise:4d756c74697365742e636172645f65726173655f65715f697465}{P141.}] \refstepcounter{paperpremise}\label{premise:4d756c74697365742e636172645f65726173655f65715f697465}\textbf{Cardinality erase equals ite.} Let \(\alpha\) be a type. Assume equality on \(\alpha\) is decidable. Let \(a\) be an element of \(\alpha\). Let \(s\) be a multiset of elements of \(\alpha\). Then \(|\operatorname{Multiset\mathord{:}erase}(s,\allowbreak{} a)| = \begin{cases}\operatorname{pred}(|s|) &\text{if } a \in s\\ |s| &\text{otherwise}\end{cases}\).
\item[\paperlineanchor{premise:4d756c74697365742e636172645f65726173655f6c65}\paperlink{premise:4d756c74697365742e636172645f65726173655f6c65}{P142.}] \refstepcounter{paperpremise}\label{premise:4d756c74697365742e636172645f65726173655f6c65}\textbf{Cardinality erase at most.} Let \(\alpha\) be a type. Assume equality on \(\alpha\) is decidable. Let \(a\) be an element of \(\alpha\). Let \(s\) be a multiset of elements of \(\alpha\). Then \(\left|\operatorname{Multiset\mathord{:}erase}\left(s, a\right)\right| \le \left|s\right|\).
\item[\paperlineanchor{premise:4d756c74697365742e636172645f65726173655f6c745f6f665f6d656d}\paperlink{premise:4d756c74697365742e636172645f65726173655f6c745f6f665f6d656d}{P143.}] \refstepcounter{paperpremise}\label{premise:4d756c74697365742e636172645f65726173655f6c745f6f665f6d656d}\textbf{Cardinality erase less than of membership.} Let \(\alpha\) be a type. Assume equality on \(\alpha\) is decidable. Let \(a\) be an element of \(\alpha\). Let \(s\) be a multiset of elements of \(\alpha\). Assume \(a \in s\). Then \(\left|\operatorname{Multiset\mathord{:}erase}\left(s, a\right)\right| < \left|s\right|\).
\item[\paperlineanchor{premise:4d756c74697365742e636172645f65726173655f6f665f6d656d}\paperlink{premise:4d756c74697365742e636172645f65726173655f6f665f6d656d}{P144.}] \refstepcounter{paperpremise}\label{premise:4d756c74697365742e636172645f65726173655f6f665f6d656d}\textbf{Cardinality erase of membership.} Let \(\alpha\) be a type. Assume equality on \(\alpha\) is decidable. Let \(a\) be an element of \(\alpha\). Let \(s\) be a multiset of elements of \(\alpha\). Assume \(a \in s\). Then \(|\operatorname{Multiset\mathord{:}erase}(s,\allowbreak{} a)| = \operatorname{pred}(|s|)\).
\item[\paperlineanchor{premise:4d756c74697365742e636172645f66696c7465725f6c655f696666}\paperlink{premise:4d756c74697365742e636172645f66696c7465725f6c655f696666}{P145.}] \refstepcounter{paperpremise}\label{premise:4d756c74697365742e636172645f66696c7465725f6c655f696666}\textbf{Cardinality filter at most if and only if.} Let \(\alpha\) be a type. Let \(s\) be a multiset of elements of \(\alpha\). Let \(P\) be a function of type \(\alpha \to \mathsf{Prop}\). Assume the predicate \(P\) is decidable. Let \(n\) be a natural number. Then \(|\operatorname{Multiset\mathord{:}filter}(P,\allowbreak{} s)| \le n \iff (\forall s^{'}\, :\, \operatorname{Multiset}(\alpha),\; s^{'} \le s \implies n < |s^{'}| \implies (\exists a_{2}\, :\, \alpha,\; a_{2} \in s^{'} \land \neg P(a_{2})))\).
\item[\paperlineanchor{premise:4d756c74697365742e636172645f6c655f63617264}\paperlink{premise:4d756c74697365742e636172645f6c655f63617264}{P146.}] \refstepcounter{paperpremise}\label{premise:4d756c74697365742e636172645f6c655f63617264}\textbf{Cardinality at most cardinality.} Let \(\alpha\) be a type. Let \(s\) and \(t\) be multisets of elements of \(\alpha\). Assume \(s \le t\). Then \(\left|s\right| \le \left|t\right|\).
\item[\paperlineanchor{premise:4d756c74697365742e636172645f6c745f63617264}\paperlink{premise:4d756c74697365742e636172645f6c745f63617264}{P147.}] \refstepcounter{paperpremise}\label{premise:4d756c74697365742e636172645f6c745f63617264}\textbf{Cardinality less than cardinality.} Let \(\alpha\) be a type. Let \(s\) and \(t\) be multisets of elements of \(\alpha\). Assume \(s < t\). Then \(\left|s\right| < \left|t\right|\).
\item[\paperlineanchor{premise:4d756c74697365742e636172645f6d6170}\paperlink{premise:4d756c74697365742e636172645f6d6170}{P148.}] \refstepcounter{paperpremise}\label{premise:4d756c74697365742e636172645f6d6170}\textbf{Cardinality map.} Let \(\alpha\) and \(\beta\) be types. Let \(f\) be a function of type \(\alpha \to \beta\). Let \(s\) be a multiset of elements of \(\alpha\). Then \(\left|\operatorname{Multiset\mathord{:}map}\left(f, s\right)\right| = \left|s\right|\).
\item[\paperlineanchor{premise:4d756c74697365742e636172645f72616e6765}\paperlink{premise:4d756c74697365742e636172645f72616e6765}{P149.}] \refstepcounter{paperpremise}\label{premise:4d756c74697365742e636172645f72616e6765}\textbf{Cardinality range.} Let \(n\) be a natural number. Then \(\left|\operatorname{Multiset\mathord{:}range}\left(n\right)\right| = n\).
\item[\paperlineanchor{premise:4d756c74697365742e636172645f73696e676c65746f6e}\paperlink{premise:4d756c74697365742e636172645f73696e676c65746f6e}{P150.}] \refstepcounter{paperpremise}\label{premise:4d756c74697365742e636172645f73696e676c65746f6e}\textbf{Cardinality singleton.} Let \(\alpha\) be a type. Let \(a\) be an element of \(\alpha\). Then \(\left|\{a\}\right| = 1\).
\item[\paperlineanchor{premise:4d756c74697365742e636f655f63617264}\paperlink{premise:4d756c74697365742e636f655f63617264}{P151.}] \refstepcounter{paperpremise}\label{premise:4d756c74697365742e636f655f63617264}\textbf{Coercion cardinality.} Let \(\alpha\) be a type. Let \(l\) be a list of elements of \(\alpha\). Then \(\left|\operatorname{Multiset\mathord{:}of\,List}\left(l\right)\right| = \operatorname{length}(l)\).
\item[\paperlineanchor{premise:4d756c74697365742e636f655f746f4c697374}\paperlink{premise:4d756c74697365742e636f655f746f4c697374}{P152.}] \refstepcounter{paperpremise}\label{premise:4d756c74697365742e636f655f746f4c697374}\textbf{Coercion to list.} Let \(\alpha\) be a type. Let \(s\) be a multiset of elements of \(\alpha\). Then \(\operatorname{Multiset\mathord{:}of\,List}(\operatorname{Multiset\mathord{:}to\,List}(s)) = s\).
\item[\paperlineanchor{premise:4d756c74697365742e64656475705f65715f73656c66}\paperlink{premise:4d756c74697365742e64656475705f65715f73656c66}{P153.}] \refstepcounter{paperpremise}\label{premise:4d756c74697365742e64656475705f65715f73656c66}\textbf{Dedup equals self.} Let \(\alpha\) be a type. Assume equality on \(\alpha\) is decidable. Let \(s\) be a multiset of elements of \(\alpha\). Then \(\operatorname{Multiset\mathord{:}dedup}\left(s\right) = s \iff s\text{ has no repeated elements}\).
\item[\paperlineanchor{premise:4d756c74697365742e64656475705f6c65}\paperlink{premise:4d756c74697365742e64656475705f6c65}{P154.}] \refstepcounter{paperpremise}\label{premise:4d756c74697365742e64656475705f6c65}\textbf{Dedup at most.} Let \(\alpha\) be a type. Assume equality on \(\alpha\) is decidable. Let \(s\) be a multiset of elements of \(\alpha\). Then \(\operatorname{Multiset\mathord{:}dedup}\left(s\right) \le s\).
\item[\paperlineanchor{premise:4d756c74697365742e65715f6f665f6c655f6f665f636172645f6c65}\paperlink{premise:4d756c74697365742e65715f6f665f6c655f6f665f636172645f6c65}{P155.}] \refstepcounter{paperpremise}\label{premise:4d756c74697365742e65715f6f665f6c655f6f665f636172645f6c65}\textbf{Equals of at most of cardinality at most.} Let \(\alpha\) be a type. Let \(s\) and \(t\) be multisets of elements of \(\alpha\). Assume \(s \le t\). Assume \(\left|t\right| \le \left|s\right|\). Then \(s = t\).
\item[\paperlineanchor{premise:4d756c74697365742e696e6a5f6f6e5f6f665f6e6f6475705f6d6170}\paperlink{premise:4d756c74697365742e696e6a5f6f6e5f6f665f6e6f6475705f6d6170}{P156.}] \refstepcounter{paperpremise}\label{premise:4d756c74697365742e696e6a5f6f6e5f6f665f6e6f6475705f6d6170}\textbf{Injectivity on of without repetitions map.} Let \(\alpha\) and \(\beta\) be types. Let \(f\) be a function of type \(\alpha \to \beta\). Let \(s\) be a multiset of elements of \(\alpha\). Assume \(\operatorname{Multiset\mathord{:}map}\left(f, s\right)\text{ has no repeated elements}\). Let \(x\) be an element of \(\alpha\). Assume \(x \in s\). Let \(y\) be an element of \(\alpha\). Assume \(y \in s\). Assume \(f\left(x\right) = f\left(y\right)\). Then \(x = y\).
\item[\paperlineanchor{premise:4d756c74697365742e7375627365745f6f665f6c65}\paperlink{premise:4d756c74697365742e7375627365745f6f665f6c65}{P157.}] \refstepcounter{paperpremise}\label{premise:4d756c74697365742e7375627365745f6f665f6c65}\textbf{Subset of at most.} Let \(\alpha\) be a type. Let \(s\) and \(t\) be multisets of elements of \(\alpha\). Assume \(s \le t\). Then \(s \subseteq t\).
\item[\paperlineanchor{premise:4e61742e6164645f6173736f63}\paperlink{premise:4e61742e6164645f6173736f63}{P158.}] \refstepcounter{paperpremise}\label{premise:4e61742e6164645f6173736f63}\textbf{Addition associativity.} Let \(n\), \(m\) and \(k\) be natural numbers. Then \(n + m + k = n + (m + k)\).
\item[\paperlineanchor{premise:4e61742e6164645f636f6d6d}\paperlink{premise:4e61742e6164645f636f6d6d}{P159.}] \refstepcounter{paperpremise}\label{premise:4e61742e6164645f636f6d6d}\textbf{Addition commutativity.} Let \(n\) and \(m\) be natural numbers. Then \(n + m = m + n\).
\item[\paperlineanchor{premise:4e61742e6164645f65715f7a65726f}\paperlink{premise:4e61742e6164645f65715f7a65726f}{P160.}] \refstepcounter{paperpremise}\label{premise:4e61742e6164645f65715f7a65726f}\textbf{Sum equals zero.} Let \(m\) and \(n\) be natural numbers. Then \(m + n = 0 \iff m = 0 \land n = 0\).
\item[\paperlineanchor{premise:4e61742e6164645f6c6566745f696e6a}\paperlink{premise:4e61742e6164645f6c6566745f696e6a}{P161.}] \refstepcounter{paperpremise}\label{premise:4e61742e6164645f6c6566745f696e6a}\textbf{Alias.} Let \(m\), \(k\) and \(n\) be natural numbers. Then \(m + n = k + n \iff m = k\).
\item[\paperlineanchor{premise:4e61742e6164645f72696768745f63616e63656c}\paperlink{premise:4e61742e6164645f72696768745f63616e63656c}{P162.}] \refstepcounter{paperpremise}\label{premise:4e61742e6164645f72696768745f63616e63656c}\textbf{Sum right cancel.} Let \(n\), \(m\) and \(k\) be natural numbers. If \(n + m = k + m\), then \(n = k\).
\item[\paperlineanchor{premise:4e61742e6164645f72696768745f636f6d6d}\paperlink{premise:4e61742e6164645f72696768745f636f6d6d}{P163.}] \refstepcounter{paperpremise}\label{premise:4e61742e6164645f72696768745f636f6d6d}\textbf{Interchanging the last two summands.} Let \(n\), \(m\) and \(k\) be natural numbers. Then \(n + m + k = n + k + m\).
\item[\paperlineanchor{premise:4e61742e6164645f7375625f63616e63656c}\paperlink{premise:4e61742e6164645f7375625f63616e63656c}{P164.}] \refstepcounter{paperpremise}\label{premise:4e61742e6164645f7375625f63616e63656c}\textbf{Sum difference cancel.} Let \(n\) and \(m\) be natural numbers. Then \(n + m - m = n\).
\item[\paperlineanchor{premise:4e61742e6164645f7375625f63616e63656c5f6c656674}\paperlink{premise:4e61742e6164645f7375625f63616e63656c5f6c656674}{P165.}] \refstepcounter{paperpremise}\label{premise:4e61742e6164645f7375625f63616e63656c5f6c656674}\textbf{Sum difference cancel left.} Let \(n\) and \(m\) be natural numbers. Then \(n + m - n = m\).
\item[\paperlineanchor{premise:4e61742e6164645f7375625f63616e63656c5f7269676874}\paperlink{premise:4e61742e6164645f7375625f63616e63656c5f7269676874}{P166.}] \refstepcounter{paperpremise}\label{premise:4e61742e6164645f7375625f63616e63656c5f7269676874}\textbf{Sum difference cancel right.} Let \(n\) and \(m\) be natural numbers. Then \(n + m - m = n\).
\item[\paperlineanchor{premise:4e61742e6164645f7a65726f}\paperlink{premise:4e61742e6164645f7a65726f}{P167.}] \refstepcounter{paperpremise}\label{premise:4e61742e6164645f7a65726f}\textbf{Adding zero on the right.} Let \(n\) be a natural number. Then \(n + 0 = n\).
\item[\paperlineanchor{premise:4e61742e696e73744e655a65726f53756363}\paperlink{premise:4e61742e696e73744e655a65726f53756363}{P168.}] \refstepcounter{paperpremise}\label{premise:4e61742e696e73744e655a65726f53756363}\textbf{Inst differs from zero successor.} Let \(n\) be a natural number. Then \(\operatorname{Ne\,Zero}\left(n + 1\right)\).
\item[\paperlineanchor{premise:4e61742e6c655f6164645f7269676874}\paperlink{premise:4e61742e6c655f6164645f7269676874}{P169.}] \refstepcounter{paperpremise}\label{premise:4e61742e6c655f6164645f7269676874}\textbf{Adding a natural number can only increase a sum.} Let \(n\) and \(k\) be natural numbers. Then \(n \le n + k\).
\item[\paperlineanchor{premise:4e61742e6c655f6f665f6e6f745f6c74}\paperlink{premise:4e61742e6c655f6f665f6e6f745f6c74}{P170.}] \refstepcounter{paperpremise}\label{premise:4e61742e6c655f6f665f6e6f745f6c74}\textbf{At most of not less than.} Let \(a\) and \(b\) be natural numbers. Assume \(\neg \left(a < b\right)\). Then \(b \le a\).
\item[\paperlineanchor{premise:4e61742e6c655f73756363}\paperlink{premise:4e61742e6c655f73756363}{P171.}] \refstepcounter{paperpremise}\label{premise:4e61742e6c655f73756363}\textbf{At most successor.} Let \(n\) be a natural number. Then \(n \le n + 1\).
\item[\paperlineanchor{premise:4e61742e6c745f6966665f6164645f6f6e655f6c65}\paperlink{premise:4e61742e6c745f6966665f6164645f6f6e655f6c65}{P172.}] \refstepcounter{paperpremise}\label{premise:4e61742e6c745f6966665f6164645f6f6e655f6c65}\textbf{Less than if and only if sum one at most.} Let \(m\) and \(n\) be natural numbers. Then \(m < n \iff m + 1 \le n\).
\item[\paperlineanchor{premise:4e61742e6c745f6f665f6e6f745f6c65}\paperlink{premise:4e61742e6c745f6f665f6e6f745f6c65}{P173.}] \refstepcounter{paperpremise}\label{premise:4e61742e6c745f6f665f6e6f745f6c65}\textbf{Less than of not at most.} Let \(a\) and \(b\) be natural numbers. Assume \(\neg \left(a \le b\right)\). Then \(b < a\).
\item[\paperlineanchor{premise:4e61742e6c745f6f725f67745f6f665f6e65}\paperlink{premise:4e61742e6c745f6f725f67745f6f665f6e65}{P174.}] \refstepcounter{paperpremise}\label{premise:4e61742e6c745f6f725f67745f6f665f6e65}\textbf{Less than or greater than of differs from.} Let \(a\) and \(b\) be natural numbers. Assume \(a \ne b\). Then \(a < b \lor a > b\).
\item[\paperlineanchor{premise:4e61742e706f735f6966665f6e655f7a65726f}\paperlink{premise:4e61742e706f735f6966665f6e655f7a65726f}{P175.}] \refstepcounter{paperpremise}\label{premise:4e61742e706f735f6966665f6e655f7a65726f}\textbf{Positivity if and only if differs from zero.} Let \(n\) be a natural number. Then \(0 < n \iff n \ne 0\).
\item[\paperlineanchor{premise:4e61742e73656c665f65715f6164645f6c656674}\paperlink{premise:4e61742e73656c665f65715f6164645f6c656674}{P176.}] \refstepcounter{paperpremise}\label{premise:4e61742e73656c665f65715f6164645f6c656674}\textbf{Self equals sum left.} Let \(a\) and \(b\) be natural numbers. Then \(a = b + a \iff b = 0\).
\item[\paperlineanchor{premise:4e61742e73656c665f65715f6164645f7269676874}\paperlink{premise:4e61742e73656c665f65715f6164645f7269676874}{P177.}] \refstepcounter{paperpremise}\label{premise:4e61742e73656c665f65715f6164645f7269676874}\textbf{Self equals sum right.} Let \(a\) and \(b\) be natural numbers. Then \(a = a + b \iff b = 0\).
\item[\paperlineanchor{premise:4e61742e7375625f65715f6966665f65715f616464}\paperlink{premise:4e61742e7375625f65715f6966665f65715f616464}{P178.}] \refstepcounter{paperpremise}\label{premise:4e61742e7375625f65715f6966665f65715f616464}\textbf{Difference equals if and only if equals sum.} Let \(b\), \(a\), and \(c\) be natural numbers. Assume \(b \le a\). Then \(a - b = c \iff a = c + b\).
\item[\paperlineanchor{premise:4e61742e7375625f6c65}\paperlink{premise:4e61742e7375625f6c65}{P179.}] \refstepcounter{paperpremise}\label{premise:4e61742e7375625f6c65}\textbf{Difference at most.} Let \(n\) and \(m\) be natural numbers. Then \(n - m \le n\).
\item[\paperlineanchor{premise:4e61742e7375625f6c655f6966665f6c655f616464}\paperlink{premise:4e61742e7375625f6c655f6966665f6c655f616464}{P180.}] \refstepcounter{paperpremise}\label{premise:4e61742e7375625f6c655f6966665f6c655f616464}\textbf{Difference at most if and only if at most sum.} Let \(a\), \(b\), and \(c\) be natural numbers. Then \(a - b \le c \iff a \le c + b\).
\item[\paperlineanchor{premise:4e61742e7375625f6c655f7375625f6c656674}\paperlink{premise:4e61742e7375625f6c655f7375625f6c656674}{P181.}] \refstepcounter{paperpremise}\label{premise:4e61742e7375625f6c655f7375625f6c656674}\textbf{Difference at most difference left.} Let \(n\) and \(m\) be natural numbers. Assume \(n \le m\). Let \(k\) be a natural number. Then \(k - m \le k - n\).
\item[\paperlineanchor{premise:4e61742e737563635f6c655f6f665f6c74}\paperlink{premise:4e61742e737563635f6c655f6f665f6c74}{P182.}] \refstepcounter{paperpremise}\label{premise:4e61742e737563635f6c655f6f665f6c74}\textbf{Successor at most of less than.} Let \(n\) and \(m\) be natural numbers. If \(n < m\), then \(n + 1 \le m\).
\item[\paperlineanchor{premise:4e61742e737563635f6c655f73756363}\paperlink{premise:4e61742e737563635f6c655f73756363}{P183.}] \refstepcounter{paperpremise}\label{premise:4e61742e737563635f6c655f73756363}\textbf{Successors preserve order.} Let \(n\) and \(m\) be natural numbers. If \(n \le m\), then \(n + 1 \le m + 1\).
\item[\paperlineanchor{premise:4e61742e737563635f6e655f73656c66}\paperlink{premise:4e61742e737563635f6e655f73656c66}{P184.}] \refstepcounter{paperpremise}\label{premise:4e61742e737563635f6e655f73656c66}\textbf{Successor differs from self.} Let \(n\) be a natural number. Then \(n + 1 \ne n\).
\item[\paperlineanchor{premise:4e61742e74776f5f6d756c}\paperlink{premise:4e61742e74776f5f6d756c}{P185.}] \refstepcounter{paperpremise}\label{premise:4e61742e74776f5f6d756c}\textbf{Doubling by addition.} Let \(n\) be a natural number. Then \(2 \cdot n = n + n\).
\item[\paperlineanchor{premise:4e61742e7a65726f5f6c65}\paperlink{premise:4e61742e7a65726f5f6c65}{P186.}] \refstepcounter{paperpremise}\label{premise:4e61742e7a65726f5f6c65}\textbf{Natural numbers are nonnegative.} Let \(n\) be a natural number. Then \(0 \le n\).
\item[\paperlineanchor{premise:4e61742e7a65726f5f6c745f73756363}\paperlink{premise:4e61742e7a65726f5f6c745f73756363}{P187.}] \refstepcounter{paperpremise}\label{premise:4e61742e7a65726f5f6c745f73756363}\textbf{Zero less than successor.} Let \(n\) be a natural number. Then \(0 < n + 1\).
\item[\paperlineanchor{premise:5365742e5375726a4f6e2e696d6167655f65715f6f665f6d617073546f}\paperlink{premise:5365742e5375726a4f6e2e696d6167655f65715f6f665f6d617073546f}{P188.}] \refstepcounter{paperpremise}\label{premise:5365742e5375726a4f6e2e696d6167655f65715f6f665f6d617073546f}\textbf{Image equals of maps to.} Let \(\alpha\) and \(\beta\) be types. Let \(s\) be a subset of \(\alpha\). Let \(t\) be a subset of \(\beta\). Let \(f\) be a function of type \(\alpha \to \beta\). Assume \(t\subseteq f\left[s\right]\). Assume \(\operatorname{Maps\,To}\left(f, s, t\right)\). Then \(f\left[s\right] = t\).
\item[\paperlineanchor{premise:5365742e696d6167655f636f6e6772}\paperlink{premise:5365742e696d6167655f636f6e6772}{P189.}] \refstepcounter{paperpremise}\label{premise:5365742e696d6167655f636f6e6772}\textbf{Image congruence.} Let \(\alpha\) and \(\beta\) be types. Let \(f\) and \(g\) be functions of type \(\alpha \to \beta\). Let \(s\) be a subset of \(\alpha\). Assume \(\forall a\, :\, \alpha,\; a \in s \implies f\left(a\right) = g\left(a\right)\). Then \(f\left[s\right] = g\left[s\right]\).
\item[\paperlineanchor{premise:5365742e696d6167655f756e6976}\paperlink{premise:5365742e696d6167655f756e6976}{P190.}] \refstepcounter{paperpremise}\label{premise:5365742e696d6167655f756e6976}\textbf{Image whole space.} Let \(\alpha\) and \(\beta\) be types. Let \(f\) be a function of type \(\alpha \to \beta\). Then \(f\left[\alpha\right] = \operatorname{Set\mathord{:}range}\left(f\right)\).
\item[\paperlineanchor{premise:5365742e696e737449735265666c537562736574}\paperlink{premise:5365742e696e737449735265666c537562736574}{P191.}] \refstepcounter{paperpremise}\label{premise:5365742e696e737449735265666c537562736574}\textbf{Inst is refl subset.} Let \(\alpha\) be a type. Then \((\mathit{x1} \mapsto \mathit{x2} \mapsto \mathit{x1} \subseteq \mathit{x2})\text{ is reflexive on }\mathcal{P}(\alpha)\).
\item[\paperlineanchor{premise:5365742e6d656d5f72616e6765}\paperlink{premise:5365742e6d656d5f72616e6765}{P192.}] \refstepcounter{paperpremise}\label{premise:5365742e6d656d5f72616e6765}\textbf{Membership range.} Let \(\alpha\) be a type. Let \(\iota\) be a sort. Let \(f\) be a function of type \(\iota \to \alpha\). Let \(x\) be an element of \(\alpha\). Then \(x \in \operatorname{Set\mathord{:}range}(f) \iff (\exists y\, :\, \iota,\; f(y) = x)\).
\item[\paperlineanchor{premise:5365742e6d656d5f756e6976}\paperlink{premise:5365742e6d656d5f756e6976}{P193.}] \refstepcounter{paperpremise}\label{premise:5365742e6d656d5f756e6976}\textbf{Membership whole space.} Let \(\alpha\) be a type. Let \(x\) be an element of \(\alpha\). Then \(x \in \alpha\).
\item[\paperlineanchor{premise:5365742e6e6f6e7472697669616c5f636f655f736f7274}\paperlink{premise:5365742e6e6f6e7472697669616c5f636f655f736f7274}{P194.}] \refstepcounter{paperpremise}\label{premise:5365742e6e6f6e7472697669616c5f636f655f736f7274}\textbf{Nontrivial coercion sort.} Let \(\alpha\) be a type. Let \(s\) be a subset of \(\alpha\). Then \(\operatorname{Nontrivial}\left(s\right) \iff s\text{ contains two distinct elements}\).
\item[\paperlineanchor{premise:5365742e73756273696e676c65746f6e5f636f65}\paperlink{premise:5365742e73756273696e676c65746f6e5f636f65}{P195.}] \refstepcounter{paperpremise}\label{premise:5365742e73756273696e676c65746f6e5f636f65}\textbf{Subsingleton coercion.} Let \(\alpha\) be a type. Let \(s\) be a subset of \(\alpha\). Then \(s\text{ has at most one element} \iff \operatorname{Set\mathord{:}Subsingleton}\left(s\right)\).
\item[\paperlineanchor{premise:53756272656c6174696f6e2e7766}\paperlink{premise:53756272656c6174696f6e2e7766}{P196.}] \refstepcounter{paperpremise}\label{premise:53756272656c6174696f6e2e7766}\textbf{Wf.} Let \(\alpha\) be a sort. Let \(r\) and \(q\) be relations from \(\alpha\) to \(\alpha\). Assume \(\operatorname{Subrelation}\left(q, r\right)\). Assume \(\operatorname{Well\,Founded}\left(r\right)\). Then \(\operatorname{Well\,Founded}\left(q\right)\).
\item[\paperlineanchor{premise:53756273696e676c65746f6e2e656c696d}\paperlink{premise:53756273696e676c65746f6e2e656c696d}{P197.}] \refstepcounter{paperpremise}\label{premise:53756273696e676c65746f6e2e656c696d}\textbf{Elim.} Let \(\alpha\) be a sort. Assume \(\alpha\text{ has at most one element}\). Let \(a\) and \(b\) be elements of \(\alpha\). Then \(a = b\).
\item[\paperlineanchor{premise:537562747970652e6571}\paperlink{premise:537562747970652e6571}{P198.}] \refstepcounter{paperpremise}\label{premise:537562747970652e6571}\textbf{Equals.} Let \(\alpha\) be a type. Let \(p\) be a function of type \(\alpha \to \mathsf{Prop}\). Let \(\mathit{a1}\) be an element of \(\{x\in \alpha \mid p\left(x\right)\}\). Let \(\mathit{a2}\) be an element of \(\{x\in \alpha \mid p\left(x\right)\}\). Assume \(\mathit{a1} = \mathit{a2}\). Then \(\mathit{a1} = \mathit{a2}\).
\item[\paperlineanchor{premise:537562747970652e657869737473}\paperlink{premise:537562747970652e657869737473}{P199.}] \refstepcounter{paperpremise}\label{premise:537562747970652e657869737473}\textbf{Exists.} Let \(\alpha\) be a sort. Let \(p\) be a function of type \(\alpha \to \mathsf{Prop}\). Let \(q\) be a function of type \(\{a\in \alpha \mid p\left(a\right)\} \to \mathsf{Prop}\). Then \((\exists x\, :\, \{a\in \alpha \mid p(a)\},\; q(x)) \iff (\exists a\, :\, \alpha,\; \exists b\text{ proving }p(a),\; q((a\, :\, \{a_{1}\in \alpha \mid p(a_{1})\})))\).
\item[\paperlineanchor{premise:537562747970652e657874}\paperlink{premise:537562747970652e657874}{P200.}] \refstepcounter{paperpremise}\label{premise:537562747970652e657874}\textbf{Ext.} Let \(\alpha\) be a sort. Let \(p\) be a function of type \(\alpha \to \mathsf{Prop}\). Let \(\mathit{a1}\) be an element of \(\{x\in \alpha \mid p\left(x\right)\}\). Let \(\mathit{a2}\) be an element of \(\{x\in \alpha \mid p\left(x\right)\}\). Assume \(\mathit{a1} = \mathit{a2}\). Then \(\mathit{a1} = \mathit{a2}\).
\item[\paperlineanchor{premise:537562747970652e6578745f696666}\paperlink{premise:537562747970652e6578745f696666}{P201.}] \refstepcounter{paperpremise}\label{premise:537562747970652e6578745f696666}\textbf{Ext if and only if.} Let \(\alpha\) be a sort. Let \(p\) be a function of type \(\alpha \to \mathsf{Prop}\). Let \(\mathit{a1}\) be an element of \(\{x\in \alpha \mid p\left(x\right)\}\). Let \(\mathit{a2}\) be an element of \(\{x\in \alpha \mid p\left(x\right)\}\). Then \(\mathit{a1} = \mathit{a2} \iff \mathit{a1} = \mathit{a2}\).
\item[\paperlineanchor{premise:537562747970652e6d6b2e696e6a4571}\paperlink{premise:537562747970652e6d6b2e696e6a4571}{P202.}] \refstepcounter{paperpremise}\label{premise:537562747970652e6d6b2e696e6a4571}\textbf{Injectivity equals.} Let \(\alpha\) be a sort. Let \(p\) be a function of type \(\alpha \to \mathsf{Prop}\). Let \(\mathit{val}\) be an element of \(\alpha\). Let \(\mathit{property}\) be a proof of \(p\left(\mathit{val}\right)\). Let \(\mathit{val}_{1}\) be an element of \(\alpha\). Let \(\mathit{property}_{1}\) be a proof of \(p\left(\mathit{val}_{1}\right)\). Then \(((\mathit{val}\, :\, \{x\in \alpha \mid p(x)\}) = (\mathit{val}_{1}\, :\, \{x\in \alpha \mid p(x)\})) = (\mathit{val} = \mathit{val}_{1})\).
\item[\paperlineanchor{premise:537562747970652e70726f7065727479}\paperlink{premise:537562747970652e70726f7065727479}{P203.}] \refstepcounter{paperpremise}\label{premise:537562747970652e70726f7065727479}\textbf{Property.} Let \(\alpha\) be a sort. Let \(p\) be a function of type \(\alpha \to \mathsf{Prop}\). Let \(\mathit{self}\) be an element of \(\{x\in \alpha \mid p(x)\}\). Then \(p\left(\mathit{self}\right)\).
\item[\paperlineanchor{premise:57656c6c466f756e64656452656c6174696f6e2e7766}\paperlink{premise:57656c6c466f756e64656452656c6174696f6e2e7766}{P204.}] \refstepcounter{paperpremise}\label{premise:57656c6c466f756e64656452656c6174696f6e2e7766}\textbf{Wf.} Let \(\alpha\) be a sort. Fix a structure \(\mathit{self}\) of type \(\operatorname{Well\,Founded\,Relation}\left(\alpha\right)\). Then \(\operatorname{Well\,Founded}\left(\operatorname{rel}\right)\).
\item[\paperlineanchor{premise:616e645f66616c7365}\paperlink{premise:616e645f66616c7365}{P205.}] \refstepcounter{paperpremise}\label{premise:616e645f66616c7365}\textbf{And false.} Let \(p\) be a proposition. Then \(\left(p \land \mathrm{False}\right) = \mathrm{False}\).
\item[\paperlineanchor{premise:616e645f696d70}\paperlink{premise:616e645f696d70}{P206.}] \refstepcounter{paperpremise}\label{premise:616e645f696d70}\textbf{And imp.} Let \(a\), \(b\), and \(c\) be propositions. Then \(a \land b \implies c \iff a \implies b \implies c\).
\item[\paperlineanchor{premise:616e645f73656c66}\paperlink{premise:616e645f73656c66}{P207.}] \refstepcounter{paperpremise}\label{premise:616e645f73656c66}\textbf{And self.} Let \(p\) be a proposition. Then \(\left(p \land p\right) = p\).
\item[\paperlineanchor{premise:636f6e67725f617267}\paperlink{premise:636f6e67725f617267}{P208.}] \refstepcounter{paperpremise}\label{premise:636f6e67725f617267}\textbf{Alias.} Let \(\alpha\) be a sort. Let \(\beta\) be a sort. Let \(a_{1}\) and \(a_{2}\) be elements of \(\alpha\). Let \(f\) be a function of type \(\alpha \to \beta\). Assume \(a_{1} = a_{2}\). Then \(f\left(a_{1}\right) = f\left(a_{2}\right)\).
\item[\paperlineanchor{premise:646974655f636f6e645f65715f74727565}\paperlink{premise:646974655f636f6e645f65715f74727565}{P209.}] \refstepcounter{paperpremise}\label{premise:646974655f636f6e645f65715f74727565}\textbf{Dite cond equals true.} Let \(\alpha\) be a sort. Let \(c\) be a proposition. Let \(x\) be an element of \(\operatorname{Decidable}\left(c\right)\). Let \(t\) be a function of type \(c \to \alpha\). Let \(e\) be a function of type \(\neg c \to \alpha\). Let \(h\) be a proof of \(c = \mathrm{True}\). Then \(\begin{cases}t(h_{1}) &\text{given } h_{1}\text{ proving } c\\ e(h_{1}) &\text{given } h_{1}\text{ proving } \neg c\end{cases} = t(\operatorname{transport}^{\leftarrow}_{h}(\mathrm{trivial}))\).
\item[\paperlineanchor{premise:646974655f636f6e6772}\paperlink{premise:646974655f636f6e6772}{P210.}] \refstepcounter{paperpremise}\label{premise:646974655f636f6e6772}\textbf{Dite congruence.} Let \(b\) and \(c\) be propositions. Let \(\alpha\) be a sort. Let \(x\) be an element of \(\operatorname{Decidable}\left(b\right)\). Assume \(c\) is decidable. Let \(x_{1}\) be a function of type \(b \to \alpha\). Let \(u\) be a function of type \(c \to \alpha\). Let \(y\) be a function of type \(\neg b \to \alpha\). Let \(v\) be a function of type \(\neg c \to \alpha\). Let \(h_{1}\) be a proof of \(b = c\). Assume \(\forall h\, :\, c,\; x_{1}\left(\operatorname{transport}^{\leftarrow}_{h_{1}}(h)\right) = u\left(h\right)\). Assume \(\forall h\, :\, \neg c,\; y\left(\operatorname{transport}^{\leftarrow}_{h_{1}}(h)\right) = v\left(h\right)\). Then \(\begin{cases}x_{1}(h) &\text{given } h\text{ proving } b\\ y(h) &\text{given } h\text{ proving } \neg b\end{cases} = \begin{cases}u(h) &\text{given } h\text{ proving } c\\ v(h) &\text{given } h\text{ proving } \neg c\end{cases}\).
\item[\paperlineanchor{premise:65715f636f6d6d}\paperlink{premise:65715f636f6d6d}{P211.}] \refstepcounter{paperpremise}\label{premise:65715f636f6d6d}\textbf{Equals commutation.} Let \(\alpha\) be a sort. Let \(a\) and \(b\) be elements of \(\alpha\). Then \(a = b \iff b = a\).
\item[\paperlineanchor{premise:65715f66616c7365}\paperlink{premise:65715f66616c7365}{P212.}] \refstepcounter{paperpremise}\label{premise:65715f66616c7365}\textbf{Equals false.} Let \(p\) be a proposition. Assume \(\neg p\). Then \(p = \mathrm{False}\).
\item[\paperlineanchor{premise:65715f66616c73655f6f665f646563696465}\paperlink{premise:65715f66616c73655f6f665f646563696465}{P213.}] \refstepcounter{paperpremise}\label{premise:65715f66616c73655f6f665f646563696465}\textbf{Equals false of decide.} Let \(p\) be a proposition. Let \(x\) be an element of \(\operatorname{Decidable}\left(p\right)\). Assume \(\operatorname{decide}\left(p\right) = \mathrm{false}\). Then \(p = \mathrm{False}\).
\item[\paperlineanchor{premise:65715f6f725f6e65}\paperlink{premise:65715f6f725f6e65}{P214.}] \refstepcounter{paperpremise}\label{premise:65715f6f725f6e65}\textbf{Equals or differs from.} Let \(\alpha\) be a sort. Let \(x\) and \(y\) be elements of \(\alpha\). Then \(x = y \lor x \ne y\).
\item[\paperlineanchor{premise:65715f73656c66}\paperlink{premise:65715f73656c66}{P215.}] \refstepcounter{paperpremise}\label{premise:65715f73656c66}\textbf{Equals self.} Let \(\alpha\) be a sort. Let \(a\) be an element of \(\alpha\). Then \(\left(a = a\right) = \mathrm{True}\).
\item[\paperlineanchor{premise:65715f74727565}\paperlink{premise:65715f74727565}{P216.}] \refstepcounter{paperpremise}\label{premise:65715f74727565}\textbf{Equals true.} Let \(p\) be a proposition. Assume \(p\). Then \(p = \mathrm{True}\).
\item[\paperlineanchor{premise:6578697374735f616e645f6c656674}\paperlink{premise:6578697374735f616e645f6c656674}{P217.}] \refstepcounter{paperpremise}\label{premise:6578697374735f616e645f6c656674}\textbf{Exists and left.} Let \(\alpha\) be a sort. Let \(p\) be a function of type \(\alpha \to \mathsf{Prop}\). Let \(b\) be a proposition. Then \((\exists x\, :\, \alpha,\; b \land p(x)) \iff b \land (\exists x\, :\, \alpha,\; p(x))\).
\item[\paperlineanchor{premise:6578697374735f636f6d6d}\paperlink{premise:6578697374735f636f6d6d}{P218.}] \refstepcounter{paperpremise}\label{premise:6578697374735f636f6d6d}\textbf{Exists commutation.} Let \(\alpha\) be a sort. Let \(\beta\) be a sort. Let \(p\) be a relation from \(\alpha\) to \(\beta\). Then \((\exists a\, :\, \alpha,\; \exists b\, :\, \beta,\; p(a,\allowbreak{} b)) \iff (\exists b\, :\, \beta,\; \exists a\, :\, \alpha,\; p(a,\allowbreak{} b))\).
\item[\paperlineanchor{premise:6578697374735f636f6e7374}\paperlink{premise:6578697374735f636f6e7374}{P219.}] \refstepcounter{paperpremise}\label{premise:6578697374735f636f6e7374}\textbf{Exists constant.} Let \(b\) be a proposition. Let \(\alpha\) be a sort. Assume \(\alpha\text{ is nonempty}\). Then \(\left(\exists x\, :\, \alpha,\; b\right) \iff b\).
\item[\paperlineanchor{premise:6578697374735f66616c7365}\paperlink{premise:6578697374735f66616c7365}{P220.}] \refstepcounter{paperpremise}\label{premise:6578697374735f66616c7365}\textbf{Exists false.} Let \(\alpha\) be a sort. Then \(\neg \left(\exists x\, :\, \alpha,\; \mathrm{False}\right)\).
\item[\paperlineanchor{premise:6578697374735f70726f70}\paperlink{premise:6578697374735f70726f70}{P221.}] \refstepcounter{paperpremise}\label{premise:6578697374735f70726f70}\textbf{Exists prop.} Let \(b\) and \(a\) be propositions. Then \(a \land b \iff a \land b\).
\item[\paperlineanchor{premise:6578697374735f70726f705f636f6e6772}\paperlink{premise:6578697374735f70726f705f636f6e6772}{P222.}] \refstepcounter{paperpremise}\label{premise:6578697374735f70726f705f636f6e6772}\textbf{Exists prop congruence.} Let \(p\) and \(p^{'}\) be propositions. Let \(q\) be a function of type \(p \to \mathsf{Prop}\). Let \(q^{'}\) be a function of type \(p \to \mathsf{Prop}\). Assume \(\forall h\, :\, p,\; q\left(h\right) \iff q^{'}\left(h\right)\). Let \(\mathit{hp}\) be a proof of \(p \iff p^{'}\). Then \((\exists x\, :\, p,\; q(x)) \iff (\exists h\text{ proving }p^{'},\; q^{'}(\operatorname{Iff\mathord{:}mpr}(\mathit{hp},\allowbreak{} h)))\).
\item[\paperlineanchor{premise:66616c73655f616e64}\paperlink{premise:66616c73655f616e64}{P223.}] \refstepcounter{paperpremise}\label{premise:66616c73655f616e64}\textbf{False and.} Let \(p\) be a proposition. Then \(\left(\mathrm{False} \land p\right) = \mathrm{False}\).
\item[\paperlineanchor{premise:666f72616c6c5f6170706c795f65715f696d705f696666e28282}\paperlink{premise:666f72616c6c5f6170706c795f65715f696d705f696666e28282}{P224.}] \refstepcounter{paperpremise}\label{premise:666f72616c6c5f6170706c795f65715f696d705f696666e28282}\textbf{Forall apply equals imp iff2.} Let \(\alpha\) be a sort. Let \(\beta\) be a sort. Let \(f\) be a function of type \(\alpha \to \beta\). Let \(p\) be a function of type \(\alpha \to \mathsf{Prop}\). Let \(q\) be a function of type \(\beta \to \mathsf{Prop}\). Then \((\forall b\, :\, \beta,\; \forall a\, :\, \alpha,\; p(a) \implies f(a) = b \implies q(b)) \iff (\forall a\, :\, \alpha,\; p(a) \implies q(f(a)))\).
\item[\paperlineanchor{premise:666f72616c6c5f636f6e6772}\paperlink{premise:666f72616c6c5f636f6e6772}{P225.}] \refstepcounter{paperpremise}\label{premise:666f72616c6c5f636f6e6772}\textbf{Forall congruence.} Let \(\alpha\) be a sort. Let \(p\) and \(q\) be functions of type \(\alpha \to \mathsf{Prop}\). Assume \(\forall a\, :\, \alpha,\; p\left(a\right) = q\left(a\right)\). Then \((\forall a\, :\, \alpha,\; p(a)) = (\forall a\, :\, \alpha,\; q(a))\).
\item[\paperlineanchor{premise:666f72616c6c5f6571}\paperlink{premise:666f72616c6c5f6571}{P226.}] \refstepcounter{paperpremise}\label{premise:666f72616c6c5f6571}\textbf{Forall equals.} Let \(\alpha\) be a sort. Let \(p\) be a function of type \(\alpha \to \mathsf{Prop}\). Let \(a^{'}\) be an element of \(\alpha\). Then \(\left(\forall a\, :\, \alpha,\; a = a^{'} \implies p\left(a\right)\right) \iff p\left(a^{'}\right)\).
\item[\paperlineanchor{premise:666f72616c6c5f6578697374735f696e646578}\paperlink{premise:666f72616c6c5f6578697374735f696e646578}{P227.}] \refstepcounter{paperpremise}\label{premise:666f72616c6c5f6578697374735f696e646578}\textbf{Forall exists index.} Let \(\alpha\) be a sort. Let \(p\) be a function of type \(\alpha \to \mathsf{Prop}\). Let \(q\) be a function of type \(\left(\exists x\, :\, \alpha,\; p\left(x\right)\right) \to \mathsf{Prop}\). Then \((\forall h\, :\, \exists x\, :\, \alpha,\; p(x),\; q(h)) \iff (\forall x\, :\, \alpha,\; \forall h\, :\, p(x),\; q(\operatorname{Exists\mathord{:}intro}(x,\allowbreak{} h)))\).
\item[\paperlineanchor{premise:66756e657874}\paperlink{premise:66756e657874}{P228.}] \refstepcounter{paperpremise}\label{premise:66756e657874}\textbf{Function extensionality.} Let \(\alpha\) be a sort. Let \(\beta\) be a function of type \(\alpha \to \mathsf{Sort}\). Let \(f\) be a function of type \(\Pi_{x\, :\, \alpha}\, \beta\left(x\right)\). Let \(g\) be a function of type \(\Pi_{x\, :\, \alpha}\, \beta\left(x\right)\). Assume \(\forall x\, :\, \alpha,\; f\left(x\right) = g\left(x\right)\). Then \(f = g\).
\item[\paperlineanchor{premise:67655f6f665f6571}\paperlink{premise:67655f6f665f6571}{P229.}] \refstepcounter{paperpremise}\label{premise:67655f6f665f6571}\textbf{At least of equals.} Let \(\alpha\) be a type. Equip \(\alpha\) with a preorder. Let \(a\) and \(b\) be elements of \(\alpha\). Assume \(a = b\). Then \(b \le a\).
\item[\paperlineanchor{premise:69665f6e6567}\paperlink{premise:69665f6e6567}{P230.}] \refstepcounter{paperpremise}\label{premise:69665f6e6567}\textbf{If neg.} Let \(c\) be a proposition. Let \(h\) be an element of \(\operatorname{Decidable}\left(c\right)\). Assume \(\neg c\). Let \(\alpha\) be a sort. Let \(t\) and \(e\) be elements of \(\alpha\). Then \(\begin{cases}t &\text{if } c\\ e &\text{otherwise}\end{cases} = e\).
\item[\paperlineanchor{premise:69665f706f73}\paperlink{premise:69665f706f73}{P231.}] \refstepcounter{paperpremise}\label{premise:69665f706f73}\textbf{If positivity.} Let \(c\) be a proposition. Let \(h\) be an element of \(\operatorname{Decidable}\left(c\right)\). Assume \(c\). Let \(\alpha\) be a sort. Let \(t\) and \(e\) be elements of \(\alpha\). Then \(\begin{cases}t &\text{if } c\\ e &\text{otherwise}\end{cases} = t\).
\item[\paperlineanchor{premise:6966665f73656c66}\paperlink{premise:6966665f73656c66}{P232.}] \refstepcounter{paperpremise}\label{premise:6966665f73656c66}\textbf{If and only if self.} Let \(p\) be a proposition. Then \(\left(p \iff p\right) = \mathrm{True}\).
\item[\paperlineanchor{premise:696d705f66616c7365}\paperlink{premise:696d705f66616c7365}{P233.}] \refstepcounter{paperpremise}\label{premise:696d705f66616c7365}\textbf{Imp false.} Let \(a\) be a proposition. Then \(a \implies \mathrm{False} \iff \neg a\).
\item[\paperlineanchor{premise:696d705f73656c66}\paperlink{premise:696d705f73656c66}{P234.}] \refstepcounter{paperpremise}\label{premise:696d705f73656c66}\textbf{Imp self.} Let \(a\) be a proposition. Then \(a \implies a \iff \mathrm{True}\).
\item[\paperlineanchor{premise:696d706c6965735f636f6e6772}\paperlink{premise:696d706c6965735f636f6e6772}{P235.}] \refstepcounter{paperpremise}\label{premise:696d706c6965735f636f6e6772}\textbf{Implies congruence.} Let \(p_{1}\) and \(p_{2}\) be sorts. Let \(q_{1}\) and \(q_{2}\) be sorts. Assume \(p_{1} = p_{2}\). Assume \(q_{1} = q_{2}\). Then \(\left(p_{1} \to q_{1}\right) = \left(p_{2} \to q_{2}\right)\).
\item[\paperlineanchor{premise:696d706c6965735f74727565}\paperlink{premise:696d706c6965735f74727565}{P236.}] \refstepcounter{paperpremise}\label{premise:696d706c6965735f74727565}\textbf{Implies true.} Let \(\alpha\) be a sort. Then \(\left(\forall a\, :\, \alpha,\; \mathrm{True}\right) = \mathrm{True}\).
\item[\paperlineanchor{premise:696e73744973456d70747946616c7365}\paperlink{premise:696e73744973456d70747946616c7365}{P237.}] \refstepcounter{paperpremise}\label{premise:696e73744973456d70747946616c7365}\textbf{Inst is empty false.} \(\mathrm{False}\text{ is empty}\).
\item[\paperlineanchor{premise:696e73744e6f6e656d7074794f66496e68616269746564}\paperlink{premise:696e73744e6f6e656d7074794f66496e68616269746564}{P238.}] \refstepcounter{paperpremise}\label{premise:696e73744e6f6e656d7074794f66496e68616269746564}\textbf{Inst nonempty of inhabited.} Let \(\alpha\) be a sort. Fix a structure \(\mathit{inst}\) of type \(\operatorname{Inhabited}\left(\alpha\right)\). Then \(\alpha\text{ is nonempty}\).
\item[\paperlineanchor{premise:6974655f636f6e6772}\paperlink{premise:6974655f636f6e6772}{P239.}] \refstepcounter{paperpremise}\label{premise:6974655f636f6e6772}\textbf{Ite congruence.} Let \(\alpha\) be a sort. Let \(b\) and \(c\) be propositions. Let \(x\), \(y\), \(u\), and \(v\) be elements of \(\alpha\). Let \(s\) be an element of \(\operatorname{Decidable}\left(b\right)\). Assume \(c\) is decidable. Assume \(b = c\). Assume \(c \implies x = u\). Assume \(\neg c \implies y = v\). Then \(\begin{cases}x &\text{if } b\\ y &\text{otherwise}\end{cases} = \begin{cases}u &\text{if } c\\ v &\text{otherwise}\end{cases}\).
\item[\paperlineanchor{premise:6974655f65715f6c6566745f696666}\paperlink{premise:6974655f65715f6c6566745f696666}{P240.}] \refstepcounter{paperpremise}\label{premise:6974655f65715f6c6566745f696666}\textbf{Ite equals left if and only if.} Let \(\alpha\) be a sort. Let \(p\) be a proposition. Assume \(p\) is decidable. Let \(x\) and \(y\) be elements of \(\alpha\). Then \(\begin{cases}x &\text{if } p\\ y &\text{otherwise}\end{cases} = x \iff \neg p \implies y = x\).
\item[\paperlineanchor{premise:6974655f65715f72696768745f696666}\paperlink{premise:6974655f65715f72696768745f696666}{P241.}] \refstepcounter{paperpremise}\label{premise:6974655f65715f72696768745f696666}\textbf{Ite equals right if and only if.} Let \(\alpha\) be a sort. Let \(p\) be a proposition. Assume \(p\) is decidable. Let \(x\) and \(y\) be elements of \(\alpha\). Then \(\begin{cases}x &\text{if } p\\ y &\text{otherwise}\end{cases} = y \iff p \implies x = y\).
\item[\paperlineanchor{premise:6c655f6f665f6c655f6f665f6571}\paperlink{premise:6c655f6f665f6c655f6f665f6571}{P242.}] \refstepcounter{paperpremise}\label{premise:6c655f6f665f6c655f6f665f6571}\textbf{At most of at most of equals.} Let \(\alpha\) be a type. Let \(a\), \(b\), and \(c\) be elements of \(\alpha\). Fix a structure \(\mathit{inst}\) of type \(\operatorname{LE}\left(\alpha\right)\). Assume \(a \le b\). Assume \(b = c\). Then \(a \le c\).
\item[\paperlineanchor{premise:6c655f7265666c}\paperlink{premise:6c655f7265666c}{P243.}] \refstepcounter{paperpremise}\label{premise:6c655f7265666c}\textbf{At most refl.} Let \(\alpha\) be a type. Equip \(\alpha\) with a preorder. Let \(a\) be an element of \(\alpha\). Then \(a \le a\).
\item[\paperlineanchor{premise:6c655f72666c}\paperlink{premise:6c655f72666c}{P244.}] \refstepcounter{paperpremise}\label{premise:6c655f72666c}\textbf{At most rfl.} Let \(\alpha\) be a type. Equip \(\alpha\) with a preorder. Let \(a\) be an element of \(\alpha\). Then \(a \le a\).
\item[\paperlineanchor{premise:6c745f6f665f6c655f6f665f6c74}\paperlink{premise:6c745f6f665f6c655f6f665f6c74}{P245.}] \refstepcounter{paperpremise}\label{premise:6c745f6f665f6c655f6f665f6c74}\textbf{Less than of at most of less than.} Let \(\alpha\) be a type. Equip \(\alpha\) with a preorder. Let \(a\), \(b\), and \(c\) be elements of \(\alpha\). Assume \(a \le b\). Assume \(b < c\). Then \(a < c\).
\item[\paperlineanchor{premise:6e655f6f665f6e655f6f665f6571}\paperlink{premise:6e655f6f665f6e655f6f665f6571}{P246.}] \refstepcounter{paperpremise}\label{premise:6e655f6f665f6e655f6f665f6571}\textbf{Differs from of differs from of equals.} Let \(\alpha\) be a sort. Let \(a\), \(b\), and \(c\) be elements of \(\alpha\). Assume \(a \ne b\). Assume \(b = c\). Then \(a \ne c\).
\item[\paperlineanchor{premise:6e6f745f636f6e6772}\paperlink{premise:6e6f745f636f6e6772}{P247.}] \refstepcounter{paperpremise}\label{premise:6e6f745f636f6e6772}\textbf{Not congruence.} Let \(a\) and \(b\) be propositions. Assume \(a \iff b\). Then \(\neg a \iff \neg b\).
\item[\paperlineanchor{premise:6e6f745f66616c73655f65715f74727565}\paperlink{premise:6e6f745f66616c73655f65715f74727565}{P248.}] \refstepcounter{paperpremise}\label{premise:6e6f745f66616c73655f65715f74727565}\textbf{Not false equals true.} \(\left(\neg \mathrm{False}\right) = \mathrm{True}\).
\item[\paperlineanchor{premise:6e6f745f66616c73655f696666}\paperlink{premise:6e6f745f66616c73655f696666}{P249.}] \refstepcounter{paperpremise}\label{premise:6e6f745f66616c73655f696666}\textbf{Not false if and only if.} \(\neg \mathrm{False} \iff \mathrm{True}\).
\item[\paperlineanchor{premise:6e6f745f6966665f6e6f74}\paperlink{premise:6e6f745f6966665f6e6f74}{P250.}] \refstepcounter{paperpremise}\label{premise:6e6f745f6966665f6e6f74}\textbf{Not if and only if not.} Let \(a\) and \(b\) be propositions. Then \(\left(\neg a \iff \neg b\right) \iff \left(a \iff b\right)\).
\item[\paperlineanchor{premise:6e6f745f6c74}\paperlink{premise:6e6f745f6c74}{P251.}] \refstepcounter{paperpremise}\label{premise:6e6f745f6c74}\textbf{Not less than.} Let \(\alpha\) be a type. Equip \(\alpha\) with a linear order. Let \(a\) and \(b\) be elements of \(\alpha\). Then \(\neg \left(a < b\right) \iff b \le a\).
\item[\paperlineanchor{premise:6e6f745f6e6f6e7472697669616c5f6966665f73756273696e676c65746f6e}\paperlink{premise:6e6f745f6e6f6e7472697669616c5f6966665f73756273696e676c65746f6e}{P252.}] \refstepcounter{paperpremise}\label{premise:6e6f745f6e6f6e7472697669616c5f6966665f73756273696e676c65746f6e}\textbf{Not nontrivial if and only if subsingleton.} Let \(\alpha\) be a type. Then \(\neg \operatorname{Nontrivial}\left(\alpha\right) \iff \alpha\text{ has at most one element}\).
\item[\paperlineanchor{premise:6e6f745f6f72}\paperlink{premise:6e6f745f6f72}{P253.}] \refstepcounter{paperpremise}\label{premise:6e6f745f6f72}\textbf{Not or.} Let \(p\) and \(q\) be propositions. Then \(\neg \left(p \lor q\right) \iff \neg p \land \neg q\).
\item[\paperlineanchor{premise:6f665f6465636964655f65715f74727565}\paperlink{premise:6f665f6465636964655f65715f74727565}{P254.}] \refstepcounter{paperpremise}\label{premise:6f665f6465636964655f65715f74727565}\textbf{Of decide equals true.} Let \(p\) be a proposition. Assume \(p\) is decidable. Assume \(\operatorname{decide}\left(p\right) = \mathrm{true}\). Then \(p\).
\item[\paperlineanchor{premise:6f665f65715f74727565}\paperlink{premise:6f665f65715f74727565}{P255.}] \refstepcounter{paperpremise}\label{premise:6f665f65715f74727565}\textbf{Of equals true.} Let \(p\) be a proposition. Assume \(p = \mathrm{True}\). Then \(p\).
\item[\paperlineanchor{premise:6f6e655f6e655f7a65726f}\paperlink{premise:6f6e655f6e655f7a65726f}{P256.}] \refstepcounter{paperpremise}\label{premise:6f6e655f6e655f7a65726f}\textbf{One differs from zero.} Let \(\alpha\) be a type. Fix a structure \(\mathit{inst}\) of type \(\operatorname{Zero}\left(\alpha\right)\). Fix a structure \(\mathit{inst}_{1}\) of type \(\operatorname{One}\left(\alpha\right)\). Let \(\mathit{inst}_{2}\) be a proof of \(\operatorname{Ne\,Zero}\left(1\right)\). Then \(1 \ne 0\).
\item[\paperlineanchor{premise:6f725f66616c7365}\paperlink{premise:6f725f66616c7365}{P257.}] \refstepcounter{paperpremise}\label{premise:6f725f66616c7365}\textbf{Or false.} Let \(p\) be a proposition. Then \(\left(p \lor \mathrm{False}\right) = p\).
\item[\paperlineanchor{premise:6f725f73656c66}\paperlink{premise:6f725f73656c66}{P258.}] \refstepcounter{paperpremise}\label{premise:6f725f73656c66}\textbf{Or self.} Let \(p\) be a proposition. Then \(\left(p \lor p\right) = p\).
\item[\paperlineanchor{premise:6f725f73656c665f696666}\paperlink{premise:6f725f73656c665f696666}{P259.}] \refstepcounter{paperpremise}\label{premise:6f725f73656c665f696666}\textbf{Or self if and only if.} Let \(a\) be a proposition. Then \(a \lor a \iff a\).
\item[\paperlineanchor{premise:7375627365745f7265666c}\paperlink{premise:7375627365745f7265666c}{P260.}] \refstepcounter{paperpremise}\label{premise:7375627365745f7265666c}\textbf{Subset refl.} Let \(\alpha\) be a type. Equip \(\alpha\) with a subset relation. Assume \((\mathit{x1} \mapsto \mathit{x2} \mapsto \mathit{x1} \subseteq \mathit{x2})\text{ is reflexive on }\alpha\). Let \(a\) be an element of \(\alpha\). Then \(a \subseteq a\).
\item[\paperlineanchor{premise:747275655f616e64}\paperlink{premise:747275655f616e64}{P261.}] \refstepcounter{paperpremise}\label{premise:747275655f616e64}\textbf{True and.} Let \(p\) be a proposition. Then \(\left(\mathrm{True} \land p\right) = p\).
\end{description}

\section{Finite sets}
\begin{lemma}\label{theorem:1}
Let \(\alpha\) be a type. Let \(s\) be a finite subset of \(\alpha\). Then \(\left|s\right| = \left|\operatorname{underlying\,multiset}(s)\right|\).
\end{lemma}

\begin{proof}
By definition, we obtain \(\left|s\right| = \left|\operatorname{underlying\,multiset}(s)\right|\).
\end{proof}

\begin{theorem}\label{theorem:2}
Let \(\alpha\) be a type. Let \(s\) be a finite subset of \(\alpha\). Then \(\left|\operatorname{underlying\,multiset}(s)\right| = \left|s\right|\).
\end{theorem}

\begin{proof}
By definition, we obtain \(\left|\operatorname{underlying\,multiset}(s)\right| = \left|s\right|\).
\end{proof}

\begin{theorem}\label{theorem:3}
Let \(\alpha\) be a type. Let \(m\) be a multiset of elements of \(\alpha\). Let \(\mathit{nodup}\) be a proof of \(m\text{ has no repeated elements}\). Then \(\left|\operatorname{as\,finite\,set}(m)\right| = \left|m\right|\).
\end{theorem}

\begin{proof}
By definition, we obtain \(\left|\operatorname{as\,finite\,set}(m)\right| = \left|m\right|\).
\end{proof}

\begin{theorem}\label{theorem:4}
Let \(\alpha\) be a type. Then \(\left|\varnothing\right| = 0\).
\end{theorem}

\begin{proof}
By definition, we obtain \(\left|\varnothing\right| = 0\).
\end{proof}

\begin{lemma}\label{theorem:5}
Let \(\alpha\) be a type. Let \(s\) and \(t\) be finite subsets of \(\alpha\). Assume \(s \subseteq t\). Then \(\left|s\right| \le \left|t\right|\).
\end{lemma}

\begin{proof}
By \paperlink{premise:46696e7365742e76616c5f6c655f696666}{Prerequisite~P\ref{premise:46696e7365742e76616c5f6c655f696666}}, we obtain \(\operatorname{underlying\,multiset}(s) \le \operatorname{underlying\,multiset}(t) \iff s \subseteq t\).
By the reverse implication, we obtain \(s \subseteq t \implies \operatorname{underlying\,multiset}(s) \le \operatorname{underlying\,multiset}(t)\).
By \paperlink{premise:4d756c74697365742e636172645f6c655f63617264}{Prerequisite~P\ref{premise:4d756c74697365742e636172645f6c655f63617264}}, we obtain \(\operatorname{underlying\,multiset}(s) \le \operatorname{underlying\,multiset}(t) \implies |\operatorname{underlying\,multiset}(s)| \le |\operatorname{underlying\,multiset}(t)|\).
Composing these implications, we obtain \(s \subseteq t \implies \left|s\right| \le \left|t\right|\).
\end{proof}

\begin{corollary}\label{theorem:6}
Let \(\alpha\) be a type. Then \(\operatorname{Monotone}\left(\operatorname{Finset\mathord{:}card}\right)\).
\end{corollary}

\begin{proof}
Apply \paperlink{theorem:5}{Lemma~\ref{theorem:5}}.
\end{proof}

\begin{lemma}\label{theorem:7}
Let \(\alpha\) be a type. Let \(s\) be a finite subset of \(\alpha\). Then \(\left|s\right| = 0 \iff s = \varnothing\).
\end{lemma}

\begin{proof}
By \paperlink{premise:4d756c74697365742e636172645f65715f7a65726f}{Prerequisite~P\ref{premise:4d756c74697365742e636172645f65715f7a65726f}}, we obtain \(\left|\operatorname{underlying\,multiset}(s)\right| = 0 \iff \operatorname{underlying\,multiset}(s) = 0\).
By \paperlink{premise:46696e7365742e76616c5f65715f7a65726f}{Prerequisite~P\ref{premise:46696e7365742e76616c5f65715f7a65726f}}, we obtain \(\operatorname{underlying\,multiset}(s) = 0 \iff s = \varnothing\).
Chaining the equivalences, we obtain \(\left|s\right| = 0 \iff s = \varnothing\).
\end{proof}

\begin{lemma}\label{theorem:8}
Let \(\alpha\) be a type. Let \(s\) be a finite subset of \(\alpha\). Then \(\left|s\right| \ne 0 \iff s\ne\varnothing\).
\end{lemma}

\begin{proof}
By \paperlink{theorem:7}{Lemma~\ref{theorem:7}}, we obtain \(\left|s\right| = 0 \iff s = \varnothing\).
By \paperlink{premise:4966662e6e65}{Prerequisite~P\ref{premise:4966662e6e65}} with \(\mathbb{N}\) as its type parameter \(\alpha\) and \(\beta = \mathcal{P}_{\mathrm{fin}}(\alpha)\), we obtain \(\left|s\right| \ne 0 \iff s \ne \varnothing\).
By \paperlink{premise:46696e7365742e6e6f6e656d7074795f6966665f6e655f656d707479}{Prerequisite~P\ref{premise:46696e7365742e6e6f6e656d7074795f6966665f6e655f656d707479}}, we obtain \(s\ne\varnothing \iff s \ne \varnothing\).
Reversing the equivalence, we obtain \(s \ne \varnothing \iff s\ne\varnothing\).
Chaining the equivalences, we obtain \(\left|s\right| \ne 0 \iff s\ne\varnothing\).
\end{proof}

\begin{lemma}\label{theorem:9}
Let \(\alpha\) be a type. Let \(s\) be a finite subset of \(\alpha\). Then \(0 < \left|s\right| \iff s\ne\varnothing\).
\end{lemma}

\begin{proof}
By \paperlink{premise:4e61742e706f735f6966665f6e655f7a65726f}{Prerequisite~P\ref{premise:4e61742e706f735f6966665f6e655f7a65726f}}, we obtain \(0 < \left|s\right| \iff \left|s\right| \ne 0\).
By \paperlink{theorem:8}{Lemma~\ref{theorem:8}}, we obtain \(\left|s\right| \ne 0 \iff s\ne\varnothing\).
Chaining the equivalences, we obtain \(0 < \left|s\right| \iff s\ne\varnothing\).
\end{proof}

\begin{corollary}\label{theorem:10}
Let \(\alpha\) be a type. Let \(s\) be a finite subset of \(\alpha\). Then \(1 \le \left|s\right| \iff s\ne\varnothing\).
\end{corollary}

\begin{proof}
Apply \paperlink{theorem:9}{Lemma~\ref{theorem:9}}.
\end{proof}

\begin{theorem}[{Alias}]\label{theorem:11}
Let \(\alpha\) be a type. Let \(s\) be a finite subset of \(\alpha\). Assume \(s\ne\varnothing\). Then \(0 < \left|s\right|\).
\end{theorem}

\begin{proof}
Apply the reverse implication of \paperlink{theorem:9}{Lemma~\ref{theorem:9}}.
\end{proof}

\begin{theorem}[{Alias}]\label{theorem:12}
Let \(\alpha\) be a type. Let \(s\) be a finite subset of \(\alpha\). Assume \(s\ne\varnothing\). Then \(\left|s\right| \ne 0\).
\end{theorem}

\begin{proof}
Apply the reverse implication of \paperlink{theorem:8}{Lemma~\ref{theorem:8}}.
\end{proof}

\begin{theorem}\label{theorem:13}
Let \(\alpha\) be a type. Let \(s\) be a finite subset of \(\alpha\). Let \(a\) be an element of \(\alpha\). Assume \(a \in s\). Then \(\left|s\right| \ne 0\).
\end{theorem}

\begin{proof}
By \paperlink{theorem:7}{Lemma~\ref{theorem:7}}, we obtain \(\left|s\right| = 0 \iff s = \varnothing\).
By \paperlink{premise:6e6f745f636f6e6772}{Prerequisite~P\ref{premise:6e6f745f636f6e6772}}, we obtain \(\neg \left(\left|s\right| = 0\right) \iff \neg \left(s = \varnothing\right)\).
Use the hypothesis \(a \in s\).
By \paperlink{premise:46696e7365742e6e655f656d7074795f6f665f6d656d}{Prerequisite~P\ref{premise:46696e7365742e6e655f656d7074795f6f665f6d656d}}, we obtain \(s \ne \varnothing\).
By the reverse implication, we obtain \(\left|s\right| \ne 0\).
\end{proof}

\begin{lemma}\label{theorem:14}
Let \(\alpha\) be a type. Let \(a\) be an element of \(\alpha\). Then \(\left|\{a\}\right| = 1\).
\end{lemma}

\begin{proof}
By \paperlink{premise:4d756c74697365742e636172645f73696e676c65746f6e}{Prerequisite~P\ref{premise:4d756c74697365742e636172645f73696e676c65746f6e}}, we obtain \(\left|\{a\}\right| = 1\).
\end{proof}

\begin{theorem}\label{theorem:15}
Let \(\alpha\) be a type. Let \(s\) be a finite subset of \(\alpha\). Let \(a\) be an element of \(\alpha\). Assume equality on \(\alpha\) is decidable. Then \(\left|\{a\} \cap s\right| \le 1\).
\end{theorem}

\begin{proof}
Consider the following case.
\begin{proofstep}
Assume \(a \notin s\).
Rewriting with \paperlink{premise:46696e7365742e73696e676c65746f6e5f696e7465725f6f665f6e6f745f6d656d}{Prerequisite~P\ref{premise:46696e7365742e73696e676c65746f6e5f696e7465725f6f665f6e6f745f6d656d}} (using \(a \notin s\)) and \paperlink{premise:4e61742e7a65726f5f6c65}{Prerequisite~P\ref{premise:4e61742e7a65726f5f6c65}} gives \(\left|\{a\} \cap s\right| \le 1\).
\end{proofstep}
Consider the following case.
\begin{proofstep}
Assume \(a \in s\).
Rewriting with \paperlink{premise:46696e7365742e73696e676c65746f6e5f696e7465725f6f665f6d656d}{Prerequisite~P\ref{premise:46696e7365742e73696e676c65746f6e5f696e7465725f6f665f6d656d}} (using \(a \in s\)), \paperlink{theorem:14}{Lemma~\ref{theorem:14}} and \paperlink{premise:6c655f7265666c}{Prerequisite~P\ref{premise:6c655f7265666c}} gives \(\left|\{a\} \cap s\right| \le 1\).
\end{proofstep}
The proposition and its negation exhaust the cases. Hence \(\left|\{a\} \cap s\right| \le 1\).
\end{proof}

\begin{lemma}\label{theorem:16}
Let \(\alpha\) be a type. Let \(s\) be a finite subset of \(\alpha\). Let \(a\) be an element of \(\alpha\). Let \(h\) be a proof of \(a \notin s\). Then \(\left|\operatorname{Finset\mathord{:}cons}\left(a, s\right)\right| = \left|s\right| + 1\).
\end{lemma}

\begin{proof}
By \paperlink{premise:4d756c74697365742e636172645f636f6e73}{Prerequisite~P\ref{premise:4d756c74697365742e636172645f636f6e73}} with \(\operatorname{underlying\,multiset}(s)\) in place of its parameter \(s\), we obtain \(\left|\operatorname{Finset\mathord{:}cons}\left(a, s\right)\right| = \left|s\right| + 1\).
\end{proof}

\section{Insert erase}
\begin{lemma}\label{theorem:17}
Let \(\alpha\) be a type. Let \(s\) be a finite subset of \(\alpha\). Let \(a\) be an element of \(\alpha\). Assume equality on \(\alpha\) is decidable. Assume \(a \notin s\). Then \(\left|\{a\} \cup s\right| = \left|s\right| + 1\).
\end{lemma}

\begin{proof}
Use the hypothesis \(a \notin s\).
By \paperlink{premise:46696e7365742e636f6e735f65715f696e73657274}{Prerequisite~P\ref{premise:46696e7365742e636f6e735f65715f696e73657274}}, we obtain \(\operatorname{Finset\mathord{:}cons}\left(a, s\right) = \{a\} \cup s\).
Reversing the equality, we obtain \(\{a\} \cup s = \operatorname{Finset\mathord{:}cons}\left(a, s\right)\).
Applying the same function to equal arguments, we obtain \((|\{a\} \cup s| = |s| + 1) = (|\operatorname{Finset\mathord{:}cons}(a,\allowbreak{} s)| = |s| + 1)\).
By the preceding argument, we obtain \((|\{a\} \cup s| = |s| + 1) = (|\operatorname{Finset\mathord{:}cons}(a,\allowbreak{} s)| = |s| + 1)\).
By \paperlink{theorem:16}{Lemma~\ref{theorem:16}}, we obtain \(\left|\operatorname{Finset\mathord{:}cons}\left(a, s\right)\right| = \left|s\right| + 1\).
Applying the same function to equal arguments, we obtain \((|\operatorname{Finset\mathord{:}cons}(a,\allowbreak{} s)| = |s| + 1) = (|s| + 1 = |s| + 1)\).
By the preceding argument, we obtain \((|\operatorname{Finset\mathord{:}cons}(a,\allowbreak{} s)| = |s| + 1) = (|s| + 1 = |s| + 1)\).
By definition, we obtain \(\left|s\right| + 1 = \left|s\right| + 1\).
Substituting along the equality, we obtain \(\left|\operatorname{Finset\mathord{:}cons}\left(a, s\right)\right| = \left|s\right| + 1\).
Substituting along the equality, we obtain \(\left|\{a\} \cup s\right| = \left|s\right| + 1\).
\end{proof}

\begin{lemma}\label{theorem:18}
Let \(\alpha\) be a type. Let \(s\) be a finite subset of \(\alpha\). Let \(a\) be an element of \(\alpha\). Assume equality on \(\alpha\) is decidable. Assume \(a \in s\). Then \(\left|\{a\} \cup s\right| = \left|s\right|\).
\end{lemma}

\begin{proof}
By \paperlink{premise:46696e7365742e696e736572745f65715f6f665f6d656d}{Prerequisite~P\ref{premise:46696e7365742e696e736572745f65715f6f665f6d656d}} applied to \(a \in s\), we have \(\{a\} \cup s = s\).
Substitution using \paperlink{premise:46696e7365742e696e736572745f65715f6f665f6d656d}{Prerequisite~P\ref{premise:46696e7365742e696e736572745f65715f6f665f6d656d}} applied to \(a \in s\) gives \(\left|\{a\} \cup s\right| = \left|s\right|\).
\end{proof}

\begin{lemma}\label{theorem:19}
Let \(\alpha\) be a type. Assume equality on \(\alpha\) is decidable. Let \(a\) be an element of \(\alpha\). Let \(s\) be a finite subset of \(\alpha\). Then \(\left|\{a\} \cup s\right| \le \left|s\right| + 1\).
\end{lemma}

\begin{proof}
Distinguish whether \(a \in s\) holds.
Consider the following case.
\begin{proofstep}
Assume \(a \in s\).
Use the hypothesis \(a \in s\).
By \paperlink{premise:46696e7365742e696e736572745f65715f6f665f6d656d}{Prerequisite~P\ref{premise:46696e7365742e696e736572745f65715f6f665f6d656d}}, we obtain \(\{a\} \cup s = s\).
Applying the same function to equal arguments, we obtain \((|\{a\} \cup s| \le |s| + 1) = (|s| \le |s| + 1)\).
By the preceding argument, we obtain \((|\{a\} \cup s| \le |s| + 1) = (|s| \le |s| + 1)\).
By \paperlink{premise:4e61742e6c655f73756363}{Prerequisite~P\ref{premise:4e61742e6c655f73756363}} with \(n = \left|s\right|\), we obtain \(\left|s\right| \le \left|s\right| + 1\).
Substituting along the equality, we obtain \(\left|\{a\} \cup s\right| \le \left|s\right| + 1\).
\end{proofstep}
Consider the following case.
\begin{proofstep}
Assume \(a \notin s\).
By \paperlink{theorem:17}{Lemma~\ref{theorem:17}}, we obtain \(\left|\{a\} \cup s\right| = \left|s\right| + 1\).
Applying the same function to equal arguments, we obtain \((|\{a\} \cup s| \le |s| + 1) = (|s| + 1 \le |s| + 1)\).
By the preceding argument, we obtain \((|\{a\} \cup s| \le |s| + 1) = (|s| + 1 \le |s| + 1)\).
By \paperlink{premise:6c655f7265666c}{Prerequisite~P\ref{premise:6c655f7265666c}} with \(\mathbb{N}\) as its type parameter \(\alpha\) and \(\left|s\right| + 1\) in place of its parameter \(a\), we obtain \(\left|s\right| + 1 \le \left|s\right| + 1\).
Substituting along the equality, we obtain \(\left|\{a\} \cup s\right| \le \left|s\right| + 1\).
\end{proofstep}
The proposition or its negation holds. Hence \(\left|\{a\} \cup s\right| \le \left|s\right| + 1\).
\end{proof}

\begin{corollary}\label{theorem:20}
Let \(\alpha\) be a type. Assume equality on \(\alpha\) is decidable. Let \(a\) and \(b\) be elements of \(\alpha\). Then \(\left|\{a\} \cup \{b\}\right| \le 2\).
\end{corollary}

\begin{proof}
Apply \paperlink{theorem:19}{Lemma~\ref{theorem:19}} with \(s = \{b\}\).
\end{proof}

\begin{lemma}\label{theorem:21}
Let \(\alpha\) be a type. Assume equality on \(\alpha\) is decidable. Let \(a\), \(b\), and \(c\) be elements of \(\alpha\). Then \(\left|\{a\} \cup \left(\{b\} \cup \{c\}\right)\right| \le 3\).
\end{lemma}

\begin{proof}
By \paperlink{theorem:19}{Lemma~\ref{theorem:19}} with \(s = \{b\} \cup \{c\}\), we obtain \(\left|\{a\} \cup \left(\{b\} \cup \{c\}\right)\right| \le \left|\{b\} \cup \{c\}\right| + 1\).
By \paperlink{theorem:20}{Corollary~\ref{theorem:20}}, we obtain \(\left|\{b\} \cup \{c\}\right| \le 2\).
By \paperlink{premise:4e61742e737563635f6c655f73756363}{Prerequisite~P\ref{premise:4e61742e737563635f6c655f73756363}}, we obtain \(\left|\{b\} \cup \{c\}\right| + 1 \le 2 + 1\).
By \paperlink{premise:4c452e6c652e7472616e73}{\emph{Alias}} (\paperlink{premise:4c452e6c652e7472616e73}{Prerequisite~P\ref{premise:4c452e6c652e7472616e73}}) with \(\mathbb{N}\) as its type parameter \(\alpha\), we obtain \(\left|\{a\} \cup \left(\{b\} \cup \{c\}\right)\right| \le 3\).
\end{proof}

\begin{lemma}\label{theorem:22}
Let \(\alpha\) be a type. Assume equality on \(\alpha\) is decidable. Let \(a\), \(b\), \(c\), and \(d\) be elements of \(\alpha\). Then \(\left|\{a\} \cup \left(\{b\} \cup \left(\{c\} \cup \{d\}\right)\right)\right| \le 4\).
\end{lemma}

\begin{proof}
By \paperlink{theorem:19}{Lemma~\ref{theorem:19}} with \(s = \{b\} \cup \left(\{c\} \cup \{d\}\right)\), we obtain \(|\{a\} \cup (\{b\} \cup (\{c\} \cup \{d\}))| \le |\{b\} \cup (\{c\} \cup \{d\})| + 1\).
By \paperlink{theorem:21}{Lemma~\ref{theorem:21}}, we obtain \(\left|\{b\} \cup \left(\{c\} \cup \{d\}\right)\right| \le 3\).
By \paperlink{premise:4e61742e737563635f6c655f73756363}{Prerequisite~P\ref{premise:4e61742e737563635f6c655f73756363}}, we obtain \(\left|\{b\} \cup \left(\{c\} \cup \{d\}\right)\right| + 1 \le 3 + 1\).
By \paperlink{premise:4c452e6c652e7472616e73}{\emph{Alias}} (\paperlink{premise:4c452e6c652e7472616e73}{Prerequisite~P\ref{premise:4c452e6c652e7472616e73}}) with \(\mathbb{N}\) as its type parameter \(\alpha\), we obtain \(\left|\{a\} \cup \left(\{b\} \cup \left(\{c\} \cup \{d\}\right)\right)\right| \le 4\).
\end{proof}

\begin{lemma}\label{theorem:23}
Let \(\alpha\) be a type. Assume equality on \(\alpha\) is decidable. Let \(a\), \(b\), \(c\), \(d\), and \(e\) be elements of \(\alpha\). Then
\begin{center}
\(\displaystyle |\{a\} \cup (\{b\} \cup (\{c\} \cup (\{d\} \cup \{e\})))| \le 5.\)
\end{center}
\end{lemma}

\begin{proof}
By \paperlink{theorem:19}{Lemma~\ref{theorem:19}} with \(s = \{b\} \cup \left(\{c\} \cup \left(\{d\} \cup \{e\}\right)\right)\), we obtain \(|\{a\} \cup (\{b\} \cup (\{c\} \cup (\{d\} \cup \{e\})))| \le |\{b\} \cup (\{c\} \cup (\{d\} \cup \{e\}))| + 1\).
By \paperlink{theorem:22}{Lemma~\ref{theorem:22}}, we obtain \(\left|\{b\} \cup \left(\{c\} \cup \left(\{d\} \cup \{e\}\right)\right)\right| \le 4\).
By \paperlink{premise:4e61742e737563635f6c655f73756363}{Prerequisite~P\ref{premise:4e61742e737563635f6c655f73756363}}, we obtain \(\left|\{b\} \cup \left(\{c\} \cup \left(\{d\} \cup \{e\}\right)\right)\right| + 1 \le 4 + 1\).
By \paperlink{premise:4c452e6c652e7472616e73}{\emph{Alias}} (\paperlink{premise:4c452e6c652e7472616e73}{Prerequisite~P\ref{premise:4c452e6c652e7472616e73}}) with \(\mathbb{N}\) as its type parameter \(\alpha\), we obtain \(\left|\{a\} \cup \left(\{b\} \cup \left(\{c\} \cup \left(\{d\} \cup \{e\}\right)\right)\right)\right| \le 5\).
\end{proof}

\begin{theorem}\label{theorem:24}
Let \(\alpha\) be a type. Assume equality on \(\alpha\) is decidable. Let \(a\), \(b\), \(c\), \(d\), \(e\), and \(f\) be elements of \(\alpha\). Then
\begin{center}
\(\displaystyle |\{a\} \cup (\{b\} \cup (\{c\} \cup (\{d\} \cup (\{e\} \cup \{f\}))))| \le 6.\)
\end{center}
\end{theorem}

\begin{proof}
By \paperlink{theorem:19}{Lemma~\ref{theorem:19}} with \(s = \{b\} \cup \left(\{c\} \cup \left(\{d\} \cup \left(\{e\} \cup \{f\}\right)\right)\right)\), we obtain \(|\{a\} \cup (\{b\} \cup (\{c\} \cup (\{d\} \cup (\{e\} \cup \{f\}))))| \le |\{b\} \cup (\{c\} \cup (\{d\} \cup (\{e\} \cup \{f\})))| + 1\).
By \paperlink{theorem:23}{Lemma~\ref{theorem:23}}, we obtain \(\left|\{b\} \cup \left(\{c\} \cup \left(\{d\} \cup \left(\{e\} \cup \{f\}\right)\right)\right)\right| \le 5\).
By \paperlink{premise:4e61742e737563635f6c655f73756363}{Prerequisite~P\ref{premise:4e61742e737563635f6c655f73756363}}, we obtain \(|\{b\} \cup (\{c\} \cup (\{d\} \cup (\{e\} \cup \{f\})))| + 1 \le 5 + 1\).
By \paperlink{premise:4c452e6c652e7472616e73}{\emph{Alias}} (\paperlink{premise:4c452e6c652e7472616e73}{Prerequisite~P\ref{premise:4c452e6c652e7472616e73}}) with \(\mathbb{N}\) as its type parameter \(\alpha\), we obtain \(|\{a\} \cup (\{b\} \cup (\{c\} \cup (\{d\} \cup (\{e\} \cup \{f\}))))| \le 6\).
\end{proof}

\begin{lemma}\label{theorem:25}
Let \(\alpha\) be a type. Let \(s\) be a finite subset of \(\alpha\). Let \(a\) be an element of \(\alpha\). Assume equality on \(\alpha\) is decidable. Then
\begin{center}
\(\displaystyle |\{a\} \cup s| = \begin{cases}|s| &\text{if } a \in s\\ |s| + 1 &\text{otherwise}\end{cases}.\)
\end{center}
\end{lemma}

\begin{proof}
Consider the following cases.
\begin{proofstep}
Assume \(a \in s\).
Starting from \paperlink{theorem:18}{Lemma~\ref{theorem:18}} applied to \(a \in s\), substitution using \paperlink{premise:69665f706f73}{Prerequisite~P\ref{premise:69665f706f73}} applied to \(a \in s\) gives \(|\{a\} \cup s| = \begin{cases}|s| &\text{if } a \in s\\ |s| + 1 &\text{otherwise}\end{cases}\).
\end{proofstep}
\begin{proofstep}
Assume \(a \notin s\).
Starting from \paperlink{theorem:17}{Lemma~\ref{theorem:17}} applied to \(a \notin s\), substitution using \paperlink{premise:69665f6e6567}{Prerequisite~P\ref{premise:69665f6e6567}} applied to \(a \notin s\) gives \(|\{a\} \cup s| = \begin{cases}|s| &\text{if } a \in s\\ |s| + 1 &\text{otherwise}\end{cases}\).
\end{proofstep}
The cases exhaust the possibilities.
\end{proof}

\begin{theorem}\label{theorem:26}
Let \(\alpha\) be a type. Let \(a\) and \(b\) be elements of \(\alpha\). Assume equality on \(\alpha\) is decidable. Then \(\left|\{a\} \cup \{b\}\right| = 1 \lor \left|\{a\} \cup \{b\}\right| = 2\).
\end{theorem}

\begin{proof}
The substitution uses the following argument.
\begin{proofstep}
Assume \(a = b\).
By \paperlink{theorem:14}{Lemma~\ref{theorem:14}} with \(b\) in place of its parameter \(a\), we obtain \(\left|\{b\}\right| = 1\).
\end{proofstep}
The substitution uses the following argument.
\begin{proofstep}
Assume \(\neg \left(a = b\right)\).
By \paperlink{theorem:14}{Lemma~\ref{theorem:14}} with \(b\) in place of its parameter \(a\), we obtain \(\left|\{b\}\right| = 1\).
Applying the same function to equal arguments, we obtain \(\left|\{b\}\right| + 1 = 1 + 1\).
\end{proofstep}
The substitution uses the following argument.
\begin{proofstep}
Assume \(a = b\).
By \paperlink{theorem:14}{Lemma~\ref{theorem:14}} with \(b\) in place of its parameter \(a\), we obtain \(\left|\{b\}\right| = 1\).
\end{proofstep}
The substitution uses the following argument.
\begin{proofstep}
Assume \(\neg \left(a = b\right)\).
By \paperlink{theorem:14}{Lemma~\ref{theorem:14}} with \(b\) in place of its parameter \(a\), we obtain \(\left|\{b\}\right| = 1\).
Applying the same function to equal arguments, we obtain \(\left|\{b\}\right| + 1 = 1 + 1\).
\end{proofstep}
By \paperlink{premise:436c6173736963616c2e6f725f6966665f6e6f745f696d705f6c656674}{Prerequisite~P\ref{premise:436c6173736963616c2e6f725f6966665f6e6f745f696d705f6c656674}}, we obtain \(a = b \lor \neg \left(a = b\right) \iff \neg \left(a = b\right) \implies \neg \left(a = b\right)\).
For the required implication, argue as follows.
\begin{proofstep}
Assume \(\neg \left(a = b\right)\).
Use the hypothesis \(\neg \left(a = b\right)\).
\end{proofstep}
By the reverse implication, we obtain \(a = b \lor \neg \left(a = b\right)\).
Rewriting with \paperlink{theorem:25}{Lemma~\ref{theorem:25}}, \paperlink{premise:6974655f636f6e6772}{Prerequisite~P\ref{premise:6974655f636f6e6772}}, \paperlink{premise:46696e7365742e6d656d5f73696e676c65746f6e}{Prerequisite~P\ref{premise:46696e7365742e6d656d5f73696e676c65746f6e}}, \paperlink{premise:6974655f65715f6c6566745f696666}{Prerequisite~P\ref{premise:6974655f65715f6c6566745f696666}}, \paperlink{premise:696d706c6965735f636f6e6772}{Prerequisite~P\ref{premise:696d706c6965735f636f6e6772}}, \paperlink{premise:65715f66616c7365}{Prerequisite~P\ref{premise:65715f66616c7365}}, \paperlink{premise:4e61742e737563635f6e655f73656c66}{Prerequisite~P\ref{premise:4e61742e737563635f6e655f73656c66}}, \paperlink{premise:696d705f66616c7365}{Prerequisite~P\ref{premise:696d705f66616c7365}}, \paperlink{premise:446563696461626c652e6e6f745f6e6f74}{Prerequisite~P\ref{premise:446563696461626c652e6e6f745f6e6f74}}, \paperlink{premise:6974655f65715f72696768745f696666}{Prerequisite~P\ref{premise:6974655f65715f72696768745f696666}}, \paperlink{premise:65715f66616c73655f6f665f646563696465}{Prerequisite~P\ref{premise:65715f66616c73655f6f665f646563696465}} and the 4 identities established above in the preceding result gives \(\left|\{a\} \cup \{b\}\right| = 1 \lor \left|\{a\} \cup \{b\}\right| = 2\).
\end{proof}

\begin{lemma}\label{theorem:27}
Let \(\alpha\) be a type. Let \(a\) and \(b\) be elements of \(\alpha\). Assume equality on \(\alpha\) is decidable. Assume \(a \ne b\). Then \(\left|\{a\} \cup \{b\}\right| = 2\).
\end{lemma}

\begin{proof}
By \paperlink{premise:46696e7365742e6e6f745f6d656d5f73696e676c65746f6e}{Prerequisite~P\ref{premise:46696e7365742e6e6f745f6d656d5f73696e676c65746f6e}}, we obtain \(a \notin \{b\} \iff a \ne b\).
Use the hypothesis \(a \ne b\).
By the reverse implication, we obtain \(a \notin \{b\}\).
By \paperlink{theorem:17}{Lemma~\ref{theorem:17}}, we obtain \(\left|\{a\} \cup \{b\}\right| = \left|\{b\}\right| + 1\).
Applying the same function to equal arguments, we obtain \(\left(\left|\{a\} \cup \{b\}\right| = 2\right) = \left(\left|\{b\}\right| + 1 = 2\right)\).
By the preceding argument, we obtain \(\left(\left|\{a\} \cup \{b\}\right| = 2\right) = \left(\left|\{b\}\right| + 1 = 2\right)\).
By \paperlink{theorem:14}{Lemma~\ref{theorem:14}} with \(b\) in place of its parameter \(a\), we obtain \(\left|\{b\}\right| = 1\).
Applying the same function to equal arguments, we obtain \(\left(\left|\{b\}\right| + 1 = 2\right) = \left(1 + 1 = 2\right)\).
By the preceding argument, we obtain \(\left(\left|\{b\}\right| + 1 = 2\right) = \left(1 + 1 = 2\right)\).
By definition, we obtain \(1 + 1 = 2\).
Substituting along the equality, we obtain \(\left|\{b\}\right| + 1 = 2\).
Substituting along the equality, we obtain \(\left|\{a\} \cup \{b\}\right| = 2\).
\end{proof}

\begin{lemma}\label{theorem:28}
Let \(\alpha\) be a type. Let \(s\) be a finite subset of \(\alpha\). Let \(a\) be an element of \(\alpha\). Assume equality on \(\alpha\) is decidable. Assume \(a \in s\). Then \(\left|s \setminus \{a\}\right| = \left|s\right| - 1\).
\end{lemma}

\begin{proof}
By \paperlink{premise:4d756c74697365742e636172645f65726173655f6f665f6d656d}{Prerequisite~P\ref{premise:4d756c74697365742e636172645f65726173655f6f665f6d656d}}, we obtain \(a \in s \implies \left|s \setminus \{a\}\right| = \left|s\right| - 1\).
\end{proof}

\begin{theorem}\label{theorem:29}
Let \(\alpha\) be a type. Let \(s\) be a finite subset of \(\alpha\). Let \(a\) be an element of \(\alpha\). Assume equality on \(\alpha\) is decidable. Assume \(a \in s\). Then \(\left|s \setminus \{a\}\right| + 1 = \left|s\right|\).
\end{theorem}

\begin{proof}
By \paperlink{premise:4d756c74697365742e636172645f65726173655f6164645f6f6e65}{Prerequisite~P\ref{premise:4d756c74697365742e636172645f65726173655f6164645f6f6e65}}, we obtain \(a \in s \implies \left|s \setminus \{a\}\right| + 1 = \left|s\right|\).
\end{proof}

\begin{theorem}\label{theorem:30}
Let \(\alpha\) be a type. Let \(s\) be a finite subset of \(\alpha\). Let \(a\) be an element of \(\alpha\). Assume equality on \(\alpha\) is decidable. Assume \(a \in s\). Then \(\left|s \setminus \{a\}\right| < \left|s\right|\).
\end{theorem}

\begin{proof}
By \paperlink{premise:4d756c74697365742e636172645f65726173655f6c745f6f665f6d656d}{Prerequisite~P\ref{premise:4d756c74697365742e636172645f65726173655f6c745f6f665f6d656d}}, we obtain \(a \in s \implies \left|s \setminus \{a\}\right| < \left|s\right|\).
\end{proof}

\begin{theorem}\label{theorem:31}
Let \(\alpha\) be a type. Let \(s\) be a finite subset of \(\alpha\). Let \(a\) be an element of \(\alpha\). Assume equality on \(\alpha\) is decidable. Then \(\left|s \setminus \{a\}\right| \le \left|s\right|\).
\end{theorem}

\begin{proof}
By \paperlink{premise:4d756c74697365742e636172645f65726173655f6c65}{Prerequisite~P\ref{premise:4d756c74697365742e636172645f65726173655f6c65}}, we obtain \(\left|s \setminus \{a\}\right| \le \left|s\right|\).
\end{proof}

\begin{theorem}\label{theorem:32}
Let \(\alpha\) be a type. Let \(s\) be a finite subset of \(\alpha\). Let \(a\) be an element of \(\alpha\). Assume equality on \(\alpha\) is decidable. Then \(\left|s\right| - 1 \le \left|s \setminus \{a\}\right|\).
\end{theorem}

\begin{proof}
Distinguish whether \(a \in s\) holds.
Consider the following case.
\begin{proofstep}
Assume \(a \in s\).
By \paperlink{theorem:28}{Lemma~\ref{theorem:28}}, we obtain \(\left|s \setminus \{a\}\right| = \left|s\right| - 1\).
By \paperlink{premise:45712e6765}{Prerequisite~P\ref{premise:45712e6765}} with \(\mathbb{N}\) as its type parameter \(\alpha\), we obtain \(\left|s\right| - 1 \le \left|s \setminus \{a\}\right|\).
\end{proofstep}
Consider the following case.
\begin{proofstep}
Assume \(a \notin s\).
Use the hypothesis \(a \notin s\).
By \paperlink{premise:46696e7365742e65726173655f65715f6f665f6e6f745f6d656d}{Prerequisite~P\ref{premise:46696e7365742e65726173655f65715f6f665f6e6f745f6d656d}}, we obtain \(s \setminus \{a\} = s\).
Applying the same function to equal arguments, we obtain \((|s| - 1 \le |s \setminus \{a\}|) = (|s| - 1 \le |s|)\).
By the preceding argument, we obtain \((|s| - 1 \le |s \setminus \{a\}|) = (|s| - 1 \le |s|)\).
By \paperlink{premise:4e61742e7375625f6c65}{Prerequisite~P\ref{premise:4e61742e7375625f6c65}} with \(n = \left|s\right|\) and \(m = 1\), we obtain \(\left|s\right| - 1 \le \left|s\right|\).
Substituting along the equality, we obtain \(\left|s\right| - 1 \le \left|s \setminus \{a\}\right|\).
\end{proofstep}
The proposition or its negation holds. Hence \(\left|s\right| - 1 \le \left|s \setminus \{a\}\right|\).
\end{proof}

\begin{theorem}\label{theorem:33}
Let \(\alpha\) be a type. Let \(s\) be a finite subset of \(\alpha\). Let \(a\) be an element of \(\alpha\). Assume equality on \(\alpha\) is decidable. Then
\begin{center}
\(\displaystyle |s \setminus \{a\}| = \begin{cases}|s| - 1 &\text{if } a \in s\\ |s| &\text{otherwise}\end{cases}.\)
\end{center}
\end{theorem}

\begin{proof}
By \paperlink{premise:4d756c74697365742e636172645f65726173655f65715f697465}{Prerequisite~P\ref{premise:4d756c74697365742e636172645f65726173655f65715f697465}}, we obtain \(|s \setminus \{a\}| = \begin{cases}|s| - 1 &\text{if } a \in s\\ |s| &\text{otherwise}\end{cases}\).
\end{proof}

\section{Finite sets}
\begin{lemma}\label{theorem:34}
Let \(n\) be a natural number. Then \(\left|\{i\in\mathbb{N}\mid i<n\}\right| = n\).
\end{lemma}

\begin{proof}
By \paperlink{premise:4d756c74697365742e636172645f72616e6765}{Prerequisite~P\ref{premise:4d756c74697365742e636172645f72616e6765}}, we obtain \(\left|\{i\in\mathbb{N}\mid i<n\}\right| = n\).
\end{proof}

\begin{lemma}\label{theorem:35}
Let \(\alpha\) be a type. Let \(s\) be a finite subset of \(\alpha\). Then \(\left|\operatorname{Finset\mathord{:}attach}\left(s\right)\right| = \left|s\right|\).
\end{lemma}

\begin{proof}
By \paperlink{premise:4d756c74697365742e636172645f617474616368}{Prerequisite~P\ref{premise:4d756c74697365742e636172645f617474616368}}, we obtain \(\left|\operatorname{Finset\mathord{:}attach}\left(s\right)\right| = \left|s\right|\).
\end{proof}

\section{To mllistultiset}
\begin{theorem}\label{theorem:36}
Let \(\alpha\) be a type. Assume equality on \(\alpha\) is decidable. Let \(m\) be a multiset of elements of \(\alpha\). Then
\begin{center}
\(\displaystyle |\operatorname{Multiset\mathord{:}to\,Finset}(m)| = |\operatorname{Multiset\mathord{:}dedup}(m)|.\)
\end{center}
\end{theorem}

\begin{proof}
By definition, we obtain \(|\operatorname{Multiset\mathord{:}to\,Finset}(m)| = |\operatorname{Multiset\mathord{:}dedup}(m)|\).
\end{proof}

\begin{lemma}\label{theorem:37}
Let \(\alpha\) be a type. Assume equality on \(\alpha\) is decidable. Let \(m\) be a multiset of elements of \(\alpha\). Then \(\left|\operatorname{Multiset\mathord{:}to\,Finset}\left(m\right)\right| \le \left|m\right|\).
\end{lemma}

\begin{proof}
By \paperlink{premise:4d756c74697365742e64656475705f6c65}{Prerequisite~P\ref{premise:4d756c74697365742e64656475705f6c65}} with \(s = m\), we obtain \(\operatorname{Multiset\mathord{:}dedup}\left(m\right) \le m\).
By \paperlink{premise:4d756c74697365742e636172645f6c655f63617264}{Prerequisite~P\ref{premise:4d756c74697365742e636172645f6c655f63617264}}, we obtain \(\left|\operatorname{Multiset\mathord{:}to\,Finset}\left(m\right)\right| \le \left|m\right|\).
\end{proof}

\begin{lemma}\label{theorem:38}
Let \(\alpha\) be a type. Assume equality on \(\alpha\) is decidable. Let \(m\) be a multiset of elements of \(\alpha\). Assume \(m\text{ has no repeated elements}\). Then \(\left|\operatorname{Multiset\mathord{:}to\,Finset}\left(m\right)\right| = \left|m\right|\).
\end{lemma}

\begin{proof}
By \paperlink{premise:4d756c74697365742e64656475705f65715f73656c66}{Prerequisite~P\ref{premise:4d756c74697365742e64656475705f65715f73656c66}}, we obtain \(\operatorname{Multiset\mathord{:}dedup}\left(m\right) = m \iff m\text{ has no repeated elements}\).
Use the hypothesis \(m\text{ has no repeated elements}\).
By the reverse implication, we obtain \(\operatorname{Multiset\mathord{:}dedup}\left(m\right) = m\).
By \paperlink{premise:636f6e67725f617267}{\emph{Alias}} (\paperlink{premise:636f6e67725f617267}{Prerequisite~P\ref{premise:636f6e67725f617267}}) with \(\operatorname{Multiset}(\alpha)\) as its type parameter \(\alpha\) and \(\beta = \mathbb{N}\) and \(f = \operatorname{Multiset\mathord{:}card}\), we obtain \(\left|\operatorname{Multiset\mathord{:}to\,Finset}\left(m\right)\right| = \left|m\right|\).
\end{proof}

\begin{lemma}\label{theorem:39}
Let \(\alpha\) be a type. Assume equality on \(\alpha\) is decidable. Let \(m\) be a multiset of elements of \(\alpha\). Then
\begin{center}
\(\displaystyle |\operatorname{Multiset\mathord{:}dedup}(m)| = |m| \iff m\text{ has no repeated elements}.\)
\end{center}
\end{lemma}

\begin{proof}
Consider the following case.
\begin{proofstep}
Assume \(\left|\operatorname{Multiset\mathord{:}dedup}\left(m\right)\right| = \left|m\right|\).
By \paperlink{premise:4d756c74697365742e64656475705f6c65}{Prerequisite~P\ref{premise:4d756c74697365742e64656475705f6c65}} with \(s = m\), we obtain \(\operatorname{Multiset\mathord{:}dedup}\left(m\right) \le m\).
Use the hypothesis \(\left|\operatorname{Multiset\mathord{:}dedup}\left(m\right)\right| = \left|m\right|\).
By \paperlink{premise:45712e6765}{Prerequisite~P\ref{premise:45712e6765}} with \(\mathbb{N}\) as its type parameter \(\alpha\), we obtain \(\left|m\right| \le \left|\operatorname{Multiset\mathord{:}dedup}\left(m\right)\right|\).
By \paperlink{premise:4d756c74697365742e65715f6f665f6c655f6f665f636172645f6c65}{Prerequisite~P\ref{premise:4d756c74697365742e65715f6f665f6c655f6f665f636172645f6c65}}, we obtain \(\operatorname{Multiset\mathord{:}dedup}\left(m\right) = m\).
\end{proofstep}
By \paperlink{premise:636f6e67725f617267}{\emph{Alias}} (\paperlink{premise:636f6e67725f617267}{Prerequisite~P\ref{premise:636f6e67725f617267}}) with \(\operatorname{Multiset}(\alpha)\) as its type parameter \(\alpha\) and \(\beta = \mathbb{N}\) and \(f = \operatorname{Multiset\mathord{:}card}\), we obtain \(\operatorname{Multiset\mathord{:}dedup}(m) = m \implies |\operatorname{Multiset\mathord{:}dedup}(m)| = |m|\).
Together the two implications give the equivalence. Hence \(|\operatorname{Multiset\mathord{:}dedup}(m)| = |m| \iff \operatorname{Multiset\mathord{:}dedup}(m) = m\).
By \paperlink{premise:4d756c74697365742e64656475705f65715f73656c66}{Prerequisite~P\ref{premise:4d756c74697365742e64656475705f65715f73656c66}}, we obtain \(\operatorname{Multiset\mathord{:}dedup}\left(m\right) = m \iff m\text{ has no repeated elements}\).
Chaining the equivalences, we obtain \(|\operatorname{Multiset\mathord{:}dedup}(m)| = |m| \iff m\text{ has no repeated elements}\).
\end{proof}

\begin{corollary}\label{theorem:40}
Let \(\alpha\) be a type. Assume equality on \(\alpha\) is decidable. Let \(m\) be a multiset of elements of \(\alpha\). Then
\begin{center}
\(\displaystyle |\operatorname{Multiset\mathord{:}to\,Finset}(m)| = |m| \iff m\text{ has no repeated elements}.\)
\end{center}
\end{corollary}

\begin{proof}
Apply \paperlink{theorem:39}{Lemma~\ref{theorem:39}}.
\end{proof}

\begin{theorem}\label{theorem:41}
Let \(\alpha\) be a type. Assume equality on \(\alpha\) is decidable. Let \(l\) be a list of elements of \(\alpha\). Then
\begin{center}
\(\displaystyle |\operatorname{List\mathord{:}to\,Finset}(l)| = \operatorname{length}(\operatorname{List\mathord{:}dedup}(l)).\)
\end{center}
\end{theorem}

\begin{proof}
By definition, we obtain \(|\operatorname{List\mathord{:}to\,Finset}(l)| = \operatorname{length}(\operatorname{List\mathord{:}dedup}(l))\).
\end{proof}

\begin{corollary}\label{theorem:42}
Let \(\alpha\) be a type. Assume equality on \(\alpha\) is decidable. Let \(l\) be a list of elements of \(\alpha\). Then
\begin{center}
\(\displaystyle |\operatorname{List\mathord{:}to\,Finset}(l)| \le \operatorname{length}(l).\)
\end{center}
\end{corollary}

\begin{proof}
Apply \paperlink{theorem:37}{Lemma~\ref{theorem:37}} with \(m = \operatorname{Quotient\mathord{:}mk}\left(\operatorname{is\,Setoid}\left(\alpha\right), l\right)\).
\end{proof}

\begin{corollary}\label{theorem:43}
Let \(\alpha\) be a type. Assume equality on \(\alpha\) is decidable. Let \(l\) be a list of elements of \(\alpha\). Assume \(l\text{ has no repeated elements}\). Then
\begin{center}
\(\displaystyle \left|\operatorname{List\mathord{:}to\,Finset}\left(l\right)\right| = \operatorname{length}(l).\)
\end{center}
\end{corollary}

\begin{proof}
Apply \paperlink{theorem:38}{Lemma~\ref{theorem:38}}.
\end{proof}

\section{Finite sets}
\begin{theorem}\label{theorem:44}
Let \(\alpha\) be a type. Let \(s\) be a finite subset of \(\alpha\). Then
\begin{center}
\(\displaystyle \operatorname{length}(\operatorname{Finset\mathord{:}to\,List}\left(s\right)) = \left|s\right|.\)
\end{center}
\end{theorem}

\begin{proof}
Rewriting with \paperlink{theorem:1}{Lemma~\ref{theorem:1}}, \paperlink{premise:4d756c74697365742e636f655f746f4c697374}{Prerequisite~P\ref{premise:4d756c74697365742e636f655f746f4c697374}} and \paperlink{premise:4d756c74697365742e636f655f63617264}{Prerequisite~P\ref{premise:4d756c74697365742e636f655f63617264}} gives \(\operatorname{length}(\operatorname{Finset\mathord{:}to\,List}\left(s\right)) = \left|s\right|\).
\end{proof}

\begin{corollary}\label{theorem:45}
Let \(\alpha\) and \(\beta\) be types. Let \(s\) be a finite subset of \(\alpha\). Let \(f\) be a function of type \(\alpha \to \beta\). Assume equality on \(\beta\) is decidable. Then \(\left|f\left[s\right]\right| \le \left|s\right|\).
\end{corollary}

\begin{proof}
By \paperlink{premise:4d756c74697365742e636172645f6d6170}{Prerequisite~P\ref{premise:4d756c74697365742e636172645f6d6170}} with \(\operatorname{underlying\,multiset}(s)\) in place of its parameter \(s\), we obtain \(|\operatorname{Multiset\mathord{:}map}(f,\allowbreak{} \operatorname{underlying\,multiset}(s))| = |\operatorname{underlying\,multiset}(s)|\).
Applying the same function to equal arguments, we obtain \((|\operatorname{Multiset\mathord{:}to\,Finset}(\operatorname{Multiset\mathord{:}map}(f,\allowbreak{} \operatorname{underlying\,multiset}(s)))| \le |\operatorname{Multiset\mathord{:}map}(f,\allowbreak{} \operatorname{underlying\,multiset}(s))|) = (|\operatorname{Multiset\mathord{:}to\,Finset}(\operatorname{Multiset\mathord{:}map}(f,\allowbreak{} \operatorname{underlying\,multiset}(s)))| \le |\operatorname{underlying\,multiset}(s)|)\).
By \paperlink{theorem:37}{Lemma~\ref{theorem:37}} with \(m = \operatorname{Multiset\mathord{:}map}\left(f, \operatorname{underlying\,multiset}(s)\right)\), we obtain \(|\operatorname{Multiset\mathord{:}to\,Finset}(\operatorname{Multiset\mathord{:}map}(f,\allowbreak{} \operatorname{underlying\,multiset}(s)))| \le |\operatorname{Multiset\mathord{:}map}(f,\allowbreak{} \operatorname{underlying\,multiset}(s))|\).
Substituting along the equality, we obtain \(\left|f\left[s\right]\right| \le \left|s\right|\).
\end{proof}

\begin{lemma}\label{theorem:46}
Let \(\alpha\) and \(\beta\) be types. Let \(s\) be a finite subset of \(\alpha\). Let \(f\) be a function of type \(\alpha \to \beta\). Assume equality on \(\beta\) is decidable. Assume \(f\text{ is injective on }\operatorname{underlying\,set}(s)\). Then \(\left|f\left[s\right]\right| = \left|s\right|\).
\end{lemma}

\begin{proof}
Use the hypothesis \(f\text{ is injective on }\operatorname{underlying\,set}(s)\).
By \paperlink{premise:46696e7365742e696d6167655f76616c5f6f665f696e6a4f6e}{Prerequisite~P\ref{premise:46696e7365742e696d6167655f76616c5f6f665f696e6a4f6e}}, we obtain \(\operatorname{underlying\,multiset}(f[s]) = \operatorname{Multiset\mathord{:}map}(f,\allowbreak{} \operatorname{underlying\,multiset}(s))\).
Applying the same function to equal arguments, we obtain \(|\operatorname{underlying\,multiset}(f[s])| = |\operatorname{Multiset\mathord{:}map}(f,\allowbreak{} \operatorname{underlying\,multiset}(s))|\).
By \paperlink{premise:4d756c74697365742e636172645f6d6170}{Prerequisite~P\ref{premise:4d756c74697365742e636172645f6d6170}} with \(\operatorname{underlying\,multiset}(s)\) in place of its parameter \(s\), we obtain \(|\operatorname{Multiset\mathord{:}map}(f,\allowbreak{} \operatorname{underlying\,multiset}(s))| = |\operatorname{underlying\,multiset}(s)|\).
Chaining the equalities, we obtain \(|\operatorname{underlying\,multiset}(f[s])| = |\operatorname{underlying\,multiset}(s)|\).
Applying the same function to equal arguments, we obtain \((|\operatorname{underlying\,multiset}(f[s])| = |\operatorname{underlying\,multiset}(s)|) = (|\operatorname{underlying\,multiset}(s)| = |\operatorname{underlying\,multiset}(s)|)\).
By \paperlink{premise:65715f73656c66}{Prerequisite~P\ref{premise:65715f73656c66}} with \(\mathbb{N}\) as its type parameter \(\alpha\) and \(a = \left|\operatorname{underlying\,multiset}(s)\right|\), we obtain \((|\operatorname{underlying\,multiset}(s)| = |\operatorname{underlying\,multiset}(s)|) = \mathrm{True}\).
Chaining the equalities, we obtain \((|\operatorname{underlying\,multiset}(f[s])| = |\operatorname{underlying\,multiset}(s)|) = \mathrm{True}\).
By \paperlink{premise:6f665f65715f74727565}{Prerequisite~P\ref{premise:6f665f65715f74727565}}, we obtain \(\left|f\left[s\right]\right| = \left|s\right|\).
\end{proof}

\begin{lemma}\label{theorem:47}
Let \(\alpha\) and \(\beta\) be types. Let \(s\) be a finite subset of \(\alpha\). Let \(f\) be a function of type \(\alpha \to \beta\). Assume equality on \(\beta\) is decidable. Assume \(\left|f\left[s\right]\right| = \left|s\right|\). Then \(f\text{ is injective on }\operatorname{underlying\,set}(s)\).
\end{lemma}

\begin{proof}
By \paperlink{premise:4d756c74697365742e64656475705f6c65}{Prerequisite~P\ref{premise:4d756c74697365742e64656475705f6c65}} with \(\operatorname{Multiset\mathord{:}map}\left(f, \operatorname{underlying\,multiset}(s)\right)\) in place of its parameter \(s\), we obtain \(\operatorname{Multiset\mathord{:}dedup}(\operatorname{Multiset\mathord{:}map}(f,\allowbreak{} \operatorname{underlying\,multiset}(s))) \le \operatorname{Multiset\mathord{:}map}(f,\allowbreak{} \operatorname{underlying\,multiset}(s))\).
Rewriting with \paperlink{premise:4d756c74697365742e636172645f6d6170}{Prerequisite~P\ref{premise:4d756c74697365742e636172645f6d6170}}, \paperlink{theorem:1}{Lemma~\ref{theorem:1}}, \(\left|f\left[s\right]\right| = \left|s\right|\), substitution of equals and \paperlink{premise:6c655f72666c}{Prerequisite~P\ref{premise:6c655f72666c}} gives \(|\operatorname{Multiset\mathord{:}map}(f,\allowbreak{} \operatorname{underlying\,multiset}(s))| \le |\operatorname{Multiset\mathord{:}dedup}(\operatorname{Multiset\mathord{:}map}(f,\allowbreak{} \operatorname{underlying\,multiset}(s)))|\).
By \paperlink{premise:4d756c74697365742e65715f6f665f6c655f6f665f636172645f6c65}{Prerequisite~P\ref{premise:4d756c74697365742e65715f6f665f6c655f6f665f636172645f6c65}}, we obtain \(\operatorname{Multiset\mathord{:}dedup}(\operatorname{Multiset\mathord{:}map}(f,\allowbreak{} \operatorname{underlying\,multiset}(s))) = \operatorname{Multiset\mathord{:}map}(f,\allowbreak{} \operatorname{underlying\,multiset}(s))\).
Use the previously proved fact \(\operatorname{Multiset\mathord{:}dedup}(\operatorname{Multiset\mathord{:}map}(f,\allowbreak{} \operatorname{underlying\,multiset}(s))) = \operatorname{Multiset\mathord{:}map}(f,\allowbreak{} \operatorname{underlying\,multiset}(s))\).
Rewriting with \paperlink{premise:4d756c74697365742e64656475705f65715f73656c66}{Prerequisite~P\ref{premise:4d756c74697365742e64656475705f65715f73656c66}} in the preceding result gives \(\operatorname{Multiset\mathord{:}map}(f,\allowbreak{} \operatorname{underlying\,multiset}(s))\text{ has no repeated elements}\).
By \paperlink{premise:4d756c74697365742e696e6a5f6f6e5f6f665f6e6f6475705f6d6170}{Prerequisite~P\ref{premise:4d756c74697365742e696e6a5f6f6e5f6f665f6e6f6475705f6d6170}}, we obtain \(f\text{ is injective on }\operatorname{underlying\,set}(s)\).
\end{proof}

\begin{corollary}\label{theorem:48}
Let \(\alpha\) and \(\beta\) be types. Let \(s\) be a finite subset of \(\alpha\). Let \(f\) be a function of type \(\alpha \to \beta\). Assume equality on \(\beta\) is decidable. Then
\begin{center}
\(\displaystyle |f[s]| = |s| \iff f\text{ is injective on }\operatorname{underlying\,set}(s).\)
\end{center}
\end{corollary}

\begin{proof}
By \paperlink{theorem:47}{Lemma~\ref{theorem:47}}, we obtain \(|f[s]| = |s| \implies f\text{ is injective on }\operatorname{underlying\,set}(s)\).
By \paperlink{theorem:46}{Lemma~\ref{theorem:46}}, we obtain \(f\text{ is injective on }\operatorname{underlying\,set}(s) \implies |f[s]| = |s|\).
Together the two implications give the equivalence. Hence \(\left|f\left[s\right]\right| = \left|s\right| \iff f\text{ is injective on }\operatorname{underlying\,set}(s)\).
\end{proof}

\begin{lemma}\label{theorem:49}
Let \(\alpha\) and \(\beta\) be types. Let \(f\) be a function of type \(\alpha \to \beta\). Assume equality on \(\beta\) is decidable. Let \(s\) be a finite subset of \(\alpha\). Assume \(f\text{ is injective}\). Then \(\left|f\left[s\right]\right| = \left|s\right|\).
\end{lemma}

\begin{proof}
For the required implication, argue as follows.
\begin{proofstep}
Fix \(x\) in \(\alpha\).
Assume \(x \in \operatorname{underlying\,set}(s)\).
Fix \(x_{2}\) in \(\alpha\).
Assume \(x_{2} \in \operatorname{underlying\,set}(s)\).
Assume \(f\left(x\right) = f\left(x_{2}\right)\).
Use the hypothesis \(f\left(x\right) = f\left(x_{2}\right)\).
Applying the hypothesis \(f\text{ is injective}\), we obtain \(x = x_{2}\).
\end{proofstep}
Apply \paperlink{theorem:46}{Lemma~\ref{theorem:46}}.
\end{proof}

\begin{theorem}\label{theorem:50}
Let \(\alpha\) and \(\beta\) be types. Let \(s\) be a finite subset of \(\alpha\). Let \(f\) be a function of type \(\alpha \to \beta\). Assume equality on \(\beta\) is decidable. Let \(y\) be an element of \(\beta\). Then \(\left|\{x\in s \mid f\left(x\right) = y\}\right| \ne 0 \iff y \in f\left[s\right]\).
\end{theorem}

\begin{proof}
Rewriting with \paperlink{premise:46696e7365742e66696265725f6e6f6e656d7074795f6966665f6d656d5f696d616765}{Prerequisite~P\ref{premise:46696e7365742e66696265725f6e6f6e656d7074795f6966665f6d656d5f696d616765}}, \paperlink{theorem:9}{Lemma~\ref{theorem:9}} and \paperlink{premise:4e61742e706f735f6966665f6e655f7a65726f}{Prerequisite~P\ref{premise:4e61742e706f735f6966665f6e655f7a65726f}} gives \(\left|\{x\in s \mid f\left(x\right) = y\}\right| \ne 0 \iff y \in f\left[s\right]\).
\end{proof}

\begin{theorem}\label{theorem:51}
Let \(\alpha\) be a type. Let \(s\) be a finite subset of \(\alpha\). Let \(P\) be a function of type \(\alpha \to \mathsf{Prop}\). Assume the predicate \(P\) is decidable. Let \(n\) be a natural number. Then
\begin{center}
\(\displaystyle |\{x\in s \mid P(x)\}| \le n \iff (\forall s^{'}\, :\, \mathcal{P}_{\mathrm{fin}}(\alpha),\; s^{'} \subseteq s \implies n < |s^{'}| \implies (\exists a_{2}\, :\, \alpha,\; a_{2} \in s^{'} \land \neg P(a_{2}))).\)
\end{center}
\end{theorem}

\begin{proof}
By \paperlink{premise:4d756c74697365742e636172645f66696c7465725f6c655f696666}{Prerequisite~P\ref{premise:4d756c74697365742e636172645f66696c7465725f6c655f696666}} with \(\operatorname{underlying\,multiset}(s)\) in place of its parameter \(s\), we obtain \(|\operatorname{Multiset\mathord{:}filter}(P,\allowbreak{} \operatorname{underlying\,multiset}(s))| \le n \iff (\forall s^{'}\, :\, \operatorname{Multiset}(\alpha),\; s^{'} \le \operatorname{underlying\,multiset}(s) \implies n < |s^{'}| \implies (\exists a_{2}\, :\, \alpha,\; a_{2} \in s^{'} \land \neg P(a_{2})))\).
Consider the following case.
\begin{proofstep}
Assume \(\forall s^{'}\, :\, \operatorname{Multiset}(\alpha),\; s^{'} \le \operatorname{underlying\,multiset}(s) \implies n < |s^{'}| \implies (\exists a_{2}\, :\, \alpha,\; a_{2} \in s^{'} \land \neg P(a_{2}))\).
Fix \(s^{'}\) in \(\mathcal{P}_{\mathrm{fin}}(\alpha)\).
Assume \(s^{'} \subseteq s\).
Assume \(n < \left|s^{'}\right|\).
Rewriting with \paperlink{premise:46696e7365742e76616c5f6c655f696666}{Prerequisite~P\ref{premise:46696e7365742e76616c5f6c655f696666}} and \(s^{'} \subseteq s\) gives \(\operatorname{underlying\,multiset}(s^{'}) \le \operatorname{underlying\,multiset}(s)\).
Use the hypothesis \(n < \left|s^{'}\right|\).
Applying the hypothesis \(\forall s^{'}\, :\, \operatorname{Multiset}(\alpha),\; s^{'} \le \operatorname{underlying\,multiset}(s) \implies n < |s^{'}| \implies (\exists a_{2}\, :\, \alpha,\; a_{2} \in s^{'} \land \neg P(a_{2}))\), we obtain \(\exists a_{2}\, :\, \alpha,\; a_{2} \in s^{'} \land \neg P\left(a_{2}\right)\).
\end{proofstep}
Consider the following case.
\begin{proofstep}
Assume \(\forall s^{'}\, :\, \mathcal{P}_{\mathrm{fin}}(\alpha),\; s^{'} \subseteq s \implies n < |s^{'}| \implies (\exists a_{2}\, :\, \alpha,\; a_{2} \in s^{'} \land \neg P(a_{2}))\).
Fix \(s^{'}\) in \(\operatorname{Multiset}(\alpha)\).
Assume \(s^{'} \le \operatorname{underlying\,multiset}(s)\).
Assume \(n < \left|s^{'}\right|\).
For the required implication, argue as follows.
\begin{proofstep}
Fix \(x\) in \(\alpha\).
Assume \(x \in \operatorname{as\,finite\,set}(s^{'})\).
Use the hypothesis \(s^{'} \le \operatorname{underlying\,multiset}(s)\).
Use the hypothesis \(x \in \operatorname{as\,finite\,set}(s^{'})\).
By \paperlink{premise:4d756c74697365742e7375627365745f6f665f6c65}{Prerequisite~P\ref{premise:4d756c74697365742e7375627365745f6f665f6c65}}, we obtain \(x \in \operatorname{underlying\,multiset}(s)\).
\end{proofstep}
Use the hypothesis \(n < \left|s^{'}\right|\).
Applying the hypothesis \(\forall s^{'}\, :\, \mathcal{P}_{\mathrm{fin}}(\alpha),\; s^{'} \subseteq s \implies n < |s^{'}| \implies (\exists a_{2}\, :\, \alpha,\; a_{2} \in s^{'} \land \neg P(a_{2}))\), we obtain \(\exists a_{2}\, :\, \alpha,\; a_{2} \in s^{'} \land \neg P\left(a_{2}\right)\).
\end{proofstep}
Together the two implications give the equivalence. Hence \((\forall s^{'}\, :\, \operatorname{Multiset}(\alpha),\; s^{'} \le \operatorname{underlying\,multiset}(s) \implies n < |s^{'}| \implies (\exists a_{2}\, :\, \alpha,\; a_{2} \in s^{'} \land \neg P(a_{2}))) \iff (\forall s^{'}\, :\, \mathcal{P}_{\mathrm{fin}}(\alpha),\; s^{'} \subseteq s \implies n < |s^{'}| \implies (\exists a_{2}\, :\, \alpha,\; a_{2} \in s^{'} \land \neg P(a_{2})))\).
Chaining the equivalences, we obtain \(|\{x\in s \mid P(x)\}| \le n \iff (\forall s^{'}\, :\, \mathcal{P}_{\mathrm{fin}}(\alpha),\; s^{'} \subseteq s \implies n < |s^{'}| \implies (\exists a_{2}\, :\, \alpha,\; a_{2} \in s^{'} \land \neg P(a_{2})))\).
\end{proof}

\begin{lemma}\label{theorem:52}
Let \(\alpha\) and \(\beta\) be types. Let \(s\) be a finite subset of \(\alpha\). Let \(f\) be an element of \(\operatorname{Embedding}\left(\alpha, \beta\right)\). Then \(\left|\operatorname{Finset\mathord{:}map}\left(f, s\right)\right| = \left|s\right|\).
\end{lemma}

\begin{proof}
By \paperlink{premise:4d756c74697365742e636172645f6d6170}{Prerequisite~P\ref{premise:4d756c74697365742e636172645f6d6170}} with \(\operatorname{coe}\left(f\right)\) in place of its parameter \(f\) and \(\operatorname{underlying\,multiset}(s)\) in place of its parameter \(s\), we obtain \(\left|\operatorname{Finset\mathord{:}map}\left(f, s\right)\right| = \left|s\right|\).
\end{proof}

\begin{theorem}\label{theorem:53}
Let \(\alpha\) be a type. Let \(p\) be a function of type \(\alpha \to \mathsf{Prop}\). Assume the predicate \(p\) is decidable. Let \(s\) be a finite subset of \(\alpha\). Then \(\left|\operatorname{subtype}\left(p, s\right)\right| = \left|\{x\in s \mid p(x)\}\right|\).
\end{theorem}

\begin{proof}
By \paperlink{theorem:52}{Lemma~\ref{theorem:52}} with \(\{x\in s \mid p(x)\}\) as its type parameter \(\alpha\) and \(\beta = \{x\in \alpha \mid p(x)\}\) and \(f = \operatorname{Function\mathord{:}Embedding\mathord{:}mk}(x \mapsto (x\, :\, \{x_{1}\in \alpha \mid p(x_{1})\}))\), we obtain \(|\operatorname{Finset\mathord{:}map}(\operatorname{Function\mathord{:}Embedding\mathord{:}mk}(x \mapsto (x\, :\, \{x_{1}\in \alpha \mid p(x_{1})\})),\allowbreak{} \operatorname{Finset\mathord{:}attach}(\{x\in s \mid p(x)\}))| = |\operatorname{Finset\mathord{:}attach}(\{x\in s \mid p(x)\})|\).
By \paperlink{theorem:35}{Lemma~\ref{theorem:35}}, we obtain \(|\operatorname{Finset\mathord{:}attach}(\{x\in s \mid p(x)\})| = |\{x\in s \mid p(x)\}|\).
Chaining the equalities, we obtain \(|\operatorname{Finset\mathord{:}map}(\operatorname{Function\mathord{:}Embedding\mathord{:}mk}(x \mapsto (x\, :\, \{x_{1}\in \alpha \mid p(x_{1})\})),\allowbreak{} \operatorname{Finset\mathord{:}attach}(\{x\in s \mid p(x)\}))| = |\{x\in s \mid p(x)\}|\).
Applying the same function to equal arguments, we obtain \((|\operatorname{Finset\mathord{:}map}(\operatorname{Function\mathord{:}Embedding\mathord{:}mk}(x \mapsto (x\, :\, \{x_{1}\in \alpha \mid p(x_{1})\})),\allowbreak{} \operatorname{Finset\mathord{:}attach}(\{x\in s \mid p(x)\}))| = |\{x\in s \mid p(x)\}|) = (|\{x\in s \mid p(x)\}| = |\{x\in s \mid p(x)\}|)\).
By \paperlink{premise:65715f73656c66}{Prerequisite~P\ref{premise:65715f73656c66}} with \(\mathbb{N}\) as its type parameter \(\alpha\) and \(a = \left|\{x\in s \mid p(x)\}\right|\), we obtain \(\left(\left|\{x\in s \mid p(x)\}\right| = \left|\{x\in s \mid p(x)\}\right|\right) = \mathrm{True}\).
Chaining the equalities, we obtain \((|\operatorname{Finset\mathord{:}map}(\operatorname{Function\mathord{:}Embedding\mathord{:}mk}(x \mapsto (x\, :\, \{x_{1}\in \alpha \mid p(x_{1})\})),\allowbreak{} \operatorname{Finset\mathord{:}attach}(\{x\in s \mid p(x)\}))| = |\{x\in s \mid p(x)\}|) = \mathrm{True}\).
By \paperlink{premise:6f665f65715f74727565}{Prerequisite~P\ref{premise:6f665f65715f74727565}}, we obtain \(\left|\operatorname{subtype}\left(p, s\right)\right| = \left|\{x\in s \mid p(x)\}\right|\).
\end{proof}

\begin{lemma}\label{theorem:54}
Let \(\alpha\) be a type. Let \(s\) be a finite subset of \(\alpha\). Let \(p\) be a function of type \(\alpha \to \mathsf{Prop}\). Assume the predicate \(p\) is decidable. Then \(\left|\{x\in s \mid p(x)\}\right| \le \left|s\right|\).
\end{lemma}

\begin{proof}
By \paperlink{premise:46696e7365742e66696c7465725f737562736574}{Prerequisite~P\ref{premise:46696e7365742e66696c7465725f737562736574}}, we obtain \(\{x\in s \mid p(x)\} \subseteq s\).
Apply \paperlink{theorem:5}{Lemma~\ref{theorem:5}}.
\end{proof}

\begin{lemma}\label{theorem:55}
Let \(\alpha\) be a type. Let \(s\) and \(t\) be finite subsets of \(\alpha\). Assume \(s \subseteq t\). Assume \(\left|t\right| \le \left|s\right|\). Then \(s = t\).
\end{lemma}

\begin{proof}
By \paperlink{premise:46696e7365742e76616c5f6c655f696666}{Prerequisite~P\ref{premise:46696e7365742e76616c5f6c655f696666}}, we obtain \(\operatorname{underlying\,multiset}(s) \le \operatorname{underlying\,multiset}(t) \iff s \subseteq t\).
Use the hypothesis \(s \subseteq t\).
By the reverse implication, we obtain \(\operatorname{underlying\,multiset}(s) \le \operatorname{underlying\,multiset}(t)\).
Use the hypothesis \(\left|t\right| \le \left|s\right|\).
By \paperlink{premise:4d756c74697365742e65715f6f665f6c655f6f665f636172645f6c65}{Prerequisite~P\ref{premise:4d756c74697365742e65715f6f665f6c655f6f665f636172645f6c65}}, we obtain \(\operatorname{underlying\,multiset}(s) = \operatorname{underlying\,multiset}(t)\).
By \paperlink{premise:46696e7365742e65715f6f665f766571}{Prerequisite~P\ref{premise:46696e7365742e65715f6f665f766571}}, we obtain \(s = t\).
\end{proof}

\begin{lemma}\label{theorem:56}
Let \(\alpha\) be a type. Let \(s\) and \(t\) be finite subsets of \(\alpha\). Assume \(s \subseteq t\). Then \(\left|t\right| \le \left|s\right| \iff s = t\).
\end{lemma}

\begin{proof}
By \paperlink{theorem:55}{Lemma~\ref{theorem:55}}, we obtain \(\left|t\right| \le \left|s\right| \implies s = t\).
Consider the following case.
\begin{proofstep}
Assume \(s = t\).
Use the hypothesis \(s = t\).
By \paperlink{premise:636f6e67725f617267}{\emph{Alias}} (\paperlink{premise:636f6e67725f617267}{Prerequisite~P\ref{premise:636f6e67725f617267}}) with \(\mathcal{P}_{\mathrm{fin}}(\alpha)\) as its type parameter \(\alpha\) and \(\beta = \mathbb{N}\) and \(f = \operatorname{Finset\mathord{:}card}\), we obtain \(\left|s\right| = \left|t\right|\).
By \paperlink{premise:67655f6f665f6571}{Prerequisite~P\ref{premise:67655f6f665f6571}} with \(\mathbb{N}\) as its type parameter \(\alpha\), we obtain \(\left|t\right| \le \left|s\right|\).
\end{proofstep}
Together the two implications give the equivalence. Hence \(\left|t\right| \le \left|s\right| \iff s = t\).
\end{proof}

\begin{corollary}\label{theorem:57}
Let \(\alpha\) be a type. Let \(s\) and \(t\) be finite subsets of \(\alpha\). Assume \(s \subseteq t\). Assume \(\left|t\right| \le \left|s\right|\). Then \(t = s\).
\end{corollary}

\begin{proof}
By \paperlink{theorem:55}{Lemma~\ref{theorem:55}}, we obtain \(s = t\).
Reversing the equality, we obtain \(t = s\).
\end{proof}

\begin{theorem}\label{theorem:58}
Let \(\alpha\) be a type. Let \(s\) and \(t\) be finite subsets of \(\alpha\). Assume \(s \subseteq t\). Then \(\left|t\right| \le \left|s\right| \iff t = s\).
\end{theorem}

\begin{proof}
By \paperlink{theorem:56}{Lemma~\ref{theorem:56}}, we obtain \(\left|t\right| \le \left|s\right| \iff s = t\).
By \paperlink{premise:65715f636f6d6d}{Prerequisite~P\ref{premise:65715f636f6d6d}} with \(\mathcal{P}_{\mathrm{fin}}(\alpha)\) as its type parameter \(\alpha\), we obtain \(s = t \iff t = s\).
Chaining the equivalences, we obtain \(\left|t\right| \le \left|s\right| \iff t = s\).
\end{proof}

\begin{theorem}\label{theorem:59}
Let \(\alpha\) be a type. Let \(s\) and \(t\) be finite subsets of \(\alpha\). Assume \(\left|t\right| \le \left|s\right|\). Then \(s \subseteq t \iff s = t\).
\end{theorem}

\begin{proof}
Consider the following case.
\begin{proofstep}
Assume \(s \subseteq t\).
By \paperlink{theorem:55}{Lemma~\ref{theorem:55}}, we obtain \(s = t\).
\end{proofstep}
By \paperlink{premise:46696e7365742e696e737449735265666c537562736574}{Prerequisite~P\ref{premise:46696e7365742e696e737449735265666c537562736574}}, we obtain \((\mathit{x1} \mapsto \mathit{x2} \mapsto \mathit{x1} \subseteq \mathit{x2})\text{ is reflexive on }\mathcal{P}_{\mathrm{fin}}(\alpha)\).
By \paperlink{premise:45712e73756273657427}{\emph{Alias}} (\paperlink{premise:45712e73756273657427}{Prerequisite~P\ref{premise:45712e73756273657427}}) with \(\mathcal{P}_{\mathrm{fin}}(\alpha)\) as its type parameter \(\alpha\), we obtain \(s = t \implies s \subseteq t\).
Together the two implications give the equivalence. Hence \(s \subseteq t \iff s = t\).
\end{proof}

\begin{theorem}\label{theorem:60}
Let \(\alpha\) be a type. Let \(s\) be a finite subset of \(\alpha\). Let \(f\) be an element of \(\operatorname{Embedding}\left(\alpha, \alpha\right)\). Assume \(\operatorname{Finset\mathord{:}map}\left(f, s\right) \subseteq s\). Then \(\operatorname{Finset\mathord{:}map}\left(f, s\right) = s\).
\end{theorem}

\begin{proof}
By \paperlink{theorem:52}{Lemma~\ref{theorem:52}}, we obtain \(\left|\operatorname{Finset\mathord{:}map}\left(f, s\right)\right| = \left|s\right|\).
By \paperlink{premise:45712e6765}{Prerequisite~P\ref{premise:45712e6765}} with \(\mathbb{N}\) as its type parameter \(\alpha\), we obtain \(\left|s\right| \le \left|\operatorname{Finset\mathord{:}map}\left(f, s\right)\right|\).
Apply \paperlink{theorem:55}{Lemma~\ref{theorem:55}}.
\end{proof}

\begin{theorem}\label{theorem:61}
Let \(\alpha\) be a type. Let \(s\) be a finite subset of \(\alpha\). Let \(p\) be a function of type \(\alpha \to \mathsf{Prop}\). Assume the predicate \(p\) is decidable. Then
\begin{center}
\(\displaystyle |\{x\in s \mid p(x)\}| = |s| \iff (\forall x\, :\, \alpha,\; x \in s \implies p(x)).\)
\end{center}
\end{theorem}

\begin{proof}
Rewriting with \paperlink{premise:46696e7365742e66696c7465725f65715f73656c66}{Prerequisite~P\ref{premise:46696e7365742e66696c7465725f65715f73656c66}}, \paperlink{theorem:56}{Lemma~\ref{theorem:56}}, \paperlink{premise:46696e7365742e66696c7465725f737562736574}{Prerequisite~P\ref{premise:46696e7365742e66696c7465725f737562736574}}, \paperlink{premise:6e6f745f6c74}{Prerequisite~P\ref{premise:6e6f745f6c74}}, \paperlink{premise:4c452e6c652e65715f6966665f6e6f745f6c74}{\emph{Alias}} (\paperlink{premise:4c452e6c652e65715f6966665f6e6f745f6c74}{Prerequisite~P\ref{premise:4c452e6c652e65715f6966665f6e6f745f6c74}}) and \paperlink{theorem:54}{Lemma~\ref{theorem:54}} gives \(|\{x\in s \mid p(x)\}| = |s| \iff (\forall x\, :\, \alpha,\; x \in s \implies p(x))\).
\end{proof}

\begin{theorem}[{Alias}]\label{theorem:62}
Let \(\alpha\) be a type. Let \(s\) be a finite subset of \(\alpha\). Let \(p\) be a function of type \(\alpha \to \mathsf{Prop}\). Assume the predicate \(p\) is decidable. Assume \(\left|\{x\in s \mid p(x)\}\right| = \left|s\right|\). Let \(x\) be an element of \(\alpha\). Assume \(x \in s\). Then \(p\left(x\right)\).
\end{theorem}

\begin{proof}
Apply the forward implication of \paperlink{theorem:61}{Theorem~\ref{theorem:61}}.
\end{proof}

\begin{theorem}\label{theorem:63}
Let \(\alpha\) be a type. Let \(s\) be a finite subset of \(\alpha\). Let \(p\) be a function of type \(\alpha \to \mathsf{Prop}\). Assume the predicate \(p\) is decidable. Then
\begin{center}
\(\displaystyle |\{x\in s \mid p(x)\}| = 0 \iff (\forall x\, :\, \alpha,\; x \in s \implies \neg p(x)).\)
\end{center}
\end{theorem}

\begin{proof}
By \paperlink{theorem:7}{Lemma~\ref{theorem:7}}, we obtain \(\left|\{x\in s \mid p(x)\}\right| = 0 \iff \{x\in s \mid p(x)\} = \varnothing\).
By propositional extensionality, we obtain \(\left(\left|\{x\in s \mid p(x)\}\right| = 0\right) = \left(\{x\in s \mid p(x)\} = \varnothing\right)\).
Applying the same function to equal arguments, we obtain \((|\{x\in s \mid p(x)\}| = 0 \iff (\forall x\, :\, \alpha,\; x \in s \implies \neg p(x))) = (\{x\in s \mid p(x)\} = \varnothing \iff (\forall x\, :\, \alpha,\; x \in s \implies \neg p(x)))\).
By the preceding argument, we obtain \((|\{x\in s \mid p(x)\}| = 0 \iff (\forall x\, :\, \alpha,\; x \in s \implies \neg p(x))) = (\{x\in s \mid p(x)\} = \varnothing \iff (\forall x\, :\, \alpha,\; x \in s \implies \neg p(x)))\).
By \paperlink{premise:46696e7365742e66696c7465725f65715f656d7074795f696666}{Prerequisite~P\ref{premise:46696e7365742e66696c7465725f65715f656d7074795f696666}}, we obtain \(\{x\in s \mid p(x)\} = \varnothing \iff (\forall x\, :\, \alpha,\; x \in s \implies \neg p(x))\).
By propositional extensionality, we obtain \((\{x\in s \mid p(x)\} = \varnothing) = (\forall x\, :\, \alpha,\; x \in s \implies \neg p(x))\).
Applying the same function to equal arguments, we obtain \((\{x\in s \mid p(x)\} = \varnothing \iff (\forall x\, :\, \alpha,\; x \in s \implies \neg p(x))) = ((\forall x\, :\, \alpha,\; x \in s \implies \neg p(x)) \iff (\forall x\, :\, \alpha,\; x \in s \implies \neg p(x)))\).
By the preceding argument, we obtain \((\{x\in s \mid p(x)\} = \varnothing \iff (\forall x\, :\, \alpha,\; x \in s \implies \neg p(x))) = ((\forall x\, :\, \alpha,\; x \in s \implies \neg p(x)) \iff (\forall x\, :\, \alpha,\; x \in s \implies \neg p(x)))\).
By \paperlink{premise:4966662e72666c}{Prerequisite~P\ref{premise:4966662e72666c}}, we obtain \((\forall x\, :\, \alpha,\; x \in s \implies \neg p(x)) \iff (\forall x\, :\, \alpha,\; x \in s \implies \neg p(x))\).
Substituting along the equality, we obtain \(\{x\in s \mid p(x)\} = \varnothing \iff (\forall x\, :\, \alpha,\; x \in s \implies \neg p(x))\).
Substituting along the equality, we obtain \(|\{x\in s \mid p(x)\}| = 0 \iff (\forall x\, :\, \alpha,\; x \in s \implies \neg p(x))\).
\end{proof}

\begin{lemma}\label{theorem:64}
Let \(\alpha\) be a type. Let \(s\) and \(t\) be finite subsets of \(\alpha\). Assume \(s \subsetneq t\). Then \(\left|s\right| < \left|t\right|\).
\end{lemma}

\begin{proof}
By \paperlink{premise:46696e7365742e76616c5f6c745f696666}{Prerequisite~P\ref{premise:46696e7365742e76616c5f6c745f696666}}, we obtain \(\operatorname{underlying\,multiset}(s) < \operatorname{underlying\,multiset}(t) \iff s \subsetneq t\).
Use the hypothesis \(s \subsetneq t\).
By the reverse implication, we obtain \(\operatorname{underlying\,multiset}(s) < \operatorname{underlying\,multiset}(t)\).
By \paperlink{premise:4d756c74697365742e636172645f6c745f63617264}{Prerequisite~P\ref{premise:4d756c74697365742e636172645f6c745f63617264}}, we obtain \(\left|s\right| < \left|t\right|\).
\end{proof}

\begin{corollary}\label{theorem:65}
Let \(\alpha\) be a type. Then \(\operatorname{Strict\,Mono}\left(\operatorname{Finset\mathord{:}card}\right)\).
\end{corollary}

\begin{proof}
Fix \(x\) in \(\mathcal{P}_{\mathrm{fin}}(\alpha)\).
Fix \(x_{1}\) in \(\mathcal{P}_{\mathrm{fin}}(\alpha)\).
Apply \paperlink{theorem:64}{Lemma~\ref{theorem:64}}.
\end{proof}

\begin{theorem}\label{theorem:66}
Let \(\alpha\) be a type. Let \(s\) be a finite subset of \(\alpha\). Let \(n\) be a natural number. Let \(f\) be a function of type \(\Pi_{i\, :\, \mathbb{N}}\, i < n \to \alpha\). Assume \(\forall a\, :\, \alpha,\; a \in s \implies (\exists i\, :\, \mathbb{N},\; \exists h\text{ proving }(i < n),\; f(i,\allowbreak{} h) = a)\). Assume \(\forall i\, :\, \mathbb{N},\; \forall h_{1}\, :\, i < n,\; f\left(i, h_{1}\right) \in s\). Assume \(\forall i\, :\, \mathbb{N},\; \forall j\, :\, \mathbb{N},\; \forall \mathit{hi}\, :\, i < n,\; \forall \mathit{hj}\, :\, j < n,\; f(i,\allowbreak{} \mathit{hi}) = f(j,\allowbreak{} \mathit{hj}) \implies i = j\). Then \(\left|s\right| = n\).
\end{theorem}

\begin{proof}
For the required implication, argue as follows.
\begin{proofstep}
Fix \(a\) in \(\alpha\).
Consider the following case.
\begin{proofstep}
Assume \(a \in s\).
Use the hypothesis \(a \in s\).
Applying the hypothesis \(\forall a\, :\, \alpha,\; a \in s \implies (\exists i\, :\, \mathbb{N},\; \exists h\text{ proving }(i < n),\; f(i,\allowbreak{} h) = a)\), we obtain \(\exists i\, :\, \mathbb{N},\; \exists h_{1}\text{ proving }\left(i < n\right),\; f\left(i, h_{1}\right) = a\).
Consider the following case.
\begin{proofstep}
Fix \(i\) in \(\mathbb{N}\).
Assume \(\exists h_{2}\text{ proving }\left(i < n\right),\; f\left(i, h_{2}\right) = a\).
Use the hypothesis \(\exists h_{2}\text{ proving }\left(i < n\right),\; f\left(i, h_{2}\right) = a\).
Consider the following case.
\begin{proofstep}
Assume \(i < n\).
Assume \(f\left(i, \mathit{hi}\right) = a\).
For the required implication, argue as follows.
\begin{proofstep}
Assume \(f\left(i, \mathit{hi}\right) \in s\).
Choose the witness \(i\).
The construction requires evidence that \(i \in \{i_{1}\in\mathbb{N}\mid i_{1}<n\}\).
By \paperlink{premise:46696e7365742e6d656d5f72616e6765}{Prerequisite~P\ref{premise:46696e7365742e6d656d5f72616e6765}}, we obtain \(i \in \{i_{1}\in\mathbb{N}\mid i_{1}<n\} \iff i < n\).
Use the hypothesis \(i < n\).
By the reverse implication, we obtain \(i \in \{i_{1}\in\mathbb{N}\mid i_{1}<n\}\).
By definition, we obtain \(f(i,\allowbreak{} \operatorname{Iff\mathord{:}mp}(\operatorname{Finset\mathord{:}mem\,range},\allowbreak{} \operatorname{Iff\mathord{:}mpr}(\operatorname{Finset\mathord{:}mem\,range},\allowbreak{} \mathit{hi}))) = f(i,\allowbreak{} \operatorname{Iff\mathord{:}mp}(\operatorname{Finset\mathord{:}mem\,range},\allowbreak{} \operatorname{Iff\mathord{:}mpr}(\operatorname{Finset\mathord{:}mem\,range},\allowbreak{} \mathit{hi})))\).
This witness proves existence. Hence \(\exists \mathit{hi}_{1}\text{ proving }(i \in \{i_{1}\in\mathbb{N}\mid i_{1}<n\}),\; f(i,\allowbreak{} \operatorname{Iff\mathord{:}mp}(\operatorname{Finset\mathord{:}mem\,range},\allowbreak{} \mathit{hi}_{1})) = f(i,\allowbreak{} \mathit{hi})\).
This witness proves existence. Hence \(\exists i_{1}\, :\, \mathbb{N},\; \exists \mathit{hi}_{1}\text{ proving }(i_{1} \in \{i_{2}\in\mathbb{N}\mid i_{2}<n\}),\; f(i_{1},\allowbreak{} \operatorname{Iff\mathord{:}mp}(\operatorname{Finset\mathord{:}mem\,range},\allowbreak{} \mathit{hi}_{1})) = f(i,\allowbreak{} \mathit{hi})\).
\end{proofstep}
Use the hypothesis \(f\left(i, \mathit{hi}\right) = a\).
Use the hypothesis \(a \in s\).
Substituting equals, we obtain \(\exists i_{1}\, :\, \mathbb{N},\; \exists \mathit{hi}_{1}\text{ proving }(i_{1} \in \{i_{2}\in\mathbb{N}\mid i_{2}<n\}),\; f(i_{1},\allowbreak{} \operatorname{Iff\mathord{:}mp}(\operatorname{Finset\mathord{:}mem\,range},\allowbreak{} \mathit{hi}_{1})) = a\).
\end{proofstep}
Using the existential hypothesis, we obtain \(\exists i_{1}\, :\, \mathbb{N},\; \exists \mathit{hi}\text{ proving }(i_{1} \in \{i_{2}\in\mathbb{N}\mid i_{2}<n\}),\; f(i_{1},\allowbreak{} \operatorname{Iff\mathord{:}mp}(\operatorname{Finset\mathord{:}mem\,range},\allowbreak{} \mathit{hi})) = a\).
\end{proofstep}
Using the existential hypothesis, we obtain \(\exists i\, :\, \mathbb{N},\; \exists \mathit{hi}\text{ proving }(i \in \{i_{1}\in\mathbb{N}\mid i_{1}<n\}),\; f(i,\allowbreak{} \operatorname{Iff\mathord{:}mp}(\operatorname{Finset\mathord{:}mem\,range},\allowbreak{} \mathit{hi})) = a\).
\end{proofstep}
Consider the following case.
\begin{proofstep}
Assume \(\exists i\, :\, \mathbb{N},\; \exists \mathit{hi}\text{ proving }(i \in \{i_{1}\in\mathbb{N}\mid i_{1}<n\}),\; f(i,\allowbreak{} \operatorname{Iff\mathord{:}mp}(\operatorname{Finset\mathord{:}mem\,range},\allowbreak{} \mathit{hi})) = a\).
Use the hypothesis \(\exists i\, :\, \mathbb{N},\; \exists \mathit{hi}\text{ proving }(i \in \{i_{1}\in\mathbb{N}\mid i_{1}<n\}),\; f(i,\allowbreak{} \operatorname{Iff\mathord{:}mp}(\operatorname{Finset\mathord{:}mem\,range},\allowbreak{} \mathit{hi})) = a\).
Consider the following case.
\begin{proofstep}
Fix \(i\) in \(\mathbb{N}\).
Assume \(\exists \mathit{hi}\text{ proving }(i \in \{i_{1}\in\mathbb{N}\mid i_{1}<n\}),\; f(i,\allowbreak{} \operatorname{Iff\mathord{:}mp}(\operatorname{Finset\mathord{:}mem\,range},\allowbreak{} \mathit{hi})) = a\).
Use the hypothesis \(\exists \mathit{hi}\text{ proving }(i \in \{i_{1}\in\mathbb{N}\mid i_{1}<n\}),\; f(i,\allowbreak{} \operatorname{Iff\mathord{:}mp}(\operatorname{Finset\mathord{:}mem\,range},\allowbreak{} \mathit{hi})) = a\).
Consider the following case.
\begin{proofstep}
Assume \(i \in \{i_{1}\in\mathbb{N}\mid i_{1}<n\}\).
Assume \(f(i,\allowbreak{} \operatorname{Iff\mathord{:}mp}(\operatorname{Finset\mathord{:}mem\,range},\allowbreak{} \mathit{hi})) = a\).
By \paperlink{premise:46696e7365742e6d656d5f72616e6765}{Prerequisite~P\ref{premise:46696e7365742e6d656d5f72616e6765}}, we obtain \(i \in \{i_{1}\in\mathbb{N}\mid i_{1}<n\} \iff i < n\).
Use the hypothesis \(i \in \{i_{1}\in\mathbb{N}\mid i_{1}<n\}\).
By the forward implication, we obtain \(i < n\).
Applying the hypothesis \(\forall i\, :\, \mathbb{N},\; \forall h_{1}\, :\, i < n,\; f\left(i, h_{1}\right) \in s\), we obtain \(f(i,\allowbreak{} \operatorname{Iff\mathord{:}mp}(\operatorname{Finset\mathord{:}mem\,range},\allowbreak{} \mathit{hi})) \in s\).
Use the hypothesis \(f(i,\allowbreak{} \operatorname{Iff\mathord{:}mp}(\operatorname{Finset\mathord{:}mem\,range},\allowbreak{} \mathit{hi})) = a\).
Substituting equals, we obtain \(a \in s\).
\end{proofstep}
Using the existential hypothesis, we obtain \(a \in s\).
\end{proofstep}
Using the existential hypothesis, we obtain \(a \in s\).
\end{proofstep}
Together the two implications give the equivalence. Hence \(a \in s \iff (\exists i\, :\, \mathbb{N},\; \exists \mathit{hi}\text{ proving }(i \in \{i_{1}\in\mathbb{N}\mid i_{1}<n\}),\; f(i,\allowbreak{} \operatorname{Iff\mathord{:}mp}(\operatorname{Finset\mathord{:}mem\,range},\allowbreak{} \mathit{hi})) = a)\).
The substitution uses the following argument.
\begin{proofstep}
Fix \(a_{1}\) in \(\{i\in\mathbb{N}\mid i<n\}\).
By \paperlink{premise:46696e7365742e6d656d5f617474616368}{Prerequisite~P\ref{premise:46696e7365742e6d656d5f617474616368}} with \(\mathbb{N}\) as its type parameter \(\alpha\) and \(\{i\in\mathbb{N}\mid i<n\}\) in place of its parameter \(s\) and \(x = a_{1}\), we obtain \(a_{1} \in \operatorname{Finset\mathord{:}attach}\left(\{i\in\mathbb{N}\mid i<n\}\right)\).
By \paperlink{premise:65715f74727565}{Prerequisite~P\ref{premise:65715f74727565}}, we obtain \((a_{1} \in \operatorname{Finset\mathord{:}attach}(\{i\in\mathbb{N}\mid i<n\})) = \mathrm{True}\).
Applying the same function to equal arguments, we obtain \((a_{1} \in \operatorname{Finset\mathord{:}attach}(\{i\in\mathbb{N}\mid i<n\}) \land f(\operatorname{val}(a_{1}),\allowbreak{} \operatorname{Iff\mathord{:}mp}(\operatorname{Finset\mathord{:}mem\,range},\allowbreak{} \operatorname{property}(a_{1}))) = a) = (\mathrm{True} \land f(\operatorname{val}(a_{1}),\allowbreak{} \operatorname{Iff\mathord{:}mp}(\operatorname{Finset\mathord{:}mem\,range},\allowbreak{} \operatorname{property}(a_{1}))) = a)\).
By \paperlink{premise:747275655f616e64}{Prerequisite~P\ref{premise:747275655f616e64}} with \(p = f(\operatorname{val}(a_{1}),\allowbreak{} \operatorname{Iff\mathord{:}mp}(\operatorname{Finset\mathord{:}mem\,range},\allowbreak{} \operatorname{property}(a_{1}))) = a\), we obtain \((\mathrm{True} \land f(\operatorname{val}(a_{1}),\allowbreak{} \operatorname{Iff\mathord{:}mp}(\operatorname{Finset\mathord{:}mem\,range},\allowbreak{} \operatorname{property}(a_{1}))) = a) = (f(\operatorname{val}(a_{1}),\allowbreak{} \operatorname{Iff\mathord{:}mp}(\operatorname{Finset\mathord{:}mem\,range},\allowbreak{} \operatorname{property}(a_{1}))) = a)\).
Chaining the equalities, we obtain \((a_{1} \in \operatorname{Finset\mathord{:}attach}(\{i\in\mathbb{N}\mid i<n\}) \land f(\operatorname{val}(a_{1}),\allowbreak{} \operatorname{Iff\mathord{:}mp}(\operatorname{Finset\mathord{:}mem\,range},\allowbreak{} \operatorname{property}(a_{1}))) = a) = (f(\operatorname{val}(a_{1}),\allowbreak{} \operatorname{Iff\mathord{:}mp}(\operatorname{Finset\mathord{:}mem\,range},\allowbreak{} \operatorname{property}(a_{1}))) = a)\).
\end{proofstep}
The substitution uses the following argument.
\begin{proofstep}
Fix \(a_{1}\) in \(\mathbb{N}\).
For the required implication, argue as follows.
\begin{proofstep}
Assume \(a_{1} \in \{i\in\mathbb{N}\mid i<n\}\).
By definition, we obtain \((f(a_{1},\allowbreak{} \operatorname{Iff\mathord{:}mp}(\operatorname{Finset\mathord{:}mem\,range},\allowbreak{} \operatorname{property}((a_{1}\, :\, \{i\in\mathbb{N}\mid i<n\})))) = a) = (f(a_{1},\allowbreak{} \operatorname{Iff\mathord{:}mp}(\operatorname{Finset\mathord{:}mem\,range},\allowbreak{} \operatorname{property}((a_{1}\, :\, \{i\in\mathbb{N}\mid i<n\})))) = a)\).
By \paperlink{premise:4966662e6f665f6571}{Prerequisite~P\ref{premise:4966662e6f665f6571}}, we obtain \(f(a_{1},\allowbreak{} \operatorname{Iff\mathord{:}mp}(\operatorname{Finset\mathord{:}mem\,range},\allowbreak{} \operatorname{property}((a_{1}\, :\, \{i\in\mathbb{N}\mid i<n\})))) = a \iff f(a_{1},\allowbreak{} \operatorname{Iff\mathord{:}mp}(\operatorname{Finset\mathord{:}mem\,range},\allowbreak{} \operatorname{property}((a_{1}\, :\, \{i\in\mathbb{N}\mid i<n\})))) = a\).
\end{proofstep}
By definition, we obtain \(\left(a_{1} \in \{i\in\mathbb{N}\mid i<n\}\right) = \left(a_{1} \in \{i\in\mathbb{N}\mid i<n\}\right)\).
By \paperlink{premise:4966662e6f665f6571}{Prerequisite~P\ref{premise:4966662e6f665f6571}}, we obtain \(a_{1} \in \{i\in\mathbb{N}\mid i<n\} \iff a_{1} \in \{i\in\mathbb{N}\mid i<n\}\).
By \paperlink{premise:6578697374735f70726f705f636f6e6772}{Prerequisite~P\ref{premise:6578697374735f70726f705f636f6e6772}}, we obtain \((\exists b\text{ proving }(a_{1} \in \{i\in\mathbb{N}\mid i<n\}),\; f(a_{1},\allowbreak{} \operatorname{Iff\mathord{:}mp}(\operatorname{Finset\mathord{:}mem\,range},\allowbreak{} \operatorname{property}((a_{1}\, :\, \{i\in\mathbb{N}\mid i<n\})))) = a) \iff (\exists h_{1}\text{ proving }(a_{1} \in \{i\in\mathbb{N}\mid i<n\}),\; f(a_{1},\allowbreak{} \operatorname{Iff\mathord{:}mp}(\operatorname{Finset\mathord{:}mem\,range},\allowbreak{} \operatorname{property}((a_{1}\, :\, \{i\in\mathbb{N}\mid i<n\})))) = a)\).
By propositional extensionality, we obtain \((\exists b\text{ proving }(a_{1} \in \{i\in\mathbb{N}\mid i<n\}),\; f(a_{1},\allowbreak{} \operatorname{Iff\mathord{:}mp}(\operatorname{Finset\mathord{:}mem\,range},\allowbreak{} \operatorname{property}((a_{1}\, :\, \{i\in\mathbb{N}\mid i<n\})))) = a) = (\exists h_{1}\text{ proving }(a_{1} \in \{i\in\mathbb{N}\mid i<n\}),\; f(a_{1},\allowbreak{} \operatorname{Iff\mathord{:}mp}(\operatorname{Finset\mathord{:}mem\,range},\allowbreak{} \operatorname{property}((a_{1}\, :\, \{i\in\mathbb{N}\mid i<n\})))) = a)\).
\end{proofstep}
Use the previously proved fact \(a \in s \iff (\exists i\, :\, \mathbb{N},\; \exists \mathit{hi}\text{ proving }(i \in \{i_{1}\in\mathbb{N}\mid i_{1}<n\}),\; f(i,\allowbreak{} \operatorname{Iff\mathord{:}mp}(\operatorname{Finset\mathord{:}mem\,range},\allowbreak{} \mathit{hi})) = a)\), which simplifies to \(a \in s \iff (\exists a_{1}\, :\, \mathbb{N},\; \exists h_{1}\text{ proving }(a_{1} \in \{i\in\mathbb{N}\mid i<n\}),\; f(a_{1},\allowbreak{} \operatorname{Iff\mathord{:}mp}(\operatorname{Finset\mathord{:}mem\,range},\allowbreak{} \operatorname{property}((a_{1}\, :\, \{i\in\mathbb{N}\mid i<n\})))) = a)\).
Rewriting with \paperlink{premise:46696e7365742e6d656d5f696d616765}{Prerequisite~P\ref{premise:46696e7365742e6d656d5f696d616765}}, \paperlink{premise:66756e657874}{\emph{Function extensionality}} (\paperlink{premise:66756e657874}{Prerequisite~P\ref{premise:66756e657874}}), \paperlink{premise:537562747970652e657869737473}{Prerequisite~P\ref{premise:537562747970652e657869737473}} and the two identities established above in the preceding result gives \(a \in s \iff a \in (i \mapsto f(i,\allowbreak{} \operatorname{Iff\mathord{:}mp}(\operatorname{Finset\mathord{:}mem\,range},\allowbreak{} \operatorname{property}(i))))[\operatorname{Finset\mathord{:}attach}(\{i\in\mathbb{N}\mid i<n\})]\).
\end{proofstep}
By \paperlink{premise:46696e7365742e657874}{Prerequisite~P\ref{premise:46696e7365742e657874}}, we obtain \(s = (i \mapsto f(i,\allowbreak{} \operatorname{Iff\mathord{:}mp}(\operatorname{Finset\mathord{:}mem\,range},\allowbreak{} \operatorname{property}(i))))[\operatorname{Finset\mathord{:}attach}(\{i\in\mathbb{N}\mid i<n\})]\).
\par\smallskip\noindent\emph{Equality 1.}\enspace
Rewriting with \(s = (i \mapsto f(i,\allowbreak{} \operatorname{Iff\mathord{:}mp}(\operatorname{Finset\mathord{:}mem\,range},\allowbreak{} \operatorname{property}(i))))[\operatorname{Finset\mathord{:}attach}(\{i\in\mathbb{N}\mid i<n\})]\) gives \(|s| = |(i \mapsto f(i,\allowbreak{} \operatorname{Iff\mathord{:}mp}(\operatorname{Finset\mathord{:}mem\,range},\allowbreak{} \operatorname{property}(i))))[\operatorname{Finset\mathord{:}attach}(\{i\in\mathbb{N}\mid i<n\})]|\).
\par\smallskip\noindent\emph{Equality 2.}\enspace
For the required implication, argue as follows.
\begin{proofstep}
Fix \(h_{1}\) in \(\{i\in\mathbb{N}\mid i<n\}\).
Let \(i\) be the underlying element of \(h_{1}\). Its defining property gives \(i \in \{i_{1}\in\mathbb{N}\mid i_{1}<n\}\).
Fix \(h_{2}\) in \(\{i_{1}\in\mathbb{N}\mid i_{1}<n\}\).
Let \(j\) be the underlying element of \(h_{2}\). Its defining property gives \(j \in \{i_{1}\in\mathbb{N}\mid i_{1}<n\}\).
Assume \(f(i,\allowbreak{} \operatorname{Iff\mathord{:}mp}(\operatorname{Finset\mathord{:}mem\,range},\allowbreak{} \operatorname{property}((i\, :\, \{i_{1}\in\mathbb{N}\mid i_{1}<n\})))) = f(j,\allowbreak{} \operatorname{Iff\mathord{:}mp}(\operatorname{Finset\mathord{:}mem\,range},\allowbreak{} \operatorname{property}((j\, :\, \{i_{1}\in\mathbb{N}\mid i_{1}<n\}))))\).
By \paperlink{premise:46696e7365742e6d656d5f72616e6765}{Prerequisite~P\ref{premise:46696e7365742e6d656d5f72616e6765}}, we obtain \(i \in \{i_{1}\in\mathbb{N}\mid i_{1}<n\} \iff i < n\).
Use the hypothesis \(i \in \{i_{1}\in\mathbb{N}\mid i_{1}<n\}\).
By the forward implication, we obtain \(i < n\).
By \paperlink{premise:46696e7365742e6d656d5f72616e6765}{Prerequisite~P\ref{premise:46696e7365742e6d656d5f72616e6765}}, we obtain \(j \in \{i_{1}\in\mathbb{N}\mid i_{1}<n\} \iff j < n\).
Use the hypothesis \(j \in \{i_{1}\in\mathbb{N}\mid i_{1}<n\}\).
By the forward implication, we obtain \(j < n\).
Use the hypothesis \(f(i,\allowbreak{} \operatorname{Iff\mathord{:}mp}(\operatorname{Finset\mathord{:}mem\,range},\allowbreak{} \operatorname{property}((i\, :\, \{i_{1}\in\mathbb{N}\mid i_{1}<n\})))) = f(j,\allowbreak{} \operatorname{Iff\mathord{:}mp}(\operatorname{Finset\mathord{:}mem\,range},\allowbreak{} \operatorname{property}((j\, :\, \{i_{1}\in\mathbb{N}\mid i_{1}<n\}))))\).
Applying the hypothesis \(\forall i\, :\, \mathbb{N},\; \forall j\, :\, \mathbb{N},\; \forall \mathit{hi}\, :\, i < n,\; \forall \mathit{hj}\, :\, j < n,\; f(i,\allowbreak{} \mathit{hi}) = f(j,\allowbreak{} \mathit{hj}) \implies i = j\), we obtain \(i = j\).
By \paperlink{premise:537562747970652e6571}{Prerequisite~P\ref{premise:537562747970652e6571}} with \(\mathbb{N}\) as its type parameter \(\alpha\), we obtain \((i\, :\, \{i_{1}\in\mathbb{N}\mid i_{1}<n\}) = (j\, :\, \{i_{1}\in\mathbb{N}\mid i_{1}<n\})\).
This proves \(f(i,\allowbreak{} \operatorname{Iff\mathord{:}mp}(\operatorname{Finset\mathord{:}mem\,range},\allowbreak{} \operatorname{property}((i\, :\, \{i_{1}\in\mathbb{N}\mid i_{1}<n\})))) = f(\operatorname{val}(h_{2}),\allowbreak{} \operatorname{Iff\mathord{:}mp}(\operatorname{Finset\mathord{:}mem\,range},\allowbreak{} \operatorname{property}(h_{2}))) \implies (i\, :\, \{i_{1}\in\mathbb{N}\mid i_{1}<n\}) = h_{2}\).
This proves \(\forall a_{2}\, :\, \{i\in\mathbb{N}\mid i<n\},\; f(\operatorname{val}(h_{1}),\allowbreak{} \operatorname{Iff\mathord{:}mp}(\operatorname{Finset\mathord{:}mem\,range},\allowbreak{} \operatorname{property}(h_{1}))) = f(a_{2},\allowbreak{} \operatorname{Iff\mathord{:}mp}(\operatorname{Finset\mathord{:}mem\,range},\allowbreak{} \operatorname{property}(a_{2}))) \implies h_{1} = a_{2}\).
\end{proofstep}
By \paperlink{theorem:49}{Lemma~\ref{theorem:49}} with \(\{i\in\mathbb{N}\mid i<n\}\) as its type parameter \(\alpha\) and \(\operatorname{Finset\mathord{:}attach}\left(\{i\in\mathbb{N}\mid i<n\}\right)\) in place of its parameter \(s\), we obtain \(|(i \mapsto f(i,\allowbreak{} \operatorname{Iff\mathord{:}mp}(\operatorname{Finset\mathord{:}mem\,range},\allowbreak{} \operatorname{property}(i))))[\operatorname{Finset\mathord{:}attach}(\{i\in\mathbb{N}\mid i<n\})]| = |\operatorname{Finset\mathord{:}attach}(\{i\in\mathbb{N}\mid i<n\})|\).
\par\smallskip\noindent\emph{Equality 3.}\enspace
By \paperlink{theorem:35}{Lemma~\ref{theorem:35}} with \(\mathbb{N}\) as its type parameter \(\alpha\), we obtain \(|\operatorname{Finset\mathord{:}attach}(\{i\in\mathbb{N}\mid i<n\})| = |\{i\in\mathbb{N}\mid i<n\}|\).
\par\smallskip\noindent\emph{Equality 4.}\enspace
By \paperlink{theorem:34}{Lemma~\ref{theorem:34}}, we obtain \(\left|\{i\in\mathbb{N}\mid i<n\}\right| = n\).
\par Combining these equalities gives \(\left|s\right| = n\).
\end{proof}

\section{Bij}
\begin{lemma}\label{theorem:67}
Let \(\alpha\) and \(\beta\) be types. Let \(s\) be a finite subset of \(\alpha\). Let \(t\) be a finite subset of \(\beta\). Let \(i\) be a function of type \(\Pi_{a\, :\, \alpha}\, a \in s \to \beta\). Assume \(\forall a\, :\, \alpha,\; \forall \mathit{ha}\, :\, a \in s,\; i\left(a, \mathit{ha}\right) \in t\). Assume \(\forall a_{1}\, :\, \alpha,\; \forall h\, :\, a_{1} \in s,\; \forall a_{2}\, :\, \alpha,\; \forall h_{1}\, :\, a_{2} \in s,\; i(a_{1},\allowbreak{} h) = i(a_{2},\allowbreak{} h_{1}) \implies a_{1} = a_{2}\). Assume \(\forall b\, :\, \beta,\; b \in t \implies (\exists a_{1}\, :\, \alpha,\; \exists \mathit{ha}\text{ proving }(a_{1} \in s),\; i(a_{1},\allowbreak{} \mathit{ha}) = b)\). Then \(\left|s\right| = \left|t\right|\).
\end{lemma}

\begin{proof}
\par\smallskip\noindent\emph{Equality 1.}\enspace
Using \paperlink{theorem:35}{Lemma~\ref{theorem:35}} in reverse gives \(\left|s\right| = \left|\operatorname{Finset\mathord{:}attach}\left(s\right)\right|\).
\par\smallskip\noindent\emph{Equality 2.}\enspace
For the required implication, argue as follows.
\begin{proofstep}
Fix \(h\) in \(s\).
Let \(\mathit{val}\) be the underlying element of \(h\). Its defining property gives \(\mathit{val} \in s\).
Fix \(h_{1}\) in \(s\).
Let \(\mathit{val}_{1}\) be the underlying element of \(h_{1}\). Its defining property gives \(\mathit{val}_{1} \in s\).
Assume \(i(\mathit{val},\allowbreak{} \operatorname{property}((\mathit{val}\, :\, s))) = i(\mathit{val}_{1},\allowbreak{} \operatorname{property}((\mathit{val}_{1}\, :\, s)))\).
By \paperlink{premise:537562747970652e70726f7065727479}{Prerequisite~P\ref{premise:537562747970652e70726f7065727479}} with \(\mathit{self} = \left(\mathit{val}\, :\, s\right)\), we obtain \(\mathit{val} \in s\).
By \paperlink{premise:537562747970652e70726f7065727479}{Prerequisite~P\ref{premise:537562747970652e70726f7065727479}} with \(\mathit{self} = \left(\mathit{val}_{1}\, :\, s\right)\), we obtain \(\mathit{val}_{1} \in s\).
Use the hypothesis \(i(\mathit{val},\allowbreak{} \operatorname{property}((\mathit{val}\, :\, s))) = i(\mathit{val}_{1},\allowbreak{} \operatorname{property}((\mathit{val}_{1}\, :\, s)))\).
Applying the hypothesis \(\forall a_{1}\, :\, \alpha,\; \forall h\, :\, a_{1} \in s,\; \forall a_{2}\, :\, \alpha,\; \forall h_{1}\, :\, a_{2} \in s,\; i(a_{1},\allowbreak{} h) = i(a_{2},\allowbreak{} h_{1}) \implies a_{1} = a_{2}\), we obtain \(\mathit{val} = \mathit{val}_{1}\).
Rewriting with \paperlink{premise:537562747970652e6d6b2e696e6a4571}{Prerequisite~P\ref{premise:537562747970652e6d6b2e696e6a4571}} (using \(\mathit{val} \in s\) and \(\mathit{val}_{1} \in s\)) in the preceding result gives \(\left(\mathit{val}\, :\, s\right) = \left(\mathit{val}_{1}\, :\, s\right)\).
This proves \(i(\mathit{val},\allowbreak{} \operatorname{property}((\mathit{val}\, :\, s))) = i(\operatorname{val}(h_{1}),\allowbreak{} \operatorname{property}(h_{1})) \implies (\mathit{val}\, :\, s) = h_{1}\).
This proves \(\forall a_{2}\, :\, s,\; i(\operatorname{val}(h),\allowbreak{} \operatorname{property}(h)) = i(a_{2},\allowbreak{} \operatorname{property}(a_{2})) \implies h = a_{2}\).
\end{proofstep}
By \paperlink{theorem:49}{Lemma~\ref{theorem:49}} with \(s\) as its type parameter \(\alpha\) and \(\operatorname{Finset\mathord{:}attach}\left(s\right)\) in place of its parameter \(s\), we obtain \(|(a \mapsto i(a,\allowbreak{} \operatorname{property}(a)))[\operatorname{Finset\mathord{:}attach}(s)]| = |\operatorname{Finset\mathord{:}attach}(s)|\).
Reversing the equality, we obtain \(|\operatorname{Finset\mathord{:}attach}(s)| = |(a \mapsto i(a,\allowbreak{} \operatorname{property}(a)))[\operatorname{Finset\mathord{:}attach}(s)]|\).
\par\smallskip\noindent\emph{Equality 3.}\enspace
By definition, we obtain \(|(a \mapsto i(a,\allowbreak{} \operatorname{property}(a)))[\operatorname{Finset\mathord{:}attach}(s)]| = |(a \mapsto i(a,\allowbreak{} \operatorname{property}(a)))[\operatorname{Finset\mathord{:}attach}(s)]|\).
For the required implication, argue as follows.
\begin{proofstep}
Fix \(b\) in \(\beta\).
Consider the following case.
\begin{proofstep}
Assume \(b \in (a \mapsto i(a,\allowbreak{} \operatorname{property}(a)))[\operatorname{Finset\mathord{:}attach}(s)]\).
By \paperlink{premise:46696e7365742e6d656d5f696d616765}{Prerequisite~P\ref{premise:46696e7365742e6d656d5f696d616765}} with \(s\) as its type parameter \(\alpha\), we obtain \(b \in (a \mapsto i(a,\allowbreak{} \operatorname{property}(a)))[\operatorname{Finset\mathord{:}attach}(s)] \iff (\exists a\, :\, s,\; a \in \operatorname{Finset\mathord{:}attach}(s) \land i(a,\allowbreak{} \operatorname{property}(a)) = b)\).
Use the hypothesis \(b \in (a \mapsto i(a,\allowbreak{} \operatorname{property}(a)))[\operatorname{Finset\mathord{:}attach}(s)]\).
By the forward implication, we obtain \(\exists a\, :\, s,\; a \in \operatorname{Finset\mathord{:}attach}(s) \land i(a,\allowbreak{} \operatorname{property}(a)) = b\).
Consider the following case.
\begin{proofstep}
Fix \(w\) in \(s\).
Assume \(w \in \operatorname{Finset\mathord{:}attach}(s) \land i(\operatorname{val}(w),\allowbreak{} \operatorname{property}(w)) = b\).
Use the hypothesis \(w \in \operatorname{Finset\mathord{:}attach}(s) \land i(\operatorname{val}(w),\allowbreak{} \operatorname{property}(w)) = b\).
Consider the following case.
\begin{proofstep}
Assume \(w \in \operatorname{Finset\mathord{:}attach}\left(s\right)\).
Assume \(i\left(\operatorname{val}(w), \operatorname{property}\left(w\right)\right) = b\).
For the required implication, argue as follows.
\begin{proofstep}
Assume \(i(\operatorname{val}(w),\allowbreak{} \operatorname{property}(w)) \in (a \mapsto i(a,\allowbreak{} \operatorname{property}(a)))[\operatorname{Finset\mathord{:}attach}(s)]\).
By \paperlink{premise:537562747970652e70726f7065727479}{Prerequisite~P\ref{premise:537562747970652e70726f7065727479}} with \(\mathit{self} = w\), we obtain \(\operatorname{val}(w) \in s\).
Applying the hypothesis \(\forall a\, :\, \alpha,\; \forall \mathit{ha}\, :\, a \in s,\; i\left(a, \mathit{ha}\right) \in t\), we obtain \(i\left(\operatorname{val}(w), \operatorname{property}\left(w\right)\right) \in t\).
\end{proofstep}
Use the hypothesis \(i\left(\operatorname{val}(w), \operatorname{property}\left(w\right)\right) = b\).
Use the hypothesis \(b \in (a \mapsto i(a,\allowbreak{} \operatorname{property}(a)))[\operatorname{Finset\mathord{:}attach}(s)]\).
Substituting equals, we obtain \(b \in t\).
\end{proofstep}
Using the two conjuncts, we obtain \(b \in t\).
\end{proofstep}
Using the existential hypothesis, we obtain \(b \in t\).
\end{proofstep}
Consider the following case.
\begin{proofstep}
Assume \(b \in t\).
Use the hypothesis \(b \in t\).
Applying the hypothesis \(\forall b\, :\, \beta,\; b \in t \implies (\exists a_{1}\, :\, \alpha,\; \exists \mathit{ha}\text{ proving }(a_{1} \in s),\; i(a_{1},\allowbreak{} \mathit{ha}) = b)\), we obtain \(\exists a\, :\, \alpha,\; \exists \mathit{ha}\text{ proving }(a \in s),\; i(a,\allowbreak{} \mathit{ha}) = b\).
Consider the following case.
\begin{proofstep}
Fix \(a\) in \(\alpha\).
Assume \(\exists \mathit{ha}\text{ proving }\left(a \in s\right),\; i\left(a, \mathit{ha}\right) = b\).
Use the hypothesis \(\exists \mathit{ha}\text{ proving }\left(a \in s\right),\; i\left(a, \mathit{ha}\right) = b\).
Consider the following case.
\begin{proofstep}
Assume \(a \in s\).
Assume \(i\left(a, \mathit{ha}\right) = b\).
For the required implication, argue as follows.
\begin{proofstep}
Assume \(i\left(a, \mathit{ha}\right) \in t\).
By \paperlink{premise:46696e7365742e6d656d5f696d616765}{Prerequisite~P\ref{premise:46696e7365742e6d656d5f696d616765}} with \(s\) as its type parameter \(\alpha\), we obtain \(i(a,\allowbreak{} \mathit{ha}) \in (a_{1} \mapsto i(a_{1},\allowbreak{} \operatorname{property}(a_{1})))[\operatorname{Finset\mathord{:}attach}(s)] \iff (\exists a_{1}\, :\, s,\; a_{1} \in \operatorname{Finset\mathord{:}attach}(s) \land i(a_{1},\allowbreak{} \operatorname{property}(a_{1})) = i(a,\allowbreak{} \mathit{ha}))\).
Choose the witness \(\left(a\, :\, s\right)\).
\par\smallskip\noindent\emph{Property.}\enspace
Use the hypothesis \(a \in s\), which simplifies to \(\left(x \mapsto x \in s\right)\left(a\right)\).
Rewriting with \paperlink{premise:46696e7365742e6d656d5f617474616368}{Prerequisite~P\ref{premise:46696e7365742e6d656d5f617474616368}} and \paperlink{premise:616e645f73656c66}{Prerequisite~P\ref{premise:616e645f73656c66}} gives \((a\, :\, s) \in \operatorname{Finset\mathord{:}attach}(s) \land i(a,\allowbreak{} \operatorname{property}((a\, :\, s))) = i(a,\allowbreak{} \operatorname{property}((a\, :\, s)))\).
This witness proves existence. Hence \(\exists a_{1}\, :\, s,\; a_{1} \in \operatorname{Finset\mathord{:}attach}(s) \land i(a_{1},\allowbreak{} \operatorname{property}(a_{1})) = i(a,\allowbreak{} \mathit{ha})\).
By the reverse implication, we obtain \(i(a,\allowbreak{} \mathit{ha}) \in (a_{1} \mapsto i(a_{1},\allowbreak{} \operatorname{property}(a_{1})))[\operatorname{Finset\mathord{:}attach}(s)]\).
\end{proofstep}
Use the hypothesis \(i\left(a, \mathit{ha}\right) = b\).
Use the hypothesis \(b \in t\).
Substituting equals, we obtain \(b \in (a_{1} \mapsto i(a_{1},\allowbreak{} \operatorname{property}(a_{1})))[\operatorname{Finset\mathord{:}attach}(s)]\).
\end{proofstep}
Using the existential hypothesis, we obtain \(b \in (a_{1} \mapsto i(a_{1},\allowbreak{} \operatorname{property}(a_{1})))[\operatorname{Finset\mathord{:}attach}(s)]\).
\end{proofstep}
Using the existential hypothesis, we obtain \(b \in (a \mapsto i(a,\allowbreak{} \operatorname{property}(a)))[\operatorname{Finset\mathord{:}attach}(s)]\).
\end{proofstep}
Together the two implications give the equivalence. Hence \(b \in (a \mapsto i(a,\allowbreak{} \operatorname{property}(a)))[\operatorname{Finset\mathord{:}attach}(s)] \iff b \in t\).
\end{proofstep}
By \paperlink{premise:46696e7365742e657874}{Prerequisite~P\ref{premise:46696e7365742e657874}}, we obtain \((a \mapsto i(a,\allowbreak{} \operatorname{property}(a)))[\operatorname{Finset\mathord{:}attach}(s)] = t\).
Substituting equals, we obtain \(|(a \mapsto i(a,\allowbreak{} \operatorname{property}(a)))[\operatorname{Finset\mathord{:}attach}(s)]| = |t|\).
\par Combining these equalities gives \(\left|s\right| = \left|t\right|\).
\end{proof}

\begin{lemma}\label{theorem:68}
Let \(\alpha\) and \(\beta\) be types. Let \(s\) be a finite subset of \(\alpha\). Let \(t\) be a finite subset of \(\beta\). Let \(i\) be a function of type \(\Pi_{a\, :\, \alpha}\, a \in s \to \beta\). Let \(j\) be a function of type \(\Pi_{a\, :\, \beta}\, a \in t \to \alpha\). Let \(\mathit{hi}\) be a proof of \(\forall a\, :\, \alpha,\; \forall \mathit{ha}\, :\, a \in s,\; i\left(a, \mathit{ha}\right) \in t\). Let \(\mathit{hj}\) be a proof of \(\forall a\, :\, \beta,\; \forall \mathit{ha}\, :\, a \in t,\; j\left(a, \mathit{ha}\right) \in s\). Assume \(\forall a\, :\, \alpha,\; \forall \mathit{ha}\, :\, a \in s,\; j(i(a,\allowbreak{} \mathit{ha}),\allowbreak{} \mathit{hi}(a,\allowbreak{} \mathit{ha})) = a\). Assume \(\forall a\, :\, \beta,\; \forall \mathit{ha}\, :\, a \in t,\; i(j(a,\allowbreak{} \mathit{ha}),\allowbreak{} \mathit{hj}(a,\allowbreak{} \mathit{ha})) = a\). Then \(\left|s\right| = \left|t\right|\).
\end{lemma}

\begin{proof}
For the required implication, argue as follows.
\begin{proofstep}
Fix \(\mathit{a1}\) in \(\alpha\).
Assume \(\mathit{a1} \in s\).
Fix \(\mathit{a2}\) in \(\alpha\).
Assume \(\mathit{a2} \in s\).
Assume \(i\left(\mathit{a1}, \mathit{h1}\right) = i\left(\mathit{a2}, \mathit{h2}\right)\).
Rewriting with \(i\left(\mathit{a1}, \mathit{h1}\right) = i\left(\mathit{a2}, \mathit{h2}\right)\), \(\forall a\, :\, \alpha,\; \forall \mathit{ha}\, :\, a \in s,\; i\left(a, \mathit{ha}\right) \in t\), \(\mathit{a1} \in s\) and substitution of equals gives \(j(i(\mathit{a1},\allowbreak{} \mathit{h1}),\allowbreak{} \mathit{hi}(\mathit{a1},\allowbreak{} \mathit{h1})) = j(i(\mathit{a2},\allowbreak{} \mathit{h2}),\allowbreak{} \mathit{hi}(\mathit{a2},\allowbreak{} \mathit{h2}))\).
Rewriting with \(\forall a\, :\, \alpha,\; \forall \mathit{ha}\, :\, a \in s,\; j(i(a,\allowbreak{} \mathit{ha}),\allowbreak{} \mathit{hi}(a,\allowbreak{} \mathit{ha})) = a\) and \(\mathit{a2} \in s\) in the preceding result gives \(j(i(\mathit{a1},\allowbreak{} \mathit{h1}),\allowbreak{} \mathit{hi}(\mathit{a1},\allowbreak{} \mathit{h1})) = \mathit{a2}\).
Rewriting with \(\forall a\, :\, \alpha,\; \forall \mathit{ha}\, :\, a \in s,\; j(i(a,\allowbreak{} \mathit{ha}),\allowbreak{} \mathit{hi}(a,\allowbreak{} \mathit{ha})) = a\) and \(\mathit{a1} \in s\) in the preceding result gives \(\mathit{a1} = \mathit{a2}\).
\end{proofstep}
For the required implication, argue as follows.
\begin{proofstep}
Fix \(b\) in \(\beta\).
Assume \(b \in t\).
Choose the witness \(j\left(b, \mathit{hb}\right)\).
The construction requires evidence that \(j\left(b, \mathit{hb}\right) \in s\).
Use the hypothesis \(b \in t\).
Applying the hypothesis \(\forall a\, :\, \beta,\; \forall \mathit{ha}\, :\, a \in t,\; j\left(a, \mathit{ha}\right) \in s\), we obtain \(j\left(b, \mathit{hb}\right) \in s\).
Use the hypothesis \(b \in t\).
Applying the hypothesis \(\forall a\, :\, \beta,\; \forall \mathit{ha}\, :\, a \in t,\; i(j(a,\allowbreak{} \mathit{ha}),\allowbreak{} \mathit{hj}(a,\allowbreak{} \mathit{ha})) = a\), we obtain \(i\left(j\left(b, \mathit{hb}\right), \mathit{hj}\left(b, \mathit{hb}\right)\right) = b\).
This witness proves existence. Hence \(\exists \mathit{ha}\text{ proving }(j(b,\allowbreak{} \mathit{hb}) \in s),\; i(j(b,\allowbreak{} \mathit{hb}),\allowbreak{} \mathit{ha}) = b\).
This witness proves existence. Hence \(\exists a\, :\, \alpha,\; \exists \mathit{ha}\text{ proving }(a \in s),\; i(a,\allowbreak{} \mathit{ha}) = b\).
\end{proofstep}
Apply \paperlink{theorem:67}{Lemma~\ref{theorem:67}}.
\end{proof}

\begin{theorem}\label{theorem:69}
Let \(\alpha\) and \(\beta\) be types. Let \(s\) be a finite subset of \(\alpha\). Let \(t\) be a finite subset of \(\beta\). Let \(i\) be a function of type \(\alpha \to \beta\). Assume \(\forall a\, :\, \alpha,\; a \in s \implies i\left(a\right) \in t\). Assume \(i\text{ is injective on }\operatorname{underlying\,set}(s)\). Assume \(\operatorname{underlying\,set}(t)\subseteq i\left[\operatorname{underlying\,set}(s)\right]\). Then \(\left|s\right| = \left|t\right|\).
\end{theorem}

\begin{proof}
The substitution uses the following argument.
\begin{proofstep}
Fix \(b\) in \(\beta\).
By definition, we obtain \(\left(b \in t\right) = \left(b \in t\right)\).
For the required implication, argue as follows.
\begin{proofstep}
Fix \(a\) in \(\alpha\).
For the required implication, argue as follows.
\begin{proofstep}
Assume \(a \in s\).
By definition, we obtain \(\left(i\left(a\right) = b\right) = \left(i\left(a\right) = b\right)\).
By \paperlink{premise:4966662e6f665f6571}{Prerequisite~P\ref{premise:4966662e6f665f6571}}, we obtain \(i\left(a\right) = b \iff i\left(a\right) = b\).
\end{proofstep}
By definition, we obtain \(\left(a \in s\right) = \left(a \in s\right)\).
By \paperlink{premise:4966662e6f665f6571}{Prerequisite~P\ref{premise:4966662e6f665f6571}}, we obtain \(a \in s \iff a \in s\).
By \paperlink{premise:6578697374735f70726f705f636f6e6772}{Prerequisite~P\ref{premise:6578697374735f70726f705f636f6e6772}}, we obtain \(a \in s \land i\left(a\right) = b \iff a \in s \land i\left(a\right) = b\).
By propositional extensionality, we obtain \(\left(a \in s \land i\left(a\right) = b\right) = \left(a \in s \land i\left(a\right) = b\right)\).
By \paperlink{premise:6578697374735f70726f70}{Prerequisite~P\ref{premise:6578697374735f70726f70}}, we obtain \(a \in s \land i\left(a\right) = b \iff a \in s \land i\left(a\right) = b\).
By propositional extensionality, we obtain \(\left(a \in s \land i\left(a\right) = b\right) = \left(a \in s \land i\left(a\right) = b\right)\).
Chaining the equalities, we obtain \(\left(a \in s \land i\left(a\right) = b\right) = \left(a \in s \land i\left(a\right) = b\right)\).
\end{proofstep}
By function extensionality, we obtain \((x \mapsto x \in s \land i(x) = b) = (x \mapsto x \in s \land i(x) = b)\).
Applying the same function to equal arguments, we obtain \((\exists x\, :\, \alpha,\; x \in s \land i(x) = b) = (\exists x\, :\, \alpha,\; x \in s \land i(x) = b)\).
By \paperlink{premise:696d706c6965735f636f6e6772}{Prerequisite~P\ref{premise:696d706c6965735f636f6e6772}} with \(p_{1} = b \in t\) and \(p_{2} = b \in t\) and \(q_{1} = \exists x\, :\, \alpha,\; \exists \mathit{ha}\text{ proving }(x \in s),\; (a \mapsto x_{1} \mapsto i(a))(x,\allowbreak{} \mathit{ha}) = b\) and \(q_{2} = \exists x\, :\, \alpha,\; x \in s \land i\left(x\right) = b\), we obtain \((b \in t \implies (\exists x\, :\, \alpha,\; x \in s \land i(x) = b)) = (b \in t \implies (\exists x\, :\, \alpha,\; x \in s \land i(x) = b))\).
\end{proofstep}
Use the hypothesis \(\operatorname{underlying\,set}(t)\subseteq i\left[\operatorname{underlying\,set}(s)\right]\), which simplifies to \(\forall b\, :\, \beta,\; b \in t \implies (\exists a_{1}\, :\, \alpha,\; a_{1} \in s \land i(a_{1}) = b)\).
Rewriting with \paperlink{premise:666f72616c6c5f636f6e6772}{Prerequisite~P\ref{premise:666f72616c6c5f636f6e6772}} and the identity established above in the preceding result gives \(\forall b\, :\, \beta,\; b \in t \implies (\exists a_{1}\, :\, \alpha,\; a_{1} \in s \land i(a_{1}) = b)\).
Apply \paperlink{theorem:67}{Lemma~\ref{theorem:67}} with \(a \mapsto x \mapsto i\left(a\right)\) in place of its parameter \(i\).
\end{proof}

\begin{corollary}\label{theorem:70}
Let \(\alpha\) and \(\beta\) be types. Let \(s\) be a finite subset of \(\alpha\). Let \(t\) be a finite subset of \(\beta\). Let \(i\) be a function of type \(\alpha \to \beta\). Let \(j\) be a function of type \(\beta \to \alpha\). Assume \(\forall a\, :\, \alpha,\; a \in s \implies i\left(a\right) \in t\). Assume \(\forall a\, :\, \beta,\; a \in t \implies j\left(a\right) \in s\). Assume \(\forall a\, :\, \alpha,\; a \in s \implies j\left(i\left(a\right)\right) = a\). Assume \(\forall a\, :\, \beta,\; a \in t \implies i\left(j\left(a\right)\right) = a\). Then \(\left|s\right| = \left|t\right|\).
\end{corollary}

\begin{proof}
Apply \paperlink{theorem:68}{Lemma~\ref{theorem:68}} with \(a \mapsto x \mapsto i\left(a\right)\) in place of its parameter \(i\) and \(b \mapsto x \mapsto j\left(b\right)\) in place of its parameter \(j\).
\end{proof}

\begin{lemma}\label{theorem:71}
Let \(\alpha\) and \(\beta\) be types. Let \(s\) be a finite subset of \(\alpha\). Let \(t\) be a finite subset of \(\beta\). Let \(e\) be a bijection from \(\alpha\) to \(\beta\). Assume \(\forall i\, :\, \alpha,\; i \in s \iff e\left(i\right) \in t\). Then \(\left|s\right| = \left|t\right|\).
\end{lemma}

\begin{proof}
The substitution uses the following argument.
\begin{proofstep}
Fix \(a\) in \(\alpha\).
Applying the hypothesis \(\forall i\, :\, \alpha,\; i \in s \iff e\left(i\right) \in t\), we obtain \(a \in s \iff e\left(a\right) \in t\).
By propositional extensionality, we obtain \(\left(a \in s\right) = \left(e\left(a\right) \in t\right)\).
By definition, we obtain \(\left(e\left(a\right) \in t\right) = \left(e\left(a\right) \in t\right)\).
By \paperlink{premise:696d706c6965735f636f6e6772}{Prerequisite~P\ref{premise:696d706c6965735f636f6e6772}} with \(p_{1} = a \in s\) and \(p_{2} = e\left(a\right) \in t\) and \(q_{1} = e\left(a\right) \in t\) and \(q_{2} = e\left(a\right) \in t\), we obtain \((a \in s \implies e(a) \in t) = (e(a) \in t \implies e(a) \in t)\).
By \paperlink{premise:696d705f73656c66}{Prerequisite~P\ref{premise:696d705f73656c66}}, we obtain \(e\left(a\right) \in t \implies e\left(a\right) \in t \iff \mathrm{True}\).
By propositional extensionality, we obtain \(\left(e\left(a\right) \in t \implies e\left(a\right) \in t\right) = \mathrm{True}\).
Chaining the equalities, we obtain \(\left(a \in s \implies e\left(a\right) \in t\right) = \mathrm{True}\).
\end{proofstep}
Rewriting with \paperlink{premise:666f72616c6c5f636f6e6772}{Prerequisite~P\ref{premise:666f72616c6c5f636f6e6772}}, \paperlink{premise:696d706c6965735f74727565}{Prerequisite~P\ref{premise:696d706c6965735f74727565}} and the identity established above gives \(\forall a\, :\, \alpha,\; a \in s \implies e\left(a\right) \in t\).
The substitution uses the following argument.
\begin{proofstep}
Fix \(a\) in \(\beta\).
By definition, we obtain \(\left(a \in t\right) = \left(a \in t\right)\).
Applying the hypothesis \(\forall i\, :\, \alpha,\; i \in s \iff e\left(i\right) \in t\), we obtain \(\left(e^{-1}\right)\left(a\right) \in s \iff e\left(\left(e^{-1}\right)\left(a\right)\right) \in t\).
By propositional extensionality, we obtain \(((e^{-1})(a) \in s) = (e((e^{-1})(a)) \in t)\).
By \paperlink{premise:45717569762e6170706c795f73796d6d5f6170706c79}{Prerequisite~P\ref{premise:45717569762e6170706c795f73796d6d5f6170706c79}} with \(x = a\), we obtain \(e\left(\left(e^{-1}\right)\left(a\right)\right) = a\).
Applying the same function to equal arguments, we obtain \(\left(e\left(\left(e^{-1}\right)\left(a\right)\right) \in t\right) = \left(a \in t\right)\).
Chaining the equalities, we obtain \(\left(\left(e^{-1}\right)\left(a\right) \in s\right) = \left(a \in t\right)\).
By \paperlink{premise:696d706c6965735f636f6e6772}{Prerequisite~P\ref{premise:696d706c6965735f636f6e6772}} with \(p_{1} = a \in t\) and \(p_{2} = a \in t\) and \(q_{1} = \left(e^{-1}\right)\left(a\right) \in s\) and \(q_{2} = a \in t\), we obtain \(\left(a \in t \implies \left(e^{-1}\right)\left(a\right) \in s\right) = \left(a \in t \implies a \in t\right)\).
By \paperlink{premise:696d705f73656c66}{Prerequisite~P\ref{premise:696d705f73656c66}}, we obtain \(a \in t \implies a \in t \iff \mathrm{True}\).
By propositional extensionality, we obtain \(\left(a \in t \implies a \in t\right) = \mathrm{True}\).
Chaining the equalities, we obtain \(\left(a \in t \implies \left(e^{-1}\right)\left(a\right) \in s\right) = \mathrm{True}\).
\end{proofstep}
Rewriting with \paperlink{premise:666f72616c6c5f636f6e6772}{Prerequisite~P\ref{premise:666f72616c6c5f636f6e6772}}, \paperlink{premise:696d706c6965735f74727565}{Prerequisite~P\ref{premise:696d706c6965735f74727565}} and the identity established above gives \(\forall a\, :\, \beta,\; a \in t \implies \left(e^{-1}\right)\left(a\right) \in s\).
The substitution uses the following argument.
\begin{proofstep}
Fix \(a\) in \(\alpha\).
Applying the hypothesis \(\forall i\, :\, \alpha,\; i \in s \iff e\left(i\right) \in t\), we obtain \(a \in s \iff e\left(a\right) \in t\).
By propositional extensionality, we obtain \(\left(a \in s\right) = \left(e\left(a\right) \in t\right)\).
By \paperlink{premise:45717569762e73796d6d5f6170706c795f6170706c79}{Prerequisite~P\ref{premise:45717569762e73796d6d5f6170706c795f6170706c79}} with \(x = a\), we obtain \(\left(e^{-1}\right)\left(e\left(a\right)\right) = a\).
Applying the same function to equal arguments, we obtain \(\left(\left(e^{-1}\right)\left(e\left(a\right)\right) = a\right) = \left(a = a\right)\).
By \paperlink{premise:65715f73656c66}{Prerequisite~P\ref{premise:65715f73656c66}}, we obtain \(\left(a = a\right) = \mathrm{True}\).
Chaining the equalities, we obtain \(\left(\left(e^{-1}\right)\left(e\left(a\right)\right) = a\right) = \mathrm{True}\).
By \paperlink{premise:696d706c6965735f636f6e6772}{Prerequisite~P\ref{premise:696d706c6965735f636f6e6772}} with \(p_{1} = a \in s\) and \(p_{2} = e\left(a\right) \in t\) and \(q_{1} = \left(e^{-1}\right)\left(e\left(a\right)\right) = a\) and \(q_{2} = \mathrm{True}\), we obtain \((a \in s \implies (e^{-1})(e(a)) = a) = (e(a) \in t \implies \mathrm{True})\).
By \paperlink{premise:696d706c6965735f74727565}{Prerequisite~P\ref{premise:696d706c6965735f74727565}} with \(e\left(a\right) \in t\) as its type parameter \(\alpha\), we obtain \(\left(e\left(a\right) \in t \implies \mathrm{True}\right) = \mathrm{True}\).
Chaining the equalities, we obtain \(\left(a \in s \implies \left(e^{-1}\right)\left(e\left(a\right)\right) = a\right) = \mathrm{True}\).
\end{proofstep}
Rewriting with \paperlink{premise:666f72616c6c5f636f6e6772}{Prerequisite~P\ref{premise:666f72616c6c5f636f6e6772}}, \paperlink{premise:696d706c6965735f74727565}{Prerequisite~P\ref{premise:696d706c6965735f74727565}} and the identity established above gives \(\forall a\, :\, \alpha,\; a \in s \implies \left(e^{-1}\right)\left(e\left(a\right)\right) = a\).
The substitution uses the following argument.
\begin{proofstep}
Fix \(a\) in \(\beta\).
By definition, we obtain \(\left(a \in t\right) = \left(a \in t\right)\).
By \paperlink{premise:45717569762e6170706c795f73796d6d5f6170706c79}{Prerequisite~P\ref{premise:45717569762e6170706c795f73796d6d5f6170706c79}} with \(x = a\), we obtain \(e\left(\left(e^{-1}\right)\left(a\right)\right) = a\).
Applying the same function to equal arguments, we obtain \(\left(e\left(\left(e^{-1}\right)\left(a\right)\right) = a\right) = \left(a = a\right)\).
By \paperlink{premise:65715f73656c66}{Prerequisite~P\ref{premise:65715f73656c66}}, we obtain \(\left(a = a\right) = \mathrm{True}\).
Chaining the equalities, we obtain \(\left(e\left(\left(e^{-1}\right)\left(a\right)\right) = a\right) = \mathrm{True}\).
By \paperlink{premise:696d706c6965735f636f6e6772}{Prerequisite~P\ref{premise:696d706c6965735f636f6e6772}} with \(p_{1} = a \in t\) and \(p_{2} = a \in t\) and \(q_{1} = e\left(\left(e^{-1}\right)\left(a\right)\right) = a\) and \(q_{2} = \mathrm{True}\), we obtain \((a \in t \implies e((e^{-1})(a)) = a) = (a \in t \implies \mathrm{True})\).
By \paperlink{premise:696d706c6965735f74727565}{Prerequisite~P\ref{premise:696d706c6965735f74727565}} with \(a \in t\) as its type parameter \(\alpha\), we obtain \(\left(a \in t \implies \mathrm{True}\right) = \mathrm{True}\).
Chaining the equalities, we obtain \(\left(a \in t \implies e\left(\left(e^{-1}\right)\left(a\right)\right) = a\right) = \mathrm{True}\).
\end{proofstep}
Rewriting with \paperlink{premise:666f72616c6c5f636f6e6772}{Prerequisite~P\ref{premise:666f72616c6c5f636f6e6772}}, \paperlink{premise:696d706c6965735f74727565}{Prerequisite~P\ref{premise:696d706c6965735f74727565}} and the identity established above gives \(\forall a\, :\, \beta,\; a \in t \implies e\left(\left(e^{-1}\right)\left(a\right)\right) = a\).
Apply \paperlink{theorem:70}{Corollary~\ref{theorem:70}} with \(i = \operatorname{coe}\left(e\right)\) and \(j = \operatorname{coe}\left(e^{-1}\right)\).
\end{proof}

\begin{corollary}\label{theorem:72}
Let \(\alpha\) and \(\beta\) be types. Let \(s\) be a finite subset of \(\alpha\). Let \(t\) be a finite subset of \(\beta\). Let \(e\) be a function of type \(\alpha \to \beta\). Assume \(e\text{ is bijective}\). Assume \(\forall i\, :\, \alpha,\; i \in s \iff e\left(i\right) \in t\). Then \(\left|s\right| = \left|t\right|\).
\end{corollary}

\begin{proof}
Apply \paperlink{theorem:71}{Lemma~\ref{theorem:71}} with \(\operatorname{of\,Bijective}\left(e\right)\) in place of its parameter \(e\).
\end{proof}

\begin{lemma}\label{theorem:73}
Let \(\alpha\) and \(\beta\) be types. Let \(s\) be a finite subset of \(\alpha\). Let \(t\) be a finite subset of \(\beta\). Let \(f\) be a function of type \(\alpha \to \beta\). Assume \(\forall a\, :\, \alpha,\; a \in s \implies f\left(a\right) \in t\). Assume \(f\text{ is injective on }\operatorname{underlying\,set}(s)\). Then \(\left|s\right| \le \left|t\right|\).
\end{lemma}

\begin{proof}
By \paperlink{theorem:46}{Lemma~\ref{theorem:46}}, we obtain \(\left|f\left[s\right]\right| = \left|s\right|\).
Reversing the equality, we obtain \(\left|s\right| = \left|f\left[s\right]\right|\).
By \paperlink{premise:46696e7365742e696d6167655f7375627365745f696666}{Prerequisite~P\ref{premise:46696e7365742e696d6167655f7375627365745f696666}}, we obtain \(f\left[s\right] \subseteq t \iff \left(\forall x\, :\, \alpha,\; x \in s \implies f\left(x\right) \in t\right)\).
Use the hypothesis \(\forall a\, :\, \alpha,\; a \in s \implies f\left(a\right) \in t\).
By the reverse implication, we obtain \(f\left[s\right] \subseteq t\).
By \paperlink{theorem:5}{Lemma~\ref{theorem:5}}, we obtain \(\left|f\left[s\right]\right| \le \left|t\right|\).
By transitivity of the displayed relations, we obtain \(\left|s\right| \le \left|t\right|\).
\end{proof}

\begin{theorem}\label{theorem:74}
Let \(\alpha\) and \(\beta\) be types. Let \(s\) be a finite subset of \(\alpha\). Let \(t\) be a finite subset of \(\beta\). Let \(f\) be a function of type \(s \to t\). Assume \(f\text{ is injective}\). Then \(\left|s\right| \le \left|t\right|\).
\end{theorem}

\begin{proof}
By \paperlink{premise:46696e7365742e65715f656d7074795f6f725f6e6f6e656d707479}{Prerequisite~P\ref{premise:46696e7365742e65715f656d7074795f6f725f6e6f6e656d707479}}, we obtain \(s = \varnothing \lor s\ne\varnothing\).
Consider the following case.
\begin{proofstep}
Assume \(s = \varnothing\).
For the required implication, argue as follows.
\begin{proofstep}
Fix \(f_{1}\) in \(\varnothing \to t\).
Assume \(f_{1}\text{ is injective}\).
Rewriting with \paperlink{premise:4e61742e7a65726f5f6c65}{Prerequisite~P\ref{premise:4e61742e7a65726f5f6c65}} gives \(0 \le \left|t\right|\).
\end{proofstep}
Use the hypothesis \(s = \varnothing\).
Reversing the equality, we obtain \(\varnothing = s\).
Use the hypothesis \(f\text{ is injective}\).
Substituting equals, we obtain \(\left|s\right| \le \left|t\right|\).
\end{proofstep}
Consider the following case.
\begin{proofstep}
Assume \(s\ne\varnothing\).
Use the hypothesis \(s\ne\varnothing\).
Consider the following case.
\begin{proofstep}
Fix \(a_{0}\) in \(\alpha\).
Assume \(a_{0} \in s\).
Set \(f^{'}\) equal to \(a \mapsto \operatorname{val}(f(\begin{cases}(a\, :\, s) &\text{given } \mathit{ha}\text{ proving } a \in s\\ (a_{0}\, :\, s) &\text{if } \neg (a \in s)\end{cases}))\).
For the required implication, argue as follows.
\begin{proofstep}
Fix \(a\) in \(\alpha\).
Assume \(a \in s\).
Rewriting with \paperlink{premise:646974655f636f6e645f65715f74727565}{Prerequisite~P\ref{premise:646974655f636f6e645f65715f74727565}}, \(a \in s\) and \paperlink{premise:46696e7365742e636f655f6d656d}{Prerequisite~P\ref{premise:46696e7365742e636f655f6d656d}} gives \(\operatorname{val}(f(\begin{cases}(a\, :\, s) &\text{given } \mathit{ha}\text{ proving } a \in s\\ (a_{0}\, :\, s) &\text{if } \neg (a \in s)\end{cases})) \in t\).
\end{proofstep}
For the required implication, argue as follows.
\begin{proofstep}
Fix \(a_{1}\) in \(\alpha\).
Assume \(a_{1} \in \operatorname{underlying\,set}(s)\).
Fix \(a_{2}\) in \(\alpha\).
Assume \(a_{2} \in \operatorname{underlying\,set}(s)\).
Assume \(f^{'}\left(a_{1}\right) = f^{'}\left(a_{2}\right)\).
By \paperlink{premise:537562747970652e6578745f696666}{Prerequisite~P\ref{premise:537562747970652e6578745f696666}}, we obtain \(\begin{cases}(a_{1}\, :\, s) &\text{given } h_{4}\text{ proving } \mathrm{True}\\ (a_{0}\, :\, s) &\text{if } \neg \mathrm{True}\end{cases} = \begin{cases}(a_{2}\, :\, s) &\text{given } h_{4}\text{ proving } \mathrm{True}\\ (a_{0}\, :\, s) &\text{if } \neg \mathrm{True}\end{cases} \iff \operatorname{val}(\begin{cases}(a_{1}\, :\, s) &\text{given } h_{4}\text{ proving } \mathrm{True}\\ (a_{0}\, :\, s) &\text{if } \neg \mathrm{True}\end{cases}) = \operatorname{val}(\begin{cases}(a_{2}\, :\, s) &\text{given } h_{4}\text{ proving } \mathrm{True}\\ (a_{0}\, :\, s) &\text{if } \neg \mathrm{True}\end{cases})\).
The substitution uses the following argument.
\begin{proofstep}
Assume \(\mathrm{True}\).
By definition, we obtain \(\left(a_{1}\, :\, s\right) = \left(a_{1}\, :\, s\right)\).
\end{proofstep}
The substitution uses the following argument.
\begin{proofstep}
Assume \(\neg \mathrm{True}\).
By definition, we obtain \(\left(a_{0}\, :\, s\right) = \left(a_{0}\, :\, s\right)\).
\end{proofstep}
The substitution uses the following argument.
\begin{proofstep}
Assume \(\mathrm{True}\).
By definition, we obtain \(\left(a_{2}\, :\, s\right) = \left(a_{2}\, :\, s\right)\).
\end{proofstep}
The substitution uses the following argument.
\begin{proofstep}
Assume \(\neg \mathrm{True}\).
By definition, we obtain \(\left(a_{0}\, :\, s\right) = \left(a_{0}\, :\, s\right)\).
\end{proofstep}
Use the hypothesis \(f^{'}\left(a_{1}\right) = f^{'}\left(a_{2}\right)\).
Rewriting with \paperlink{premise:646974655f636f6e6772}{Prerequisite~P\ref{premise:646974655f636f6e6772}}, propositional extensionality, \paperlink{premise:46696e7365742e6d656d5f636f65}{Prerequisite~P\ref{premise:46696e7365742e6d656d5f636f65}}, \(a_{1} \in \operatorname{underlying\,set}(s)\), substitution of equals, \(a_{2} \in \operatorname{underlying\,set}(s)\), \paperlink{premise:537562747970652e6578745f696666}{Prerequisite~P\ref{premise:537562747970652e6578745f696666}} and the 4 identities established above in the preceding result gives \(f(\begin{cases}(a_{1}\, :\, s) &\text{given } h_{4}\text{ proving } \mathrm{True}\\ (a_{0}\, :\, s) &\text{if } \neg \mathrm{True}\end{cases}) = f(\begin{cases}(a_{2}\, :\, s) &\text{given } h_{4}\text{ proving } \mathrm{True}\\ (a_{0}\, :\, s) &\text{if } \neg \mathrm{True}\end{cases})\).
Applying the hypothesis \(f\text{ is injective}\), we obtain \(\begin{cases}(a_{1}\, :\, s) &\text{given } h_{4}\text{ proving } \mathrm{True}\\ (a_{0}\, :\, s) &\text{if } \neg \mathrm{True}\end{cases} = \begin{cases}(a_{2}\, :\, s) &\text{given } h_{4}\text{ proving } \mathrm{True}\\ (a_{0}\, :\, s) &\text{if } \neg \mathrm{True}\end{cases}\).
By the forward implication, we obtain \(\operatorname{val}(\begin{cases}(a_{1}\, :\, s) &\text{given } h_{4}\text{ proving } \mathrm{True}\\ (a_{0}\, :\, s) &\text{if } \neg \mathrm{True}\end{cases}) = \operatorname{val}(\begin{cases}(a_{2}\, :\, s) &\text{given } h_{4}\text{ proving } \mathrm{True}\\ (a_{0}\, :\, s) &\text{if } \neg \mathrm{True}\end{cases})\).
\end{proofstep}
By \paperlink{theorem:73}{Lemma~\ref{theorem:73}} with \(f^{'}\) in place of its parameter \(f\), we obtain \(\left|s\right| \le \left|t\right|\).
\end{proofstep}
Using the existential hypothesis, we obtain \(\left|s\right| \le \left|t\right|\).
\end{proofstep}
The two alternatives exhaust the cases. Hence \(\left|s\right| \le \left|t\right|\).
\end{proof}

\begin{theorem}\label{theorem:75}
Let \(\alpha\) and \(\beta\) be types. Let \(s\) be a finite subset of \(\alpha\). Let \(t\) be a finite subset of \(\beta\). Let \(f\) be a function of type \(\alpha \to \beta\). Assume \(\operatorname{underlying\,set}(t)\subseteq f\left[\operatorname{underlying\,set}(s)\right]\). Then \(\left|t\right| \le \left|s\right|\).
\end{theorem}

\begin{proof}
Using \paperlink{premise:46696e7365742e636f655f696d616765}{Prerequisite~P\ref{premise:46696e7365742e636f655f696d616765}} in reverse gives \(f\left[\operatorname{underlying\,set}(s)\right] = \operatorname{underlying\,set}(f\left[s\right])\).
Applying the same function to equal arguments, we obtain \((\operatorname{underlying\,set}(t) \subseteq f[\operatorname{underlying\,set}(s)]) = (\operatorname{underlying\,set}(t) \subseteq \operatorname{underlying\,set}(f[s]))\).
By \paperlink{premise:46696e7365742e636f655f737562736574}{Prerequisite~P\ref{premise:46696e7365742e636f655f737562736574}}, we obtain \(\operatorname{underlying\,set}(t) \subseteq \operatorname{underlying\,set}(f[s]) \iff t \subseteq f[s]\).
By propositional extensionality, we obtain \((\operatorname{underlying\,set}(t) \subseteq \operatorname{underlying\,set}(f[s])) = (t \subseteq f[s])\).
Chaining the equalities, we obtain \((\operatorname{underlying\,set}(t) \subseteq f[\operatorname{underlying\,set}(s)]) = (t \subseteq f[s])\).
Use the hypothesis \(\operatorname{underlying\,set}(t)\subseteq f\left[\operatorname{underlying\,set}(s)\right]\).
Transporting the result along the equality, we obtain \(t \subseteq f\left[s\right]\).
By \paperlink{theorem:5}{Lemma~\ref{theorem:5}}, we obtain \(\left|t\right| \le \left|f\left[s\right]\right|\).
By \paperlink{theorem:45}{Corollary~\ref{theorem:45}}, we obtain \(\left|f\left[s\right]\right| \le \left|s\right|\).
By \paperlink{premise:4c452e6c652e7472616e73}{\emph{Alias}} (\paperlink{premise:4c452e6c652e7472616e73}{Prerequisite~P\ref{premise:4c452e6c652e7472616e73}}) with \(\mathbb{N}\) as its type parameter \(\alpha\), we obtain \(\left|t\right| \le \left|s\right|\).
\end{proof}

\begin{theorem}\label{theorem:76}
Let \(\alpha\) and \(\beta\) be types. Let \(s\) be a finite subset of \(\alpha\). Let \(t\) be a finite subset of \(\beta\). Assume \(\left|t\right| < \left|s\right|\). Let \(f\) be a function of type \(\alpha \to \beta\). Assume \(\forall a\, :\, \alpha,\; a \in s \implies f\left(a\right) \in t\). Then
\begin{center}
\(\displaystyle \exists x\, :\, \alpha,\; x \in s \land (\exists y\, :\, \alpha,\; y \in s \land (x \ne y \land f(x) = f(y))).\)
\end{center}
\end{theorem}

\begin{proof}
Consider the following case.
\begin{proofstep}
Assume \(\neg (\exists x\, :\, \alpha,\; x \in s \land (\exists y\, :\, \alpha,\; y \in s \land (x \ne y \land f(x) = f(y))))\).
Use the hypothesis \(\left|t\right| < \left|s\right|\).
For the required implication, argue as follows.
\begin{proofstep}
Fix \(x\) in \(\alpha\).
Assume \(x \in \operatorname{underlying\,set}(s)\).
Fix \(y\) in \(\alpha\).
Assume \(y \in \operatorname{underlying\,set}(s)\).
By \paperlink{premise:4d6174686c69622e5461637469632e507573684e65672e6e6f745f6578697374735f6571}{Prerequisite~P\ref{premise:4d6174686c69622e5461637469632e507573684e65672e6e6f745f6578697374735f6571}} with \(x_{1} \mapsto x_{1} \in s \land (\exists y_{1}\, :\, \alpha,\; y_{1} \in s \land (x_{1} \ne y_{1} \land f(x_{1}) = f(y_{1})))\) in place of its parameter \(s\), we obtain \((\neg (\exists x_{1}\, :\, \alpha,\; x_{1} \in s \land (\exists y_{1}\, :\, \alpha,\; y_{1} \in s \land (x_{1} \ne y_{1} \land f(x_{1}) = f(y_{1}))))) = (\forall x_{1}\, :\, \alpha,\; \neg (x_{1} \in s \land (\exists y_{1}\, :\, \alpha,\; y_{1} \in s \land (x_{1} \ne y_{1} \land f(x_{1}) = f(y_{1})))))\).
For the required implication, argue as follows.
\begin{proofstep}
Fix \(x_{1}\) in \(\alpha\).
By \paperlink{premise:4d6174686c69622e5461637469632e507573684e65672e6e6f745f616e645f6571}{Prerequisite~P\ref{premise:4d6174686c69622e5461637469632e507573684e65672e6e6f745f616e645f6571}} with \(p = x_{1} \in s\) and \(q = \exists y_{1}\, :\, \alpha,\; y_{1} \in s \land (x_{1} \ne y_{1} \land f(x_{1}) = f(y_{1}))\), we obtain \((\neg (x_{1} \in s \land (\exists y_{1}\, :\, \alpha,\; y_{1} \in s \land (x_{1} \ne y_{1} \land f(x_{1}) = f(y_{1}))))) = (x_{1} \in s \implies \neg (\exists y_{1}\, :\, \alpha,\; y_{1} \in s \land (x_{1} \ne y_{1} \land f(x_{1}) = f(y_{1}))))\).
By definition, we obtain \(\left(x_{1} \in s\right) = \left(x_{1} \in s\right)\).
By \paperlink{premise:4d6174686c69622e5461637469632e507573684e65672e6e6f745f6578697374735f6571}{Prerequisite~P\ref{premise:4d6174686c69622e5461637469632e507573684e65672e6e6f745f6578697374735f6571}} with \(y_{1} \mapsto y_{1} \in s \land \left(x_{1} \ne y_{1} \land f\left(x_{1}\right) = f\left(y_{1}\right)\right)\) in place of its parameter \(s\), we obtain \((\neg (\exists x_{2}\, :\, \alpha,\; x_{2} \in s \land (x_{1} \ne x_{2} \land f(x_{1}) = f(x_{2})))) = (\forall x_{2}\, :\, \alpha,\; \neg (x_{2} \in s \land (x_{1} \ne x_{2} \land f(x_{1}) = f(x_{2}))))\).
For the required implication, argue as follows.
\begin{proofstep}
Fix \(y_{1}\) in \(\alpha\).
By \paperlink{premise:4d6174686c69622e5461637469632e507573684e65672e6e6f745f616e645f6571}{Prerequisite~P\ref{premise:4d6174686c69622e5461637469632e507573684e65672e6e6f745f616e645f6571}} with \(p = y_{1} \in s\) and \(q = x_{1} \ne y_{1} \land f\left(x_{1}\right) = f\left(y_{1}\right)\), we obtain \((\neg (y_{1} \in s \land (x_{1} \ne y_{1} \land f(x_{1}) = f(y_{1})))) = (y_{1} \in s \implies \neg (x_{1} \ne y_{1} \land f(x_{1}) = f(y_{1})))\).
By definition, we obtain \(\left(y_{1} \in s\right) = \left(y_{1} \in s\right)\).
By \paperlink{premise:4d6174686c69622e5461637469632e507573684e65672e6e6f745f616e645f6571}{Prerequisite~P\ref{premise:4d6174686c69622e5461637469632e507573684e65672e6e6f745f616e645f6571}} with \(p = x_{1} \ne y_{1}\) and \(q = f\left(x_{1}\right) = f\left(y_{1}\right)\), we obtain \((\neg (x_{1} \ne y_{1} \land f(x_{1}) = f(y_{1}))) = (x_{1} \ne y_{1} \implies \neg (f(x_{1}) = f(y_{1})))\).
By \paperlink{premise:696d706c6965735f636f6e6772}{Prerequisite~P\ref{premise:696d706c6965735f636f6e6772}} with \(p_{1} = y_{1} \in s\) and \(p_{2} = y_{1} \in s\) and \(q_{1} = \neg \left(x_{1} \ne y_{1} \land f\left(x_{1}\right) = f\left(y_{1}\right)\right)\) and \(q_{2} = x_{1} \ne y_{1} \implies \neg \left(f\left(x_{1}\right) = f\left(y_{1}\right)\right)\), we obtain \((y_{1} \in s \implies \neg (x_{1} \ne y_{1} \land f(x_{1}) = f(y_{1}))) = (y_{1} \in s \implies x_{1} \ne y_{1} \implies \neg (f(x_{1}) = f(y_{1})))\).
Chaining the equalities, we obtain \((\neg (y_{1} \in s \land (x_{1} \ne y_{1} \land f(x_{1}) = f(y_{1})))) = (y_{1} \in s \implies x_{1} \ne y_{1} \implies \neg (f(x_{1}) = f(y_{1})))\).
\end{proofstep}
By \paperlink{premise:666f72616c6c5f636f6e6772}{Prerequisite~P\ref{premise:666f72616c6c5f636f6e6772}}, we obtain \((\forall a\, :\, \alpha,\; \neg (a \in s \land (x_{1} \ne a \land f(x_{1}) = f(a)))) = (\forall a\, :\, \alpha,\; a \in s \implies x_{1} \ne a \implies \neg (f(x_{1}) = f(a)))\).
Chaining the equalities, we obtain \((\neg (\exists x_{2}\, :\, \alpha,\; x_{2} \in s \land (x_{1} \ne x_{2} \land f(x_{1}) = f(x_{2})))) = (\forall a\, :\, \alpha,\; a \in s \implies x_{1} \ne a \implies \neg (f(x_{1}) = f(a)))\).
By \paperlink{premise:696d706c6965735f636f6e6772}{Prerequisite~P\ref{premise:696d706c6965735f636f6e6772}} with \(p_{1} = x_{1} \in s\) and \(p_{2} = x_{1} \in s\) and \(q_{1} = \neg (\exists x_{2}\, :\, \alpha,\; x_{2} \in s \land (x_{1} \ne x_{2} \land f(x_{1}) = f(x_{2})))\) and \(q_{2} = \forall a\, :\, \alpha,\; a \in s \implies x_{1} \ne a \implies \neg (f(x_{1}) = f(a))\), we obtain \((x_{1} \in s \implies \neg (\exists x_{2}\, :\, \alpha,\; x_{2} \in s \land (x_{1} \ne x_{2} \land f(x_{1}) = f(x_{2})))) = (x_{1} \in s \implies (\forall a_{1}\, :\, \alpha,\; a_{1} \in s \implies x_{1} \ne a_{1} \implies \neg (f(x_{1}) = f(a_{1}))))\).
Chaining the equalities, we obtain \((\neg (x_{1} \in s \land (\exists y_{1}\, :\, \alpha,\; y_{1} \in s \land (x_{1} \ne y_{1} \land f(x_{1}) = f(y_{1}))))) = (x_{1} \in s \implies (\forall a_{1}\, :\, \alpha,\; a_{1} \in s \implies x_{1} \ne a_{1} \implies \neg (f(x_{1}) = f(a_{1}))))\).
\end{proofstep}
By \paperlink{premise:666f72616c6c5f636f6e6772}{Prerequisite~P\ref{premise:666f72616c6c5f636f6e6772}}, we obtain \((\forall a\, :\, \alpha,\; \neg (a \in s \land (\exists y_{1}\, :\, \alpha,\; y_{1} \in s \land (a \ne y_{1} \land f(a) = f(y_{1}))))) = (\forall a\, :\, \alpha,\; a \in s \implies (\forall a_{2}\, :\, \alpha,\; a_{2} \in s \implies a \ne a_{2} \implies \neg (f(a) = f(a_{2}))))\).
Chaining the equalities, we obtain \((\neg (\exists x_{1}\, :\, \alpha,\; x_{1} \in s \land (\exists y_{1}\, :\, \alpha,\; y_{1} \in s \land (x_{1} \ne y_{1} \land f(x_{1}) = f(y_{1}))))) = (\forall a\, :\, \alpha,\; a \in s \implies (\forall a_{2}\, :\, \alpha,\; a_{2} \in s \implies a \ne a_{2} \implies \neg (f(a) = f(a_{2}))))\).
Use the hypothesis \(\neg (\exists x\, :\, \alpha,\; x \in s \land (\exists y\, :\, \alpha,\; y \in s \land (x \ne y \land f(x) = f(y))))\).
Use the hypothesis \(x \in \operatorname{underlying\,set}(s)\).
Use the hypothesis \(y \in \operatorname{underlying\,set}(s)\).
Substituting along the equality, we obtain \(x \ne y \implies f\left(x\right) \ne f\left(y\right)\).
By \paperlink{premise:4d6174686c69622e5461637469632e436f6e747261706f73652e6d7472}{Prerequisite~P\ref{premise:4d6174686c69622e5461637469632e436f6e747261706f73652e6d7472}}, we obtain \(f\left(x\right) = f\left(y\right) \implies x = y\).
\end{proofstep}
By \paperlink{theorem:73}{Lemma~\ref{theorem:73}}, we obtain \(\left|s\right| \le \left|t\right|\).
By \paperlink{premise:4c542e6c742e6e6f745f6c65}{\emph{Alias}} (\paperlink{premise:4c542e6c742e6e6f745f6c65}{Prerequisite~P\ref{premise:4c542e6c742e6e6f745f6c65}}) with \(\mathbb{N}\) as its type parameter \(\alpha\), we obtain \(\mathrm{False}\).
\end{proofstep}
The contrary assumption leads to a contradiction. Hence \(\exists x\, :\, \alpha,\; x \in s \land (\exists y\, :\, \alpha,\; y \in s \land (x \ne y \land f(x) = f(y)))\).
\end{proof}

\begin{theorem}\label{theorem:77}
Let \(\alpha\) be a type. Let \(s\) be a finite subset of \(\alpha\). Let \(n\) be a natural number. Let \(f\) be a function of type \(\mathbb{N} \to \alpha\). Assume \(\forall i\, :\, \mathbb{N},\; i < n \implies f\left(i\right) \in s\). Assume \(\forall i\, :\, \mathbb{N},\; i < n \implies (\forall j\, :\, \mathbb{N},\; j < n \implies f(i) = f(j) \implies i = j)\). Then \(n \le \left|s\right|\).
\end{theorem}

\begin{proof}
Using \paperlink{theorem:34}{Lemma~\ref{theorem:34}} in reverse gives \(n = \left|\{i\in\mathbb{N}\mid i<n\}\right|\).
The substitution uses the following argument.
\begin{proofstep}
Fix \(a\) in \(\mathbb{N}\).
By \paperlink{premise:46696e7365742e6d656d5f72616e6765}{Prerequisite~P\ref{premise:46696e7365742e6d656d5f72616e6765}}, we obtain \(a \in \{i\in\mathbb{N}\mid i<n\} \iff a < n\).
By propositional extensionality, we obtain \(\left(a \in \{i\in\mathbb{N}\mid i<n\}\right) = \left(a < n\right)\).
By definition, we obtain \(\left(f\left(a\right) \in s\right) = \left(f\left(a\right) \in s\right)\).
By \paperlink{premise:696d706c6965735f636f6e6772}{Prerequisite~P\ref{premise:696d706c6965735f636f6e6772}} with \(p_{1} = a \in \{i\in\mathbb{N}\mid i<n\}\) and \(p_{2} = a < n\) and \(q_{1} = f\left(a\right) \in s\) and \(q_{2} = f\left(a\right) \in s\), we obtain \((a \in \{i\in\mathbb{N}\mid i<n\} \implies f(a) \in s) = (a < n \implies f(a) \in s)\).
\end{proofstep}
Use the hypothesis \(\forall i\, :\, \mathbb{N},\; i < n \implies f\left(i\right) \in s\).
Rewriting with \paperlink{premise:666f72616c6c5f636f6e6772}{Prerequisite~P\ref{premise:666f72616c6c5f636f6e6772}} and the identity established above in the preceding result gives \(\forall a\, :\, \mathbb{N},\; a \in \{i\in\mathbb{N}\mid i<n\} \implies f\left(a\right) \in s\).
Use the hypothesis \(\forall i\, :\, \mathbb{N},\; i < n \implies (\forall j\, :\, \mathbb{N},\; j < n \implies f(i) = f(j) \implies i = j)\), which simplifies to \(f\text{ is injective on }\operatorname{Iio}\left(n\right)\).
Rewriting with \paperlink{premise:46696e7365742e636f655f72616e6765}{Prerequisite~P\ref{premise:46696e7365742e636f655f72616e6765}} in the preceding result gives \(f\text{ is injective on }\operatorname{underlying\,set}(\{i\in\mathbb{N}\mid i<n\})\).
By \paperlink{theorem:73}{Lemma~\ref{theorem:73}} with \(\mathbb{N}\) as its type parameter \(\alpha\), we obtain \(\left|\{i\in\mathbb{N}\mid i<n\}\right| \le \left|s\right|\).
By transitivity of the displayed relations, we obtain \(n \le \left|s\right|\).
\end{proof}

\begin{lemma}\label{theorem:78}
Let \(\alpha\) and \(\beta\) be types. Let \(s\) be a finite subset of \(\alpha\). Let \(t\) be a finite subset of \(\beta\). Let \(f\) be a function of type \(\alpha \to \beta\). Assume \(\operatorname{Maps\,To}\left(f, \operatorname{underlying\,set}(s), \operatorname{underlying\,set}(t)\right)\). Assume \(f\text{ is injective on }\operatorname{underlying\,set}(s)\). Assume \(\left|t\right| \le \left|s\right|\). Then \(\operatorname{underlying\,set}(t)\subseteq f\left[\operatorname{underlying\,set}(s)\right]\).
\end{lemma}

\begin{proof}
The substitution uses the following argument.
\begin{proofstep}
Fix \(x\) in \(\beta\).
By \paperlink{premise:46696e7365742e6d656d5f696d616765}{Prerequisite~P\ref{premise:46696e7365742e6d656d5f696d616765}}, we obtain \(x \in f\left[s\right] \iff \left(\exists a\, :\, \alpha,\; a \in s \land f\left(a\right) = x\right)\).
By propositional extensionality, we obtain \(\left(x \in f\left[s\right]\right) = \left(\exists a\, :\, \alpha,\; a \in s \land f\left(a\right) = x\right)\).
By definition, we obtain \(\left(x \in t\right) = \left(x \in t\right)\).
By \paperlink{premise:696d706c6965735f636f6e6772}{Prerequisite~P\ref{premise:696d706c6965735f636f6e6772}} with \(p_{1} = x \in f\left[s\right]\) and \(p_{2} = \exists a\, :\, \alpha,\; a \in s \land f\left(a\right) = x\) and \(q_{1} = x \in t\) and \(q_{2} = x \in t\), we obtain \((x \in f[s] \implies x \in t) = ((\exists a\, :\, \alpha,\; a \in s \land f(a) = x) \implies x \in t)\).
By \paperlink{premise:666f72616c6c5f6578697374735f696e646578}{Prerequisite~P\ref{premise:666f72616c6c5f6578697374735f696e646578}}, we obtain \((\forall h\, :\, \exists x_{1}\, :\, \alpha,\; (x_{2} \mapsto x_{2} \in s \land f(x_{2}) = x)(x_{1}),\; (h_{1} \mapsto x \in t)(h)) \iff (\forall x_{1}\, :\, \alpha,\; \forall h\, :\, (x_{2} \mapsto x_{2} \in s \land f(x_{2}) = x)(x_{1}),\; (h_{1} \mapsto x \in t)(\operatorname{Exists\mathord{:}intro}(x_{1},\allowbreak{} h)))\).
By propositional extensionality, we obtain \(((\exists x_{1}\, :\, \alpha,\; x_{1} \in s \land f(x_{1}) = x) \implies x \in t) = (\forall x_{1}\, :\, \alpha,\; x_{1} \in s \land f(x_{1}) = x \implies x \in t)\).
Chaining the equalities, we obtain \((x \in f[s] \implies x \in t) = (\forall x_{1}\, :\, \alpha,\; x_{1} \in s \land f(x_{1}) = x \implies x \in t)\).
For the required implication, argue as follows.
\begin{proofstep}
Fix \(x_{1}\) in \(\alpha\).
By \paperlink{premise:616e645f696d70}{Prerequisite~P\ref{premise:616e645f696d70}}, we obtain \(x_{1} \in s \land f(x_{1}) = x \implies x \in t \iff x_{1} \in s \implies f(x_{1}) = x \implies x \in t\).
By propositional extensionality, we obtain \((x_{1} \in s \land f(x_{1}) = x \implies x \in t) = (x_{1} \in s \implies f(x_{1}) = x \implies x \in t)\).
\end{proofstep}
By \paperlink{premise:666f72616c6c5f636f6e6772}{Prerequisite~P\ref{premise:666f72616c6c5f636f6e6772}}, we obtain \((\forall a\, :\, \alpha,\; a \in s \land f(a) = x \implies x \in t) = (\forall a\, :\, \alpha,\; a \in s \implies f(a) = x \implies x \in t)\).
Chaining the equalities, we obtain \((x \in f[s] \implies x \in t) = (\forall a\, :\, \alpha,\; a \in s \implies f(a) = x \implies x \in t)\).
\end{proofstep}
Fix \(a\) in \(\alpha\).
Assume \(a \in s\).
Rewriting with \paperlink{premise:46696e7365742e6d656d5f636f65}{Prerequisite~P\ref{premise:46696e7365742e6d656d5f636f65}} and \(a \in s\) gives \(a \in \operatorname{underlying\,set}(s)\).
Applying the hypothesis \(\operatorname{Maps\,To}\left(f, \operatorname{underlying\,set}(s), \operatorname{underlying\,set}(t)\right)\), we obtain \(f\left(a\right) \in t\).
Rewriting with \paperlink{premise:46696e7365742e7375627365745f696666}{Prerequisite~P\ref{premise:46696e7365742e7375627365745f696666}}, \paperlink{premise:666f72616c6c5f636f6e6772}{Prerequisite~P\ref{premise:666f72616c6c5f636f6e6772}}, \paperlink{premise:666f72616c6c5f6170706c795f65715f696d705f696666e28282}{Prerequisite~P\ref{premise:666f72616c6c5f6170706c795f65715f696d705f696666e28282}} and the identity established above in the preceding result gives \(f\left[s\right] \subseteq t\).
Use the hypothesis \(\left|t\right| \le \left|s\right|\).
Using \paperlink{theorem:46}{Lemma~\ref{theorem:46}} in reverse under \(f\text{ is injective on }\operatorname{underlying\,set}(s)\) gives \(\left|s\right| = \left|f\left[s\right]\right|\).
By \paperlink{premise:4c452e6c652e7472616e735f6571}{\emph{Alias}} (\paperlink{premise:4c452e6c652e7472616e735f6571}{Prerequisite~P\ref{premise:4c452e6c652e7472616e735f6571}}) with \(\mathbb{N}\) as its type parameter \(\alpha\), we obtain \(\left|t\right| \le \left|f\left[s\right]\right|\).
By \paperlink{theorem:55}{Lemma~\ref{theorem:55}}, we obtain \(f\left[s\right] = t\).
Rewriting with \(f\left[s\right] = t\), \paperlink{premise:46696e7365742e636f655f696d616765}{Prerequisite~P\ref{premise:46696e7365742e636f655f696d616765}}, \paperlink{premise:7375627365745f7265666c}{Prerequisite~P\ref{premise:7375627365745f7265666c}} and \paperlink{premise:5365742e696e737449735265666c537562736574}{Prerequisite~P\ref{premise:5365742e696e737449735265666c537562736574}} gives \(\operatorname{underlying\,set}(t)\subseteq f\left[\operatorname{underlying\,set}(s)\right]\).
\end{proof}

\begin{theorem}\label{theorem:79}
Let \(\alpha\) and \(\beta\) be types. Let \(s\) be a finite subset of \(\alpha\). Let \(t\) be a finite subset of \(\beta\). Let \(f\) be a function of type \(\Pi_{a\, :\, \alpha}\, a \in s \to \beta\). Assume \(\forall a\, :\, \alpha,\; \forall \mathit{ha}\, :\, a \in s,\; f\left(a, \mathit{ha}\right) \in t\). Assume \(\forall a_{1}\, :\, \alpha,\; \forall a_{2}\, :\, \alpha,\; \forall h\, :\, a_{1} \in s,\; \forall h_{1}\, :\, a_{2} \in s,\; f(a_{1},\allowbreak{} h) = f(a_{2},\allowbreak{} h_{1}) \implies a_{1} = a_{2}\). Assume \(\left|t\right| \le \left|s\right|\). Let \(b\) be an element of \(\beta\). Assume \(b \in t\). Then
\begin{center}
\(\displaystyle \exists a_{1}\, :\, \alpha,\; \exists \mathit{ha}\text{ proving }(a_{1} \in s),\; b = f(a_{1},\allowbreak{} \mathit{ha}).\)
\end{center}
\end{theorem}

\begin{proof}
Set \(f^{'}\) equal to \(a \mapsto f\left(a, \operatorname{property}\left(a\right)\right)\).
Fix \(a\) in \(s\).
Fix \(x\) in \(s\).
Assume \(x \in \operatorname{underlying\,set}(\operatorname{Finset\mathord{:}attach}\left(s\right))\).
Fix \(y\) in \(s\).
Assume \(y \in \operatorname{underlying\,set}(\operatorname{Finset\mathord{:}attach}\left(s\right))\).
Assume \(f^{'}\left(x\right) = f^{'}\left(y\right)\).
By \paperlink{premise:537562747970652e70726f7065727479}{Prerequisite~P\ref{premise:537562747970652e70726f7065727479}} with \(\mathit{self} = x\), we obtain \(\operatorname{val}(x) \in s\).
By \paperlink{premise:537562747970652e70726f7065727479}{Prerequisite~P\ref{premise:537562747970652e70726f7065727479}} with \(\mathit{self} = y\), we obtain \(\operatorname{val}(y) \in s\).
Use the hypothesis \(f^{'}\left(x\right) = f^{'}\left(y\right)\).
Applying the hypothesis \(\forall a_{1}\, :\, \alpha,\; \forall a_{2}\, :\, \alpha,\; \forall h\, :\, a_{1} \in s,\; \forall h_{1}\, :\, a_{2} \in s,\; f(a_{1},\allowbreak{} h) = f(a_{2},\allowbreak{} h_{1}) \implies a_{1} = a_{2}\), we obtain \(\operatorname{val}(x) = \operatorname{val}(y)\).
By \paperlink{premise:537562747970652e657874}{Prerequisite~P\ref{premise:537562747970652e657874}}, we obtain \(x = y\).
Fix \(x\) in \(s\).
Assume \(x \in \operatorname{underlying\,set}(\operatorname{Finset\mathord{:}attach}\left(s\right))\).
By \paperlink{premise:537562747970652e70726f7065727479}{Prerequisite~P\ref{premise:537562747970652e70726f7065727479}} with \(\mathit{self} = x\), we obtain \(\operatorname{val}(x) \in s\).
Applying the hypothesis \(\forall a\, :\, \alpha,\; \forall \mathit{ha}\, :\, a \in s,\; f\left(a, \mathit{ha}\right) \in t\), we obtain \(f\left(\operatorname{val}(x), \operatorname{property}\left(x\right)\right) \in t\).
Use the hypothesis \(\left|t\right| \le \left|s\right|\).
Rewriting with \paperlink{theorem:35}{Lemma~\ref{theorem:35}} in the preceding result gives \(\left|t\right| \le \left|\operatorname{Finset\mathord{:}attach}\left(s\right)\right|\).
By \paperlink{theorem:78}{Lemma~\ref{theorem:78}} with \(s\) as its type parameter \(\alpha\) and \(f^{'}\) in place of its parameter \(f\), we obtain \(b \in f^{'}\left[\operatorname{underlying\,set}(\operatorname{Finset\mathord{:}attach}\left(s\right))\right]\).
Consider the following case.
\begin{proofstep}
Fix \(a_{1}\) in \(s\).
Assume \(a_{1} \in \operatorname{underlying\,set}(\operatorname{Finset\mathord{:}attach}(s)) \land f^{'}(a_{1}) = b\).
Use the hypothesis \(a_{1} \in \operatorname{underlying\,set}(\operatorname{Finset\mathord{:}attach}(s)) \land f^{'}(a_{1}) = b\).
Consider the following case.
\begin{proofstep}
Assume \(a_{1} \in \operatorname{underlying\,set}(\operatorname{Finset\mathord{:}attach}\left(s\right))\).
Assume \(f^{'}\left(a_{1}\right) = b\).
For the required implication, argue as follows.
\begin{proofstep}
Assume \(f^{'}\left(a_{1}\right) \in t\).
Choose the witness \(\operatorname{val}(a_{1})\).
The construction requires evidence that \(\operatorname{val}(a_{1}) \in s\).
By \paperlink{premise:537562747970652e70726f7065727479}{Prerequisite~P\ref{premise:537562747970652e70726f7065727479}} with \(\mathit{self} = a_{1}\), we obtain \(\operatorname{val}(a_{1}) \in s\).
By definition, we obtain \(f^{'}\left(a_{1}\right) = f^{'}\left(a_{1}\right)\).
This witness proves existence. Hence \(\exists \mathit{ha}_{1}\text{ proving }(\operatorname{val}(a_{1}) \in s),\; f^{'}(a_{1}) = f(\operatorname{val}(a_{1}),\allowbreak{} \mathit{ha}_{1})\).
This witness proves existence. Hence \(\exists a_{2}\, :\, \alpha,\; \exists \mathit{ha}_{1}\text{ proving }(a_{2} \in s),\; f^{'}(a_{1}) = f(a_{2},\allowbreak{} \mathit{ha}_{1})\).
\end{proofstep}
Use the hypothesis \(f^{'}\left(a_{1}\right) = b\).
Use the hypothesis \(b \in t\).
Substituting equals, we obtain \(\exists a_{2}\, :\, \alpha,\; \exists \mathit{ha}_{1}\text{ proving }(a_{2} \in s),\; b = f(a_{2},\allowbreak{} \mathit{ha}_{1})\).
\end{proofstep}
Using the two conjuncts, we obtain \(\exists a_{2}\, :\, \alpha,\; \exists \mathit{ha}\text{ proving }(a_{2} \in s),\; b = f(a_{2},\allowbreak{} \mathit{ha})\).
\end{proofstep}
Using the existential hypothesis, we obtain \(\exists a_{1}\, :\, \alpha,\; \exists \mathit{ha}\text{ proving }(a_{1} \in s),\; b = f(a_{1},\allowbreak{} \mathit{ha})\).
\end{proof}

\begin{lemma}\label{theorem:80}
Let \(\alpha\) and \(\beta\) be types. Let \(s\) be a finite subset of \(\alpha\). Let \(t\) be a finite subset of \(\beta\). Let \(f\) be a function of type \(\alpha \to \beta\). Assume \(\operatorname{Maps\,To}\left(f, \operatorname{underlying\,set}(s), \operatorname{underlying\,set}(t)\right)\). Assume \(\operatorname{underlying\,set}(t)\subseteq f\left[\operatorname{underlying\,set}(s)\right]\). Assume \(\left|s\right| \le \left|t\right|\). Then \(f\text{ is injective on }\operatorname{underlying\,set}(s)\).
\end{lemma}

\begin{proof}
Rewriting with \paperlink{premise:46696e7365742e636f655f696d616765}{Prerequisite~P\ref{premise:46696e7365742e636f655f696d616765}} and \paperlink{premise:5365742e5375726a4f6e2e696d6167655f65715f6f665f6d617073546f}{Prerequisite~P\ref{premise:5365742e5375726a4f6e2e696d6167655f65715f6f665f6d617073546f}} (using \(\operatorname{underlying\,set}(t)\subseteq f\left[\operatorname{underlying\,set}(s)\right]\) and \(\operatorname{Maps\,To}\left(f, \operatorname{underlying\,set}(s), \operatorname{underlying\,set}(t)\right)\)) gives \(\operatorname{underlying\,set}(f\left[s\right]) = \operatorname{underlying\,set}(t)\).
By \paperlink{premise:46696e7365742e636f655f696e6a656374697665}{Prerequisite~P\ref{premise:46696e7365742e636f655f696e6a656374697665}}, we obtain \(f\left[s\right] = t\).
Rewriting with \(f\left[s\right] = t\) gives \(\left|f\left[s\right]\right| = \left|t\right|\).
By \paperlink{theorem:45}{Corollary~\ref{theorem:45}}, we obtain \(\left|f\left[s\right]\right| \le \left|s\right|\).
Consider the following case.
\begin{proofstep}
Assume \(\neg \left(\left|f\left[s\right]\right| = \left|s\right|\right)\).
Use the hypothesis \(\neg \left(\left|f\left[s\right]\right| = \left|s\right|\right)\).
By \paperlink{premise:4e61742e6c745f6f725f67745f6f665f6e65}{Prerequisite~P\ref{premise:4e61742e6c745f6f725f67745f6f665f6e65}}, we obtain \(\left|f\left[s\right]\right| < \left|s\right| \lor \left|f\left[s\right]\right| > \left|s\right|\).
Consider the following case.
\begin{proofstep}
Assume \(\left|f\left[s\right]\right| < \left|s\right|\).
By definition, we obtain \(\operatorname{decide}(\operatorname{is\,Impossible}(\operatorname{Lean\mathord{:}Omega\mathord{:}Constraint\mathord{:}mk}(\operatorname{some}(1),\allowbreak{} \operatorname{some}(0))) = \mathrm{true}) = \mathrm{true}\).
By \paperlink{premise:6f665f6465636964655f65715f74727565}{Prerequisite~P\ref{premise:6f665f6465636964655f65715f74727565}}, we obtain \(\operatorname{is\,Impossible}(\operatorname{Lean\mathord{:}Omega\mathord{:}Constraint\mathord{:}mk}(\operatorname{some}(1),\allowbreak{} \operatorname{some}(0))) = \mathrm{true}\).
Use the previously proved fact \(\left|f\left[s\right]\right| \le \left|s\right|\).
By \paperlink{premise:4c65616e2e4f6d6567612e496e742e6f664e61745f6c655f6f665f6c65}{Prerequisite~P\ref{premise:4c65616e2e4f6d6567612e496e742e6f664e61745f6c655f6f665f6c65}}, we obtain \(\left|f\left[s\right]\right| \le \left|s\right|\).
By \paperlink{premise:496e742e7375625f6e6f6e6e65675f6f665f6c65}{Prerequisite~P\ref{premise:496e742e7375625f6e6f6e6e65675f6f665f6c65}}, we obtain \(0 \le \left|s\right| - \left|f\left[s\right]\right|\).
Using \paperlink{premise:4c65616e2e4f6d6567612e4c696e656172436f6d626f2e636f6f7264696e6174655f6576616c5f31}{Prerequisite~P\ref{premise:4c65616e2e4f6d6567612e4c696e656172436f6d626f2e636f6f7264696e6174655f6576616c5f31}} in reverse gives \(|s| = \operatorname{eval}(\operatorname{coordinate}(1),\allowbreak{} \operatorname{Lean\mathord{:}Omega\mathord{:}Coeffs\mathord{:}of\,List}(\langle |t|\rangle \mathbin{\Vert} (\langle |s|\rangle \mathbin{\Vert} \langle |f[s]|\rangle)))\).
Using \paperlink{premise:4c65616e2e4f6d6567612e4c696e656172436f6d626f2e636f6f7264696e6174655f6576616c5f32}{Prerequisite~P\ref{premise:4c65616e2e4f6d6567612e4c696e656172436f6d626f2e636f6f7264696e6174655f6576616c5f32}} in reverse gives \(|f[s]| = \operatorname{eval}(\operatorname{coordinate}(2),\allowbreak{} \operatorname{Lean\mathord{:}Omega\mathord{:}Coeffs\mathord{:}of\,List}(\langle |t|\rangle \mathbin{\Vert} (\langle |s|\rangle \mathbin{\Vert} \langle |f[s]|\rangle)))\).
By \paperlink{premise:4c65616e2e4f6d6567612e496e742e7375625f636f6e6772}{Prerequisite~P\ref{premise:4c65616e2e4f6d6567612e496e742e7375625f636f6e6772}}, we obtain \(|s| - |f[s]| = \operatorname{eval}(\operatorname{coordinate}(1),\allowbreak{} \operatorname{Lean\mathord{:}Omega\mathord{:}Coeffs\mathord{:}of\,List}(\langle |t|\rangle \mathbin{\Vert} (\langle |s|\rangle \mathbin{\Vert} \langle |f[s]|\rangle))) - \operatorname{eval}(\operatorname{coordinate}(2),\allowbreak{} \operatorname{Lean\mathord{:}Omega\mathord{:}Coeffs\mathord{:}of\,List}(\langle |t|\rangle \mathbin{\Vert} (\langle |s|\rangle \mathbin{\Vert} \langle |f[s]|\rangle)))\).
Using \paperlink{premise:4c65616e2e4f6d6567612e4c696e656172436f6d626f2e7375625f6576616c}{Prerequisite~P\ref{premise:4c65616e2e4f6d6567612e4c696e656172436f6d626f2e7375625f6576616c}} with \(l_{1} = \operatorname{Lean\mathord{:}Omega\mathord{:}Linear\,Combo\mathord{:}mk}(0,\allowbreak{} \langle 0\rangle \mathbin{\Vert} \langle 1\rangle)\) and \(l_{2} = \operatorname{Lean\mathord{:}Omega\mathord{:}Linear\,Combo\mathord{:}mk}(0,\allowbreak{} \langle 0\rangle \mathbin{\Vert} (\langle 0\rangle \mathbin{\Vert} \langle 1\rangle))\) and \(v = \operatorname{Lean\mathord{:}Omega\mathord{:}Coeffs\mathord{:}of\,List}(\langle |t|\rangle \mathbin{\Vert} (\langle |s|\rangle \mathbin{\Vert} \langle |f[s]|\rangle))\) in reverse gives \(\operatorname{eval}(\operatorname{Lean\mathord{:}Omega\mathord{:}Linear\,Combo\mathord{:}mk}(0,\allowbreak{} \langle 0\rangle \mathbin{\Vert} \langle 1\rangle),\allowbreak{} \operatorname{Lean\mathord{:}Omega\mathord{:}Coeffs\mathord{:}of\,List}(\langle |t|\rangle \mathbin{\Vert} (\langle |s|\rangle \mathbin{\Vert} \langle |f[s]|\rangle))) - \operatorname{eval}(\operatorname{Lean\mathord{:}Omega\mathord{:}Linear\,Combo\mathord{:}mk}(0,\allowbreak{} \langle 0\rangle \mathbin{\Vert} (\langle 0\rangle \mathbin{\Vert} \langle 1\rangle)),\allowbreak{} \operatorname{Lean\mathord{:}Omega\mathord{:}Coeffs\mathord{:}of\,List}(\langle |t|\rangle \mathbin{\Vert} (\langle |s|\rangle \mathbin{\Vert} \langle |f[s]|\rangle))) = \operatorname{eval}(\operatorname{Lean\mathord{:}Omega\mathord{:}Linear\,Combo\mathord{:}mk}(0,\allowbreak{} \langle 0\rangle \mathbin{\Vert} \langle 1\rangle) - \operatorname{Lean\mathord{:}Omega\mathord{:}Linear\,Combo\mathord{:}mk}(0,\allowbreak{} \langle 0\rangle \mathbin{\Vert} (\langle 0\rangle \mathbin{\Vert} \langle 1\rangle)),\allowbreak{} \operatorname{Lean\mathord{:}Omega\mathord{:}Coeffs\mathord{:}of\,List}(\langle |t|\rangle \mathbin{\Vert} (\langle |s|\rangle \mathbin{\Vert} \langle |f[s]|\rangle)))\).
Chaining the equalities, we obtain \(|s| - |f[s]| = \operatorname{eval}(\operatorname{Lean\mathord{:}Omega\mathord{:}Linear\,Combo\mathord{:}mk}(0,\allowbreak{} \langle 0\rangle \mathbin{\Vert} \langle 1\rangle) - \operatorname{Lean\mathord{:}Omega\mathord{:}Linear\,Combo\mathord{:}mk}(0,\allowbreak{} \langle 0\rangle \mathbin{\Vert} (\langle 0\rangle \mathbin{\Vert} \langle 1\rangle)),\allowbreak{} \operatorname{Lean\mathord{:}Omega\mathord{:}Coeffs\mathord{:}of\,List}(\langle |t|\rangle \mathbin{\Vert} (\langle |s|\rangle \mathbin{\Vert} \langle |f[s]|\rangle)))\).
By \paperlink{premise:6c655f6f665f6c655f6f665f6571}{Prerequisite~P\ref{premise:6c655f6f665f6c655f6f665f6571}} with \(\mathbb{Z}\) as its type parameter \(\alpha\), we obtain \(0 \le \operatorname{eval}(\operatorname{Lean\mathord{:}Omega\mathord{:}Linear\,Combo\mathord{:}mk}(0,\allowbreak{} \langle 0\rangle \mathbin{\Vert} \langle 1\rangle) - \operatorname{Lean\mathord{:}Omega\mathord{:}Linear\,Combo\mathord{:}mk}(0,\allowbreak{} \langle 0\rangle \mathbin{\Vert} (\langle 0\rangle \mathbin{\Vert} \langle 1\rangle)),\allowbreak{} \operatorname{Lean\mathord{:}Omega\mathord{:}Coeffs\mathord{:}of\,List}(\langle |t|\rangle \mathbin{\Vert} (\langle |s|\rangle \mathbin{\Vert} \langle |f[s]|\rangle)))\).
By \paperlink{premise:4c65616e2e4f6d6567612e436f6e73747261696e742e616464496e657175616c6974795f736174}{Prerequisite~P\ref{premise:4c65616e2e4f6d6567612e436f6e73747261696e742e616464496e657175616c6974795f736174}}, we obtain \(\operatorname{sat^{\prime}}(\operatorname{Lean\mathord{:}Omega\mathord{:}Constraint\mathord{:}mk}(\operatorname{some}(-0),\allowbreak{} \operatorname{none}),\allowbreak{} \langle 0\rangle \mathbin{\Vert} (\langle 1\rangle \mathbin{\Vert} \langle -1\rangle),\allowbreak{} \operatorname{Lean\mathord{:}Omega\mathord{:}Coeffs\mathord{:}of\,List}(\langle |t|\rangle \mathbin{\Vert} (\langle |s|\rangle \mathbin{\Vert} \langle |f[s]|\rangle))) = \mathrm{true}\).
Using \paperlink{premise:4c65616e2e4f6d6567612e4c696e656172436f6d626f2e636f6f7264696e6174655f6576616c5f32}{Prerequisite~P\ref{premise:4c65616e2e4f6d6567612e4c696e656172436f6d626f2e636f6f7264696e6174655f6576616c5f32}} in reverse gives \(|f[s]| = \operatorname{eval}(\operatorname{coordinate}(2),\allowbreak{} \operatorname{Lean\mathord{:}Omega\mathord{:}Coeffs\mathord{:}of\,List}(\langle |t|\rangle \mathbin{\Vert} (\langle |s|\rangle \mathbin{\Vert} \langle |f[s]|\rangle)))\).
Using \paperlink{premise:4c65616e2e4f6d6567612e4c696e656172436f6d626f2e636f6f7264696e6174655f6576616c5f30}{Prerequisite~P\ref{premise:4c65616e2e4f6d6567612e4c696e656172436f6d626f2e636f6f7264696e6174655f6576616c5f30}} in reverse gives \(|t| = \operatorname{eval}(\operatorname{coordinate}(0),\allowbreak{} \operatorname{Lean\mathord{:}Omega\mathord{:}Coeffs\mathord{:}of\,List}(\langle |t|\rangle \mathbin{\Vert} (\langle |s|\rangle \mathbin{\Vert} \langle |f[s]|\rangle)))\).
By \paperlink{premise:4c65616e2e4f6d6567612e496e742e7375625f636f6e6772}{Prerequisite~P\ref{premise:4c65616e2e4f6d6567612e496e742e7375625f636f6e6772}}, we obtain \(|f[s]| - |t| = \operatorname{eval}(\operatorname{coordinate}(2),\allowbreak{} \operatorname{Lean\mathord{:}Omega\mathord{:}Coeffs\mathord{:}of\,List}(\langle |t|\rangle \mathbin{\Vert} (\langle |s|\rangle \mathbin{\Vert} \langle |f[s]|\rangle))) - \operatorname{eval}(\operatorname{coordinate}(0),\allowbreak{} \operatorname{Lean\mathord{:}Omega\mathord{:}Coeffs\mathord{:}of\,List}(\langle |t|\rangle \mathbin{\Vert} (\langle |s|\rangle \mathbin{\Vert} \langle |f[s]|\rangle)))\).
Using \paperlink{premise:4c65616e2e4f6d6567612e4c696e656172436f6d626f2e7375625f6576616c}{Prerequisite~P\ref{premise:4c65616e2e4f6d6567612e4c696e656172436f6d626f2e7375625f6576616c}} with \(l_{1} = \operatorname{Lean\mathord{:}Omega\mathord{:}Linear\,Combo\mathord{:}mk}(0,\allowbreak{} \langle 0\rangle \mathbin{\Vert} (\langle 0\rangle \mathbin{\Vert} \langle 1\rangle))\) and \(l_{2} = \operatorname{Lean\mathord{:}Omega\mathord{:}Linear\,Combo\mathord{:}mk}(0,\allowbreak{} \langle 1\rangle)\) and \(v = \operatorname{Lean\mathord{:}Omega\mathord{:}Coeffs\mathord{:}of\,List}(\langle |t|\rangle \mathbin{\Vert} (\langle |s|\rangle \mathbin{\Vert} \langle |f[s]|\rangle))\) in reverse gives \(\operatorname{eval}(\operatorname{Lean\mathord{:}Omega\mathord{:}Linear\,Combo\mathord{:}mk}(0,\allowbreak{} \langle 0\rangle \mathbin{\Vert} (\langle 0\rangle \mathbin{\Vert} \langle 1\rangle)),\allowbreak{} \operatorname{Lean\mathord{:}Omega\mathord{:}Coeffs\mathord{:}of\,List}(\langle |t|\rangle \mathbin{\Vert} (\langle |s|\rangle \mathbin{\Vert} \langle |f[s]|\rangle))) - \operatorname{eval}(\operatorname{Lean\mathord{:}Omega\mathord{:}Linear\,Combo\mathord{:}mk}(0,\allowbreak{} \langle 1\rangle),\allowbreak{} \operatorname{Lean\mathord{:}Omega\mathord{:}Coeffs\mathord{:}of\,List}(\langle |t|\rangle \mathbin{\Vert} (\langle |s|\rangle \mathbin{\Vert} \langle |f[s]|\rangle))) = \operatorname{eval}(\operatorname{Lean\mathord{:}Omega\mathord{:}Linear\,Combo\mathord{:}mk}(0,\allowbreak{} \langle 0\rangle \mathbin{\Vert} (\langle 0\rangle \mathbin{\Vert} \langle 1\rangle)) - \operatorname{Lean\mathord{:}Omega\mathord{:}Linear\,Combo\mathord{:}mk}(0,\allowbreak{} \langle 1\rangle),\allowbreak{} \operatorname{Lean\mathord{:}Omega\mathord{:}Coeffs\mathord{:}of\,List}(\langle |t|\rangle \mathbin{\Vert} (\langle |s|\rangle \mathbin{\Vert} \langle |f[s]|\rangle)))\).
Chaining the equalities, we obtain \(|f[s]| - |t| = \operatorname{eval}(\operatorname{Lean\mathord{:}Omega\mathord{:}Linear\,Combo\mathord{:}mk}(0,\allowbreak{} \langle 0\rangle \mathbin{\Vert} (\langle 0\rangle \mathbin{\Vert} \langle 1\rangle)) - \operatorname{Lean\mathord{:}Omega\mathord{:}Linear\,Combo\mathord{:}mk}(0,\allowbreak{} \langle 1\rangle),\allowbreak{} \operatorname{Lean\mathord{:}Omega\mathord{:}Coeffs\mathord{:}of\,List}(\langle |t|\rangle \mathbin{\Vert} (\langle |s|\rangle \mathbin{\Vert} \langle |f[s]|\rangle)))\).
Reversing the equality, we obtain \(\operatorname{eval}(\operatorname{Lean\mathord{:}Omega\mathord{:}Linear\,Combo\mathord{:}mk}(0,\allowbreak{} \langle 0\rangle \mathbin{\Vert} (\langle 0\rangle \mathbin{\Vert} \langle 1\rangle)) - \operatorname{Lean\mathord{:}Omega\mathord{:}Linear\,Combo\mathord{:}mk}(0,\allowbreak{} \langle 1\rangle),\allowbreak{} \operatorname{Lean\mathord{:}Omega\mathord{:}Coeffs\mathord{:}of\,List}(\langle |t|\rangle \mathbin{\Vert} (\langle |s|\rangle \mathbin{\Vert} \langle |f[s]|\rangle))) = |f[s]| - |t|\).
Use the previously proved fact \(\left|f\left[s\right]\right| = \left|t\right|\).
By \paperlink{premise:4c65616e2e4f6d6567612e496e742e6f664e61745f636f6e6772}{Prerequisite~P\ref{premise:4c65616e2e4f6d6567612e496e742e6f664e61745f636f6e6772}}, we obtain \(\left|f\left[s\right]\right| = \left|t\right|\).
By \paperlink{premise:496e742e7375625f65715f7a65726f5f6f665f6571}{Prerequisite~P\ref{premise:496e742e7375625f65715f7a65726f5f6f665f6571}}, we obtain \(\left|f\left[s\right]\right| - \left|t\right| = 0\).
Chaining the equalities, we obtain \(\operatorname{eval}(\operatorname{Lean\mathord{:}Omega\mathord{:}Linear\,Combo\mathord{:}mk}(0,\allowbreak{} \langle 0\rangle \mathbin{\Vert} (\langle 0\rangle \mathbin{\Vert} \langle 1\rangle)) - \operatorname{Lean\mathord{:}Omega\mathord{:}Linear\,Combo\mathord{:}mk}(0,\allowbreak{} \langle 1\rangle),\allowbreak{} \operatorname{Lean\mathord{:}Omega\mathord{:}Coeffs\mathord{:}of\,List}(\langle |t|\rangle \mathbin{\Vert} (\langle |s|\rangle \mathbin{\Vert} \langle |f[s]|\rangle))) = 0\).
By \paperlink{premise:4c65616e2e4f6d6567612e436f6e73747261696e742e616464457175616c6974795f736174}{Prerequisite~P\ref{premise:4c65616e2e4f6d6567612e436f6e73747261696e742e616464457175616c6974795f736174}}, we obtain \(\operatorname{sat^{\prime}}(\operatorname{Lean\mathord{:}Omega\mathord{:}Constraint\mathord{:}mk}(\operatorname{some}(-0),\allowbreak{} \operatorname{some}(-0)),\allowbreak{} \langle -1\rangle \mathbin{\Vert} (\langle 0\rangle \mathbin{\Vert} \langle 1\rangle),\allowbreak{} \operatorname{Lean\mathord{:}Omega\mathord{:}Coeffs\mathord{:}of\,List}(\langle |t|\rangle \mathbin{\Vert} (\langle |s|\rangle \mathbin{\Vert} \langle |f[s]|\rangle))) = \mathrm{true}\).
By \paperlink{premise:4c65616e2e4f6d6567612e746964795f736174}{Prerequisite~P\ref{premise:4c65616e2e4f6d6567612e746964795f736174}}, we obtain \(\operatorname{sat^{\prime}}(\operatorname{tidy\,Constraint}(\operatorname{Lean\mathord{:}Omega\mathord{:}Constraint\mathord{:}mk}(\operatorname{some}(0),\allowbreak{} \operatorname{some}(0)),\allowbreak{} \langle -1\rangle \mathbin{\Vert} (\langle 0\rangle \mathbin{\Vert} \langle 1\rangle)),\allowbreak{} \operatorname{tidy\,Coeffs}(\operatorname{Lean\mathord{:}Omega\mathord{:}Constraint\mathord{:}mk}(\operatorname{some}(0),\allowbreak{} \operatorname{some}(0)),\allowbreak{} \langle -1\rangle \mathbin{\Vert} (\langle 0\rangle \mathbin{\Vert} \langle 1\rangle)),\allowbreak{} \operatorname{Lean\mathord{:}Omega\mathord{:}Coeffs\mathord{:}of\,List}(\langle |t|\rangle \mathbin{\Vert} (\langle |s|\rangle \mathbin{\Vert} \langle |f[s]|\rangle))) = \mathrm{true}\).
Use the hypothesis \(\left|s\right| \le \left|t\right|\).
By \paperlink{premise:4c65616e2e4f6d6567612e496e742e6f664e61745f6c655f6f665f6c65}{Prerequisite~P\ref{premise:4c65616e2e4f6d6567612e496e742e6f664e61745f6c655f6f665f6c65}}, we obtain \(\left|s\right| \le \left|t\right|\).
By \paperlink{premise:496e742e7375625f6e6f6e6e65675f6f665f6c65}{Prerequisite~P\ref{premise:496e742e7375625f6e6f6e6e65675f6f665f6c65}}, we obtain \(0 \le \left|t\right| - \left|s\right|\).
Using \paperlink{premise:4c65616e2e4f6d6567612e4c696e656172436f6d626f2e636f6f7264696e6174655f6576616c5f30}{Prerequisite~P\ref{premise:4c65616e2e4f6d6567612e4c696e656172436f6d626f2e636f6f7264696e6174655f6576616c5f30}} in reverse gives \(|t| = \operatorname{eval}(\operatorname{coordinate}(0),\allowbreak{} \operatorname{Lean\mathord{:}Omega\mathord{:}Coeffs\mathord{:}of\,List}(\langle |t|\rangle \mathbin{\Vert} (\langle |s|\rangle \mathbin{\Vert} \langle |f[s]|\rangle)))\).
Using \paperlink{premise:4c65616e2e4f6d6567612e4c696e656172436f6d626f2e636f6f7264696e6174655f6576616c5f31}{Prerequisite~P\ref{premise:4c65616e2e4f6d6567612e4c696e656172436f6d626f2e636f6f7264696e6174655f6576616c5f31}} in reverse gives \(|s| = \operatorname{eval}(\operatorname{coordinate}(1),\allowbreak{} \operatorname{Lean\mathord{:}Omega\mathord{:}Coeffs\mathord{:}of\,List}(\langle |t|\rangle \mathbin{\Vert} (\langle |s|\rangle \mathbin{\Vert} \langle |f[s]|\rangle)))\).
By \paperlink{premise:4c65616e2e4f6d6567612e496e742e7375625f636f6e6772}{Prerequisite~P\ref{premise:4c65616e2e4f6d6567612e496e742e7375625f636f6e6772}}, we obtain \(|t| - |s| = \operatorname{eval}(\operatorname{coordinate}(0),\allowbreak{} \operatorname{Lean\mathord{:}Omega\mathord{:}Coeffs\mathord{:}of\,List}(\langle |t|\rangle \mathbin{\Vert} (\langle |s|\rangle \mathbin{\Vert} \langle |f[s]|\rangle))) - \operatorname{eval}(\operatorname{coordinate}(1),\allowbreak{} \operatorname{Lean\mathord{:}Omega\mathord{:}Coeffs\mathord{:}of\,List}(\langle |t|\rangle \mathbin{\Vert} (\langle |s|\rangle \mathbin{\Vert} \langle |f[s]|\rangle)))\).
Using \paperlink{premise:4c65616e2e4f6d6567612e4c696e656172436f6d626f2e7375625f6576616c}{Prerequisite~P\ref{premise:4c65616e2e4f6d6567612e4c696e656172436f6d626f2e7375625f6576616c}} with \(l_{1} = \operatorname{Lean\mathord{:}Omega\mathord{:}Linear\,Combo\mathord{:}mk}(0,\allowbreak{} \langle 1\rangle)\) and \(l_{2} = \operatorname{Lean\mathord{:}Omega\mathord{:}Linear\,Combo\mathord{:}mk}(0,\allowbreak{} \langle 0\rangle \mathbin{\Vert} \langle 1\rangle)\) and \(v = \operatorname{Lean\mathord{:}Omega\mathord{:}Coeffs\mathord{:}of\,List}(\langle |t|\rangle \mathbin{\Vert} (\langle |s|\rangle \mathbin{\Vert} \langle |f[s]|\rangle))\) in reverse gives \(\operatorname{eval}(\operatorname{Lean\mathord{:}Omega\mathord{:}Linear\,Combo\mathord{:}mk}(0,\allowbreak{} \langle 1\rangle),\allowbreak{} \operatorname{Lean\mathord{:}Omega\mathord{:}Coeffs\mathord{:}of\,List}(\langle |t|\rangle \mathbin{\Vert} (\langle |s|\rangle \mathbin{\Vert} \langle |f[s]|\rangle))) - \operatorname{eval}(\operatorname{Lean\mathord{:}Omega\mathord{:}Linear\,Combo\mathord{:}mk}(0,\allowbreak{} \langle 0\rangle \mathbin{\Vert} \langle 1\rangle),\allowbreak{} \operatorname{Lean\mathord{:}Omega\mathord{:}Coeffs\mathord{:}of\,List}(\langle |t|\rangle \mathbin{\Vert} (\langle |s|\rangle \mathbin{\Vert} \langle |f[s]|\rangle))) = \operatorname{eval}(\operatorname{Lean\mathord{:}Omega\mathord{:}Linear\,Combo\mathord{:}mk}(0,\allowbreak{} \langle 1\rangle) - \operatorname{Lean\mathord{:}Omega\mathord{:}Linear\,Combo\mathord{:}mk}(0,\allowbreak{} \langle 0\rangle \mathbin{\Vert} \langle 1\rangle),\allowbreak{} \operatorname{Lean\mathord{:}Omega\mathord{:}Coeffs\mathord{:}of\,List}(\langle |t|\rangle \mathbin{\Vert} (\langle |s|\rangle \mathbin{\Vert} \langle |f[s]|\rangle)))\).
Chaining the equalities, we obtain \(|t| - |s| = \operatorname{eval}(\operatorname{Lean\mathord{:}Omega\mathord{:}Linear\,Combo\mathord{:}mk}(0,\allowbreak{} \langle 1\rangle) - \operatorname{Lean\mathord{:}Omega\mathord{:}Linear\,Combo\mathord{:}mk}(0,\allowbreak{} \langle 0\rangle \mathbin{\Vert} \langle 1\rangle),\allowbreak{} \operatorname{Lean\mathord{:}Omega\mathord{:}Coeffs\mathord{:}of\,List}(\langle |t|\rangle \mathbin{\Vert} (\langle |s|\rangle \mathbin{\Vert} \langle |f[s]|\rangle)))\).
By \paperlink{premise:6c655f6f665f6c655f6f665f6571}{Prerequisite~P\ref{premise:6c655f6f665f6c655f6f665f6571}} with \(\mathbb{Z}\) as its type parameter \(\alpha\), we obtain \(0 \le \operatorname{eval}(\operatorname{Lean\mathord{:}Omega\mathord{:}Linear\,Combo\mathord{:}mk}(0,\allowbreak{} \langle 1\rangle) - \operatorname{Lean\mathord{:}Omega\mathord{:}Linear\,Combo\mathord{:}mk}(0,\allowbreak{} \langle 0\rangle \mathbin{\Vert} \langle 1\rangle),\allowbreak{} \operatorname{Lean\mathord{:}Omega\mathord{:}Coeffs\mathord{:}of\,List}(\langle |t|\rangle \mathbin{\Vert} (\langle |s|\rangle \mathbin{\Vert} \langle |f[s]|\rangle)))\).
By \paperlink{premise:4c65616e2e4f6d6567612e436f6e73747261696e742e616464496e657175616c6974795f736174}{Prerequisite~P\ref{premise:4c65616e2e4f6d6567612e436f6e73747261696e742e616464496e657175616c6974795f736174}}, we obtain \(\operatorname{sat^{\prime}}(\operatorname{Lean\mathord{:}Omega\mathord{:}Constraint\mathord{:}mk}(\operatorname{some}(-0),\allowbreak{} \operatorname{none}),\allowbreak{} \langle 1\rangle \mathbin{\Vert} \langle -1\rangle,\allowbreak{} \operatorname{Lean\mathord{:}Omega\mathord{:}Coeffs\mathord{:}of\,List}(\langle |t|\rangle \mathbin{\Vert} (\langle |s|\rangle \mathbin{\Vert} \langle |f[s]|\rangle))) = \mathrm{true}\).
By \paperlink{premise:4c65616e2e4f6d6567612e636f6d626f5f73617427}{Prerequisite~P\ref{premise:4c65616e2e4f6d6567612e636f6d626f5f73617427}} with \(\operatorname{Lean\mathord{:}Omega\mathord{:}Constraint\mathord{:}mk}(\operatorname{some}(0),\allowbreak{} \operatorname{some}(0))\) in place of its parameter \(s\) and \(\operatorname{Lean\mathord{:}Omega\mathord{:}Constraint\mathord{:}mk}(\operatorname{some}(0),\allowbreak{} \operatorname{none})\) in place of its parameter \(t\) and \(-1\) in place of its parameter \(a\) and \(x = \langle 1\rangle \mathbin{\Vert} \left(\langle 0\rangle \mathbin{\Vert} \langle -1\rangle\right)\) and \(b = 1\) and \(y = \langle 1\rangle \mathbin{\Vert} \langle -1\rangle\) and \(v = \operatorname{Lean\mathord{:}Omega\mathord{:}Coeffs\mathord{:}of\,List}(\langle |t|\rangle \mathbin{\Vert} (\langle |s|\rangle \mathbin{\Vert} \langle |f[s]|\rangle))\), we obtain \(\operatorname{sat^{\prime}}(\operatorname{Lean\mathord{:}Omega\mathord{:}Constraint\mathord{:}combo}(-1,\allowbreak{} \operatorname{Lean\mathord{:}Omega\mathord{:}Constraint\mathord{:}mk}(\operatorname{some}(0),\allowbreak{} \operatorname{some}(0)),\allowbreak{} 1,\allowbreak{} \operatorname{Lean\mathord{:}Omega\mathord{:}Constraint\mathord{:}mk}(\operatorname{some}(0),\allowbreak{} \operatorname{none})),\allowbreak{} \operatorname{Lean\mathord{:}Omega\mathord{:}Coeffs\mathord{:}combo}(-1,\allowbreak{} \langle 1\rangle \mathbin{\Vert} (\langle 0\rangle \mathbin{\Vert} \langle -1\rangle),\allowbreak{} 1,\allowbreak{} \langle 1\rangle \mathbin{\Vert} \langle -1\rangle),\allowbreak{} \operatorname{Lean\mathord{:}Omega\mathord{:}Coeffs\mathord{:}of\,List}(\langle |t|\rangle \mathbin{\Vert} (\langle |s|\rangle \mathbin{\Vert} \langle |f[s]|\rangle))) = \mathrm{true}\).
By \paperlink{premise:4c65616e2e4f6d6567612e746964795f736174}{Prerequisite~P\ref{premise:4c65616e2e4f6d6567612e746964795f736174}}, we obtain \(\operatorname{sat^{\prime}}(\operatorname{tidy\,Constraint}(\operatorname{Lean\mathord{:}Omega\mathord{:}Constraint\mathord{:}mk}(\operatorname{some}(0),\allowbreak{} \operatorname{none}),\allowbreak{} \langle 0\rangle \mathbin{\Vert} (\langle -1\rangle \mathbin{\Vert} \langle 1\rangle)),\allowbreak{} \operatorname{tidy\,Coeffs}(\operatorname{Lean\mathord{:}Omega\mathord{:}Constraint\mathord{:}mk}(\operatorname{some}(0),\allowbreak{} \operatorname{none}),\allowbreak{} \langle 0\rangle \mathbin{\Vert} (\langle -1\rangle \mathbin{\Vert} \langle 1\rangle)),\allowbreak{} \operatorname{Lean\mathord{:}Omega\mathord{:}Coeffs\mathord{:}of\,List}(\langle |t|\rangle \mathbin{\Vert} (\langle |s|\rangle \mathbin{\Vert} \langle |f[s]|\rangle))) = \mathrm{true}\).
By \paperlink{premise:4c65616e2e4f6d6567612e436f6e73747261696e742e636f6d62696e655f73617427}{Prerequisite~P\ref{premise:4c65616e2e4f6d6567612e436f6e73747261696e742e636f6d62696e655f73617427}}, we obtain \(\operatorname{sat^{\prime}}(\operatorname{combine}(\operatorname{Lean\mathord{:}Omega\mathord{:}Constraint\mathord{:}mk}(\operatorname{some}(0),\allowbreak{} \operatorname{none}),\allowbreak{} \operatorname{Lean\mathord{:}Omega\mathord{:}Constraint\mathord{:}mk}(\operatorname{none},\allowbreak{} \operatorname{some}(0))),\allowbreak{} \langle 0\rangle \mathbin{\Vert} (\langle 1\rangle \mathbin{\Vert} \langle -1\rangle),\allowbreak{} \operatorname{Lean\mathord{:}Omega\mathord{:}Coeffs\mathord{:}of\,List}(\langle |t|\rangle \mathbin{\Vert} (\langle |s|\rangle \mathbin{\Vert} \langle |f[s]|\rangle))) = \mathrm{true}\).
Use the hypothesis \(\left|f\left[s\right]\right| < \left|s\right|\).
By \paperlink{premise:4c65616e2e4f6d6567612e496e742e6f664e61745f6c745f6f665f6c74}{Prerequisite~P\ref{premise:4c65616e2e4f6d6567612e496e742e6f664e61745f6c745f6f665f6c74}}, we obtain \(\left|f\left[s\right]\right| < \left|s\right|\).
By \paperlink{premise:496e742e6164645f6f6e655f6c655f6f665f6c74}{Prerequisite~P\ref{premise:496e742e6164645f6f6e655f6c655f6f665f6c74}}, we obtain \(\left|f\left[s\right]\right| + 1 \le \left|s\right|\).
By \paperlink{premise:496e742e7375625f6e6f6e6e65675f6f665f6c65}{Prerequisite~P\ref{premise:496e742e7375625f6e6f6e6e65675f6f665f6c65}}, we obtain \(0 \le \left|s\right| - \left(\left|f\left[s\right]\right| + 1\right)\).
Using \paperlink{premise:4c65616e2e4f6d6567612e4c696e656172436f6d626f2e636f6f7264696e6174655f6576616c5f31}{Prerequisite~P\ref{premise:4c65616e2e4f6d6567612e4c696e656172436f6d626f2e636f6f7264696e6174655f6576616c5f31}} in reverse gives \(|s| = \operatorname{eval}(\operatorname{coordinate}(1),\allowbreak{} \operatorname{Lean\mathord{:}Omega\mathord{:}Coeffs\mathord{:}of\,List}(\langle |t|\rangle \mathbin{\Vert} (\langle |s|\rangle \mathbin{\Vert} \langle |f[s]|\rangle)))\).
Using \paperlink{premise:4c65616e2e4f6d6567612e4c696e656172436f6d626f2e636f6f7264696e6174655f6576616c5f32}{Prerequisite~P\ref{premise:4c65616e2e4f6d6567612e4c696e656172436f6d626f2e636f6f7264696e6174655f6576616c5f32}} in reverse gives \(|f[s]| = \operatorname{eval}(\operatorname{coordinate}(2),\allowbreak{} \operatorname{Lean\mathord{:}Omega\mathord{:}Coeffs\mathord{:}of\,List}(\langle |t|\rangle \mathbin{\Vert} (\langle |s|\rangle \mathbin{\Vert} \langle |f[s]|\rangle)))\).
By definition, we obtain \(1 = \operatorname{eval}(\operatorname{Lean\mathord{:}Omega\mathord{:}Linear\,Combo\mathord{:}mk}(1,\allowbreak{} \langle\,\rangle),\allowbreak{} \operatorname{Lean\mathord{:}Omega\mathord{:}Coeffs\mathord{:}of\,List}(\langle |t|\rangle \mathbin{\Vert} (\langle |s|\rangle \mathbin{\Vert} \langle |f[s]|\rangle)))\).
By \paperlink{premise:4c65616e2e4f6d6567612e496e742e6164645f636f6e6772}{Prerequisite~P\ref{premise:4c65616e2e4f6d6567612e496e742e6164645f636f6e6772}}, we obtain \(|f[s]| + 1 = \operatorname{eval}(\operatorname{coordinate}(2),\allowbreak{} \operatorname{Lean\mathord{:}Omega\mathord{:}Coeffs\mathord{:}of\,List}(\langle |t|\rangle \mathbin{\Vert} (\langle |s|\rangle \mathbin{\Vert} \langle |f[s]|\rangle))) + \operatorname{eval}(\operatorname{Lean\mathord{:}Omega\mathord{:}Linear\,Combo\mathord{:}mk}(1,\allowbreak{} \langle\,\rangle),\allowbreak{} \operatorname{Lean\mathord{:}Omega\mathord{:}Coeffs\mathord{:}of\,List}(\langle |t|\rangle \mathbin{\Vert} (\langle |s|\rangle \mathbin{\Vert} \langle |f[s]|\rangle)))\).
Using \paperlink{premise:4c65616e2e4f6d6567612e4c696e656172436f6d626f2e6164645f6576616c}{Prerequisite~P\ref{premise:4c65616e2e4f6d6567612e4c696e656172436f6d626f2e6164645f6576616c}} with \(l_{1} = \operatorname{Lean\mathord{:}Omega\mathord{:}Linear\,Combo\mathord{:}mk}(0,\allowbreak{} \langle 0\rangle \mathbin{\Vert} (\langle 0\rangle \mathbin{\Vert} \langle 1\rangle))\) and \(l_{2} = \operatorname{Lean\mathord{:}Omega\mathord{:}Linear\,Combo\mathord{:}mk}(1,\allowbreak{} \langle\,\rangle)\) and \(v = \operatorname{Lean\mathord{:}Omega\mathord{:}Coeffs\mathord{:}of\,List}(\langle |t|\rangle \mathbin{\Vert} (\langle |s|\rangle \mathbin{\Vert} \langle |f[s]|\rangle))\) in reverse gives \(\operatorname{eval}(\operatorname{Lean\mathord{:}Omega\mathord{:}Linear\,Combo\mathord{:}mk}(0,\allowbreak{} \langle 0\rangle \mathbin{\Vert} (\langle 0\rangle \mathbin{\Vert} \langle 1\rangle)),\allowbreak{} \operatorname{Lean\mathord{:}Omega\mathord{:}Coeffs\mathord{:}of\,List}(\langle |t|\rangle \mathbin{\Vert} (\langle |s|\rangle \mathbin{\Vert} \langle |f[s]|\rangle))) + \operatorname{eval}(\operatorname{Lean\mathord{:}Omega\mathord{:}Linear\,Combo\mathord{:}mk}(1,\allowbreak{} \langle\,\rangle),\allowbreak{} \operatorname{Lean\mathord{:}Omega\mathord{:}Coeffs\mathord{:}of\,List}(\langle |t|\rangle \mathbin{\Vert} (\langle |s|\rangle \mathbin{\Vert} \langle |f[s]|\rangle))) = \operatorname{eval}(\operatorname{Lean\mathord{:}Omega\mathord{:}Linear\,Combo\mathord{:}mk}(0,\allowbreak{} \langle 0\rangle \mathbin{\Vert} (\langle 0\rangle \mathbin{\Vert} \langle 1\rangle)) + \operatorname{Lean\mathord{:}Omega\mathord{:}Linear\,Combo\mathord{:}mk}(1,\allowbreak{} \langle\,\rangle),\allowbreak{} \operatorname{Lean\mathord{:}Omega\mathord{:}Coeffs\mathord{:}of\,List}(\langle |t|\rangle \mathbin{\Vert} (\langle |s|\rangle \mathbin{\Vert} \langle |f[s]|\rangle)))\).
Chaining the equalities, we obtain \(|f[s]| + 1 = \operatorname{eval}(\operatorname{Lean\mathord{:}Omega\mathord{:}Linear\,Combo\mathord{:}mk}(0,\allowbreak{} \langle 0\rangle \mathbin{\Vert} (\langle 0\rangle \mathbin{\Vert} \langle 1\rangle)) + \operatorname{Lean\mathord{:}Omega\mathord{:}Linear\,Combo\mathord{:}mk}(1,\allowbreak{} \langle\,\rangle),\allowbreak{} \operatorname{Lean\mathord{:}Omega\mathord{:}Coeffs\mathord{:}of\,List}(\langle |t|\rangle \mathbin{\Vert} (\langle |s|\rangle \mathbin{\Vert} \langle |f[s]|\rangle)))\).
By \paperlink{premise:4c65616e2e4f6d6567612e496e742e7375625f636f6e6772}{Prerequisite~P\ref{premise:4c65616e2e4f6d6567612e496e742e7375625f636f6e6772}}, we obtain \(|s| - (|f[s]| + 1) = \operatorname{eval}(\operatorname{coordinate}(1),\allowbreak{} \operatorname{Lean\mathord{:}Omega\mathord{:}Coeffs\mathord{:}of\,List}(\langle |t|\rangle \mathbin{\Vert} (\langle |s|\rangle \mathbin{\Vert} \langle |f[s]|\rangle))) - \operatorname{eval}(\operatorname{Lean\mathord{:}Omega\mathord{:}Linear\,Combo\mathord{:}mk}(0,\allowbreak{} \langle 0\rangle \mathbin{\Vert} (\langle 0\rangle \mathbin{\Vert} \langle 1\rangle)) + \operatorname{Lean\mathord{:}Omega\mathord{:}Linear\,Combo\mathord{:}mk}(1,\allowbreak{} \langle\,\rangle),\allowbreak{} \operatorname{Lean\mathord{:}Omega\mathord{:}Coeffs\mathord{:}of\,List}(\langle |t|\rangle \mathbin{\Vert} (\langle |s|\rangle \mathbin{\Vert} \langle |f[s]|\rangle)))\).
Using \paperlink{premise:4c65616e2e4f6d6567612e4c696e656172436f6d626f2e7375625f6576616c}{Prerequisite~P\ref{premise:4c65616e2e4f6d6567612e4c696e656172436f6d626f2e7375625f6576616c}} with \(l_{1} = \operatorname{Lean\mathord{:}Omega\mathord{:}Linear\,Combo\mathord{:}mk}(0,\allowbreak{} \langle 0\rangle \mathbin{\Vert} \langle 1\rangle)\) and \(l_{2} = \operatorname{Lean\mathord{:}Omega\mathord{:}Linear\,Combo\mathord{:}mk}(1,\allowbreak{} \langle 0\rangle \mathbin{\Vert} (\langle 0\rangle \mathbin{\Vert} \langle 1\rangle))\) and \(v = \operatorname{Lean\mathord{:}Omega\mathord{:}Coeffs\mathord{:}of\,List}(\langle |t|\rangle \mathbin{\Vert} (\langle |s|\rangle \mathbin{\Vert} \langle |f[s]|\rangle))\) in reverse gives \(\operatorname{eval}(\operatorname{Lean\mathord{:}Omega\mathord{:}Linear\,Combo\mathord{:}mk}(0,\allowbreak{} \langle 0\rangle \mathbin{\Vert} \langle 1\rangle),\allowbreak{} \operatorname{Lean\mathord{:}Omega\mathord{:}Coeffs\mathord{:}of\,List}(\langle |t|\rangle \mathbin{\Vert} (\langle |s|\rangle \mathbin{\Vert} \langle |f[s]|\rangle))) - \operatorname{eval}(\operatorname{Lean\mathord{:}Omega\mathord{:}Linear\,Combo\mathord{:}mk}(1,\allowbreak{} \langle 0\rangle \mathbin{\Vert} (\langle 0\rangle \mathbin{\Vert} \langle 1\rangle)),\allowbreak{} \operatorname{Lean\mathord{:}Omega\mathord{:}Coeffs\mathord{:}of\,List}(\langle |t|\rangle \mathbin{\Vert} (\langle |s|\rangle \mathbin{\Vert} \langle |f[s]|\rangle))) = \operatorname{eval}(\operatorname{Lean\mathord{:}Omega\mathord{:}Linear\,Combo\mathord{:}mk}(0,\allowbreak{} \langle 0\rangle \mathbin{\Vert} \langle 1\rangle) - \operatorname{Lean\mathord{:}Omega\mathord{:}Linear\,Combo\mathord{:}mk}(1,\allowbreak{} \langle 0\rangle \mathbin{\Vert} (\langle 0\rangle \mathbin{\Vert} \langle 1\rangle)),\allowbreak{} \operatorname{Lean\mathord{:}Omega\mathord{:}Coeffs\mathord{:}of\,List}(\langle |t|\rangle \mathbin{\Vert} (\langle |s|\rangle \mathbin{\Vert} \langle |f[s]|\rangle)))\).
Chaining the equalities, we obtain \(|s| - (|f[s]| + 1) = \operatorname{eval}(\operatorname{Lean\mathord{:}Omega\mathord{:}Linear\,Combo\mathord{:}mk}(0,\allowbreak{} \langle 0\rangle \mathbin{\Vert} \langle 1\rangle) - \operatorname{Lean\mathord{:}Omega\mathord{:}Linear\,Combo\mathord{:}mk}(1,\allowbreak{} \langle 0\rangle \mathbin{\Vert} (\langle 0\rangle \mathbin{\Vert} \langle 1\rangle)),\allowbreak{} \operatorname{Lean\mathord{:}Omega\mathord{:}Coeffs\mathord{:}of\,List}(\langle |t|\rangle \mathbin{\Vert} (\langle |s|\rangle \mathbin{\Vert} \langle |f[s]|\rangle)))\).
By \paperlink{premise:6c655f6f665f6c655f6f665f6571}{Prerequisite~P\ref{premise:6c655f6f665f6c655f6f665f6571}} with \(\mathbb{Z}\) as its type parameter \(\alpha\), we obtain \(0 \le \operatorname{eval}(\operatorname{Lean\mathord{:}Omega\mathord{:}Linear\,Combo\mathord{:}mk}(0,\allowbreak{} \langle 0\rangle \mathbin{\Vert} \langle 1\rangle) - \operatorname{Lean\mathord{:}Omega\mathord{:}Linear\,Combo\mathord{:}mk}(1,\allowbreak{} \langle 0\rangle \mathbin{\Vert} (\langle 0\rangle \mathbin{\Vert} \langle 1\rangle)),\allowbreak{} \operatorname{Lean\mathord{:}Omega\mathord{:}Coeffs\mathord{:}of\,List}(\langle |t|\rangle \mathbin{\Vert} (\langle |s|\rangle \mathbin{\Vert} \langle |f[s]|\rangle)))\).
By \paperlink{premise:4c65616e2e4f6d6567612e436f6e73747261696e742e616464496e657175616c6974795f736174}{Prerequisite~P\ref{premise:4c65616e2e4f6d6567612e436f6e73747261696e742e616464496e657175616c6974795f736174}}, we obtain \(\operatorname{sat^{\prime}}(\operatorname{Lean\mathord{:}Omega\mathord{:}Constraint\mathord{:}mk}(\operatorname{some}(-(-1)),\allowbreak{} \operatorname{none}),\allowbreak{} \langle 0\rangle \mathbin{\Vert} (\langle 1\rangle \mathbin{\Vert} \langle -1\rangle),\allowbreak{} \operatorname{Lean\mathord{:}Omega\mathord{:}Coeffs\mathord{:}of\,List}(\langle |t|\rangle \mathbin{\Vert} (\langle |s|\rangle \mathbin{\Vert} \langle |f[s]|\rangle))) = \mathrm{true}\).
By \paperlink{premise:4c65616e2e4f6d6567612e636f6d626f5f73617427}{Prerequisite~P\ref{premise:4c65616e2e4f6d6567612e636f6d626f5f73617427}} with \(\operatorname{Lean\mathord{:}Omega\mathord{:}Constraint\mathord{:}mk}(\operatorname{some}(0),\allowbreak{} \operatorname{some}(0))\) in place of its parameter \(s\) and \(\operatorname{Lean\mathord{:}Omega\mathord{:}Constraint\mathord{:}mk}(\operatorname{some}(1),\allowbreak{} \operatorname{none})\) in place of its parameter \(t\) and \(-1\) in place of its parameter \(a\) and \(x = \langle 0\rangle \mathbin{\Vert} \left(\langle 1\rangle \mathbin{\Vert} \langle -1\rangle\right)\) and \(b = 1\) and \(y = \langle 0\rangle \mathbin{\Vert} \left(\langle 1\rangle \mathbin{\Vert} \langle -1\rangle\right)\) and \(v = \operatorname{Lean\mathord{:}Omega\mathord{:}Coeffs\mathord{:}of\,List}(\langle |t|\rangle \mathbin{\Vert} (\langle |s|\rangle \mathbin{\Vert} \langle |f[s]|\rangle))\), we obtain \(\operatorname{sat^{\prime}}(\operatorname{Lean\mathord{:}Omega\mathord{:}Constraint\mathord{:}combo}(-1,\allowbreak{} \operatorname{Lean\mathord{:}Omega\mathord{:}Constraint\mathord{:}mk}(\operatorname{some}(0),\allowbreak{} \operatorname{some}(0)),\allowbreak{} 1,\allowbreak{} \operatorname{Lean\mathord{:}Omega\mathord{:}Constraint\mathord{:}mk}(\operatorname{some}(1),\allowbreak{} \operatorname{none})),\allowbreak{} \operatorname{Lean\mathord{:}Omega\mathord{:}Coeffs\mathord{:}combo}(-1,\allowbreak{} \langle 0\rangle \mathbin{\Vert} (\langle 1\rangle \mathbin{\Vert} \langle -1\rangle),\allowbreak{} 1,\allowbreak{} \langle 0\rangle \mathbin{\Vert} (\langle 1\rangle \mathbin{\Vert} \langle -1\rangle)),\allowbreak{} \operatorname{Lean\mathord{:}Omega\mathord{:}Coeffs\mathord{:}of\,List}(\langle |t|\rangle \mathbin{\Vert} (\langle |s|\rangle \mathbin{\Vert} \langle |f[s]|\rangle))) = \mathrm{true}\).
By \paperlink{premise:4c65616e2e4f6d6567612e746964795f736174}{Prerequisite~P\ref{premise:4c65616e2e4f6d6567612e746964795f736174}}, we obtain \(\operatorname{sat^{\prime}}(\operatorname{tidy\,Constraint}(\operatorname{Lean\mathord{:}Omega\mathord{:}Constraint\mathord{:}mk}(\operatorname{some}(1),\allowbreak{} \operatorname{none}),\allowbreak{} \langle 0\rangle \mathbin{\Vert} (\langle 0\rangle \mathbin{\Vert} \langle 0\rangle)),\allowbreak{} \operatorname{tidy\,Coeffs}(\operatorname{Lean\mathord{:}Omega\mathord{:}Constraint\mathord{:}mk}(\operatorname{some}(1),\allowbreak{} \operatorname{none}),\allowbreak{} \langle 0\rangle \mathbin{\Vert} (\langle 0\rangle \mathbin{\Vert} \langle 0\rangle)),\allowbreak{} \operatorname{Lean\mathord{:}Omega\mathord{:}Coeffs\mathord{:}of\,List}(\langle |t|\rangle \mathbin{\Vert} (\langle |s|\rangle \mathbin{\Vert} \langle |f[s]|\rangle))) = \mathrm{true}\).
By \paperlink{premise:4c65616e2e4f6d6567612e436f6e73747261696e742e6e6f745f736174275f6f665f6973496d706f737369626c65}{Prerequisite~P\ref{premise:4c65616e2e4f6d6567612e436f6e73747261696e742e6e6f745f736174275f6f665f6973496d706f737369626c65}}, we obtain \(\mathrm{False}\).
\end{proofstep}
Consider the following case.
\begin{proofstep}
Assume \(\left|f\left[s\right]\right| > \left|s\right|\).
By definition, we obtain \(\operatorname{decide}(\operatorname{is\,Impossible}(\operatorname{Lean\mathord{:}Omega\mathord{:}Constraint\mathord{:}mk}(\operatorname{some}(1),\allowbreak{} \operatorname{some}(0))) = \mathrm{true}) = \mathrm{true}\).
By \paperlink{premise:6f665f6465636964655f65715f74727565}{Prerequisite~P\ref{premise:6f665f6465636964655f65715f74727565}}, we obtain \(\operatorname{is\,Impossible}(\operatorname{Lean\mathord{:}Omega\mathord{:}Constraint\mathord{:}mk}(\operatorname{some}(1),\allowbreak{} \operatorname{some}(0))) = \mathrm{true}\).
Use the previously proved fact \(\left|f\left[s\right]\right| \le \left|s\right|\).
By \paperlink{premise:4c65616e2e4f6d6567612e496e742e6f664e61745f6c655f6f665f6c65}{Prerequisite~P\ref{premise:4c65616e2e4f6d6567612e496e742e6f664e61745f6c655f6f665f6c65}}, we obtain \(\left|f\left[s\right]\right| \le \left|s\right|\).
By \paperlink{premise:496e742e7375625f6e6f6e6e65675f6f665f6c65}{Prerequisite~P\ref{premise:496e742e7375625f6e6f6e6e65675f6f665f6c65}}, we obtain \(0 \le \left|s\right| - \left|f\left[s\right]\right|\).
Using \paperlink{premise:4c65616e2e4f6d6567612e4c696e656172436f6d626f2e636f6f7264696e6174655f6576616c5f31}{Prerequisite~P\ref{premise:4c65616e2e4f6d6567612e4c696e656172436f6d626f2e636f6f7264696e6174655f6576616c5f31}} in reverse gives \(|s| = \operatorname{eval}(\operatorname{coordinate}(1),\allowbreak{} \operatorname{Lean\mathord{:}Omega\mathord{:}Coeffs\mathord{:}of\,List}(\langle |t|\rangle \mathbin{\Vert} (\langle |s|\rangle \mathbin{\Vert} \langle |f[s]|\rangle)))\).
Using \paperlink{premise:4c65616e2e4f6d6567612e4c696e656172436f6d626f2e636f6f7264696e6174655f6576616c5f32}{Prerequisite~P\ref{premise:4c65616e2e4f6d6567612e4c696e656172436f6d626f2e636f6f7264696e6174655f6576616c5f32}} in reverse gives \(|f[s]| = \operatorname{eval}(\operatorname{coordinate}(2),\allowbreak{} \operatorname{Lean\mathord{:}Omega\mathord{:}Coeffs\mathord{:}of\,List}(\langle |t|\rangle \mathbin{\Vert} (\langle |s|\rangle \mathbin{\Vert} \langle |f[s]|\rangle)))\).
By \paperlink{premise:4c65616e2e4f6d6567612e496e742e7375625f636f6e6772}{Prerequisite~P\ref{premise:4c65616e2e4f6d6567612e496e742e7375625f636f6e6772}}, we obtain \(|s| - |f[s]| = \operatorname{eval}(\operatorname{coordinate}(1),\allowbreak{} \operatorname{Lean\mathord{:}Omega\mathord{:}Coeffs\mathord{:}of\,List}(\langle |t|\rangle \mathbin{\Vert} (\langle |s|\rangle \mathbin{\Vert} \langle |f[s]|\rangle))) - \operatorname{eval}(\operatorname{coordinate}(2),\allowbreak{} \operatorname{Lean\mathord{:}Omega\mathord{:}Coeffs\mathord{:}of\,List}(\langle |t|\rangle \mathbin{\Vert} (\langle |s|\rangle \mathbin{\Vert} \langle |f[s]|\rangle)))\).
Using \paperlink{premise:4c65616e2e4f6d6567612e4c696e656172436f6d626f2e7375625f6576616c}{Prerequisite~P\ref{premise:4c65616e2e4f6d6567612e4c696e656172436f6d626f2e7375625f6576616c}} with \(l_{1} = \operatorname{Lean\mathord{:}Omega\mathord{:}Linear\,Combo\mathord{:}mk}(0,\allowbreak{} \langle 0\rangle \mathbin{\Vert} \langle 1\rangle)\) and \(l_{2} = \operatorname{Lean\mathord{:}Omega\mathord{:}Linear\,Combo\mathord{:}mk}(0,\allowbreak{} \langle 0\rangle \mathbin{\Vert} (\langle 0\rangle \mathbin{\Vert} \langle 1\rangle))\) and \(v = \operatorname{Lean\mathord{:}Omega\mathord{:}Coeffs\mathord{:}of\,List}(\langle |t|\rangle \mathbin{\Vert} (\langle |s|\rangle \mathbin{\Vert} \langle |f[s]|\rangle))\) in reverse gives \(\operatorname{eval}(\operatorname{Lean\mathord{:}Omega\mathord{:}Linear\,Combo\mathord{:}mk}(0,\allowbreak{} \langle 0\rangle \mathbin{\Vert} \langle 1\rangle),\allowbreak{} \operatorname{Lean\mathord{:}Omega\mathord{:}Coeffs\mathord{:}of\,List}(\langle |t|\rangle \mathbin{\Vert} (\langle |s|\rangle \mathbin{\Vert} \langle |f[s]|\rangle))) - \operatorname{eval}(\operatorname{Lean\mathord{:}Omega\mathord{:}Linear\,Combo\mathord{:}mk}(0,\allowbreak{} \langle 0\rangle \mathbin{\Vert} (\langle 0\rangle \mathbin{\Vert} \langle 1\rangle)),\allowbreak{} \operatorname{Lean\mathord{:}Omega\mathord{:}Coeffs\mathord{:}of\,List}(\langle |t|\rangle \mathbin{\Vert} (\langle |s|\rangle \mathbin{\Vert} \langle |f[s]|\rangle))) = \operatorname{eval}(\operatorname{Lean\mathord{:}Omega\mathord{:}Linear\,Combo\mathord{:}mk}(0,\allowbreak{} \langle 0\rangle \mathbin{\Vert} \langle 1\rangle) - \operatorname{Lean\mathord{:}Omega\mathord{:}Linear\,Combo\mathord{:}mk}(0,\allowbreak{} \langle 0\rangle \mathbin{\Vert} (\langle 0\rangle \mathbin{\Vert} \langle 1\rangle)),\allowbreak{} \operatorname{Lean\mathord{:}Omega\mathord{:}Coeffs\mathord{:}of\,List}(\langle |t|\rangle \mathbin{\Vert} (\langle |s|\rangle \mathbin{\Vert} \langle |f[s]|\rangle)))\).
Chaining the equalities, we obtain \(|s| - |f[s]| = \operatorname{eval}(\operatorname{Lean\mathord{:}Omega\mathord{:}Linear\,Combo\mathord{:}mk}(0,\allowbreak{} \langle 0\rangle \mathbin{\Vert} \langle 1\rangle) - \operatorname{Lean\mathord{:}Omega\mathord{:}Linear\,Combo\mathord{:}mk}(0,\allowbreak{} \langle 0\rangle \mathbin{\Vert} (\langle 0\rangle \mathbin{\Vert} \langle 1\rangle)),\allowbreak{} \operatorname{Lean\mathord{:}Omega\mathord{:}Coeffs\mathord{:}of\,List}(\langle |t|\rangle \mathbin{\Vert} (\langle |s|\rangle \mathbin{\Vert} \langle |f[s]|\rangle)))\).
By \paperlink{premise:6c655f6f665f6c655f6f665f6571}{Prerequisite~P\ref{premise:6c655f6f665f6c655f6f665f6571}} with \(\mathbb{Z}\) as its type parameter \(\alpha\), we obtain \(0 \le \operatorname{eval}(\operatorname{Lean\mathord{:}Omega\mathord{:}Linear\,Combo\mathord{:}mk}(0,\allowbreak{} \langle 0\rangle \mathbin{\Vert} \langle 1\rangle) - \operatorname{Lean\mathord{:}Omega\mathord{:}Linear\,Combo\mathord{:}mk}(0,\allowbreak{} \langle 0\rangle \mathbin{\Vert} (\langle 0\rangle \mathbin{\Vert} \langle 1\rangle)),\allowbreak{} \operatorname{Lean\mathord{:}Omega\mathord{:}Coeffs\mathord{:}of\,List}(\langle |t|\rangle \mathbin{\Vert} (\langle |s|\rangle \mathbin{\Vert} \langle |f[s]|\rangle)))\).
By \paperlink{premise:4c65616e2e4f6d6567612e436f6e73747261696e742e616464496e657175616c6974795f736174}{Prerequisite~P\ref{premise:4c65616e2e4f6d6567612e436f6e73747261696e742e616464496e657175616c6974795f736174}}, we obtain \(\operatorname{sat^{\prime}}(\operatorname{Lean\mathord{:}Omega\mathord{:}Constraint\mathord{:}mk}(\operatorname{some}(-0),\allowbreak{} \operatorname{none}),\allowbreak{} \langle 0\rangle \mathbin{\Vert} (\langle 1\rangle \mathbin{\Vert} \langle -1\rangle),\allowbreak{} \operatorname{Lean\mathord{:}Omega\mathord{:}Coeffs\mathord{:}of\,List}(\langle |t|\rangle \mathbin{\Vert} (\langle |s|\rangle \mathbin{\Vert} \langle |f[s]|\rangle))) = \mathrm{true}\).
Using \paperlink{premise:4c65616e2e4f6d6567612e4c696e656172436f6d626f2e636f6f7264696e6174655f6576616c5f32}{Prerequisite~P\ref{premise:4c65616e2e4f6d6567612e4c696e656172436f6d626f2e636f6f7264696e6174655f6576616c5f32}} in reverse gives \(|f[s]| = \operatorname{eval}(\operatorname{coordinate}(2),\allowbreak{} \operatorname{Lean\mathord{:}Omega\mathord{:}Coeffs\mathord{:}of\,List}(\langle |t|\rangle \mathbin{\Vert} (\langle |s|\rangle \mathbin{\Vert} \langle |f[s]|\rangle)))\).
Using \paperlink{premise:4c65616e2e4f6d6567612e4c696e656172436f6d626f2e636f6f7264696e6174655f6576616c5f30}{Prerequisite~P\ref{premise:4c65616e2e4f6d6567612e4c696e656172436f6d626f2e636f6f7264696e6174655f6576616c5f30}} in reverse gives \(|t| = \operatorname{eval}(\operatorname{coordinate}(0),\allowbreak{} \operatorname{Lean\mathord{:}Omega\mathord{:}Coeffs\mathord{:}of\,List}(\langle |t|\rangle \mathbin{\Vert} (\langle |s|\rangle \mathbin{\Vert} \langle |f[s]|\rangle)))\).
By \paperlink{premise:4c65616e2e4f6d6567612e496e742e7375625f636f6e6772}{Prerequisite~P\ref{premise:4c65616e2e4f6d6567612e496e742e7375625f636f6e6772}}, we obtain \(|f[s]| - |t| = \operatorname{eval}(\operatorname{coordinate}(2),\allowbreak{} \operatorname{Lean\mathord{:}Omega\mathord{:}Coeffs\mathord{:}of\,List}(\langle |t|\rangle \mathbin{\Vert} (\langle |s|\rangle \mathbin{\Vert} \langle |f[s]|\rangle))) - \operatorname{eval}(\operatorname{coordinate}(0),\allowbreak{} \operatorname{Lean\mathord{:}Omega\mathord{:}Coeffs\mathord{:}of\,List}(\langle |t|\rangle \mathbin{\Vert} (\langle |s|\rangle \mathbin{\Vert} \langle |f[s]|\rangle)))\).
Using \paperlink{premise:4c65616e2e4f6d6567612e4c696e656172436f6d626f2e7375625f6576616c}{Prerequisite~P\ref{premise:4c65616e2e4f6d6567612e4c696e656172436f6d626f2e7375625f6576616c}} with \(l_{1} = \operatorname{Lean\mathord{:}Omega\mathord{:}Linear\,Combo\mathord{:}mk}(0,\allowbreak{} \langle 0\rangle \mathbin{\Vert} (\langle 0\rangle \mathbin{\Vert} \langle 1\rangle))\) and \(l_{2} = \operatorname{Lean\mathord{:}Omega\mathord{:}Linear\,Combo\mathord{:}mk}(0,\allowbreak{} \langle 1\rangle)\) and \(v = \operatorname{Lean\mathord{:}Omega\mathord{:}Coeffs\mathord{:}of\,List}(\langle |t|\rangle \mathbin{\Vert} (\langle |s|\rangle \mathbin{\Vert} \langle |f[s]|\rangle))\) in reverse gives \(\operatorname{eval}(\operatorname{Lean\mathord{:}Omega\mathord{:}Linear\,Combo\mathord{:}mk}(0,\allowbreak{} \langle 0\rangle \mathbin{\Vert} (\langle 0\rangle \mathbin{\Vert} \langle 1\rangle)),\allowbreak{} \operatorname{Lean\mathord{:}Omega\mathord{:}Coeffs\mathord{:}of\,List}(\langle |t|\rangle \mathbin{\Vert} (\langle |s|\rangle \mathbin{\Vert} \langle |f[s]|\rangle))) - \operatorname{eval}(\operatorname{Lean\mathord{:}Omega\mathord{:}Linear\,Combo\mathord{:}mk}(0,\allowbreak{} \langle 1\rangle),\allowbreak{} \operatorname{Lean\mathord{:}Omega\mathord{:}Coeffs\mathord{:}of\,List}(\langle |t|\rangle \mathbin{\Vert} (\langle |s|\rangle \mathbin{\Vert} \langle |f[s]|\rangle))) = \operatorname{eval}(\operatorname{Lean\mathord{:}Omega\mathord{:}Linear\,Combo\mathord{:}mk}(0,\allowbreak{} \langle 0\rangle \mathbin{\Vert} (\langle 0\rangle \mathbin{\Vert} \langle 1\rangle)) - \operatorname{Lean\mathord{:}Omega\mathord{:}Linear\,Combo\mathord{:}mk}(0,\allowbreak{} \langle 1\rangle),\allowbreak{} \operatorname{Lean\mathord{:}Omega\mathord{:}Coeffs\mathord{:}of\,List}(\langle |t|\rangle \mathbin{\Vert} (\langle |s|\rangle \mathbin{\Vert} \langle |f[s]|\rangle)))\).
Chaining the equalities, we obtain \(|f[s]| - |t| = \operatorname{eval}(\operatorname{Lean\mathord{:}Omega\mathord{:}Linear\,Combo\mathord{:}mk}(0,\allowbreak{} \langle 0\rangle \mathbin{\Vert} (\langle 0\rangle \mathbin{\Vert} \langle 1\rangle)) - \operatorname{Lean\mathord{:}Omega\mathord{:}Linear\,Combo\mathord{:}mk}(0,\allowbreak{} \langle 1\rangle),\allowbreak{} \operatorname{Lean\mathord{:}Omega\mathord{:}Coeffs\mathord{:}of\,List}(\langle |t|\rangle \mathbin{\Vert} (\langle |s|\rangle \mathbin{\Vert} \langle |f[s]|\rangle)))\).
Reversing the equality, we obtain \(\operatorname{eval}(\operatorname{Lean\mathord{:}Omega\mathord{:}Linear\,Combo\mathord{:}mk}(0,\allowbreak{} \langle 0\rangle \mathbin{\Vert} (\langle 0\rangle \mathbin{\Vert} \langle 1\rangle)) - \operatorname{Lean\mathord{:}Omega\mathord{:}Linear\,Combo\mathord{:}mk}(0,\allowbreak{} \langle 1\rangle),\allowbreak{} \operatorname{Lean\mathord{:}Omega\mathord{:}Coeffs\mathord{:}of\,List}(\langle |t|\rangle \mathbin{\Vert} (\langle |s|\rangle \mathbin{\Vert} \langle |f[s]|\rangle))) = |f[s]| - |t|\).
Use the previously proved fact \(\left|f\left[s\right]\right| = \left|t\right|\).
By \paperlink{premise:4c65616e2e4f6d6567612e496e742e6f664e61745f636f6e6772}{Prerequisite~P\ref{premise:4c65616e2e4f6d6567612e496e742e6f664e61745f636f6e6772}}, we obtain \(\left|f\left[s\right]\right| = \left|t\right|\).
By \paperlink{premise:496e742e7375625f65715f7a65726f5f6f665f6571}{Prerequisite~P\ref{premise:496e742e7375625f65715f7a65726f5f6f665f6571}}, we obtain \(\left|f\left[s\right]\right| - \left|t\right| = 0\).
Chaining the equalities, we obtain \(\operatorname{eval}(\operatorname{Lean\mathord{:}Omega\mathord{:}Linear\,Combo\mathord{:}mk}(0,\allowbreak{} \langle 0\rangle \mathbin{\Vert} (\langle 0\rangle \mathbin{\Vert} \langle 1\rangle)) - \operatorname{Lean\mathord{:}Omega\mathord{:}Linear\,Combo\mathord{:}mk}(0,\allowbreak{} \langle 1\rangle),\allowbreak{} \operatorname{Lean\mathord{:}Omega\mathord{:}Coeffs\mathord{:}of\,List}(\langle |t|\rangle \mathbin{\Vert} (\langle |s|\rangle \mathbin{\Vert} \langle |f[s]|\rangle))) = 0\).
By \paperlink{premise:4c65616e2e4f6d6567612e436f6e73747261696e742e616464457175616c6974795f736174}{Prerequisite~P\ref{premise:4c65616e2e4f6d6567612e436f6e73747261696e742e616464457175616c6974795f736174}}, we obtain \(\operatorname{sat^{\prime}}(\operatorname{Lean\mathord{:}Omega\mathord{:}Constraint\mathord{:}mk}(\operatorname{some}(-0),\allowbreak{} \operatorname{some}(-0)),\allowbreak{} \langle -1\rangle \mathbin{\Vert} (\langle 0\rangle \mathbin{\Vert} \langle 1\rangle),\allowbreak{} \operatorname{Lean\mathord{:}Omega\mathord{:}Coeffs\mathord{:}of\,List}(\langle |t|\rangle \mathbin{\Vert} (\langle |s|\rangle \mathbin{\Vert} \langle |f[s]|\rangle))) = \mathrm{true}\).
By \paperlink{premise:4c65616e2e4f6d6567612e746964795f736174}{Prerequisite~P\ref{premise:4c65616e2e4f6d6567612e746964795f736174}}, we obtain \(\operatorname{sat^{\prime}}(\operatorname{tidy\,Constraint}(\operatorname{Lean\mathord{:}Omega\mathord{:}Constraint\mathord{:}mk}(\operatorname{some}(0),\allowbreak{} \operatorname{some}(0)),\allowbreak{} \langle -1\rangle \mathbin{\Vert} (\langle 0\rangle \mathbin{\Vert} \langle 1\rangle)),\allowbreak{} \operatorname{tidy\,Coeffs}(\operatorname{Lean\mathord{:}Omega\mathord{:}Constraint\mathord{:}mk}(\operatorname{some}(0),\allowbreak{} \operatorname{some}(0)),\allowbreak{} \langle -1\rangle \mathbin{\Vert} (\langle 0\rangle \mathbin{\Vert} \langle 1\rangle)),\allowbreak{} \operatorname{Lean\mathord{:}Omega\mathord{:}Coeffs\mathord{:}of\,List}(\langle |t|\rangle \mathbin{\Vert} (\langle |s|\rangle \mathbin{\Vert} \langle |f[s]|\rangle))) = \mathrm{true}\).
Use the hypothesis \(\left|s\right| \le \left|t\right|\).
By \paperlink{premise:4c65616e2e4f6d6567612e496e742e6f664e61745f6c655f6f665f6c65}{Prerequisite~P\ref{premise:4c65616e2e4f6d6567612e496e742e6f664e61745f6c655f6f665f6c65}}, we obtain \(\left|s\right| \le \left|t\right|\).
By \paperlink{premise:496e742e7375625f6e6f6e6e65675f6f665f6c65}{Prerequisite~P\ref{premise:496e742e7375625f6e6f6e6e65675f6f665f6c65}}, we obtain \(0 \le \left|t\right| - \left|s\right|\).
Using \paperlink{premise:4c65616e2e4f6d6567612e4c696e656172436f6d626f2e636f6f7264696e6174655f6576616c5f30}{Prerequisite~P\ref{premise:4c65616e2e4f6d6567612e4c696e656172436f6d626f2e636f6f7264696e6174655f6576616c5f30}} in reverse gives \(|t| = \operatorname{eval}(\operatorname{coordinate}(0),\allowbreak{} \operatorname{Lean\mathord{:}Omega\mathord{:}Coeffs\mathord{:}of\,List}(\langle |t|\rangle \mathbin{\Vert} (\langle |s|\rangle \mathbin{\Vert} \langle |f[s]|\rangle)))\).
Using \paperlink{premise:4c65616e2e4f6d6567612e4c696e656172436f6d626f2e636f6f7264696e6174655f6576616c5f31}{Prerequisite~P\ref{premise:4c65616e2e4f6d6567612e4c696e656172436f6d626f2e636f6f7264696e6174655f6576616c5f31}} in reverse gives \(|s| = \operatorname{eval}(\operatorname{coordinate}(1),\allowbreak{} \operatorname{Lean\mathord{:}Omega\mathord{:}Coeffs\mathord{:}of\,List}(\langle |t|\rangle \mathbin{\Vert} (\langle |s|\rangle \mathbin{\Vert} \langle |f[s]|\rangle)))\).
By \paperlink{premise:4c65616e2e4f6d6567612e496e742e7375625f636f6e6772}{Prerequisite~P\ref{premise:4c65616e2e4f6d6567612e496e742e7375625f636f6e6772}}, we obtain \(|t| - |s| = \operatorname{eval}(\operatorname{coordinate}(0),\allowbreak{} \operatorname{Lean\mathord{:}Omega\mathord{:}Coeffs\mathord{:}of\,List}(\langle |t|\rangle \mathbin{\Vert} (\langle |s|\rangle \mathbin{\Vert} \langle |f[s]|\rangle))) - \operatorname{eval}(\operatorname{coordinate}(1),\allowbreak{} \operatorname{Lean\mathord{:}Omega\mathord{:}Coeffs\mathord{:}of\,List}(\langle |t|\rangle \mathbin{\Vert} (\langle |s|\rangle \mathbin{\Vert} \langle |f[s]|\rangle)))\).
Using \paperlink{premise:4c65616e2e4f6d6567612e4c696e656172436f6d626f2e7375625f6576616c}{Prerequisite~P\ref{premise:4c65616e2e4f6d6567612e4c696e656172436f6d626f2e7375625f6576616c}} with \(l_{1} = \operatorname{Lean\mathord{:}Omega\mathord{:}Linear\,Combo\mathord{:}mk}(0,\allowbreak{} \langle 1\rangle)\) and \(l_{2} = \operatorname{Lean\mathord{:}Omega\mathord{:}Linear\,Combo\mathord{:}mk}(0,\allowbreak{} \langle 0\rangle \mathbin{\Vert} \langle 1\rangle)\) and \(v = \operatorname{Lean\mathord{:}Omega\mathord{:}Coeffs\mathord{:}of\,List}(\langle |t|\rangle \mathbin{\Vert} (\langle |s|\rangle \mathbin{\Vert} \langle |f[s]|\rangle))\) in reverse gives \(\operatorname{eval}(\operatorname{Lean\mathord{:}Omega\mathord{:}Linear\,Combo\mathord{:}mk}(0,\allowbreak{} \langle 1\rangle),\allowbreak{} \operatorname{Lean\mathord{:}Omega\mathord{:}Coeffs\mathord{:}of\,List}(\langle |t|\rangle \mathbin{\Vert} (\langle |s|\rangle \mathbin{\Vert} \langle |f[s]|\rangle))) - \operatorname{eval}(\operatorname{Lean\mathord{:}Omega\mathord{:}Linear\,Combo\mathord{:}mk}(0,\allowbreak{} \langle 0\rangle \mathbin{\Vert} \langle 1\rangle),\allowbreak{} \operatorname{Lean\mathord{:}Omega\mathord{:}Coeffs\mathord{:}of\,List}(\langle |t|\rangle \mathbin{\Vert} (\langle |s|\rangle \mathbin{\Vert} \langle |f[s]|\rangle))) = \operatorname{eval}(\operatorname{Lean\mathord{:}Omega\mathord{:}Linear\,Combo\mathord{:}mk}(0,\allowbreak{} \langle 1\rangle) - \operatorname{Lean\mathord{:}Omega\mathord{:}Linear\,Combo\mathord{:}mk}(0,\allowbreak{} \langle 0\rangle \mathbin{\Vert} \langle 1\rangle),\allowbreak{} \operatorname{Lean\mathord{:}Omega\mathord{:}Coeffs\mathord{:}of\,List}(\langle |t|\rangle \mathbin{\Vert} (\langle |s|\rangle \mathbin{\Vert} \langle |f[s]|\rangle)))\).
Chaining the equalities, we obtain \(|t| - |s| = \operatorname{eval}(\operatorname{Lean\mathord{:}Omega\mathord{:}Linear\,Combo\mathord{:}mk}(0,\allowbreak{} \langle 1\rangle) - \operatorname{Lean\mathord{:}Omega\mathord{:}Linear\,Combo\mathord{:}mk}(0,\allowbreak{} \langle 0\rangle \mathbin{\Vert} \langle 1\rangle),\allowbreak{} \operatorname{Lean\mathord{:}Omega\mathord{:}Coeffs\mathord{:}of\,List}(\langle |t|\rangle \mathbin{\Vert} (\langle |s|\rangle \mathbin{\Vert} \langle |f[s]|\rangle)))\).
By \paperlink{premise:6c655f6f665f6c655f6f665f6571}{Prerequisite~P\ref{premise:6c655f6f665f6c655f6f665f6571}} with \(\mathbb{Z}\) as its type parameter \(\alpha\), we obtain \(0 \le \operatorname{eval}(\operatorname{Lean\mathord{:}Omega\mathord{:}Linear\,Combo\mathord{:}mk}(0,\allowbreak{} \langle 1\rangle) - \operatorname{Lean\mathord{:}Omega\mathord{:}Linear\,Combo\mathord{:}mk}(0,\allowbreak{} \langle 0\rangle \mathbin{\Vert} \langle 1\rangle),\allowbreak{} \operatorname{Lean\mathord{:}Omega\mathord{:}Coeffs\mathord{:}of\,List}(\langle |t|\rangle \mathbin{\Vert} (\langle |s|\rangle \mathbin{\Vert} \langle |f[s]|\rangle)))\).
By \paperlink{premise:4c65616e2e4f6d6567612e436f6e73747261696e742e616464496e657175616c6974795f736174}{Prerequisite~P\ref{premise:4c65616e2e4f6d6567612e436f6e73747261696e742e616464496e657175616c6974795f736174}}, we obtain \(\operatorname{sat^{\prime}}(\operatorname{Lean\mathord{:}Omega\mathord{:}Constraint\mathord{:}mk}(\operatorname{some}(-0),\allowbreak{} \operatorname{none}),\allowbreak{} \langle 1\rangle \mathbin{\Vert} \langle -1\rangle,\allowbreak{} \operatorname{Lean\mathord{:}Omega\mathord{:}Coeffs\mathord{:}of\,List}(\langle |t|\rangle \mathbin{\Vert} (\langle |s|\rangle \mathbin{\Vert} \langle |f[s]|\rangle))) = \mathrm{true}\).
By \paperlink{premise:4c65616e2e4f6d6567612e636f6d626f5f73617427}{Prerequisite~P\ref{premise:4c65616e2e4f6d6567612e636f6d626f5f73617427}} with \(\operatorname{Lean\mathord{:}Omega\mathord{:}Constraint\mathord{:}mk}(\operatorname{some}(0),\allowbreak{} \operatorname{some}(0))\) in place of its parameter \(s\) and \(\operatorname{Lean\mathord{:}Omega\mathord{:}Constraint\mathord{:}mk}(\operatorname{some}(0),\allowbreak{} \operatorname{none})\) in place of its parameter \(t\) and \(-1\) in place of its parameter \(a\) and \(x = \langle 1\rangle \mathbin{\Vert} \left(\langle 0\rangle \mathbin{\Vert} \langle -1\rangle\right)\) and \(b = 1\) and \(y = \langle 1\rangle \mathbin{\Vert} \langle -1\rangle\) and \(v = \operatorname{Lean\mathord{:}Omega\mathord{:}Coeffs\mathord{:}of\,List}(\langle |t|\rangle \mathbin{\Vert} (\langle |s|\rangle \mathbin{\Vert} \langle |f[s]|\rangle))\), we obtain \(\operatorname{sat^{\prime}}(\operatorname{Lean\mathord{:}Omega\mathord{:}Constraint\mathord{:}combo}(-1,\allowbreak{} \operatorname{Lean\mathord{:}Omega\mathord{:}Constraint\mathord{:}mk}(\operatorname{some}(0),\allowbreak{} \operatorname{some}(0)),\allowbreak{} 1,\allowbreak{} \operatorname{Lean\mathord{:}Omega\mathord{:}Constraint\mathord{:}mk}(\operatorname{some}(0),\allowbreak{} \operatorname{none})),\allowbreak{} \operatorname{Lean\mathord{:}Omega\mathord{:}Coeffs\mathord{:}combo}(-1,\allowbreak{} \langle 1\rangle \mathbin{\Vert} (\langle 0\rangle \mathbin{\Vert} \langle -1\rangle),\allowbreak{} 1,\allowbreak{} \langle 1\rangle \mathbin{\Vert} \langle -1\rangle),\allowbreak{} \operatorname{Lean\mathord{:}Omega\mathord{:}Coeffs\mathord{:}of\,List}(\langle |t|\rangle \mathbin{\Vert} (\langle |s|\rangle \mathbin{\Vert} \langle |f[s]|\rangle))) = \mathrm{true}\).
By \paperlink{premise:4c65616e2e4f6d6567612e746964795f736174}{Prerequisite~P\ref{premise:4c65616e2e4f6d6567612e746964795f736174}}, we obtain \(\operatorname{sat^{\prime}}(\operatorname{tidy\,Constraint}(\operatorname{Lean\mathord{:}Omega\mathord{:}Constraint\mathord{:}mk}(\operatorname{some}(0),\allowbreak{} \operatorname{none}),\allowbreak{} \langle 0\rangle \mathbin{\Vert} (\langle -1\rangle \mathbin{\Vert} \langle 1\rangle)),\allowbreak{} \operatorname{tidy\,Coeffs}(\operatorname{Lean\mathord{:}Omega\mathord{:}Constraint\mathord{:}mk}(\operatorname{some}(0),\allowbreak{} \operatorname{none}),\allowbreak{} \langle 0\rangle \mathbin{\Vert} (\langle -1\rangle \mathbin{\Vert} \langle 1\rangle)),\allowbreak{} \operatorname{Lean\mathord{:}Omega\mathord{:}Coeffs\mathord{:}of\,List}(\langle |t|\rangle \mathbin{\Vert} (\langle |s|\rangle \mathbin{\Vert} \langle |f[s]|\rangle))) = \mathrm{true}\).
By \paperlink{premise:4c65616e2e4f6d6567612e436f6e73747261696e742e636f6d62696e655f73617427}{Prerequisite~P\ref{premise:4c65616e2e4f6d6567612e436f6e73747261696e742e636f6d62696e655f73617427}}, we obtain \(\operatorname{sat^{\prime}}(\operatorname{combine}(\operatorname{Lean\mathord{:}Omega\mathord{:}Constraint\mathord{:}mk}(\operatorname{some}(0),\allowbreak{} \operatorname{none}),\allowbreak{} \operatorname{Lean\mathord{:}Omega\mathord{:}Constraint\mathord{:}mk}(\operatorname{none},\allowbreak{} \operatorname{some}(0))),\allowbreak{} \langle 0\rangle \mathbin{\Vert} (\langle 1\rangle \mathbin{\Vert} \langle -1\rangle),\allowbreak{} \operatorname{Lean\mathord{:}Omega\mathord{:}Coeffs\mathord{:}of\,List}(\langle |t|\rangle \mathbin{\Vert} (\langle |s|\rangle \mathbin{\Vert} \langle |f[s]|\rangle))) = \mathrm{true}\).
Use the hypothesis \(\left|f\left[s\right]\right| > \left|s\right|\).
By \paperlink{premise:4c65616e2e4f6d6567612e496e742e6f664e61745f6c745f6f665f6c74}{Prerequisite~P\ref{premise:4c65616e2e4f6d6567612e496e742e6f664e61745f6c745f6f665f6c74}}, we obtain \(\left|s\right| < \left|f\left[s\right]\right|\).
By \paperlink{premise:496e742e6164645f6f6e655f6c655f6f665f6c74}{Prerequisite~P\ref{premise:496e742e6164645f6f6e655f6c655f6f665f6c74}}, we obtain \(\left|s\right| + 1 \le \left|f\left[s\right]\right|\).
By \paperlink{premise:496e742e7375625f6e6f6e6e65675f6f665f6c65}{Prerequisite~P\ref{premise:496e742e7375625f6e6f6e6e65675f6f665f6c65}}, we obtain \(0 \le \left|f\left[s\right]\right| - \left(\left|s\right| + 1\right)\).
Using \paperlink{premise:4c65616e2e4f6d6567612e4c696e656172436f6d626f2e636f6f7264696e6174655f6576616c5f32}{Prerequisite~P\ref{premise:4c65616e2e4f6d6567612e4c696e656172436f6d626f2e636f6f7264696e6174655f6576616c5f32}} in reverse gives \(|f[s]| = \operatorname{eval}(\operatorname{coordinate}(2),\allowbreak{} \operatorname{Lean\mathord{:}Omega\mathord{:}Coeffs\mathord{:}of\,List}(\langle |t|\rangle \mathbin{\Vert} (\langle |s|\rangle \mathbin{\Vert} \langle |f[s]|\rangle)))\).
Using \paperlink{premise:4c65616e2e4f6d6567612e4c696e656172436f6d626f2e636f6f7264696e6174655f6576616c5f31}{Prerequisite~P\ref{premise:4c65616e2e4f6d6567612e4c696e656172436f6d626f2e636f6f7264696e6174655f6576616c5f31}} in reverse gives \(|s| = \operatorname{eval}(\operatorname{coordinate}(1),\allowbreak{} \operatorname{Lean\mathord{:}Omega\mathord{:}Coeffs\mathord{:}of\,List}(\langle |t|\rangle \mathbin{\Vert} (\langle |s|\rangle \mathbin{\Vert} \langle |f[s]|\rangle)))\).
By definition, we obtain \(1 = \operatorname{eval}(\operatorname{Lean\mathord{:}Omega\mathord{:}Linear\,Combo\mathord{:}mk}(1,\allowbreak{} \langle\,\rangle),\allowbreak{} \operatorname{Lean\mathord{:}Omega\mathord{:}Coeffs\mathord{:}of\,List}(\langle |t|\rangle \mathbin{\Vert} (\langle |s|\rangle \mathbin{\Vert} \langle |f[s]|\rangle)))\).
By \paperlink{premise:4c65616e2e4f6d6567612e496e742e6164645f636f6e6772}{Prerequisite~P\ref{premise:4c65616e2e4f6d6567612e496e742e6164645f636f6e6772}}, we obtain \(|s| + 1 = \operatorname{eval}(\operatorname{coordinate}(1),\allowbreak{} \operatorname{Lean\mathord{:}Omega\mathord{:}Coeffs\mathord{:}of\,List}(\langle |t|\rangle \mathbin{\Vert} (\langle |s|\rangle \mathbin{\Vert} \langle |f[s]|\rangle))) + \operatorname{eval}(\operatorname{Lean\mathord{:}Omega\mathord{:}Linear\,Combo\mathord{:}mk}(1,\allowbreak{} \langle\,\rangle),\allowbreak{} \operatorname{Lean\mathord{:}Omega\mathord{:}Coeffs\mathord{:}of\,List}(\langle |t|\rangle \mathbin{\Vert} (\langle |s|\rangle \mathbin{\Vert} \langle |f[s]|\rangle)))\).
Using \paperlink{premise:4c65616e2e4f6d6567612e4c696e656172436f6d626f2e6164645f6576616c}{Prerequisite~P\ref{premise:4c65616e2e4f6d6567612e4c696e656172436f6d626f2e6164645f6576616c}} with \(l_{1} = \operatorname{Lean\mathord{:}Omega\mathord{:}Linear\,Combo\mathord{:}mk}(0,\allowbreak{} \langle 0\rangle \mathbin{\Vert} \langle 1\rangle)\) and \(l_{2} = \operatorname{Lean\mathord{:}Omega\mathord{:}Linear\,Combo\mathord{:}mk}(1,\allowbreak{} \langle\,\rangle)\) and \(v = \operatorname{Lean\mathord{:}Omega\mathord{:}Coeffs\mathord{:}of\,List}(\langle |t|\rangle \mathbin{\Vert} (\langle |s|\rangle \mathbin{\Vert} \langle |f[s]|\rangle))\) in reverse gives \(\operatorname{eval}(\operatorname{Lean\mathord{:}Omega\mathord{:}Linear\,Combo\mathord{:}mk}(0,\allowbreak{} \langle 0\rangle \mathbin{\Vert} \langle 1\rangle),\allowbreak{} \operatorname{Lean\mathord{:}Omega\mathord{:}Coeffs\mathord{:}of\,List}(\langle |t|\rangle \mathbin{\Vert} (\langle |s|\rangle \mathbin{\Vert} \langle |f[s]|\rangle))) + \operatorname{eval}(\operatorname{Lean\mathord{:}Omega\mathord{:}Linear\,Combo\mathord{:}mk}(1,\allowbreak{} \langle\,\rangle),\allowbreak{} \operatorname{Lean\mathord{:}Omega\mathord{:}Coeffs\mathord{:}of\,List}(\langle |t|\rangle \mathbin{\Vert} (\langle |s|\rangle \mathbin{\Vert} \langle |f[s]|\rangle))) = \operatorname{eval}(\operatorname{Lean\mathord{:}Omega\mathord{:}Linear\,Combo\mathord{:}mk}(0,\allowbreak{} \langle 0\rangle \mathbin{\Vert} \langle 1\rangle) + \operatorname{Lean\mathord{:}Omega\mathord{:}Linear\,Combo\mathord{:}mk}(1,\allowbreak{} \langle\,\rangle),\allowbreak{} \operatorname{Lean\mathord{:}Omega\mathord{:}Coeffs\mathord{:}of\,List}(\langle |t|\rangle \mathbin{\Vert} (\langle |s|\rangle \mathbin{\Vert} \langle |f[s]|\rangle)))\).
Chaining the equalities, we obtain \(|s| + 1 = \operatorname{eval}(\operatorname{Lean\mathord{:}Omega\mathord{:}Linear\,Combo\mathord{:}mk}(0,\allowbreak{} \langle 0\rangle \mathbin{\Vert} \langle 1\rangle) + \operatorname{Lean\mathord{:}Omega\mathord{:}Linear\,Combo\mathord{:}mk}(1,\allowbreak{} \langle\,\rangle),\allowbreak{} \operatorname{Lean\mathord{:}Omega\mathord{:}Coeffs\mathord{:}of\,List}(\langle |t|\rangle \mathbin{\Vert} (\langle |s|\rangle \mathbin{\Vert} \langle |f[s]|\rangle)))\).
By \paperlink{premise:4c65616e2e4f6d6567612e496e742e7375625f636f6e6772}{Prerequisite~P\ref{premise:4c65616e2e4f6d6567612e496e742e7375625f636f6e6772}}, we obtain \(|f[s]| - (|s| + 1) = \operatorname{eval}(\operatorname{coordinate}(2),\allowbreak{} \operatorname{Lean\mathord{:}Omega\mathord{:}Coeffs\mathord{:}of\,List}(\langle |t|\rangle \mathbin{\Vert} (\langle |s|\rangle \mathbin{\Vert} \langle |f[s]|\rangle))) - \operatorname{eval}(\operatorname{Lean\mathord{:}Omega\mathord{:}Linear\,Combo\mathord{:}mk}(0,\allowbreak{} \langle 0\rangle \mathbin{\Vert} \langle 1\rangle) + \operatorname{Lean\mathord{:}Omega\mathord{:}Linear\,Combo\mathord{:}mk}(1,\allowbreak{} \langle\,\rangle),\allowbreak{} \operatorname{Lean\mathord{:}Omega\mathord{:}Coeffs\mathord{:}of\,List}(\langle |t|\rangle \mathbin{\Vert} (\langle |s|\rangle \mathbin{\Vert} \langle |f[s]|\rangle)))\).
Using \paperlink{premise:4c65616e2e4f6d6567612e4c696e656172436f6d626f2e7375625f6576616c}{Prerequisite~P\ref{premise:4c65616e2e4f6d6567612e4c696e656172436f6d626f2e7375625f6576616c}} with \(l_{1} = \operatorname{Lean\mathord{:}Omega\mathord{:}Linear\,Combo\mathord{:}mk}(0,\allowbreak{} \langle 0\rangle \mathbin{\Vert} (\langle 0\rangle \mathbin{\Vert} \langle 1\rangle))\) and \(l_{2} = \operatorname{Lean\mathord{:}Omega\mathord{:}Linear\,Combo\mathord{:}mk}(1,\allowbreak{} \langle 0\rangle \mathbin{\Vert} \langle 1\rangle)\) and \(v = \operatorname{Lean\mathord{:}Omega\mathord{:}Coeffs\mathord{:}of\,List}(\langle |t|\rangle \mathbin{\Vert} (\langle |s|\rangle \mathbin{\Vert} \langle |f[s]|\rangle))\) in reverse gives \(\operatorname{eval}(\operatorname{Lean\mathord{:}Omega\mathord{:}Linear\,Combo\mathord{:}mk}(0,\allowbreak{} \langle 0\rangle \mathbin{\Vert} (\langle 0\rangle \mathbin{\Vert} \langle 1\rangle)),\allowbreak{} \operatorname{Lean\mathord{:}Omega\mathord{:}Coeffs\mathord{:}of\,List}(\langle |t|\rangle \mathbin{\Vert} (\langle |s|\rangle \mathbin{\Vert} \langle |f[s]|\rangle))) - \operatorname{eval}(\operatorname{Lean\mathord{:}Omega\mathord{:}Linear\,Combo\mathord{:}mk}(1,\allowbreak{} \langle 0\rangle \mathbin{\Vert} \langle 1\rangle),\allowbreak{} \operatorname{Lean\mathord{:}Omega\mathord{:}Coeffs\mathord{:}of\,List}(\langle |t|\rangle \mathbin{\Vert} (\langle |s|\rangle \mathbin{\Vert} \langle |f[s]|\rangle))) = \operatorname{eval}(\operatorname{Lean\mathord{:}Omega\mathord{:}Linear\,Combo\mathord{:}mk}(0,\allowbreak{} \langle 0\rangle \mathbin{\Vert} (\langle 0\rangle \mathbin{\Vert} \langle 1\rangle)) - \operatorname{Lean\mathord{:}Omega\mathord{:}Linear\,Combo\mathord{:}mk}(1,\allowbreak{} \langle 0\rangle \mathbin{\Vert} \langle 1\rangle),\allowbreak{} \operatorname{Lean\mathord{:}Omega\mathord{:}Coeffs\mathord{:}of\,List}(\langle |t|\rangle \mathbin{\Vert} (\langle |s|\rangle \mathbin{\Vert} \langle |f[s]|\rangle)))\).
Chaining the equalities, we obtain \(|f[s]| - (|s| + 1) = \operatorname{eval}(\operatorname{Lean\mathord{:}Omega\mathord{:}Linear\,Combo\mathord{:}mk}(0,\allowbreak{} \langle 0\rangle \mathbin{\Vert} (\langle 0\rangle \mathbin{\Vert} \langle 1\rangle)) - \operatorname{Lean\mathord{:}Omega\mathord{:}Linear\,Combo\mathord{:}mk}(1,\allowbreak{} \langle 0\rangle \mathbin{\Vert} \langle 1\rangle),\allowbreak{} \operatorname{Lean\mathord{:}Omega\mathord{:}Coeffs\mathord{:}of\,List}(\langle |t|\rangle \mathbin{\Vert} (\langle |s|\rangle \mathbin{\Vert} \langle |f[s]|\rangle)))\).
By \paperlink{premise:6c655f6f665f6c655f6f665f6571}{Prerequisite~P\ref{premise:6c655f6f665f6c655f6f665f6571}} with \(\mathbb{Z}\) as its type parameter \(\alpha\), we obtain \(0 \le \operatorname{eval}(\operatorname{Lean\mathord{:}Omega\mathord{:}Linear\,Combo\mathord{:}mk}(0,\allowbreak{} \langle 0\rangle \mathbin{\Vert} (\langle 0\rangle \mathbin{\Vert} \langle 1\rangle)) - \operatorname{Lean\mathord{:}Omega\mathord{:}Linear\,Combo\mathord{:}mk}(1,\allowbreak{} \langle 0\rangle \mathbin{\Vert} \langle 1\rangle),\allowbreak{} \operatorname{Lean\mathord{:}Omega\mathord{:}Coeffs\mathord{:}of\,List}(\langle |t|\rangle \mathbin{\Vert} (\langle |s|\rangle \mathbin{\Vert} \langle |f[s]|\rangle)))\).
By \paperlink{premise:4c65616e2e4f6d6567612e436f6e73747261696e742e616464496e657175616c6974795f736174}{Prerequisite~P\ref{premise:4c65616e2e4f6d6567612e436f6e73747261696e742e616464496e657175616c6974795f736174}}, we obtain \(\operatorname{sat^{\prime}}(\operatorname{Lean\mathord{:}Omega\mathord{:}Constraint\mathord{:}mk}(\operatorname{some}(-(-1)),\allowbreak{} \operatorname{none}),\allowbreak{} \langle 0\rangle \mathbin{\Vert} (\langle -1\rangle \mathbin{\Vert} \langle 1\rangle),\allowbreak{} \operatorname{Lean\mathord{:}Omega\mathord{:}Coeffs\mathord{:}of\,List}(\langle |t|\rangle \mathbin{\Vert} (\langle |s|\rangle \mathbin{\Vert} \langle |f[s]|\rangle))) = \mathrm{true}\).
By \paperlink{premise:4c65616e2e4f6d6567612e636f6d626f5f73617427}{Prerequisite~P\ref{premise:4c65616e2e4f6d6567612e636f6d626f5f73617427}} with \(\operatorname{Lean\mathord{:}Omega\mathord{:}Constraint\mathord{:}mk}(\operatorname{some}(0),\allowbreak{} \operatorname{some}(0))\) in place of its parameter \(s\) and \(\operatorname{Lean\mathord{:}Omega\mathord{:}Constraint\mathord{:}mk}(\operatorname{some}(1),\allowbreak{} \operatorname{none})\) in place of its parameter \(t\) and \(1\) in place of its parameter \(a\) and \(x = \langle 0\rangle \mathbin{\Vert} \left(\langle 1\rangle \mathbin{\Vert} \langle -1\rangle\right)\) and \(b = 1\) and \(y = \langle 0\rangle \mathbin{\Vert} \left(\langle -1\rangle \mathbin{\Vert} \langle 1\rangle\right)\) and \(v = \operatorname{Lean\mathord{:}Omega\mathord{:}Coeffs\mathord{:}of\,List}(\langle |t|\rangle \mathbin{\Vert} (\langle |s|\rangle \mathbin{\Vert} \langle |f[s]|\rangle))\), we obtain \(\operatorname{sat^{\prime}}(\operatorname{Lean\mathord{:}Omega\mathord{:}Constraint\mathord{:}combo}(1,\allowbreak{} \operatorname{Lean\mathord{:}Omega\mathord{:}Constraint\mathord{:}mk}(\operatorname{some}(0),\allowbreak{} \operatorname{some}(0)),\allowbreak{} 1,\allowbreak{} \operatorname{Lean\mathord{:}Omega\mathord{:}Constraint\mathord{:}mk}(\operatorname{some}(1),\allowbreak{} \operatorname{none})),\allowbreak{} \operatorname{Lean\mathord{:}Omega\mathord{:}Coeffs\mathord{:}combo}(1,\allowbreak{} \langle 0\rangle \mathbin{\Vert} (\langle 1\rangle \mathbin{\Vert} \langle -1\rangle),\allowbreak{} 1,\allowbreak{} \langle 0\rangle \mathbin{\Vert} (\langle -1\rangle \mathbin{\Vert} \langle 1\rangle)),\allowbreak{} \operatorname{Lean\mathord{:}Omega\mathord{:}Coeffs\mathord{:}of\,List}(\langle |t|\rangle \mathbin{\Vert} (\langle |s|\rangle \mathbin{\Vert} \langle |f[s]|\rangle))) = \mathrm{true}\).
By \paperlink{premise:4c65616e2e4f6d6567612e746964795f736174}{Prerequisite~P\ref{premise:4c65616e2e4f6d6567612e746964795f736174}}, we obtain \(\operatorname{sat^{\prime}}(\operatorname{tidy\,Constraint}(\operatorname{Lean\mathord{:}Omega\mathord{:}Constraint\mathord{:}mk}(\operatorname{some}(1),\allowbreak{} \operatorname{none}),\allowbreak{} \langle 0\rangle \mathbin{\Vert} (\langle 0\rangle \mathbin{\Vert} \langle 0\rangle)),\allowbreak{} \operatorname{tidy\,Coeffs}(\operatorname{Lean\mathord{:}Omega\mathord{:}Constraint\mathord{:}mk}(\operatorname{some}(1),\allowbreak{} \operatorname{none}),\allowbreak{} \langle 0\rangle \mathbin{\Vert} (\langle 0\rangle \mathbin{\Vert} \langle 0\rangle)),\allowbreak{} \operatorname{Lean\mathord{:}Omega\mathord{:}Coeffs\mathord{:}of\,List}(\langle |t|\rangle \mathbin{\Vert} (\langle |s|\rangle \mathbin{\Vert} \langle |f[s]|\rangle))) = \mathrm{true}\).
By \paperlink{premise:4c65616e2e4f6d6567612e436f6e73747261696e742e6e6f745f736174275f6f665f6973496d706f737369626c65}{Prerequisite~P\ref{premise:4c65616e2e4f6d6567612e436f6e73747261696e742e6e6f745f736174275f6f665f6973496d706f737369626c65}}, we obtain \(\mathrm{False}\).
\end{proofstep}
The two alternatives exhaust the cases. Hence \(\mathrm{False}\).
\end{proofstep}
The contrary assumption leads to a contradiction. Hence \(\left|f\left[s\right]\right| = \left|s\right|\).
Rewriting with \paperlink{theorem:48}{Corollary~\ref{theorem:48}} in the preceding result gives \(f\text{ is injective on }\operatorname{underlying\,set}(s)\).
\end{proof}

\begin{theorem}[{Uniqueness of the second coordinate for a fixed first coordinate}]\label{theorem:81}
Let \(\alpha\) and \(\beta\) be types. Let \(s\) be a finite subset of \(\alpha\). Let \(t\) be a finite subset of \(\beta\). Let \(f\) be a function of type \(\Pi_{a\, :\, \alpha}\, a \in s \to \beta\). Assume \(\forall a\, :\, \alpha,\; \forall \mathit{ha}\, :\, a \in s,\; f\left(a, \mathit{ha}\right) \in t\). Assume \(\forall b\, :\, \beta,\; b \in t \implies (\exists a_{1}\, :\, \alpha,\; \exists \mathit{ha}\text{ proving }(a_{1} \in s),\; f(a_{1},\allowbreak{} \mathit{ha}) = b)\). Assume \(\left|s\right| \le \left|t\right|\). Let \(a_{1}\) be an element of \(\alpha\). Let \(h\) be a proof of \(a_{1} \in s\). Let \(a_{2}\) be an element of \(\alpha\). Let \(h_{1}\) be a proof of \(a_{2} \in s\). Assume \(f\left(a_{1}, h\right) = f\left(a_{2}, h_{1}\right)\). Then \(a_{1} = a_{2}\).
\end{theorem}

\begin{proof}
Set \(f^{'}\) equal to \(a \mapsto f\left(a, \operatorname{property}\left(a\right)\right)\).
Fix \(a\) in \(s\).
Fix \(x\) in \(\beta\).
Assume \(x \in \operatorname{underlying\,set}(t)\).
The substitution uses the following argument.
\begin{proofstep}
Fix \(a\) in \(s\).
Assume \(a \in \operatorname{underlying\,set}(\operatorname{Finset\mathord{:}attach}\left(s\right))\).
By definition, we obtain \(f(\operatorname{val}(a),\allowbreak{} \operatorname{property}(a)) = f(\operatorname{val}(a),\allowbreak{} \operatorname{property}(a))\).
\end{proofstep}
The substitution uses the following argument.
\begin{proofstep}
Fix \(a\) in \(\alpha\).
For the required implication, argue as follows.
\begin{proofstep}
Assume \(a \in s\).
By definition, we obtain \((f(a,\allowbreak{} \operatorname{property}((a\, :\, s))) = x) = (f(a,\allowbreak{} \operatorname{property}((a\, :\, s))) = x)\).
By \paperlink{premise:4966662e6f665f6571}{Prerequisite~P\ref{premise:4966662e6f665f6571}}, we obtain \(f(a,\allowbreak{} \operatorname{property}((a\, :\, s))) = x \iff f(a,\allowbreak{} \operatorname{property}((a\, :\, s))) = x\).
\end{proofstep}
By definition, we obtain \(\left(a \in s\right) = \left(a \in s\right)\).
By \paperlink{premise:4966662e6f665f6571}{Prerequisite~P\ref{premise:4966662e6f665f6571}}, we obtain \(a \in s \iff a \in s\).
By \paperlink{premise:6578697374735f70726f705f636f6e6772}{Prerequisite~P\ref{premise:6578697374735f70726f705f636f6e6772}}, we obtain \((\exists b\text{ proving }(a \in s),\; f(a,\allowbreak{} \operatorname{property}((a\, :\, s))) = x) \iff (\exists h_{3}\text{ proving }(a \in s),\; f(a,\allowbreak{} \operatorname{property}((a\, :\, s))) = x)\).
By propositional extensionality, we obtain \((\exists b\text{ proving }(a \in s),\; f(a,\allowbreak{} \operatorname{property}((a\, :\, s))) = x) = (\exists h_{3}\text{ proving }(a \in s),\; f(a,\allowbreak{} \operatorname{property}((a\, :\, s))) = x)\).
\end{proofstep}
Use the hypothesis \(x \in \operatorname{underlying\,set}(t)\).
Applying the hypothesis \(\forall b\, :\, \beta,\; b \in t \implies (\exists a_{1}\, :\, \alpha,\; \exists \mathit{ha}\text{ proving }(a_{1} \in s),\; f(a_{1},\allowbreak{} \mathit{ha}) = b)\), we obtain \(\exists a\, :\, \alpha,\; \exists h_{3}\text{ proving }(a \in s),\; f(a,\allowbreak{} \operatorname{property}((a\, :\, s))) = x\).
Rewriting with \paperlink{premise:5365742e696d6167655f636f6e6772}{Prerequisite~P\ref{premise:5365742e696d6167655f636f6e6772}}, \paperlink{premise:46696e7365742e636f655f617474616368}{Prerequisite~P\ref{premise:46696e7365742e636f655f617474616368}}, \paperlink{premise:5365742e696d6167655f756e6976}{Prerequisite~P\ref{premise:5365742e696d6167655f756e6976}}, \paperlink{premise:5365742e6d656d5f72616e6765}{Prerequisite~P\ref{premise:5365742e6d656d5f72616e6765}}, \paperlink{premise:537562747970652e657869737473}{Prerequisite~P\ref{premise:537562747970652e657869737473}}, \paperlink{premise:66756e657874}{\emph{Function extensionality}} (\paperlink{premise:66756e657874}{Prerequisite~P\ref{premise:66756e657874}}) and the two identities established above in the preceding result gives \(x \in f^{'}\left[\operatorname{underlying\,set}(\operatorname{Finset\mathord{:}attach}\left(s\right))\right]\).
For the required implication, argue as follows.
\begin{proofstep}
Fix \(x\) in \(s\).
Assume \(x \in \operatorname{underlying\,set}(\operatorname{Finset\mathord{:}attach}\left(s\right))\).
By \paperlink{premise:537562747970652e70726f7065727479}{Prerequisite~P\ref{premise:537562747970652e70726f7065727479}} with \(\mathit{self} = x\), we obtain \(\operatorname{val}(x) \in s\).
Applying the hypothesis \(\forall a\, :\, \alpha,\; \forall \mathit{ha}\, :\, a \in s,\; f\left(a, \mathit{ha}\right) \in t\), we obtain \(f\left(\operatorname{val}(x), \operatorname{property}\left(x\right)\right) \in t\).
\end{proofstep}
Use the hypothesis \(\left|s\right| \le \left|t\right|\).
Rewriting with \paperlink{theorem:35}{Lemma~\ref{theorem:35}} in the preceding result gives \(\left|\operatorname{Finset\mathord{:}attach}\left(s\right)\right| \le \left|t\right|\).
By \paperlink{theorem:80}{Lemma~\ref{theorem:80}} with \(s\) as its type parameter \(\alpha\) and \(f^{'}\) in place of its parameter \(f\), we obtain \(f^{'}\text{ is injective on }\operatorname{underlying\,set}(\operatorname{Finset\mathord{:}attach}(s))\).
Rewriting with \paperlink{premise:46696e7365742e636f655f617474616368}{Prerequisite~P\ref{premise:46696e7365742e636f655f617474616368}} and \paperlink{premise:5365742e6d656d5f756e6976}{Prerequisite~P\ref{premise:5365742e6d656d5f756e6976}} gives \((a_{1}\, :\, s) \in \operatorname{underlying\,set}(\operatorname{Finset\mathord{:}attach}(s))\).
Rewriting with \paperlink{premise:46696e7365742e636f655f617474616368}{Prerequisite~P\ref{premise:46696e7365742e636f655f617474616368}} and \paperlink{premise:5365742e6d656d5f756e6976}{Prerequisite~P\ref{premise:5365742e6d656d5f756e6976}} gives \((a_{2}\, :\, s) \in \operatorname{underlying\,set}(\operatorname{Finset\mathord{:}attach}(s))\).
Use the hypothesis \(f\left(a_{1}, h\right) = f\left(a_{2}, h_{1}\right)\).
Applying the hypothesis \(f^{'}\text{ is injective on }\operatorname{underlying\,set}(\operatorname{Finset\mathord{:}attach}(s))\), we obtain \(\left(a_{1}\, :\, s\right) = \left(a_{2}\, :\, s\right)\).
Applying the same function to equal arguments, we obtain \(a_{1} = a_{2}\).
\end{proof}

\section{Finite sets}
\begin{theorem}\label{theorem:82}
Let \(\alpha\) be a type. Let \(s\) and \(t\) be finite subsets of \(\alpha\). Let \(h\) be a proof of \(\operatorname{Disjoint}\left(s, t\right)\). Then \(\left|\operatorname{disj\,Union}\left(s, t\right)\right| = \left|s\right| + \left|t\right|\).
\end{theorem}

\begin{proof}
By \paperlink{premise:4d756c74697365742e636172645f616464}{Prerequisite~P\ref{premise:4d756c74697365742e636172645f616464}} with \(\operatorname{underlying\,multiset}(s)\) in place of its parameter \(s\) and \(\operatorname{underlying\,multiset}(t)\) in place of its parameter \(t\), we obtain \(\left|\operatorname{disj\,Union}\left(s, t\right)\right| = \left|s\right| + \left|t\right|\).
\end{proof}

\section{Lattice}
\begin{lemma}\label{theorem:83}
Let \(\alpha\) be a type. Assume equality on \(\alpha\) is decidable. Let \(s\) and \(t\) be finite subsets of \(\alpha\). Then \(\left|s \cup t\right| + \left|s \cap t\right| = \left|s\right| + \left|t\right|\).
\end{lemma}

\begin{proof}
Rewriting with \paperlink{premise:46696e7365742e756e696f6e5f656d707479}{Prerequisite~P\ref{premise:46696e7365742e756e696f6e5f656d707479}} and \paperlink{premise:46696e7365742e696e7465725f656d707479}{Prerequisite~P\ref{premise:46696e7365742e696e7465725f656d707479}} gives \(\left|s \cup \varnothing\right| + \left|s \cap \varnothing\right| = \left|s\right|\).
For the required implication, argue as follows.
\begin{proofstep}
Fix \(a\) in \(\alpha\).
Fix \(r\) in \(\mathcal{P}_{\mathrm{fin}}(\alpha)\).
Assume \(a \notin r\).
Assume \(\left|s \cup r\right| + \left|s \cap r\right| = \left|s\right| + \left|r\right|\).
Distinguish whether \(a \in s\) holds.
Consider the following case.
\begin{proofstep}
Assume \(a \in s\).
Rewriting with \paperlink{premise:46696e7365742e756e696f6e5f696e73657274}{Prerequisite~P\ref{premise:46696e7365742e756e696f6e5f696e73657274}}, \paperlink{premise:46696e7365742e696e736572745f65715f6f665f6d656d}{Prerequisite~P\ref{premise:46696e7365742e696e736572745f65715f6f665f6d656d}}, \paperlink{premise:46696e7365742e6d656d5f756e696f6e}{Prerequisite~P\ref{premise:46696e7365742e6d656d5f756e696f6e}}, \(a \in s\), \paperlink{premise:65715f66616c7365}{Prerequisite~P\ref{premise:65715f66616c7365}} (using \(a \notin r\)), \paperlink{premise:6f725f66616c7365}{Prerequisite~P\ref{premise:6f725f66616c7365}}, \paperlink{premise:46696e7365742e696e7465725f696e736572745f6f665f6d656d}{Prerequisite~P\ref{premise:46696e7365742e696e7465725f696e736572745f6f665f6d656d}} (using \(a \in s\)), \paperlink{theorem:17}{Lemma~\ref{theorem:17}}, \paperlink{premise:46696e7365742e6d656d5f696e746572}{Prerequisite~P\ref{premise:46696e7365742e6d656d5f696e746572}}, \paperlink{premise:616e645f66616c7365}{Prerequisite~P\ref{premise:616e645f66616c7365}}, \paperlink{premise:6e6f745f66616c73655f65715f74727565}{Prerequisite~P\ref{premise:6e6f745f66616c73655f65715f74727565}}, \paperlink{premise:4e61742e6164645f6173736f63}{Prerequisite~P\ref{premise:4e61742e6164645f6173736f63}}, \(\left|s \cup r\right| + \left|s \cap r\right| = \left|s\right| + \left|r\right|\) and \paperlink{premise:4e61742e6164645f6c6566745f696e6a}{\emph{Alias}} (\paperlink{premise:4e61742e6164645f6c6566745f696e6a}{Prerequisite~P\ref{premise:4e61742e6164645f6c6566745f696e6a}}) gives \(|s \cup (\{a\} \cup r)| + |s \cap (\{a\} \cup r)| = |s| + |\{a\} \cup r|\).
\end{proofstep}
Consider the following case.
\begin{proofstep}
Assume \(a \notin s\).
Rewriting with \paperlink{premise:46696e7365742e756e696f6e5f696e73657274}{Prerequisite~P\ref{premise:46696e7365742e756e696f6e5f696e73657274}}, \paperlink{theorem:17}{Lemma~\ref{theorem:17}}, \paperlink{premise:46696e7365742e6d656d5f756e696f6e}{Prerequisite~P\ref{premise:46696e7365742e6d656d5f756e696f6e}}, \paperlink{premise:65715f66616c7365}{Prerequisite~P\ref{premise:65715f66616c7365}} (using \(a \notin s\)), \paperlink{premise:65715f66616c7365}{Prerequisite~P\ref{premise:65715f66616c7365}} (using \(a \notin r\)), \paperlink{premise:6f725f73656c66}{Prerequisite~P\ref{premise:6f725f73656c66}}, \paperlink{premise:6e6f745f66616c73655f65715f74727565}{Prerequisite~P\ref{premise:6e6f745f66616c73655f65715f74727565}}, \paperlink{premise:46696e7365742e696e7465725f696e736572745f6f665f6e6f745f6d656d}{Prerequisite~P\ref{premise:46696e7365742e696e7465725f696e736572745f6f665f6e6f745f6d656d}}, \paperlink{premise:4e61742e6164645f72696768745f636f6d6d}{Prerequisite~P\ref{premise:4e61742e6164645f72696768745f636f6d6d}}, \(\left|s \cup r\right| + \left|s \cap r\right| = \left|s\right| + \left|r\right|\), \paperlink{premise:4e61742e6164645f6173736f63}{Prerequisite~P\ref{premise:4e61742e6164645f6173736f63}} and \paperlink{premise:4e61742e6164645f6c6566745f696e6a}{\emph{Alias}} (\paperlink{premise:4e61742e6164645f6c6566745f696e6a}{Prerequisite~P\ref{premise:4e61742e6164645f6c6566745f696e6a}}) gives \(|s \cup (\{a\} \cup r)| + |s \cap (\{a\} \cup r)| = |s| + |\{a\} \cup r|\).
\end{proofstep}
The proposition or its negation holds. Hence \(|s \cup (\{a\} \cup r)| + |s \cap (\{a\} \cup r)| = |s| + |\{a\} \cup r|\).
\end{proofstep}
By \paperlink{premise:46696e7365742e696e64756374696f6e5f6f6e}{Prerequisite~P\ref{premise:46696e7365742e696e64756374696f6e5f6f6e}} with \(t\) in place of its parameter \(s\), we obtain \(\left|s \cup t\right| + \left|s \cap t\right| = \left|s\right| + \left|t\right|\).
\end{proof}

\begin{lemma}\label{theorem:84}
Let \(\alpha\) be a type. Assume equality on \(\alpha\) is decidable. Let \(s\) and \(t\) be finite subsets of \(\alpha\). Then \(\left|s \cap t\right| + \left|s \cup t\right| = \left|s\right| + \left|t\right|\).
\end{lemma}

\begin{proof}
By \paperlink{premise:4e61742e6164645f636f6d6d}{Prerequisite~P\ref{premise:4e61742e6164645f636f6d6d}} with \(n = \left|s \cap t\right|\) and \(m = \left|s \cup t\right|\), we obtain \(\left|s \cap t\right| + \left|s \cup t\right| = \left|s \cup t\right| + \left|s \cap t\right|\).
Applying the same function to equal arguments, we obtain \((|s \cap t| + |s \cup t| = |s| + |t|) = (|s \cup t| + |s \cap t| = |s| + |t|)\).
By the preceding argument, we obtain \((|s \cap t| + |s \cup t| = |s| + |t|) = (|s \cup t| + |s \cap t| = |s| + |t|)\).
By \paperlink{theorem:83}{Lemma~\ref{theorem:83}}, we obtain \(\left|s \cup t\right| + \left|s \cap t\right| = \left|s\right| + \left|t\right|\).
Applying the same function to equal arguments, we obtain \((|s \cup t| + |s \cap t| = |s| + |t|) = (|s| + |t| = |s| + |t|)\).
By the preceding argument, we obtain \((|s \cup t| + |s \cap t| = |s| + |t|) = (|s| + |t| = |s| + |t|)\).
By definition, we obtain \(\left|s\right| + \left|t\right| = \left|s\right| + \left|t\right|\).
Substituting along the equality, we obtain \(\left|s \cup t\right| + \left|s \cap t\right| = \left|s\right| + \left|t\right|\).
Substituting along the equality, we obtain \(\left|s \cap t\right| + \left|s \cup t\right| = \left|s\right| + \left|t\right|\).
\end{proof}

\begin{theorem}\label{theorem:85}
Let \(\alpha\) be a type. Assume equality on \(\alpha\) is decidable. Let \(s\) and \(t\) be finite subsets of \(\alpha\). Then \(\left|s \cup t\right| = \left|s\right| + \left|t\right| - \left|s \cap t\right|\).
\end{theorem}

\begin{proof}
By \paperlink{theorem:83}{Lemma~\ref{theorem:83}}, we obtain \(\left|s \cup t\right| + \left|s \cap t\right| = \left|s\right| + \left|t\right|\).
Reversing the equality, we obtain \(\left|s\right| + \left|t\right| = \left|s \cup t\right| + \left|s \cap t\right|\).
Applying the same function to equal arguments, we obtain \((|s \cup t| = |s| + |t| - |s \cap t|) = (|s \cup t| = |s \cup t| + |s \cap t| - |s \cap t|)\).
By the preceding argument, we obtain \((|s \cup t| = |s| + |t| - |s \cap t|) = (|s \cup t| = |s \cup t| + |s \cap t| - |s \cap t|)\).
By \paperlink{premise:4e61742e6164645f7375625f63616e63656c}{Prerequisite~P\ref{premise:4e61742e6164645f7375625f63616e63656c}} with \(n = \left|s \cup t\right|\) and \(m = \left|s \cap t\right|\), we obtain \(\left|s \cup t\right| + \left|s \cap t\right| - \left|s \cap t\right| = \left|s \cup t\right|\).
Applying the same function to equal arguments, we obtain \((|s \cup t| = |s \cup t| + |s \cap t| - |s \cap t|) = (|s \cup t| = |s \cup t|)\).
By the preceding argument, we obtain \((|s \cup t| = |s \cup t| + |s \cap t| - |s \cap t|) = (|s \cup t| = |s \cup t|)\).
By definition, we obtain \(\left|s \cup t\right| = \left|s \cup t\right|\).
Substituting along the equality, we obtain \(\left|s \cup t\right| = \left|s \cup t\right| + \left|s \cap t\right| - \left|s \cap t\right|\).
Substituting along the equality, we obtain \(\left|s \cup t\right| = \left|s\right| + \left|t\right| - \left|s \cap t\right|\).
\end{proof}

\begin{theorem}\label{theorem:86}
Let \(\alpha\) be a type. Assume equality on \(\alpha\) is decidable. Let \(s\) and \(t\) be finite subsets of \(\alpha\). Then \(\left|s \cap t\right| = \left|s\right| + \left|t\right| - \left|s \cup t\right|\).
\end{theorem}

\begin{proof}
By \paperlink{theorem:84}{Lemma~\ref{theorem:84}}, we obtain \(\left|s \cap t\right| + \left|s \cup t\right| = \left|s\right| + \left|t\right|\).
Reversing the equality, we obtain \(\left|s\right| + \left|t\right| = \left|s \cap t\right| + \left|s \cup t\right|\).
Applying the same function to equal arguments, we obtain \((|s \cap t| = |s| + |t| - |s \cup t|) = (|s \cap t| = |s \cap t| + |s \cup t| - |s \cup t|)\).
By the preceding argument, we obtain \((|s \cap t| = |s| + |t| - |s \cup t|) = (|s \cap t| = |s \cap t| + |s \cup t| - |s \cup t|)\).
By \paperlink{premise:4e61742e6164645f7375625f63616e63656c}{Prerequisite~P\ref{premise:4e61742e6164645f7375625f63616e63656c}} with \(n = \left|s \cap t\right|\) and \(m = \left|s \cup t\right|\), we obtain \(\left|s \cap t\right| + \left|s \cup t\right| - \left|s \cup t\right| = \left|s \cap t\right|\).
Applying the same function to equal arguments, we obtain \((|s \cap t| = |s \cap t| + |s \cup t| - |s \cup t|) = (|s \cap t| = |s \cap t|)\).
By the preceding argument, we obtain \((|s \cap t| = |s \cap t| + |s \cup t| - |s \cup t|) = (|s \cap t| = |s \cap t|)\).
By definition, we obtain \(\left|s \cap t\right| = \left|s \cap t\right|\).
Substituting along the equality, we obtain \(\left|s \cap t\right| = \left|s \cap t\right| + \left|s \cup t\right| - \left|s \cup t\right|\).
Substituting along the equality, we obtain \(\left|s \cap t\right| = \left|s\right| + \left|t\right| - \left|s \cup t\right|\).
\end{proof}

\begin{theorem}\label{theorem:87}
Let \(\alpha\) be a type. Assume equality on \(\alpha\) is decidable. Let \(s\) and \(t\) be finite subsets of \(\alpha\). Then \(\left|s \cup t\right| \le \left|s\right| + \left|t\right|\).
\end{theorem}

\begin{proof}
By \paperlink{premise:4e61742e6c655f6164645f7269676874}{Prerequisite~P\ref{premise:4e61742e6c655f6164645f7269676874}} with \(n = \left|s \cup t\right|\) and \(k = \left|s \cap t\right|\), we obtain \(\left|s \cup t\right| \le \left|s \cup t\right| + \left|s \cap t\right|\).
By \paperlink{theorem:83}{Lemma~\ref{theorem:83}}, we obtain \(\left|s \cup t\right| + \left|s \cap t\right| = \left|s\right| + \left|t\right|\).
Substituting equals, we obtain \(\left|s \cup t\right| \le \left|s\right| + \left|t\right|\).
\end{proof}

\begin{theorem}\label{theorem:88}
Let \(\alpha\) be a type. Let \(s\) and \(t\) be finite subsets of \(\alpha\). Assume equality on \(\alpha\) is decidable. Then
\begin{center}
\(\displaystyle |s \cup t| = |s| + |t| \iff \operatorname{Disjoint}(s,\allowbreak{} t).\)
\end{center}
\end{theorem}

\begin{proof}
Rewriting with \paperlink{premise:4e61742e73656c665f65715f6164645f7269676874}{Prerequisite~P\ref{premise:4e61742e73656c665f65715f6164645f7269676874}}, \paperlink{theorem:7}{Lemma~\ref{theorem:7}}, \paperlink{premise:46696e7365742e6469736a6f696e745f6966665f696e7465725f65715f656d707479}{Prerequisite~P\ref{premise:46696e7365742e6469736a6f696e745f6966665f696e7465725f65715f656d707479}} and \paperlink{premise:6966665f73656c66}{Prerequisite~P\ref{premise:6966665f73656c66}} gives \(|s \cup t| = |s \cup t| + |s \cap t| \iff \operatorname{Disjoint}(s,\allowbreak{} t)\).
Rewriting with \paperlink{theorem:83}{Lemma~\ref{theorem:83}} in the preceding result gives \(\left|s \cup t\right| = \left|s\right| + \left|t\right| \iff \operatorname{Disjoint}\left(s, t\right)\).
\end{proof}

\begin{lemma}[{Alias}]\label{theorem:89}
Let \(\alpha\) be a type. Let \(s\) and \(t\) be finite subsets of \(\alpha\). Assume equality on \(\alpha\) is decidable. Assume \(\operatorname{Disjoint}\left(s, t\right)\). Then \(\left|s \cup t\right| = \left|s\right| + \left|t\right|\).
\end{lemma}

\begin{proof}
Apply the reverse implication of \paperlink{theorem:88}{Theorem~\ref{theorem:88}}.
\end{proof}

\begin{lemma}\label{theorem:90}
Let \(\alpha\) be a type. Let \(s\) and \(t\) be finite subsets of \(\alpha\). Assume equality on \(\alpha\) is decidable. Assume \(s \subseteq t\). Then \(\left|t \setminus s\right| = \left|t\right| - \left|s\right|\).
\end{lemma}

\begin{proof}
Rewriting with \paperlink{premise:4e61742e6164645f7375625f63616e63656c5f7269676874}{Prerequisite~P\ref{premise:4e61742e6164645f7375625f63616e63656c5f7269676874}}, \paperlink{theorem:89}{Lemma~\ref{theorem:89}} and \paperlink{premise:46696e7365742e73646966665f6469736a6f696e74}{Prerequisite~P\ref{premise:46696e7365742e73646966665f6469736a6f696e74}} gives \(\left|t \setminus s\right| = \left|\left(t \setminus s\right) \cup s\right| - \left|s\right|\).
Use the previously proved fact \(\left|t \setminus s\right| = \left|\left(t \setminus s\right) \cup s\right| - \left|s\right|\).
Rewriting with \paperlink{premise:46696e7365742e73646966665f756e696f6e5f6f665f737562736574}{Prerequisite~P\ref{premise:46696e7365742e73646966665f756e696f6e5f6f665f737562736574}} (using \(s \subseteq t\)) in the preceding result gives \(\left|t \setminus s\right| = \left|t\right| - \left|s\right|\).
\end{proof}

\begin{theorem}\label{theorem:91}
Let \(\alpha\) be a type. Assume equality on \(\alpha\) is decidable. Let \(s\) and \(t\) be finite subsets of \(\alpha\). Assume \(s \subseteq t\). Then \(\left|t \setminus s\right| + \left|s\right| = \left|t\right|\).
\end{theorem}

\begin{proof}
By \paperlink{theorem:5}{Lemma~\ref{theorem:5}}, we obtain \(\left|s\right| \le \left|t\right|\).
By \paperlink{premise:4e61742e7375625f65715f6966665f65715f616464}{Prerequisite~P\ref{premise:4e61742e7375625f65715f6966665f65715f616464}}, we obtain \(|t| - |s| = |t \setminus s| \iff |t| = |t \setminus s| + |s|\).
By \paperlink{theorem:90}{Lemma~\ref{theorem:90}}, we obtain \(\left|t \setminus s\right| = \left|t\right| - \left|s\right|\).
Reversing the equality, we obtain \(\left|t\right| - \left|s\right| = \left|t \setminus s\right|\).
By the forward implication, we obtain \(\left|t\right| = \left|t \setminus s\right| + \left|s\right|\).
Reversing the equality, we obtain \(\left|t \setminus s\right| + \left|s\right| = \left|t\right|\).
\end{proof}

\begin{lemma}\label{theorem:92}
Let \(\alpha\) be a type. Assume equality on \(\alpha\) is decidable. Let \(s\) and \(t\) be finite subsets of \(\alpha\). Then \(\left|t\right| - \left|s\right| \le \left|t \setminus s\right|\).
\end{lemma}

\begin{proof}
By \paperlink{premise:46696e7365742e696e7465725f7375627365745f6c656674}{Prerequisite~P\ref{premise:46696e7365742e696e7465725f7375627365745f6c656674}}, we obtain \(s \cap t \subseteq s\).
By \paperlink{theorem:5}{Lemma~\ref{theorem:5}}, we obtain \(\left|s \cap t\right| \le \left|s\right|\).
By \paperlink{premise:4e61742e7375625f6c655f7375625f6c656674}{Prerequisite~P\ref{premise:4e61742e7375625f6c655f7375625f6c656674}} with \(k = \left|t\right|\), we obtain \(\left|t\right| - \left|s\right| \le \left|t\right| - \left|s \cap t\right|\).
By \paperlink{premise:46696e7365742e696e7465725f7375627365745f7269676874}{Prerequisite~P\ref{premise:46696e7365742e696e7465725f7375627365745f7269676874}}, we obtain \(s \cap t \subseteq t\).
By \paperlink{theorem:90}{Lemma~\ref{theorem:90}}, we obtain \(\left|t \setminus \left(s \cap t\right)\right| = \left|t\right| - \left|s \cap t\right|\).
Reversing the equality, we obtain \(\left|t\right| - \left|s \cap t\right| = \left|t \setminus \left(s \cap t\right)\right|\).
By transitivity of the displayed relations, we obtain \(\left|t\right| - \left|s\right| \le \left|t \setminus \left(s \cap t\right)\right|\).
By \paperlink{premise:46696e7365742e73646966665f696e7465725f73656c665f7269676874}{Prerequisite~P\ref{premise:46696e7365742e73646966665f696e7465725f73656c665f7269676874}} with \(t\) in place of its parameter \(s\) and \(s\) in place of its parameter \(t\), we obtain \(t \setminus \left(s \cap t\right) = t \setminus s\).
Applying the same function to equal arguments, we obtain \((|t \setminus (s \cap t)| \le |t \setminus s|) = (|t \setminus s| \le |t \setminus s|)\).
By the preceding argument, we obtain \((|t \setminus (s \cap t)| \le |t \setminus s|) = (|t \setminus s| \le |t \setminus s|)\).
By \paperlink{premise:6c655f7265666c}{Prerequisite~P\ref{premise:6c655f7265666c}} with \(\mathbb{N}\) as its type parameter \(\alpha\) and \(a = \left|t \setminus s\right|\), we obtain \(\left|t \setminus s\right| \le \left|t \setminus s\right|\).
Substituting along the equality, we obtain \(\left|t \setminus \left(s \cap t\right)\right| \le \left|t \setminus s\right|\).
By transitivity of the displayed relations, we obtain \(\left|t\right| - \left|s\right| \le \left|t \setminus s\right|\).
\end{proof}

\begin{theorem}\label{theorem:93}
Let \(\alpha\) be a type. Let \(s\) and \(t\) be finite subsets of \(\alpha\). Assume equality on \(\alpha\) is decidable. Then \(\left|s\right| \le \left|s \setminus t\right| + \left|t\right|\).
\end{theorem}

\begin{proof}
By \paperlink{premise:4e61742e7375625f6c655f6966665f6c655f616464}{Prerequisite~P\ref{premise:4e61742e7375625f6c655f6966665f6c655f616464}}, we obtain \(|s| - |t| \le |s \setminus t| \iff |s| \le |s \setminus t| + |t|\).
By \paperlink{theorem:92}{Lemma~\ref{theorem:92}} with \(t\) in place of its parameter \(s\) and \(s\) in place of its parameter \(t\), we obtain \(\left|s\right| - \left|t\right| \le \left|s \setminus t\right|\).
By the forward implication, we obtain \(\left|s\right| \le \left|s \setminus t\right| + \left|t\right|\).
\end{proof}

\begin{lemma}\label{theorem:94}
Let \(\alpha\) be a type. Assume equality on \(\alpha\) is decidable. Let \(s\) and \(t\) be finite subsets of \(\alpha\). Then \(\left|s \setminus t\right| + \left|t\right| = \left|s \cup t\right|\).
\end{lemma}

\begin{proof}
By \paperlink{premise:46696e7365742e73646966665f756e696f6e5f73656c665f65715f756e696f6e}{Prerequisite~P\ref{premise:46696e7365742e73646966665f756e696f6e5f73656c665f65715f756e696f6e}}, we have \(\left(s \setminus t\right) \cup t = s \cup t\).
By \paperlink{premise:46696e7365742e73646966665f6469736a6f696e74}{Prerequisite~P\ref{premise:46696e7365742e73646966665f6469736a6f696e74}}, we have \(\operatorname{Disjoint}\left(s \setminus t, t\right)\).
By \paperlink{theorem:89}{Lemma~\ref{theorem:89}} applied to \(\operatorname{Disjoint}\left(s \setminus t, t\right)\), we have \(\left|\left(s \setminus t\right) \cup t\right| = \left|s \setminus t\right| + \left|t\right|\).
Starting from \paperlink{premise:45712e73796d6d}{Prerequisite~P\ref{premise:45712e73796d6d}} applied to \(\left|\left(s \setminus t\right) \cup t\right| = \left|s \setminus t\right| + \left|t\right|\), substitution using \paperlink{premise:46696e7365742e73646966665f756e696f6e5f73656c665f65715f756e696f6e}{Prerequisite~P\ref{premise:46696e7365742e73646966665f756e696f6e5f73656c665f65715f756e696f6e}} gives \(\left|s \setminus t\right| + \left|t\right| = \left|s \cup t\right|\).
\end{proof}

\begin{theorem}\label{theorem:95}
Let \(\alpha\) be a type. Let \(s\) and \(t\) be finite subsets of \(\alpha\). Assume equality on \(\alpha\) is decidable. Assume \(\left|s\right| = \left|t\right|\). Then \(\left|s \setminus t\right| = \left|t \setminus s\right|\).
\end{theorem}

\begin{proof}
Rewriting with \paperlink{premise:46696e7365742e756e696f6e5f636f6d6d}{Prerequisite~P\ref{premise:46696e7365742e756e696f6e5f636f6d6d}} gives \(\left|s \cup t\right| = \left|t \cup s\right|\).
Rewriting with \paperlink{theorem:94}{Lemma~\ref{theorem:94}} in the preceding result gives \(\left|s \cup t\right| = \left|t \setminus s\right| + \left|s\right|\).
Rewriting with \(\left|s\right| = \left|t\right|\) in the preceding result gives \(\left|s \cup t\right| = \left|t \setminus s\right| + \left|t\right|\).
Rewriting with \paperlink{theorem:94}{Lemma~\ref{theorem:94}} in the preceding result gives \(\left|s \setminus t\right| + \left|t\right| = \left|t \setminus s\right| + \left|t\right|\).
By \paperlink{premise:4e61742e6164645f72696768745f63616e63656c}{Prerequisite~P\ref{premise:4e61742e6164645f72696768745f63616e63656c}}, we obtain \(\left|s \setminus t\right| = \left|t \setminus s\right|\).
\end{proof}

\begin{lemma}\label{theorem:96}
Let \(\alpha\) be a type. Let \(s\) and \(t\) be finite subsets of \(\alpha\). Assume equality on \(\alpha\) is decidable. Assume \(\left|s\right| < \left|t\right|\). Then \(t \setminus s\ne\varnothing\).
\end{lemma}

\begin{proof}
Assume \(t \subseteq s\).
Use the hypothesis \(\left|s\right| < \left|t\right|\).
By \paperlink{theorem:5}{Lemma~\ref{theorem:5}}, we obtain \(\left|t\right| \le \left|s\right|\).
By \paperlink{premise:4c542e6c742e6e6f745f6c65}{\emph{Alias}} (\paperlink{premise:4c542e6c742e6e6f745f6c65}{Prerequisite~P\ref{premise:4c542e6c742e6e6f745f6c65}}) with \(\mathbb{N}\) as its type parameter \(\alpha\), we obtain \(\mathrm{False}\).
Rewriting with \paperlink{premise:46696e7365742e73646966665f65715f656d7074795f6966665f737562736574}{Prerequisite~P\ref{premise:46696e7365742e73646966665f65715f656d7074795f6966665f737562736574}} in the preceding result gives \(\neg \left(t \setminus s = \varnothing\right)\).
Simplifying by definition in the preceding result gives \(t \setminus s \ne \varnothing\).
Rewriting with \paperlink{premise:46696e7365742e6e6f6e656d7074795f6966665f6e655f656d707479}{Prerequisite~P\ref{premise:46696e7365742e6e6f6e656d7074795f6966665f6e655f656d707479}} in the preceding result gives \(t \setminus s\ne\varnothing\).
\end{proof}

\begin{corollary}\label{theorem:97}
Let \(\alpha\) be a type. Let \(s\) and \(t\) be finite subsets of \(\alpha\). Assume \(\left|s\right| < \left|t\right|\). Then \(\exists e\, :\, \alpha,\; e \in t \land e \notin s\).
\end{corollary}

\begin{proof}
For the required implication, argue as follows.
\begin{proofstep}
Fix \(x\) in \(\alpha\).
By \paperlink{premise:46696e7365742e6d656d5f7364696666}{Prerequisite~P\ref{premise:46696e7365742e6d656d5f7364696666}}, we obtain \(x \in t \setminus s \iff x \in t \land x \notin s\).
By propositional extensionality, we obtain \(\left(x \in t \setminus s\right) = \left(x \in t \land x \notin s\right)\).
\end{proofstep}
By function extensionality, we obtain \(\left(y \mapsto y \in t \setminus s\right) = \left(x \mapsto x \in t \land x \notin s\right)\).
Applying the same function to equal arguments, we obtain \((\exists x\, :\, \alpha,\; (y \mapsto y \in t \setminus s)(x)) = (\exists x\, :\, \alpha,\; x \in t \land x \notin s)\).
By \paperlink{theorem:96}{Lemma~\ref{theorem:96}}, we obtain \(t \setminus s\ne\varnothing\).
Substituting along the equality, we obtain \(\exists e\, :\, \alpha,\; e \in t \land e \notin s\).
\end{proof}

\begin{lemma}\label{theorem:98}
Let \(\alpha\) be a type. Assume equality on \(\alpha\) is decidable. Let \(s\) and \(t\) be finite subsets of \(\alpha\). Then \(\left|s \setminus t\right| + \left|s \cap t\right| = \left|s\right|\).
\end{lemma}

\begin{proof}
By \paperlink{premise:46696e7365742e73646966665f756e696f6e5f696e746572}{Prerequisite~P\ref{premise:46696e7365742e73646966665f756e696f6e5f696e746572}}, we have \(\left(s \setminus t\right) \cup \left(s \cap t\right) = s\).
By \paperlink{premise:46696e7365742e6469736a6f696e745f73646966665f696e746572}{Prerequisite~P\ref{premise:46696e7365742e6469736a6f696e745f73646966665f696e746572}}, we have \(\operatorname{Disjoint}\left(s \setminus t, s \cap t\right)\).
By \paperlink{theorem:89}{Lemma~\ref{theorem:89}} applied to \(\operatorname{Disjoint}\left(s \setminus t, s \cap t\right)\), we have \(|(s \setminus t) \cup (s \cap t)| = |s \setminus t| + |s \cap t|\).
Starting from \paperlink{premise:45712e73796d6d}{Prerequisite~P\ref{premise:45712e73796d6d}} applied to \(|(s \setminus t) \cup (s \cap t)| = |s \setminus t| + |s \cap t|\), substitution using \paperlink{premise:46696e7365742e73646966665f756e696f6e5f696e746572}{Prerequisite~P\ref{premise:46696e7365742e73646966665f756e696f6e5f696e746572}} gives \(\left|s \setminus t\right| + \left|s \cap t\right| = \left|s\right|\).
\end{proof}

\begin{theorem}\label{theorem:99}
Let \(\alpha\) be a type. Assume equality on \(\alpha\) is decidable. Let \(s\) and \(t\) be finite subsets of \(\alpha\). Then \(\left|s \cap t\right| + \left|s \setminus t\right| = \left|s\right|\).
\end{theorem}

\begin{proof}
By \paperlink{premise:4e61742e6164645f636f6d6d}{Prerequisite~P\ref{premise:4e61742e6164645f636f6d6d}} with \(n = \left|s \cap t\right|\) and \(m = \left|s \setminus t\right|\), we obtain \(\left|s \cap t\right| + \left|s \setminus t\right| = \left|s \setminus t\right| + \left|s \cap t\right|\).
Applying the same function to equal arguments, we obtain \((|s \cap t| + |s \setminus t| = |s|) = (|s \setminus t| + |s \cap t| = |s|)\).
By the preceding argument, we obtain \((|s \cap t| + |s \setminus t| = |s|) = (|s \setminus t| + |s \cap t| = |s|)\).
By \paperlink{theorem:98}{Lemma~\ref{theorem:98}}, we obtain \(\left|s \setminus t\right| + \left|s \cap t\right| = \left|s\right|\).
Applying the same function to equal arguments, we obtain \((|s \setminus t| + |s \cap t| = |s|) = (|s| = |s|)\).
By the preceding argument, we obtain \((|s \setminus t| + |s \cap t| = |s|) = (|s| = |s|)\).
By definition, we obtain \(\left|s\right| = \left|s\right|\).
Substituting along the equality, we obtain \(\left|s \setminus t\right| + \left|s \cap t\right| = \left|s\right|\).
Substituting along the equality, we obtain \(\left|s \cap t\right| + \left|s \setminus t\right| = \left|s\right|\).
\end{proof}

\begin{theorem}[{Pigeonhole principle for two finsets inside an ambient finset}]\label{theorem:100}
Let \(\alpha\) be a type. Let \(s\), \(t\), and \(u\) be finite subsets of \(\alpha\). Assume equality on \(\alpha\) is decidable. Assume \(t \subseteq s\). Assume \(u \subseteq s\). Assume \(\left|s\right| < \left|t\right| + \left|u\right|\). Then \(t \cap u\ne\varnothing\).
\end{theorem}

\begin{proof}
Assume \(\neg \left(t \cap u\ne\varnothing\right)\).
Rewriting with \paperlink{theorem:83}{Lemma~\ref{theorem:83}}, \paperlink{premise:4966662e6d70}{Prerequisite~P\ref{premise:4966662e6d70}} (using \(\neg \left(t \cap u\ne\varnothing\right)\)) and \paperlink{premise:46696e7365742e6e6f745f6e6f6e656d7074795f6966665f65715f656d707479}{Prerequisite~P\ref{premise:46696e7365742e6e6f745f6e6f6e656d7074795f6966665f65715f656d707479}} gives \(\left|t\right| + \left|u\right| = \left|t \cup u\right|\).
Use the hypothesis \(t \subseteq s\).
Use the hypothesis \(u \subseteq s\).
By \paperlink{premise:46696e7365742e756e696f6e5f737562736574}{Prerequisite~P\ref{premise:46696e7365742e756e696f6e5f737562736574}}, we obtain \(t \cup u \subseteq s\).
By \paperlink{theorem:5}{Lemma~\ref{theorem:5}}, we obtain \(\left|t \cup u\right| \le \left|s\right|\).
By transitivity of the displayed relations, we obtain \(\left|t\right| + \left|u\right| \le \left|s\right|\).
Rewriting with \paperlink{premise:696d706c6965735f636f6e6772}{Prerequisite~P\ref{premise:696d706c6965735f636f6e6772}} and \paperlink{premise:4d6174686c69622e5461637469632e507573684e65672e6e6f745f6c745f6571}{Prerequisite~P\ref{premise:4d6174686c69622e5461637469632e507573684e65672e6e6f745f6c745f6571}} in the preceding result gives \(\neg (t \cap u\ne\varnothing) \implies \neg (|s| < |t| + |u|)\).
Use the hypothesis \(\left|s\right| < \left|t\right| + \left|u\right|\).
By \paperlink{premise:4d6174686c69622e5461637469632e436f6e747261706f73652e6d7472}{Prerequisite~P\ref{premise:4d6174686c69622e5461637469632e436f6e747261706f73652e6d7472}}, we obtain \(t \cap u\ne\varnothing\).
\end{proof}

\section{Finite sets}
\begin{theorem}\label{theorem:101}
Let \(\alpha\) be a type. Let \(s\) be a finite subset of \(\alpha\). Let \(p\) be a function of type \(\alpha \to \mathsf{Prop}\). Assume the predicate \(p\) is decidable. Fix a structure \(\mathit{inst}_{1}\) of type \(\Pi_{x\, :\, \alpha}\, \operatorname{Decidable}\left(\neg p\left(x\right)\right)\). Then
\begin{center}
\(\displaystyle |\{x\in s \mid p(x)\}| + |\{a\in s \mid \neg p(a)\}| = |s|.\)
\end{center}
\end{theorem}

\begin{proof}
By \paperlink{premise:46696e7365742e66696c7465725f756e696f6e5f66696c7465725f6e65675f6571}{Prerequisite~P\ref{premise:46696e7365742e66696c7465725f756e696f6e5f66696c7465725f6e65675f6571}}, we have \(\{x\in s \mid p(x)\} \cup \{a\in s \mid \neg p\left(a\right)\} = s\).
By \paperlink{premise:46696e7365742e6469736a6f696e745f66696c7465725f66696c7465725f6e6567}{Prerequisite~P\ref{premise:46696e7365742e6469736a6f696e745f66696c7465725f66696c7465725f6e6567}}, we have \(\operatorname{Disjoint}\left(\{x\in s \mid p(x)\}, \{a\in s \mid \neg p\left(a\right)\}\right)\).
By \paperlink{theorem:89}{Lemma~\ref{theorem:89}} applied to \(\operatorname{Disjoint}\left(\{x\in s \mid p(x)\}, \{a\in s \mid \neg p\left(a\right)\}\right)\), we have \(|\{x\in s \mid p(x)\} \cup \{a\in s \mid \neg p(a)\}| = |\{x\in s \mid p(x)\}| + |\{a\in s \mid \neg p(a)\}|\).
Starting from \paperlink{premise:45712e73796d6d}{Prerequisite~P\ref{premise:45712e73796d6d}} applied to \(|\{x\in s \mid p(x)\} \cup \{a\in s \mid \neg p(a)\}| = |\{x\in s \mid p(x)\}| + |\{a\in s \mid \neg p(a)\}|\), substitution using \paperlink{premise:46696e7365742e66696c7465725f756e696f6e5f66696c7465725f6e65675f6571}{Prerequisite~P\ref{premise:46696e7365742e66696c7465725f756e696f6e5f66696c7465725f6e65675f6571}} gives \(\left|\{x\in s \mid p(x)\}\right| + \left|\{a\in s \mid \neg p\left(a\right)\}\right| = \left|s\right|\).
\end{proof}

\begin{lemma}\label{theorem:102}
Let \(\alpha\) be a type. Let \(s\) and \(t\) be finite subsets of \(\alpha\). Let \(n\) be a natural number. Assume \(s \subseteq t\). Assume \(\left|s\right| \le n\). Assume \(n \le \left|t\right|\). Then
\begin{center}
\(\displaystyle \exists u\, :\, \mathcal{P}_{\mathrm{fin}}(\alpha),\; s \subseteq u \land (u \subseteq t \land |u| = n).\)
\end{center}
\end{lemma}

\begin{proof}
Start at \(\left|t\right|\) and work downward to \(n\). Each step uses the result for the next larger integer.
Use the hypothesis \(n \le \left|t\right|\).
Choose the witness \(t\).
Rewriting with \(s \subseteq t\), \paperlink{premise:7375627365745f7265666c}{Prerequisite~P\ref{premise:7375627365745f7265666c}}, \paperlink{premise:46696e7365742e696e737449735265666c537562736574}{Prerequisite~P\ref{premise:46696e7365742e696e737449735265666c537562736574}} and \paperlink{premise:616e645f73656c66}{Prerequisite~P\ref{premise:616e645f73656c66}} gives \(s \subseteq t \land \left(t \subseteq t \land \left|t\right| = \left|t\right|\right)\).
This witness proves existence. Hence \(\exists u\, :\, \mathcal{P}_{\mathrm{fin}}(\alpha),\; s \subseteq u \land (u \subseteq t \land |u| = |t|)\).
Consider the following case.
\begin{proofstep}
Fix \(k\) in \(\mathbb{N}\).
Assume \(k < \left|t\right|\).
Assume \(n \le k\).
Assume \(\exists u\, :\, \mathcal{P}_{\mathrm{fin}}(\alpha),\; s \subseteq u \land (u \subseteq t \land |u| = k + 1)\).
Use the hypothesis \(\exists u\, :\, \mathcal{P}_{\mathrm{fin}}(\alpha),\; s \subseteq u \land (u \subseteq t \land |u| = k + 1)\).
Take the witnesses from the preceding result.
\begin{proofstep}
Fix \(u\) in \(\mathcal{P}_{\mathrm{fin}}(\alpha)\).
Assume \(s \subseteq u\).
Assume \(u \subseteq t\).
Assume \(\left|u\right| = k + 1\).
Consider the following case.
\begin{proofstep}
Assume \(\neg \left(0 < \left|u\right| - \left|s\right|\right)\).
By \paperlink{premise:4c65616e2e4f6d6567612e496e742e6f664e61745f7375625f646963686f746f6d79}{Prerequisite~P\ref{premise:4c65616e2e4f6d6567612e496e742e6f664e61745f7375625f646963686f746f6d79}}, we obtain \(|s| \le |u| \land |u| - |s| = |u| - |s| \lor |u| < |s| \land |u| - |s| = 0\).
Consider the following case.
\begin{proofstep}
Assume \(\left|s\right| \le \left|u\right| \land \left|u\right| - \left|s\right| = \left|u\right| - \left|s\right|\).
By definition, we obtain \(\operatorname{decide}(\operatorname{is\,Impossible}(\operatorname{Lean\mathord{:}Omega\mathord{:}Constraint\mathord{:}mk}(\operatorname{some}(1),\allowbreak{} \operatorname{some}(0))) = \mathrm{true}) = \mathrm{true}\).
By \paperlink{premise:6f665f6465636964655f65715f74727565}{Prerequisite~P\ref{premise:6f665f6465636964655f65715f74727565}}, we obtain \(\operatorname{is\,Impossible}(\operatorname{Lean\mathord{:}Omega\mathord{:}Constraint\mathord{:}mk}(\operatorname{some}(1),\allowbreak{} \operatorname{some}(0))) = \mathrm{true}\).
By \paperlink{premise:496e742e6f664e61745f6e6f6e6e6567}{Prerequisite~P\ref{premise:496e742e6f664e61745f6e6f6e6e6567}} with \(\left|u\right| - \left|s\right|\) in place of its parameter \(n\), we obtain \(0 \le \left|u\right| - \left|s\right|\).
Using \paperlink{premise:4c65616e2e4f6d6567612e4c696e656172436f6d626f2e636f6f7264696e6174655f6576616c5f35}{Prerequisite~P\ref{premise:4c65616e2e4f6d6567612e4c696e656172436f6d626f2e636f6f7264696e6174655f6576616c5f35}} in reverse gives \(|u| - |s| = \operatorname{eval}(\operatorname{coordinate}(5),\allowbreak{} \operatorname{Lean\mathord{:}Omega\mathord{:}Coeffs\mathord{:}of\,List}(\langle n\rangle \mathbin{\Vert} (\langle |s|\rangle \mathbin{\Vert} (\langle |t|\rangle \mathbin{\Vert} (\langle k\rangle \mathbin{\Vert} (\langle |u|\rangle \mathbin{\Vert} \langle |u| - |s|\rangle))))))\).
By \paperlink{premise:6c655f6f665f6c655f6f665f6571}{Prerequisite~P\ref{premise:6c655f6f665f6c655f6f665f6571}} with \(\mathbb{Z}\) as its type parameter \(\alpha\), we obtain \(0 \le \operatorname{eval}(\operatorname{coordinate}(5),\allowbreak{} \operatorname{Lean\mathord{:}Omega\mathord{:}Coeffs\mathord{:}of\,List}(\langle n\rangle \mathbin{\Vert} (\langle |s|\rangle \mathbin{\Vert} (\langle |t|\rangle \mathbin{\Vert} (\langle k\rangle \mathbin{\Vert} (\langle |u|\rangle \mathbin{\Vert} \langle |u| - |s|\rangle))))))\).
By \paperlink{premise:4c65616e2e4f6d6567612e436f6e73747261696e742e616464496e657175616c6974795f736174}{Prerequisite~P\ref{premise:4c65616e2e4f6d6567612e436f6e73747261696e742e616464496e657175616c6974795f736174}}, we obtain \(\operatorname{sat^{\prime}}(\operatorname{Lean\mathord{:}Omega\mathord{:}Constraint\mathord{:}mk}(\operatorname{some}(-0),\allowbreak{} \operatorname{none}),\allowbreak{} \langle 0\rangle \mathbin{\Vert} (\langle 0\rangle \mathbin{\Vert} (\langle 0\rangle \mathbin{\Vert} (\langle 0\rangle \mathbin{\Vert} (\langle 0\rangle \mathbin{\Vert} \langle 1\rangle)))),\allowbreak{} \operatorname{Lean\mathord{:}Omega\mathord{:}Coeffs\mathord{:}of\,List}(\langle n\rangle \mathbin{\Vert} (\langle |s|\rangle \mathbin{\Vert} (\langle |t|\rangle \mathbin{\Vert} (\langle k\rangle \mathbin{\Vert} (\langle |u|\rangle \mathbin{\Vert} \langle |u| - |s|\rangle)))))) = \mathrm{true}\).
Use the hypothesis \(\neg \left(0 < \left|u\right| - \left|s\right|\right)\).
By \paperlink{premise:4e61742e6c655f6f665f6e6f745f6c74}{Prerequisite~P\ref{premise:4e61742e6c655f6f665f6e6f745f6c74}}, we obtain \(\left|u\right| - \left|s\right| \le 0\).
By \paperlink{premise:4c65616e2e4f6d6567612e496e742e6f664e61745f6c655f6f665f6c65}{Prerequisite~P\ref{premise:4c65616e2e4f6d6567612e496e742e6f664e61745f6c655f6f665f6c65}}, we obtain \(\left|u\right| - \left|s\right| \le 0\).
By \paperlink{premise:496e742e7375625f6e6f6e6e65675f6f665f6c65}{Prerequisite~P\ref{premise:496e742e7375625f6e6f6e6e65675f6f665f6c65}}, we obtain \(0 \le 0 - \left(\left|u\right| - \left|s\right|\right)\).
By definition, we obtain \(0 = \operatorname{eval}(\operatorname{Lean\mathord{:}Omega\mathord{:}Linear\,Combo\mathord{:}mk}(0,\allowbreak{} \langle\,\rangle),\allowbreak{} \operatorname{Lean\mathord{:}Omega\mathord{:}Coeffs\mathord{:}of\,List}(\langle n\rangle \mathbin{\Vert} (\langle |s|\rangle \mathbin{\Vert} (\langle |t|\rangle \mathbin{\Vert} (\langle k\rangle \mathbin{\Vert} (\langle |u|\rangle \mathbin{\Vert} \langle |u| - |s|\rangle))))))\).
Using \paperlink{premise:4c65616e2e4f6d6567612e4c696e656172436f6d626f2e636f6f7264696e6174655f6576616c5f35}{Prerequisite~P\ref{premise:4c65616e2e4f6d6567612e4c696e656172436f6d626f2e636f6f7264696e6174655f6576616c5f35}} in reverse gives \(|u| - |s| = \operatorname{eval}(\operatorname{coordinate}(5),\allowbreak{} \operatorname{Lean\mathord{:}Omega\mathord{:}Coeffs\mathord{:}of\,List}(\langle n\rangle \mathbin{\Vert} (\langle |s|\rangle \mathbin{\Vert} (\langle |t|\rangle \mathbin{\Vert} (\langle k\rangle \mathbin{\Vert} (\langle |u|\rangle \mathbin{\Vert} \langle |u| - |s|\rangle))))))\).
By \paperlink{premise:4c65616e2e4f6d6567612e496e742e7375625f636f6e6772}{Prerequisite~P\ref{premise:4c65616e2e4f6d6567612e496e742e7375625f636f6e6772}}, we obtain \(0 - (|u| - |s|) = \operatorname{eval}(\operatorname{Lean\mathord{:}Omega\mathord{:}Linear\,Combo\mathord{:}mk}(0,\allowbreak{} \langle\,\rangle),\allowbreak{} \operatorname{Lean\mathord{:}Omega\mathord{:}Coeffs\mathord{:}of\,List}(\langle n\rangle \mathbin{\Vert} (\langle |s|\rangle \mathbin{\Vert} (\langle |t|\rangle \mathbin{\Vert} (\langle k\rangle \mathbin{\Vert} (\langle |u|\rangle \mathbin{\Vert} \langle |u| - |s|\rangle)))))) - \operatorname{eval}(\operatorname{coordinate}(5),\allowbreak{} \operatorname{Lean\mathord{:}Omega\mathord{:}Coeffs\mathord{:}of\,List}(\langle n\rangle \mathbin{\Vert} (\langle |s|\rangle \mathbin{\Vert} (\langle |t|\rangle \mathbin{\Vert} (\langle k\rangle \mathbin{\Vert} (\langle |u|\rangle \mathbin{\Vert} \langle |u| - |s|\rangle))))))\).
Using \paperlink{premise:4c65616e2e4f6d6567612e4c696e656172436f6d626f2e7375625f6576616c}{Prerequisite~P\ref{premise:4c65616e2e4f6d6567612e4c696e656172436f6d626f2e7375625f6576616c}} with \(l_{1} = \operatorname{Lean\mathord{:}Omega\mathord{:}Linear\,Combo\mathord{:}mk}(0,\allowbreak{} \langle\,\rangle)\) and \(l_{2} = \operatorname{Lean\mathord{:}Omega\mathord{:}Linear\,Combo\mathord{:}mk}(0,\allowbreak{} \langle 0\rangle \mathbin{\Vert} (\langle 0\rangle \mathbin{\Vert} (\langle 0\rangle \mathbin{\Vert} (\langle 0\rangle \mathbin{\Vert} (\langle 0\rangle \mathbin{\Vert} \langle 1\rangle)))))\) and \(v = \operatorname{Lean\mathord{:}Omega\mathord{:}Coeffs\mathord{:}of\,List}(\langle n\rangle \mathbin{\Vert} (\langle |s|\rangle \mathbin{\Vert} (\langle |t|\rangle \mathbin{\Vert} (\langle k\rangle \mathbin{\Vert} (\langle |u|\rangle \mathbin{\Vert} \langle |u| - |s|\rangle)))))\) in reverse gives \(\operatorname{eval}(\operatorname{Lean\mathord{:}Omega\mathord{:}Linear\,Combo\mathord{:}mk}(0,\allowbreak{} \langle\,\rangle),\allowbreak{} \operatorname{Lean\mathord{:}Omega\mathord{:}Coeffs\mathord{:}of\,List}(\langle n\rangle \mathbin{\Vert} (\langle |s|\rangle \mathbin{\Vert} (\langle |t|\rangle \mathbin{\Vert} (\langle k\rangle \mathbin{\Vert} (\langle |u|\rangle \mathbin{\Vert} \langle |u| - |s|\rangle)))))) - \operatorname{eval}(\operatorname{Lean\mathord{:}Omega\mathord{:}Linear\,Combo\mathord{:}mk}(0,\allowbreak{} \langle 0\rangle \mathbin{\Vert} (\langle 0\rangle \mathbin{\Vert} (\langle 0\rangle \mathbin{\Vert} (\langle 0\rangle \mathbin{\Vert} (\langle 0\rangle \mathbin{\Vert} \langle 1\rangle))))),\allowbreak{} \operatorname{Lean\mathord{:}Omega\mathord{:}Coeffs\mathord{:}of\,List}(\langle n\rangle \mathbin{\Vert} (\langle |s|\rangle \mathbin{\Vert} (\langle |t|\rangle \mathbin{\Vert} (\langle k\rangle \mathbin{\Vert} (\langle |u|\rangle \mathbin{\Vert} \langle |u| - |s|\rangle)))))) = \operatorname{eval}(\operatorname{Lean\mathord{:}Omega\mathord{:}Linear\,Combo\mathord{:}mk}(0,\allowbreak{} \langle\,\rangle) - \operatorname{Lean\mathord{:}Omega\mathord{:}Linear\,Combo\mathord{:}mk}(0,\allowbreak{} \langle 0\rangle \mathbin{\Vert} (\langle 0\rangle \mathbin{\Vert} (\langle 0\rangle \mathbin{\Vert} (\langle 0\rangle \mathbin{\Vert} (\langle 0\rangle \mathbin{\Vert} \langle 1\rangle))))),\allowbreak{} \operatorname{Lean\mathord{:}Omega\mathord{:}Coeffs\mathord{:}of\,List}(\langle n\rangle \mathbin{\Vert} (\langle |s|\rangle \mathbin{\Vert} (\langle |t|\rangle \mathbin{\Vert} (\langle k\rangle \mathbin{\Vert} (\langle |u|\rangle \mathbin{\Vert} \langle |u| - |s|\rangle))))))\).
Chaining the equalities, we obtain \(0 - (|u| - |s|) = \operatorname{eval}(\operatorname{Lean\mathord{:}Omega\mathord{:}Linear\,Combo\mathord{:}mk}(0,\allowbreak{} \langle\,\rangle) - \operatorname{Lean\mathord{:}Omega\mathord{:}Linear\,Combo\mathord{:}mk}(0,\allowbreak{} \langle 0\rangle \mathbin{\Vert} (\langle 0\rangle \mathbin{\Vert} (\langle 0\rangle \mathbin{\Vert} (\langle 0\rangle \mathbin{\Vert} (\langle 0\rangle \mathbin{\Vert} \langle 1\rangle))))),\allowbreak{} \operatorname{Lean\mathord{:}Omega\mathord{:}Coeffs\mathord{:}of\,List}(\langle n\rangle \mathbin{\Vert} (\langle |s|\rangle \mathbin{\Vert} (\langle |t|\rangle \mathbin{\Vert} (\langle k\rangle \mathbin{\Vert} (\langle |u|\rangle \mathbin{\Vert} \langle |u| - |s|\rangle))))))\).
By \paperlink{premise:6c655f6f665f6c655f6f665f6571}{Prerequisite~P\ref{premise:6c655f6f665f6c655f6f665f6571}} with \(\mathbb{Z}\) as its type parameter \(\alpha\), we obtain \(0 \le \operatorname{eval}(\operatorname{Lean\mathord{:}Omega\mathord{:}Linear\,Combo\mathord{:}mk}(0,\allowbreak{} \langle\,\rangle) - \operatorname{Lean\mathord{:}Omega\mathord{:}Linear\,Combo\mathord{:}mk}(0,\allowbreak{} \langle 0\rangle \mathbin{\Vert} (\langle 0\rangle \mathbin{\Vert} (\langle 0\rangle \mathbin{\Vert} (\langle 0\rangle \mathbin{\Vert} (\langle 0\rangle \mathbin{\Vert} \langle 1\rangle))))),\allowbreak{} \operatorname{Lean\mathord{:}Omega\mathord{:}Coeffs\mathord{:}of\,List}(\langle n\rangle \mathbin{\Vert} (\langle |s|\rangle \mathbin{\Vert} (\langle |t|\rangle \mathbin{\Vert} (\langle k\rangle \mathbin{\Vert} (\langle |u|\rangle \mathbin{\Vert} \langle |u| - |s|\rangle))))))\).
By \paperlink{premise:4c65616e2e4f6d6567612e436f6e73747261696e742e616464496e657175616c6974795f736174}{Prerequisite~P\ref{premise:4c65616e2e4f6d6567612e436f6e73747261696e742e616464496e657175616c6974795f736174}}, we obtain \(\operatorname{sat^{\prime}}(\operatorname{Lean\mathord{:}Omega\mathord{:}Constraint\mathord{:}mk}(\operatorname{some}(-0),\allowbreak{} \operatorname{none}),\allowbreak{} \langle 0\rangle \mathbin{\Vert} (\langle 0\rangle \mathbin{\Vert} (\langle 0\rangle \mathbin{\Vert} (\langle 0\rangle \mathbin{\Vert} (\langle 0\rangle \mathbin{\Vert} \langle -1\rangle)))),\allowbreak{} \operatorname{Lean\mathord{:}Omega\mathord{:}Coeffs\mathord{:}of\,List}(\langle n\rangle \mathbin{\Vert} (\langle |s|\rangle \mathbin{\Vert} (\langle |t|\rangle \mathbin{\Vert} (\langle k\rangle \mathbin{\Vert} (\langle |u|\rangle \mathbin{\Vert} \langle |u| - |s|\rangle)))))) = \mathrm{true}\).
By \paperlink{premise:4c65616e2e4f6d6567612e746964795f736174}{Prerequisite~P\ref{premise:4c65616e2e4f6d6567612e746964795f736174}}, we obtain \(\operatorname{sat^{\prime}}(\operatorname{tidy\,Constraint}(\operatorname{Lean\mathord{:}Omega\mathord{:}Constraint\mathord{:}mk}(\operatorname{some}(0),\allowbreak{} \operatorname{none}),\allowbreak{} \langle 0\rangle \mathbin{\Vert} (\langle 0\rangle \mathbin{\Vert} (\langle 0\rangle \mathbin{\Vert} (\langle 0\rangle \mathbin{\Vert} (\langle 0\rangle \mathbin{\Vert} \langle -1\rangle))))),\allowbreak{} \operatorname{tidy\,Coeffs}(\operatorname{Lean\mathord{:}Omega\mathord{:}Constraint\mathord{:}mk}(\operatorname{some}(0),\allowbreak{} \operatorname{none}),\allowbreak{} \langle 0\rangle \mathbin{\Vert} (\langle 0\rangle \mathbin{\Vert} (\langle 0\rangle \mathbin{\Vert} (\langle 0\rangle \mathbin{\Vert} (\langle 0\rangle \mathbin{\Vert} \langle -1\rangle))))),\allowbreak{} \operatorname{Lean\mathord{:}Omega\mathord{:}Coeffs\mathord{:}of\,List}(\langle n\rangle \mathbin{\Vert} (\langle |s|\rangle \mathbin{\Vert} (\langle |t|\rangle \mathbin{\Vert} (\langle k\rangle \mathbin{\Vert} (\langle |u|\rangle \mathbin{\Vert} \langle |u| - |s|\rangle)))))) = \mathrm{true}\).
By \paperlink{premise:4c65616e2e4f6d6567612e436f6e73747261696e742e636f6d62696e655f73617427}{Prerequisite~P\ref{premise:4c65616e2e4f6d6567612e436f6e73747261696e742e636f6d62696e655f73617427}}, we obtain \(\operatorname{sat^{\prime}}(\operatorname{combine}(\operatorname{Lean\mathord{:}Omega\mathord{:}Constraint\mathord{:}mk}(\operatorname{some}(0),\allowbreak{} \operatorname{none}),\allowbreak{} \operatorname{Lean\mathord{:}Omega\mathord{:}Constraint\mathord{:}mk}(\operatorname{none},\allowbreak{} \operatorname{some}(0))),\allowbreak{} \langle 0\rangle \mathbin{\Vert} (\langle 0\rangle \mathbin{\Vert} (\langle 0\rangle \mathbin{\Vert} (\langle 0\rangle \mathbin{\Vert} (\langle 0\rangle \mathbin{\Vert} \langle 1\rangle)))),\allowbreak{} \operatorname{Lean\mathord{:}Omega\mathord{:}Coeffs\mathord{:}of\,List}(\langle n\rangle \mathbin{\Vert} (\langle |s|\rangle \mathbin{\Vert} (\langle |t|\rangle \mathbin{\Vert} (\langle k\rangle \mathbin{\Vert} (\langle |u|\rangle \mathbin{\Vert} \langle |u| - |s|\rangle)))))) = \mathrm{true}\).
Using \paperlink{premise:4c65616e2e4f6d6567612e4c696e656172436f6d626f2e636f6f7264696e6174655f6576616c5f35}{Prerequisite~P\ref{premise:4c65616e2e4f6d6567612e4c696e656172436f6d626f2e636f6f7264696e6174655f6576616c5f35}} in reverse gives \(|u| - |s| = \operatorname{eval}(\operatorname{coordinate}(5),\allowbreak{} \operatorname{Lean\mathord{:}Omega\mathord{:}Coeffs\mathord{:}of\,List}(\langle n\rangle \mathbin{\Vert} (\langle |s|\rangle \mathbin{\Vert} (\langle |t|\rangle \mathbin{\Vert} (\langle k\rangle \mathbin{\Vert} (\langle |u|\rangle \mathbin{\Vert} \langle |u| - |s|\rangle))))))\).
Using \paperlink{premise:4c65616e2e4f6d6567612e4c696e656172436f6d626f2e636f6f7264696e6174655f6576616c5f34}{Prerequisite~P\ref{premise:4c65616e2e4f6d6567612e4c696e656172436f6d626f2e636f6f7264696e6174655f6576616c5f34}} in reverse gives \(|u| = \operatorname{eval}(\operatorname{coordinate}(4),\allowbreak{} \operatorname{Lean\mathord{:}Omega\mathord{:}Coeffs\mathord{:}of\,List}(\langle n\rangle \mathbin{\Vert} (\langle |s|\rangle \mathbin{\Vert} (\langle |t|\rangle \mathbin{\Vert} (\langle k\rangle \mathbin{\Vert} (\langle |u|\rangle \mathbin{\Vert} \langle |u| - |s|\rangle))))))\).
Using \paperlink{premise:4c65616e2e4f6d6567612e4c696e656172436f6d626f2e636f6f7264696e6174655f6576616c5f31}{Prerequisite~P\ref{premise:4c65616e2e4f6d6567612e4c696e656172436f6d626f2e636f6f7264696e6174655f6576616c5f31}} in reverse gives \(|s| = \operatorname{eval}(\operatorname{coordinate}(1),\allowbreak{} \operatorname{Lean\mathord{:}Omega\mathord{:}Coeffs\mathord{:}of\,List}(\langle n\rangle \mathbin{\Vert} (\langle |s|\rangle \mathbin{\Vert} (\langle |t|\rangle \mathbin{\Vert} (\langle k\rangle \mathbin{\Vert} (\langle |u|\rangle \mathbin{\Vert} \langle |u| - |s|\rangle))))))\).
By \paperlink{premise:4c65616e2e4f6d6567612e496e742e7375625f636f6e6772}{Prerequisite~P\ref{premise:4c65616e2e4f6d6567612e496e742e7375625f636f6e6772}}, we obtain \(|u| - |s| = \operatorname{eval}(\operatorname{coordinate}(4),\allowbreak{} \operatorname{Lean\mathord{:}Omega\mathord{:}Coeffs\mathord{:}of\,List}(\langle n\rangle \mathbin{\Vert} (\langle |s|\rangle \mathbin{\Vert} (\langle |t|\rangle \mathbin{\Vert} (\langle k\rangle \mathbin{\Vert} (\langle |u|\rangle \mathbin{\Vert} \langle |u| - |s|\rangle)))))) - \operatorname{eval}(\operatorname{coordinate}(1),\allowbreak{} \operatorname{Lean\mathord{:}Omega\mathord{:}Coeffs\mathord{:}of\,List}(\langle n\rangle \mathbin{\Vert} (\langle |s|\rangle \mathbin{\Vert} (\langle |t|\rangle \mathbin{\Vert} (\langle k\rangle \mathbin{\Vert} (\langle |u|\rangle \mathbin{\Vert} \langle |u| - |s|\rangle))))))\).
Using \paperlink{premise:4c65616e2e4f6d6567612e4c696e656172436f6d626f2e7375625f6576616c}{Prerequisite~P\ref{premise:4c65616e2e4f6d6567612e4c696e656172436f6d626f2e7375625f6576616c}} with \(l_{1} = \operatorname{Lean\mathord{:}Omega\mathord{:}Linear\,Combo\mathord{:}mk}(0,\allowbreak{} \langle 0\rangle \mathbin{\Vert} (\langle 0\rangle \mathbin{\Vert} (\langle 0\rangle \mathbin{\Vert} (\langle 0\rangle \mathbin{\Vert} \langle 1\rangle))))\) and \(l_{2} = \operatorname{Lean\mathord{:}Omega\mathord{:}Linear\,Combo\mathord{:}mk}(0,\allowbreak{} \langle 0\rangle \mathbin{\Vert} \langle 1\rangle)\) and \(v = \operatorname{Lean\mathord{:}Omega\mathord{:}Coeffs\mathord{:}of\,List}(\langle n\rangle \mathbin{\Vert} (\langle |s|\rangle \mathbin{\Vert} (\langle |t|\rangle \mathbin{\Vert} (\langle k\rangle \mathbin{\Vert} (\langle |u|\rangle \mathbin{\Vert} \langle |u| - |s|\rangle)))))\) in reverse gives \(\operatorname{eval}(\operatorname{Lean\mathord{:}Omega\mathord{:}Linear\,Combo\mathord{:}mk}(0,\allowbreak{} \langle 0\rangle \mathbin{\Vert} (\langle 0\rangle \mathbin{\Vert} (\langle 0\rangle \mathbin{\Vert} (\langle 0\rangle \mathbin{\Vert} \langle 1\rangle)))),\allowbreak{} \operatorname{Lean\mathord{:}Omega\mathord{:}Coeffs\mathord{:}of\,List}(\langle n\rangle \mathbin{\Vert} (\langle |s|\rangle \mathbin{\Vert} (\langle |t|\rangle \mathbin{\Vert} (\langle k\rangle \mathbin{\Vert} (\langle |u|\rangle \mathbin{\Vert} \langle |u| - |s|\rangle)))))) - \operatorname{eval}(\operatorname{Lean\mathord{:}Omega\mathord{:}Linear\,Combo\mathord{:}mk}(0,\allowbreak{} \langle 0\rangle \mathbin{\Vert} \langle 1\rangle),\allowbreak{} \operatorname{Lean\mathord{:}Omega\mathord{:}Coeffs\mathord{:}of\,List}(\langle n\rangle \mathbin{\Vert} (\langle |s|\rangle \mathbin{\Vert} (\langle |t|\rangle \mathbin{\Vert} (\langle k\rangle \mathbin{\Vert} (\langle |u|\rangle \mathbin{\Vert} \langle |u| - |s|\rangle)))))) = \operatorname{eval}(\operatorname{Lean\mathord{:}Omega\mathord{:}Linear\,Combo\mathord{:}mk}(0,\allowbreak{} \langle 0\rangle \mathbin{\Vert} (\langle 0\rangle \mathbin{\Vert} (\langle 0\rangle \mathbin{\Vert} (\langle 0\rangle \mathbin{\Vert} \langle 1\rangle)))) - \operatorname{Lean\mathord{:}Omega\mathord{:}Linear\,Combo\mathord{:}mk}(0,\allowbreak{} \langle 0\rangle \mathbin{\Vert} \langle 1\rangle),\allowbreak{} \operatorname{Lean\mathord{:}Omega\mathord{:}Coeffs\mathord{:}of\,List}(\langle n\rangle \mathbin{\Vert} (\langle |s|\rangle \mathbin{\Vert} (\langle |t|\rangle \mathbin{\Vert} (\langle k\rangle \mathbin{\Vert} (\langle |u|\rangle \mathbin{\Vert} \langle |u| - |s|\rangle))))))\).
Chaining the equalities, we obtain \(|u| - |s| = \operatorname{eval}(\operatorname{Lean\mathord{:}Omega\mathord{:}Linear\,Combo\mathord{:}mk}(0,\allowbreak{} \langle 0\rangle \mathbin{\Vert} (\langle 0\rangle \mathbin{\Vert} (\langle 0\rangle \mathbin{\Vert} (\langle 0\rangle \mathbin{\Vert} \langle 1\rangle)))) - \operatorname{Lean\mathord{:}Omega\mathord{:}Linear\,Combo\mathord{:}mk}(0,\allowbreak{} \langle 0\rangle \mathbin{\Vert} \langle 1\rangle),\allowbreak{} \operatorname{Lean\mathord{:}Omega\mathord{:}Coeffs\mathord{:}of\,List}(\langle n\rangle \mathbin{\Vert} (\langle |s|\rangle \mathbin{\Vert} (\langle |t|\rangle \mathbin{\Vert} (\langle k\rangle \mathbin{\Vert} (\langle |u|\rangle \mathbin{\Vert} \langle |u| - |s|\rangle))))))\).
By \paperlink{premise:4c65616e2e4f6d6567612e496e742e7375625f636f6e6772}{Prerequisite~P\ref{premise:4c65616e2e4f6d6567612e496e742e7375625f636f6e6772}}, we obtain \(|u| - |s| - (|u| - |s|) = \operatorname{eval}(\operatorname{coordinate}(5),\allowbreak{} \operatorname{Lean\mathord{:}Omega\mathord{:}Coeffs\mathord{:}of\,List}(\langle n\rangle \mathbin{\Vert} (\langle |s|\rangle \mathbin{\Vert} (\langle |t|\rangle \mathbin{\Vert} (\langle k\rangle \mathbin{\Vert} (\langle |u|\rangle \mathbin{\Vert} \langle |u| - |s|\rangle)))))) - \operatorname{eval}(\operatorname{Lean\mathord{:}Omega\mathord{:}Linear\,Combo\mathord{:}mk}(0,\allowbreak{} \langle 0\rangle \mathbin{\Vert} (\langle 0\rangle \mathbin{\Vert} (\langle 0\rangle \mathbin{\Vert} (\langle 0\rangle \mathbin{\Vert} \langle 1\rangle)))) - \operatorname{Lean\mathord{:}Omega\mathord{:}Linear\,Combo\mathord{:}mk}(0,\allowbreak{} \langle 0\rangle \mathbin{\Vert} \langle 1\rangle),\allowbreak{} \operatorname{Lean\mathord{:}Omega\mathord{:}Coeffs\mathord{:}of\,List}(\langle n\rangle \mathbin{\Vert} (\langle |s|\rangle \mathbin{\Vert} (\langle |t|\rangle \mathbin{\Vert} (\langle k\rangle \mathbin{\Vert} (\langle |u|\rangle \mathbin{\Vert} \langle |u| - |s|\rangle))))))\).
Using \paperlink{premise:4c65616e2e4f6d6567612e4c696e656172436f6d626f2e7375625f6576616c}{Prerequisite~P\ref{premise:4c65616e2e4f6d6567612e4c696e656172436f6d626f2e7375625f6576616c}} with \(l_{1} = \operatorname{Lean\mathord{:}Omega\mathord{:}Linear\,Combo\mathord{:}mk}(0,\allowbreak{} \langle 0\rangle \mathbin{\Vert} (\langle 0\rangle \mathbin{\Vert} (\langle 0\rangle \mathbin{\Vert} (\langle 0\rangle \mathbin{\Vert} (\langle 0\rangle \mathbin{\Vert} \langle 1\rangle)))))\) and \(l_{2} = \operatorname{Lean\mathord{:}Omega\mathord{:}Linear\,Combo\mathord{:}mk}(0,\allowbreak{} \langle 0\rangle \mathbin{\Vert} (\langle -1\rangle \mathbin{\Vert} (\langle 0\rangle \mathbin{\Vert} (\langle 0\rangle \mathbin{\Vert} \langle 1\rangle))))\) and \(v = \operatorname{Lean\mathord{:}Omega\mathord{:}Coeffs\mathord{:}of\,List}(\langle n\rangle \mathbin{\Vert} (\langle |s|\rangle \mathbin{\Vert} (\langle |t|\rangle \mathbin{\Vert} (\langle k\rangle \mathbin{\Vert} (\langle |u|\rangle \mathbin{\Vert} \langle |u| - |s|\rangle)))))\) in reverse gives \(\operatorname{eval}(\operatorname{Lean\mathord{:}Omega\mathord{:}Linear\,Combo\mathord{:}mk}(0,\allowbreak{} \langle 0\rangle \mathbin{\Vert} (\langle 0\rangle \mathbin{\Vert} (\langle 0\rangle \mathbin{\Vert} (\langle 0\rangle \mathbin{\Vert} (\langle 0\rangle \mathbin{\Vert} \langle 1\rangle))))),\allowbreak{} \operatorname{Lean\mathord{:}Omega\mathord{:}Coeffs\mathord{:}of\,List}(\langle n\rangle \mathbin{\Vert} (\langle |s|\rangle \mathbin{\Vert} (\langle |t|\rangle \mathbin{\Vert} (\langle k\rangle \mathbin{\Vert} (\langle |u|\rangle \mathbin{\Vert} \langle |u| - |s|\rangle)))))) - \operatorname{eval}(\operatorname{Lean\mathord{:}Omega\mathord{:}Linear\,Combo\mathord{:}mk}(0,\allowbreak{} \langle 0\rangle \mathbin{\Vert} (\langle -1\rangle \mathbin{\Vert} (\langle 0\rangle \mathbin{\Vert} (\langle 0\rangle \mathbin{\Vert} \langle 1\rangle)))),\allowbreak{} \operatorname{Lean\mathord{:}Omega\mathord{:}Coeffs\mathord{:}of\,List}(\langle n\rangle \mathbin{\Vert} (\langle |s|\rangle \mathbin{\Vert} (\langle |t|\rangle \mathbin{\Vert} (\langle k\rangle \mathbin{\Vert} (\langle |u|\rangle \mathbin{\Vert} \langle |u| - |s|\rangle)))))) = \operatorname{eval}(\operatorname{Lean\mathord{:}Omega\mathord{:}Linear\,Combo\mathord{:}mk}(0,\allowbreak{} \langle 0\rangle \mathbin{\Vert} (\langle 0\rangle \mathbin{\Vert} (\langle 0\rangle \mathbin{\Vert} (\langle 0\rangle \mathbin{\Vert} (\langle 0\rangle \mathbin{\Vert} \langle 1\rangle))))) - \operatorname{Lean\mathord{:}Omega\mathord{:}Linear\,Combo\mathord{:}mk}(0,\allowbreak{} \langle 0\rangle \mathbin{\Vert} (\langle -1\rangle \mathbin{\Vert} (\langle 0\rangle \mathbin{\Vert} (\langle 0\rangle \mathbin{\Vert} \langle 1\rangle)))),\allowbreak{} \operatorname{Lean\mathord{:}Omega\mathord{:}Coeffs\mathord{:}of\,List}(\langle n\rangle \mathbin{\Vert} (\langle |s|\rangle \mathbin{\Vert} (\langle |t|\rangle \mathbin{\Vert} (\langle k\rangle \mathbin{\Vert} (\langle |u|\rangle \mathbin{\Vert} \langle |u| - |s|\rangle))))))\).
Chaining the equalities, we obtain \(|u| - |s| - (|u| - |s|) = \operatorname{eval}(\operatorname{Lean\mathord{:}Omega\mathord{:}Linear\,Combo\mathord{:}mk}(0,\allowbreak{} \langle 0\rangle \mathbin{\Vert} (\langle 0\rangle \mathbin{\Vert} (\langle 0\rangle \mathbin{\Vert} (\langle 0\rangle \mathbin{\Vert} (\langle 0\rangle \mathbin{\Vert} \langle 1\rangle))))) - \operatorname{Lean\mathord{:}Omega\mathord{:}Linear\,Combo\mathord{:}mk}(0,\allowbreak{} \langle 0\rangle \mathbin{\Vert} (\langle -1\rangle \mathbin{\Vert} (\langle 0\rangle \mathbin{\Vert} (\langle 0\rangle \mathbin{\Vert} \langle 1\rangle)))),\allowbreak{} \operatorname{Lean\mathord{:}Omega\mathord{:}Coeffs\mathord{:}of\,List}(\langle n\rangle \mathbin{\Vert} (\langle |s|\rangle \mathbin{\Vert} (\langle |t|\rangle \mathbin{\Vert} (\langle k\rangle \mathbin{\Vert} (\langle |u|\rangle \mathbin{\Vert} \langle |u| - |s|\rangle))))))\).
Reversing the equality, we obtain \(\operatorname{eval}(\operatorname{Lean\mathord{:}Omega\mathord{:}Linear\,Combo\mathord{:}mk}(0,\allowbreak{} \langle 0\rangle \mathbin{\Vert} (\langle 0\rangle \mathbin{\Vert} (\langle 0\rangle \mathbin{\Vert} (\langle 0\rangle \mathbin{\Vert} (\langle 0\rangle \mathbin{\Vert} \langle 1\rangle))))) - \operatorname{Lean\mathord{:}Omega\mathord{:}Linear\,Combo\mathord{:}mk}(0,\allowbreak{} \langle 0\rangle \mathbin{\Vert} (\langle -1\rangle \mathbin{\Vert} (\langle 0\rangle \mathbin{\Vert} (\langle 0\rangle \mathbin{\Vert} \langle 1\rangle)))),\allowbreak{} \operatorname{Lean\mathord{:}Omega\mathord{:}Coeffs\mathord{:}of\,List}(\langle n\rangle \mathbin{\Vert} (\langle |s|\rangle \mathbin{\Vert} (\langle |t|\rangle \mathbin{\Vert} (\langle k\rangle \mathbin{\Vert} (\langle |u|\rangle \mathbin{\Vert} \langle |u| - |s|\rangle)))))) = |u| - |s| - (|u| - |s|)\).
Use the hypothesis \(\left|s\right| \le \left|u\right| \land \left|u\right| - \left|s\right| = \left|u\right| - \left|s\right|\).
Taking the second conjunct, we obtain \(\left|u\right| - \left|s\right| = \left|u\right| - \left|s\right|\).
By \paperlink{premise:496e742e7375625f65715f7a65726f5f6f665f6571}{Prerequisite~P\ref{premise:496e742e7375625f65715f7a65726f5f6f665f6571}}, we obtain \(\left|u\right| - \left|s\right| - \left(\left|u\right| - \left|s\right|\right) = 0\).
Chaining the equalities, we obtain \(\operatorname{eval}(\operatorname{Lean\mathord{:}Omega\mathord{:}Linear\,Combo\mathord{:}mk}(0,\allowbreak{} \langle 0\rangle \mathbin{\Vert} (\langle 0\rangle \mathbin{\Vert} (\langle 0\rangle \mathbin{\Vert} (\langle 0\rangle \mathbin{\Vert} (\langle 0\rangle \mathbin{\Vert} \langle 1\rangle))))) - \operatorname{Lean\mathord{:}Omega\mathord{:}Linear\,Combo\mathord{:}mk}(0,\allowbreak{} \langle 0\rangle \mathbin{\Vert} (\langle -1\rangle \mathbin{\Vert} (\langle 0\rangle \mathbin{\Vert} (\langle 0\rangle \mathbin{\Vert} \langle 1\rangle)))),\allowbreak{} \operatorname{Lean\mathord{:}Omega\mathord{:}Coeffs\mathord{:}of\,List}(\langle n\rangle \mathbin{\Vert} (\langle |s|\rangle \mathbin{\Vert} (\langle |t|\rangle \mathbin{\Vert} (\langle k\rangle \mathbin{\Vert} (\langle |u|\rangle \mathbin{\Vert} \langle |u| - |s|\rangle)))))) = 0\).
By \paperlink{premise:4c65616e2e4f6d6567612e436f6e73747261696e742e616464457175616c6974795f736174}{Prerequisite~P\ref{premise:4c65616e2e4f6d6567612e436f6e73747261696e742e616464457175616c6974795f736174}}, we obtain \(\operatorname{sat^{\prime}}(\operatorname{Lean\mathord{:}Omega\mathord{:}Constraint\mathord{:}mk}(\operatorname{some}(-0),\allowbreak{} \operatorname{some}(-0)),\allowbreak{} \langle 0\rangle \mathbin{\Vert} (\langle 1\rangle \mathbin{\Vert} (\langle 0\rangle \mathbin{\Vert} (\langle 0\rangle \mathbin{\Vert} (\langle -1\rangle \mathbin{\Vert} \langle 1\rangle)))),\allowbreak{} \operatorname{Lean\mathord{:}Omega\mathord{:}Coeffs\mathord{:}of\,List}(\langle n\rangle \mathbin{\Vert} (\langle |s|\rangle \mathbin{\Vert} (\langle |t|\rangle \mathbin{\Vert} (\langle k\rangle \mathbin{\Vert} (\langle |u|\rangle \mathbin{\Vert} \langle |u| - |s|\rangle)))))) = \mathrm{true}\).
By \paperlink{premise:4c65616e2e4f6d6567612e636f6d626f5f73617427}{Prerequisite~P\ref{premise:4c65616e2e4f6d6567612e636f6d626f5f73617427}} with \(\operatorname{Lean\mathord{:}Omega\mathord{:}Constraint\mathord{:}mk}(\operatorname{some}(0),\allowbreak{} \operatorname{some}(0))\) in place of its parameter \(s\) and \(\operatorname{Lean\mathord{:}Omega\mathord{:}Constraint\mathord{:}mk}(\operatorname{some}(0),\allowbreak{} \operatorname{some}(0))\) in place of its parameter \(t\) and \(-1\) in place of its parameter \(a\) and \(x = \langle 0\rangle \mathbin{\Vert} (\langle 0\rangle \mathbin{\Vert} (\langle 0\rangle \mathbin{\Vert} (\langle 0\rangle \mathbin{\Vert} (\langle 0\rangle \mathbin{\Vert} \langle 1\rangle))))\) and \(b = 1\) and \(y = \langle 0\rangle \mathbin{\Vert} (\langle 1\rangle \mathbin{\Vert} (\langle 0\rangle \mathbin{\Vert} (\langle 0\rangle \mathbin{\Vert} (\langle -1\rangle \mathbin{\Vert} \langle 1\rangle))))\) and \(v = \operatorname{Lean\mathord{:}Omega\mathord{:}Coeffs\mathord{:}of\,List}(\langle n\rangle \mathbin{\Vert} (\langle |s|\rangle \mathbin{\Vert} (\langle |t|\rangle \mathbin{\Vert} (\langle k\rangle \mathbin{\Vert} (\langle |u|\rangle \mathbin{\Vert} \langle |u| - |s|\rangle)))))\), we obtain \(\operatorname{sat^{\prime}}(\operatorname{Lean\mathord{:}Omega\mathord{:}Constraint\mathord{:}combo}(-1,\allowbreak{} \operatorname{Lean\mathord{:}Omega\mathord{:}Constraint\mathord{:}mk}(\operatorname{some}(0),\allowbreak{} \operatorname{some}(0)),\allowbreak{} 1,\allowbreak{} \operatorname{Lean\mathord{:}Omega\mathord{:}Constraint\mathord{:}mk}(\operatorname{some}(0),\allowbreak{} \operatorname{some}(0))),\allowbreak{} \operatorname{Lean\mathord{:}Omega\mathord{:}Coeffs\mathord{:}combo}(-1,\allowbreak{} \langle 0\rangle \mathbin{\Vert} (\langle 0\rangle \mathbin{\Vert} (\langle 0\rangle \mathbin{\Vert} (\langle 0\rangle \mathbin{\Vert} (\langle 0\rangle \mathbin{\Vert} \langle 1\rangle)))),\allowbreak{} 1,\allowbreak{} \langle 0\rangle \mathbin{\Vert} (\langle 1\rangle \mathbin{\Vert} (\langle 0\rangle \mathbin{\Vert} (\langle 0\rangle \mathbin{\Vert} (\langle -1\rangle \mathbin{\Vert} \langle 1\rangle))))),\allowbreak{} \operatorname{Lean\mathord{:}Omega\mathord{:}Coeffs\mathord{:}of\,List}(\langle n\rangle \mathbin{\Vert} (\langle |s|\rangle \mathbin{\Vert} (\langle |t|\rangle \mathbin{\Vert} (\langle k\rangle \mathbin{\Vert} (\langle |u|\rangle \mathbin{\Vert} \langle |u| - |s|\rangle)))))) = \mathrm{true}\).
Use the hypothesis \(\left|s\right| \le n\).
By \paperlink{premise:4c65616e2e4f6d6567612e496e742e6f664e61745f6c655f6f665f6c65}{Prerequisite~P\ref{premise:4c65616e2e4f6d6567612e496e742e6f664e61745f6c655f6f665f6c65}}, we obtain \(\left|s\right| \le n\).
By \paperlink{premise:496e742e7375625f6e6f6e6e65675f6f665f6c65}{Prerequisite~P\ref{premise:496e742e7375625f6e6f6e6e65675f6f665f6c65}}, we obtain \(0 \le n - \left|s\right|\).
Using \paperlink{premise:4c65616e2e4f6d6567612e4c696e656172436f6d626f2e636f6f7264696e6174655f6576616c5f30}{Prerequisite~P\ref{premise:4c65616e2e4f6d6567612e4c696e656172436f6d626f2e636f6f7264696e6174655f6576616c5f30}} in reverse gives \(n = \operatorname{eval}(\operatorname{coordinate}(0),\allowbreak{} \operatorname{Lean\mathord{:}Omega\mathord{:}Coeffs\mathord{:}of\,List}(\langle n\rangle \mathbin{\Vert} (\langle |s|\rangle \mathbin{\Vert} (\langle |t|\rangle \mathbin{\Vert} (\langle k\rangle \mathbin{\Vert} (\langle |u|\rangle \mathbin{\Vert} \langle |u| - |s|\rangle))))))\).
Using \paperlink{premise:4c65616e2e4f6d6567612e4c696e656172436f6d626f2e636f6f7264696e6174655f6576616c5f31}{Prerequisite~P\ref{premise:4c65616e2e4f6d6567612e4c696e656172436f6d626f2e636f6f7264696e6174655f6576616c5f31}} in reverse gives \(|s| = \operatorname{eval}(\operatorname{coordinate}(1),\allowbreak{} \operatorname{Lean\mathord{:}Omega\mathord{:}Coeffs\mathord{:}of\,List}(\langle n\rangle \mathbin{\Vert} (\langle |s|\rangle \mathbin{\Vert} (\langle |t|\rangle \mathbin{\Vert} (\langle k\rangle \mathbin{\Vert} (\langle |u|\rangle \mathbin{\Vert} \langle |u| - |s|\rangle))))))\).
By \paperlink{premise:4c65616e2e4f6d6567612e496e742e7375625f636f6e6772}{Prerequisite~P\ref{premise:4c65616e2e4f6d6567612e496e742e7375625f636f6e6772}}, we obtain \(n - |s| = \operatorname{eval}(\operatorname{coordinate}(0),\allowbreak{} \operatorname{Lean\mathord{:}Omega\mathord{:}Coeffs\mathord{:}of\,List}(\langle n\rangle \mathbin{\Vert} (\langle |s|\rangle \mathbin{\Vert} (\langle |t|\rangle \mathbin{\Vert} (\langle k\rangle \mathbin{\Vert} (\langle |u|\rangle \mathbin{\Vert} \langle |u| - |s|\rangle)))))) - \operatorname{eval}(\operatorname{coordinate}(1),\allowbreak{} \operatorname{Lean\mathord{:}Omega\mathord{:}Coeffs\mathord{:}of\,List}(\langle n\rangle \mathbin{\Vert} (\langle |s|\rangle \mathbin{\Vert} (\langle |t|\rangle \mathbin{\Vert} (\langle k\rangle \mathbin{\Vert} (\langle |u|\rangle \mathbin{\Vert} \langle |u| - |s|\rangle))))))\).
Using \paperlink{premise:4c65616e2e4f6d6567612e4c696e656172436f6d626f2e7375625f6576616c}{Prerequisite~P\ref{premise:4c65616e2e4f6d6567612e4c696e656172436f6d626f2e7375625f6576616c}} with \(l_{1} = \operatorname{Lean\mathord{:}Omega\mathord{:}Linear\,Combo\mathord{:}mk}(0,\allowbreak{} \langle 1\rangle)\) and \(l_{2} = \operatorname{Lean\mathord{:}Omega\mathord{:}Linear\,Combo\mathord{:}mk}(0,\allowbreak{} \langle 0\rangle \mathbin{\Vert} \langle 1\rangle)\) and \(v = \operatorname{Lean\mathord{:}Omega\mathord{:}Coeffs\mathord{:}of\,List}(\langle n\rangle \mathbin{\Vert} (\langle |s|\rangle \mathbin{\Vert} (\langle |t|\rangle \mathbin{\Vert} (\langle k\rangle \mathbin{\Vert} (\langle |u|\rangle \mathbin{\Vert} \langle |u| - |s|\rangle)))))\) in reverse gives \(\operatorname{eval}(\operatorname{Lean\mathord{:}Omega\mathord{:}Linear\,Combo\mathord{:}mk}(0,\allowbreak{} \langle 1\rangle),\allowbreak{} \operatorname{Lean\mathord{:}Omega\mathord{:}Coeffs\mathord{:}of\,List}(\langle n\rangle \mathbin{\Vert} (\langle |s|\rangle \mathbin{\Vert} (\langle |t|\rangle \mathbin{\Vert} (\langle k\rangle \mathbin{\Vert} (\langle |u|\rangle \mathbin{\Vert} \langle |u| - |s|\rangle)))))) - \operatorname{eval}(\operatorname{Lean\mathord{:}Omega\mathord{:}Linear\,Combo\mathord{:}mk}(0,\allowbreak{} \langle 0\rangle \mathbin{\Vert} \langle 1\rangle),\allowbreak{} \operatorname{Lean\mathord{:}Omega\mathord{:}Coeffs\mathord{:}of\,List}(\langle n\rangle \mathbin{\Vert} (\langle |s|\rangle \mathbin{\Vert} (\langle |t|\rangle \mathbin{\Vert} (\langle k\rangle \mathbin{\Vert} (\langle |u|\rangle \mathbin{\Vert} \langle |u| - |s|\rangle)))))) = \operatorname{eval}(\operatorname{Lean\mathord{:}Omega\mathord{:}Linear\,Combo\mathord{:}mk}(0,\allowbreak{} \langle 1\rangle) - \operatorname{Lean\mathord{:}Omega\mathord{:}Linear\,Combo\mathord{:}mk}(0,\allowbreak{} \langle 0\rangle \mathbin{\Vert} \langle 1\rangle),\allowbreak{} \operatorname{Lean\mathord{:}Omega\mathord{:}Coeffs\mathord{:}of\,List}(\langle n\rangle \mathbin{\Vert} (\langle |s|\rangle \mathbin{\Vert} (\langle |t|\rangle \mathbin{\Vert} (\langle k\rangle \mathbin{\Vert} (\langle |u|\rangle \mathbin{\Vert} \langle |u| - |s|\rangle))))))\).
Chaining the equalities, we obtain \(n - |s| = \operatorname{eval}(\operatorname{Lean\mathord{:}Omega\mathord{:}Linear\,Combo\mathord{:}mk}(0,\allowbreak{} \langle 1\rangle) - \operatorname{Lean\mathord{:}Omega\mathord{:}Linear\,Combo\mathord{:}mk}(0,\allowbreak{} \langle 0\rangle \mathbin{\Vert} \langle 1\rangle),\allowbreak{} \operatorname{Lean\mathord{:}Omega\mathord{:}Coeffs\mathord{:}of\,List}(\langle n\rangle \mathbin{\Vert} (\langle |s|\rangle \mathbin{\Vert} (\langle |t|\rangle \mathbin{\Vert} (\langle k\rangle \mathbin{\Vert} (\langle |u|\rangle \mathbin{\Vert} \langle |u| - |s|\rangle))))))\).
By \paperlink{premise:6c655f6f665f6c655f6f665f6571}{Prerequisite~P\ref{premise:6c655f6f665f6c655f6f665f6571}} with \(\mathbb{Z}\) as its type parameter \(\alpha\), we obtain \(0 \le \operatorname{eval}(\operatorname{Lean\mathord{:}Omega\mathord{:}Linear\,Combo\mathord{:}mk}(0,\allowbreak{} \langle 1\rangle) - \operatorname{Lean\mathord{:}Omega\mathord{:}Linear\,Combo\mathord{:}mk}(0,\allowbreak{} \langle 0\rangle \mathbin{\Vert} \langle 1\rangle),\allowbreak{} \operatorname{Lean\mathord{:}Omega\mathord{:}Coeffs\mathord{:}of\,List}(\langle n\rangle \mathbin{\Vert} (\langle |s|\rangle \mathbin{\Vert} (\langle |t|\rangle \mathbin{\Vert} (\langle k\rangle \mathbin{\Vert} (\langle |u|\rangle \mathbin{\Vert} \langle |u| - |s|\rangle))))))\).
By \paperlink{premise:4c65616e2e4f6d6567612e436f6e73747261696e742e616464496e657175616c6974795f736174}{Prerequisite~P\ref{premise:4c65616e2e4f6d6567612e436f6e73747261696e742e616464496e657175616c6974795f736174}}, we obtain \(\operatorname{sat^{\prime}}(\operatorname{Lean\mathord{:}Omega\mathord{:}Constraint\mathord{:}mk}(\operatorname{some}(-0),\allowbreak{} \operatorname{none}),\allowbreak{} \langle 1\rangle \mathbin{\Vert} \langle -1\rangle,\allowbreak{} \operatorname{Lean\mathord{:}Omega\mathord{:}Coeffs\mathord{:}of\,List}(\langle n\rangle \mathbin{\Vert} (\langle |s|\rangle \mathbin{\Vert} (\langle |t|\rangle \mathbin{\Vert} (\langle k\rangle \mathbin{\Vert} (\langle |u|\rangle \mathbin{\Vert} \langle |u| - |s|\rangle)))))) = \mathrm{true}\).
By \paperlink{premise:4c65616e2e4f6d6567612e636f6d626f5f73617427}{Prerequisite~P\ref{premise:4c65616e2e4f6d6567612e636f6d626f5f73617427}} with \(\operatorname{Lean\mathord{:}Omega\mathord{:}Constraint\mathord{:}mk}(\operatorname{some}(0),\allowbreak{} \operatorname{some}(0))\) in place of its parameter \(s\) and \(\operatorname{Lean\mathord{:}Omega\mathord{:}Constraint\mathord{:}mk}(\operatorname{some}(0),\allowbreak{} \operatorname{none})\) in place of its parameter \(t\) and \(1\) in place of its parameter \(a\) and \(x = \langle 0\rangle \mathbin{\Vert} (\langle 1\rangle \mathbin{\Vert} (\langle 0\rangle \mathbin{\Vert} (\langle 0\rangle \mathbin{\Vert} (\langle -1\rangle \mathbin{\Vert} \langle 0\rangle))))\) and \(b = 1\) and \(y = \langle 1\rangle \mathbin{\Vert} \langle -1\rangle\) and \(v = \operatorname{Lean\mathord{:}Omega\mathord{:}Coeffs\mathord{:}of\,List}(\langle n\rangle \mathbin{\Vert} (\langle |s|\rangle \mathbin{\Vert} (\langle |t|\rangle \mathbin{\Vert} (\langle k\rangle \mathbin{\Vert} (\langle |u|\rangle \mathbin{\Vert} \langle |u| - |s|\rangle)))))\), we obtain \(\operatorname{sat^{\prime}}(\operatorname{Lean\mathord{:}Omega\mathord{:}Constraint\mathord{:}combo}(1,\allowbreak{} \operatorname{Lean\mathord{:}Omega\mathord{:}Constraint\mathord{:}mk}(\operatorname{some}(0),\allowbreak{} \operatorname{some}(0)),\allowbreak{} 1,\allowbreak{} \operatorname{Lean\mathord{:}Omega\mathord{:}Constraint\mathord{:}mk}(\operatorname{some}(0),\allowbreak{} \operatorname{none})),\allowbreak{} \operatorname{Lean\mathord{:}Omega\mathord{:}Coeffs\mathord{:}combo}(1,\allowbreak{} \langle 0\rangle \mathbin{\Vert} (\langle 1\rangle \mathbin{\Vert} (\langle 0\rangle \mathbin{\Vert} (\langle 0\rangle \mathbin{\Vert} (\langle -1\rangle \mathbin{\Vert} \langle 0\rangle)))),\allowbreak{} 1,\allowbreak{} \langle 1\rangle \mathbin{\Vert} \langle -1\rangle),\allowbreak{} \operatorname{Lean\mathord{:}Omega\mathord{:}Coeffs\mathord{:}of\,List}(\langle n\rangle \mathbin{\Vert} (\langle |s|\rangle \mathbin{\Vert} (\langle |t|\rangle \mathbin{\Vert} (\langle k\rangle \mathbin{\Vert} (\langle |u|\rangle \mathbin{\Vert} \langle |u| - |s|\rangle)))))) = \mathrm{true}\).
Using \paperlink{premise:4c65616e2e4f6d6567612e4c696e656172436f6d626f2e636f6f7264696e6174655f6576616c5f34}{Prerequisite~P\ref{premise:4c65616e2e4f6d6567612e4c696e656172436f6d626f2e636f6f7264696e6174655f6576616c5f34}} in reverse gives \(|u| = \operatorname{eval}(\operatorname{coordinate}(4),\allowbreak{} \operatorname{Lean\mathord{:}Omega\mathord{:}Coeffs\mathord{:}of\,List}(\langle n\rangle \mathbin{\Vert} (\langle |s|\rangle \mathbin{\Vert} (\langle |t|\rangle \mathbin{\Vert} (\langle k\rangle \mathbin{\Vert} (\langle |u|\rangle \mathbin{\Vert} \langle |u| - |s|\rangle))))))\).
By \paperlink{premise:496e742e6f664e61745f616464}{Prerequisite~P\ref{premise:496e742e6f664e61745f616464}} with \(k\) in place of its parameter \(n\) and \(m = 1\), we obtain \(k + 1 = k + 1\).
Using \paperlink{premise:4c65616e2e4f6d6567612e4c696e656172436f6d626f2e636f6f7264696e6174655f6576616c5f33}{Prerequisite~P\ref{premise:4c65616e2e4f6d6567612e4c696e656172436f6d626f2e636f6f7264696e6174655f6576616c5f33}} in reverse gives \(k = \operatorname{eval}(\operatorname{coordinate}(3),\allowbreak{} \operatorname{Lean\mathord{:}Omega\mathord{:}Coeffs\mathord{:}of\,List}(\langle n\rangle \mathbin{\Vert} (\langle |s|\rangle \mathbin{\Vert} (\langle |t|\rangle \mathbin{\Vert} (\langle k\rangle \mathbin{\Vert} (\langle |u|\rangle \mathbin{\Vert} \langle |u| - |s|\rangle))))))\).
By definition, we obtain \(1 = \operatorname{eval}(\operatorname{Lean\mathord{:}Omega\mathord{:}Linear\,Combo\mathord{:}mk}(1,\allowbreak{} \langle\,\rangle),\allowbreak{} \operatorname{Lean\mathord{:}Omega\mathord{:}Coeffs\mathord{:}of\,List}(\langle n\rangle \mathbin{\Vert} (\langle |s|\rangle \mathbin{\Vert} (\langle |t|\rangle \mathbin{\Vert} (\langle k\rangle \mathbin{\Vert} (\langle |u|\rangle \mathbin{\Vert} \langle |u| - |s|\rangle))))))\).
By \paperlink{premise:4c65616e2e4f6d6567612e496e742e6164645f636f6e6772}{Prerequisite~P\ref{premise:4c65616e2e4f6d6567612e496e742e6164645f636f6e6772}}, we obtain \(k + 1 = \operatorname{eval}(\operatorname{coordinate}(3),\allowbreak{} \operatorname{Lean\mathord{:}Omega\mathord{:}Coeffs\mathord{:}of\,List}(\langle n\rangle \mathbin{\Vert} (\langle |s|\rangle \mathbin{\Vert} (\langle |t|\rangle \mathbin{\Vert} (\langle k\rangle \mathbin{\Vert} (\langle |u|\rangle \mathbin{\Vert} \langle |u| - |s|\rangle)))))) + \operatorname{eval}(\operatorname{Lean\mathord{:}Omega\mathord{:}Linear\,Combo\mathord{:}mk}(1,\allowbreak{} \langle\,\rangle),\allowbreak{} \operatorname{Lean\mathord{:}Omega\mathord{:}Coeffs\mathord{:}of\,List}(\langle n\rangle \mathbin{\Vert} (\langle |s|\rangle \mathbin{\Vert} (\langle |t|\rangle \mathbin{\Vert} (\langle k\rangle \mathbin{\Vert} (\langle |u|\rangle \mathbin{\Vert} \langle |u| - |s|\rangle))))))\).
Using \paperlink{premise:4c65616e2e4f6d6567612e4c696e656172436f6d626f2e6164645f6576616c}{Prerequisite~P\ref{premise:4c65616e2e4f6d6567612e4c696e656172436f6d626f2e6164645f6576616c}} with \(l_{1} = \operatorname{Lean\mathord{:}Omega\mathord{:}Linear\,Combo\mathord{:}mk}(0,\allowbreak{} \langle 0\rangle \mathbin{\Vert} (\langle 0\rangle \mathbin{\Vert} (\langle 0\rangle \mathbin{\Vert} \langle 1\rangle)))\) and \(l_{2} = \operatorname{Lean\mathord{:}Omega\mathord{:}Linear\,Combo\mathord{:}mk}(1,\allowbreak{} \langle\,\rangle)\) and \(v = \operatorname{Lean\mathord{:}Omega\mathord{:}Coeffs\mathord{:}of\,List}(\langle n\rangle \mathbin{\Vert} (\langle |s|\rangle \mathbin{\Vert} (\langle |t|\rangle \mathbin{\Vert} (\langle k\rangle \mathbin{\Vert} (\langle |u|\rangle \mathbin{\Vert} \langle |u| - |s|\rangle)))))\) in reverse gives \(\operatorname{eval}(\operatorname{Lean\mathord{:}Omega\mathord{:}Linear\,Combo\mathord{:}mk}(0,\allowbreak{} \langle 0\rangle \mathbin{\Vert} (\langle 0\rangle \mathbin{\Vert} (\langle 0\rangle \mathbin{\Vert} \langle 1\rangle))),\allowbreak{} \operatorname{Lean\mathord{:}Omega\mathord{:}Coeffs\mathord{:}of\,List}(\langle n\rangle \mathbin{\Vert} (\langle |s|\rangle \mathbin{\Vert} (\langle |t|\rangle \mathbin{\Vert} (\langle k\rangle \mathbin{\Vert} (\langle |u|\rangle \mathbin{\Vert} \langle |u| - |s|\rangle)))))) + \operatorname{eval}(\operatorname{Lean\mathord{:}Omega\mathord{:}Linear\,Combo\mathord{:}mk}(1,\allowbreak{} \langle\,\rangle),\allowbreak{} \operatorname{Lean\mathord{:}Omega\mathord{:}Coeffs\mathord{:}of\,List}(\langle n\rangle \mathbin{\Vert} (\langle |s|\rangle \mathbin{\Vert} (\langle |t|\rangle \mathbin{\Vert} (\langle k\rangle \mathbin{\Vert} (\langle |u|\rangle \mathbin{\Vert} \langle |u| - |s|\rangle)))))) = \operatorname{eval}(\operatorname{Lean\mathord{:}Omega\mathord{:}Linear\,Combo\mathord{:}mk}(0,\allowbreak{} \langle 0\rangle \mathbin{\Vert} (\langle 0\rangle \mathbin{\Vert} (\langle 0\rangle \mathbin{\Vert} \langle 1\rangle))) + \operatorname{Lean\mathord{:}Omega\mathord{:}Linear\,Combo\mathord{:}mk}(1,\allowbreak{} \langle\,\rangle),\allowbreak{} \operatorname{Lean\mathord{:}Omega\mathord{:}Coeffs\mathord{:}of\,List}(\langle n\rangle \mathbin{\Vert} (\langle |s|\rangle \mathbin{\Vert} (\langle |t|\rangle \mathbin{\Vert} (\langle k\rangle \mathbin{\Vert} (\langle |u|\rangle \mathbin{\Vert} \langle |u| - |s|\rangle))))))\).
Chaining the equalities, we obtain \(k + 1 = \operatorname{eval}(\operatorname{Lean\mathord{:}Omega\mathord{:}Linear\,Combo\mathord{:}mk}(0,\allowbreak{} \langle 0\rangle \mathbin{\Vert} (\langle 0\rangle \mathbin{\Vert} (\langle 0\rangle \mathbin{\Vert} \langle 1\rangle))) + \operatorname{Lean\mathord{:}Omega\mathord{:}Linear\,Combo\mathord{:}mk}(1,\allowbreak{} \langle\,\rangle),\allowbreak{} \operatorname{Lean\mathord{:}Omega\mathord{:}Coeffs\mathord{:}of\,List}(\langle n\rangle \mathbin{\Vert} (\langle |s|\rangle \mathbin{\Vert} (\langle |t|\rangle \mathbin{\Vert} (\langle k\rangle \mathbin{\Vert} (\langle |u|\rangle \mathbin{\Vert} \langle |u| - |s|\rangle))))))\).
Chaining the equalities, we obtain \(k + 1 = \operatorname{eval}(\operatorname{Lean\mathord{:}Omega\mathord{:}Linear\,Combo\mathord{:}mk}(0,\allowbreak{} \langle 0\rangle \mathbin{\Vert} (\langle 0\rangle \mathbin{\Vert} (\langle 0\rangle \mathbin{\Vert} \langle 1\rangle))) + \operatorname{Lean\mathord{:}Omega\mathord{:}Linear\,Combo\mathord{:}mk}(1,\allowbreak{} \langle\,\rangle),\allowbreak{} \operatorname{Lean\mathord{:}Omega\mathord{:}Coeffs\mathord{:}of\,List}(\langle n\rangle \mathbin{\Vert} (\langle |s|\rangle \mathbin{\Vert} (\langle |t|\rangle \mathbin{\Vert} (\langle k\rangle \mathbin{\Vert} (\langle |u|\rangle \mathbin{\Vert} \langle |u| - |s|\rangle))))))\).
By \paperlink{premise:4c65616e2e4f6d6567612e496e742e7375625f636f6e6772}{Prerequisite~P\ref{premise:4c65616e2e4f6d6567612e496e742e7375625f636f6e6772}}, we obtain \(|u| - (k + 1) = \operatorname{eval}(\operatorname{coordinate}(4),\allowbreak{} \operatorname{Lean\mathord{:}Omega\mathord{:}Coeffs\mathord{:}of\,List}(\langle n\rangle \mathbin{\Vert} (\langle |s|\rangle \mathbin{\Vert} (\langle |t|\rangle \mathbin{\Vert} (\langle k\rangle \mathbin{\Vert} (\langle |u|\rangle \mathbin{\Vert} \langle |u| - |s|\rangle)))))) - \operatorname{eval}(\operatorname{Lean\mathord{:}Omega\mathord{:}Linear\,Combo\mathord{:}mk}(0,\allowbreak{} \langle 0\rangle \mathbin{\Vert} (\langle 0\rangle \mathbin{\Vert} (\langle 0\rangle \mathbin{\Vert} \langle 1\rangle))) + \operatorname{Lean\mathord{:}Omega\mathord{:}Linear\,Combo\mathord{:}mk}(1,\allowbreak{} \langle\,\rangle),\allowbreak{} \operatorname{Lean\mathord{:}Omega\mathord{:}Coeffs\mathord{:}of\,List}(\langle n\rangle \mathbin{\Vert} (\langle |s|\rangle \mathbin{\Vert} (\langle |t|\rangle \mathbin{\Vert} (\langle k\rangle \mathbin{\Vert} (\langle |u|\rangle \mathbin{\Vert} \langle |u| - |s|\rangle))))))\).
Using \paperlink{premise:4c65616e2e4f6d6567612e4c696e656172436f6d626f2e7375625f6576616c}{Prerequisite~P\ref{premise:4c65616e2e4f6d6567612e4c696e656172436f6d626f2e7375625f6576616c}} with \(l_{1} = \operatorname{Lean\mathord{:}Omega\mathord{:}Linear\,Combo\mathord{:}mk}(0,\allowbreak{} \langle 0\rangle \mathbin{\Vert} (\langle 0\rangle \mathbin{\Vert} (\langle 0\rangle \mathbin{\Vert} (\langle 0\rangle \mathbin{\Vert} \langle 1\rangle))))\) and \(l_{2} = \operatorname{Lean\mathord{:}Omega\mathord{:}Linear\,Combo\mathord{:}mk}(1,\allowbreak{} \langle 0\rangle \mathbin{\Vert} (\langle 0\rangle \mathbin{\Vert} (\langle 0\rangle \mathbin{\Vert} \langle 1\rangle)))\) and \(v = \operatorname{Lean\mathord{:}Omega\mathord{:}Coeffs\mathord{:}of\,List}(\langle n\rangle \mathbin{\Vert} (\langle |s|\rangle \mathbin{\Vert} (\langle |t|\rangle \mathbin{\Vert} (\langle k\rangle \mathbin{\Vert} (\langle |u|\rangle \mathbin{\Vert} \langle |u| - |s|\rangle)))))\) in reverse gives \(\operatorname{eval}(\operatorname{Lean\mathord{:}Omega\mathord{:}Linear\,Combo\mathord{:}mk}(0,\allowbreak{} \langle 0\rangle \mathbin{\Vert} (\langle 0\rangle \mathbin{\Vert} (\langle 0\rangle \mathbin{\Vert} (\langle 0\rangle \mathbin{\Vert} \langle 1\rangle)))),\allowbreak{} \operatorname{Lean\mathord{:}Omega\mathord{:}Coeffs\mathord{:}of\,List}(\langle n\rangle \mathbin{\Vert} (\langle |s|\rangle \mathbin{\Vert} (\langle |t|\rangle \mathbin{\Vert} (\langle k\rangle \mathbin{\Vert} (\langle |u|\rangle \mathbin{\Vert} \langle |u| - |s|\rangle)))))) - \operatorname{eval}(\operatorname{Lean\mathord{:}Omega\mathord{:}Linear\,Combo\mathord{:}mk}(1,\allowbreak{} \langle 0\rangle \mathbin{\Vert} (\langle 0\rangle \mathbin{\Vert} (\langle 0\rangle \mathbin{\Vert} \langle 1\rangle))),\allowbreak{} \operatorname{Lean\mathord{:}Omega\mathord{:}Coeffs\mathord{:}of\,List}(\langle n\rangle \mathbin{\Vert} (\langle |s|\rangle \mathbin{\Vert} (\langle |t|\rangle \mathbin{\Vert} (\langle k\rangle \mathbin{\Vert} (\langle |u|\rangle \mathbin{\Vert} \langle |u| - |s|\rangle)))))) = \operatorname{eval}(\operatorname{Lean\mathord{:}Omega\mathord{:}Linear\,Combo\mathord{:}mk}(0,\allowbreak{} \langle 0\rangle \mathbin{\Vert} (\langle 0\rangle \mathbin{\Vert} (\langle 0\rangle \mathbin{\Vert} (\langle 0\rangle \mathbin{\Vert} \langle 1\rangle)))) - \operatorname{Lean\mathord{:}Omega\mathord{:}Linear\,Combo\mathord{:}mk}(1,\allowbreak{} \langle 0\rangle \mathbin{\Vert} (\langle 0\rangle \mathbin{\Vert} (\langle 0\rangle \mathbin{\Vert} \langle 1\rangle))),\allowbreak{} \operatorname{Lean\mathord{:}Omega\mathord{:}Coeffs\mathord{:}of\,List}(\langle n\rangle \mathbin{\Vert} (\langle |s|\rangle \mathbin{\Vert} (\langle |t|\rangle \mathbin{\Vert} (\langle k\rangle \mathbin{\Vert} (\langle |u|\rangle \mathbin{\Vert} \langle |u| - |s|\rangle))))))\).
Chaining the equalities, we obtain \(|u| - (k + 1) = \operatorname{eval}(\operatorname{Lean\mathord{:}Omega\mathord{:}Linear\,Combo\mathord{:}mk}(0,\allowbreak{} \langle 0\rangle \mathbin{\Vert} (\langle 0\rangle \mathbin{\Vert} (\langle 0\rangle \mathbin{\Vert} (\langle 0\rangle \mathbin{\Vert} \langle 1\rangle)))) - \operatorname{Lean\mathord{:}Omega\mathord{:}Linear\,Combo\mathord{:}mk}(1,\allowbreak{} \langle 0\rangle \mathbin{\Vert} (\langle 0\rangle \mathbin{\Vert} (\langle 0\rangle \mathbin{\Vert} \langle 1\rangle))),\allowbreak{} \operatorname{Lean\mathord{:}Omega\mathord{:}Coeffs\mathord{:}of\,List}(\langle n\rangle \mathbin{\Vert} (\langle |s|\rangle \mathbin{\Vert} (\langle |t|\rangle \mathbin{\Vert} (\langle k\rangle \mathbin{\Vert} (\langle |u|\rangle \mathbin{\Vert} \langle |u| - |s|\rangle))))))\).
Reversing the equality, we obtain \(\operatorname{eval}(\operatorname{Lean\mathord{:}Omega\mathord{:}Linear\,Combo\mathord{:}mk}(0,\allowbreak{} \langle 0\rangle \mathbin{\Vert} (\langle 0\rangle \mathbin{\Vert} (\langle 0\rangle \mathbin{\Vert} (\langle 0\rangle \mathbin{\Vert} \langle 1\rangle)))) - \operatorname{Lean\mathord{:}Omega\mathord{:}Linear\,Combo\mathord{:}mk}(1,\allowbreak{} \langle 0\rangle \mathbin{\Vert} (\langle 0\rangle \mathbin{\Vert} (\langle 0\rangle \mathbin{\Vert} \langle 1\rangle))),\allowbreak{} \operatorname{Lean\mathord{:}Omega\mathord{:}Coeffs\mathord{:}of\,List}(\langle n\rangle \mathbin{\Vert} (\langle |s|\rangle \mathbin{\Vert} (\langle |t|\rangle \mathbin{\Vert} (\langle k\rangle \mathbin{\Vert} (\langle |u|\rangle \mathbin{\Vert} \langle |u| - |s|\rangle)))))) = |u| - (k + 1)\).
Use the hypothesis \(\left|u\right| = k + 1\).
By \paperlink{premise:4c65616e2e4f6d6567612e496e742e6f664e61745f636f6e6772}{Prerequisite~P\ref{premise:4c65616e2e4f6d6567612e496e742e6f664e61745f636f6e6772}}, we obtain \(\left|u\right| = k + 1\).
By \paperlink{premise:496e742e7375625f65715f7a65726f5f6f665f6571}{Prerequisite~P\ref{premise:496e742e7375625f65715f7a65726f5f6f665f6571}}, we obtain \(\left|u\right| - \left(k + 1\right) = 0\).
Chaining the equalities, we obtain \(\operatorname{eval}(\operatorname{Lean\mathord{:}Omega\mathord{:}Linear\,Combo\mathord{:}mk}(0,\allowbreak{} \langle 0\rangle \mathbin{\Vert} (\langle 0\rangle \mathbin{\Vert} (\langle 0\rangle \mathbin{\Vert} (\langle 0\rangle \mathbin{\Vert} \langle 1\rangle)))) - \operatorname{Lean\mathord{:}Omega\mathord{:}Linear\,Combo\mathord{:}mk}(1,\allowbreak{} \langle 0\rangle \mathbin{\Vert} (\langle 0\rangle \mathbin{\Vert} (\langle 0\rangle \mathbin{\Vert} \langle 1\rangle))),\allowbreak{} \operatorname{Lean\mathord{:}Omega\mathord{:}Coeffs\mathord{:}of\,List}(\langle n\rangle \mathbin{\Vert} (\langle |s|\rangle \mathbin{\Vert} (\langle |t|\rangle \mathbin{\Vert} (\langle k\rangle \mathbin{\Vert} (\langle |u|\rangle \mathbin{\Vert} \langle |u| - |s|\rangle)))))) = 0\).
By \paperlink{premise:4c65616e2e4f6d6567612e436f6e73747261696e742e616464457175616c6974795f736174}{Prerequisite~P\ref{premise:4c65616e2e4f6d6567612e436f6e73747261696e742e616464457175616c6974795f736174}}, we obtain \(\operatorname{sat^{\prime}}(\operatorname{Lean\mathord{:}Omega\mathord{:}Constraint\mathord{:}mk}(\operatorname{some}(-(-1)),\allowbreak{} \operatorname{some}(-(-1))),\allowbreak{} \langle 0\rangle \mathbin{\Vert} (\langle 0\rangle \mathbin{\Vert} (\langle 0\rangle \mathbin{\Vert} (\langle -1\rangle \mathbin{\Vert} \langle 1\rangle))),\allowbreak{} \operatorname{Lean\mathord{:}Omega\mathord{:}Coeffs\mathord{:}of\,List}(\langle n\rangle \mathbin{\Vert} (\langle |s|\rangle \mathbin{\Vert} (\langle |t|\rangle \mathbin{\Vert} (\langle k\rangle \mathbin{\Vert} (\langle |u|\rangle \mathbin{\Vert} \langle |u| - |s|\rangle)))))) = \mathrm{true}\).
By \paperlink{premise:4c65616e2e4f6d6567612e746964795f736174}{Prerequisite~P\ref{premise:4c65616e2e4f6d6567612e746964795f736174}}, we obtain \(\operatorname{sat^{\prime}}(\operatorname{tidy\,Constraint}(\operatorname{Lean\mathord{:}Omega\mathord{:}Constraint\mathord{:}mk}(\operatorname{some}(1),\allowbreak{} \operatorname{some}(1)),\allowbreak{} \langle 0\rangle \mathbin{\Vert} (\langle 0\rangle \mathbin{\Vert} (\langle 0\rangle \mathbin{\Vert} (\langle -1\rangle \mathbin{\Vert} \langle 1\rangle)))),\allowbreak{} \operatorname{tidy\,Coeffs}(\operatorname{Lean\mathord{:}Omega\mathord{:}Constraint\mathord{:}mk}(\operatorname{some}(1),\allowbreak{} \operatorname{some}(1)),\allowbreak{} \langle 0\rangle \mathbin{\Vert} (\langle 0\rangle \mathbin{\Vert} (\langle 0\rangle \mathbin{\Vert} (\langle -1\rangle \mathbin{\Vert} \langle 1\rangle)))),\allowbreak{} \operatorname{Lean\mathord{:}Omega\mathord{:}Coeffs\mathord{:}of\,List}(\langle n\rangle \mathbin{\Vert} (\langle |s|\rangle \mathbin{\Vert} (\langle |t|\rangle \mathbin{\Vert} (\langle k\rangle \mathbin{\Vert} (\langle |u|\rangle \mathbin{\Vert} \langle |u| - |s|\rangle)))))) = \mathrm{true}\).
Use the hypothesis \(n \le k\).
By \paperlink{premise:4c65616e2e4f6d6567612e496e742e6f664e61745f6c655f6f665f6c65}{Prerequisite~P\ref{premise:4c65616e2e4f6d6567612e496e742e6f664e61745f6c655f6f665f6c65}}, we obtain \(n \le k\).
By \paperlink{premise:496e742e7375625f6e6f6e6e65675f6f665f6c65}{Prerequisite~P\ref{premise:496e742e7375625f6e6f6e6e65675f6f665f6c65}}, we obtain \(0 \le k - n\).
Using \paperlink{premise:4c65616e2e4f6d6567612e4c696e656172436f6d626f2e636f6f7264696e6174655f6576616c5f33}{Prerequisite~P\ref{premise:4c65616e2e4f6d6567612e4c696e656172436f6d626f2e636f6f7264696e6174655f6576616c5f33}} in reverse gives \(k = \operatorname{eval}(\operatorname{coordinate}(3),\allowbreak{} \operatorname{Lean\mathord{:}Omega\mathord{:}Coeffs\mathord{:}of\,List}(\langle n\rangle \mathbin{\Vert} (\langle |s|\rangle \mathbin{\Vert} (\langle |t|\rangle \mathbin{\Vert} (\langle k\rangle \mathbin{\Vert} (\langle |u|\rangle \mathbin{\Vert} \langle |u| - |s|\rangle))))))\).
Using \paperlink{premise:4c65616e2e4f6d6567612e4c696e656172436f6d626f2e636f6f7264696e6174655f6576616c5f30}{Prerequisite~P\ref{premise:4c65616e2e4f6d6567612e4c696e656172436f6d626f2e636f6f7264696e6174655f6576616c5f30}} in reverse gives \(n = \operatorname{eval}(\operatorname{coordinate}(0),\allowbreak{} \operatorname{Lean\mathord{:}Omega\mathord{:}Coeffs\mathord{:}of\,List}(\langle n\rangle \mathbin{\Vert} (\langle |s|\rangle \mathbin{\Vert} (\langle |t|\rangle \mathbin{\Vert} (\langle k\rangle \mathbin{\Vert} (\langle |u|\rangle \mathbin{\Vert} \langle |u| - |s|\rangle))))))\).
By \paperlink{premise:4c65616e2e4f6d6567612e496e742e7375625f636f6e6772}{Prerequisite~P\ref{premise:4c65616e2e4f6d6567612e496e742e7375625f636f6e6772}}, we obtain \(k - n = \operatorname{eval}(\operatorname{coordinate}(3),\allowbreak{} \operatorname{Lean\mathord{:}Omega\mathord{:}Coeffs\mathord{:}of\,List}(\langle n\rangle \mathbin{\Vert} (\langle |s|\rangle \mathbin{\Vert} (\langle |t|\rangle \mathbin{\Vert} (\langle k\rangle \mathbin{\Vert} (\langle |u|\rangle \mathbin{\Vert} \langle |u| - |s|\rangle)))))) - \operatorname{eval}(\operatorname{coordinate}(0),\allowbreak{} \operatorname{Lean\mathord{:}Omega\mathord{:}Coeffs\mathord{:}of\,List}(\langle n\rangle \mathbin{\Vert} (\langle |s|\rangle \mathbin{\Vert} (\langle |t|\rangle \mathbin{\Vert} (\langle k\rangle \mathbin{\Vert} (\langle |u|\rangle \mathbin{\Vert} \langle |u| - |s|\rangle))))))\).
Using \paperlink{premise:4c65616e2e4f6d6567612e4c696e656172436f6d626f2e7375625f6576616c}{Prerequisite~P\ref{premise:4c65616e2e4f6d6567612e4c696e656172436f6d626f2e7375625f6576616c}} with \(l_{1} = \operatorname{Lean\mathord{:}Omega\mathord{:}Linear\,Combo\mathord{:}mk}(0,\allowbreak{} \langle 0\rangle \mathbin{\Vert} (\langle 0\rangle \mathbin{\Vert} (\langle 0\rangle \mathbin{\Vert} \langle 1\rangle)))\) and \(l_{2} = \operatorname{Lean\mathord{:}Omega\mathord{:}Linear\,Combo\mathord{:}mk}(0,\allowbreak{} \langle 1\rangle)\) and \(v = \operatorname{Lean\mathord{:}Omega\mathord{:}Coeffs\mathord{:}of\,List}(\langle n\rangle \mathbin{\Vert} (\langle |s|\rangle \mathbin{\Vert} (\langle |t|\rangle \mathbin{\Vert} (\langle k\rangle \mathbin{\Vert} (\langle |u|\rangle \mathbin{\Vert} \langle |u| - |s|\rangle)))))\) in reverse gives \(\operatorname{eval}(\operatorname{Lean\mathord{:}Omega\mathord{:}Linear\,Combo\mathord{:}mk}(0,\allowbreak{} \langle 0\rangle \mathbin{\Vert} (\langle 0\rangle \mathbin{\Vert} (\langle 0\rangle \mathbin{\Vert} \langle 1\rangle))),\allowbreak{} \operatorname{Lean\mathord{:}Omega\mathord{:}Coeffs\mathord{:}of\,List}(\langle n\rangle \mathbin{\Vert} (\langle |s|\rangle \mathbin{\Vert} (\langle |t|\rangle \mathbin{\Vert} (\langle k\rangle \mathbin{\Vert} (\langle |u|\rangle \mathbin{\Vert} \langle |u| - |s|\rangle)))))) - \operatorname{eval}(\operatorname{Lean\mathord{:}Omega\mathord{:}Linear\,Combo\mathord{:}mk}(0,\allowbreak{} \langle 1\rangle),\allowbreak{} \operatorname{Lean\mathord{:}Omega\mathord{:}Coeffs\mathord{:}of\,List}(\langle n\rangle \mathbin{\Vert} (\langle |s|\rangle \mathbin{\Vert} (\langle |t|\rangle \mathbin{\Vert} (\langle k\rangle \mathbin{\Vert} (\langle |u|\rangle \mathbin{\Vert} \langle |u| - |s|\rangle)))))) = \operatorname{eval}(\operatorname{Lean\mathord{:}Omega\mathord{:}Linear\,Combo\mathord{:}mk}(0,\allowbreak{} \langle 0\rangle \mathbin{\Vert} (\langle 0\rangle \mathbin{\Vert} (\langle 0\rangle \mathbin{\Vert} \langle 1\rangle))) - \operatorname{Lean\mathord{:}Omega\mathord{:}Linear\,Combo\mathord{:}mk}(0,\allowbreak{} \langle 1\rangle),\allowbreak{} \operatorname{Lean\mathord{:}Omega\mathord{:}Coeffs\mathord{:}of\,List}(\langle n\rangle \mathbin{\Vert} (\langle |s|\rangle \mathbin{\Vert} (\langle |t|\rangle \mathbin{\Vert} (\langle k\rangle \mathbin{\Vert} (\langle |u|\rangle \mathbin{\Vert} \langle |u| - |s|\rangle))))))\).
Chaining the equalities, we obtain \(k - n = \operatorname{eval}(\operatorname{Lean\mathord{:}Omega\mathord{:}Linear\,Combo\mathord{:}mk}(0,\allowbreak{} \langle 0\rangle \mathbin{\Vert} (\langle 0\rangle \mathbin{\Vert} (\langle 0\rangle \mathbin{\Vert} \langle 1\rangle))) - \operatorname{Lean\mathord{:}Omega\mathord{:}Linear\,Combo\mathord{:}mk}(0,\allowbreak{} \langle 1\rangle),\allowbreak{} \operatorname{Lean\mathord{:}Omega\mathord{:}Coeffs\mathord{:}of\,List}(\langle n\rangle \mathbin{\Vert} (\langle |s|\rangle \mathbin{\Vert} (\langle |t|\rangle \mathbin{\Vert} (\langle k\rangle \mathbin{\Vert} (\langle |u|\rangle \mathbin{\Vert} \langle |u| - |s|\rangle))))))\).
By \paperlink{premise:6c655f6f665f6c655f6f665f6571}{Prerequisite~P\ref{premise:6c655f6f665f6c655f6f665f6571}} with \(\mathbb{Z}\) as its type parameter \(\alpha\), we obtain \(0 \le \operatorname{eval}(\operatorname{Lean\mathord{:}Omega\mathord{:}Linear\,Combo\mathord{:}mk}(0,\allowbreak{} \langle 0\rangle \mathbin{\Vert} (\langle 0\rangle \mathbin{\Vert} (\langle 0\rangle \mathbin{\Vert} \langle 1\rangle))) - \operatorname{Lean\mathord{:}Omega\mathord{:}Linear\,Combo\mathord{:}mk}(0,\allowbreak{} \langle 1\rangle),\allowbreak{} \operatorname{Lean\mathord{:}Omega\mathord{:}Coeffs\mathord{:}of\,List}(\langle n\rangle \mathbin{\Vert} (\langle |s|\rangle \mathbin{\Vert} (\langle |t|\rangle \mathbin{\Vert} (\langle k\rangle \mathbin{\Vert} (\langle |u|\rangle \mathbin{\Vert} \langle |u| - |s|\rangle))))))\).
By \paperlink{premise:4c65616e2e4f6d6567612e436f6e73747261696e742e616464496e657175616c6974795f736174}{Prerequisite~P\ref{premise:4c65616e2e4f6d6567612e436f6e73747261696e742e616464496e657175616c6974795f736174}}, we obtain \(\operatorname{sat^{\prime}}(\operatorname{Lean\mathord{:}Omega\mathord{:}Constraint\mathord{:}mk}(\operatorname{some}(-0),\allowbreak{} \operatorname{none}),\allowbreak{} \langle -1\rangle \mathbin{\Vert} (\langle 0\rangle \mathbin{\Vert} (\langle 0\rangle \mathbin{\Vert} \langle 1\rangle)),\allowbreak{} \operatorname{Lean\mathord{:}Omega\mathord{:}Coeffs\mathord{:}of\,List}(\langle n\rangle \mathbin{\Vert} (\langle |s|\rangle \mathbin{\Vert} (\langle |t|\rangle \mathbin{\Vert} (\langle k\rangle \mathbin{\Vert} (\langle |u|\rangle \mathbin{\Vert} \langle |u| - |s|\rangle)))))) = \mathrm{true}\).
By \paperlink{premise:4c65616e2e4f6d6567612e746964795f736174}{Prerequisite~P\ref{premise:4c65616e2e4f6d6567612e746964795f736174}}, we obtain \(\operatorname{sat^{\prime}}(\operatorname{tidy\,Constraint}(\operatorname{Lean\mathord{:}Omega\mathord{:}Constraint\mathord{:}mk}(\operatorname{some}(0),\allowbreak{} \operatorname{none}),\allowbreak{} \langle -1\rangle \mathbin{\Vert} (\langle 0\rangle \mathbin{\Vert} (\langle 0\rangle \mathbin{\Vert} \langle 1\rangle))),\allowbreak{} \operatorname{tidy\,Coeffs}(\operatorname{Lean\mathord{:}Omega\mathord{:}Constraint\mathord{:}mk}(\operatorname{some}(0),\allowbreak{} \operatorname{none}),\allowbreak{} \langle -1\rangle \mathbin{\Vert} (\langle 0\rangle \mathbin{\Vert} (\langle 0\rangle \mathbin{\Vert} \langle 1\rangle))),\allowbreak{} \operatorname{Lean\mathord{:}Omega\mathord{:}Coeffs\mathord{:}of\,List}(\langle n\rangle \mathbin{\Vert} (\langle |s|\rangle \mathbin{\Vert} (\langle |t|\rangle \mathbin{\Vert} (\langle k\rangle \mathbin{\Vert} (\langle |u|\rangle \mathbin{\Vert} \langle |u| - |s|\rangle)))))) = \mathrm{true}\).
By \paperlink{premise:4c65616e2e4f6d6567612e636f6d626f5f73617427}{Prerequisite~P\ref{premise:4c65616e2e4f6d6567612e636f6d626f5f73617427}} with \(\operatorname{Lean\mathord{:}Omega\mathord{:}Constraint\mathord{:}mk}(\operatorname{some}(-1),\allowbreak{} \operatorname{some}(-1))\) in place of its parameter \(s\) and \(\operatorname{Lean\mathord{:}Omega\mathord{:}Constraint\mathord{:}mk}(\operatorname{none},\allowbreak{} \operatorname{some}(0))\) in place of its parameter \(t\) and \(1\) in place of its parameter \(a\) and \(x = \langle 0\rangle \mathbin{\Vert} (\langle 0\rangle \mathbin{\Vert} (\langle 0\rangle \mathbin{\Vert} (\langle 1\rangle \mathbin{\Vert} \langle -1\rangle)))\) and \(b = 1\) and \(y = \langle 1\rangle \mathbin{\Vert} (\langle 0\rangle \mathbin{\Vert} (\langle 0\rangle \mathbin{\Vert} \langle -1\rangle))\) and \(v = \operatorname{Lean\mathord{:}Omega\mathord{:}Coeffs\mathord{:}of\,List}(\langle n\rangle \mathbin{\Vert} (\langle |s|\rangle \mathbin{\Vert} (\langle |t|\rangle \mathbin{\Vert} (\langle k\rangle \mathbin{\Vert} (\langle |u|\rangle \mathbin{\Vert} \langle |u| - |s|\rangle)))))\), we obtain \(\operatorname{sat^{\prime}}(\operatorname{Lean\mathord{:}Omega\mathord{:}Constraint\mathord{:}combo}(1,\allowbreak{} \operatorname{Lean\mathord{:}Omega\mathord{:}Constraint\mathord{:}mk}(\operatorname{some}(-1),\allowbreak{} \operatorname{some}(-1)),\allowbreak{} 1,\allowbreak{} \operatorname{Lean\mathord{:}Omega\mathord{:}Constraint\mathord{:}mk}(\operatorname{none},\allowbreak{} \operatorname{some}(0))),\allowbreak{} \operatorname{Lean\mathord{:}Omega\mathord{:}Coeffs\mathord{:}combo}(1,\allowbreak{} \langle 0\rangle \mathbin{\Vert} (\langle 0\rangle \mathbin{\Vert} (\langle 0\rangle \mathbin{\Vert} (\langle 1\rangle \mathbin{\Vert} \langle -1\rangle))),\allowbreak{} 1,\allowbreak{} \langle 1\rangle \mathbin{\Vert} (\langle 0\rangle \mathbin{\Vert} (\langle 0\rangle \mathbin{\Vert} \langle -1\rangle))),\allowbreak{} \operatorname{Lean\mathord{:}Omega\mathord{:}Coeffs\mathord{:}of\,List}(\langle n\rangle \mathbin{\Vert} (\langle |s|\rangle \mathbin{\Vert} (\langle |t|\rangle \mathbin{\Vert} (\langle k\rangle \mathbin{\Vert} (\langle |u|\rangle \mathbin{\Vert} \langle |u| - |s|\rangle)))))) = \mathrm{true}\).
By \paperlink{premise:4c65616e2e4f6d6567612e636f6d626f5f73617427}{Prerequisite~P\ref{premise:4c65616e2e4f6d6567612e636f6d626f5f73617427}} with \(\operatorname{Lean\mathord{:}Omega\mathord{:}Constraint\mathord{:}mk}(\operatorname{some}(0),\allowbreak{} \operatorname{none})\) in place of its parameter \(s\) and \(\operatorname{Lean\mathord{:}Omega\mathord{:}Constraint\mathord{:}mk}(\operatorname{none},\allowbreak{} \operatorname{some}(-1))\) in place of its parameter \(t\) and \(1\) in place of its parameter \(a\) and \(x = \langle 1\rangle \mathbin{\Vert} (\langle 0\rangle \mathbin{\Vert} (\langle 0\rangle \mathbin{\Vert} (\langle 0\rangle \mathbin{\Vert} (\langle -1\rangle \mathbin{\Vert} \langle 0\rangle))))\) and \(b = -1\) and \(y = \langle 1\rangle \mathbin{\Vert} (\langle 0\rangle \mathbin{\Vert} (\langle 0\rangle \mathbin{\Vert} (\langle 0\rangle \mathbin{\Vert} \langle -1\rangle)))\) and \(v = \operatorname{Lean\mathord{:}Omega\mathord{:}Coeffs\mathord{:}of\,List}(\langle n\rangle \mathbin{\Vert} (\langle |s|\rangle \mathbin{\Vert} (\langle |t|\rangle \mathbin{\Vert} (\langle k\rangle \mathbin{\Vert} (\langle |u|\rangle \mathbin{\Vert} \langle |u| - |s|\rangle)))))\), we obtain \(\operatorname{sat^{\prime}}(\operatorname{Lean\mathord{:}Omega\mathord{:}Constraint\mathord{:}combo}(1,\allowbreak{} \operatorname{Lean\mathord{:}Omega\mathord{:}Constraint\mathord{:}mk}(\operatorname{some}(0),\allowbreak{} \operatorname{none}),\allowbreak{} -1,\allowbreak{} \operatorname{Lean\mathord{:}Omega\mathord{:}Constraint\mathord{:}mk}(\operatorname{none},\allowbreak{} \operatorname{some}(-1))),\allowbreak{} \operatorname{Lean\mathord{:}Omega\mathord{:}Coeffs\mathord{:}combo}(1,\allowbreak{} \langle 1\rangle \mathbin{\Vert} (\langle 0\rangle \mathbin{\Vert} (\langle 0\rangle \mathbin{\Vert} (\langle 0\rangle \mathbin{\Vert} (\langle -1\rangle \mathbin{\Vert} \langle 0\rangle)))),\allowbreak{} -1,\allowbreak{} \langle 1\rangle \mathbin{\Vert} (\langle 0\rangle \mathbin{\Vert} (\langle 0\rangle \mathbin{\Vert} (\langle 0\rangle \mathbin{\Vert} \langle -1\rangle)))),\allowbreak{} \operatorname{Lean\mathord{:}Omega\mathord{:}Coeffs\mathord{:}of\,List}(\langle n\rangle \mathbin{\Vert} (\langle |s|\rangle \mathbin{\Vert} (\langle |t|\rangle \mathbin{\Vert} (\langle k\rangle \mathbin{\Vert} (\langle |u|\rangle \mathbin{\Vert} \langle |u| - |s|\rangle)))))) = \mathrm{true}\).
By \paperlink{premise:4c65616e2e4f6d6567612e746964795f736174}{Prerequisite~P\ref{premise:4c65616e2e4f6d6567612e746964795f736174}}, we obtain \(\operatorname{sat^{\prime}}(\operatorname{tidy\,Constraint}(\operatorname{Lean\mathord{:}Omega\mathord{:}Constraint\mathord{:}mk}(\operatorname{some}(1),\allowbreak{} \operatorname{none}),\allowbreak{} \langle 0\rangle \mathbin{\Vert} (\langle 0\rangle \mathbin{\Vert} (\langle 0\rangle \mathbin{\Vert} (\langle 0\rangle \mathbin{\Vert} (\langle 0\rangle \mathbin{\Vert} \langle 0\rangle))))),\allowbreak{} \operatorname{tidy\,Coeffs}(\operatorname{Lean\mathord{:}Omega\mathord{:}Constraint\mathord{:}mk}(\operatorname{some}(1),\allowbreak{} \operatorname{none}),\allowbreak{} \langle 0\rangle \mathbin{\Vert} (\langle 0\rangle \mathbin{\Vert} (\langle 0\rangle \mathbin{\Vert} (\langle 0\rangle \mathbin{\Vert} (\langle 0\rangle \mathbin{\Vert} \langle 0\rangle))))),\allowbreak{} \operatorname{Lean\mathord{:}Omega\mathord{:}Coeffs\mathord{:}of\,List}(\langle n\rangle \mathbin{\Vert} (\langle |s|\rangle \mathbin{\Vert} (\langle |t|\rangle \mathbin{\Vert} (\langle k\rangle \mathbin{\Vert} (\langle |u|\rangle \mathbin{\Vert} \langle |u| - |s|\rangle)))))) = \mathrm{true}\).
By \paperlink{premise:4c65616e2e4f6d6567612e436f6e73747261696e742e6e6f745f736174275f6f665f6973496d706f737369626c65}{Prerequisite~P\ref{premise:4c65616e2e4f6d6567612e436f6e73747261696e742e6e6f745f736174275f6f665f6973496d706f737369626c65}}, we obtain \(\mathrm{False}\).
\end{proofstep}
Consider the following case.
\begin{proofstep}
Assume \(\left|u\right| < \left|s\right| \land \left|u\right| - \left|s\right| = 0\).
By definition, we obtain \(\operatorname{decide}(\operatorname{is\,Impossible}(\operatorname{Lean\mathord{:}Omega\mathord{:}Constraint\mathord{:}mk}(\operatorname{some}(1),\allowbreak{} \operatorname{some}(-1))) = \mathrm{true}) = \mathrm{true}\).
By \paperlink{premise:6f665f6465636964655f65715f74727565}{Prerequisite~P\ref{premise:6f665f6465636964655f65715f74727565}}, we obtain \(\operatorname{is\,Impossible}(\operatorname{Lean\mathord{:}Omega\mathord{:}Constraint\mathord{:}mk}(\operatorname{some}(1),\allowbreak{} \operatorname{some}(-1))) = \mathrm{true}\).
Use the hypothesis \(\left|s\right| \le n\).
By \paperlink{premise:4c65616e2e4f6d6567612e496e742e6f664e61745f6c655f6f665f6c65}{Prerequisite~P\ref{premise:4c65616e2e4f6d6567612e496e742e6f664e61745f6c655f6f665f6c65}}, we obtain \(\left|s\right| \le n\).
By \paperlink{premise:496e742e7375625f6e6f6e6e65675f6f665f6c65}{Prerequisite~P\ref{premise:496e742e7375625f6e6f6e6e65675f6f665f6c65}}, we obtain \(0 \le n - \left|s\right|\).
Using \paperlink{premise:4c65616e2e4f6d6567612e4c696e656172436f6d626f2e636f6f7264696e6174655f6576616c5f30}{Prerequisite~P\ref{premise:4c65616e2e4f6d6567612e4c696e656172436f6d626f2e636f6f7264696e6174655f6576616c5f30}} in reverse gives \(n = \operatorname{eval}(\operatorname{coordinate}(0),\allowbreak{} \operatorname{Lean\mathord{:}Omega\mathord{:}Coeffs\mathord{:}of\,List}(\langle n\rangle \mathbin{\Vert} (\langle |s|\rangle \mathbin{\Vert} (\langle |t|\rangle \mathbin{\Vert} (\langle k\rangle \mathbin{\Vert} (\langle |u|\rangle \mathbin{\Vert} \langle |u| - |s|\rangle))))))\).
Using \paperlink{premise:4c65616e2e4f6d6567612e4c696e656172436f6d626f2e636f6f7264696e6174655f6576616c5f31}{Prerequisite~P\ref{premise:4c65616e2e4f6d6567612e4c696e656172436f6d626f2e636f6f7264696e6174655f6576616c5f31}} in reverse gives \(|s| = \operatorname{eval}(\operatorname{coordinate}(1),\allowbreak{} \operatorname{Lean\mathord{:}Omega\mathord{:}Coeffs\mathord{:}of\,List}(\langle n\rangle \mathbin{\Vert} (\langle |s|\rangle \mathbin{\Vert} (\langle |t|\rangle \mathbin{\Vert} (\langle k\rangle \mathbin{\Vert} (\langle |u|\rangle \mathbin{\Vert} \langle |u| - |s|\rangle))))))\).
By \paperlink{premise:4c65616e2e4f6d6567612e496e742e7375625f636f6e6772}{Prerequisite~P\ref{premise:4c65616e2e4f6d6567612e496e742e7375625f636f6e6772}}, we obtain \(n - |s| = \operatorname{eval}(\operatorname{coordinate}(0),\allowbreak{} \operatorname{Lean\mathord{:}Omega\mathord{:}Coeffs\mathord{:}of\,List}(\langle n\rangle \mathbin{\Vert} (\langle |s|\rangle \mathbin{\Vert} (\langle |t|\rangle \mathbin{\Vert} (\langle k\rangle \mathbin{\Vert} (\langle |u|\rangle \mathbin{\Vert} \langle |u| - |s|\rangle)))))) - \operatorname{eval}(\operatorname{coordinate}(1),\allowbreak{} \operatorname{Lean\mathord{:}Omega\mathord{:}Coeffs\mathord{:}of\,List}(\langle n\rangle \mathbin{\Vert} (\langle |s|\rangle \mathbin{\Vert} (\langle |t|\rangle \mathbin{\Vert} (\langle k\rangle \mathbin{\Vert} (\langle |u|\rangle \mathbin{\Vert} \langle |u| - |s|\rangle))))))\).
Using \paperlink{premise:4c65616e2e4f6d6567612e4c696e656172436f6d626f2e7375625f6576616c}{Prerequisite~P\ref{premise:4c65616e2e4f6d6567612e4c696e656172436f6d626f2e7375625f6576616c}} with \(l_{1} = \operatorname{Lean\mathord{:}Omega\mathord{:}Linear\,Combo\mathord{:}mk}(0,\allowbreak{} \langle 1\rangle)\) and \(l_{2} = \operatorname{Lean\mathord{:}Omega\mathord{:}Linear\,Combo\mathord{:}mk}(0,\allowbreak{} \langle 0\rangle \mathbin{\Vert} \langle 1\rangle)\) and \(v = \operatorname{Lean\mathord{:}Omega\mathord{:}Coeffs\mathord{:}of\,List}(\langle n\rangle \mathbin{\Vert} (\langle |s|\rangle \mathbin{\Vert} (\langle |t|\rangle \mathbin{\Vert} (\langle k\rangle \mathbin{\Vert} (\langle |u|\rangle \mathbin{\Vert} \langle |u| - |s|\rangle)))))\) in reverse gives \(\operatorname{eval}(\operatorname{Lean\mathord{:}Omega\mathord{:}Linear\,Combo\mathord{:}mk}(0,\allowbreak{} \langle 1\rangle),\allowbreak{} \operatorname{Lean\mathord{:}Omega\mathord{:}Coeffs\mathord{:}of\,List}(\langle n\rangle \mathbin{\Vert} (\langle |s|\rangle \mathbin{\Vert} (\langle |t|\rangle \mathbin{\Vert} (\langle k\rangle \mathbin{\Vert} (\langle |u|\rangle \mathbin{\Vert} \langle |u| - |s|\rangle)))))) - \operatorname{eval}(\operatorname{Lean\mathord{:}Omega\mathord{:}Linear\,Combo\mathord{:}mk}(0,\allowbreak{} \langle 0\rangle \mathbin{\Vert} \langle 1\rangle),\allowbreak{} \operatorname{Lean\mathord{:}Omega\mathord{:}Coeffs\mathord{:}of\,List}(\langle n\rangle \mathbin{\Vert} (\langle |s|\rangle \mathbin{\Vert} (\langle |t|\rangle \mathbin{\Vert} (\langle k\rangle \mathbin{\Vert} (\langle |u|\rangle \mathbin{\Vert} \langle |u| - |s|\rangle)))))) = \operatorname{eval}(\operatorname{Lean\mathord{:}Omega\mathord{:}Linear\,Combo\mathord{:}mk}(0,\allowbreak{} \langle 1\rangle) - \operatorname{Lean\mathord{:}Omega\mathord{:}Linear\,Combo\mathord{:}mk}(0,\allowbreak{} \langle 0\rangle \mathbin{\Vert} \langle 1\rangle),\allowbreak{} \operatorname{Lean\mathord{:}Omega\mathord{:}Coeffs\mathord{:}of\,List}(\langle n\rangle \mathbin{\Vert} (\langle |s|\rangle \mathbin{\Vert} (\langle |t|\rangle \mathbin{\Vert} (\langle k\rangle \mathbin{\Vert} (\langle |u|\rangle \mathbin{\Vert} \langle |u| - |s|\rangle))))))\).
Chaining the equalities, we obtain \(n - |s| = \operatorname{eval}(\operatorname{Lean\mathord{:}Omega\mathord{:}Linear\,Combo\mathord{:}mk}(0,\allowbreak{} \langle 1\rangle) - \operatorname{Lean\mathord{:}Omega\mathord{:}Linear\,Combo\mathord{:}mk}(0,\allowbreak{} \langle 0\rangle \mathbin{\Vert} \langle 1\rangle),\allowbreak{} \operatorname{Lean\mathord{:}Omega\mathord{:}Coeffs\mathord{:}of\,List}(\langle n\rangle \mathbin{\Vert} (\langle |s|\rangle \mathbin{\Vert} (\langle |t|\rangle \mathbin{\Vert} (\langle k\rangle \mathbin{\Vert} (\langle |u|\rangle \mathbin{\Vert} \langle |u| - |s|\rangle))))))\).
By \paperlink{premise:6c655f6f665f6c655f6f665f6571}{Prerequisite~P\ref{premise:6c655f6f665f6c655f6f665f6571}} with \(\mathbb{Z}\) as its type parameter \(\alpha\), we obtain \(0 \le \operatorname{eval}(\operatorname{Lean\mathord{:}Omega\mathord{:}Linear\,Combo\mathord{:}mk}(0,\allowbreak{} \langle 1\rangle) - \operatorname{Lean\mathord{:}Omega\mathord{:}Linear\,Combo\mathord{:}mk}(0,\allowbreak{} \langle 0\rangle \mathbin{\Vert} \langle 1\rangle),\allowbreak{} \operatorname{Lean\mathord{:}Omega\mathord{:}Coeffs\mathord{:}of\,List}(\langle n\rangle \mathbin{\Vert} (\langle |s|\rangle \mathbin{\Vert} (\langle |t|\rangle \mathbin{\Vert} (\langle k\rangle \mathbin{\Vert} (\langle |u|\rangle \mathbin{\Vert} \langle |u| - |s|\rangle))))))\).
By \paperlink{premise:4c65616e2e4f6d6567612e436f6e73747261696e742e616464496e657175616c6974795f736174}{Prerequisite~P\ref{premise:4c65616e2e4f6d6567612e436f6e73747261696e742e616464496e657175616c6974795f736174}}, we obtain \(\operatorname{sat^{\prime}}(\operatorname{Lean\mathord{:}Omega\mathord{:}Constraint\mathord{:}mk}(\operatorname{some}(-0),\allowbreak{} \operatorname{none}),\allowbreak{} \langle 1\rangle \mathbin{\Vert} \langle -1\rangle,\allowbreak{} \operatorname{Lean\mathord{:}Omega\mathord{:}Coeffs\mathord{:}of\,List}(\langle n\rangle \mathbin{\Vert} (\langle |s|\rangle \mathbin{\Vert} (\langle |t|\rangle \mathbin{\Vert} (\langle k\rangle \mathbin{\Vert} (\langle |u|\rangle \mathbin{\Vert} \langle |u| - |s|\rangle)))))) = \mathrm{true}\).
Using \paperlink{premise:4c65616e2e4f6d6567612e4c696e656172436f6d626f2e636f6f7264696e6174655f6576616c5f34}{Prerequisite~P\ref{premise:4c65616e2e4f6d6567612e4c696e656172436f6d626f2e636f6f7264696e6174655f6576616c5f34}} in reverse gives \(|u| = \operatorname{eval}(\operatorname{coordinate}(4),\allowbreak{} \operatorname{Lean\mathord{:}Omega\mathord{:}Coeffs\mathord{:}of\,List}(\langle n\rangle \mathbin{\Vert} (\langle |s|\rangle \mathbin{\Vert} (\langle |t|\rangle \mathbin{\Vert} (\langle k\rangle \mathbin{\Vert} (\langle |u|\rangle \mathbin{\Vert} \langle |u| - |s|\rangle))))))\).
By \paperlink{premise:496e742e6f664e61745f616464}{Prerequisite~P\ref{premise:496e742e6f664e61745f616464}} with \(k\) in place of its parameter \(n\) and \(m = 1\), we obtain \(k + 1 = k + 1\).
Using \paperlink{premise:4c65616e2e4f6d6567612e4c696e656172436f6d626f2e636f6f7264696e6174655f6576616c5f33}{Prerequisite~P\ref{premise:4c65616e2e4f6d6567612e4c696e656172436f6d626f2e636f6f7264696e6174655f6576616c5f33}} in reverse gives \(k = \operatorname{eval}(\operatorname{coordinate}(3),\allowbreak{} \operatorname{Lean\mathord{:}Omega\mathord{:}Coeffs\mathord{:}of\,List}(\langle n\rangle \mathbin{\Vert} (\langle |s|\rangle \mathbin{\Vert} (\langle |t|\rangle \mathbin{\Vert} (\langle k\rangle \mathbin{\Vert} (\langle |u|\rangle \mathbin{\Vert} \langle |u| - |s|\rangle))))))\).
By definition, we obtain \(1 = \operatorname{eval}(\operatorname{Lean\mathord{:}Omega\mathord{:}Linear\,Combo\mathord{:}mk}(1,\allowbreak{} \langle\,\rangle),\allowbreak{} \operatorname{Lean\mathord{:}Omega\mathord{:}Coeffs\mathord{:}of\,List}(\langle n\rangle \mathbin{\Vert} (\langle |s|\rangle \mathbin{\Vert} (\langle |t|\rangle \mathbin{\Vert} (\langle k\rangle \mathbin{\Vert} (\langle |u|\rangle \mathbin{\Vert} \langle |u| - |s|\rangle))))))\).
By \paperlink{premise:4c65616e2e4f6d6567612e496e742e6164645f636f6e6772}{Prerequisite~P\ref{premise:4c65616e2e4f6d6567612e496e742e6164645f636f6e6772}}, we obtain \(k + 1 = \operatorname{eval}(\operatorname{coordinate}(3),\allowbreak{} \operatorname{Lean\mathord{:}Omega\mathord{:}Coeffs\mathord{:}of\,List}(\langle n\rangle \mathbin{\Vert} (\langle |s|\rangle \mathbin{\Vert} (\langle |t|\rangle \mathbin{\Vert} (\langle k\rangle \mathbin{\Vert} (\langle |u|\rangle \mathbin{\Vert} \langle |u| - |s|\rangle)))))) + \operatorname{eval}(\operatorname{Lean\mathord{:}Omega\mathord{:}Linear\,Combo\mathord{:}mk}(1,\allowbreak{} \langle\,\rangle),\allowbreak{} \operatorname{Lean\mathord{:}Omega\mathord{:}Coeffs\mathord{:}of\,List}(\langle n\rangle \mathbin{\Vert} (\langle |s|\rangle \mathbin{\Vert} (\langle |t|\rangle \mathbin{\Vert} (\langle k\rangle \mathbin{\Vert} (\langle |u|\rangle \mathbin{\Vert} \langle |u| - |s|\rangle))))))\).
Using \paperlink{premise:4c65616e2e4f6d6567612e4c696e656172436f6d626f2e6164645f6576616c}{Prerequisite~P\ref{premise:4c65616e2e4f6d6567612e4c696e656172436f6d626f2e6164645f6576616c}} with \(l_{1} = \operatorname{Lean\mathord{:}Omega\mathord{:}Linear\,Combo\mathord{:}mk}(0,\allowbreak{} \langle 0\rangle \mathbin{\Vert} (\langle 0\rangle \mathbin{\Vert} (\langle 0\rangle \mathbin{\Vert} \langle 1\rangle)))\) and \(l_{2} = \operatorname{Lean\mathord{:}Omega\mathord{:}Linear\,Combo\mathord{:}mk}(1,\allowbreak{} \langle\,\rangle)\) and \(v = \operatorname{Lean\mathord{:}Omega\mathord{:}Coeffs\mathord{:}of\,List}(\langle n\rangle \mathbin{\Vert} (\langle |s|\rangle \mathbin{\Vert} (\langle |t|\rangle \mathbin{\Vert} (\langle k\rangle \mathbin{\Vert} (\langle |u|\rangle \mathbin{\Vert} \langle |u| - |s|\rangle)))))\) in reverse gives \(\operatorname{eval}(\operatorname{Lean\mathord{:}Omega\mathord{:}Linear\,Combo\mathord{:}mk}(0,\allowbreak{} \langle 0\rangle \mathbin{\Vert} (\langle 0\rangle \mathbin{\Vert} (\langle 0\rangle \mathbin{\Vert} \langle 1\rangle))),\allowbreak{} \operatorname{Lean\mathord{:}Omega\mathord{:}Coeffs\mathord{:}of\,List}(\langle n\rangle \mathbin{\Vert} (\langle |s|\rangle \mathbin{\Vert} (\langle |t|\rangle \mathbin{\Vert} (\langle k\rangle \mathbin{\Vert} (\langle |u|\rangle \mathbin{\Vert} \langle |u| - |s|\rangle)))))) + \operatorname{eval}(\operatorname{Lean\mathord{:}Omega\mathord{:}Linear\,Combo\mathord{:}mk}(1,\allowbreak{} \langle\,\rangle),\allowbreak{} \operatorname{Lean\mathord{:}Omega\mathord{:}Coeffs\mathord{:}of\,List}(\langle n\rangle \mathbin{\Vert} (\langle |s|\rangle \mathbin{\Vert} (\langle |t|\rangle \mathbin{\Vert} (\langle k\rangle \mathbin{\Vert} (\langle |u|\rangle \mathbin{\Vert} \langle |u| - |s|\rangle)))))) = \operatorname{eval}(\operatorname{Lean\mathord{:}Omega\mathord{:}Linear\,Combo\mathord{:}mk}(0,\allowbreak{} \langle 0\rangle \mathbin{\Vert} (\langle 0\rangle \mathbin{\Vert} (\langle 0\rangle \mathbin{\Vert} \langle 1\rangle))) + \operatorname{Lean\mathord{:}Omega\mathord{:}Linear\,Combo\mathord{:}mk}(1,\allowbreak{} \langle\,\rangle),\allowbreak{} \operatorname{Lean\mathord{:}Omega\mathord{:}Coeffs\mathord{:}of\,List}(\langle n\rangle \mathbin{\Vert} (\langle |s|\rangle \mathbin{\Vert} (\langle |t|\rangle \mathbin{\Vert} (\langle k\rangle \mathbin{\Vert} (\langle |u|\rangle \mathbin{\Vert} \langle |u| - |s|\rangle))))))\).
Chaining the equalities, we obtain \(k + 1 = \operatorname{eval}(\operatorname{Lean\mathord{:}Omega\mathord{:}Linear\,Combo\mathord{:}mk}(0,\allowbreak{} \langle 0\rangle \mathbin{\Vert} (\langle 0\rangle \mathbin{\Vert} (\langle 0\rangle \mathbin{\Vert} \langle 1\rangle))) + \operatorname{Lean\mathord{:}Omega\mathord{:}Linear\,Combo\mathord{:}mk}(1,\allowbreak{} \langle\,\rangle),\allowbreak{} \operatorname{Lean\mathord{:}Omega\mathord{:}Coeffs\mathord{:}of\,List}(\langle n\rangle \mathbin{\Vert} (\langle |s|\rangle \mathbin{\Vert} (\langle |t|\rangle \mathbin{\Vert} (\langle k\rangle \mathbin{\Vert} (\langle |u|\rangle \mathbin{\Vert} \langle |u| - |s|\rangle))))))\).
Chaining the equalities, we obtain \(k + 1 = \operatorname{eval}(\operatorname{Lean\mathord{:}Omega\mathord{:}Linear\,Combo\mathord{:}mk}(0,\allowbreak{} \langle 0\rangle \mathbin{\Vert} (\langle 0\rangle \mathbin{\Vert} (\langle 0\rangle \mathbin{\Vert} \langle 1\rangle))) + \operatorname{Lean\mathord{:}Omega\mathord{:}Linear\,Combo\mathord{:}mk}(1,\allowbreak{} \langle\,\rangle),\allowbreak{} \operatorname{Lean\mathord{:}Omega\mathord{:}Coeffs\mathord{:}of\,List}(\langle n\rangle \mathbin{\Vert} (\langle |s|\rangle \mathbin{\Vert} (\langle |t|\rangle \mathbin{\Vert} (\langle k\rangle \mathbin{\Vert} (\langle |u|\rangle \mathbin{\Vert} \langle |u| - |s|\rangle))))))\).
By \paperlink{premise:4c65616e2e4f6d6567612e496e742e7375625f636f6e6772}{Prerequisite~P\ref{premise:4c65616e2e4f6d6567612e496e742e7375625f636f6e6772}}, we obtain \(|u| - (k + 1) = \operatorname{eval}(\operatorname{coordinate}(4),\allowbreak{} \operatorname{Lean\mathord{:}Omega\mathord{:}Coeffs\mathord{:}of\,List}(\langle n\rangle \mathbin{\Vert} (\langle |s|\rangle \mathbin{\Vert} (\langle |t|\rangle \mathbin{\Vert} (\langle k\rangle \mathbin{\Vert} (\langle |u|\rangle \mathbin{\Vert} \langle |u| - |s|\rangle)))))) - \operatorname{eval}(\operatorname{Lean\mathord{:}Omega\mathord{:}Linear\,Combo\mathord{:}mk}(0,\allowbreak{} \langle 0\rangle \mathbin{\Vert} (\langle 0\rangle \mathbin{\Vert} (\langle 0\rangle \mathbin{\Vert} \langle 1\rangle))) + \operatorname{Lean\mathord{:}Omega\mathord{:}Linear\,Combo\mathord{:}mk}(1,\allowbreak{} \langle\,\rangle),\allowbreak{} \operatorname{Lean\mathord{:}Omega\mathord{:}Coeffs\mathord{:}of\,List}(\langle n\rangle \mathbin{\Vert} (\langle |s|\rangle \mathbin{\Vert} (\langle |t|\rangle \mathbin{\Vert} (\langle k\rangle \mathbin{\Vert} (\langle |u|\rangle \mathbin{\Vert} \langle |u| - |s|\rangle))))))\).
Using \paperlink{premise:4c65616e2e4f6d6567612e4c696e656172436f6d626f2e7375625f6576616c}{Prerequisite~P\ref{premise:4c65616e2e4f6d6567612e4c696e656172436f6d626f2e7375625f6576616c}} with \(l_{1} = \operatorname{Lean\mathord{:}Omega\mathord{:}Linear\,Combo\mathord{:}mk}(0,\allowbreak{} \langle 0\rangle \mathbin{\Vert} (\langle 0\rangle \mathbin{\Vert} (\langle 0\rangle \mathbin{\Vert} (\langle 0\rangle \mathbin{\Vert} \langle 1\rangle))))\) and \(l_{2} = \operatorname{Lean\mathord{:}Omega\mathord{:}Linear\,Combo\mathord{:}mk}(1,\allowbreak{} \langle 0\rangle \mathbin{\Vert} (\langle 0\rangle \mathbin{\Vert} (\langle 0\rangle \mathbin{\Vert} \langle 1\rangle)))\) and \(v = \operatorname{Lean\mathord{:}Omega\mathord{:}Coeffs\mathord{:}of\,List}(\langle n\rangle \mathbin{\Vert} (\langle |s|\rangle \mathbin{\Vert} (\langle |t|\rangle \mathbin{\Vert} (\langle k\rangle \mathbin{\Vert} (\langle |u|\rangle \mathbin{\Vert} \langle |u| - |s|\rangle)))))\) in reverse gives \(\operatorname{eval}(\operatorname{Lean\mathord{:}Omega\mathord{:}Linear\,Combo\mathord{:}mk}(0,\allowbreak{} \langle 0\rangle \mathbin{\Vert} (\langle 0\rangle \mathbin{\Vert} (\langle 0\rangle \mathbin{\Vert} (\langle 0\rangle \mathbin{\Vert} \langle 1\rangle)))),\allowbreak{} \operatorname{Lean\mathord{:}Omega\mathord{:}Coeffs\mathord{:}of\,List}(\langle n\rangle \mathbin{\Vert} (\langle |s|\rangle \mathbin{\Vert} (\langle |t|\rangle \mathbin{\Vert} (\langle k\rangle \mathbin{\Vert} (\langle |u|\rangle \mathbin{\Vert} \langle |u| - |s|\rangle)))))) - \operatorname{eval}(\operatorname{Lean\mathord{:}Omega\mathord{:}Linear\,Combo\mathord{:}mk}(1,\allowbreak{} \langle 0\rangle \mathbin{\Vert} (\langle 0\rangle \mathbin{\Vert} (\langle 0\rangle \mathbin{\Vert} \langle 1\rangle))),\allowbreak{} \operatorname{Lean\mathord{:}Omega\mathord{:}Coeffs\mathord{:}of\,List}(\langle n\rangle \mathbin{\Vert} (\langle |s|\rangle \mathbin{\Vert} (\langle |t|\rangle \mathbin{\Vert} (\langle k\rangle \mathbin{\Vert} (\langle |u|\rangle \mathbin{\Vert} \langle |u| - |s|\rangle)))))) = \operatorname{eval}(\operatorname{Lean\mathord{:}Omega\mathord{:}Linear\,Combo\mathord{:}mk}(0,\allowbreak{} \langle 0\rangle \mathbin{\Vert} (\langle 0\rangle \mathbin{\Vert} (\langle 0\rangle \mathbin{\Vert} (\langle 0\rangle \mathbin{\Vert} \langle 1\rangle)))) - \operatorname{Lean\mathord{:}Omega\mathord{:}Linear\,Combo\mathord{:}mk}(1,\allowbreak{} \langle 0\rangle \mathbin{\Vert} (\langle 0\rangle \mathbin{\Vert} (\langle 0\rangle \mathbin{\Vert} \langle 1\rangle))),\allowbreak{} \operatorname{Lean\mathord{:}Omega\mathord{:}Coeffs\mathord{:}of\,List}(\langle n\rangle \mathbin{\Vert} (\langle |s|\rangle \mathbin{\Vert} (\langle |t|\rangle \mathbin{\Vert} (\langle k\rangle \mathbin{\Vert} (\langle |u|\rangle \mathbin{\Vert} \langle |u| - |s|\rangle))))))\).
Chaining the equalities, we obtain \(|u| - (k + 1) = \operatorname{eval}(\operatorname{Lean\mathord{:}Omega\mathord{:}Linear\,Combo\mathord{:}mk}(0,\allowbreak{} \langle 0\rangle \mathbin{\Vert} (\langle 0\rangle \mathbin{\Vert} (\langle 0\rangle \mathbin{\Vert} (\langle 0\rangle \mathbin{\Vert} \langle 1\rangle)))) - \operatorname{Lean\mathord{:}Omega\mathord{:}Linear\,Combo\mathord{:}mk}(1,\allowbreak{} \langle 0\rangle \mathbin{\Vert} (\langle 0\rangle \mathbin{\Vert} (\langle 0\rangle \mathbin{\Vert} \langle 1\rangle))),\allowbreak{} \operatorname{Lean\mathord{:}Omega\mathord{:}Coeffs\mathord{:}of\,List}(\langle n\rangle \mathbin{\Vert} (\langle |s|\rangle \mathbin{\Vert} (\langle |t|\rangle \mathbin{\Vert} (\langle k\rangle \mathbin{\Vert} (\langle |u|\rangle \mathbin{\Vert} \langle |u| - |s|\rangle))))))\).
Reversing the equality, we obtain \(\operatorname{eval}(\operatorname{Lean\mathord{:}Omega\mathord{:}Linear\,Combo\mathord{:}mk}(0,\allowbreak{} \langle 0\rangle \mathbin{\Vert} (\langle 0\rangle \mathbin{\Vert} (\langle 0\rangle \mathbin{\Vert} (\langle 0\rangle \mathbin{\Vert} \langle 1\rangle)))) - \operatorname{Lean\mathord{:}Omega\mathord{:}Linear\,Combo\mathord{:}mk}(1,\allowbreak{} \langle 0\rangle \mathbin{\Vert} (\langle 0\rangle \mathbin{\Vert} (\langle 0\rangle \mathbin{\Vert} \langle 1\rangle))),\allowbreak{} \operatorname{Lean\mathord{:}Omega\mathord{:}Coeffs\mathord{:}of\,List}(\langle n\rangle \mathbin{\Vert} (\langle |s|\rangle \mathbin{\Vert} (\langle |t|\rangle \mathbin{\Vert} (\langle k\rangle \mathbin{\Vert} (\langle |u|\rangle \mathbin{\Vert} \langle |u| - |s|\rangle)))))) = |u| - (k + 1)\).
Use the hypothesis \(\left|u\right| = k + 1\).
By \paperlink{premise:4c65616e2e4f6d6567612e496e742e6f664e61745f636f6e6772}{Prerequisite~P\ref{premise:4c65616e2e4f6d6567612e496e742e6f664e61745f636f6e6772}}, we obtain \(\left|u\right| = k + 1\).
By \paperlink{premise:496e742e7375625f65715f7a65726f5f6f665f6571}{Prerequisite~P\ref{premise:496e742e7375625f65715f7a65726f5f6f665f6571}}, we obtain \(\left|u\right| - \left(k + 1\right) = 0\).
Chaining the equalities, we obtain \(\operatorname{eval}(\operatorname{Lean\mathord{:}Omega\mathord{:}Linear\,Combo\mathord{:}mk}(0,\allowbreak{} \langle 0\rangle \mathbin{\Vert} (\langle 0\rangle \mathbin{\Vert} (\langle 0\rangle \mathbin{\Vert} (\langle 0\rangle \mathbin{\Vert} \langle 1\rangle)))) - \operatorname{Lean\mathord{:}Omega\mathord{:}Linear\,Combo\mathord{:}mk}(1,\allowbreak{} \langle 0\rangle \mathbin{\Vert} (\langle 0\rangle \mathbin{\Vert} (\langle 0\rangle \mathbin{\Vert} \langle 1\rangle))),\allowbreak{} \operatorname{Lean\mathord{:}Omega\mathord{:}Coeffs\mathord{:}of\,List}(\langle n\rangle \mathbin{\Vert} (\langle |s|\rangle \mathbin{\Vert} (\langle |t|\rangle \mathbin{\Vert} (\langle k\rangle \mathbin{\Vert} (\langle |u|\rangle \mathbin{\Vert} \langle |u| - |s|\rangle)))))) = 0\).
By \paperlink{premise:4c65616e2e4f6d6567612e436f6e73747261696e742e616464457175616c6974795f736174}{Prerequisite~P\ref{premise:4c65616e2e4f6d6567612e436f6e73747261696e742e616464457175616c6974795f736174}}, we obtain \(\operatorname{sat^{\prime}}(\operatorname{Lean\mathord{:}Omega\mathord{:}Constraint\mathord{:}mk}(\operatorname{some}(-(-1)),\allowbreak{} \operatorname{some}(-(-1))),\allowbreak{} \langle 0\rangle \mathbin{\Vert} (\langle 0\rangle \mathbin{\Vert} (\langle 0\rangle \mathbin{\Vert} (\langle -1\rangle \mathbin{\Vert} \langle 1\rangle))),\allowbreak{} \operatorname{Lean\mathord{:}Omega\mathord{:}Coeffs\mathord{:}of\,List}(\langle n\rangle \mathbin{\Vert} (\langle |s|\rangle \mathbin{\Vert} (\langle |t|\rangle \mathbin{\Vert} (\langle k\rangle \mathbin{\Vert} (\langle |u|\rangle \mathbin{\Vert} \langle |u| - |s|\rangle)))))) = \mathrm{true}\).
By \paperlink{premise:4c65616e2e4f6d6567612e746964795f736174}{Prerequisite~P\ref{premise:4c65616e2e4f6d6567612e746964795f736174}}, we obtain \(\operatorname{sat^{\prime}}(\operatorname{tidy\,Constraint}(\operatorname{Lean\mathord{:}Omega\mathord{:}Constraint\mathord{:}mk}(\operatorname{some}(1),\allowbreak{} \operatorname{some}(1)),\allowbreak{} \langle 0\rangle \mathbin{\Vert} (\langle 0\rangle \mathbin{\Vert} (\langle 0\rangle \mathbin{\Vert} (\langle -1\rangle \mathbin{\Vert} \langle 1\rangle)))),\allowbreak{} \operatorname{tidy\,Coeffs}(\operatorname{Lean\mathord{:}Omega\mathord{:}Constraint\mathord{:}mk}(\operatorname{some}(1),\allowbreak{} \operatorname{some}(1)),\allowbreak{} \langle 0\rangle \mathbin{\Vert} (\langle 0\rangle \mathbin{\Vert} (\langle 0\rangle \mathbin{\Vert} (\langle -1\rangle \mathbin{\Vert} \langle 1\rangle)))),\allowbreak{} \operatorname{Lean\mathord{:}Omega\mathord{:}Coeffs\mathord{:}of\,List}(\langle n\rangle \mathbin{\Vert} (\langle |s|\rangle \mathbin{\Vert} (\langle |t|\rangle \mathbin{\Vert} (\langle k\rangle \mathbin{\Vert} (\langle |u|\rangle \mathbin{\Vert} \langle |u| - |s|\rangle)))))) = \mathrm{true}\).
Use the hypothesis \(n \le k\).
By \paperlink{premise:4c65616e2e4f6d6567612e496e742e6f664e61745f6c655f6f665f6c65}{Prerequisite~P\ref{premise:4c65616e2e4f6d6567612e496e742e6f664e61745f6c655f6f665f6c65}}, we obtain \(n \le k\).
By \paperlink{premise:496e742e7375625f6e6f6e6e65675f6f665f6c65}{Prerequisite~P\ref{premise:496e742e7375625f6e6f6e6e65675f6f665f6c65}}, we obtain \(0 \le k - n\).
Using \paperlink{premise:4c65616e2e4f6d6567612e4c696e656172436f6d626f2e636f6f7264696e6174655f6576616c5f33}{Prerequisite~P\ref{premise:4c65616e2e4f6d6567612e4c696e656172436f6d626f2e636f6f7264696e6174655f6576616c5f33}} in reverse gives \(k = \operatorname{eval}(\operatorname{coordinate}(3),\allowbreak{} \operatorname{Lean\mathord{:}Omega\mathord{:}Coeffs\mathord{:}of\,List}(\langle n\rangle \mathbin{\Vert} (\langle |s|\rangle \mathbin{\Vert} (\langle |t|\rangle \mathbin{\Vert} (\langle k\rangle \mathbin{\Vert} (\langle |u|\rangle \mathbin{\Vert} \langle |u| - |s|\rangle))))))\).
Using \paperlink{premise:4c65616e2e4f6d6567612e4c696e656172436f6d626f2e636f6f7264696e6174655f6576616c5f30}{Prerequisite~P\ref{premise:4c65616e2e4f6d6567612e4c696e656172436f6d626f2e636f6f7264696e6174655f6576616c5f30}} in reverse gives \(n = \operatorname{eval}(\operatorname{coordinate}(0),\allowbreak{} \operatorname{Lean\mathord{:}Omega\mathord{:}Coeffs\mathord{:}of\,List}(\langle n\rangle \mathbin{\Vert} (\langle |s|\rangle \mathbin{\Vert} (\langle |t|\rangle \mathbin{\Vert} (\langle k\rangle \mathbin{\Vert} (\langle |u|\rangle \mathbin{\Vert} \langle |u| - |s|\rangle))))))\).
By \paperlink{premise:4c65616e2e4f6d6567612e496e742e7375625f636f6e6772}{Prerequisite~P\ref{premise:4c65616e2e4f6d6567612e496e742e7375625f636f6e6772}}, we obtain \(k - n = \operatorname{eval}(\operatorname{coordinate}(3),\allowbreak{} \operatorname{Lean\mathord{:}Omega\mathord{:}Coeffs\mathord{:}of\,List}(\langle n\rangle \mathbin{\Vert} (\langle |s|\rangle \mathbin{\Vert} (\langle |t|\rangle \mathbin{\Vert} (\langle k\rangle \mathbin{\Vert} (\langle |u|\rangle \mathbin{\Vert} \langle |u| - |s|\rangle)))))) - \operatorname{eval}(\operatorname{coordinate}(0),\allowbreak{} \operatorname{Lean\mathord{:}Omega\mathord{:}Coeffs\mathord{:}of\,List}(\langle n\rangle \mathbin{\Vert} (\langle |s|\rangle \mathbin{\Vert} (\langle |t|\rangle \mathbin{\Vert} (\langle k\rangle \mathbin{\Vert} (\langle |u|\rangle \mathbin{\Vert} \langle |u| - |s|\rangle))))))\).
Using \paperlink{premise:4c65616e2e4f6d6567612e4c696e656172436f6d626f2e7375625f6576616c}{Prerequisite~P\ref{premise:4c65616e2e4f6d6567612e4c696e656172436f6d626f2e7375625f6576616c}} with \(l_{1} = \operatorname{Lean\mathord{:}Omega\mathord{:}Linear\,Combo\mathord{:}mk}(0,\allowbreak{} \langle 0\rangle \mathbin{\Vert} (\langle 0\rangle \mathbin{\Vert} (\langle 0\rangle \mathbin{\Vert} \langle 1\rangle)))\) and \(l_{2} = \operatorname{Lean\mathord{:}Omega\mathord{:}Linear\,Combo\mathord{:}mk}(0,\allowbreak{} \langle 1\rangle)\) and \(v = \operatorname{Lean\mathord{:}Omega\mathord{:}Coeffs\mathord{:}of\,List}(\langle n\rangle \mathbin{\Vert} (\langle |s|\rangle \mathbin{\Vert} (\langle |t|\rangle \mathbin{\Vert} (\langle k\rangle \mathbin{\Vert} (\langle |u|\rangle \mathbin{\Vert} \langle |u| - |s|\rangle)))))\) in reverse gives \(\operatorname{eval}(\operatorname{Lean\mathord{:}Omega\mathord{:}Linear\,Combo\mathord{:}mk}(0,\allowbreak{} \langle 0\rangle \mathbin{\Vert} (\langle 0\rangle \mathbin{\Vert} (\langle 0\rangle \mathbin{\Vert} \langle 1\rangle))),\allowbreak{} \operatorname{Lean\mathord{:}Omega\mathord{:}Coeffs\mathord{:}of\,List}(\langle n\rangle \mathbin{\Vert} (\langle |s|\rangle \mathbin{\Vert} (\langle |t|\rangle \mathbin{\Vert} (\langle k\rangle \mathbin{\Vert} (\langle |u|\rangle \mathbin{\Vert} \langle |u| - |s|\rangle)))))) - \operatorname{eval}(\operatorname{Lean\mathord{:}Omega\mathord{:}Linear\,Combo\mathord{:}mk}(0,\allowbreak{} \langle 1\rangle),\allowbreak{} \operatorname{Lean\mathord{:}Omega\mathord{:}Coeffs\mathord{:}of\,List}(\langle n\rangle \mathbin{\Vert} (\langle |s|\rangle \mathbin{\Vert} (\langle |t|\rangle \mathbin{\Vert} (\langle k\rangle \mathbin{\Vert} (\langle |u|\rangle \mathbin{\Vert} \langle |u| - |s|\rangle)))))) = \operatorname{eval}(\operatorname{Lean\mathord{:}Omega\mathord{:}Linear\,Combo\mathord{:}mk}(0,\allowbreak{} \langle 0\rangle \mathbin{\Vert} (\langle 0\rangle \mathbin{\Vert} (\langle 0\rangle \mathbin{\Vert} \langle 1\rangle))) - \operatorname{Lean\mathord{:}Omega\mathord{:}Linear\,Combo\mathord{:}mk}(0,\allowbreak{} \langle 1\rangle),\allowbreak{} \operatorname{Lean\mathord{:}Omega\mathord{:}Coeffs\mathord{:}of\,List}(\langle n\rangle \mathbin{\Vert} (\langle |s|\rangle \mathbin{\Vert} (\langle |t|\rangle \mathbin{\Vert} (\langle k\rangle \mathbin{\Vert} (\langle |u|\rangle \mathbin{\Vert} \langle |u| - |s|\rangle))))))\).
Chaining the equalities, we obtain \(k - n = \operatorname{eval}(\operatorname{Lean\mathord{:}Omega\mathord{:}Linear\,Combo\mathord{:}mk}(0,\allowbreak{} \langle 0\rangle \mathbin{\Vert} (\langle 0\rangle \mathbin{\Vert} (\langle 0\rangle \mathbin{\Vert} \langle 1\rangle))) - \operatorname{Lean\mathord{:}Omega\mathord{:}Linear\,Combo\mathord{:}mk}(0,\allowbreak{} \langle 1\rangle),\allowbreak{} \operatorname{Lean\mathord{:}Omega\mathord{:}Coeffs\mathord{:}of\,List}(\langle n\rangle \mathbin{\Vert} (\langle |s|\rangle \mathbin{\Vert} (\langle |t|\rangle \mathbin{\Vert} (\langle k\rangle \mathbin{\Vert} (\langle |u|\rangle \mathbin{\Vert} \langle |u| - |s|\rangle))))))\).
By \paperlink{premise:6c655f6f665f6c655f6f665f6571}{Prerequisite~P\ref{premise:6c655f6f665f6c655f6f665f6571}} with \(\mathbb{Z}\) as its type parameter \(\alpha\), we obtain \(0 \le \operatorname{eval}(\operatorname{Lean\mathord{:}Omega\mathord{:}Linear\,Combo\mathord{:}mk}(0,\allowbreak{} \langle 0\rangle \mathbin{\Vert} (\langle 0\rangle \mathbin{\Vert} (\langle 0\rangle \mathbin{\Vert} \langle 1\rangle))) - \operatorname{Lean\mathord{:}Omega\mathord{:}Linear\,Combo\mathord{:}mk}(0,\allowbreak{} \langle 1\rangle),\allowbreak{} \operatorname{Lean\mathord{:}Omega\mathord{:}Coeffs\mathord{:}of\,List}(\langle n\rangle \mathbin{\Vert} (\langle |s|\rangle \mathbin{\Vert} (\langle |t|\rangle \mathbin{\Vert} (\langle k\rangle \mathbin{\Vert} (\langle |u|\rangle \mathbin{\Vert} \langle |u| - |s|\rangle))))))\).
By \paperlink{premise:4c65616e2e4f6d6567612e436f6e73747261696e742e616464496e657175616c6974795f736174}{Prerequisite~P\ref{premise:4c65616e2e4f6d6567612e436f6e73747261696e742e616464496e657175616c6974795f736174}}, we obtain \(\operatorname{sat^{\prime}}(\operatorname{Lean\mathord{:}Omega\mathord{:}Constraint\mathord{:}mk}(\operatorname{some}(-0),\allowbreak{} \operatorname{none}),\allowbreak{} \langle -1\rangle \mathbin{\Vert} (\langle 0\rangle \mathbin{\Vert} (\langle 0\rangle \mathbin{\Vert} \langle 1\rangle)),\allowbreak{} \operatorname{Lean\mathord{:}Omega\mathord{:}Coeffs\mathord{:}of\,List}(\langle n\rangle \mathbin{\Vert} (\langle |s|\rangle \mathbin{\Vert} (\langle |t|\rangle \mathbin{\Vert} (\langle k\rangle \mathbin{\Vert} (\langle |u|\rangle \mathbin{\Vert} \langle |u| - |s|\rangle)))))) = \mathrm{true}\).
By \paperlink{premise:4c65616e2e4f6d6567612e746964795f736174}{Prerequisite~P\ref{premise:4c65616e2e4f6d6567612e746964795f736174}}, we obtain \(\operatorname{sat^{\prime}}(\operatorname{tidy\,Constraint}(\operatorname{Lean\mathord{:}Omega\mathord{:}Constraint\mathord{:}mk}(\operatorname{some}(0),\allowbreak{} \operatorname{none}),\allowbreak{} \langle -1\rangle \mathbin{\Vert} (\langle 0\rangle \mathbin{\Vert} (\langle 0\rangle \mathbin{\Vert} \langle 1\rangle))),\allowbreak{} \operatorname{tidy\,Coeffs}(\operatorname{Lean\mathord{:}Omega\mathord{:}Constraint\mathord{:}mk}(\operatorname{some}(0),\allowbreak{} \operatorname{none}),\allowbreak{} \langle -1\rangle \mathbin{\Vert} (\langle 0\rangle \mathbin{\Vert} (\langle 0\rangle \mathbin{\Vert} \langle 1\rangle))),\allowbreak{} \operatorname{Lean\mathord{:}Omega\mathord{:}Coeffs\mathord{:}of\,List}(\langle n\rangle \mathbin{\Vert} (\langle |s|\rangle \mathbin{\Vert} (\langle |t|\rangle \mathbin{\Vert} (\langle k\rangle \mathbin{\Vert} (\langle |u|\rangle \mathbin{\Vert} \langle |u| - |s|\rangle)))))) = \mathrm{true}\).
By \paperlink{premise:4c65616e2e4f6d6567612e636f6d626f5f73617427}{Prerequisite~P\ref{premise:4c65616e2e4f6d6567612e636f6d626f5f73617427}} with \(\operatorname{Lean\mathord{:}Omega\mathord{:}Constraint\mathord{:}mk}(\operatorname{some}(-1),\allowbreak{} \operatorname{some}(-1))\) in place of its parameter \(s\) and \(\operatorname{Lean\mathord{:}Omega\mathord{:}Constraint\mathord{:}mk}(\operatorname{none},\allowbreak{} \operatorname{some}(0))\) in place of its parameter \(t\) and \(1\) in place of its parameter \(a\) and \(x = \langle 0\rangle \mathbin{\Vert} (\langle 0\rangle \mathbin{\Vert} (\langle 0\rangle \mathbin{\Vert} (\langle 1\rangle \mathbin{\Vert} \langle -1\rangle)))\) and \(b = 1\) and \(y = \langle 1\rangle \mathbin{\Vert} (\langle 0\rangle \mathbin{\Vert} (\langle 0\rangle \mathbin{\Vert} \langle -1\rangle))\) and \(v = \operatorname{Lean\mathord{:}Omega\mathord{:}Coeffs\mathord{:}of\,List}(\langle n\rangle \mathbin{\Vert} (\langle |s|\rangle \mathbin{\Vert} (\langle |t|\rangle \mathbin{\Vert} (\langle k\rangle \mathbin{\Vert} (\langle |u|\rangle \mathbin{\Vert} \langle |u| - |s|\rangle)))))\), we obtain \(\operatorname{sat^{\prime}}(\operatorname{Lean\mathord{:}Omega\mathord{:}Constraint\mathord{:}combo}(1,\allowbreak{} \operatorname{Lean\mathord{:}Omega\mathord{:}Constraint\mathord{:}mk}(\operatorname{some}(-1),\allowbreak{} \operatorname{some}(-1)),\allowbreak{} 1,\allowbreak{} \operatorname{Lean\mathord{:}Omega\mathord{:}Constraint\mathord{:}mk}(\operatorname{none},\allowbreak{} \operatorname{some}(0))),\allowbreak{} \operatorname{Lean\mathord{:}Omega\mathord{:}Coeffs\mathord{:}combo}(1,\allowbreak{} \langle 0\rangle \mathbin{\Vert} (\langle 0\rangle \mathbin{\Vert} (\langle 0\rangle \mathbin{\Vert} (\langle 1\rangle \mathbin{\Vert} \langle -1\rangle))),\allowbreak{} 1,\allowbreak{} \langle 1\rangle \mathbin{\Vert} (\langle 0\rangle \mathbin{\Vert} (\langle 0\rangle \mathbin{\Vert} \langle -1\rangle))),\allowbreak{} \operatorname{Lean\mathord{:}Omega\mathord{:}Coeffs\mathord{:}of\,List}(\langle n\rangle \mathbin{\Vert} (\langle |s|\rangle \mathbin{\Vert} (\langle |t|\rangle \mathbin{\Vert} (\langle k\rangle \mathbin{\Vert} (\langle |u|\rangle \mathbin{\Vert} \langle |u| - |s|\rangle)))))) = \mathrm{true}\).
By \paperlink{premise:4c65616e2e4f6d6567612e636f6d626f5f73617427}{Prerequisite~P\ref{premise:4c65616e2e4f6d6567612e636f6d626f5f73617427}} with \(\operatorname{Lean\mathord{:}Omega\mathord{:}Constraint\mathord{:}mk}(\operatorname{some}(0),\allowbreak{} \operatorname{none})\) in place of its parameter \(s\) and \(\operatorname{Lean\mathord{:}Omega\mathord{:}Constraint\mathord{:}mk}(\operatorname{none},\allowbreak{} \operatorname{some}(-1))\) in place of its parameter \(t\) and \(1\) in place of its parameter \(a\) and \(x = \langle 1\rangle \mathbin{\Vert} \langle -1\rangle\) and \(b = -1\) and \(y = \langle 1\rangle \mathbin{\Vert} (\langle 0\rangle \mathbin{\Vert} (\langle 0\rangle \mathbin{\Vert} (\langle 0\rangle \mathbin{\Vert} \langle -1\rangle)))\) and \(v = \operatorname{Lean\mathord{:}Omega\mathord{:}Coeffs\mathord{:}of\,List}(\langle n\rangle \mathbin{\Vert} (\langle |s|\rangle \mathbin{\Vert} (\langle |t|\rangle \mathbin{\Vert} (\langle k\rangle \mathbin{\Vert} (\langle |u|\rangle \mathbin{\Vert} \langle |u| - |s|\rangle)))))\), we obtain \(\operatorname{sat^{\prime}}(\operatorname{Lean\mathord{:}Omega\mathord{:}Constraint\mathord{:}combo}(1,\allowbreak{} \operatorname{Lean\mathord{:}Omega\mathord{:}Constraint\mathord{:}mk}(\operatorname{some}(0),\allowbreak{} \operatorname{none}),\allowbreak{} -1,\allowbreak{} \operatorname{Lean\mathord{:}Omega\mathord{:}Constraint\mathord{:}mk}(\operatorname{none},\allowbreak{} \operatorname{some}(-1))),\allowbreak{} \operatorname{Lean\mathord{:}Omega\mathord{:}Coeffs\mathord{:}combo}(1,\allowbreak{} \langle 1\rangle \mathbin{\Vert} \langle -1\rangle,\allowbreak{} -1,\allowbreak{} \langle 1\rangle \mathbin{\Vert} (\langle 0\rangle \mathbin{\Vert} (\langle 0\rangle \mathbin{\Vert} (\langle 0\rangle \mathbin{\Vert} \langle -1\rangle)))),\allowbreak{} \operatorname{Lean\mathord{:}Omega\mathord{:}Coeffs\mathord{:}of\,List}(\langle n\rangle \mathbin{\Vert} (\langle |s|\rangle \mathbin{\Vert} (\langle |t|\rangle \mathbin{\Vert} (\langle k\rangle \mathbin{\Vert} (\langle |u|\rangle \mathbin{\Vert} \langle |u| - |s|\rangle)))))) = \mathrm{true}\).
By \paperlink{premise:4c65616e2e4f6d6567612e746964795f736174}{Prerequisite~P\ref{premise:4c65616e2e4f6d6567612e746964795f736174}}, we obtain \(\operatorname{sat^{\prime}}(\operatorname{tidy\,Constraint}(\operatorname{Lean\mathord{:}Omega\mathord{:}Constraint\mathord{:}mk}(\operatorname{some}(1),\allowbreak{} \operatorname{none}),\allowbreak{} \langle 0\rangle \mathbin{\Vert} (\langle -1\rangle \mathbin{\Vert} (\langle 0\rangle \mathbin{\Vert} (\langle 0\rangle \mathbin{\Vert} \langle 1\rangle)))),\allowbreak{} \operatorname{tidy\,Coeffs}(\operatorname{Lean\mathord{:}Omega\mathord{:}Constraint\mathord{:}mk}(\operatorname{some}(1),\allowbreak{} \operatorname{none}),\allowbreak{} \langle 0\rangle \mathbin{\Vert} (\langle -1\rangle \mathbin{\Vert} (\langle 0\rangle \mathbin{\Vert} (\langle 0\rangle \mathbin{\Vert} \langle 1\rangle)))),\allowbreak{} \operatorname{Lean\mathord{:}Omega\mathord{:}Coeffs\mathord{:}of\,List}(\langle n\rangle \mathbin{\Vert} (\langle |s|\rangle \mathbin{\Vert} (\langle |t|\rangle \mathbin{\Vert} (\langle k\rangle \mathbin{\Vert} (\langle |u|\rangle \mathbin{\Vert} \langle |u| - |s|\rangle)))))) = \mathrm{true}\).
Use the hypothesis \(\left|u\right| < \left|s\right| \land \left|u\right| - \left|s\right| = 0\).
Taking the first conjunct, we obtain \(\left|u\right| < \left|s\right|\).
By \paperlink{premise:4c65616e2e4f6d6567612e496e742e6f664e61745f6c745f6f665f6c74}{Prerequisite~P\ref{premise:4c65616e2e4f6d6567612e496e742e6f664e61745f6c745f6f665f6c74}}, we obtain \(\left|u\right| < \left|s\right|\).
By \paperlink{premise:496e742e6164645f6f6e655f6c655f6f665f6c74}{Prerequisite~P\ref{premise:496e742e6164645f6f6e655f6c655f6f665f6c74}}, we obtain \(\left|u\right| + 1 \le \left|s\right|\).
By \paperlink{premise:496e742e7375625f6e6f6e6e65675f6f665f6c65}{Prerequisite~P\ref{premise:496e742e7375625f6e6f6e6e65675f6f665f6c65}}, we obtain \(0 \le \left|s\right| - \left(\left|u\right| + 1\right)\).
Using \paperlink{premise:4c65616e2e4f6d6567612e4c696e656172436f6d626f2e636f6f7264696e6174655f6576616c5f31}{Prerequisite~P\ref{premise:4c65616e2e4f6d6567612e4c696e656172436f6d626f2e636f6f7264696e6174655f6576616c5f31}} in reverse gives \(|s| = \operatorname{eval}(\operatorname{coordinate}(1),\allowbreak{} \operatorname{Lean\mathord{:}Omega\mathord{:}Coeffs\mathord{:}of\,List}(\langle n\rangle \mathbin{\Vert} (\langle |s|\rangle \mathbin{\Vert} (\langle |t|\rangle \mathbin{\Vert} (\langle k\rangle \mathbin{\Vert} (\langle |u|\rangle \mathbin{\Vert} \langle |u| - |s|\rangle))))))\).
Using \paperlink{premise:4c65616e2e4f6d6567612e4c696e656172436f6d626f2e636f6f7264696e6174655f6576616c5f34}{Prerequisite~P\ref{premise:4c65616e2e4f6d6567612e4c696e656172436f6d626f2e636f6f7264696e6174655f6576616c5f34}} in reverse gives \(|u| = \operatorname{eval}(\operatorname{coordinate}(4),\allowbreak{} \operatorname{Lean\mathord{:}Omega\mathord{:}Coeffs\mathord{:}of\,List}(\langle n\rangle \mathbin{\Vert} (\langle |s|\rangle \mathbin{\Vert} (\langle |t|\rangle \mathbin{\Vert} (\langle k\rangle \mathbin{\Vert} (\langle |u|\rangle \mathbin{\Vert} \langle |u| - |s|\rangle))))))\).
By definition, we obtain \(1 = \operatorname{eval}(\operatorname{Lean\mathord{:}Omega\mathord{:}Linear\,Combo\mathord{:}mk}(1,\allowbreak{} \langle\,\rangle),\allowbreak{} \operatorname{Lean\mathord{:}Omega\mathord{:}Coeffs\mathord{:}of\,List}(\langle n\rangle \mathbin{\Vert} (\langle |s|\rangle \mathbin{\Vert} (\langle |t|\rangle \mathbin{\Vert} (\langle k\rangle \mathbin{\Vert} (\langle |u|\rangle \mathbin{\Vert} \langle |u| - |s|\rangle))))))\).
By \paperlink{premise:4c65616e2e4f6d6567612e496e742e6164645f636f6e6772}{Prerequisite~P\ref{premise:4c65616e2e4f6d6567612e496e742e6164645f636f6e6772}}, we obtain \(|u| + 1 = \operatorname{eval}(\operatorname{coordinate}(4),\allowbreak{} \operatorname{Lean\mathord{:}Omega\mathord{:}Coeffs\mathord{:}of\,List}(\langle n\rangle \mathbin{\Vert} (\langle |s|\rangle \mathbin{\Vert} (\langle |t|\rangle \mathbin{\Vert} (\langle k\rangle \mathbin{\Vert} (\langle |u|\rangle \mathbin{\Vert} \langle |u| - |s|\rangle)))))) + \operatorname{eval}(\operatorname{Lean\mathord{:}Omega\mathord{:}Linear\,Combo\mathord{:}mk}(1,\allowbreak{} \langle\,\rangle),\allowbreak{} \operatorname{Lean\mathord{:}Omega\mathord{:}Coeffs\mathord{:}of\,List}(\langle n\rangle \mathbin{\Vert} (\langle |s|\rangle \mathbin{\Vert} (\langle |t|\rangle \mathbin{\Vert} (\langle k\rangle \mathbin{\Vert} (\langle |u|\rangle \mathbin{\Vert} \langle |u| - |s|\rangle))))))\).
Using \paperlink{premise:4c65616e2e4f6d6567612e4c696e656172436f6d626f2e6164645f6576616c}{Prerequisite~P\ref{premise:4c65616e2e4f6d6567612e4c696e656172436f6d626f2e6164645f6576616c}} with \(l_{1} = \operatorname{Lean\mathord{:}Omega\mathord{:}Linear\,Combo\mathord{:}mk}(0,\allowbreak{} \langle 0\rangle \mathbin{\Vert} (\langle 0\rangle \mathbin{\Vert} (\langle 0\rangle \mathbin{\Vert} (\langle 0\rangle \mathbin{\Vert} \langle 1\rangle))))\) and \(l_{2} = \operatorname{Lean\mathord{:}Omega\mathord{:}Linear\,Combo\mathord{:}mk}(1,\allowbreak{} \langle\,\rangle)\) and \(v = \operatorname{Lean\mathord{:}Omega\mathord{:}Coeffs\mathord{:}of\,List}(\langle n\rangle \mathbin{\Vert} (\langle |s|\rangle \mathbin{\Vert} (\langle |t|\rangle \mathbin{\Vert} (\langle k\rangle \mathbin{\Vert} (\langle |u|\rangle \mathbin{\Vert} \langle |u| - |s|\rangle)))))\) in reverse gives \(\operatorname{eval}(\operatorname{Lean\mathord{:}Omega\mathord{:}Linear\,Combo\mathord{:}mk}(0,\allowbreak{} \langle 0\rangle \mathbin{\Vert} (\langle 0\rangle \mathbin{\Vert} (\langle 0\rangle \mathbin{\Vert} (\langle 0\rangle \mathbin{\Vert} \langle 1\rangle)))),\allowbreak{} \operatorname{Lean\mathord{:}Omega\mathord{:}Coeffs\mathord{:}of\,List}(\langle n\rangle \mathbin{\Vert} (\langle |s|\rangle \mathbin{\Vert} (\langle |t|\rangle \mathbin{\Vert} (\langle k\rangle \mathbin{\Vert} (\langle |u|\rangle \mathbin{\Vert} \langle |u| - |s|\rangle)))))) + \operatorname{eval}(\operatorname{Lean\mathord{:}Omega\mathord{:}Linear\,Combo\mathord{:}mk}(1,\allowbreak{} \langle\,\rangle),\allowbreak{} \operatorname{Lean\mathord{:}Omega\mathord{:}Coeffs\mathord{:}of\,List}(\langle n\rangle \mathbin{\Vert} (\langle |s|\rangle \mathbin{\Vert} (\langle |t|\rangle \mathbin{\Vert} (\langle k\rangle \mathbin{\Vert} (\langle |u|\rangle \mathbin{\Vert} \langle |u| - |s|\rangle)))))) = \operatorname{eval}(\operatorname{Lean\mathord{:}Omega\mathord{:}Linear\,Combo\mathord{:}mk}(0,\allowbreak{} \langle 0\rangle \mathbin{\Vert} (\langle 0\rangle \mathbin{\Vert} (\langle 0\rangle \mathbin{\Vert} (\langle 0\rangle \mathbin{\Vert} \langle 1\rangle)))) + \operatorname{Lean\mathord{:}Omega\mathord{:}Linear\,Combo\mathord{:}mk}(1,\allowbreak{} \langle\,\rangle),\allowbreak{} \operatorname{Lean\mathord{:}Omega\mathord{:}Coeffs\mathord{:}of\,List}(\langle n\rangle \mathbin{\Vert} (\langle |s|\rangle \mathbin{\Vert} (\langle |t|\rangle \mathbin{\Vert} (\langle k\rangle \mathbin{\Vert} (\langle |u|\rangle \mathbin{\Vert} \langle |u| - |s|\rangle))))))\).
Chaining the equalities, we obtain \(|u| + 1 = \operatorname{eval}(\operatorname{Lean\mathord{:}Omega\mathord{:}Linear\,Combo\mathord{:}mk}(0,\allowbreak{} \langle 0\rangle \mathbin{\Vert} (\langle 0\rangle \mathbin{\Vert} (\langle 0\rangle \mathbin{\Vert} (\langle 0\rangle \mathbin{\Vert} \langle 1\rangle)))) + \operatorname{Lean\mathord{:}Omega\mathord{:}Linear\,Combo\mathord{:}mk}(1,\allowbreak{} \langle\,\rangle),\allowbreak{} \operatorname{Lean\mathord{:}Omega\mathord{:}Coeffs\mathord{:}of\,List}(\langle n\rangle \mathbin{\Vert} (\langle |s|\rangle \mathbin{\Vert} (\langle |t|\rangle \mathbin{\Vert} (\langle k\rangle \mathbin{\Vert} (\langle |u|\rangle \mathbin{\Vert} \langle |u| - |s|\rangle))))))\).
By \paperlink{premise:4c65616e2e4f6d6567612e496e742e7375625f636f6e6772}{Prerequisite~P\ref{premise:4c65616e2e4f6d6567612e496e742e7375625f636f6e6772}}, we obtain \(|s| - (|u| + 1) = \operatorname{eval}(\operatorname{coordinate}(1),\allowbreak{} \operatorname{Lean\mathord{:}Omega\mathord{:}Coeffs\mathord{:}of\,List}(\langle n\rangle \mathbin{\Vert} (\langle |s|\rangle \mathbin{\Vert} (\langle |t|\rangle \mathbin{\Vert} (\langle k\rangle \mathbin{\Vert} (\langle |u|\rangle \mathbin{\Vert} \langle |u| - |s|\rangle)))))) - \operatorname{eval}(\operatorname{Lean\mathord{:}Omega\mathord{:}Linear\,Combo\mathord{:}mk}(0,\allowbreak{} \langle 0\rangle \mathbin{\Vert} (\langle 0\rangle \mathbin{\Vert} (\langle 0\rangle \mathbin{\Vert} (\langle 0\rangle \mathbin{\Vert} \langle 1\rangle)))) + \operatorname{Lean\mathord{:}Omega\mathord{:}Linear\,Combo\mathord{:}mk}(1,\allowbreak{} \langle\,\rangle),\allowbreak{} \operatorname{Lean\mathord{:}Omega\mathord{:}Coeffs\mathord{:}of\,List}(\langle n\rangle \mathbin{\Vert} (\langle |s|\rangle \mathbin{\Vert} (\langle |t|\rangle \mathbin{\Vert} (\langle k\rangle \mathbin{\Vert} (\langle |u|\rangle \mathbin{\Vert} \langle |u| - |s|\rangle))))))\).
Using \paperlink{premise:4c65616e2e4f6d6567612e4c696e656172436f6d626f2e7375625f6576616c}{Prerequisite~P\ref{premise:4c65616e2e4f6d6567612e4c696e656172436f6d626f2e7375625f6576616c}} with \(l_{1} = \operatorname{Lean\mathord{:}Omega\mathord{:}Linear\,Combo\mathord{:}mk}(0,\allowbreak{} \langle 0\rangle \mathbin{\Vert} \langle 1\rangle)\) and \(l_{2} = \operatorname{Lean\mathord{:}Omega\mathord{:}Linear\,Combo\mathord{:}mk}(1,\allowbreak{} \langle 0\rangle \mathbin{\Vert} (\langle 0\rangle \mathbin{\Vert} (\langle 0\rangle \mathbin{\Vert} (\langle 0\rangle \mathbin{\Vert} \langle 1\rangle))))\) and \(v = \operatorname{Lean\mathord{:}Omega\mathord{:}Coeffs\mathord{:}of\,List}(\langle n\rangle \mathbin{\Vert} (\langle |s|\rangle \mathbin{\Vert} (\langle |t|\rangle \mathbin{\Vert} (\langle k\rangle \mathbin{\Vert} (\langle |u|\rangle \mathbin{\Vert} \langle |u| - |s|\rangle)))))\) in reverse gives \(\operatorname{eval}(\operatorname{Lean\mathord{:}Omega\mathord{:}Linear\,Combo\mathord{:}mk}(0,\allowbreak{} \langle 0\rangle \mathbin{\Vert} \langle 1\rangle),\allowbreak{} \operatorname{Lean\mathord{:}Omega\mathord{:}Coeffs\mathord{:}of\,List}(\langle n\rangle \mathbin{\Vert} (\langle |s|\rangle \mathbin{\Vert} (\langle |t|\rangle \mathbin{\Vert} (\langle k\rangle \mathbin{\Vert} (\langle |u|\rangle \mathbin{\Vert} \langle |u| - |s|\rangle)))))) - \operatorname{eval}(\operatorname{Lean\mathord{:}Omega\mathord{:}Linear\,Combo\mathord{:}mk}(1,\allowbreak{} \langle 0\rangle \mathbin{\Vert} (\langle 0\rangle \mathbin{\Vert} (\langle 0\rangle \mathbin{\Vert} (\langle 0\rangle \mathbin{\Vert} \langle 1\rangle)))),\allowbreak{} \operatorname{Lean\mathord{:}Omega\mathord{:}Coeffs\mathord{:}of\,List}(\langle n\rangle \mathbin{\Vert} (\langle |s|\rangle \mathbin{\Vert} (\langle |t|\rangle \mathbin{\Vert} (\langle k\rangle \mathbin{\Vert} (\langle |u|\rangle \mathbin{\Vert} \langle |u| - |s|\rangle)))))) = \operatorname{eval}(\operatorname{Lean\mathord{:}Omega\mathord{:}Linear\,Combo\mathord{:}mk}(0,\allowbreak{} \langle 0\rangle \mathbin{\Vert} \langle 1\rangle) - \operatorname{Lean\mathord{:}Omega\mathord{:}Linear\,Combo\mathord{:}mk}(1,\allowbreak{} \langle 0\rangle \mathbin{\Vert} (\langle 0\rangle \mathbin{\Vert} (\langle 0\rangle \mathbin{\Vert} (\langle 0\rangle \mathbin{\Vert} \langle 1\rangle)))),\allowbreak{} \operatorname{Lean\mathord{:}Omega\mathord{:}Coeffs\mathord{:}of\,List}(\langle n\rangle \mathbin{\Vert} (\langle |s|\rangle \mathbin{\Vert} (\langle |t|\rangle \mathbin{\Vert} (\langle k\rangle \mathbin{\Vert} (\langle |u|\rangle \mathbin{\Vert} \langle |u| - |s|\rangle))))))\).
Chaining the equalities, we obtain \(|s| - (|u| + 1) = \operatorname{eval}(\operatorname{Lean\mathord{:}Omega\mathord{:}Linear\,Combo\mathord{:}mk}(0,\allowbreak{} \langle 0\rangle \mathbin{\Vert} \langle 1\rangle) - \operatorname{Lean\mathord{:}Omega\mathord{:}Linear\,Combo\mathord{:}mk}(1,\allowbreak{} \langle 0\rangle \mathbin{\Vert} (\langle 0\rangle \mathbin{\Vert} (\langle 0\rangle \mathbin{\Vert} (\langle 0\rangle \mathbin{\Vert} \langle 1\rangle)))),\allowbreak{} \operatorname{Lean\mathord{:}Omega\mathord{:}Coeffs\mathord{:}of\,List}(\langle n\rangle \mathbin{\Vert} (\langle |s|\rangle \mathbin{\Vert} (\langle |t|\rangle \mathbin{\Vert} (\langle k\rangle \mathbin{\Vert} (\langle |u|\rangle \mathbin{\Vert} \langle |u| - |s|\rangle))))))\).
By \paperlink{premise:6c655f6f665f6c655f6f665f6571}{Prerequisite~P\ref{premise:6c655f6f665f6c655f6f665f6571}} with \(\mathbb{Z}\) as its type parameter \(\alpha\), we obtain \(0 \le \operatorname{eval}(\operatorname{Lean\mathord{:}Omega\mathord{:}Linear\,Combo\mathord{:}mk}(0,\allowbreak{} \langle 0\rangle \mathbin{\Vert} \langle 1\rangle) - \operatorname{Lean\mathord{:}Omega\mathord{:}Linear\,Combo\mathord{:}mk}(1,\allowbreak{} \langle 0\rangle \mathbin{\Vert} (\langle 0\rangle \mathbin{\Vert} (\langle 0\rangle \mathbin{\Vert} (\langle 0\rangle \mathbin{\Vert} \langle 1\rangle)))),\allowbreak{} \operatorname{Lean\mathord{:}Omega\mathord{:}Coeffs\mathord{:}of\,List}(\langle n\rangle \mathbin{\Vert} (\langle |s|\rangle \mathbin{\Vert} (\langle |t|\rangle \mathbin{\Vert} (\langle k\rangle \mathbin{\Vert} (\langle |u|\rangle \mathbin{\Vert} \langle |u| - |s|\rangle))))))\).
By \paperlink{premise:4c65616e2e4f6d6567612e436f6e73747261696e742e616464496e657175616c6974795f736174}{Prerequisite~P\ref{premise:4c65616e2e4f6d6567612e436f6e73747261696e742e616464496e657175616c6974795f736174}}, we obtain \(\operatorname{sat^{\prime}}(\operatorname{Lean\mathord{:}Omega\mathord{:}Constraint\mathord{:}mk}(\operatorname{some}(-(-1)),\allowbreak{} \operatorname{none}),\allowbreak{} \langle 0\rangle \mathbin{\Vert} (\langle 1\rangle \mathbin{\Vert} (\langle 0\rangle \mathbin{\Vert} (\langle 0\rangle \mathbin{\Vert} \langle -1\rangle))),\allowbreak{} \operatorname{Lean\mathord{:}Omega\mathord{:}Coeffs\mathord{:}of\,List}(\langle n\rangle \mathbin{\Vert} (\langle |s|\rangle \mathbin{\Vert} (\langle |t|\rangle \mathbin{\Vert} (\langle k\rangle \mathbin{\Vert} (\langle |u|\rangle \mathbin{\Vert} \langle |u| - |s|\rangle)))))) = \mathrm{true}\).
By \paperlink{premise:4c65616e2e4f6d6567612e436f6e73747261696e742e636f6d62696e655f73617427}{Prerequisite~P\ref{premise:4c65616e2e4f6d6567612e436f6e73747261696e742e636f6d62696e655f73617427}}, we obtain \(\operatorname{sat^{\prime}}(\operatorname{combine}(\operatorname{Lean\mathord{:}Omega\mathord{:}Constraint\mathord{:}mk}(\operatorname{none},\allowbreak{} \operatorname{some}(-1)),\allowbreak{} \operatorname{Lean\mathord{:}Omega\mathord{:}Constraint\mathord{:}mk}(\operatorname{some}(1),\allowbreak{} \operatorname{none})),\allowbreak{} \langle 0\rangle \mathbin{\Vert} (\langle 1\rangle \mathbin{\Vert} (\langle 0\rangle \mathbin{\Vert} (\langle 0\rangle \mathbin{\Vert} \langle -1\rangle))),\allowbreak{} \operatorname{Lean\mathord{:}Omega\mathord{:}Coeffs\mathord{:}of\,List}(\langle n\rangle \mathbin{\Vert} (\langle |s|\rangle \mathbin{\Vert} (\langle |t|\rangle \mathbin{\Vert} (\langle k\rangle \mathbin{\Vert} (\langle |u|\rangle \mathbin{\Vert} \langle |u| - |s|\rangle)))))) = \mathrm{true}\).
By \paperlink{premise:4c65616e2e4f6d6567612e436f6e73747261696e742e6e6f745f736174275f6f665f6973496d706f737369626c65}{Prerequisite~P\ref{premise:4c65616e2e4f6d6567612e436f6e73747261696e742e6e6f745f736174275f6f665f6973496d706f737369626c65}}, we obtain \(\mathrm{False}\).
\end{proofstep}
The two alternatives exhaust the cases. Hence \(\mathrm{False}\).
\end{proofstep}
The contrary assumption leads to a contradiction. Hence \(0 < \left|u\right| - \left|s\right|\).
Rewriting with \paperlink{theorem:90}{Lemma~\ref{theorem:90}} (using \(s \subseteq u\)) in the preceding result gives \(0 < \left|u \setminus s\right|\).
Rewriting with \paperlink{theorem:9}{Lemma~\ref{theorem:9}} in the preceding result gives \(u \setminus s\ne\varnothing\).
Consider the following case.
\begin{proofstep}
Fix \(a_{1}\) in \(\alpha\).
Assume \(a_{1} \in u \setminus s\).
Choose the witness \(u \setminus \{a_{1}\}\).
Use the hypothesis \(a_{1} \in u \setminus s\).
Rewriting with \paperlink{premise:46696e7365742e6d656d5f7364696666}{Prerequisite~P\ref{premise:46696e7365742e6d656d5f7364696666}} in the preceding result gives \(a_{1} \in u \land a_{1} \notin s\).
Taking the second conjunct gives \(a_{1} \notin s\).
Use the hypothesis \(a_{1} \in u \setminus s\).
Rewriting with \paperlink{premise:46696e7365742e6d656d5f7364696666}{Prerequisite~P\ref{premise:46696e7365742e6d656d5f7364696666}} in the preceding result gives \(a_{1} \in u \land a_{1} \notin s\).
Taking the first conjunct gives \(a_{1} \in u\).
Use the hypothesis \(a_{1} \in u \setminus s\).
Rewriting with \paperlink{premise:46696e7365742e6d656d5f7364696666}{Prerequisite~P\ref{premise:46696e7365742e6d656d5f7364696666}} in the preceding result gives \(a_{1} \in u \land a_{1} \notin s\).
Taking the first conjunct gives \(a_{1} \in u\).
Rewriting with \paperlink{premise:46696e7365742e7375627365745f6572617365}{Prerequisite~P\ref{premise:46696e7365742e7375627365745f6572617365}}, \(s \subseteq u\), \paperlink{premise:65715f66616c7365}{Prerequisite~P\ref{premise:65715f66616c7365}}, \(a_{1} \notin s\), \paperlink{premise:6e6f745f66616c73655f65715f74727565}{Prerequisite~P\ref{premise:6e6f745f66616c73655f65715f74727565}}, \paperlink{premise:616e645f73656c66}{Prerequisite~P\ref{premise:616e645f73656c66}}, \paperlink{premise:46696e7365742e65726173655f7375627365745f6966665f6f665f6d656d}{Prerequisite~P\ref{premise:46696e7365742e65726173655f7375627365745f6966665f6f665f6d656d}}, \(u \subseteq t\), \(a_{1} \in u\), \paperlink{theorem:28}{Lemma~\ref{theorem:28}} and \(\left|u\right| = k + 1\) gives \(s \subseteq u \setminus \{a_{1}\} \land (u \setminus \{a_{1}\} \subseteq t \land |u \setminus \{a_{1}\}| = k)\).
This witness proves existence. Hence \(\exists u_{1}\, :\, \mathcal{P}_{\mathrm{fin}}(\alpha),\; s \subseteq u_{1} \land (u_{1} \subseteq t \land |u_{1}| = k)\).
\end{proofstep}
Using the existential hypothesis, we obtain \(\exists u_{1}\, :\, \mathcal{P}_{\mathrm{fin}}(\alpha),\; s \subseteq u_{1} \land (u_{1} \subseteq t \land |u_{1}| = k)\).
\end{proofstep}
Using the witnesses and their properties, we obtain \(\exists u\, :\, \mathcal{P}_{\mathrm{fin}}(\alpha),\; s \subseteq u \land (u \subseteq t \land |u| = k)\).
\end{proofstep}
By downward induction through this interval, we obtain \(\exists u\, :\, \mathcal{P}_{\mathrm{fin}}(\alpha),\; s \subseteq u \land (u \subseteq t \land |u| = n)\).
\end{proof}

\begin{lemma}\label{theorem:103}
Let \(\alpha\) be a type. Let \(s\) be a finite subset of \(\alpha\). Let \(n\) be a natural number. Assume \(n \le \left|s\right|\). Then
\begin{center}
\(\displaystyle \exists t\, :\, \mathcal{P}_{\mathrm{fin}}(\alpha),\; t \subseteq s \land \left|t\right| = n.\)
\end{center}
\end{lemma}

\begin{proof}
The substitution uses the following argument.
\begin{proofstep}
Fix \(u\) in \(\mathcal{P}_{\mathrm{fin}}(\alpha)\).
By \paperlink{premise:46696e7365742e656d7074795f737562736574}{Prerequisite~P\ref{premise:46696e7365742e656d7074795f737562736574}} with \(u\) in place of its parameter \(s\), we obtain \(\varnothing \subseteq u\).
By \paperlink{premise:65715f74727565}{Prerequisite~P\ref{premise:65715f74727565}}, we obtain \(\left(\varnothing \subseteq u\right) = \mathrm{True}\).
Applying the same function to equal arguments, we obtain \((\varnothing \subseteq u \land (u \subseteq s \land |u| = n)) = (\mathrm{True} \land (u \subseteq s \land |u| = n))\).
By \paperlink{premise:747275655f616e64}{Prerequisite~P\ref{premise:747275655f616e64}} with \(p = u \subseteq s \land \left|u\right| = n\), we obtain \((\mathrm{True} \land (u \subseteq s \land |u| = n)) = (u \subseteq s \land |u| = n)\).
Chaining the equalities, we obtain \((\varnothing \subseteq u \land (u \subseteq s \land |u| = n)) = (u \subseteq s \land |u| = n)\).
\end{proofstep}
By \paperlink{premise:46696e7365742e656d7074795f737562736574}{Prerequisite~P\ref{premise:46696e7365742e656d7074795f737562736574}}, we obtain \(\varnothing \subseteq s\).
Rewriting with \paperlink{premise:4e61742e7a65726f5f6c65}{Prerequisite~P\ref{premise:4e61742e7a65726f5f6c65}} gives \(0 \le n\).
By \paperlink{theorem:102}{Lemma~\ref{theorem:102}}, we obtain \(\exists u\, :\, \mathcal{P}_{\mathrm{fin}}(\alpha),\; \varnothing \subseteq u \land (u \subseteq s \land |u| = n)\).
Rewriting with \paperlink{premise:66756e657874}{\emph{Function extensionality}} (\paperlink{premise:66756e657874}{Prerequisite~P\ref{premise:66756e657874}}) and the identity established above in the preceding result gives \(\exists t\, :\, \mathcal{P}_{\mathrm{fin}}(\alpha),\; t \subseteq s \land \left|t\right| = n\).
\end{proof}

\begin{lemma}\label{theorem:104}
Let \(\alpha\) be a type. Let \(s\) be a finite subset of \(\alpha\). Let \(n\) be a natural number. Then
\begin{center}
\(\displaystyle n \le |s| \iff (\exists t\, :\, \mathcal{P}_{\mathrm{fin}}(\alpha),\; t \subseteq s \land |t| = n).\)
\end{center}
\end{lemma}

\begin{proof}
Consider the following case.
\begin{proofstep}
Assume \(n \le \left|s\right|\).
By \paperlink{theorem:103}{Lemma~\ref{theorem:103}}, we obtain \(\exists t\, :\, \mathcal{P}_{\mathrm{fin}}(\alpha),\; t \subseteq s \land \left|t\right| = n\).
\end{proofstep}
Consider the following case.
\begin{proofstep}
Assume \(\exists t\, :\, \mathcal{P}_{\mathrm{fin}}(\alpha),\; t \subseteq s \land \left|t\right| = n\).
Use the hypothesis \(\exists t\, :\, \mathcal{P}_{\mathrm{fin}}(\alpha),\; t \subseteq s \land \left|t\right| = n\).
Take the witnesses from the preceding result.
\begin{proofstep}
Fix \(t\) in \(\mathcal{P}_{\mathrm{fin}}(\alpha)\).
Assume \(t \subseteq s\).
Assume \(\left|t\right| = n\).
By \paperlink{theorem:5}{Lemma~\ref{theorem:5}}, we obtain \(\left|t\right| \le \left|s\right|\).
Use the hypothesis \(\left|t\right| = n\).
Substituting equals, we obtain \(n \le \left|s\right|\).
\end{proofstep}
Using the witnesses and their properties, we obtain \(n \le \left|s\right|\).
\end{proofstep}
Together the two implications give the equivalence. Hence \(n \le |s| \iff (\exists t\, :\, \mathcal{P}_{\mathrm{fin}}(\alpha),\; t \subseteq s \land |t| = n)\).
\end{proof}

\begin{theorem}\label{theorem:105}
Let \(\alpha\) be a type. Assume equality on \(\alpha\) is decidable. Let \(X\) and \(Y\) be finite subsets of \(\alpha\). Let \(n\) be a natural number. Assume \(2 \cdot n < \left|X \cup Y\right|\). Then
\begin{center}
\(\displaystyle \exists C\, :\, \mathcal{P}_{\mathrm{fin}}(\alpha),\; n < |C| \land (C \subseteq X \lor C \subseteq Y).\)
\end{center}
\end{theorem}

\begin{proof}
By \paperlink{theorem:7}{Lemma~\ref{theorem:7}}, we obtain \(\left|X \cap \left(Y \setminus X\right)\right| = 0 \iff X \cap \left(Y \setminus X\right) = \varnothing\).
By \paperlink{premise:46696e7365742e696e7465725f73646966665f73656c66}{Prerequisite~P\ref{premise:46696e7365742e696e7465725f73646966665f73656c66}} with \(s_{1} = X\) and \(s_{2} = Y\), we obtain \(X \cap \left(Y \setminus X\right) = \varnothing\).
By the reverse implication, we obtain \(\left|X \cap \left(Y \setminus X\right)\right| = 0\).
Rewriting with \paperlink{premise:4e61742e6164645f7a65726f}{Prerequisite~P\ref{premise:4e61742e6164645f7a65726f}}, \(\left|X \cap \left(Y \setminus X\right)\right| = 0\), \paperlink{premise:46696e7365742e756e696f6e5f73646966665f73656c665f65715f756e696f6e}{Prerequisite~P\ref{premise:46696e7365742e756e696f6e5f73646966665f73656c665f65715f756e696f6e}} and \paperlink{theorem:83}{Lemma~\ref{theorem:83}} gives \(\left|X \cup Y\right| = \left|X\right| + \left|Y \setminus X\right|\).
Assume \(\left|X\right| \le n \land \left|Y \setminus X\right| \le n\).
Consider the following case.
\begin{proofstep}
Assume \(\neg \left(\left|X\right| + \left|Y \setminus X\right| \le n + n\right)\).
By definition, we obtain \(\operatorname{decide}(\operatorname{is\,Impossible}(\operatorname{Lean\mathord{:}Omega\mathord{:}Constraint\mathord{:}mk}(\operatorname{some}(1),\allowbreak{} \operatorname{some}(0))) = \mathrm{true}) = \mathrm{true}\).
By \paperlink{premise:6f665f6465636964655f65715f74727565}{Prerequisite~P\ref{premise:6f665f6465636964655f65715f74727565}}, we obtain \(\operatorname{is\,Impossible}(\operatorname{Lean\mathord{:}Omega\mathord{:}Constraint\mathord{:}mk}(\operatorname{some}(1),\allowbreak{} \operatorname{some}(0))) = \mathrm{true}\).
Use the hypothesis \(\neg \left(\left|X\right| + \left|Y \setminus X\right| \le n + n\right)\).
By \paperlink{premise:4e61742e6c745f6f665f6e6f745f6c65}{Prerequisite~P\ref{premise:4e61742e6c745f6f665f6e6f745f6c65}}, we obtain \(n + n < \left|X\right| + \left|Y \setminus X\right|\).
By \paperlink{premise:4c65616e2e4f6d6567612e496e742e6f664e61745f6c745f6f665f6c74}{Prerequisite~P\ref{premise:4c65616e2e4f6d6567612e496e742e6f664e61745f6c745f6f665f6c74}}, we obtain \(n + n < \left|X\right| + \left|Y \setminus X\right|\).
By \paperlink{premise:496e742e6164645f6f6e655f6c655f6f665f6c74}{Prerequisite~P\ref{premise:496e742e6164645f6f6e655f6c655f6f665f6c74}}, we obtain \(n + n + 1 \le \left|X\right| + \left|Y \setminus X\right|\).
By \paperlink{premise:496e742e7375625f6e6f6e6e65675f6f665f6c65}{Prerequisite~P\ref{premise:496e742e7375625f6e6f6e6e65675f6f665f6c65}}, we obtain \(0 \le \left|X\right| + \left|Y \setminus X\right| - \left(n + n + 1\right)\).
By \paperlink{premise:496e742e6f664e61745f616464}{Prerequisite~P\ref{premise:496e742e6f664e61745f616464}} with \(\left|X\right|\) in place of its parameter \(n\) and \(m = \left|Y \setminus X\right|\), we obtain \(\left|X\right| + \left|Y \setminus X\right| = \left|X\right| + \left|Y \setminus X\right|\).
Using \paperlink{premise:4c65616e2e4f6d6567612e4c696e656172436f6d626f2e636f6f7264696e6174655f6576616c5f32}{Prerequisite~P\ref{premise:4c65616e2e4f6d6567612e4c696e656172436f6d626f2e636f6f7264696e6174655f6576616c5f32}} in reverse gives \(|X| = \operatorname{eval}(\operatorname{coordinate}(2),\allowbreak{} \operatorname{Lean\mathord{:}Omega\mathord{:}Coeffs\mathord{:}of\,List}(\langle |X \cap (Y \setminus X)|\rangle \mathbin{\Vert} (\langle |X \cup Y|\rangle \mathbin{\Vert} (\langle |X|\rangle \mathbin{\Vert} (\langle |Y \setminus X|\rangle \mathbin{\Vert} \langle n\rangle)))))\).
Using \paperlink{premise:4c65616e2e4f6d6567612e4c696e656172436f6d626f2e636f6f7264696e6174655f6576616c5f33}{Prerequisite~P\ref{premise:4c65616e2e4f6d6567612e4c696e656172436f6d626f2e636f6f7264696e6174655f6576616c5f33}} in reverse gives \(|Y \setminus X| = \operatorname{eval}(\operatorname{coordinate}(3),\allowbreak{} \operatorname{Lean\mathord{:}Omega\mathord{:}Coeffs\mathord{:}of\,List}(\langle |X \cap (Y \setminus X)|\rangle \mathbin{\Vert} (\langle |X \cup Y|\rangle \mathbin{\Vert} (\langle |X|\rangle \mathbin{\Vert} (\langle |Y \setminus X|\rangle \mathbin{\Vert} \langle n\rangle)))))\).
By \paperlink{premise:4c65616e2e4f6d6567612e496e742e6164645f636f6e6772}{Prerequisite~P\ref{premise:4c65616e2e4f6d6567612e496e742e6164645f636f6e6772}}, we obtain \(|X| + |Y \setminus X| = \operatorname{eval}(\operatorname{coordinate}(2),\allowbreak{} \operatorname{Lean\mathord{:}Omega\mathord{:}Coeffs\mathord{:}of\,List}(\langle |X \cap (Y \setminus X)|\rangle \mathbin{\Vert} (\langle |X \cup Y|\rangle \mathbin{\Vert} (\langle |X|\rangle \mathbin{\Vert} (\langle |Y \setminus X|\rangle \mathbin{\Vert} \langle n\rangle))))) + \operatorname{eval}(\operatorname{coordinate}(3),\allowbreak{} \operatorname{Lean\mathord{:}Omega\mathord{:}Coeffs\mathord{:}of\,List}(\langle |X \cap (Y \setminus X)|\rangle \mathbin{\Vert} (\langle |X \cup Y|\rangle \mathbin{\Vert} (\langle |X|\rangle \mathbin{\Vert} (\langle |Y \setminus X|\rangle \mathbin{\Vert} \langle n\rangle)))))\).
Using \paperlink{premise:4c65616e2e4f6d6567612e4c696e656172436f6d626f2e6164645f6576616c}{Prerequisite~P\ref{premise:4c65616e2e4f6d6567612e4c696e656172436f6d626f2e6164645f6576616c}} with \(l_{1} = \operatorname{Lean\mathord{:}Omega\mathord{:}Linear\,Combo\mathord{:}mk}(0,\allowbreak{} \langle 0\rangle \mathbin{\Vert} (\langle 0\rangle \mathbin{\Vert} \langle 1\rangle))\) and \(l_{2} = \operatorname{Lean\mathord{:}Omega\mathord{:}Linear\,Combo\mathord{:}mk}(0,\allowbreak{} \langle 0\rangle \mathbin{\Vert} (\langle 0\rangle \mathbin{\Vert} (\langle 0\rangle \mathbin{\Vert} \langle 1\rangle)))\) and \(v = \operatorname{Lean\mathord{:}Omega\mathord{:}Coeffs\mathord{:}of\,List}(\langle |X \cap (Y \setminus X)|\rangle \mathbin{\Vert} (\langle |X \cup Y|\rangle \mathbin{\Vert} (\langle |X|\rangle \mathbin{\Vert} (\langle |Y \setminus X|\rangle \mathbin{\Vert} \langle n\rangle))))\) in reverse gives \(\operatorname{eval}(\operatorname{Lean\mathord{:}Omega\mathord{:}Linear\,Combo\mathord{:}mk}(0,\allowbreak{} \langle 0\rangle \mathbin{\Vert} (\langle 0\rangle \mathbin{\Vert} \langle 1\rangle)),\allowbreak{} \operatorname{Lean\mathord{:}Omega\mathord{:}Coeffs\mathord{:}of\,List}(\langle |X \cap (Y \setminus X)|\rangle \mathbin{\Vert} (\langle |X \cup Y|\rangle \mathbin{\Vert} (\langle |X|\rangle \mathbin{\Vert} (\langle |Y \setminus X|\rangle \mathbin{\Vert} \langle n\rangle))))) + \operatorname{eval}(\operatorname{Lean\mathord{:}Omega\mathord{:}Linear\,Combo\mathord{:}mk}(0,\allowbreak{} \langle 0\rangle \mathbin{\Vert} (\langle 0\rangle \mathbin{\Vert} (\langle 0\rangle \mathbin{\Vert} \langle 1\rangle))),\allowbreak{} \operatorname{Lean\mathord{:}Omega\mathord{:}Coeffs\mathord{:}of\,List}(\langle |X \cap (Y \setminus X)|\rangle \mathbin{\Vert} (\langle |X \cup Y|\rangle \mathbin{\Vert} (\langle |X|\rangle \mathbin{\Vert} (\langle |Y \setminus X|\rangle \mathbin{\Vert} \langle n\rangle))))) = \operatorname{eval}(\operatorname{Lean\mathord{:}Omega\mathord{:}Linear\,Combo\mathord{:}mk}(0,\allowbreak{} \langle 0\rangle \mathbin{\Vert} (\langle 0\rangle \mathbin{\Vert} \langle 1\rangle)) + \operatorname{Lean\mathord{:}Omega\mathord{:}Linear\,Combo\mathord{:}mk}(0,\allowbreak{} \langle 0\rangle \mathbin{\Vert} (\langle 0\rangle \mathbin{\Vert} (\langle 0\rangle \mathbin{\Vert} \langle 1\rangle))),\allowbreak{} \operatorname{Lean\mathord{:}Omega\mathord{:}Coeffs\mathord{:}of\,List}(\langle |X \cap (Y \setminus X)|\rangle \mathbin{\Vert} (\langle |X \cup Y|\rangle \mathbin{\Vert} (\langle |X|\rangle \mathbin{\Vert} (\langle |Y \setminus X|\rangle \mathbin{\Vert} \langle n\rangle)))))\).
Chaining the equalities, we obtain \(|X| + |Y \setminus X| = \operatorname{eval}(\operatorname{Lean\mathord{:}Omega\mathord{:}Linear\,Combo\mathord{:}mk}(0,\allowbreak{} \langle 0\rangle \mathbin{\Vert} (\langle 0\rangle \mathbin{\Vert} \langle 1\rangle)) + \operatorname{Lean\mathord{:}Omega\mathord{:}Linear\,Combo\mathord{:}mk}(0,\allowbreak{} \langle 0\rangle \mathbin{\Vert} (\langle 0\rangle \mathbin{\Vert} (\langle 0\rangle \mathbin{\Vert} \langle 1\rangle))),\allowbreak{} \operatorname{Lean\mathord{:}Omega\mathord{:}Coeffs\mathord{:}of\,List}(\langle |X \cap (Y \setminus X)|\rangle \mathbin{\Vert} (\langle |X \cup Y|\rangle \mathbin{\Vert} (\langle |X|\rangle \mathbin{\Vert} (\langle |Y \setminus X|\rangle \mathbin{\Vert} \langle n\rangle)))))\).
Chaining the equalities, we obtain \(|X| + |Y \setminus X| = \operatorname{eval}(\operatorname{Lean\mathord{:}Omega\mathord{:}Linear\,Combo\mathord{:}mk}(0,\allowbreak{} \langle 0\rangle \mathbin{\Vert} (\langle 0\rangle \mathbin{\Vert} \langle 1\rangle)) + \operatorname{Lean\mathord{:}Omega\mathord{:}Linear\,Combo\mathord{:}mk}(0,\allowbreak{} \langle 0\rangle \mathbin{\Vert} (\langle 0\rangle \mathbin{\Vert} (\langle 0\rangle \mathbin{\Vert} \langle 1\rangle))),\allowbreak{} \operatorname{Lean\mathord{:}Omega\mathord{:}Coeffs\mathord{:}of\,List}(\langle |X \cap (Y \setminus X)|\rangle \mathbin{\Vert} (\langle |X \cup Y|\rangle \mathbin{\Vert} (\langle |X|\rangle \mathbin{\Vert} (\langle |Y \setminus X|\rangle \mathbin{\Vert} \langle n\rangle)))))\).
By \paperlink{premise:496e742e6f664e61745f616464}{Prerequisite~P\ref{premise:496e742e6f664e61745f616464}} with \(m = n\), we obtain \(n + n = n + n\).
Using \paperlink{premise:4c65616e2e4f6d6567612e4c696e656172436f6d626f2e636f6f7264696e6174655f6576616c5f34}{Prerequisite~P\ref{premise:4c65616e2e4f6d6567612e4c696e656172436f6d626f2e636f6f7264696e6174655f6576616c5f34}} in reverse gives \(n = \operatorname{eval}(\operatorname{coordinate}(4),\allowbreak{} \operatorname{Lean\mathord{:}Omega\mathord{:}Coeffs\mathord{:}of\,List}(\langle |X \cap (Y \setminus X)|\rangle \mathbin{\Vert} (\langle |X \cup Y|\rangle \mathbin{\Vert} (\langle |X|\rangle \mathbin{\Vert} (\langle |Y \setminus X|\rangle \mathbin{\Vert} \langle n\rangle)))))\).
Using \paperlink{premise:4c65616e2e4f6d6567612e4c696e656172436f6d626f2e636f6f7264696e6174655f6576616c5f34}{Prerequisite~P\ref{premise:4c65616e2e4f6d6567612e4c696e656172436f6d626f2e636f6f7264696e6174655f6576616c5f34}} in reverse gives \(n = \operatorname{eval}(\operatorname{coordinate}(4),\allowbreak{} \operatorname{Lean\mathord{:}Omega\mathord{:}Coeffs\mathord{:}of\,List}(\langle |X \cap (Y \setminus X)|\rangle \mathbin{\Vert} (\langle |X \cup Y|\rangle \mathbin{\Vert} (\langle |X|\rangle \mathbin{\Vert} (\langle |Y \setminus X|\rangle \mathbin{\Vert} \langle n\rangle)))))\).
By \paperlink{premise:4c65616e2e4f6d6567612e496e742e6164645f636f6e6772}{Prerequisite~P\ref{premise:4c65616e2e4f6d6567612e496e742e6164645f636f6e6772}}, we obtain \(n + n = \operatorname{eval}(\operatorname{coordinate}(4),\allowbreak{} \operatorname{Lean\mathord{:}Omega\mathord{:}Coeffs\mathord{:}of\,List}(\langle |X \cap (Y \setminus X)|\rangle \mathbin{\Vert} (\langle |X \cup Y|\rangle \mathbin{\Vert} (\langle |X|\rangle \mathbin{\Vert} (\langle |Y \setminus X|\rangle \mathbin{\Vert} \langle n\rangle))))) + \operatorname{eval}(\operatorname{coordinate}(4),\allowbreak{} \operatorname{Lean\mathord{:}Omega\mathord{:}Coeffs\mathord{:}of\,List}(\langle |X \cap (Y \setminus X)|\rangle \mathbin{\Vert} (\langle |X \cup Y|\rangle \mathbin{\Vert} (\langle |X|\rangle \mathbin{\Vert} (\langle |Y \setminus X|\rangle \mathbin{\Vert} \langle n\rangle)))))\).
Using \paperlink{premise:4c65616e2e4f6d6567612e4c696e656172436f6d626f2e6164645f6576616c}{Prerequisite~P\ref{premise:4c65616e2e4f6d6567612e4c696e656172436f6d626f2e6164645f6576616c}} with \(l_{1} = \operatorname{Lean\mathord{:}Omega\mathord{:}Linear\,Combo\mathord{:}mk}(0,\allowbreak{} \langle 0\rangle \mathbin{\Vert} (\langle 0\rangle \mathbin{\Vert} (\langle 0\rangle \mathbin{\Vert} (\langle 0\rangle \mathbin{\Vert} \langle 1\rangle))))\) and \(l_{2} = \operatorname{Lean\mathord{:}Omega\mathord{:}Linear\,Combo\mathord{:}mk}(0,\allowbreak{} \langle 0\rangle \mathbin{\Vert} (\langle 0\rangle \mathbin{\Vert} (\langle 0\rangle \mathbin{\Vert} (\langle 0\rangle \mathbin{\Vert} \langle 1\rangle))))\) and \(v = \operatorname{Lean\mathord{:}Omega\mathord{:}Coeffs\mathord{:}of\,List}(\langle |X \cap (Y \setminus X)|\rangle \mathbin{\Vert} (\langle |X \cup Y|\rangle \mathbin{\Vert} (\langle |X|\rangle \mathbin{\Vert} (\langle |Y \setminus X|\rangle \mathbin{\Vert} \langle n\rangle))))\) in reverse gives \(\operatorname{eval}(\operatorname{Lean\mathord{:}Omega\mathord{:}Linear\,Combo\mathord{:}mk}(0,\allowbreak{} \langle 0\rangle \mathbin{\Vert} (\langle 0\rangle \mathbin{\Vert} (\langle 0\rangle \mathbin{\Vert} (\langle 0\rangle \mathbin{\Vert} \langle 1\rangle)))),\allowbreak{} \operatorname{Lean\mathord{:}Omega\mathord{:}Coeffs\mathord{:}of\,List}(\langle |X \cap (Y \setminus X)|\rangle \mathbin{\Vert} (\langle |X \cup Y|\rangle \mathbin{\Vert} (\langle |X|\rangle \mathbin{\Vert} (\langle |Y \setminus X|\rangle \mathbin{\Vert} \langle n\rangle))))) + \operatorname{eval}(\operatorname{Lean\mathord{:}Omega\mathord{:}Linear\,Combo\mathord{:}mk}(0,\allowbreak{} \langle 0\rangle \mathbin{\Vert} (\langle 0\rangle \mathbin{\Vert} (\langle 0\rangle \mathbin{\Vert} (\langle 0\rangle \mathbin{\Vert} \langle 1\rangle)))),\allowbreak{} \operatorname{Lean\mathord{:}Omega\mathord{:}Coeffs\mathord{:}of\,List}(\langle |X \cap (Y \setminus X)|\rangle \mathbin{\Vert} (\langle |X \cup Y|\rangle \mathbin{\Vert} (\langle |X|\rangle \mathbin{\Vert} (\langle |Y \setminus X|\rangle \mathbin{\Vert} \langle n\rangle))))) = \operatorname{eval}(\operatorname{Lean\mathord{:}Omega\mathord{:}Linear\,Combo\mathord{:}mk}(0,\allowbreak{} \langle 0\rangle \mathbin{\Vert} (\langle 0\rangle \mathbin{\Vert} (\langle 0\rangle \mathbin{\Vert} (\langle 0\rangle \mathbin{\Vert} \langle 1\rangle)))) + \operatorname{Lean\mathord{:}Omega\mathord{:}Linear\,Combo\mathord{:}mk}(0,\allowbreak{} \langle 0\rangle \mathbin{\Vert} (\langle 0\rangle \mathbin{\Vert} (\langle 0\rangle \mathbin{\Vert} (\langle 0\rangle \mathbin{\Vert} \langle 1\rangle)))),\allowbreak{} \operatorname{Lean\mathord{:}Omega\mathord{:}Coeffs\mathord{:}of\,List}(\langle |X \cap (Y \setminus X)|\rangle \mathbin{\Vert} (\langle |X \cup Y|\rangle \mathbin{\Vert} (\langle |X|\rangle \mathbin{\Vert} (\langle |Y \setminus X|\rangle \mathbin{\Vert} \langle n\rangle)))))\).
Chaining the equalities, we obtain \(n + n = \operatorname{eval}(\operatorname{Lean\mathord{:}Omega\mathord{:}Linear\,Combo\mathord{:}mk}(0,\allowbreak{} \langle 0\rangle \mathbin{\Vert} (\langle 0\rangle \mathbin{\Vert} (\langle 0\rangle \mathbin{\Vert} (\langle 0\rangle \mathbin{\Vert} \langle 1\rangle)))) + \operatorname{Lean\mathord{:}Omega\mathord{:}Linear\,Combo\mathord{:}mk}(0,\allowbreak{} \langle 0\rangle \mathbin{\Vert} (\langle 0\rangle \mathbin{\Vert} (\langle 0\rangle \mathbin{\Vert} (\langle 0\rangle \mathbin{\Vert} \langle 1\rangle)))),\allowbreak{} \operatorname{Lean\mathord{:}Omega\mathord{:}Coeffs\mathord{:}of\,List}(\langle |X \cap (Y \setminus X)|\rangle \mathbin{\Vert} (\langle |X \cup Y|\rangle \mathbin{\Vert} (\langle |X|\rangle \mathbin{\Vert} (\langle |Y \setminus X|\rangle \mathbin{\Vert} \langle n\rangle)))))\).
Chaining the equalities, we obtain \(n + n = \operatorname{eval}(\operatorname{Lean\mathord{:}Omega\mathord{:}Linear\,Combo\mathord{:}mk}(0,\allowbreak{} \langle 0\rangle \mathbin{\Vert} (\langle 0\rangle \mathbin{\Vert} (\langle 0\rangle \mathbin{\Vert} (\langle 0\rangle \mathbin{\Vert} \langle 1\rangle)))) + \operatorname{Lean\mathord{:}Omega\mathord{:}Linear\,Combo\mathord{:}mk}(0,\allowbreak{} \langle 0\rangle \mathbin{\Vert} (\langle 0\rangle \mathbin{\Vert} (\langle 0\rangle \mathbin{\Vert} (\langle 0\rangle \mathbin{\Vert} \langle 1\rangle)))),\allowbreak{} \operatorname{Lean\mathord{:}Omega\mathord{:}Coeffs\mathord{:}of\,List}(\langle |X \cap (Y \setminus X)|\rangle \mathbin{\Vert} (\langle |X \cup Y|\rangle \mathbin{\Vert} (\langle |X|\rangle \mathbin{\Vert} (\langle |Y \setminus X|\rangle \mathbin{\Vert} \langle n\rangle)))))\).
By definition, we obtain \(1 = \operatorname{eval}(\operatorname{Lean\mathord{:}Omega\mathord{:}Linear\,Combo\mathord{:}mk}(1,\allowbreak{} \langle\,\rangle),\allowbreak{} \operatorname{Lean\mathord{:}Omega\mathord{:}Coeffs\mathord{:}of\,List}(\langle |X \cap (Y \setminus X)|\rangle \mathbin{\Vert} (\langle |X \cup Y|\rangle \mathbin{\Vert} (\langle |X|\rangle \mathbin{\Vert} (\langle |Y \setminus X|\rangle \mathbin{\Vert} \langle n\rangle)))))\).
By \paperlink{premise:4c65616e2e4f6d6567612e496e742e6164645f636f6e6772}{Prerequisite~P\ref{premise:4c65616e2e4f6d6567612e496e742e6164645f636f6e6772}}, we obtain \(n + n + 1 = \operatorname{eval}(\operatorname{Lean\mathord{:}Omega\mathord{:}Linear\,Combo\mathord{:}mk}(0,\allowbreak{} \langle 0\rangle \mathbin{\Vert} (\langle 0\rangle \mathbin{\Vert} (\langle 0\rangle \mathbin{\Vert} (\langle 0\rangle \mathbin{\Vert} \langle 1\rangle)))) + \operatorname{Lean\mathord{:}Omega\mathord{:}Linear\,Combo\mathord{:}mk}(0,\allowbreak{} \langle 0\rangle \mathbin{\Vert} (\langle 0\rangle \mathbin{\Vert} (\langle 0\rangle \mathbin{\Vert} (\langle 0\rangle \mathbin{\Vert} \langle 1\rangle)))),\allowbreak{} \operatorname{Lean\mathord{:}Omega\mathord{:}Coeffs\mathord{:}of\,List}(\langle |X \cap (Y \setminus X)|\rangle \mathbin{\Vert} (\langle |X \cup Y|\rangle \mathbin{\Vert} (\langle |X|\rangle \mathbin{\Vert} (\langle |Y \setminus X|\rangle \mathbin{\Vert} \langle n\rangle))))) + \operatorname{eval}(\operatorname{Lean\mathord{:}Omega\mathord{:}Linear\,Combo\mathord{:}mk}(1,\allowbreak{} \langle\,\rangle),\allowbreak{} \operatorname{Lean\mathord{:}Omega\mathord{:}Coeffs\mathord{:}of\,List}(\langle |X \cap (Y \setminus X)|\rangle \mathbin{\Vert} (\langle |X \cup Y|\rangle \mathbin{\Vert} (\langle |X|\rangle \mathbin{\Vert} (\langle |Y \setminus X|\rangle \mathbin{\Vert} \langle n\rangle)))))\).
Using \paperlink{premise:4c65616e2e4f6d6567612e4c696e656172436f6d626f2e6164645f6576616c}{Prerequisite~P\ref{premise:4c65616e2e4f6d6567612e4c696e656172436f6d626f2e6164645f6576616c}} with \(l_{1} = \operatorname{Lean\mathord{:}Omega\mathord{:}Linear\,Combo\mathord{:}mk}(0,\allowbreak{} \langle 0\rangle \mathbin{\Vert} (\langle 0\rangle \mathbin{\Vert} (\langle 0\rangle \mathbin{\Vert} (\langle 0\rangle \mathbin{\Vert} \langle 2\rangle))))\) and \(l_{2} = \operatorname{Lean\mathord{:}Omega\mathord{:}Linear\,Combo\mathord{:}mk}(1,\allowbreak{} \langle\,\rangle)\) and \(v = \operatorname{Lean\mathord{:}Omega\mathord{:}Coeffs\mathord{:}of\,List}(\langle |X \cap (Y \setminus X)|\rangle \mathbin{\Vert} (\langle |X \cup Y|\rangle \mathbin{\Vert} (\langle |X|\rangle \mathbin{\Vert} (\langle |Y \setminus X|\rangle \mathbin{\Vert} \langle n\rangle))))\) in reverse gives \(\operatorname{eval}(\operatorname{Lean\mathord{:}Omega\mathord{:}Linear\,Combo\mathord{:}mk}(0,\allowbreak{} \langle 0\rangle \mathbin{\Vert} (\langle 0\rangle \mathbin{\Vert} (\langle 0\rangle \mathbin{\Vert} (\langle 0\rangle \mathbin{\Vert} \langle 2\rangle)))),\allowbreak{} \operatorname{Lean\mathord{:}Omega\mathord{:}Coeffs\mathord{:}of\,List}(\langle |X \cap (Y \setminus X)|\rangle \mathbin{\Vert} (\langle |X \cup Y|\rangle \mathbin{\Vert} (\langle |X|\rangle \mathbin{\Vert} (\langle |Y \setminus X|\rangle \mathbin{\Vert} \langle n\rangle))))) + \operatorname{eval}(\operatorname{Lean\mathord{:}Omega\mathord{:}Linear\,Combo\mathord{:}mk}(1,\allowbreak{} \langle\,\rangle),\allowbreak{} \operatorname{Lean\mathord{:}Omega\mathord{:}Coeffs\mathord{:}of\,List}(\langle |X \cap (Y \setminus X)|\rangle \mathbin{\Vert} (\langle |X \cup Y|\rangle \mathbin{\Vert} (\langle |X|\rangle \mathbin{\Vert} (\langle |Y \setminus X|\rangle \mathbin{\Vert} \langle n\rangle))))) = \operatorname{eval}(\operatorname{Lean\mathord{:}Omega\mathord{:}Linear\,Combo\mathord{:}mk}(0,\allowbreak{} \langle 0\rangle \mathbin{\Vert} (\langle 0\rangle \mathbin{\Vert} (\langle 0\rangle \mathbin{\Vert} (\langle 0\rangle \mathbin{\Vert} \langle 2\rangle)))) + \operatorname{Lean\mathord{:}Omega\mathord{:}Linear\,Combo\mathord{:}mk}(1,\allowbreak{} \langle\,\rangle),\allowbreak{} \operatorname{Lean\mathord{:}Omega\mathord{:}Coeffs\mathord{:}of\,List}(\langle |X \cap (Y \setminus X)|\rangle \mathbin{\Vert} (\langle |X \cup Y|\rangle \mathbin{\Vert} (\langle |X|\rangle \mathbin{\Vert} (\langle |Y \setminus X|\rangle \mathbin{\Vert} \langle n\rangle)))))\).
Chaining the equalities, we obtain \(n + n + 1 = \operatorname{eval}(\operatorname{Lean\mathord{:}Omega\mathord{:}Linear\,Combo\mathord{:}mk}(0,\allowbreak{} \langle 0\rangle \mathbin{\Vert} (\langle 0\rangle \mathbin{\Vert} (\langle 0\rangle \mathbin{\Vert} (\langle 0\rangle \mathbin{\Vert} \langle 2\rangle)))) + \operatorname{Lean\mathord{:}Omega\mathord{:}Linear\,Combo\mathord{:}mk}(1,\allowbreak{} \langle\,\rangle),\allowbreak{} \operatorname{Lean\mathord{:}Omega\mathord{:}Coeffs\mathord{:}of\,List}(\langle |X \cap (Y \setminus X)|\rangle \mathbin{\Vert} (\langle |X \cup Y|\rangle \mathbin{\Vert} (\langle |X|\rangle \mathbin{\Vert} (\langle |Y \setminus X|\rangle \mathbin{\Vert} \langle n\rangle)))))\).
By \paperlink{premise:4c65616e2e4f6d6567612e496e742e7375625f636f6e6772}{Prerequisite~P\ref{premise:4c65616e2e4f6d6567612e496e742e7375625f636f6e6772}}, we obtain \(|X| + |Y \setminus X| - (n + n + 1) = \operatorname{eval}(\operatorname{Lean\mathord{:}Omega\mathord{:}Linear\,Combo\mathord{:}mk}(0,\allowbreak{} \langle 0\rangle \mathbin{\Vert} (\langle 0\rangle \mathbin{\Vert} \langle 1\rangle)) + \operatorname{Lean\mathord{:}Omega\mathord{:}Linear\,Combo\mathord{:}mk}(0,\allowbreak{} \langle 0\rangle \mathbin{\Vert} (\langle 0\rangle \mathbin{\Vert} (\langle 0\rangle \mathbin{\Vert} \langle 1\rangle))),\allowbreak{} \operatorname{Lean\mathord{:}Omega\mathord{:}Coeffs\mathord{:}of\,List}(\langle |X \cap (Y \setminus X)|\rangle \mathbin{\Vert} (\langle |X \cup Y|\rangle \mathbin{\Vert} (\langle |X|\rangle \mathbin{\Vert} (\langle |Y \setminus X|\rangle \mathbin{\Vert} \langle n\rangle))))) - \operatorname{eval}(\operatorname{Lean\mathord{:}Omega\mathord{:}Linear\,Combo\mathord{:}mk}(0,\allowbreak{} \langle 0\rangle \mathbin{\Vert} (\langle 0\rangle \mathbin{\Vert} (\langle 0\rangle \mathbin{\Vert} (\langle 0\rangle \mathbin{\Vert} \langle 2\rangle)))) + \operatorname{Lean\mathord{:}Omega\mathord{:}Linear\,Combo\mathord{:}mk}(1,\allowbreak{} \langle\,\rangle),\allowbreak{} \operatorname{Lean\mathord{:}Omega\mathord{:}Coeffs\mathord{:}of\,List}(\langle |X \cap (Y \setminus X)|\rangle \mathbin{\Vert} (\langle |X \cup Y|\rangle \mathbin{\Vert} (\langle |X|\rangle \mathbin{\Vert} (\langle |Y \setminus X|\rangle \mathbin{\Vert} \langle n\rangle)))))\).
Using \paperlink{premise:4c65616e2e4f6d6567612e4c696e656172436f6d626f2e7375625f6576616c}{Prerequisite~P\ref{premise:4c65616e2e4f6d6567612e4c696e656172436f6d626f2e7375625f6576616c}} with \(l_{1} = \operatorname{Lean\mathord{:}Omega\mathord{:}Linear\,Combo\mathord{:}mk}(0,\allowbreak{} \langle 0\rangle \mathbin{\Vert} (\langle 0\rangle \mathbin{\Vert} (\langle 1\rangle \mathbin{\Vert} \langle 1\rangle)))\) and \(l_{2} = \operatorname{Lean\mathord{:}Omega\mathord{:}Linear\,Combo\mathord{:}mk}(1,\allowbreak{} \langle 0\rangle \mathbin{\Vert} (\langle 0\rangle \mathbin{\Vert} (\langle 0\rangle \mathbin{\Vert} (\langle 0\rangle \mathbin{\Vert} \langle 2\rangle))))\) and \(v = \operatorname{Lean\mathord{:}Omega\mathord{:}Coeffs\mathord{:}of\,List}(\langle |X \cap (Y \setminus X)|\rangle \mathbin{\Vert} (\langle |X \cup Y|\rangle \mathbin{\Vert} (\langle |X|\rangle \mathbin{\Vert} (\langle |Y \setminus X|\rangle \mathbin{\Vert} \langle n\rangle))))\) in reverse gives \(\operatorname{eval}(\operatorname{Lean\mathord{:}Omega\mathord{:}Linear\,Combo\mathord{:}mk}(0,\allowbreak{} \langle 0\rangle \mathbin{\Vert} (\langle 0\rangle \mathbin{\Vert} (\langle 1\rangle \mathbin{\Vert} \langle 1\rangle))),\allowbreak{} \operatorname{Lean\mathord{:}Omega\mathord{:}Coeffs\mathord{:}of\,List}(\langle |X \cap (Y \setminus X)|\rangle \mathbin{\Vert} (\langle |X \cup Y|\rangle \mathbin{\Vert} (\langle |X|\rangle \mathbin{\Vert} (\langle |Y \setminus X|\rangle \mathbin{\Vert} \langle n\rangle))))) - \operatorname{eval}(\operatorname{Lean\mathord{:}Omega\mathord{:}Linear\,Combo\mathord{:}mk}(1,\allowbreak{} \langle 0\rangle \mathbin{\Vert} (\langle 0\rangle \mathbin{\Vert} (\langle 0\rangle \mathbin{\Vert} (\langle 0\rangle \mathbin{\Vert} \langle 2\rangle)))),\allowbreak{} \operatorname{Lean\mathord{:}Omega\mathord{:}Coeffs\mathord{:}of\,List}(\langle |X \cap (Y \setminus X)|\rangle \mathbin{\Vert} (\langle |X \cup Y|\rangle \mathbin{\Vert} (\langle |X|\rangle \mathbin{\Vert} (\langle |Y \setminus X|\rangle \mathbin{\Vert} \langle n\rangle))))) = \operatorname{eval}(\operatorname{Lean\mathord{:}Omega\mathord{:}Linear\,Combo\mathord{:}mk}(0,\allowbreak{} \langle 0\rangle \mathbin{\Vert} (\langle 0\rangle \mathbin{\Vert} (\langle 1\rangle \mathbin{\Vert} \langle 1\rangle))) - \operatorname{Lean\mathord{:}Omega\mathord{:}Linear\,Combo\mathord{:}mk}(1,\allowbreak{} \langle 0\rangle \mathbin{\Vert} (\langle 0\rangle \mathbin{\Vert} (\langle 0\rangle \mathbin{\Vert} (\langle 0\rangle \mathbin{\Vert} \langle 2\rangle)))),\allowbreak{} \operatorname{Lean\mathord{:}Omega\mathord{:}Coeffs\mathord{:}of\,List}(\langle |X \cap (Y \setminus X)|\rangle \mathbin{\Vert} (\langle |X \cup Y|\rangle \mathbin{\Vert} (\langle |X|\rangle \mathbin{\Vert} (\langle |Y \setminus X|\rangle \mathbin{\Vert} \langle n\rangle)))))\).
Chaining the equalities, we obtain \(|X| + |Y \setminus X| - (n + n + 1) = \operatorname{eval}(\operatorname{Lean\mathord{:}Omega\mathord{:}Linear\,Combo\mathord{:}mk}(0,\allowbreak{} \langle 0\rangle \mathbin{\Vert} (\langle 0\rangle \mathbin{\Vert} (\langle 1\rangle \mathbin{\Vert} \langle 1\rangle))) - \operatorname{Lean\mathord{:}Omega\mathord{:}Linear\,Combo\mathord{:}mk}(1,\allowbreak{} \langle 0\rangle \mathbin{\Vert} (\langle 0\rangle \mathbin{\Vert} (\langle 0\rangle \mathbin{\Vert} (\langle 0\rangle \mathbin{\Vert} \langle 2\rangle)))),\allowbreak{} \operatorname{Lean\mathord{:}Omega\mathord{:}Coeffs\mathord{:}of\,List}(\langle |X \cap (Y \setminus X)|\rangle \mathbin{\Vert} (\langle |X \cup Y|\rangle \mathbin{\Vert} (\langle |X|\rangle \mathbin{\Vert} (\langle |Y \setminus X|\rangle \mathbin{\Vert} \langle n\rangle)))))\).
By \paperlink{premise:6c655f6f665f6c655f6f665f6571}{Prerequisite~P\ref{premise:6c655f6f665f6c655f6f665f6571}} with \(\mathbb{Z}\) as its type parameter \(\alpha\), we obtain \(0 \le \operatorname{eval}(\operatorname{Lean\mathord{:}Omega\mathord{:}Linear\,Combo\mathord{:}mk}(0,\allowbreak{} \langle 0\rangle \mathbin{\Vert} (\langle 0\rangle \mathbin{\Vert} (\langle 1\rangle \mathbin{\Vert} \langle 1\rangle))) - \operatorname{Lean\mathord{:}Omega\mathord{:}Linear\,Combo\mathord{:}mk}(1,\allowbreak{} \langle 0\rangle \mathbin{\Vert} (\langle 0\rangle \mathbin{\Vert} (\langle 0\rangle \mathbin{\Vert} (\langle 0\rangle \mathbin{\Vert} \langle 2\rangle)))),\allowbreak{} \operatorname{Lean\mathord{:}Omega\mathord{:}Coeffs\mathord{:}of\,List}(\langle |X \cap (Y \setminus X)|\rangle \mathbin{\Vert} (\langle |X \cup Y|\rangle \mathbin{\Vert} (\langle |X|\rangle \mathbin{\Vert} (\langle |Y \setminus X|\rangle \mathbin{\Vert} \langle n\rangle)))))\).
By \paperlink{premise:4c65616e2e4f6d6567612e436f6e73747261696e742e616464496e657175616c6974795f736174}{Prerequisite~P\ref{premise:4c65616e2e4f6d6567612e436f6e73747261696e742e616464496e657175616c6974795f736174}}, we obtain \(\operatorname{sat^{\prime}}(\operatorname{Lean\mathord{:}Omega\mathord{:}Constraint\mathord{:}mk}(\operatorname{some}(-(-1)),\allowbreak{} \operatorname{none}),\allowbreak{} \langle 0\rangle \mathbin{\Vert} (\langle 0\rangle \mathbin{\Vert} (\langle 1\rangle \mathbin{\Vert} (\langle 1\rangle \mathbin{\Vert} \langle -2\rangle))),\allowbreak{} \operatorname{Lean\mathord{:}Omega\mathord{:}Coeffs\mathord{:}of\,List}(\langle |X \cap (Y \setminus X)|\rangle \mathbin{\Vert} (\langle |X \cup Y|\rangle \mathbin{\Vert} (\langle |X|\rangle \mathbin{\Vert} (\langle |Y \setminus X|\rangle \mathbin{\Vert} \langle n\rangle))))) = \mathrm{true}\).
Use the hypothesis \(\left|X\right| \le n \land \left|Y \setminus X\right| \le n\).
Taking the first conjunct, we obtain \(\left|X\right| \le n\).
By \paperlink{premise:4c65616e2e4f6d6567612e496e742e6f664e61745f6c655f6f665f6c65}{Prerequisite~P\ref{premise:4c65616e2e4f6d6567612e496e742e6f664e61745f6c655f6f665f6c65}}, we obtain \(\left|X\right| \le n\).
By \paperlink{premise:496e742e7375625f6e6f6e6e65675f6f665f6c65}{Prerequisite~P\ref{premise:496e742e7375625f6e6f6e6e65675f6f665f6c65}}, we obtain \(0 \le n - \left|X\right|\).
Using \paperlink{premise:4c65616e2e4f6d6567612e4c696e656172436f6d626f2e636f6f7264696e6174655f6576616c5f34}{Prerequisite~P\ref{premise:4c65616e2e4f6d6567612e4c696e656172436f6d626f2e636f6f7264696e6174655f6576616c5f34}} in reverse gives \(n = \operatorname{eval}(\operatorname{coordinate}(4),\allowbreak{} \operatorname{Lean\mathord{:}Omega\mathord{:}Coeffs\mathord{:}of\,List}(\langle |X \cap (Y \setminus X)|\rangle \mathbin{\Vert} (\langle |X \cup Y|\rangle \mathbin{\Vert} (\langle |X|\rangle \mathbin{\Vert} (\langle |Y \setminus X|\rangle \mathbin{\Vert} \langle n\rangle)))))\).
Using \paperlink{premise:4c65616e2e4f6d6567612e4c696e656172436f6d626f2e636f6f7264696e6174655f6576616c5f32}{Prerequisite~P\ref{premise:4c65616e2e4f6d6567612e4c696e656172436f6d626f2e636f6f7264696e6174655f6576616c5f32}} in reverse gives \(|X| = \operatorname{eval}(\operatorname{coordinate}(2),\allowbreak{} \operatorname{Lean\mathord{:}Omega\mathord{:}Coeffs\mathord{:}of\,List}(\langle |X \cap (Y \setminus X)|\rangle \mathbin{\Vert} (\langle |X \cup Y|\rangle \mathbin{\Vert} (\langle |X|\rangle \mathbin{\Vert} (\langle |Y \setminus X|\rangle \mathbin{\Vert} \langle n\rangle)))))\).
By \paperlink{premise:4c65616e2e4f6d6567612e496e742e7375625f636f6e6772}{Prerequisite~P\ref{premise:4c65616e2e4f6d6567612e496e742e7375625f636f6e6772}}, we obtain \(n - |X| = \operatorname{eval}(\operatorname{coordinate}(4),\allowbreak{} \operatorname{Lean\mathord{:}Omega\mathord{:}Coeffs\mathord{:}of\,List}(\langle |X \cap (Y \setminus X)|\rangle \mathbin{\Vert} (\langle |X \cup Y|\rangle \mathbin{\Vert} (\langle |X|\rangle \mathbin{\Vert} (\langle |Y \setminus X|\rangle \mathbin{\Vert} \langle n\rangle))))) - \operatorname{eval}(\operatorname{coordinate}(2),\allowbreak{} \operatorname{Lean\mathord{:}Omega\mathord{:}Coeffs\mathord{:}of\,List}(\langle |X \cap (Y \setminus X)|\rangle \mathbin{\Vert} (\langle |X \cup Y|\rangle \mathbin{\Vert} (\langle |X|\rangle \mathbin{\Vert} (\langle |Y \setminus X|\rangle \mathbin{\Vert} \langle n\rangle)))))\).
Using \paperlink{premise:4c65616e2e4f6d6567612e4c696e656172436f6d626f2e7375625f6576616c}{Prerequisite~P\ref{premise:4c65616e2e4f6d6567612e4c696e656172436f6d626f2e7375625f6576616c}} with \(l_{1} = \operatorname{Lean\mathord{:}Omega\mathord{:}Linear\,Combo\mathord{:}mk}(0,\allowbreak{} \langle 0\rangle \mathbin{\Vert} (\langle 0\rangle \mathbin{\Vert} (\langle 0\rangle \mathbin{\Vert} (\langle 0\rangle \mathbin{\Vert} \langle 1\rangle))))\) and \(l_{2} = \operatorname{Lean\mathord{:}Omega\mathord{:}Linear\,Combo\mathord{:}mk}(0,\allowbreak{} \langle 0\rangle \mathbin{\Vert} (\langle 0\rangle \mathbin{\Vert} \langle 1\rangle))\) and \(v = \operatorname{Lean\mathord{:}Omega\mathord{:}Coeffs\mathord{:}of\,List}(\langle |X \cap (Y \setminus X)|\rangle \mathbin{\Vert} (\langle |X \cup Y|\rangle \mathbin{\Vert} (\langle |X|\rangle \mathbin{\Vert} (\langle |Y \setminus X|\rangle \mathbin{\Vert} \langle n\rangle))))\) in reverse gives \(\operatorname{eval}(\operatorname{Lean\mathord{:}Omega\mathord{:}Linear\,Combo\mathord{:}mk}(0,\allowbreak{} \langle 0\rangle \mathbin{\Vert} (\langle 0\rangle \mathbin{\Vert} (\langle 0\rangle \mathbin{\Vert} (\langle 0\rangle \mathbin{\Vert} \langle 1\rangle)))),\allowbreak{} \operatorname{Lean\mathord{:}Omega\mathord{:}Coeffs\mathord{:}of\,List}(\langle |X \cap (Y \setminus X)|\rangle \mathbin{\Vert} (\langle |X \cup Y|\rangle \mathbin{\Vert} (\langle |X|\rangle \mathbin{\Vert} (\langle |Y \setminus X|\rangle \mathbin{\Vert} \langle n\rangle))))) - \operatorname{eval}(\operatorname{Lean\mathord{:}Omega\mathord{:}Linear\,Combo\mathord{:}mk}(0,\allowbreak{} \langle 0\rangle \mathbin{\Vert} (\langle 0\rangle \mathbin{\Vert} \langle 1\rangle)),\allowbreak{} \operatorname{Lean\mathord{:}Omega\mathord{:}Coeffs\mathord{:}of\,List}(\langle |X \cap (Y \setminus X)|\rangle \mathbin{\Vert} (\langle |X \cup Y|\rangle \mathbin{\Vert} (\langle |X|\rangle \mathbin{\Vert} (\langle |Y \setminus X|\rangle \mathbin{\Vert} \langle n\rangle))))) = \operatorname{eval}(\operatorname{Lean\mathord{:}Omega\mathord{:}Linear\,Combo\mathord{:}mk}(0,\allowbreak{} \langle 0\rangle \mathbin{\Vert} (\langle 0\rangle \mathbin{\Vert} (\langle 0\rangle \mathbin{\Vert} (\langle 0\rangle \mathbin{\Vert} \langle 1\rangle)))) - \operatorname{Lean\mathord{:}Omega\mathord{:}Linear\,Combo\mathord{:}mk}(0,\allowbreak{} \langle 0\rangle \mathbin{\Vert} (\langle 0\rangle \mathbin{\Vert} \langle 1\rangle)),\allowbreak{} \operatorname{Lean\mathord{:}Omega\mathord{:}Coeffs\mathord{:}of\,List}(\langle |X \cap (Y \setminus X)|\rangle \mathbin{\Vert} (\langle |X \cup Y|\rangle \mathbin{\Vert} (\langle |X|\rangle \mathbin{\Vert} (\langle |Y \setminus X|\rangle \mathbin{\Vert} \langle n\rangle)))))\).
Chaining the equalities, we obtain \(n - |X| = \operatorname{eval}(\operatorname{Lean\mathord{:}Omega\mathord{:}Linear\,Combo\mathord{:}mk}(0,\allowbreak{} \langle 0\rangle \mathbin{\Vert} (\langle 0\rangle \mathbin{\Vert} (\langle 0\rangle \mathbin{\Vert} (\langle 0\rangle \mathbin{\Vert} \langle 1\rangle)))) - \operatorname{Lean\mathord{:}Omega\mathord{:}Linear\,Combo\mathord{:}mk}(0,\allowbreak{} \langle 0\rangle \mathbin{\Vert} (\langle 0\rangle \mathbin{\Vert} \langle 1\rangle)),\allowbreak{} \operatorname{Lean\mathord{:}Omega\mathord{:}Coeffs\mathord{:}of\,List}(\langle |X \cap (Y \setminus X)|\rangle \mathbin{\Vert} (\langle |X \cup Y|\rangle \mathbin{\Vert} (\langle |X|\rangle \mathbin{\Vert} (\langle |Y \setminus X|\rangle \mathbin{\Vert} \langle n\rangle)))))\).
By \paperlink{premise:6c655f6f665f6c655f6f665f6571}{Prerequisite~P\ref{premise:6c655f6f665f6c655f6f665f6571}} with \(\mathbb{Z}\) as its type parameter \(\alpha\), we obtain \(0 \le \operatorname{eval}(\operatorname{Lean\mathord{:}Omega\mathord{:}Linear\,Combo\mathord{:}mk}(0,\allowbreak{} \langle 0\rangle \mathbin{\Vert} (\langle 0\rangle \mathbin{\Vert} (\langle 0\rangle \mathbin{\Vert} (\langle 0\rangle \mathbin{\Vert} \langle 1\rangle)))) - \operatorname{Lean\mathord{:}Omega\mathord{:}Linear\,Combo\mathord{:}mk}(0,\allowbreak{} \langle 0\rangle \mathbin{\Vert} (\langle 0\rangle \mathbin{\Vert} \langle 1\rangle)),\allowbreak{} \operatorname{Lean\mathord{:}Omega\mathord{:}Coeffs\mathord{:}of\,List}(\langle |X \cap (Y \setminus X)|\rangle \mathbin{\Vert} (\langle |X \cup Y|\rangle \mathbin{\Vert} (\langle |X|\rangle \mathbin{\Vert} (\langle |Y \setminus X|\rangle \mathbin{\Vert} \langle n\rangle)))))\).
By \paperlink{premise:4c65616e2e4f6d6567612e436f6e73747261696e742e616464496e657175616c6974795f736174}{Prerequisite~P\ref{premise:4c65616e2e4f6d6567612e436f6e73747261696e742e616464496e657175616c6974795f736174}}, we obtain \(\operatorname{sat^{\prime}}(\operatorname{Lean\mathord{:}Omega\mathord{:}Constraint\mathord{:}mk}(\operatorname{some}(-0),\allowbreak{} \operatorname{none}),\allowbreak{} \langle 0\rangle \mathbin{\Vert} (\langle 0\rangle \mathbin{\Vert} (\langle -1\rangle \mathbin{\Vert} (\langle 0\rangle \mathbin{\Vert} \langle 1\rangle))),\allowbreak{} \operatorname{Lean\mathord{:}Omega\mathord{:}Coeffs\mathord{:}of\,List}(\langle |X \cap (Y \setminus X)|\rangle \mathbin{\Vert} (\langle |X \cup Y|\rangle \mathbin{\Vert} (\langle |X|\rangle \mathbin{\Vert} (\langle |Y \setminus X|\rangle \mathbin{\Vert} \langle n\rangle))))) = \mathrm{true}\).
By \paperlink{premise:4c65616e2e4f6d6567612e746964795f736174}{Prerequisite~P\ref{premise:4c65616e2e4f6d6567612e746964795f736174}}, we obtain \(\operatorname{sat^{\prime}}(\operatorname{tidy\,Constraint}(\operatorname{Lean\mathord{:}Omega\mathord{:}Constraint\mathord{:}mk}(\operatorname{some}(0),\allowbreak{} \operatorname{none}),\allowbreak{} \langle 0\rangle \mathbin{\Vert} (\langle 0\rangle \mathbin{\Vert} (\langle -1\rangle \mathbin{\Vert} (\langle 0\rangle \mathbin{\Vert} \langle 1\rangle)))),\allowbreak{} \operatorname{tidy\,Coeffs}(\operatorname{Lean\mathord{:}Omega\mathord{:}Constraint\mathord{:}mk}(\operatorname{some}(0),\allowbreak{} \operatorname{none}),\allowbreak{} \langle 0\rangle \mathbin{\Vert} (\langle 0\rangle \mathbin{\Vert} (\langle -1\rangle \mathbin{\Vert} (\langle 0\rangle \mathbin{\Vert} \langle 1\rangle)))),\allowbreak{} \operatorname{Lean\mathord{:}Omega\mathord{:}Coeffs\mathord{:}of\,List}(\langle |X \cap (Y \setminus X)|\rangle \mathbin{\Vert} (\langle |X \cup Y|\rangle \mathbin{\Vert} (\langle |X|\rangle \mathbin{\Vert} (\langle |Y \setminus X|\rangle \mathbin{\Vert} \langle n\rangle))))) = \mathrm{true}\).
By \paperlink{premise:4c65616e2e4f6d6567612e636f6d626f5f73617427}{Prerequisite~P\ref{premise:4c65616e2e4f6d6567612e636f6d626f5f73617427}} with \(s = \operatorname{Lean\mathord{:}Omega\mathord{:}Constraint\mathord{:}mk}(\operatorname{some}(1),\allowbreak{} \operatorname{none})\) and \(t = \operatorname{Lean\mathord{:}Omega\mathord{:}Constraint\mathord{:}mk}(\operatorname{none},\allowbreak{} \operatorname{some}(0))\) and \(1\) in place of its parameter \(a\) and \(x = \langle 0\rangle \mathbin{\Vert} (\langle 0\rangle \mathbin{\Vert} (\langle 1\rangle \mathbin{\Vert} (\langle 1\rangle \mathbin{\Vert} \langle -2\rangle)))\) and \(b = -1\) and \(y = \langle 0\rangle \mathbin{\Vert} (\langle 0\rangle \mathbin{\Vert} (\langle 1\rangle \mathbin{\Vert} (\langle 0\rangle \mathbin{\Vert} \langle -1\rangle)))\) and \(v = \operatorname{Lean\mathord{:}Omega\mathord{:}Coeffs\mathord{:}of\,List}(\langle |X \cap (Y \setminus X)|\rangle \mathbin{\Vert} (\langle |X \cup Y|\rangle \mathbin{\Vert} (\langle |X|\rangle \mathbin{\Vert} (\langle |Y \setminus X|\rangle \mathbin{\Vert} \langle n\rangle))))\), we obtain \(\operatorname{sat^{\prime}}(\operatorname{Lean\mathord{:}Omega\mathord{:}Constraint\mathord{:}combo}(1,\allowbreak{} \operatorname{Lean\mathord{:}Omega\mathord{:}Constraint\mathord{:}mk}(\operatorname{some}(1),\allowbreak{} \operatorname{none}),\allowbreak{} -1,\allowbreak{} \operatorname{Lean\mathord{:}Omega\mathord{:}Constraint\mathord{:}mk}(\operatorname{none},\allowbreak{} \operatorname{some}(0))),\allowbreak{} \operatorname{Lean\mathord{:}Omega\mathord{:}Coeffs\mathord{:}combo}(1,\allowbreak{} \langle 0\rangle \mathbin{\Vert} (\langle 0\rangle \mathbin{\Vert} (\langle 1\rangle \mathbin{\Vert} (\langle 1\rangle \mathbin{\Vert} \langle -2\rangle))),\allowbreak{} -1,\allowbreak{} \langle 0\rangle \mathbin{\Vert} (\langle 0\rangle \mathbin{\Vert} (\langle 1\rangle \mathbin{\Vert} (\langle 0\rangle \mathbin{\Vert} \langle -1\rangle)))),\allowbreak{} \operatorname{Lean\mathord{:}Omega\mathord{:}Coeffs\mathord{:}of\,List}(\langle |X \cap (Y \setminus X)|\rangle \mathbin{\Vert} (\langle |X \cup Y|\rangle \mathbin{\Vert} (\langle |X|\rangle \mathbin{\Vert} (\langle |Y \setminus X|\rangle \mathbin{\Vert} \langle n\rangle))))) = \mathrm{true}\).
Use the hypothesis \(\left|X\right| \le n \land \left|Y \setminus X\right| \le n\).
Taking the second conjunct, we obtain \(\left|Y \setminus X\right| \le n\).
By \paperlink{premise:4c65616e2e4f6d6567612e496e742e6f664e61745f6c655f6f665f6c65}{Prerequisite~P\ref{premise:4c65616e2e4f6d6567612e496e742e6f664e61745f6c655f6f665f6c65}}, we obtain \(\left|Y \setminus X\right| \le n\).
By \paperlink{premise:496e742e7375625f6e6f6e6e65675f6f665f6c65}{Prerequisite~P\ref{premise:496e742e7375625f6e6f6e6e65675f6f665f6c65}}, we obtain \(0 \le n - \left|Y \setminus X\right|\).
Using \paperlink{premise:4c65616e2e4f6d6567612e4c696e656172436f6d626f2e636f6f7264696e6174655f6576616c5f34}{Prerequisite~P\ref{premise:4c65616e2e4f6d6567612e4c696e656172436f6d626f2e636f6f7264696e6174655f6576616c5f34}} in reverse gives \(n = \operatorname{eval}(\operatorname{coordinate}(4),\allowbreak{} \operatorname{Lean\mathord{:}Omega\mathord{:}Coeffs\mathord{:}of\,List}(\langle |X \cap (Y \setminus X)|\rangle \mathbin{\Vert} (\langle |X \cup Y|\rangle \mathbin{\Vert} (\langle |X|\rangle \mathbin{\Vert} (\langle |Y \setminus X|\rangle \mathbin{\Vert} \langle n\rangle)))))\).
Using \paperlink{premise:4c65616e2e4f6d6567612e4c696e656172436f6d626f2e636f6f7264696e6174655f6576616c5f33}{Prerequisite~P\ref{premise:4c65616e2e4f6d6567612e4c696e656172436f6d626f2e636f6f7264696e6174655f6576616c5f33}} in reverse gives \(|Y \setminus X| = \operatorname{eval}(\operatorname{coordinate}(3),\allowbreak{} \operatorname{Lean\mathord{:}Omega\mathord{:}Coeffs\mathord{:}of\,List}(\langle |X \cap (Y \setminus X)|\rangle \mathbin{\Vert} (\langle |X \cup Y|\rangle \mathbin{\Vert} (\langle |X|\rangle \mathbin{\Vert} (\langle |Y \setminus X|\rangle \mathbin{\Vert} \langle n\rangle)))))\).
By \paperlink{premise:4c65616e2e4f6d6567612e496e742e7375625f636f6e6772}{Prerequisite~P\ref{premise:4c65616e2e4f6d6567612e496e742e7375625f636f6e6772}}, we obtain \(n - |Y \setminus X| = \operatorname{eval}(\operatorname{coordinate}(4),\allowbreak{} \operatorname{Lean\mathord{:}Omega\mathord{:}Coeffs\mathord{:}of\,List}(\langle |X \cap (Y \setminus X)|\rangle \mathbin{\Vert} (\langle |X \cup Y|\rangle \mathbin{\Vert} (\langle |X|\rangle \mathbin{\Vert} (\langle |Y \setminus X|\rangle \mathbin{\Vert} \langle n\rangle))))) - \operatorname{eval}(\operatorname{coordinate}(3),\allowbreak{} \operatorname{Lean\mathord{:}Omega\mathord{:}Coeffs\mathord{:}of\,List}(\langle |X \cap (Y \setminus X)|\rangle \mathbin{\Vert} (\langle |X \cup Y|\rangle \mathbin{\Vert} (\langle |X|\rangle \mathbin{\Vert} (\langle |Y \setminus X|\rangle \mathbin{\Vert} \langle n\rangle)))))\).
Using \paperlink{premise:4c65616e2e4f6d6567612e4c696e656172436f6d626f2e7375625f6576616c}{Prerequisite~P\ref{premise:4c65616e2e4f6d6567612e4c696e656172436f6d626f2e7375625f6576616c}} with \(l_{1} = \operatorname{Lean\mathord{:}Omega\mathord{:}Linear\,Combo\mathord{:}mk}(0,\allowbreak{} \langle 0\rangle \mathbin{\Vert} (\langle 0\rangle \mathbin{\Vert} (\langle 0\rangle \mathbin{\Vert} (\langle 0\rangle \mathbin{\Vert} \langle 1\rangle))))\) and \(l_{2} = \operatorname{Lean\mathord{:}Omega\mathord{:}Linear\,Combo\mathord{:}mk}(0,\allowbreak{} \langle 0\rangle \mathbin{\Vert} (\langle 0\rangle \mathbin{\Vert} (\langle 0\rangle \mathbin{\Vert} \langle 1\rangle)))\) and \(v = \operatorname{Lean\mathord{:}Omega\mathord{:}Coeffs\mathord{:}of\,List}(\langle |X \cap (Y \setminus X)|\rangle \mathbin{\Vert} (\langle |X \cup Y|\rangle \mathbin{\Vert} (\langle |X|\rangle \mathbin{\Vert} (\langle |Y \setminus X|\rangle \mathbin{\Vert} \langle n\rangle))))\) in reverse gives \(\operatorname{eval}(\operatorname{Lean\mathord{:}Omega\mathord{:}Linear\,Combo\mathord{:}mk}(0,\allowbreak{} \langle 0\rangle \mathbin{\Vert} (\langle 0\rangle \mathbin{\Vert} (\langle 0\rangle \mathbin{\Vert} (\langle 0\rangle \mathbin{\Vert} \langle 1\rangle)))),\allowbreak{} \operatorname{Lean\mathord{:}Omega\mathord{:}Coeffs\mathord{:}of\,List}(\langle |X \cap (Y \setminus X)|\rangle \mathbin{\Vert} (\langle |X \cup Y|\rangle \mathbin{\Vert} (\langle |X|\rangle \mathbin{\Vert} (\langle |Y \setminus X|\rangle \mathbin{\Vert} \langle n\rangle))))) - \operatorname{eval}(\operatorname{Lean\mathord{:}Omega\mathord{:}Linear\,Combo\mathord{:}mk}(0,\allowbreak{} \langle 0\rangle \mathbin{\Vert} (\langle 0\rangle \mathbin{\Vert} (\langle 0\rangle \mathbin{\Vert} \langle 1\rangle))),\allowbreak{} \operatorname{Lean\mathord{:}Omega\mathord{:}Coeffs\mathord{:}of\,List}(\langle |X \cap (Y \setminus X)|\rangle \mathbin{\Vert} (\langle |X \cup Y|\rangle \mathbin{\Vert} (\langle |X|\rangle \mathbin{\Vert} (\langle |Y \setminus X|\rangle \mathbin{\Vert} \langle n\rangle))))) = \operatorname{eval}(\operatorname{Lean\mathord{:}Omega\mathord{:}Linear\,Combo\mathord{:}mk}(0,\allowbreak{} \langle 0\rangle \mathbin{\Vert} (\langle 0\rangle \mathbin{\Vert} (\langle 0\rangle \mathbin{\Vert} (\langle 0\rangle \mathbin{\Vert} \langle 1\rangle)))) - \operatorname{Lean\mathord{:}Omega\mathord{:}Linear\,Combo\mathord{:}mk}(0,\allowbreak{} \langle 0\rangle \mathbin{\Vert} (\langle 0\rangle \mathbin{\Vert} (\langle 0\rangle \mathbin{\Vert} \langle 1\rangle))),\allowbreak{} \operatorname{Lean\mathord{:}Omega\mathord{:}Coeffs\mathord{:}of\,List}(\langle |X \cap (Y \setminus X)|\rangle \mathbin{\Vert} (\langle |X \cup Y|\rangle \mathbin{\Vert} (\langle |X|\rangle \mathbin{\Vert} (\langle |Y \setminus X|\rangle \mathbin{\Vert} \langle n\rangle)))))\).
Chaining the equalities, we obtain \(n - |Y \setminus X| = \operatorname{eval}(\operatorname{Lean\mathord{:}Omega\mathord{:}Linear\,Combo\mathord{:}mk}(0,\allowbreak{} \langle 0\rangle \mathbin{\Vert} (\langle 0\rangle \mathbin{\Vert} (\langle 0\rangle \mathbin{\Vert} (\langle 0\rangle \mathbin{\Vert} \langle 1\rangle)))) - \operatorname{Lean\mathord{:}Omega\mathord{:}Linear\,Combo\mathord{:}mk}(0,\allowbreak{} \langle 0\rangle \mathbin{\Vert} (\langle 0\rangle \mathbin{\Vert} (\langle 0\rangle \mathbin{\Vert} \langle 1\rangle))),\allowbreak{} \operatorname{Lean\mathord{:}Omega\mathord{:}Coeffs\mathord{:}of\,List}(\langle |X \cap (Y \setminus X)|\rangle \mathbin{\Vert} (\langle |X \cup Y|\rangle \mathbin{\Vert} (\langle |X|\rangle \mathbin{\Vert} (\langle |Y \setminus X|\rangle \mathbin{\Vert} \langle n\rangle)))))\).
By \paperlink{premise:6c655f6f665f6c655f6f665f6571}{Prerequisite~P\ref{premise:6c655f6f665f6c655f6f665f6571}} with \(\mathbb{Z}\) as its type parameter \(\alpha\), we obtain \(0 \le \operatorname{eval}(\operatorname{Lean\mathord{:}Omega\mathord{:}Linear\,Combo\mathord{:}mk}(0,\allowbreak{} \langle 0\rangle \mathbin{\Vert} (\langle 0\rangle \mathbin{\Vert} (\langle 0\rangle \mathbin{\Vert} (\langle 0\rangle \mathbin{\Vert} \langle 1\rangle)))) - \operatorname{Lean\mathord{:}Omega\mathord{:}Linear\,Combo\mathord{:}mk}(0,\allowbreak{} \langle 0\rangle \mathbin{\Vert} (\langle 0\rangle \mathbin{\Vert} (\langle 0\rangle \mathbin{\Vert} \langle 1\rangle))),\allowbreak{} \operatorname{Lean\mathord{:}Omega\mathord{:}Coeffs\mathord{:}of\,List}(\langle |X \cap (Y \setminus X)|\rangle \mathbin{\Vert} (\langle |X \cup Y|\rangle \mathbin{\Vert} (\langle |X|\rangle \mathbin{\Vert} (\langle |Y \setminus X|\rangle \mathbin{\Vert} \langle n\rangle)))))\).
By \paperlink{premise:4c65616e2e4f6d6567612e436f6e73747261696e742e616464496e657175616c6974795f736174}{Prerequisite~P\ref{premise:4c65616e2e4f6d6567612e436f6e73747261696e742e616464496e657175616c6974795f736174}}, we obtain \(\operatorname{sat^{\prime}}(\operatorname{Lean\mathord{:}Omega\mathord{:}Constraint\mathord{:}mk}(\operatorname{some}(-0),\allowbreak{} \operatorname{none}),\allowbreak{} \langle 0\rangle \mathbin{\Vert} (\langle 0\rangle \mathbin{\Vert} (\langle 0\rangle \mathbin{\Vert} (\langle -1\rangle \mathbin{\Vert} \langle 1\rangle))),\allowbreak{} \operatorname{Lean\mathord{:}Omega\mathord{:}Coeffs\mathord{:}of\,List}(\langle |X \cap (Y \setminus X)|\rangle \mathbin{\Vert} (\langle |X \cup Y|\rangle \mathbin{\Vert} (\langle |X|\rangle \mathbin{\Vert} (\langle |Y \setminus X|\rangle \mathbin{\Vert} \langle n\rangle))))) = \mathrm{true}\).
By \paperlink{premise:4c65616e2e4f6d6567612e746964795f736174}{Prerequisite~P\ref{premise:4c65616e2e4f6d6567612e746964795f736174}}, we obtain \(\operatorname{sat^{\prime}}(\operatorname{tidy\,Constraint}(\operatorname{Lean\mathord{:}Omega\mathord{:}Constraint\mathord{:}mk}(\operatorname{some}(0),\allowbreak{} \operatorname{none}),\allowbreak{} \langle 0\rangle \mathbin{\Vert} (\langle 0\rangle \mathbin{\Vert} (\langle 0\rangle \mathbin{\Vert} (\langle -1\rangle \mathbin{\Vert} \langle 1\rangle)))),\allowbreak{} \operatorname{tidy\,Coeffs}(\operatorname{Lean\mathord{:}Omega\mathord{:}Constraint\mathord{:}mk}(\operatorname{some}(0),\allowbreak{} \operatorname{none}),\allowbreak{} \langle 0\rangle \mathbin{\Vert} (\langle 0\rangle \mathbin{\Vert} (\langle 0\rangle \mathbin{\Vert} (\langle -1\rangle \mathbin{\Vert} \langle 1\rangle)))),\allowbreak{} \operatorname{Lean\mathord{:}Omega\mathord{:}Coeffs\mathord{:}of\,List}(\langle |X \cap (Y \setminus X)|\rangle \mathbin{\Vert} (\langle |X \cup Y|\rangle \mathbin{\Vert} (\langle |X|\rangle \mathbin{\Vert} (\langle |Y \setminus X|\rangle \mathbin{\Vert} \langle n\rangle))))) = \mathrm{true}\).
By \paperlink{premise:4c65616e2e4f6d6567612e436f6e73747261696e742e636f6d62696e655f73617427}{Prerequisite~P\ref{premise:4c65616e2e4f6d6567612e436f6e73747261696e742e636f6d62696e655f73617427}}, we obtain \(\operatorname{sat^{\prime}}(\operatorname{combine}(\operatorname{Lean\mathord{:}Omega\mathord{:}Constraint\mathord{:}mk}(\operatorname{some}(1),\allowbreak{} \operatorname{none}),\allowbreak{} \operatorname{Lean\mathord{:}Omega\mathord{:}Constraint\mathord{:}mk}(\operatorname{none},\allowbreak{} \operatorname{some}(0))),\allowbreak{} \langle 0\rangle \mathbin{\Vert} (\langle 0\rangle \mathbin{\Vert} (\langle 0\rangle \mathbin{\Vert} (\langle 1\rangle \mathbin{\Vert} \langle -1\rangle))),\allowbreak{} \operatorname{Lean\mathord{:}Omega\mathord{:}Coeffs\mathord{:}of\,List}(\langle |X \cap (Y \setminus X)|\rangle \mathbin{\Vert} (\langle |X \cup Y|\rangle \mathbin{\Vert} (\langle |X|\rangle \mathbin{\Vert} (\langle |Y \setminus X|\rangle \mathbin{\Vert} \langle n\rangle))))) = \mathrm{true}\).
By \paperlink{premise:4c65616e2e4f6d6567612e436f6e73747261696e742e6e6f745f736174275f6f665f6973496d706f737369626c65}{Prerequisite~P\ref{premise:4c65616e2e4f6d6567612e436f6e73747261696e742e6e6f745f736174275f6f665f6973496d706f737369626c65}}, we obtain \(\mathrm{False}\).
\end{proofstep}
The contrary assumption leads to a contradiction. Hence \(\left|X\right| + \left|Y \setminus X\right| \le n + n\).
Rewriting with \paperlink{premise:696d706c6965735f636f6e6772}{Prerequisite~P\ref{premise:696d706c6965735f636f6e6772}}, \paperlink{premise:4d6174686c69622e5461637469632e507573684e65672e6e6f745f6f725f6571}{Prerequisite~P\ref{premise:4d6174686c69622e5461637469632e507573684e65672e6e6f745f6f725f6571}} and \paperlink{premise:4d6174686c69622e5461637469632e507573684e65672e6e6f745f6c745f6571}{Prerequisite~P\ref{premise:4d6174686c69622e5461637469632e507573684e65672e6e6f745f6c745f6571}} in the preceding result gives \(\neg (n < |X| \lor n < |Y \setminus X|) \implies \neg (n + n < |X| + |Y \setminus X|)\).
Use the hypothesis \(2 \cdot n < \left|X \cup Y\right|\).
Rewriting with \(\left|X \cup Y\right| = \left|X\right| + \left|Y \setminus X\right|\) in the preceding result gives \(2 \cdot n < \left|X\right| + \left|Y \setminus X\right|\).
Rewriting with \paperlink{premise:4e61742e74776f5f6d756c}{Prerequisite~P\ref{premise:4e61742e74776f5f6d756c}} in the preceding result gives \(n + n < \left|X\right| + \left|Y \setminus X\right|\).
By \paperlink{premise:4d6174686c69622e5461637469632e436f6e747261706f73652e6d7472}{Prerequisite~P\ref{premise:4d6174686c69622e5461637469632e436f6e747261706f73652e6d7472}}, we obtain \(n < \left|X\right| \lor n < \left|Y \setminus X\right|\).
Consider the following case.
\begin{proofstep}
Assume \(n < \left|X\right|\).
Choose the witness \(X\).
Use the hypothesis \(n < \left|X\right|\).
By \paperlink{premise:46696e7365742e5375627365742e7265666c}{Prerequisite~P\ref{premise:46696e7365742e5375627365742e7265666c}} with \(s = X\), we obtain \(X \subseteq X\).
The first alternative suffices. Hence \(X \subseteq X \lor X \subseteq Y\).
Combining the two conclusions, we obtain \(n < \left|X\right| \land \left(X \subseteq X \lor X \subseteq Y\right)\).
This witness proves existence. Hence \(\exists C\, :\, \mathcal{P}_{\mathrm{fin}}(\alpha),\; n < |C| \land (C \subseteq X \lor C \subseteq Y)\).
\end{proofstep}
Consider the following case.
\begin{proofstep}
Assume \(n < \left|Y \setminus X\right|\).
Choose the witness \(Y \setminus X\).
Use the hypothesis \(n < \left|Y \setminus X\right|\).
By \paperlink{premise:46696e7365742e73646966665f737562736574}{Prerequisite~P\ref{premise:46696e7365742e73646966665f737562736574}}, we obtain \(Y \setminus X \subseteq Y\).
The second alternative suffices. Hence \(Y \setminus X \subseteq X \lor Y \setminus X \subseteq Y\).
Combining the two conclusions, we obtain \(n < \left|Y \setminus X\right| \land \left(Y \setminus X \subseteq X \lor Y \setminus X \subseteq Y\right)\).
This witness proves existence. Hence \(\exists C\, :\, \mathcal{P}_{\mathrm{fin}}(\alpha),\; n < |C| \land (C \subseteq X \lor C \subseteq Y)\).
\end{proofstep}
The two alternatives exhaust the cases. Hence \(\exists C\, :\, \mathcal{P}_{\mathrm{fin}}(\alpha),\; n < |C| \land (C \subseteq X \lor C \subseteq Y)\).
\end{proof}

\begin{lemma}\label{theorem:106}
Let \(\alpha\) be a type. Let \(s\) be a finite subset of \(\alpha\). Then \(\left|s\right| = 1 \iff \left(\exists a\, :\, \alpha,\; s = \{a\}\right)\).
\end{lemma}

\begin{proof}
Consider the following case.
\begin{proofstep}
Fix \(\mathit{val}\) in \(\operatorname{Multiset}(\alpha)\).
Assume \(\mathit{val}\text{ has no repeated elements}\).
Assume \(s = \operatorname{as\,finite\,set}(\mathit{val})\).
The substitution uses the following argument.
\begin{proofstep}
Fix \(a\) in \(\alpha\).
By \paperlink{premise:46696e7365742e76616c5f696e6a}{Prerequisite~P\ref{premise:46696e7365742e76616c5f696e6a}}, we obtain \(\operatorname{underlying\,multiset}(\operatorname{as\,finite\,set}(\mathit{val})) = \operatorname{underlying\,multiset}(\{a\}) \iff \operatorname{as\,finite\,set}(\mathit{val}) = \{a\}\).
By propositional extensionality, we obtain \((\operatorname{underlying\,multiset}(\operatorname{as\,finite\,set}(\mathit{val})) = \operatorname{underlying\,multiset}(\{a\})) = (\operatorname{as\,finite\,set}(\mathit{val}) = \{a\})\).
Reversing the equality, we obtain \((\operatorname{as\,finite\,set}(\mathit{val}) = \{a\}) = (\operatorname{underlying\,multiset}(\operatorname{as\,finite\,set}(\mathit{val})) = \operatorname{underlying\,multiset}(\{a\}))\).
\end{proofstep}
Rewriting with \paperlink{premise:4d756c74697365742e636172645f65715f6f6e65}{Prerequisite~P\ref{premise:4d756c74697365742e636172645f65715f6f6e65}}, \paperlink{premise:66756e657874}{\emph{Function extensionality}} (\paperlink{premise:66756e657874}{Prerequisite~P\ref{premise:66756e657874}}), \paperlink{premise:6966665f73656c66}{Prerequisite~P\ref{premise:6966665f73656c66}} and the identity established above gives \(|\mathit{val}| = 1 \iff (\exists x\, :\, \alpha,\; \operatorname{as\,finite\,set}(\mathit{val}) = \{x\})\).
Use the hypothesis \(s = \operatorname{as\,finite\,set}(\mathit{val})\).
Reversing the equality, we obtain \(\operatorname{as\,finite\,set}(\mathit{val}) = s\).
Substituting equals, we obtain \(\left|s\right| = 1 \iff \left(\exists a\, :\, \alpha,\; s = \{a\}\right)\).
\end{proofstep}
By definition, we obtain \(s = s\).
A finite set is represented by a multiset without repetitions, we obtain \(\left|s\right| = 1 \iff \left(\exists a\, :\, \alpha,\; s = \{a\}\right)\).
\end{proof}

\begin{theorem}\label{theorem:107}
Let \(\alpha\) be a type. Assume equality on \(\alpha\) is decidable. Let \(s\) and \(t\) be finite subsets of \(\alpha\). Then
\begin{center}
\(\displaystyle (\exists a\, :\, \alpha,\; a \notin s \land \{a\} \cup s = t) \iff s \subseteq t \land |s| + 1 = |t|.\)
\end{center}
\end{theorem}

\begin{proof}
Consider the following case.
\begin{proofstep}
Assume \(\exists a_{1}\, :\, \alpha,\; a_{1} \notin s \land \{a_{1}\} \cup s = t\).
Use the hypothesis \(\exists a_{1}\, :\, \alpha,\; a_{1} \notin s \land \{a_{1}\} \cup s = t\).
Consider the following case.
\begin{proofstep}
Fix \(a_{1}\) in \(\alpha\).
Assume \(a_{1} \notin s \land \{a_{1}\} \cup s = t\).
Use the hypothesis \(a_{1} \notin s \land \{a_{1}\} \cup s = t\).
Consider the following case.
\begin{proofstep}
Assume \(a_{1} \notin s\).
Assume \(\{a_{1}\} \cup s = t\).
By \paperlink{premise:46696e7365742e7375627365745f696e73657274}{Prerequisite~P\ref{premise:46696e7365742e7375627365745f696e73657274}} with \(a_{1}\) in place of its parameter \(a\), we obtain \(s \subseteq \{a_{1}\} \cup s\).
Using \paperlink{theorem:17}{Lemma~\ref{theorem:17}} in reverse under \(a_{1} \notin s\) gives \(\left|s\right| + 1 = \left|\{a_{1}\} \cup s\right|\).
Combining the two conclusions, we obtain \(s \subseteq \{a_{1}\} \cup s \land \left|s\right| + 1 = \left|\{a_{1}\} \cup s\right|\).
Use the hypothesis \(\{a_{1}\} \cup s = t\).
Substituting equals, we obtain \(s \subseteq t \land \left|s\right| + 1 = \left|t\right|\).
\end{proofstep}
Using the two conjuncts, we obtain \(s \subseteq t \land \left|s\right| + 1 = \left|t\right|\).
\end{proofstep}
Using the existential hypothesis, we obtain \(s \subseteq t \land \left|s\right| + 1 = \left|t\right|\).
\end{proofstep}
Consider the following case.
\begin{proofstep}
Assume \(s \subseteq t \land \left|s\right| + 1 = \left|t\right|\).
Use the hypothesis \(s \subseteq t \land \left|s\right| + 1 = \left|t\right|\).
Consider the following case.
\begin{proofstep}
Assume \(s \subseteq t\).
Assume \(\left|s\right| + 1 = \left|t\right|\).
By \paperlink{theorem:106}{Lemma~\ref{theorem:106}}, we obtain \(\left|t \setminus s\right| = 1 \iff \left(\exists a_{1}\, :\, \alpha,\; t \setminus s = \{a_{1}\}\right)\).
Rewriting with \paperlink{theorem:90}{Lemma~\ref{theorem:90}} (using \(s \subseteq t\)) gives \(\left|t \setminus s\right| = \left|t\right| - \left|s\right|\).
Rewriting with \(\left|s\right| + 1 = \left|t\right|\) gives \(\left|t\right| - \left|s\right| = \left|s\right| + 1 - \left|s\right|\).
Rewriting with \paperlink{premise:4e61742e6164645f7375625f63616e63656c5f6c656674}{Prerequisite~P\ref{premise:4e61742e6164645f7375625f63616e63656c5f6c656674}} gives \(\left|s\right| + 1 - \left|s\right| = 1\).
By the forward implication, we obtain \(\exists a_{1}\, :\, \alpha,\; t \setminus s = \{a_{1}\}\).
Consider the following case.
\begin{proofstep}
Fix \(a_{1}\) in \(\alpha\).
Assume \(t \setminus s = \{a_{1}\}\).
Choose the witness \(a_{1}\).
For the required implication, argue as follows.
\begin{proofstep}
Assume \(a_{1} \in s\).
Use the hypothesis \(t \setminus s = \{a_{1}\}\).
Reversing the equality, we obtain \(\{a_{1}\} = t \setminus s\).
Applying the same function to equal arguments, we obtain \(\left(a_{1} \notin \{a_{1}\}\right) = \left(a_{1} \notin t \setminus s\right)\).
Use the hypothesis \(a_{1} \in s\).
By \paperlink{premise:46696e7365742e6e6f745f6d656d5f73646966665f6f665f6d656d5f7269676874}{Prerequisite~P\ref{premise:46696e7365742e6e6f745f6d656d5f73646966665f6f665f6d656d5f7269676874}}, we obtain \(a_{1} \notin t \setminus s\).
By \paperlink{premise:46696e7365742e6d656d5f73696e676c65746f6e5f73656c66}{Prerequisite~P\ref{premise:46696e7365742e6d656d5f73696e676c65746f6e5f73656c66}} with \(a_{1}\) in place of its parameter \(a\), we obtain \(a_{1} \in \{a_{1}\}\).
Substituting along the equality, we obtain \(\mathrm{False}\).
\end{proofstep}
Rewriting with \paperlink{premise:46696e7365742e73646966665f756e696f6e5f6f665f737562736574}{Prerequisite~P\ref{premise:46696e7365742e73646966665f756e696f6e5f6f665f737562736574}} (using \(s \subseteq t\)), \(t \setminus s = \{a_{1}\}\) and \paperlink{premise:46696e7365742e696e736572745f6571}{Prerequisite~P\ref{premise:46696e7365742e696e736572745f6571}} gives \(\{a_{1}\} \cup s = t\).
Combining the two conclusions, we obtain \(a_{1} \notin s \land \{a_{1}\} \cup s = t\).
This witness proves existence. Hence \(\exists a_{2}\, :\, \alpha,\; a_{2} \notin s \land \{a_{2}\} \cup s = t\).
\end{proofstep}
Using the existential hypothesis, we obtain \(\exists a_{1}\, :\, \alpha,\; a_{1} \notin s \land \{a_{1}\} \cup s = t\).
\end{proofstep}
Using the two conjuncts, we obtain \(\exists a_{1}\, :\, \alpha,\; a_{1} \notin s \land \{a_{1}\} \cup s = t\).
\end{proofstep}
Together the two implications give the equivalence. Hence \((\exists a\, :\, \alpha,\; a \notin s \land \{a\} \cup s = t) \iff s \subseteq t \land |s| + 1 = |t|\).
\end{proof}

\begin{lemma}\label{theorem:108}
Let \(\alpha\) be a type. Let \(s\) be a finite subset of \(\alpha\). Then
\begin{center}
\(\displaystyle |s| \le 1 \iff (\forall a\, :\, \alpha,\; a \in s \implies (\forall b\, :\, \alpha,\; b \in s \implies a = b)).\)
\end{center}
\end{lemma}

\begin{proof}
By \paperlink{premise:46696e7365742e65715f656d7074795f6f725f6e6f6e656d707479}{Prerequisite~P\ref{premise:46696e7365742e65715f656d7074795f6f725f6e6f6e656d707479}}, we obtain \(s = \varnothing \lor s\ne\varnothing\).
Consider the following case.
\begin{proofstep}
Assume \(s = \varnothing\).
The substitution uses the following argument.
\begin{proofstep}
Fix \(a\) in \(\alpha\).
By \paperlink{premise:46696e7365742e6e6f745f6d656d5f656d707479}{Prerequisite~P\ref{premise:46696e7365742e6e6f745f6d656d5f656d707479}}, we obtain \(a \notin \varnothing\).
By \paperlink{premise:65715f66616c7365}{Prerequisite~P\ref{premise:65715f66616c7365}}, we obtain \(\left(a \in \varnothing\right) = \mathrm{False}\).
For the required implication, argue as follows.
\begin{proofstep}
Fix \(b\) in \(\alpha\).
By \paperlink{premise:46696e7365742e6e6f745f6d656d5f656d707479}{Prerequisite~P\ref{premise:46696e7365742e6e6f745f6d656d5f656d707479}} with \(b\) in place of its parameter \(a\), we obtain \(b \notin \varnothing\).
By \paperlink{premise:65715f66616c7365}{Prerequisite~P\ref{premise:65715f66616c7365}}, we obtain \(\left(b \in \varnothing\right) = \mathrm{False}\).
By definition, we obtain \(\left(a = b\right) = \left(a = b\right)\).
By \paperlink{premise:696d706c6965735f636f6e6772}{Prerequisite~P\ref{premise:696d706c6965735f636f6e6772}} with \(p_{1} = b \in \varnothing\) and \(p_{2} = \mathrm{False}\) and \(q_{1} = a = b\) and \(q_{2} = a = b\), we obtain \(\left(b \in \varnothing \implies a = b\right) = \left(\mathrm{False} \implies a = b\right)\).
By \paperlink{premise:696e73744973456d70747946616c7365}{Prerequisite~P\ref{premise:696e73744973456d70747946616c7365}}, we obtain \(\mathrm{False}\text{ is empty}\).
By \paperlink{premise:4973456d7074792e666f72616c6c5f696666}{Prerequisite~P\ref{premise:4973456d7074792e666f72616c6c5f696666}} with \(\mathrm{False}\) as its type parameter \(\alpha\), we obtain \((\forall a_{1}\, :\, \mathrm{False},\; (a_{2} \mapsto a = b)(a_{1})) \iff \mathrm{True}\).
By propositional extensionality, we obtain \(\left(\mathrm{False} \implies a = b\right) = \mathrm{True}\).
Chaining the equalities, we obtain \(\left(b \in \varnothing \implies a = b\right) = \mathrm{True}\).
\end{proofstep}
By \paperlink{premise:666f72616c6c5f636f6e6772}{Prerequisite~P\ref{premise:666f72616c6c5f636f6e6772}}, we obtain \((\forall a_{1}\, :\, \alpha,\; a_{1} \in \varnothing \implies a = a_{1}) = (\forall a_{1}\, :\, \alpha,\; \mathrm{True})\).
By \paperlink{premise:696d706c6965735f74727565}{Prerequisite~P\ref{premise:696d706c6965735f74727565}}, we obtain \(\left(\forall a_{1}\, :\, \alpha,\; \mathrm{True}\right) = \mathrm{True}\).
Chaining the equalities, we obtain \(\left(\forall a_{1}\, :\, \alpha,\; a_{1} \in \varnothing \implies a = a_{1}\right) = \mathrm{True}\).
By \paperlink{premise:696d706c6965735f636f6e6772}{Prerequisite~P\ref{premise:696d706c6965735f636f6e6772}} with \(p_{1} = a \in \varnothing\) and \(p_{2} = \mathrm{False}\) and \(q_{1} = \forall a_{1}\, :\, \alpha,\; a_{1} \in \varnothing \implies a = a_{1}\) and \(q_{2} = \mathrm{True}\), we obtain \((a \in \varnothing \implies (\forall a_{2}\, :\, \alpha,\; a_{2} \in \varnothing \implies a = a_{2})) = (\mathrm{False} \implies \mathrm{True})\).
By \paperlink{premise:696d706c6965735f74727565}{Prerequisite~P\ref{premise:696d706c6965735f74727565}} with \(\mathrm{False}\) as its type parameter \(\alpha\), we obtain \(\left(\mathrm{False} \implies \mathrm{True}\right) = \mathrm{True}\).
Chaining the equalities, we obtain \((a \in \varnothing \implies (\forall a_{2}\, :\, \alpha,\; a_{2} \in \varnothing \implies a = a_{2})) = \mathrm{True}\).
\end{proofstep}
Rewriting with \paperlink{premise:4e61742e7a65726f5f6c65}{Prerequisite~P\ref{premise:4e61742e7a65726f5f6c65}}, \paperlink{premise:666f72616c6c5f636f6e6772}{Prerequisite~P\ref{premise:666f72616c6c5f636f6e6772}}, \paperlink{premise:696d706c6965735f74727565}{Prerequisite~P\ref{premise:696d706c6965735f74727565}}, \paperlink{premise:6966665f73656c66}{Prerequisite~P\ref{premise:6966665f73656c66}} and the identity established above gives \(0 \le 1 \iff (\forall a\, :\, \alpha,\; a \in \varnothing \implies (\forall a_{2}\, :\, \alpha,\; a_{2} \in \varnothing \implies a = a_{2}))\).
Use the hypothesis \(s = \varnothing\).
Reversing the equality, we obtain \(\varnothing = s\).
Substituting equals, we obtain \(|s| \le 1 \iff (\forall a\, :\, \alpha,\; a \in s \implies (\forall b\, :\, \alpha,\; b \in s \implies a = b))\).
\end{proofstep}
Consider the following case.
\begin{proofstep}
Assume \(s\ne\varnothing\).
Use the hypothesis \(s\ne\varnothing\).
Consider the following case.
\begin{proofstep}
Fix \(x\) in \(\alpha\).
Assume \(x \in s\).
By \paperlink{theorem:9}{Lemma~\ref{theorem:9}}, we obtain \(0 < \left|s\right| \iff s\ne\varnothing\).
Choose the witness \(x\).
Use the hypothesis \(x \in s\).
This witness proves existence. Hence \(\exists x_{1}\, :\, \alpha,\; x_{1} \in s\).
By the reverse implication, we obtain \(0 < \left|s\right|\).
By \paperlink{premise:4e61742e737563635f6c655f6f665f6c74}{Prerequisite~P\ref{premise:4e61742e737563635f6c655f6f665f6c74}}, we obtain \(0 + 1 \le \left|s\right|\).
By \paperlink{premise:4c452e6c652e6c655f6966665f6571}{Prerequisite~P\ref{premise:4c452e6c652e6c655f6966665f6571}} with \(\mathbb{N}\) as its type parameter \(\alpha\), we obtain \(\left|s\right| \le 0 + 1 \iff \left|s\right| = 0 + 1\).
By \paperlink{theorem:106}{Lemma~\ref{theorem:106}}, we obtain \(\left|s\right| = 1 \iff \left(\exists a\, :\, \alpha,\; s = \{a\}\right)\).
Consider the following case.
\begin{proofstep}
Assume \(\exists a_{1}\, :\, \alpha,\; s = \{a_{1}\}\).
Use the hypothesis \(\exists a_{1}\, :\, \alpha,\; s = \{a_{1}\}\).
Consider the following case.
\begin{proofstep}
Fix \(y\) in \(\alpha\).
Assume \(s = \{y\}\).
For the required implication, argue as follows.
\begin{proofstep}
Assume \(x \in \{y\}\).
The substitution uses the following argument.
\begin{proofstep}
Fix \(a_{1}\) in \(\alpha\).
By \paperlink{premise:46696e7365742e6d656d5f73696e676c65746f6e}{Prerequisite~P\ref{premise:46696e7365742e6d656d5f73696e676c65746f6e}}, we obtain \(a_{1} \in \{y\} \iff a_{1} = y\).
By propositional extensionality, we obtain \(\left(a_{1} \in \{y\}\right) = \left(a_{1} = y\right)\).
For the required implication, argue as follows.
\begin{proofstep}
Fix \(b\) in \(\alpha\).
By \paperlink{premise:46696e7365742e6d656d5f73696e676c65746f6e}{Prerequisite~P\ref{premise:46696e7365742e6d656d5f73696e676c65746f6e}}, we obtain \(b \in \{y\} \iff b = y\).
By propositional extensionality, we obtain \(\left(b \in \{y\}\right) = \left(b = y\right)\).
By definition, we obtain \(\left(a_{1} = b\right) = \left(a_{1} = b\right)\).
By \paperlink{premise:696d706c6965735f636f6e6772}{Prerequisite~P\ref{premise:696d706c6965735f636f6e6772}} with \(p_{1} = b \in \{y\}\) and \(p_{2} = b = y\) and \(q_{1} = a_{1} = b\) and \(q_{2} = a_{1} = b\), we obtain \(\left(b \in \{y\} \implies a_{1} = b\right) = \left(b = y \implies a_{1} = b\right)\).
\end{proofstep}
By \paperlink{premise:666f72616c6c5f636f6e6772}{Prerequisite~P\ref{premise:666f72616c6c5f636f6e6772}}, we obtain \((\forall a_{2}\, :\, \alpha,\; a_{2} \in \{y\} \implies a_{1} = a_{2}) = (\forall a_{2}\, :\, \alpha,\; a_{2} = y \implies a_{1} = a_{2})\).
By \paperlink{premise:666f72616c6c5f6571}{Prerequisite~P\ref{premise:666f72616c6c5f6571}}, we obtain \(\left(\forall a_{2}\, :\, \alpha,\; a_{2} = y \implies a_{1} = a_{2}\right) \iff a_{1} = y\).
By propositional extensionality, we obtain \(\left(\forall a_{2}\, :\, \alpha,\; a_{2} = y \implies a_{1} = a_{2}\right) = \left(a_{1} = y\right)\).
Chaining the equalities, we obtain \(\left(\forall a_{2}\, :\, \alpha,\; a_{2} \in \{y\} \implies a_{1} = a_{2}\right) = \left(a_{1} = y\right)\).
By \paperlink{premise:696d706c6965735f636f6e6772}{Prerequisite~P\ref{premise:696d706c6965735f636f6e6772}} with \(p_{1} = a_{1} \in \{y\}\) and \(p_{2} = a_{1} = y\) and \(q_{1} = \forall a_{2}\, :\, \alpha,\; a_{2} \in \{y\} \implies a_{1} = a_{2}\) and \(q_{2} = a_{1} = y\), we obtain \((a_{1} \in \{y\} \implies (\forall a_{3}\, :\, \alpha,\; a_{3} \in \{y\} \implies a_{1} = a_{3})) = (a_{1} = y \implies a_{1} = y)\).
By \paperlink{premise:696d705f73656c66}{Prerequisite~P\ref{premise:696d705f73656c66}}, we obtain \(a_{1} = y \implies a_{1} = y \iff \mathrm{True}\).
By propositional extensionality, we obtain \(\left(a_{1} = y \implies a_{1} = y\right) = \mathrm{True}\).
Chaining the equalities, we obtain \((a_{1} \in \{y\} \implies (\forall a_{3}\, :\, \alpha,\; a_{3} \in \{y\} \implies a_{1} = a_{3})) = \mathrm{True}\).
\end{proofstep}
Rewriting with \paperlink{premise:666f72616c6c5f636f6e6772}{Prerequisite~P\ref{premise:666f72616c6c5f636f6e6772}}, \paperlink{premise:696d706c6965735f74727565}{Prerequisite~P\ref{premise:696d706c6965735f74727565}} and the identity established above gives \(\forall a_{1}\, :\, \alpha,\; a_{1} \in \{y\} \implies (\forall a_{3}\, :\, \alpha,\; a_{3} \in \{y\} \implies a_{1} = a_{3})\).
\end{proofstep}
Use the hypothesis \(s = \{y\}\).
Reversing the equality, we obtain \(\{y\} = s\).
Use the hypothesis \(x \in s\).
Substituting equals, we obtain \(\forall a_{1}\, :\, \alpha,\; a_{1} \in s \implies (\forall b\, :\, \alpha,\; b \in s \implies a_{1} = b)\).
\end{proofstep}
Using the existential hypothesis, we obtain \(\forall a_{1}\, :\, \alpha,\; a_{1} \in s \implies (\forall b\, :\, \alpha,\; b \in s \implies a_{1} = b)\).
\end{proofstep}
Consider the following case.
\begin{proofstep}
Assume \(\forall a\, :\, \alpha,\; a \in s \implies \left(\forall b\, :\, \alpha,\; b \in s \implies a = b\right)\).
Choose the witness \(x\).
By \paperlink{premise:46696e7365742e65715f73696e676c65746f6e5f6966665f756e697175655f6d656d}{Prerequisite~P\ref{premise:46696e7365742e65715f73696e676c65746f6e5f6966665f756e697175655f6d656d}}, we obtain \(s = \{x\} \iff x \in s \land \left(\forall x_{1}\, :\, \alpha,\; x_{1} \in s \implies x_{1} = x\right)\).
Use the hypothesis \(x \in s\).
For the required implication, argue as follows.
\begin{proofstep}
Fix \(y\) in \(\alpha\).
Assume \(y \in s\).
Use the hypothesis \(y \in s\).
Use the hypothesis \(x \in s\).
Applying the hypothesis \(\forall a\, :\, \alpha,\; a \in s \implies \left(\forall b\, :\, \alpha,\; b \in s \implies a = b\right)\), we obtain \(y = x\).
\end{proofstep}
Combining the two conclusions, we obtain \(x \in s \land \left(\forall x_{1}\, :\, \alpha,\; x_{1} \in s \implies x_{1} = x\right)\).
By the reverse implication, we obtain \(s = \{x\}\).
This witness proves existence. Hence \(\exists a\, :\, \alpha,\; s = \{a\}\).
\end{proofstep}
Together the two implications give the equivalence. Hence \((\exists a\, :\, \alpha,\; s = \{a\}) \iff (\forall a\, :\, \alpha,\; a \in s \implies (\forall b\, :\, \alpha,\; b \in s \implies a = b))\).
Chaining the equivalences, we obtain \(|s| = 1 \iff (\forall a\, :\, \alpha,\; a \in s \implies (\forall b\, :\, \alpha,\; b \in s \implies a = b))\).
Chaining the equivalences, we obtain \(|s| \le 0 + 1 \iff (\forall a\, :\, \alpha,\; a \in s \implies (\forall b\, :\, \alpha,\; b \in s \implies a = b))\).
\end{proofstep}
Using the existential hypothesis, we obtain \(|s| \le 1 \iff (\forall a\, :\, \alpha,\; a \in s \implies (\forall b\, :\, \alpha,\; b \in s \implies a = b))\).
\end{proofstep}
The two alternatives exhaust the cases. Hence \(|s| \le 1 \iff (\forall a\, :\, \alpha,\; a \in s \implies (\forall b\, :\, \alpha,\; b \in s \implies a = b))\).
\end{proof}

\begin{lemma}\label{theorem:109}
Let \(\alpha\) be a type. Let \(s\) be a finite subset of \(\alpha\). Then
\begin{center}
\(\displaystyle |s| \le 1 \iff (\forall a\, :\, \alpha,\; \forall b\, :\, \alpha,\; a \in s \implies b \in s \implies a = b).\)
\end{center}
\end{lemma}

\begin{proof}
Consider the following case.
\begin{proofstep}
Assume \(\forall a_{1}\, :\, \alpha,\; a_{1} \in s \implies (\forall b\, :\, \alpha,\; b \in s \implies a_{1} = b)\).
Fix \(a_{1}\) in \(\alpha\).
Fix \(b\) in \(\alpha\).
Assume \(a_{1} \in s\).
Assume \(b \in s\).
Use the hypothesis \(a_{1} \in s\).
Use the hypothesis \(b \in s\).
Applying the hypothesis \(\forall a_{1}\, :\, \alpha,\; a_{1} \in s \implies (\forall b\, :\, \alpha,\; b \in s \implies a_{1} = b)\), we obtain \(a_{1} = b\).
\end{proofstep}
Consider the following case.
\begin{proofstep}
Assume \(\forall a_{1}\, :\, \alpha,\; \forall b\, :\, \alpha,\; a_{1} \in s \implies b \in s \implies a_{1} = b\).
Fix \(a_{1}\) in \(\alpha\).
Assume \(a_{1} \in s\).
Fix \(b\) in \(\alpha\).
Assume \(b \in s\).
Use the hypothesis \(a_{1} \in s\).
Use the hypothesis \(b \in s\).
Applying the hypothesis \(\forall a_{1}\, :\, \alpha,\; \forall b\, :\, \alpha,\; a_{1} \in s \implies b \in s \implies a_{1} = b\), we obtain \(a_{1} = b\).
\end{proofstep}
Together the two implications give the equivalence. Hence \((\forall a\, :\, \alpha,\; a \in s \implies (\forall b\, :\, \alpha,\; b \in s \implies a = b)) \iff (\forall a\, :\, \alpha,\; \forall b\, :\, \alpha,\; a \in s \implies b \in s \implies a = b)\).
Rewriting with \paperlink{theorem:108}{Lemma~\ref{theorem:108}} in the preceding result gives \(|s| \le 1 \iff (\forall a\, :\, \alpha,\; \forall b\, :\, \alpha,\; a \in s \implies b \in s \implies a = b)\).
\end{proof}

\begin{lemma}\label{theorem:110}
Let \(\alpha\) be a type. Let \(s\) be a finite subset of \(\alpha\). Then \(\left|s\right| \le 1 \iff s\text{ has at most one element}\).
\end{lemma}

\begin{proof}
By \paperlink{theorem:108}{Lemma~\ref{theorem:108}}, we obtain \(|s| \le 1 \iff (\forall a\, :\, \alpha,\; a \in s \implies (\forall b\, :\, \alpha,\; b \in s \implies a = b))\).
By \paperlink{premise:5365742e73756273696e676c65746f6e5f636f65}{Prerequisite~P\ref{premise:5365742e73756273696e676c65746f6e5f636f65}} with \(\operatorname{underlying\,set}(s)\) in place of its parameter \(s\), we obtain \(s\text{ has at most one element} \iff \operatorname{Set\mathord{:}Subsingleton}(\operatorname{underlying\,set}(s))\).
Reversing the equivalence, we obtain \(\operatorname{Set\mathord{:}Subsingleton}(\operatorname{underlying\,set}(s)) \iff s\text{ has at most one element}\).
Chaining the equivalences, we obtain \(\left|s\right| \le 1 \iff s\text{ has at most one element}\).
\end{proof}

\begin{lemma}\label{theorem:111}
Let \(\alpha\) be a type. Let \(s\) be a finite subset of \(\alpha\). Assume \(\alpha\text{ is nonempty}\). Then \(\left|s\right| \le 1 \iff \left(\exists x\, :\, \alpha,\; s \subseteq \{x\}\right)\).
\end{lemma}

\begin{proof}
Consider the following case.
\begin{proofstep}
Assume \(\left|s\right| \le 1\).
By \paperlink{premise:46696e7365742e65715f656d7074795f6f725f6e6f6e656d707479}{Prerequisite~P\ref{premise:46696e7365742e65715f656d7074795f6f725f6e6f6e656d707479}}, we obtain \(s = \varnothing \lor s\ne\varnothing\).
Consider the following case.
\begin{proofstep}
Assume \(s = \varnothing\).
For the required implication, argue as follows.
\begin{proofstep}
Assume \(\left|\varnothing\right| \le 1\).
Choose the witness \(\text{a fixed element of }\alpha\).
Use the hypothesis \(\alpha\text{ is nonempty}\).
By \paperlink{premise:46696e7365742e656d7074795f737562736574}{Prerequisite~P\ref{premise:46696e7365742e656d7074795f737562736574}} with \(\{\text{a fixed element of }\alpha\}\) in place of its parameter \(s\), we obtain \(\varnothing \subseteq \{\text{a fixed element of }\alpha\}\).
This witness proves existence. Hence \(\exists x\, :\, \alpha,\; \varnothing \subseteq \{x\}\).
\end{proofstep}
Use the hypothesis \(s = \varnothing\).
Reversing the equality, we obtain \(\varnothing = s\).
Use the hypothesis \(\left|s\right| \le 1\).
Substituting equals, we obtain \(\exists x\, :\, \alpha,\; s \subseteq \{x\}\).
\end{proofstep}
Consider the following case.
\begin{proofstep}
Assume \(s\ne\varnothing\).
Use the hypothesis \(s\ne\varnothing\).
Consider the following case.
\begin{proofstep}
Fix \(x\) in \(\alpha\).
Assume \(x \in s\).
Choose the witness \(x\).
For the required implication, argue as follows.
\begin{proofstep}
Fix \(y\) in \(\alpha\).
Assume \(y \in s\).
Rewriting with \paperlink{premise:46696e7365742e6d656d5f73696e676c65746f6e}{Prerequisite~P\ref{premise:46696e7365742e6d656d5f73696e676c65746f6e}}, \paperlink{premise:4966662e6d70}{Prerequisite~P\ref{premise:4966662e6d70}} (using \(\left|s\right| \le 1\) and \(y \in s\) and \(x \in s\)) and \paperlink{theorem:108}{Lemma~\ref{theorem:108}} gives \(y \in \{x\}\).
\end{proofstep}
This witness proves existence. Hence \(\exists x_{1}\, :\, \alpha,\; s \subseteq \{x_{1}\}\).
\end{proofstep}
Using the existential hypothesis, we obtain \(\exists x\, :\, \alpha,\; s \subseteq \{x\}\).
\end{proofstep}
The two alternatives exhaust the cases. Hence \(\exists x\, :\, \alpha,\; s \subseteq \{x\}\).
\end{proofstep}
Consider the following case.
\begin{proofstep}
Assume \(\exists x\, :\, \alpha,\; s \subseteq \{x\}\).
Use the hypothesis \(\exists x\, :\, \alpha,\; s \subseteq \{x\}\).
Consider the following case.
\begin{proofstep}
Fix \(x\) in \(\alpha\).
Assume \(s \subseteq \{x\}\).
By \paperlink{theorem:5}{Lemma~\ref{theorem:5}}, we obtain \(\left|s\right| \le \left|\{x\}\right|\).
Rewriting with \paperlink{theorem:14}{Lemma~\ref{theorem:14}} in the preceding result gives \(\left|s\right| \le 1\).
\end{proofstep}
Using the existential hypothesis, we obtain \(\left|s\right| \le 1\).
\end{proofstep}
Together the two implications give the equivalence. Hence \(\left|s\right| \le 1 \iff \left(\exists x\, :\, \alpha,\; s \subseteq \{x\}\right)\).
\end{proof}

\begin{theorem}\label{theorem:112}
Let \(\alpha\) be a type. Let \(s\) be a finite subset of \(\alpha\). Assume \(1 < \left|s\right|\). Let \(a\) be an element of \(\alpha\). Then \(\exists b\, :\, \alpha,\; b \in s \land b \ne a\).
\end{theorem}

\begin{proof}
Choose the witness \(a\).
This element witnesses nonemptiness. Hence \(\alpha\text{ is nonempty}\).
Consider the following case.
\begin{proofstep}
Assume \(\neg \left(\exists b\, :\, \alpha,\; b \in s \land b \ne a\right)\).
Use the hypothesis \(1 < \left|s\right|\).
By \paperlink{theorem:111}{Lemma~\ref{theorem:111}}, we obtain \(\left|s\right| \le 1 \iff \left(\exists x\, :\, \alpha,\; s \subseteq \{x\}\right)\).
Choose the witness \(a\).
By \paperlink{premise:46696e7365742e7375627365745f73696e676c65746f6e5f69666627}{Prerequisite~P\ref{premise:46696e7365742e7375627365745f73696e676c65746f6e5f69666627}}, we obtain \(s \subseteq \{a\} \iff \left(\forall b\, :\, \alpha,\; b \in s \implies b = a\right)\).
The substitution uses the following argument.
\begin{proofstep}
Fix \(b\) in \(\alpha\).
By \paperlink{premise:4d6174686c69622e5461637469632e507573684e65672e6e6f745f616e645f6571}{Prerequisite~P\ref{premise:4d6174686c69622e5461637469632e507573684e65672e6e6f745f616e645f6571}} with \(p = b \in s\) and \(q = b \ne a\), we obtain \(\left(\neg \left(b \in s \land b \ne a\right)\right) = \left(b \in s \implies \neg \left(b \ne a\right)\right)\).
By definition, we obtain \(\left(b \in s\right) = \left(b \in s\right)\).
By \paperlink{premise:4d6174686c69622e5461637469632e507573684e65672e6e6f745f6e655f6571}{Prerequisite~P\ref{premise:4d6174686c69622e5461637469632e507573684e65672e6e6f745f6e655f6571}} with \(x = b\) and \(y = a\), we obtain \(\left(\neg \left(b \ne a\right)\right) = \left(b = a\right)\).
By \paperlink{premise:696d706c6965735f636f6e6772}{Prerequisite~P\ref{premise:696d706c6965735f636f6e6772}} with \(p_{1} = b \in s\) and \(p_{2} = b \in s\) and \(q_{1} = \neg \left(b \ne a\right)\) and \(q_{2} = b = a\), we obtain \(\left(b \in s \implies \neg \left(b \ne a\right)\right) = \left(b \in s \implies b = a\right)\).
Chaining the equalities, we obtain \(\left(\neg \left(b \in s \land b \ne a\right)\right) = \left(b \in s \implies b = a\right)\).
\end{proofstep}
Use the hypothesis \(\neg \left(\exists b\, :\, \alpha,\; b \in s \land b \ne a\right)\).
Rewriting with \paperlink{premise:4d6174686c69622e5461637469632e507573684e65672e6e6f745f6578697374735f6571}{Prerequisite~P\ref{premise:4d6174686c69622e5461637469632e507573684e65672e6e6f745f6578697374735f6571}}, \paperlink{premise:666f72616c6c5f636f6e6772}{Prerequisite~P\ref{premise:666f72616c6c5f636f6e6772}} and the identity established above in the preceding result gives \(\forall b\, :\, \alpha,\; b \in s \implies b = a\).
By the reverse implication, we obtain \(s \subseteq \{a\}\).
This witness proves existence. Hence \(\exists x\, :\, \alpha,\; s \subseteq \{x\}\).
By the reverse implication, we obtain \(\left|s\right| \le 1\).
By \paperlink{premise:4c542e6c742e6e6f745f6c65}{\emph{Alias}} (\paperlink{premise:4c542e6c742e6e6f745f6c65}{Prerequisite~P\ref{premise:4c542e6c742e6e6f745f6c65}}) with \(\mathbb{N}\) as its type parameter \(\alpha\), we obtain \(\mathrm{False}\).
\end{proofstep}
The contrary assumption leads to a contradiction. Hence \(\exists b\, :\, \alpha,\; b \in s \land b \ne a\).
\end{proof}

\begin{theorem}\label{theorem:113}
Let \(\alpha\) be a type. Assume \(\alpha\text{ has at most one element}\). Let \(s\) be a finite subset of \(\alpha\). Then \(\left|s\right| \le 1\).
\end{theorem}

\begin{proof}
By \paperlink{theorem:109}{Lemma~\ref{theorem:109}}, we obtain \(|s| \le 1 \iff (\forall a\, :\, \alpha,\; \forall b\, :\, \alpha,\; a \in s \implies b \in s \implies a = b)\).
For the required implication, argue as follows.
\begin{proofstep}
Fix \(x\) in \(\alpha\).
Fix \(x_{1}\) in \(\alpha\).
Assume \(x \in s\).
Assume \(x_{1} \in s\).
Use the hypothesis \(\alpha\text{ has at most one element}\).
By \paperlink{premise:53756273696e676c65746f6e2e656c696d}{Prerequisite~P\ref{premise:53756273696e676c65746f6e2e656c696d}} with \(a = x\) and \(b = x_{1}\), we obtain \(x = x_{1}\).
\end{proofstep}
By the reverse implication, we obtain \(\left|s\right| \le 1\).
\end{proof}

\begin{lemma}\label{theorem:114}
Let \(\alpha\) be a type. Let \(s\) be a finite subset of \(\alpha\). Then
\begin{center}
\(\displaystyle 1 < |s| \iff (\exists a\, :\, \alpha,\; a \in s \land (\exists b\, :\, \alpha,\; b \in s \land a \ne b)).\)
\end{center}
\end{lemma}

\begin{proof}
The substitution uses the following argument.
\begin{proofstep}
Fix \(a\) in \(\alpha\).
By \paperlink{premise:4d6174686c69622e5461637469632e507573684e65672e6e6f745f616e645f6571}{Prerequisite~P\ref{premise:4d6174686c69622e5461637469632e507573684e65672e6e6f745f616e645f6571}} with \(p = a \in s\) and \(q = \exists b\, :\, \alpha,\; b \in s \land a \ne b\), we obtain \((\neg (a \in s \land (\exists b\, :\, \alpha,\; b \in s \land a \ne b))) = (a \in s \implies \neg (\exists b\, :\, \alpha,\; b \in s \land a \ne b))\).
By definition, we obtain \(\left(a \in s\right) = \left(a \in s\right)\).
By \paperlink{premise:4d6174686c69622e5461637469632e507573684e65672e6e6f745f6578697374735f6571}{Prerequisite~P\ref{premise:4d6174686c69622e5461637469632e507573684e65672e6e6f745f6578697374735f6571}} with \(b \mapsto b \in s \land a \ne b\) in place of its parameter \(s\), we obtain \((\neg (\exists x\, :\, \alpha,\; x \in s \land a \ne x)) = (\forall x\, :\, \alpha,\; \neg (x \in s \land a \ne x))\).
For the required implication, argue as follows.
\begin{proofstep}
Fix \(b\) in \(\alpha\).
By \paperlink{premise:4d6174686c69622e5461637469632e507573684e65672e6e6f745f616e645f6571}{Prerequisite~P\ref{premise:4d6174686c69622e5461637469632e507573684e65672e6e6f745f616e645f6571}} with \(p = b \in s\) and \(q = a \ne b\), we obtain \(\left(\neg \left(b \in s \land a \ne b\right)\right) = \left(b \in s \implies \neg \left(a \ne b\right)\right)\).
By definition, we obtain \(\left(b \in s\right) = \left(b \in s\right)\).
By \paperlink{premise:4d6174686c69622e5461637469632e507573684e65672e6e6f745f6e655f6571}{Prerequisite~P\ref{premise:4d6174686c69622e5461637469632e507573684e65672e6e6f745f6e655f6571}} with \(x = a\) and \(y = b\), we obtain \(\left(\neg \left(a \ne b\right)\right) = \left(a = b\right)\).
By \paperlink{premise:696d706c6965735f636f6e6772}{Prerequisite~P\ref{premise:696d706c6965735f636f6e6772}} with \(p_{1} = b \in s\) and \(p_{2} = b \in s\) and \(q_{1} = \neg \left(a \ne b\right)\) and \(q_{2} = a = b\), we obtain \(\left(b \in s \implies \neg \left(a \ne b\right)\right) = \left(b \in s \implies a = b\right)\).
Chaining the equalities, we obtain \(\left(\neg \left(b \in s \land a \ne b\right)\right) = \left(b \in s \implies a = b\right)\).
\end{proofstep}
By \paperlink{premise:666f72616c6c5f636f6e6772}{Prerequisite~P\ref{premise:666f72616c6c5f636f6e6772}}, we obtain \((\forall a_{1}\, :\, \alpha,\; \neg (a_{1} \in s \land a \ne a_{1})) = (\forall a_{1}\, :\, \alpha,\; a_{1} \in s \implies a = a_{1})\).
Chaining the equalities, we obtain \((\neg (\exists x\, :\, \alpha,\; x \in s \land a \ne x)) = (\forall a_{1}\, :\, \alpha,\; a_{1} \in s \implies a = a_{1})\).
By \paperlink{premise:696d706c6965735f636f6e6772}{Prerequisite~P\ref{premise:696d706c6965735f636f6e6772}} with \(p_{1} = a \in s\) and \(p_{2} = a \in s\) and \(q_{1} = \neg \left(\exists x\, :\, \alpha,\; x \in s \land a \ne x\right)\) and \(q_{2} = \forall a_{1}\, :\, \alpha,\; a_{1} \in s \implies a = a_{1}\), we obtain \((a \in s \implies \neg (\exists x\, :\, \alpha,\; x \in s \land a \ne x)) = (a \in s \implies (\forall a_{2}\, :\, \alpha,\; a_{2} \in s \implies a = a_{2}))\).
Chaining the equalities, we obtain \((\neg (a \in s \land (\exists b\, :\, \alpha,\; b \in s \land a \ne b))) = (a \in s \implies (\forall a_{2}\, :\, \alpha,\; a_{2} \in s \implies a = a_{2}))\).
\end{proofstep}
By \paperlink{theorem:108}{Lemma~\ref{theorem:108}}, we obtain \(|s| \le 1 \iff (\forall a\, :\, \alpha,\; a \in s \implies (\forall b\, :\, \alpha,\; b \in s \implies a = b))\).
Rewriting with \paperlink{premise:4d6174686c69622e5461637469632e507573684e65672e6e6f745f6c745f6571}{Prerequisite~P\ref{premise:4d6174686c69622e5461637469632e507573684e65672e6e6f745f6c745f6571}}, \paperlink{premise:4d6174686c69622e5461637469632e507573684e65672e6e6f745f6578697374735f6571}{Prerequisite~P\ref{premise:4d6174686c69622e5461637469632e507573684e65672e6e6f745f6578697374735f6571}}, \paperlink{premise:666f72616c6c5f636f6e6772}{Prerequisite~P\ref{premise:666f72616c6c5f636f6e6772}} and the identity established above in the preceding result gives \(\neg (1 < |s|) \iff \neg (\exists a\, :\, \alpha,\; a \in s \land (\exists b\, :\, \alpha,\; b \in s \land a \ne b))\).
Rewriting with \paperlink{premise:6e6f745f6966665f6e6f74}{Prerequisite~P\ref{premise:6e6f745f6966665f6e6f74}} in the preceding result gives \(1 < |s| \iff (\exists a\, :\, \alpha,\; a \in s \land (\exists b\, :\, \alpha,\; b \in s \land a \ne b))\).
\end{proof}

\begin{lemma}\label{theorem:115}
Let \(\alpha\) be a type. Let \(s\) be a finite subset of \(\alpha\). Then
\begin{center}
\(\displaystyle 1 < |s| \iff (\exists a\, :\, \alpha,\; \exists b\, :\, \alpha,\; a \in s \land (b \in s \land a \ne b)).\)
\end{center}
\end{lemma}

\begin{proof}
The substitution uses the following argument.
\begin{proofstep}
Fix \(a\) in \(\alpha\).
By \paperlink{premise:6578697374735f616e645f6c656674}{Prerequisite~P\ref{premise:6578697374735f616e645f6c656674}}, we obtain \((\exists x\, :\, \alpha,\; a \in s \land (x_{1} \mapsto x_{1} \in s \land a \ne x_{1})(x)) \iff a \in s \land (\exists x\, :\, \alpha,\; (x_{1} \mapsto x_{1} \in s \land a \ne x_{1})(x))\).
By propositional extensionality, we obtain \((\exists x\, :\, \alpha,\; a \in s \land (x \in s \land a \ne x)) = (a \in s \land (\exists x\, :\, \alpha,\; x \in s \land a \ne x))\).
\end{proofstep}
Rewriting with \paperlink{premise:66756e657874}{\emph{Function extensionality}} (\paperlink{premise:66756e657874}{Prerequisite~P\ref{premise:66756e657874}}), \paperlink{premise:6966665f73656c66}{Prerequisite~P\ref{premise:6966665f73656c66}} and the identity established above gives \((\exists a\, :\, \alpha,\; a \in s \land (\exists b\, :\, \alpha,\; b \in s \land a \ne b)) \iff (\exists a\, :\, \alpha,\; \exists b\, :\, \alpha,\; a \in s \land (b \in s \land a \ne b))\).
Rewriting with \paperlink{theorem:114}{Lemma~\ref{theorem:114}} in the preceding result gives \(1 < |s| \iff (\exists a\, :\, \alpha,\; \exists b\, :\, \alpha,\; a \in s \land (b \in s \land a \ne b))\).
\end{proof}

\begin{theorem}\label{theorem:116}
Let \(\alpha\) be a type. Let \(s\) be a finite subset of \(\alpha\). Then \(1 < \left|s\right| \iff s\text{ contains two distinct elements}\).
\end{theorem}

\begin{proof}
Rewriting with \paperlink{premise:46696e7365742e636f655f736f72745f636f65}{Prerequisite~P\ref{premise:46696e7365742e636f655f736f72745f636f65}}, \paperlink{theorem:110}{Lemma~\ref{theorem:110}}, \paperlink{premise:6e6f745f6e6f6e7472697669616c5f6966665f73756273696e676c65746f6e}{Prerequisite~P\ref{premise:6e6f745f6e6f6e7472697669616c5f6966665f73756273696e676c65746f6e}}, \paperlink{premise:5365742e6e6f6e7472697669616c5f636f655f736f7274}{Prerequisite~P\ref{premise:5365742e6e6f6e7472697669616c5f636f655f736f7274}}, \paperlink{premise:6e6f745f6c74}{Prerequisite~P\ref{premise:6e6f745f6c74}} and \paperlink{premise:6e6f745f6966665f6e6f74}{Prerequisite~P\ref{premise:6e6f745f6966665f6e6f74}} gives \(1 < \left|s\right| \iff s\text{ contains two distinct elements}\).
\end{proof}

\begin{theorem}\label{theorem:117}
Let \(\alpha\) be a type. Let \(s\) be a finite subset of \(\alpha\). Assume \(1 < \left|s\right|\). Let \(a\) be an element of \(\alpha\). Then \(\exists b\, :\, \alpha,\; b \in s \land b \ne a\).
\end{theorem}

\begin{proof}
By \paperlink{theorem:114}{Lemma~\ref{theorem:114}}, we obtain \(1 < |s| \iff (\exists a_{1}\, :\, \alpha,\; a_{1} \in s \land (\exists b\, :\, \alpha,\; b \in s \land a_{1} \ne b))\).
Use the hypothesis \(1 < \left|s\right|\).
By the forward implication, we obtain \(\exists a_{1}\, :\, \alpha,\; a_{1} \in s \land (\exists b\, :\, \alpha,\; b \in s \land a_{1} \ne b)\).
Consider the following case.
\begin{proofstep}
Fix \(x\) in \(\alpha\).
Assume \(x \in s \land \left(\exists b\, :\, \alpha,\; b \in s \land x \ne b\right)\).
Use the hypothesis \(x \in s \land \left(\exists b\, :\, \alpha,\; b \in s \land x \ne b\right)\).
Consider the following case.
\begin{proofstep}
Assume \(x \in s\).
Assume \(\exists b\, :\, \alpha,\; b \in s \land x \ne b\).
Use the hypothesis \(\exists b\, :\, \alpha,\; b \in s \land x \ne b\).
Consider the following case.
\begin{proofstep}
Fix \(y\) in \(\alpha\).
Assume \(y \in s \land x \ne y\).
Use the hypothesis \(y \in s \land x \ne y\).
Consider the following case.
\begin{proofstep}
Assume \(y \in s\).
Assume \(x \ne y\).
Distinguish whether \(y = a\) holds.
Consider the following case.
\begin{proofstep}
Assume \(y = a\).
Choose the witness \(x\).
Use the hypothesis \(x \in s\).
Use the hypothesis \(x \ne y\).
Use the hypothesis \(y = a\).
By \paperlink{premise:6e655f6f665f6e655f6f665f6571}{Prerequisite~P\ref{premise:6e655f6f665f6e655f6f665f6571}}, we obtain \(x \ne a\).
Combining the two conclusions, we obtain \(x \in s \land x \ne a\).
This witness proves existence. Hence \(\exists b\, :\, \alpha,\; b \in s \land b \ne a\).
\end{proofstep}
Consider the following case.
\begin{proofstep}
Assume \(\neg \left(y = a\right)\).
Choose the witness \(y\).
Use the hypothesis \(y \in s\).
Use the hypothesis \(\neg \left(y = a\right)\).
Combining the two conclusions, we obtain \(y \in s \land y \ne a\).
This witness proves existence. Hence \(\exists b\, :\, \alpha,\; b \in s \land b \ne a\).
\end{proofstep}
The proposition or its negation holds. Hence \(\exists b\, :\, \alpha,\; b \in s \land b \ne a\).
\end{proofstep}
Using the two conjuncts, we obtain \(\exists b\, :\, \alpha,\; b \in s \land b \ne a\).
\end{proofstep}
Using the existential hypothesis, we obtain \(\exists b\, :\, \alpha,\; b \in s \land b \ne a\).
\end{proofstep}
Using the two conjuncts, we obtain \(\exists b\, :\, \alpha,\; b \in s \land b \ne a\).
\end{proofstep}
Using the existential hypothesis, we obtain \(\exists b\, :\, \alpha,\; b \in s \land b \ne a\).
\end{proof}

\begin{theorem}\label{theorem:118}
Let \(\iota\) be a type. Let \(\alpha\) be a function of type \(\iota \to \mathsf{Type}\). Fix a structure \(\mathit{inst}\) of type \(\Pi_{i\, :\, \iota}\, \operatorname{Decidable\,Eq}\left(\alpha\left(i\right)\right)\). Let \(s\) be a finite subset of \(\Pi_{i\, :\, \iota}\, \alpha\left(i\right)\). Assume \(1 < \left|s\right|\). Then
\begin{center}
\(\displaystyle \exists i\, :\, \iota,\; 1 < |(x \mapsto x(i))[s]| \land (\forall \mathit{ai}\, :\, \alpha(i),\; \{x\in s \mid x(i) = \mathit{ai}\} \subsetneq s).\)
\end{center}
\end{theorem}

\begin{proof}
The substitution uses the following argument.
\begin{proofstep}
Fix \(i\) in \(\iota\).
By \paperlink{theorem:115}{Lemma~\ref{theorem:115}} with \(\alpha\left(i\right)\) as its type parameter \(\alpha\), we obtain \(1 < |(x \mapsto x(i))[s]| \iff (\exists a\, :\, \alpha(i),\; \exists b\, :\, \alpha(i),\; a \in (x \mapsto x(i))[s] \land (b \in (x \mapsto x(i))[s] \land a \ne b))\).
By propositional extensionality, we obtain \((1 < |(x \mapsto x(i))[s]|) = (\exists a\, :\, \alpha(i),\; \exists b\, :\, \alpha(i),\; a \in (x \mapsto x(i))[s] \land (b \in (x \mapsto x(i))[s] \land a \ne b))\).
Applying the same function to equal arguments, we obtain \((1 < |(x \mapsto x(i))[s]| \land (\forall \mathit{ai}\, :\, \alpha(i),\; \{x\in s \mid x(i) = \mathit{ai}\} \subsetneq s)) = ((\exists a\, :\, \alpha(i),\; \exists b\, :\, \alpha(i),\; a \in (x \mapsto x(i))[s] \land (b \in (x \mapsto x(i))[s] \land a \ne b)) \land (\forall \mathit{ai}\, :\, \alpha(i),\; \{x\in s \mid x(i) = \mathit{ai}\} \subsetneq s))\).
\end{proofstep}
The substitution uses the following argument.
\begin{proofstep}
Fix \(a\) in \(\Pi_{i\, :\, \iota}\, \alpha\left(i\right)\).
For the required implication, argue as follows.
\begin{proofstep}
Fix \(b\) in \(\Pi_{i\, :\, \iota}\, \alpha\left(i\right)\).
By \paperlink{premise:46756e6374696f6e2e6e655f696666}{Prerequisite~P\ref{premise:46756e6374696f6e2e6e655f696666}}, we obtain \(a \ne b \iff \left(\exists a_{1}\, :\, \iota,\; a\left(a_{1}\right) \ne b\left(a_{1}\right)\right)\).
By propositional extensionality, we obtain \(\left(a \ne b\right) = \left(\exists a_{1}\, :\, \iota,\; a\left(a_{1}\right) \ne b\left(a_{1}\right)\right)\).
Applying the same function to equal arguments, we obtain \((a \in s \land (b \in s \land a \ne b)) = (a \in s \land (b \in s \land (\exists a_{1}\, :\, \iota,\; a(a_{1}) \ne b(a_{1}))))\).
\end{proofstep}
By function extensionality, we obtain \((x \mapsto a \in s \land (x \in s \land a \ne x)) = (x \mapsto a \in s \land (x \in s \land (\exists a_{1}\, :\, \iota,\; a(a_{1}) \ne x(a_{1}))))\).
Applying the same function to equal arguments, we obtain \((\exists x\, :\, \Pi_{i\, :\, \iota}\, \alpha(i),\; a \in s \land (x \in s \land a \ne x)) = (\exists x\, :\, \Pi_{i\, :\, \iota}\, \alpha(i),\; a \in s \land (x \in s \land (\exists a_{1}\, :\, \iota,\; a(a_{1}) \ne x(a_{1}))))\).
\end{proofstep}
Use the hypothesis \(1 < \left|s\right|\).
Rewriting with \paperlink{theorem:115}{Lemma~\ref{theorem:115}} in the preceding result gives \(\exists a\, :\, \Pi_{i\, :\, \iota}\, \alpha(i),\; \exists b\, :\, \Pi_{i\, :\, \iota}\, \alpha(i),\; a \in s \land (b \in s \land a \ne b)\).
Rewriting with \paperlink{premise:66756e657874}{\emph{Function extensionality}} (\paperlink{premise:66756e657874}{Prerequisite~P\ref{premise:66756e657874}}) and the identity established above in the preceding result gives \(\exists a\, :\, \Pi_{i\, :\, \iota}\, \alpha(i),\; \exists b\, :\, \Pi_{i\, :\, \iota}\, \alpha(i),\; a \in s \land (b \in s \land (\exists a_{1}\, :\, \iota,\; a(a_{1}) \ne b(a_{1})))\).
Consider the following case.
\begin{proofstep}
Fix \(\mathit{a1}\) in \(\Pi_{i\, :\, \iota}\, \alpha\left(i\right)\).
Assume \(\exists b\, :\, \Pi_{i\, :\, \iota}\, \alpha(i),\; \mathit{a1} \in s \land (b \in s \land (\exists a\, :\, \iota,\; \mathit{a1}(a) \ne b(a)))\).
Use the hypothesis \(\exists b\, :\, \Pi_{i\, :\, \iota}\, \alpha(i),\; \mathit{a1} \in s \land (b \in s \land (\exists a\, :\, \iota,\; \mathit{a1}(a) \ne b(a)))\).
Consider the following case.
\begin{proofstep}
Fix \(\mathit{a2}\) in \(\Pi_{i\, :\, \iota}\, \alpha\left(i\right)\).
Assume \(\mathit{a1} \in s \land (\mathit{a2} \in s \land (\exists a\, :\, \iota,\; \mathit{a1}(a) \ne \mathit{a2}(a)))\).
Use the hypothesis \(\mathit{a1} \in s \land (\mathit{a2} \in s \land (\exists a\, :\, \iota,\; \mathit{a1}(a) \ne \mathit{a2}(a)))\).
Consider the following case.
\begin{proofstep}
Assume \(\mathit{a1} \in s\).
Assume \(\mathit{a2} \in s \land (\exists a\, :\, \iota,\; \mathit{a1}(a) \ne \mathit{a2}(a))\).
Use the hypothesis \(\mathit{a2} \in s \land (\exists a\, :\, \iota,\; \mathit{a1}(a) \ne \mathit{a2}(a))\).
Consider the following case.
\begin{proofstep}
Assume \(\mathit{a2} \in s\).
Assume \(\exists a\, :\, \iota,\; \mathit{a1}\left(a\right) \ne \mathit{a2}\left(a\right)\).
Use the hypothesis \(\exists a\, :\, \iota,\; \mathit{a1}\left(a\right) \ne \mathit{a2}\left(a\right)\).
Consider the following case.
\begin{proofstep}
Fix \(i\) in \(\iota\).
Assume \(\mathit{a1}\left(i\right) \ne \mathit{a2}\left(i\right)\).
Choose the witness \(i\).
Choose the witness \(\mathit{a1}\left(i\right)\).
Choose the witness \(\mathit{a2}\left(i\right)\).
Use the hypothesis \(\mathit{a1} \in s\).
By \paperlink{premise:46696e7365742e6d656d5f696d6167655f6f665f6d656d}{Prerequisite~P\ref{premise:46696e7365742e6d656d5f696d6167655f6f665f6d656d}} with \(\Pi_{i_{1}\, :\, \iota}\, \alpha\left(i_{1}\right)\) as its type parameter \(\alpha\) and \(\beta = \alpha\left(i\right)\) and \(f = x \mapsto x\left(i\right)\), we obtain \(\mathit{a1}\left(i\right) \in \left(x \mapsto x\left(i\right)\right)\left[s\right]\).
Use the hypothesis \(\mathit{a2} \in s\).
By \paperlink{premise:46696e7365742e6d656d5f696d6167655f6f665f6d656d}{Prerequisite~P\ref{premise:46696e7365742e6d656d5f696d6167655f6f665f6d656d}} with \(\Pi_{i_{1}\, :\, \iota}\, \alpha\left(i_{1}\right)\) as its type parameter \(\alpha\) and \(\beta = \alpha\left(i\right)\) and \(f = x \mapsto x\left(i\right)\), we obtain \(\mathit{a2}\left(i\right) \in \left(x \mapsto x\left(i\right)\right)\left[s\right]\).
Use the hypothesis \(\mathit{a1}\left(i\right) \ne \mathit{a2}\left(i\right)\).
Combining the two conclusions, we obtain \(\mathit{a2}(i) \in (x \mapsto x(i))[s] \land \mathit{a1}(i) \ne \mathit{a2}(i)\).
Combining the two conclusions, we obtain \(\mathit{a1}(i) \in (x \mapsto x(i))[s] \land (\mathit{a2}(i) \in (x \mapsto x(i))[s] \land \mathit{a1}(i) \ne \mathit{a2}(i))\).
This witness proves existence. Hence \(\exists b\, :\, \alpha(i),\; \mathit{a1}(i) \in (x \mapsto x(i))[s] \land (b \in (x \mapsto x(i))[s] \land \mathit{a1}(i) \ne b)\).
This witness proves existence. Hence \(\exists a\, :\, \alpha(i),\; \exists b\, :\, \alpha(i),\; a \in (x \mapsto x(i))[s] \land (b \in (x \mapsto x(i))[s] \land a \ne b)\).
For the required implication, argue as follows.
\begin{proofstep}
Fix \(\mathit{ai}\) in \(\alpha\left(i\right)\).
By \paperlink{premise:65715f6f725f6e65}{Prerequisite~P\ref{premise:65715f6f725f6e65}} with \(\alpha\left(i\right)\) as its type parameter \(\alpha\) and \(x = \mathit{a2}\left(i\right)\) and \(y = \mathit{ai}\), we obtain \(\mathit{a2}\left(i\right) = \mathit{ai} \lor \mathit{a2}\left(i\right) \ne \mathit{ai}\).
Consider the following case.
\begin{proofstep}
Assume \(\mathit{a2}\left(i\right) = \mathit{ai}\).
Choose the witness \(\mathit{a1}\).
Use the hypothesis \(\mathit{a1} \in s\).
Use the hypothesis \(\mathit{a1}\left(i\right) \ne \mathit{a2}\left(i\right)\).
Combining the two conclusions, we obtain \(\mathit{a1} \in s \land \neg \left(\mathit{a1}\left(i\right) = \mathit{a2}\left(i\right)\right)\).
This witness proves existence. Hence \(\exists x\, :\, \Pi_{i_{1}\, :\, \iota}\, \alpha(i_{1}),\; x \in s \land \neg (x(i) = \mathit{a2}(i))\).
Use the hypothesis \(\mathit{a2}\left(i\right) = \mathit{ai}\).
Substituting equals, we obtain \(\exists x\, :\, \Pi_{i_{1}\, :\, \iota}\, \alpha(i_{1}),\; x \in s \land \neg (x(i) = \mathit{ai})\).
\end{proofstep}
Consider the following case.
\begin{proofstep}
Assume \(\mathit{a2}\left(i\right) \ne \mathit{ai}\).
Choose the witness \(\mathit{a2}\).
Use the hypothesis \(\mathit{a2} \in s\).
Use the hypothesis \(\mathit{a2}\left(i\right) \ne \mathit{ai}\).
Combining the two conclusions, we obtain \(\mathit{a2} \in s \land \neg \left(\mathit{a2}\left(i\right) = \mathit{ai}\right)\).
This witness proves existence. Hence \(\exists x\, :\, \Pi_{i_{1}\, :\, \iota}\, \alpha(i_{1}),\; x \in s \land \neg (x(i) = \mathit{ai})\).
\end{proofstep}
The two alternatives exhaust the cases. Hence \(\exists x\, :\, \Pi_{i_{1}\, :\, \iota}\, \alpha(i_{1}),\; x \in s \land \neg (x(i) = \mathit{ai})\).
Rewriting with \paperlink{premise:46696e7365742e66696c7465725f73737562736574}{Prerequisite~P\ref{premise:46696e7365742e66696c7465725f73737562736574}} in the preceding result gives \(\{x\in s \mid x\left(i\right) = \mathit{ai}\} \subsetneq s\).
\end{proofstep}
Combining the two conclusions, we obtain \((\exists a\, :\, \alpha(i),\; \exists b\, :\, \alpha(i),\; a \in (x \mapsto x(i))[s] \land (b \in (x \mapsto x(i))[s] \land a \ne b)) \land (\forall \mathit{ai}\, :\, \alpha(i),\; \{x\in s \mid x(i) = \mathit{ai}\} \subsetneq s)\).
This witness proves existence. Hence \(\exists i_{1}\, :\, \iota,\; (\exists a\, :\, \alpha(i_{1}),\; \exists b\, :\, \alpha(i_{1}),\; a \in (x \mapsto x(i_{1}))[s] \land (b \in (x \mapsto x(i_{1}))[s] \land a \ne b)) \land (\forall \mathit{ai}\, :\, \alpha(i_{1}),\; \{x\in s \mid x(i_{1}) = \mathit{ai}\} \subsetneq s)\).
\end{proofstep}
Using the existential hypothesis, we obtain \(\exists i\, :\, \iota,\; (\exists a\, :\, \alpha(i),\; \exists b\, :\, \alpha(i),\; a \in (x \mapsto x(i))[s] \land (b \in (x \mapsto x(i))[s] \land a \ne b)) \land (\forall \mathit{ai}\, :\, \alpha(i),\; \{x\in s \mid x(i) = \mathit{ai}\} \subsetneq s)\).
\end{proofstep}
Using the two conjuncts, we obtain \(\exists i\, :\, \iota,\; (\exists a\, :\, \alpha(i),\; \exists b\, :\, \alpha(i),\; a \in (x \mapsto x(i))[s] \land (b \in (x \mapsto x(i))[s] \land a \ne b)) \land (\forall \mathit{ai}\, :\, \alpha(i),\; \{x\in s \mid x(i) = \mathit{ai}\} \subsetneq s)\).
\end{proofstep}
Using the two conjuncts, we obtain \(\exists i\, :\, \iota,\; (\exists a\, :\, \alpha(i),\; \exists b\, :\, \alpha(i),\; a \in (x \mapsto x(i))[s] \land (b \in (x \mapsto x(i))[s] \land a \ne b)) \land (\forall \mathit{ai}\, :\, \alpha(i),\; \{x\in s \mid x(i) = \mathit{ai}\} \subsetneq s)\).
\end{proofstep}
Using the existential hypothesis, we obtain \(\exists i\, :\, \iota,\; (\exists a\, :\, \alpha(i),\; \exists b\, :\, \alpha(i),\; a \in (x \mapsto x(i))[s] \land (b \in (x \mapsto x(i))[s] \land a \ne b)) \land (\forall \mathit{ai}\, :\, \alpha(i),\; \{x\in s \mid x(i) = \mathit{ai}\} \subsetneq s)\).
\end{proofstep}
Using the existential hypothesis, we obtain \(\exists i\, :\, \iota,\; (\exists a\, :\, \alpha(i),\; \exists b\, :\, \alpha(i),\; a \in (x \mapsto x(i))[s] \land (b \in (x \mapsto x(i))[s] \land a \ne b)) \land (\forall \mathit{ai}\, :\, \alpha(i),\; \{x\in s \mid x(i) = \mathit{ai}\} \subsetneq s)\).
Rewriting with \paperlink{premise:66756e657874}{\emph{Function extensionality}} (\paperlink{premise:66756e657874}{Prerequisite~P\ref{premise:66756e657874}}) and the identity established above in the preceding result gives \(\exists i\, :\, \iota,\; 1 < |(x \mapsto x(i))[s]| \land (\forall \mathit{ai}\, :\, \alpha(i),\; \{x\in s \mid x(i) = \mathit{ai}\} \subsetneq s)\).
\end{proof}

\begin{theorem}\label{theorem:119}
Let \(\alpha\) be a type. Let \(s\) be a finite subset of \(\alpha\). Let \(n\) be a natural number. Then
\begin{center}
\(\displaystyle |s| = n + 1 \iff (\exists a\, :\, \alpha,\; \exists t\, :\, \mathcal{P}_{\mathrm{fin}}(\alpha),\; \exists h\text{ proving }(a \notin t),\; \operatorname{Finset\mathord{:}cons}(a,\allowbreak{} t) = s \land |t| = n).\)
\end{center}
\end{theorem}

\begin{proof}
The substitution uses the following argument.
\begin{proofstep}
Fix \(a\) in \(\alpha\).
For the required implication, argue as follows.
\begin{proofstep}
Fix \(t\) in \(\mathcal{P}_{\mathrm{fin}}(\alpha)\).
For the required implication, argue as follows.
\begin{proofstep}
Assume \(a \notin t\).
Use the hypothesis \(a \notin t\).
By \paperlink{premise:46696e7365742e636f6e735f6e655f656d707479}{Prerequisite~P\ref{premise:46696e7365742e636f6e735f6e655f656d707479}}, we obtain \(\operatorname{Finset\mathord{:}cons}\left(a, t\right) \ne \varnothing\).
By \paperlink{premise:65715f66616c7365}{Prerequisite~P\ref{premise:65715f66616c7365}}, we obtain \(\left(\operatorname{Finset\mathord{:}cons}\left(a, t\right) = \varnothing\right) = \mathrm{False}\).
Applying the same function to equal arguments, we obtain \((\operatorname{Finset\mathord{:}cons}(a,\allowbreak{} t) = \varnothing \land |t| = n) = (\mathrm{False} \land |t| = n)\).
By \paperlink{premise:66616c73655f616e64}{Prerequisite~P\ref{premise:66616c73655f616e64}} with \(p = \left|t\right| = n\), we obtain \(\left(\mathrm{False} \land \left|t\right| = n\right) = \mathrm{False}\).
Chaining the equalities, we obtain \((\operatorname{Finset\mathord{:}cons}(a,\allowbreak{} t) = \varnothing \land |t| = n) = \mathrm{False}\).
By \paperlink{premise:4966662e6f665f6571}{Prerequisite~P\ref{premise:4966662e6f665f6571}}, we obtain \(\operatorname{Finset\mathord{:}cons}(a,\allowbreak{} t) = \varnothing \land |t| = n \iff \mathrm{False}\).
\end{proofstep}
By definition, we obtain \(\left(a \notin t\right) = \left(a \notin t\right)\).
By \paperlink{premise:4966662e6f665f6571}{Prerequisite~P\ref{premise:4966662e6f665f6571}}, we obtain \(a \notin t \iff a \notin t\).
By \paperlink{premise:6578697374735f70726f705f636f6e6772}{Prerequisite~P\ref{premise:6578697374735f70726f705f636f6e6772}}, we obtain \((\exists h\text{ proving }(a \notin t),\; \operatorname{Finset\mathord{:}cons}(a,\allowbreak{} t) = \varnothing \land |t| = n) \iff a \notin t \land \mathrm{False}\).
By propositional extensionality, we obtain \((\exists h\text{ proving }(a \notin t),\; \operatorname{Finset\mathord{:}cons}(a,\allowbreak{} t) = \varnothing \land |t| = n) = (a \notin t \land \mathrm{False})\).
By \paperlink{premise:6578697374735f66616c7365}{Prerequisite~P\ref{premise:6578697374735f66616c7365}} with \(a \notin t\) as its type parameter \(\alpha\), we obtain \(\neg \left(\exists x\, :\, a \notin t,\; \mathrm{False}\right)\).
By \paperlink{premise:65715f66616c7365}{Prerequisite~P\ref{premise:65715f66616c7365}}, we obtain \(\left(a \notin t \land \mathrm{False}\right) = \mathrm{False}\).
Chaining the equalities, we obtain \((\exists h\text{ proving }(a \notin t),\; \operatorname{Finset\mathord{:}cons}(a,\allowbreak{} t) = \varnothing \land |t| = n) = \mathrm{False}\).
\end{proofstep}
By function extensionality, we obtain \((x \mapsto \exists h\text{ proving }(a \notin x),\; \operatorname{Finset\mathord{:}cons}(a,\allowbreak{} x) = \varnothing \land |x| = n) = (x \mapsto \mathrm{False})\).
Applying the same function to equal arguments, we obtain \((\exists x\, :\, \mathcal{P}_{\mathrm{fin}}(\alpha),\; \exists h\text{ proving }(a \notin x),\; \operatorname{Finset\mathord{:}cons}(a,\allowbreak{} x) = \varnothing \land |x| = n) = (\exists x\, :\, \mathcal{P}_{\mathrm{fin}}(\alpha),\; \mathrm{False})\).
By \paperlink{premise:696e73744e6f6e656d7074794f66496e68616269746564}{Prerequisite~P\ref{premise:696e73744e6f6e656d7074794f66496e68616269746564}} with \(\mathcal{P}_{\mathrm{fin}}(\alpha)\) as its type parameter \(\alpha\), we obtain \(\mathcal{P}_{\mathrm{fin}}(\alpha)\text{ is nonempty}\).
By \paperlink{premise:6578697374735f636f6e7374}{Prerequisite~P\ref{premise:6578697374735f636f6e7374}} with \(\mathcal{P}_{\mathrm{fin}}(\alpha)\) as its type parameter \(\alpha\), we obtain \(\left(\exists x\, :\, \mathcal{P}_{\mathrm{fin}}(\alpha),\; \mathrm{False}\right) \iff \mathrm{False}\).
By propositional extensionality, we obtain \(\left(\exists x\, :\, \mathcal{P}_{\mathrm{fin}}(\alpha),\; \mathrm{False}\right) = \mathrm{False}\).
Chaining the equalities, we obtain \((\exists x\, :\, \mathcal{P}_{\mathrm{fin}}(\alpha),\; \exists h\text{ proving }(a \notin x),\; \operatorname{Finset\mathord{:}cons}(a,\allowbreak{} x) = \varnothing \land |x| = n) = \mathrm{False}\).
\end{proofstep}
Rewriting with \paperlink{premise:696d706c6965735f636f6e6772}{Prerequisite~P\ref{premise:696d706c6965735f636f6e6772}}, \paperlink{premise:4e61742e73656c665f65715f6164645f6c656674}{Prerequisite~P\ref{premise:4e61742e73656c665f65715f6164645f6c656674}}, \paperlink{premise:4e61742e6164645f65715f7a65726f}{Prerequisite~P\ref{premise:4e61742e6164645f65715f7a65726f}}, \paperlink{premise:65715f66616c7365}{Prerequisite~P\ref{premise:65715f66616c7365}}, \paperlink{premise:6f6e655f6e655f7a65726f}{Prerequisite~P\ref{premise:6f6e655f6e655f7a65726f}}, \paperlink{premise:4e61742e696e73744e655a65726f53756363}{Prerequisite~P\ref{premise:4e61742e696e73744e655a65726f53756363}}, \paperlink{premise:616e645f66616c7365}{Prerequisite~P\ref{premise:616e645f66616c7365}}, \paperlink{premise:66756e657874}{\emph{Function extensionality}} (\paperlink{premise:66756e657874}{Prerequisite~P\ref{premise:66756e657874}}), \paperlink{premise:6578697374735f66616c7365}{Prerequisite~P\ref{premise:6578697374735f66616c7365}}, \paperlink{premise:696d705f73656c66}{Prerequisite~P\ref{premise:696d705f73656c66}} and the identity established above gives \(0 = n + 1 + 0 \implies (\exists x\, :\, \alpha,\; \exists x_{1}\, :\, \mathcal{P}_{\mathrm{fin}}(\alpha),\; \exists h\text{ proving }(x \notin x_{1}),\; \operatorname{Finset\mathord{:}cons}(x,\allowbreak{} x_{1}) = \varnothing \land |x_{1}| = n)\).
For the required implication, argue as follows.
\begin{proofstep}
Fix \(a\) in \(\alpha\).
Fix \(s_{1}\) in \(\mathcal{P}_{\mathrm{fin}}(\alpha)\).
Assume \(a \notin s_{1}\).
Assume \(|s_{1}| = n + 1 \implies (\exists a_{2}\, :\, \alpha,\; \exists t\, :\, \mathcal{P}_{\mathrm{fin}}(\alpha),\; \exists h\text{ proving }(a_{2} \notin t),\; \operatorname{Finset\mathord{:}cons}(a_{2},\allowbreak{} t) = s_{1} \land |t| = n)\).
Assume \(\left|\operatorname{Finset\mathord{:}cons}\left(a, s_{1}\right)\right| = n + 1\).
Choose the witness \(a\).
Choose the witness \(s_{1}\).
The substitution uses the following argument.
\begin{proofstep}
Assume \(a \notin s_{1}\).
By \paperlink{premise:65715f73656c66}{Prerequisite~P\ref{premise:65715f73656c66}} with \(\mathcal{P}_{\mathrm{fin}}(\alpha)\) as its type parameter \(\alpha\) and \(\operatorname{Finset\mathord{:}cons}\left(a, s_{1}\right)\) in place of its parameter \(a\), we obtain \((\operatorname{Finset\mathord{:}cons}(a,\allowbreak{} s_{1}) = \operatorname{Finset\mathord{:}cons}(a,\allowbreak{} s_{1})) = \mathrm{True}\).
Applying the same function to equal arguments, we obtain \((y \mapsto \operatorname{Finset\mathord{:}cons}(a,\allowbreak{} s_{1}) = \operatorname{Finset\mathord{:}cons}(a,\allowbreak{} s_{1}) \land y) = (y \mapsto \mathrm{True} \land y)\).
Use the hypothesis \(\left|\operatorname{Finset\mathord{:}cons}\left(a, s_{1}\right)\right| = n + 1\).
Rewriting with \paperlink{theorem:16}{Lemma~\ref{theorem:16}} (using \(a \notin s_{1}\)) and \paperlink{premise:4e61742e6164645f6c6566745f696e6a}{\emph{Alias}} (\paperlink{premise:4e61742e6164645f6c6566745f696e6a}{Prerequisite~P\ref{premise:4e61742e6164645f6c6566745f696e6a}}) in the preceding result gives \(\left|s_{1}\right| = n\).
Applying the same function to equal arguments, we obtain \(\left(\left|s_{1}\right| = n\right) = \left(n = n\right)\).
By \paperlink{premise:65715f73656c66}{Prerequisite~P\ref{premise:65715f73656c66}} with \(\mathbb{N}\) as its type parameter \(\alpha\) and \(n\) in place of its parameter \(a\), we obtain \(\left(n = n\right) = \mathrm{True}\).
Chaining the equalities, we obtain \(\left(\left|s_{1}\right| = n\right) = \mathrm{True}\).
By congruence, we obtain \((\operatorname{Finset\mathord{:}cons}(a,\allowbreak{} s_{1}) = \operatorname{Finset\mathord{:}cons}(a,\allowbreak{} s_{1}) \land |s_{1}| = n) = (\mathrm{True} \land \mathrm{True})\).
By \paperlink{premise:616e645f73656c66}{Prerequisite~P\ref{premise:616e645f73656c66}} with \(p = \mathrm{True}\), we obtain \(\left(\mathrm{True} \land \mathrm{True}\right) = \mathrm{True}\).
Chaining the equalities, we obtain \((\operatorname{Finset\mathord{:}cons}(a,\allowbreak{} s_{1}) = \operatorname{Finset\mathord{:}cons}(a,\allowbreak{} s_{1}) \land |s_{1}| = n) = \mathrm{True}\).
By \paperlink{premise:4966662e6f665f6571}{Prerequisite~P\ref{premise:4966662e6f665f6571}}, we obtain \(\operatorname{Finset\mathord{:}cons}(a,\allowbreak{} s_{1}) = \operatorname{Finset\mathord{:}cons}(a,\allowbreak{} s_{1}) \land |s_{1}| = n \iff \mathrm{True}\).
\end{proofstep}
Rewriting with \paperlink{premise:6578697374735f70726f705f636f6e6772}{Prerequisite~P\ref{premise:6578697374735f70726f705f636f6e6772}}, \paperlink{premise:4966662e6f665f6571}{Prerequisite~P\ref{premise:4966662e6f665f6571}}, \paperlink{premise:65715f66616c7365}{Prerequisite~P\ref{premise:65715f66616c7365}} (using \(a \notin s_{1}\)), \paperlink{premise:6e6f745f66616c73655f65715f74727565}{Prerequisite~P\ref{premise:6e6f745f66616c73655f65715f74727565}}, \paperlink{premise:6578697374735f636f6e7374}{Prerequisite~P\ref{premise:6578697374735f636f6e7374}}, \paperlink{premise:696e73744e6f6e656d7074794f66496e68616269746564}{Prerequisite~P\ref{premise:696e73744e6f6e656d7074794f66496e68616269746564}} and the auxiliary result established above gives \(\exists h\text{ proving }(a \notin s_{1}),\; \operatorname{Finset\mathord{:}cons}(a,\allowbreak{} s_{1}) = \operatorname{Finset\mathord{:}cons}(a,\allowbreak{} s_{1}) \land |s_{1}| = n\).
This witness proves existence. Hence \(\exists t\, :\, \mathcal{P}_{\mathrm{fin}}(\alpha),\; \exists h\text{ proving }(a \notin t),\; \operatorname{Finset\mathord{:}cons}(a,\allowbreak{} t) = \operatorname{Finset\mathord{:}cons}(a,\allowbreak{} s_{1}) \land |t| = n\).
This witness proves existence. Hence \(\exists a_{1}\, :\, \alpha,\; \exists t\, :\, \mathcal{P}_{\mathrm{fin}}(\alpha),\; \exists h\text{ proving }(a_{1} \notin t),\; \operatorname{Finset\mathord{:}cons}(a_{1},\allowbreak{} t) = \operatorname{Finset\mathord{:}cons}(a,\allowbreak{} s_{1}) \land |t| = n\).
\end{proofstep}
By \paperlink{premise:46696e7365742e636f6e735f696e64756374696f6e5f6f6e}{Prerequisite~P\ref{premise:46696e7365742e636f6e735f696e64756374696f6e5f6f6e}}, we obtain \(|s| = n + 1 \implies (\exists a_{1}\, :\, \alpha,\; \exists t\, :\, \mathcal{P}_{\mathrm{fin}}(\alpha),\; \exists h\text{ proving }(a_{1} \notin t),\; \operatorname{Finset\mathord{:}cons}(a_{1},\allowbreak{} t) = s \land |t| = n)\).
Consider the following case.
\begin{proofstep}
Assume \(\exists a\, :\, \alpha,\; \exists t\, :\, \mathcal{P}_{\mathrm{fin}}(\alpha),\; \exists h\text{ proving }(a \notin t),\; \operatorname{Finset\mathord{:}cons}(a,\allowbreak{} t) = s \land |t| = n\).
Use the hypothesis \(\exists a\, :\, \alpha,\; \exists t\, :\, \mathcal{P}_{\mathrm{fin}}(\alpha),\; \exists h\text{ proving }(a \notin t),\; \operatorname{Finset\mathord{:}cons}(a,\allowbreak{} t) = s \land |t| = n\).
Take the witnesses from the preceding result.
\begin{proofstep}
Fix \(a\) in \(\alpha\).
Fix \(t\) in \(\mathcal{P}_{\mathrm{fin}}(\alpha)\).
Assume \(a \notin t\).
Assume \(\operatorname{Finset\mathord{:}cons}\left(a, t\right) = s\).
Assume \(\left|t\right| = n\).
Use the hypothesis \(\left|t\right| = n\).
Rewriting with \(\operatorname{Finset\mathord{:}cons}\left(a, t\right) = s\), \paperlink{theorem:16}{Lemma~\ref{theorem:16}} (using \(a \notin t\)) and \paperlink{premise:4e61742e6164645f6c6566745f696e6a}{\emph{Alias}} (\paperlink{premise:4e61742e6164645f6c6566745f696e6a}{Prerequisite~P\ref{premise:4e61742e6164645f6c6566745f696e6a}}) in the preceding result gives \(\left|s\right| = n + 1\).
\end{proofstep}
Using the witnesses and their properties, we obtain \(\left|s\right| = n + 1\).
\end{proofstep}
Together the two implications give the equivalence. Hence \(|s| = n + 1 \iff (\exists a\, :\, \alpha,\; \exists t\, :\, \mathcal{P}_{\mathrm{fin}}(\alpha),\; \exists h\text{ proving }(a \notin t),\; \operatorname{Finset\mathord{:}cons}(a,\allowbreak{} t) = s \land |t| = n)\).
\end{proof}

\section{Decidable equals}
\begin{lemma}\label{theorem:120}
Let \(\alpha\) be a type. Let \(s\) be a finite subset of \(\alpha\). Let \(n\) be a natural number. Assume equality on \(\alpha\) is decidable. Then
\begin{center}
\(\displaystyle |s| = n + 1 \iff (\exists a\, :\, \alpha,\; \exists t\, :\, \mathcal{P}_{\mathrm{fin}}(\alpha),\; a \notin t \land (\{a\} \cup t = s \land |t| = n)).\)
\end{center}
\end{lemma}

\begin{proof}
Consider the following case.
\begin{proofstep}
Assume \(\left|s\right| = n + 1\).
By \paperlink{theorem:9}{Lemma~\ref{theorem:9}}, we obtain \(0 < \left|s\right| \iff s\ne\varnothing\).
By \paperlink{premise:4e61742e7a65726f5f6c745f73756363}{Prerequisite~P\ref{premise:4e61742e7a65726f5f6c745f73756363}}, we obtain \(0 < n + 1\).
Use the hypothesis \(\left|s\right| = n + 1\).
Reversing the equality, we obtain \(n + 1 = \left|s\right|\).
Substituting equals, we obtain \(0 < \left|s\right|\).
By the forward implication, we obtain \(s\ne\varnothing\).
Take the witnesses from the preceding result.
\begin{proofstep}
Fix \(a\) in \(\alpha\).
Assume \(a \in s\).
Choose the witness \(a\).
Choose the witness \(s \setminus \{a\}\).
By \paperlink{premise:46696e7365742e6e6f745f6d656d5f6572617365}{Prerequisite~P\ref{premise:46696e7365742e6e6f745f6d656d5f6572617365}}, we obtain \(a \notin s \setminus \{a\}\).
Use the hypothesis \(a \in s\).
By \paperlink{premise:46696e7365742e696e736572745f6572617365}{Prerequisite~P\ref{premise:46696e7365742e696e736572745f6572617365}}, we obtain \(\{a\} \cup \left(s \setminus \{a\}\right) = s\).
Rewriting with \paperlink{theorem:28}{Lemma~\ref{theorem:28}} (using \(a \in s\)), \(\left|s\right| = n + 1\) and \paperlink{premise:4e61742e6164645f7375625f63616e63656c5f7269676874}{Prerequisite~P\ref{premise:4e61742e6164645f7375625f63616e63656c5f7269676874}} gives \(\left|s \setminus \{a\}\right| = n\).
Combining the two conclusions, we obtain \(\{a\} \cup \left(s \setminus \{a\}\right) = s \land \left|s \setminus \{a\}\right| = n\).
Combining the two conclusions, we obtain \(a \notin s \setminus \{a\} \land (\{a\} \cup (s \setminus \{a\}) = s \land |s \setminus \{a\}| = n)\).
This witness proves existence. Hence \(\exists t\, :\, \mathcal{P}_{\mathrm{fin}}(\alpha),\; a \notin t \land (\{a\} \cup t = s \land |t| = n)\).
This witness proves existence. Hence \(\exists a_{1}\, :\, \alpha,\; \exists t\, :\, \mathcal{P}_{\mathrm{fin}}(\alpha),\; a_{1} \notin t \land (\{a_{1}\} \cup t = s \land |t| = n)\).
\end{proofstep}
Using the witnesses and their properties, we obtain \(\exists a\, :\, \alpha,\; \exists t\, :\, \mathcal{P}_{\mathrm{fin}}(\alpha),\; a \notin t \land (\{a\} \cup t = s \land |t| = n)\).
\end{proofstep}
Consider the following case.
\begin{proofstep}
Assume \(\exists a\, :\, \alpha,\; \exists t\, :\, \mathcal{P}_{\mathrm{fin}}(\alpha),\; a \notin t \land (\{a\} \cup t = s \land |t| = n)\).
Use the hypothesis \(\exists a\, :\, \alpha,\; \exists t\, :\, \mathcal{P}_{\mathrm{fin}}(\alpha),\; a \notin t \land (\{a\} \cup t = s \land |t| = n)\).
Take the witnesses from the preceding result.
\begin{proofstep}
Fix \(w\) in \(\alpha\).
Fix \(w_{1}\) in \(\mathcal{P}_{\mathrm{fin}}(\alpha)\).
Assume \(w \notin w_{1}\).
Assume \(\{w\} \cup w_{1} = s\).
Assume \(\left|w_{1}\right| = n\).
By \paperlink{theorem:17}{Lemma~\ref{theorem:17}}, we obtain \(\left|\{w\} \cup w_{1}\right| = \left|w_{1}\right| + 1\).
Use the hypothesis \(\left|w_{1}\right| = n\).
Substituting equals, we obtain \(\left|\{w\} \cup w_{1}\right| = n + 1\).
Use the hypothesis \(\{w\} \cup w_{1} = s\).
Substituting equals, we obtain \(\left|s\right| = n + 1\).
\end{proofstep}
Using the witnesses and their properties, we obtain \(\left|s\right| = n + 1\).
\end{proofstep}
Together the two implications give the equivalence. Hence \(|s| = n + 1 \iff (\exists a\, :\, \alpha,\; \exists t\, :\, \mathcal{P}_{\mathrm{fin}}(\alpha),\; a \notin t \land (\{a\} \cup t = s \land |t| = n))\).
\end{proof}

\begin{lemma}\label{theorem:121}
Let \(\alpha\) be a type. Let \(s\) be a finite subset of \(\alpha\). Assume equality on \(\alpha\) is decidable. Then
\begin{center}
\(\displaystyle |s| = 2 \iff (\exists x\, :\, \alpha,\; \exists y\, :\, \alpha,\; x \ne y \land s = \{x\} \cup \{y\}).\)
\end{center}
\end{lemma}

\begin{proof}
The substitution uses the following argument.
\begin{proofstep}
Fix \(a\) in \(\alpha\).
For the required implication, argue as follows.
\begin{proofstep}
Fix \(t\) in \(\mathcal{P}_{\mathrm{fin}}(\alpha)\).
By \paperlink{theorem:106}{Lemma~\ref{theorem:106}}, we obtain \(\left|t\right| = 1 \iff \left(\exists a_{1}\, :\, \alpha,\; t = \{a_{1}\}\right)\).
By propositional extensionality, we obtain \(\left(\left|t\right| = 1\right) = \left(\exists a_{1}\, :\, \alpha,\; t = \{a_{1}\}\right)\).
Applying the same function to equal arguments, we obtain \((a \notin t \land (\{a\} \cup t = s \land |t| = 1)) = (a \notin t \land (\{a\} \cup t = s \land (\exists a_{1}\, :\, \alpha,\; t = \{a_{1}\})))\).
\end{proofstep}
By function extensionality, we obtain \((x \mapsto a \notin x \land (\{a\} \cup x = s \land |x| = 1)) = (x \mapsto a \notin x \land (\{a\} \cup x = s \land (\exists a_{1}\, :\, \alpha,\; x = \{a_{1}\})))\).
Applying the same function to equal arguments, we obtain \((\exists x\, :\, \mathcal{P}_{\mathrm{fin}}(\alpha),\; a \notin x \land (\{a\} \cup x = s \land |x| = 1)) = (\exists x\, :\, \mathcal{P}_{\mathrm{fin}}(\alpha),\; a \notin x \land (\{a\} \cup x = s \land (\exists a_{1}\, :\, \alpha,\; x = \{a_{1}\})))\).
\end{proofstep}
Assume \(\exists a_{1}\, :\, \alpha,\; \exists t\, :\, \mathcal{P}_{\mathrm{fin}}(\alpha),\; a_{1} \notin t \land (\{a_{1}\} \cup t = s \land (\exists a_{2}\, :\, \alpha,\; t = \{a_{2}\}))\).
Use the hypothesis \(\exists a_{1}\, :\, \alpha,\; \exists t\, :\, \mathcal{P}_{\mathrm{fin}}(\alpha),\; a_{1} \notin t \land (\{a_{1}\} \cup t = s \land (\exists a_{2}\, :\, \alpha,\; t = \{a_{2}\}))\).
Consider the following case.
\begin{proofstep}
Fix \(a_{1}\) in \(\alpha\).
Assume \(\exists t\, :\, \mathcal{P}_{\mathrm{fin}}(\alpha),\; a_{1} \notin t \land (\{a_{1}\} \cup t = s \land (\exists a_{2}\, :\, \alpha,\; t = \{a_{2}\}))\).
Use the hypothesis \(\exists t\, :\, \mathcal{P}_{\mathrm{fin}}(\alpha),\; a_{1} \notin t \land (\{a_{1}\} \cup t = s \land (\exists a_{2}\, :\, \alpha,\; t = \{a_{2}\}))\).
Consider the following case.
\begin{proofstep}
Fix \(w\) in \(\mathcal{P}_{\mathrm{fin}}(\alpha)\).
Assume \(a_{1} \notin w \land (\{a_{1}\} \cup w = s \land (\exists a_{2}\, :\, \alpha,\; w = \{a_{2}\}))\).
Use the hypothesis \(a_{1} \notin w \land (\{a_{1}\} \cup w = s \land (\exists a_{2}\, :\, \alpha,\; w = \{a_{2}\}))\).
Consider the following case.
\begin{proofstep}
Assume \(a_{1} \notin w\).
Assume \(\{a_{1}\} \cup w = s \land \left(\exists a_{2}\, :\, \alpha,\; w = \{a_{2}\}\right)\).
Use the hypothesis \(\{a_{1}\} \cup w = s \land \left(\exists a_{2}\, :\, \alpha,\; w = \{a_{2}\}\right)\).
Consider the following case.
\begin{proofstep}
Assume \(\{a_{1}\} \cup w = s\).
Assume \(\exists a_{2}\, :\, \alpha,\; w = \{a_{2}\}\).
Use the hypothesis \(\exists a_{2}\, :\, \alpha,\; w = \{a_{2}\}\).
Consider the following case.
\begin{proofstep}
Fix \(b\) in \(\alpha\).
Assume \(w = \{b\}\).
For the required implication, argue as follows.
\begin{proofstep}
Assume \(a_{1} \notin \{b\}\).
Choose the witness \(a_{1}\).
Choose the witness \(b\).
By \paperlink{premise:46696e7365742e6e6f745f6d656d5f73696e676c65746f6e}{Prerequisite~P\ref{premise:46696e7365742e6e6f745f6d656d5f73696e676c65746f6e}}, we obtain \(a_{1} \notin \{b\} \iff a_{1} \ne b\).
Use the hypothesis \(a_{1} \notin \{b\}\).
By the forward implication, we obtain \(a_{1} \ne b\).
By definition, we obtain \(\{a_{1}\} \cup \{b\} = \{a_{1}\} \cup \{b\}\).
Combining the two conclusions, we obtain \(a_{1} \ne b \land \{a_{1}\} \cup \{b\} = \{a_{1}\} \cup \{b\}\).
This witness proves existence. Hence \(\exists y\, :\, \alpha,\; a_{1} \ne y \land \{a_{1}\} \cup \{b\} = \{a_{1}\} \cup \{y\}\).
This witness proves existence. Hence \(\exists x\, :\, \alpha,\; \exists y\, :\, \alpha,\; x \ne y \land \{a_{1}\} \cup \{b\} = \{x\} \cup \{y\}\).
\end{proofstep}
Use the hypothesis \(w = \{b\}\).
Reversing the equality, we obtain \(\{b\} = w\).
Use the hypothesis \(a_{1} \notin w\).
Substituting equals, we obtain \(\exists x\, :\, \alpha,\; \exists y\, :\, \alpha,\; x \ne y \land \{a_{1}\} \cup w = \{x\} \cup \{y\}\).
\end{proofstep}
Using the existential hypothesis, we obtain \(\exists x\, :\, \alpha,\; \exists y\, :\, \alpha,\; x \ne y \land \{a_{1}\} \cup w = \{x\} \cup \{y\}\).
Use the hypothesis \(\{a_{1}\} \cup w = s\).
Substituting equals, we obtain \(\exists x\, :\, \alpha,\; \exists y\, :\, \alpha,\; x \ne y \land s = \{x\} \cup \{y\}\).
\end{proofstep}
Using the two conjuncts, we obtain \(\exists x\, :\, \alpha,\; \exists y\, :\, \alpha,\; x \ne y \land s = \{x\} \cup \{y\}\).
\end{proofstep}
Using the two conjuncts, we obtain \(\exists x\, :\, \alpha,\; \exists y\, :\, \alpha,\; x \ne y \land s = \{x\} \cup \{y\}\).
\end{proofstep}
Using the existential hypothesis, we obtain \(\exists x\, :\, \alpha,\; \exists y\, :\, \alpha,\; x \ne y \land s = \{x\} \cup \{y\}\).
\end{proofstep}
Using the existential hypothesis, we obtain \(\exists x\, :\, \alpha,\; \exists y\, :\, \alpha,\; x \ne y \land s = \{x\} \cup \{y\}\).
Rewriting with \paperlink{premise:696d706c6965735f636f6e6772}{Prerequisite~P\ref{premise:696d706c6965735f636f6e6772}}, \paperlink{premise:66756e657874}{\emph{Function extensionality}} (\paperlink{premise:66756e657874}{Prerequisite~P\ref{premise:66756e657874}}) and the identity established above in the preceding result gives \((\exists a\, :\, \alpha,\; \exists t\, :\, \mathcal{P}_{\mathrm{fin}}(\alpha),\; a \notin t \land (\{a\} \cup t = s \land |t| = 1)) \implies (\exists x\, :\, \alpha,\; \exists y\, :\, \alpha,\; x \ne y \land s = \{x\} \cup \{y\})\).
Rewriting with \paperlink{theorem:120}{Lemma~\ref{theorem:120}} in the preceding result gives \(|s| = 2 \implies (\exists x\, :\, \alpha,\; \exists y\, :\, \alpha,\; x \ne y \land s = \{x\} \cup \{y\})\).
Consider the following case.
\begin{proofstep}
Assume \(\exists x\, :\, \alpha,\; \exists y\, :\, \alpha,\; x \ne y \land s = \{x\} \cup \{y\}\).
Use the hypothesis \(\exists x\, :\, \alpha,\; \exists y\, :\, \alpha,\; x \ne y \land s = \{x\} \cup \{y\}\).
Consider the following case.
\begin{proofstep}
Fix \(x\) in \(\alpha\).
Assume \(\exists y\, :\, \alpha,\; x \ne y \land s = \{x\} \cup \{y\}\).
Use the hypothesis \(\exists y\, :\, \alpha,\; x \ne y \land s = \{x\} \cup \{y\}\).
Consider the following case.
\begin{proofstep}
Fix \(y\) in \(\alpha\).
Assume \(x \ne y \land s = \{x\} \cup \{y\}\).
Use the hypothesis \(x \ne y \land s = \{x\} \cup \{y\}\).
Consider the following case.
\begin{proofstep}
Assume \(x \ne y\).
Assume \(s = \{x\} \cup \{y\}\).
By \paperlink{theorem:27}{Lemma~\ref{theorem:27}}, we obtain \(\left|\{x\} \cup \{y\}\right| = 2\).
Use the hypothesis \(s = \{x\} \cup \{y\}\).
Reversing the equality, we obtain \(\{x\} \cup \{y\} = s\).
Substituting equals, we obtain \(\left|s\right| = 2\).
\end{proofstep}
Using the two conjuncts, we obtain \(\left|s\right| = 2\).
\end{proofstep}
Using the existential hypothesis, we obtain \(\left|s\right| = 2\).
\end{proofstep}
Using the existential hypothesis, we obtain \(\left|s\right| = 2\).
\end{proofstep}
Together the two implications give the equivalence. Hence \(|s| = 2 \iff (\exists x\, :\, \alpha,\; \exists y\, :\, \alpha,\; x \ne y \land s = \{x\} \cup \{y\})\).
\end{proof}

\begin{lemma}\label{theorem:122}
Let \(\alpha\) be a type. Let \(s\) be a finite subset of \(\alpha\). Assume equality on \(\alpha\) is decidable. Then
\begin{center}
\(\displaystyle |s| = 3 \iff (\exists x\, :\, \alpha,\; \exists y\, :\, \alpha,\; \exists z\, :\, \alpha,\; x \ne y \land (x \ne z \land (y \ne z \land s = \{x\} \cup (\{y\} \cup \{z\})))).\)
\end{center}
\end{lemma}

\begin{proof}
The substitution uses the following argument.
\begin{proofstep}
Fix \(a\) in \(\alpha\).
For the required implication, argue as follows.
\begin{proofstep}
Fix \(t\) in \(\mathcal{P}_{\mathrm{fin}}(\alpha)\).
By \paperlink{theorem:121}{Lemma~\ref{theorem:121}}, we obtain \(|t| = 2 \iff (\exists x\, :\, \alpha,\; \exists y\, :\, \alpha,\; x \ne y \land t = \{x\} \cup \{y\})\).
By propositional extensionality, we obtain \((|t| = 2) = (\exists x\, :\, \alpha,\; \exists y\, :\, \alpha,\; x \ne y \land t = \{x\} \cup \{y\})\).
Applying the same function to equal arguments, we obtain \((a \notin t \land (\{a\} \cup t = s \land |t| = 2)) = (a \notin t \land (\{a\} \cup t = s \land (\exists x\, :\, \alpha,\; \exists y\, :\, \alpha,\; x \ne y \land t = \{x\} \cup \{y\})))\).
\end{proofstep}
By function extensionality, we obtain \((x \mapsto a \notin x \land (\{a\} \cup x = s \land |x| = 2)) = (x \mapsto a \notin x \land (\{a\} \cup x = s \land (\exists x_{1}\, :\, \alpha,\; \exists y\, :\, \alpha,\; x_{1} \ne y \land x = \{x_{1}\} \cup \{y\})))\).
Applying the same function to equal arguments, we obtain \((\exists x\, :\, \mathcal{P}_{\mathrm{fin}}(\alpha),\; a \notin x \land (\{a\} \cup x = s \land |x| = 2)) = (\exists x\, :\, \mathcal{P}_{\mathrm{fin}}(\alpha),\; a \notin x \land (\{a\} \cup x = s \land (\exists x_{1}\, :\, \alpha,\; \exists y\, :\, \alpha,\; x_{1} \ne y \land x = \{x_{1}\} \cup \{y\})))\).
\end{proofstep}
Assume \(\exists a_{1}\, :\, \alpha,\; \exists t\, :\, \mathcal{P}_{\mathrm{fin}}(\alpha),\; a_{1} \notin t \land (\{a_{1}\} \cup t = s \land (\exists x\, :\, \alpha,\; \exists y\, :\, \alpha,\; x \ne y \land t = \{x\} \cup \{y\}))\).
Use the hypothesis \(\exists a_{1}\, :\, \alpha,\; \exists t\, :\, \mathcal{P}_{\mathrm{fin}}(\alpha),\; a_{1} \notin t \land (\{a_{1}\} \cup t = s \land (\exists x\, :\, \alpha,\; \exists y\, :\, \alpha,\; x \ne y \land t = \{x\} \cup \{y\}))\).
Consider the following case.
\begin{proofstep}
Fix \(a_{1}\) in \(\alpha\).
Assume \(\exists t\, :\, \mathcal{P}_{\mathrm{fin}}(\alpha),\; a_{1} \notin t \land (\{a_{1}\} \cup t = s \land (\exists x\, :\, \alpha,\; \exists y\, :\, \alpha,\; x \ne y \land t = \{x\} \cup \{y\}))\).
Use the hypothesis \(\exists t\, :\, \mathcal{P}_{\mathrm{fin}}(\alpha),\; a_{1} \notin t \land (\{a_{1}\} \cup t = s \land (\exists x\, :\, \alpha,\; \exists y\, :\, \alpha,\; x \ne y \land t = \{x\} \cup \{y\}))\).
Consider the following case.
\begin{proofstep}
Fix \(w\) in \(\mathcal{P}_{\mathrm{fin}}(\alpha)\).
Assume \(a_{1} \notin w \land (\{a_{1}\} \cup w = s \land (\exists x\, :\, \alpha,\; \exists y\, :\, \alpha,\; x \ne y \land w = \{x\} \cup \{y\}))\).
Use the hypothesis \(a_{1} \notin w \land (\{a_{1}\} \cup w = s \land (\exists x\, :\, \alpha,\; \exists y\, :\, \alpha,\; x \ne y \land w = \{x\} \cup \{y\}))\).
Consider the following case.
\begin{proofstep}
Assume \(a_{1} \notin w\).
Assume \(\{a_{1}\} \cup w = s \land (\exists x\, :\, \alpha,\; \exists y\, :\, \alpha,\; x \ne y \land w = \{x\} \cup \{y\})\).
Use the hypothesis \(\{a_{1}\} \cup w = s \land (\exists x\, :\, \alpha,\; \exists y\, :\, \alpha,\; x \ne y \land w = \{x\} \cup \{y\})\).
Consider the following case.
\begin{proofstep}
Assume \(\{a_{1}\} \cup w = s\).
Assume \(\exists x\, :\, \alpha,\; \exists y\, :\, \alpha,\; x \ne y \land w = \{x\} \cup \{y\}\).
Use the hypothesis \(\exists x\, :\, \alpha,\; \exists y\, :\, \alpha,\; x \ne y \land w = \{x\} \cup \{y\}\).
Consider the following case.
\begin{proofstep}
Fix \(b\) in \(\alpha\).
Assume \(\exists y\, :\, \alpha,\; b \ne y \land w = \{b\} \cup \{y\}\).
Use the hypothesis \(\exists y\, :\, \alpha,\; b \ne y \land w = \{b\} \cup \{y\}\).
Consider the following case.
\begin{proofstep}
Fix \(c\) in \(\alpha\).
Assume \(b \ne c \land w = \{b\} \cup \{c\}\).
Use the hypothesis \(b \ne c \land w = \{b\} \cup \{c\}\).
Consider the following case.
\begin{proofstep}
Assume \(b \ne c\).
Assume \(w = \{b\} \cup \{c\}\).
For the required implication, argue as follows.
\begin{proofstep}
Assume \(a_{1} \notin \{b\} \cup \{c\}\).
Choose the witness \(a_{1}\).
Choose the witness \(b\).
Choose the witness \(c\).
Use the hypothesis \(a_{1} \notin \{b\} \cup \{c\}\).
Rewriting with \paperlink{premise:46696e7365742e6d656d5f696e73657274}{Prerequisite~P\ref{premise:46696e7365742e6d656d5f696e73657274}} in the preceding result gives \(\neg \left(a_{1} = b \lor a_{1} \in \{c\}\right)\).
Rewriting with \paperlink{premise:46696e7365742e6d656d5f73696e676c65746f6e}{Prerequisite~P\ref{premise:46696e7365742e6d656d5f73696e676c65746f6e}} in the preceding result gives \(\neg \left(a_{1} = b \lor a_{1} = c\right)\).
Rewriting with \paperlink{premise:6e6f745f6f72}{Prerequisite~P\ref{premise:6e6f745f6f72}} in the preceding result gives \(\neg \left(a_{1} = b\right) \land \neg \left(a_{1} = c\right)\).
Taking the first conjunct, we obtain \(\neg \left(a_{1} = b\right)\).
Use the hypothesis \(a_{1} \notin \{b\} \cup \{c\}\).
Rewriting with \paperlink{premise:46696e7365742e6d656d5f696e73657274}{Prerequisite~P\ref{premise:46696e7365742e6d656d5f696e73657274}} in the preceding result gives \(\neg \left(a_{1} = b \lor a_{1} \in \{c\}\right)\).
Rewriting with \paperlink{premise:46696e7365742e6d656d5f73696e676c65746f6e}{Prerequisite~P\ref{premise:46696e7365742e6d656d5f73696e676c65746f6e}} in the preceding result gives \(\neg \left(a_{1} = b \lor a_{1} = c\right)\).
Rewriting with \paperlink{premise:6e6f745f6f72}{Prerequisite~P\ref{premise:6e6f745f6f72}} in the preceding result gives \(\neg \left(a_{1} = b\right) \land \neg \left(a_{1} = c\right)\).
Taking the second conjunct, we obtain \(\neg \left(a_{1} = c\right)\).
Use the hypothesis \(b \ne c\).
By definition, we obtain \(\{a_{1}\} \cup \left(\{b\} \cup \{c\}\right) = \{a_{1}\} \cup \left(\{b\} \cup \{c\}\right)\).
Combining the two conclusions, we obtain \(b \ne c \land \{a_{1}\} \cup \left(\{b\} \cup \{c\}\right) = \{a_{1}\} \cup \left(\{b\} \cup \{c\}\right)\).
Combining the two conclusions, we obtain \(a_{1} \ne c \land (b \ne c \land \{a_{1}\} \cup (\{b\} \cup \{c\}) = \{a_{1}\} \cup (\{b\} \cup \{c\}))\).
Combining the two conclusions, we obtain \(a_{1} \ne b \land (a_{1} \ne c \land (b \ne c \land \{a_{1}\} \cup (\{b\} \cup \{c\}) = \{a_{1}\} \cup (\{b\} \cup \{c\})))\).
This witness proves existence. Hence \(\exists z\, :\, \alpha,\; a_{1} \ne b \land (a_{1} \ne z \land (b \ne z \land \{a_{1}\} \cup (\{b\} \cup \{c\}) = \{a_{1}\} \cup (\{b\} \cup \{z\})))\).
This witness proves existence. Hence \(\exists y\, :\, \alpha,\; \exists z\, :\, \alpha,\; a_{1} \ne y \land (a_{1} \ne z \land (y \ne z \land \{a_{1}\} \cup (\{b\} \cup \{c\}) = \{a_{1}\} \cup (\{y\} \cup \{z\})))\).
This witness proves existence. Hence \(\exists x\, :\, \alpha,\; \exists y\, :\, \alpha,\; \exists z\, :\, \alpha,\; x \ne y \land (x \ne z \land (y \ne z \land \{a_{1}\} \cup (\{b\} \cup \{c\}) = \{x\} \cup (\{y\} \cup \{z\})))\).
\end{proofstep}
Use the hypothesis \(w = \{b\} \cup \{c\}\).
Reversing the equality, we obtain \(\{b\} \cup \{c\} = w\).
Use the hypothesis \(a_{1} \notin w\).
Substituting equals, we obtain \(\exists x\, :\, \alpha,\; \exists y\, :\, \alpha,\; \exists z\, :\, \alpha,\; x \ne y \land (x \ne z \land (y \ne z \land \{a_{1}\} \cup w = \{x\} \cup (\{y\} \cup \{z\})))\).
\end{proofstep}
Using the two conjuncts, we obtain \(\exists x\, :\, \alpha,\; \exists y\, :\, \alpha,\; \exists z\, :\, \alpha,\; x \ne y \land (x \ne z \land (y \ne z \land \{a_{1}\} \cup w = \{x\} \cup (\{y\} \cup \{z\})))\).
\end{proofstep}
Using the existential hypothesis, we obtain \(\exists x\, :\, \alpha,\; \exists y\, :\, \alpha,\; \exists z\, :\, \alpha,\; x \ne y \land (x \ne z \land (y \ne z \land \{a_{1}\} \cup w = \{x\} \cup (\{y\} \cup \{z\})))\).
\end{proofstep}
Using the existential hypothesis, we obtain \(\exists x\, :\, \alpha,\; \exists y\, :\, \alpha,\; \exists z\, :\, \alpha,\; x \ne y \land (x \ne z \land (y \ne z \land \{a_{1}\} \cup w = \{x\} \cup (\{y\} \cup \{z\})))\).
Use the hypothesis \(\{a_{1}\} \cup w = s\).
Substituting equals, we obtain \(\exists x\, :\, \alpha,\; \exists y\, :\, \alpha,\; \exists z\, :\, \alpha,\; x \ne y \land (x \ne z \land (y \ne z \land s = \{x\} \cup (\{y\} \cup \{z\})))\).
\end{proofstep}
Using the two conjuncts, we obtain \(\exists x\, :\, \alpha,\; \exists y\, :\, \alpha,\; \exists z\, :\, \alpha,\; x \ne y \land (x \ne z \land (y \ne z \land s = \{x\} \cup (\{y\} \cup \{z\})))\).
\end{proofstep}
Using the two conjuncts, we obtain \(\exists x\, :\, \alpha,\; \exists y\, :\, \alpha,\; \exists z\, :\, \alpha,\; x \ne y \land (x \ne z \land (y \ne z \land s = \{x\} \cup (\{y\} \cup \{z\})))\).
\end{proofstep}
Using the existential hypothesis, we obtain \(\exists x\, :\, \alpha,\; \exists y\, :\, \alpha,\; \exists z\, :\, \alpha,\; x \ne y \land (x \ne z \land (y \ne z \land s = \{x\} \cup (\{y\} \cup \{z\})))\).
\end{proofstep}
Using the existential hypothesis, we obtain \(\exists x\, :\, \alpha,\; \exists y\, :\, \alpha,\; \exists z\, :\, \alpha,\; x \ne y \land (x \ne z \land (y \ne z \land s = \{x\} \cup (\{y\} \cup \{z\})))\).
Rewriting with \paperlink{premise:696d706c6965735f636f6e6772}{Prerequisite~P\ref{premise:696d706c6965735f636f6e6772}}, \paperlink{premise:66756e657874}{\emph{Function extensionality}} (\paperlink{premise:66756e657874}{Prerequisite~P\ref{premise:66756e657874}}) and the identity established above in the preceding result gives \((\exists a\, :\, \alpha,\; \exists t\, :\, \mathcal{P}_{\mathrm{fin}}(\alpha),\; a \notin t \land (\{a\} \cup t = s \land |t| = 2)) \implies (\exists x\, :\, \alpha,\; \exists y\, :\, \alpha,\; \exists z\, :\, \alpha,\; x \ne y \land (x \ne z \land (y \ne z \land s = \{x\} \cup (\{y\} \cup \{z\}))))\).
Rewriting with \paperlink{theorem:120}{Lemma~\ref{theorem:120}} in the preceding result gives \(|s| = 3 \implies (\exists x\, :\, \alpha,\; \exists y\, :\, \alpha,\; \exists z\, :\, \alpha,\; x \ne y \land (x \ne z \land (y \ne z \land s = \{x\} \cup (\{y\} \cup \{z\}))))\).
Consider the following case.
\begin{proofstep}
Assume \(\exists x\, :\, \alpha,\; \exists y\, :\, \alpha,\; \exists z\, :\, \alpha,\; x \ne y \land (x \ne z \land (y \ne z \land s = \{x\} \cup (\{y\} \cup \{z\})))\).
Use the hypothesis \(\exists x\, :\, \alpha,\; \exists y\, :\, \alpha,\; \exists z\, :\, \alpha,\; x \ne y \land (x \ne z \land (y \ne z \land s = \{x\} \cup (\{y\} \cup \{z\})))\).
Consider the following case.
\begin{proofstep}
Fix \(x\) in \(\alpha\).
Assume \(\exists y\, :\, \alpha,\; \exists z\, :\, \alpha,\; x \ne y \land (x \ne z \land (y \ne z \land s = \{x\} \cup (\{y\} \cup \{z\})))\).
Use the hypothesis \(\exists y\, :\, \alpha,\; \exists z\, :\, \alpha,\; x \ne y \land (x \ne z \land (y \ne z \land s = \{x\} \cup (\{y\} \cup \{z\})))\).
Consider the following case.
\begin{proofstep}
Fix \(y\) in \(\alpha\).
Assume \(\exists z\, :\, \alpha,\; x \ne y \land (x \ne z \land (y \ne z \land s = \{x\} \cup (\{y\} \cup \{z\})))\).
Use the hypothesis \(\exists z\, :\, \alpha,\; x \ne y \land (x \ne z \land (y \ne z \land s = \{x\} \cup (\{y\} \cup \{z\})))\).
Consider the following case.
\begin{proofstep}
Fix \(z\) in \(\alpha\).
Assume \(x \ne y \land (x \ne z \land (y \ne z \land s = \{x\} \cup (\{y\} \cup \{z\})))\).
Use the hypothesis \(x \ne y \land (x \ne z \land (y \ne z \land s = \{x\} \cup (\{y\} \cup \{z\})))\).
Consider the following case.
\begin{proofstep}
Assume \(x \ne y\).
Assume \(x \ne z \land \left(y \ne z \land s = \{x\} \cup \left(\{y\} \cup \{z\}\right)\right)\).
Use the hypothesis \(x \ne z \land \left(y \ne z \land s = \{x\} \cup \left(\{y\} \cup \{z\}\right)\right)\).
Consider the following case.
\begin{proofstep}
Assume \(x \ne z\).
Assume \(y \ne z \land s = \{x\} \cup \left(\{y\} \cup \{z\}\right)\).
Use the hypothesis \(y \ne z \land s = \{x\} \cup \left(\{y\} \cup \{z\}\right)\).
Consider the following case.
\begin{proofstep}
Assume \(y \ne z\).
Assume \(s = \{x\} \cup \left(\{y\} \cup \{z\}\right)\).
Rewriting with \paperlink{theorem:17}{Lemma~\ref{theorem:17}}, \paperlink{premise:46696e7365742e6d656d5f696e73657274}{Prerequisite~P\ref{premise:46696e7365742e6d656d5f696e73657274}}, \paperlink{premise:65715f66616c7365}{Prerequisite~P\ref{premise:65715f66616c7365}} (using \(x \ne y\)), \paperlink{premise:46696e7365742e6d656d5f73696e676c65746f6e}{Prerequisite~P\ref{premise:46696e7365742e6d656d5f73696e676c65746f6e}}, \paperlink{premise:65715f66616c7365}{Prerequisite~P\ref{premise:65715f66616c7365}} (using \(x \ne z\)), \paperlink{premise:6f725f73656c665f696666}{Prerequisite~P\ref{premise:6f725f73656c665f696666}}, \paperlink{premise:6e6f745f66616c73655f696666}{Prerequisite~P\ref{premise:6e6f745f66616c73655f696666}}, \paperlink{premise:65715f66616c7365}{Prerequisite~P\ref{premise:65715f66616c7365}} (using \(y \ne z\)) and \paperlink{theorem:14}{Lemma~\ref{theorem:14}} gives \(\left|\{x\} \cup \left(\{y\} \cup \{z\}\right)\right| = 3\).
Use the hypothesis \(s = \{x\} \cup \left(\{y\} \cup \{z\}\right)\).
Reversing the equality, we obtain \(\{x\} \cup \left(\{y\} \cup \{z\}\right) = s\).
Substituting equals, we obtain \(\left|s\right| = 3\).
\end{proofstep}
Using the two conjuncts, we obtain \(\left|s\right| = 3\).
\end{proofstep}
Using the two conjuncts, we obtain \(\left|s\right| = 3\).
\end{proofstep}
Using the two conjuncts, we obtain \(\left|s\right| = 3\).
\end{proofstep}
Using the existential hypothesis, we obtain \(\left|s\right| = 3\).
\end{proofstep}
Using the existential hypothesis, we obtain \(\left|s\right| = 3\).
\end{proofstep}
Using the existential hypothesis, we obtain \(\left|s\right| = 3\).
\end{proofstep}
Together the two implications give the equivalence. Hence \(|s| = 3 \iff (\exists x\, :\, \alpha,\; \exists y\, :\, \alpha,\; \exists z\, :\, \alpha,\; x \ne y \land (x \ne z \land (y \ne z \land s = \{x\} \cup (\{y\} \cup \{z\}))))\).
\end{proof}

\section{Finite sets}
\begin{lemma}\label{theorem:123}
Let \(\alpha\) be a type. Let \(s\) be a finite subset of \(\alpha\). Then
\begin{center}
\(\displaystyle 2 < |s| \iff (\exists a\, :\, \alpha,\; \exists b\, :\, \alpha,\; \exists c\, :\, \alpha,\; a \in s \land (b \in s \land (c \in s \land (a \ne b \land (a \ne c \land b \ne c))))).\)
\end{center}
\end{lemma}

\begin{proof}
The substitution uses the following argument.
\begin{proofstep}
Fix \(t\) in \(\mathcal{P}_{\mathrm{fin}}(\alpha)\).
By \paperlink{theorem:122}{Lemma~\ref{theorem:122}}, we obtain \(|t| = 3 \iff (\exists x\, :\, \alpha,\; \exists y\, :\, \alpha,\; \exists z\, :\, \alpha,\; x \ne y \land (x \ne z \land (y \ne z \land t = \{x\} \cup (\{y\} \cup \{z\}))))\).
By propositional extensionality, we obtain \((|t| = 3) = (\exists x\, :\, \alpha,\; \exists y\, :\, \alpha,\; \exists z\, :\, \alpha,\; x \ne y \land (x \ne z \land (y \ne z \land t = \{x\} \cup (\{y\} \cup \{z\}))))\).
Applying the same function to equal arguments, we obtain \((t \subseteq s \land |t| = 3) = (t \subseteq s \land (\exists x\, :\, \alpha,\; \exists y\, :\, \alpha,\; \exists z\, :\, \alpha,\; x \ne y \land (x \ne z \land (y \ne z \land t = \{x\} \cup (\{y\} \cup \{z\})))))\).
\end{proofstep}
The substitution uses the following argument.
\begin{proofstep}
Fix \(t\) in \(\mathcal{P}_{\mathrm{fin}}(\alpha)\).
By \paperlink{premise:6578697374735f616e645f6c656674}{Prerequisite~P\ref{premise:6578697374735f616e645f6c656674}}, we obtain \((\exists x\, :\, \alpha,\; t \subseteq s \land (x_{1} \mapsto \exists y\, :\, \alpha,\; \exists z\, :\, \alpha,\; x_{1} \ne y \land (x_{1} \ne z \land (y \ne z \land t = \{x_{1}\} \cup (\{y\} \cup \{z\}))))(x)) \iff t \subseteq s \land (\exists x\, :\, \alpha,\; (x_{1} \mapsto \exists y\, :\, \alpha,\; \exists z\, :\, \alpha,\; x_{1} \ne y \land (x_{1} \ne z \land (y \ne z \land t = \{x_{1}\} \cup (\{y\} \cup \{z\}))))(x))\).
By propositional extensionality, we obtain \((\exists x\, :\, \alpha,\; t \subseteq s \land (x_{1} \mapsto \exists y\, :\, \alpha,\; \exists z\, :\, \alpha,\; x_{1} \ne y \land (x_{1} \ne z \land (y \ne z \land t = \{x_{1}\} \cup (\{y\} \cup \{z\}))))(x)) = (t \subseteq s \land (\exists x\, :\, \alpha,\; (x_{1} \mapsto \exists y\, :\, \alpha,\; \exists z\, :\, \alpha,\; x_{1} \ne y \land (x_{1} \ne z \land (y \ne z \land t = \{x_{1}\} \cup (\{y\} \cup \{z\}))))(x)))\).
Reversing the equality, we obtain \((t \subseteq s \land (\exists x\, :\, \alpha,\; \exists y\, :\, \alpha,\; \exists z\, :\, \alpha,\; x \ne y \land (x \ne z \land (y \ne z \land t = \{x\} \cup (\{y\} \cup \{z\}))))) = (\exists x\, :\, \alpha,\; t \subseteq s \land (\exists y\, :\, \alpha,\; \exists z\, :\, \alpha,\; x \ne y \land (x \ne z \land (y \ne z \land t = \{x\} \cup (\{y\} \cup \{z\})))))\).
For the required implication, argue as follows.
\begin{proofstep}
Fix \(x\) in \(\alpha\).
By \paperlink{premise:6578697374735f616e645f6c656674}{Prerequisite~P\ref{premise:6578697374735f616e645f6c656674}}, we obtain \((\exists x_{1}\, :\, \alpha,\; t \subseteq s \land (x_{2} \mapsto \exists z\, :\, \alpha,\; x \ne x_{2} \land (x \ne z \land (x_{2} \ne z \land t = \{x\} \cup (\{x_{2}\} \cup \{z\}))))(x_{1})) \iff t \subseteq s \land (\exists x_{1}\, :\, \alpha,\; (x_{2} \mapsto \exists z\, :\, \alpha,\; x \ne x_{2} \land (x \ne z \land (x_{2} \ne z \land t = \{x\} \cup (\{x_{2}\} \cup \{z\}))))(x_{1}))\).
By propositional extensionality, we obtain \((\exists x_{1}\, :\, \alpha,\; t \subseteq s \land (x_{2} \mapsto \exists z\, :\, \alpha,\; x \ne x_{2} \land (x \ne z \land (x_{2} \ne z \land t = \{x\} \cup (\{x_{2}\} \cup \{z\}))))(x_{1})) = (t \subseteq s \land (\exists x_{1}\, :\, \alpha,\; (x_{2} \mapsto \exists z\, :\, \alpha,\; x \ne x_{2} \land (x \ne z \land (x_{2} \ne z \land t = \{x\} \cup (\{x_{2}\} \cup \{z\}))))(x_{1})))\).
Reversing the equality, we obtain \((t \subseteq s \land (\exists x_{1}\, :\, \alpha,\; \exists z\, :\, \alpha,\; x \ne x_{1} \land (x \ne z \land (x_{1} \ne z \land t = \{x\} \cup (\{x_{1}\} \cup \{z\}))))) = (\exists x_{1}\, :\, \alpha,\; t \subseteq s \land (\exists z\, :\, \alpha,\; x \ne x_{1} \land (x \ne z \land (x_{1} \ne z \land t = \{x\} \cup (\{x_{1}\} \cup \{z\})))))\).
For the required implication, argue as follows.
\begin{proofstep}
Fix \(x_{1}\) in \(\alpha\).
By \paperlink{premise:6578697374735f616e645f6c656674}{Prerequisite~P\ref{premise:6578697374735f616e645f6c656674}}, we obtain \((\exists x_{2}\, :\, \alpha,\; t \subseteq s \land (x_{3} \mapsto x \ne x_{1} \land (x \ne x_{3} \land (x_{1} \ne x_{3} \land t = \{x\} \cup (\{x_{1}\} \cup \{x_{3}\}))))(x_{2})) \iff t \subseteq s \land (\exists x_{2}\, :\, \alpha,\; (x_{3} \mapsto x \ne x_{1} \land (x \ne x_{3} \land (x_{1} \ne x_{3} \land t = \{x\} \cup (\{x_{1}\} \cup \{x_{3}\}))))(x_{2}))\).
By propositional extensionality, we obtain \((\exists x_{2}\, :\, \alpha,\; t \subseteq s \land (x_{3} \mapsto x \ne x_{1} \land (x \ne x_{3} \land (x_{1} \ne x_{3} \land t = \{x\} \cup (\{x_{1}\} \cup \{x_{3}\}))))(x_{2})) = (t \subseteq s \land (\exists x_{2}\, :\, \alpha,\; (x_{3} \mapsto x \ne x_{1} \land (x \ne x_{3} \land (x_{1} \ne x_{3} \land t = \{x\} \cup (\{x_{1}\} \cup \{x_{3}\}))))(x_{2})))\).
Reversing the equality, we obtain \((t \subseteq s \land (\exists x_{2}\, :\, \alpha,\; x \ne x_{1} \land (x \ne x_{2} \land (x_{1} \ne x_{2} \land t = \{x\} \cup (\{x_{1}\} \cup \{x_{2}\}))))) = (\exists x_{2}\, :\, \alpha,\; t \subseteq s \land (x \ne x_{1} \land (x \ne x_{2} \land (x_{1} \ne x_{2} \land t = \{x\} \cup (\{x_{1}\} \cup \{x_{2}\})))))\).
\end{proofstep}
By function extensionality, we obtain \((x_{1} \mapsto t \subseteq s \land (\exists x_{2}\, :\, \alpha,\; x \ne x_{1} \land (x \ne x_{2} \land (x_{1} \ne x_{2} \land t = \{x\} \cup (\{x_{1}\} \cup \{x_{2}\}))))) = (x_{1} \mapsto \exists x_{2}\, :\, \alpha,\; t \subseteq s \land (x \ne x_{1} \land (x \ne x_{2} \land (x_{1} \ne x_{2} \land t = \{x\} \cup (\{x_{1}\} \cup \{x_{2}\})))))\).
Applying the same function to equal arguments, we obtain \((\exists x_{1}\, :\, \alpha,\; t \subseteq s \land (\exists x_{2}\, :\, \alpha,\; x \ne x_{1} \land (x \ne x_{2} \land (x_{1} \ne x_{2} \land t = \{x\} \cup (\{x_{1}\} \cup \{x_{2}\}))))) = (\exists x_{1}\, :\, \alpha,\; \exists x_{2}\, :\, \alpha,\; t \subseteq s \land (x \ne x_{1} \land (x \ne x_{2} \land (x_{1} \ne x_{2} \land t = \{x\} \cup (\{x_{1}\} \cup \{x_{2}\})))))\).
Chaining the equalities, we obtain \((t \subseteq s \land (\exists x_{1}\, :\, \alpha,\; \exists z\, :\, \alpha,\; x \ne x_{1} \land (x \ne z \land (x_{1} \ne z \land t = \{x\} \cup (\{x_{1}\} \cup \{z\}))))) = (\exists x_{1}\, :\, \alpha,\; \exists x_{2}\, :\, \alpha,\; t \subseteq s \land (x \ne x_{1} \land (x \ne x_{2} \land (x_{1} \ne x_{2} \land t = \{x\} \cup (\{x_{1}\} \cup \{x_{2}\})))))\).
\end{proofstep}
By function extensionality, we obtain \((x \mapsto t \subseteq s \land (\exists x_{1}\, :\, \alpha,\; \exists z\, :\, \alpha,\; x \ne x_{1} \land (x \ne z \land (x_{1} \ne z \land t = \{x\} \cup (\{x_{1}\} \cup \{z\}))))) = (x \mapsto \exists x_{1}\, :\, \alpha,\; \exists x_{2}\, :\, \alpha,\; t \subseteq s \land (x \ne x_{1} \land (x \ne x_{2} \land (x_{1} \ne x_{2} \land t = \{x\} \cup (\{x_{1}\} \cup \{x_{2}\})))))\).
Applying the same function to equal arguments, we obtain \((\exists x\, :\, \alpha,\; t \subseteq s \land (\exists x_{1}\, :\, \alpha,\; \exists z\, :\, \alpha,\; x \ne x_{1} \land (x \ne z \land (x_{1} \ne z \land t = \{x\} \cup (\{x_{1}\} \cup \{z\}))))) = (\exists x\, :\, \alpha,\; \exists x_{1}\, :\, \alpha,\; \exists x_{2}\, :\, \alpha,\; t \subseteq s \land (x \ne x_{1} \land (x \ne x_{2} \land (x_{1} \ne x_{2} \land t = \{x\} \cup (\{x_{1}\} \cup \{x_{2}\})))))\).
Chaining the equalities, we obtain \((t \subseteq s \land (\exists x\, :\, \alpha,\; \exists y\, :\, \alpha,\; \exists z\, :\, \alpha,\; x \ne y \land (x \ne z \land (y \ne z \land t = \{x\} \cup (\{y\} \cup \{z\}))))) = (\exists x\, :\, \alpha,\; \exists x_{1}\, :\, \alpha,\; \exists x_{2}\, :\, \alpha,\; t \subseteq s \land (x \ne x_{1} \land (x \ne x_{2} \land (x_{1} \ne x_{2} \land t = \{x\} \cup (\{x_{1}\} \cup \{x_{2}\})))))\).
\end{proofstep}
The substitution uses the following argument.
\begin{proofstep}
Fix \(b\) in \(\alpha\).
By \paperlink{premise:6578697374735f636f6d6d}{Prerequisite~P\ref{premise:6578697374735f636f6d6d}} with \(\mathcal{P}_{\mathrm{fin}}(\alpha)\) as its type parameter \(\alpha\), we obtain \((\exists a\, :\, \mathcal{P}_{\mathrm{fin}}(\alpha),\; \exists b_{1}\, :\, \alpha,\; (a_{1} \mapsto b_{2} \mapsto \exists x\, :\, \alpha,\; a_{1} \subseteq s \land (b \ne b_{2} \land (b \ne x \land (b_{2} \ne x \land a_{1} = \{b\} \cup (\{b_{2}\} \cup \{x\})))))(a,\allowbreak{} b_{1})) \iff (\exists b_{1}\, :\, \alpha,\; \exists a\, :\, \mathcal{P}_{\mathrm{fin}}(\alpha),\; (a_{1} \mapsto b_{2} \mapsto \exists x\, :\, \alpha,\; a_{1} \subseteq s \land (b \ne b_{2} \land (b \ne x \land (b_{2} \ne x \land a_{1} = \{b\} \cup (\{b_{2}\} \cup \{x\})))))(a,\allowbreak{} b_{1}))\).
By propositional extensionality, we obtain \((\exists a\, :\, \mathcal{P}_{\mathrm{fin}}(\alpha),\; \exists b_{1}\, :\, \alpha,\; \exists x\, :\, \alpha,\; a \subseteq s \land (b \ne b_{1} \land (b \ne x \land (b_{1} \ne x \land a = \{b\} \cup (\{b_{1}\} \cup \{x\}))))) = (\exists b_{1}\, :\, \alpha,\; \exists a\, :\, \mathcal{P}_{\mathrm{fin}}(\alpha),\; \exists x\, :\, \alpha,\; a \subseteq s \land (b \ne b_{1} \land (b \ne x \land (b_{1} \ne x \land a = \{b\} \cup (\{b_{1}\} \cup \{x\})))))\).
For the required implication, argue as follows.
\begin{proofstep}
Fix \(b_{1}\) in \(\alpha\).
By \paperlink{premise:6578697374735f636f6d6d}{Prerequisite~P\ref{premise:6578697374735f636f6d6d}} with \(\mathcal{P}_{\mathrm{fin}}(\alpha)\) as its type parameter \(\alpha\), we obtain \((\exists a\, :\, \mathcal{P}_{\mathrm{fin}}(\alpha),\; \exists b_{2}\, :\, \alpha,\; (a_{1} \mapsto b_{3} \mapsto a_{1} \subseteq s \land (b \ne b_{1} \land (b \ne b_{3} \land (b_{1} \ne b_{3} \land a_{1} = \{b\} \cup (\{b_{1}\} \cup \{b_{3}\})))))(a,\allowbreak{} b_{2})) \iff (\exists b_{2}\, :\, \alpha,\; \exists a\, :\, \mathcal{P}_{\mathrm{fin}}(\alpha),\; (a_{1} \mapsto b_{3} \mapsto a_{1} \subseteq s \land (b \ne b_{1} \land (b \ne b_{3} \land (b_{1} \ne b_{3} \land a_{1} = \{b\} \cup (\{b_{1}\} \cup \{b_{3}\})))))(a,\allowbreak{} b_{2}))\).
By propositional extensionality, we obtain \((\exists a\, :\, \mathcal{P}_{\mathrm{fin}}(\alpha),\; \exists b_{2}\, :\, \alpha,\; a \subseteq s \land (b \ne b_{1} \land (b \ne b_{2} \land (b_{1} \ne b_{2} \land a = \{b\} \cup (\{b_{1}\} \cup \{b_{2}\}))))) = (\exists b_{2}\, :\, \alpha,\; \exists a\, :\, \mathcal{P}_{\mathrm{fin}}(\alpha),\; a \subseteq s \land (b \ne b_{1} \land (b \ne b_{2} \land (b_{1} \ne b_{2} \land a = \{b\} \cup (\{b_{1}\} \cup \{b_{2}\})))))\).
\end{proofstep}
By function extensionality, we obtain \((x \mapsto \exists a\, :\, \mathcal{P}_{\mathrm{fin}}(\alpha),\; \exists b_{1}\, :\, \alpha,\; a \subseteq s \land (b \ne x \land (b \ne b_{1} \land (x \ne b_{1} \land a = \{b\} \cup (\{x\} \cup \{b_{1}\}))))) = (x \mapsto \exists b_{1}\, :\, \alpha,\; \exists a\, :\, \mathcal{P}_{\mathrm{fin}}(\alpha),\; a \subseteq s \land (b \ne x \land (b \ne b_{1} \land (x \ne b_{1} \land a = \{b\} \cup (\{x\} \cup \{b_{1}\})))))\).
Applying the same function to equal arguments, we obtain \((\exists x\, :\, \alpha,\; \exists a\, :\, \mathcal{P}_{\mathrm{fin}}(\alpha),\; \exists b_{1}\, :\, \alpha,\; a \subseteq s \land (b \ne x \land (b \ne b_{1} \land (x \ne b_{1} \land a = \{b\} \cup (\{x\} \cup \{b_{1}\}))))) = (\exists x\, :\, \alpha,\; \exists b_{1}\, :\, \alpha,\; \exists a\, :\, \mathcal{P}_{\mathrm{fin}}(\alpha),\; a \subseteq s \land (b \ne x \land (b \ne b_{1} \land (x \ne b_{1} \land a = \{b\} \cup (\{x\} \cup \{b_{1}\})))))\).
Chaining the equalities, we obtain \((\exists a\, :\, \mathcal{P}_{\mathrm{fin}}(\alpha),\; \exists b_{1}\, :\, \alpha,\; \exists x\, :\, \alpha,\; a \subseteq s \land (b \ne b_{1} \land (b \ne x \land (b_{1} \ne x \land a = \{b\} \cup (\{b_{1}\} \cup \{x\}))))) = (\exists x\, :\, \alpha,\; \exists b_{1}\, :\, \alpha,\; \exists a\, :\, \mathcal{P}_{\mathrm{fin}}(\alpha),\; a \subseteq s \land (b \ne x \land (b \ne b_{1} \land (x \ne b_{1} \land a = \{b\} \cup (\{x\} \cup \{b_{1}\})))))\).
\end{proofstep}
Consider the following case.
\begin{proofstep}
Assume \(\exists b\, :\, \alpha,\; \exists b_{1}\, :\, \alpha,\; \exists b_{2}\, :\, \alpha,\; \exists a_{1}\, :\, \mathcal{P}_{\mathrm{fin}}(\alpha),\; a_{1} \subseteq s \land (b \ne b_{1} \land (b \ne b_{2} \land (b_{1} \ne b_{2} \land a_{1} = \{b\} \cup (\{b_{1}\} \cup \{b_{2}\}))))\).
Use the hypothesis \(\exists b\, :\, \alpha,\; \exists b_{1}\, :\, \alpha,\; \exists b_{2}\, :\, \alpha,\; \exists a_{1}\, :\, \mathcal{P}_{\mathrm{fin}}(\alpha),\; a_{1} \subseteq s \land (b \ne b_{1} \land (b \ne b_{2} \land (b_{1} \ne b_{2} \land a_{1} = \{b\} \cup (\{b_{1}\} \cup \{b_{2}\}))))\).
Consider the following case.
\begin{proofstep}
Fix \(a_{1}\) in \(\alpha\).
Assume \(\exists b\, :\, \alpha,\; \exists b_{1}\, :\, \alpha,\; \exists a_{2}\, :\, \mathcal{P}_{\mathrm{fin}}(\alpha),\; a_{2} \subseteq s \land (a_{1} \ne b \land (a_{1} \ne b_{1} \land (b \ne b_{1} \land a_{2} = \{a_{1}\} \cup (\{b\} \cup \{b_{1}\}))))\).
Use the hypothesis \(\exists b\, :\, \alpha,\; \exists b_{1}\, :\, \alpha,\; \exists a_{2}\, :\, \mathcal{P}_{\mathrm{fin}}(\alpha),\; a_{2} \subseteq s \land (a_{1} \ne b \land (a_{1} \ne b_{1} \land (b \ne b_{1} \land a_{2} = \{a_{1}\} \cup (\{b\} \cup \{b_{1}\}))))\).
Consider the following case.
\begin{proofstep}
Fix \(b\) in \(\alpha\).
Assume \(\exists b_{1}\, :\, \alpha,\; \exists a_{2}\, :\, \mathcal{P}_{\mathrm{fin}}(\alpha),\; a_{2} \subseteq s \land (a_{1} \ne b \land (a_{1} \ne b_{1} \land (b \ne b_{1} \land a_{2} = \{a_{1}\} \cup (\{b\} \cup \{b_{1}\}))))\).
Use the hypothesis \(\exists b_{1}\, :\, \alpha,\; \exists a_{2}\, :\, \mathcal{P}_{\mathrm{fin}}(\alpha),\; a_{2} \subseteq s \land (a_{1} \ne b \land (a_{1} \ne b_{1} \land (b \ne b_{1} \land a_{2} = \{a_{1}\} \cup (\{b\} \cup \{b_{1}\}))))\).
Consider the following case.
\begin{proofstep}
Fix \(c\) in \(\alpha\).
Assume \(\exists a_{2}\, :\, \mathcal{P}_{\mathrm{fin}}(\alpha),\; a_{2} \subseteq s \land (a_{1} \ne b \land (a_{1} \ne c \land (b \ne c \land a_{2} = \{a_{1}\} \cup (\{b\} \cup \{c\}))))\).
Use the hypothesis \(\exists a_{2}\, :\, \mathcal{P}_{\mathrm{fin}}(\alpha),\; a_{2} \subseteq s \land (a_{1} \ne b \land (a_{1} \ne c \land (b \ne c \land a_{2} = \{a_{1}\} \cup (\{b\} \cup \{c\}))))\).
Consider the following case.
\begin{proofstep}
Fix \(t\) in \(\mathcal{P}_{\mathrm{fin}}(\alpha)\).
Assume \(t \subseteq s \land (a_{1} \ne b \land (a_{1} \ne c \land (b \ne c \land t = \{a_{1}\} \cup (\{b\} \cup \{c\}))))\).
Use the hypothesis \(t \subseteq s \land (a_{1} \ne b \land (a_{1} \ne c \land (b \ne c \land t = \{a_{1}\} \cup (\{b\} \cup \{c\}))))\).
Consider the following case.
\begin{proofstep}
Assume \(t \subseteq s\).
Assume \(a_{1} \ne b \land (a_{1} \ne c \land (b \ne c \land t = \{a_{1}\} \cup (\{b\} \cup \{c\})))\).
Use the hypothesis \(a_{1} \ne b \land (a_{1} \ne c \land (b \ne c \land t = \{a_{1}\} \cup (\{b\} \cup \{c\})))\).
Consider the following case.
\begin{proofstep}
Assume \(a_{1} \ne b\).
Assume \(a_{1} \ne c \land \left(b \ne c \land t = \{a_{1}\} \cup \left(\{b\} \cup \{c\}\right)\right)\).
Use the hypothesis \(a_{1} \ne c \land \left(b \ne c \land t = \{a_{1}\} \cup \left(\{b\} \cup \{c\}\right)\right)\).
Consider the following case.
\begin{proofstep}
Assume \(a_{1} \ne c\).
Assume \(b \ne c \land t = \{a_{1}\} \cup \left(\{b\} \cup \{c\}\right)\).
Use the hypothesis \(b \ne c \land t = \{a_{1}\} \cup \left(\{b\} \cup \{c\}\right)\).
Consider the following case.
\begin{proofstep}
Assume \(b \ne c\).
Assume \(t = \{a_{1}\} \cup \left(\{b\} \cup \{c\}\right)\).
For the required implication, argue as follows.
\begin{proofstep}
Assume \(\{a_{1}\} \cup \left(\{b\} \cup \{c\}\right) \subseteq s\).
Choose the witness \(a_{1}\).
Choose the witness \(b\).
Choose the witness \(c\).
Use the hypothesis \(\{a_{1}\} \cup \left(\{b\} \cup \{c\}\right) \subseteq s\).
Rewriting with \paperlink{premise:46696e7365742e696e736572745f7375627365745f696666}{Prerequisite~P\ref{premise:46696e7365742e696e736572745f7375627365745f696666}} and \paperlink{premise:46696e7365742e73696e676c65746f6e5f7375627365745f696666}{Prerequisite~P\ref{premise:46696e7365742e73696e676c65746f6e5f7375627365745f696666}} in the preceding result gives \(a_{1} \in s \land \left(b \in s \land c \in s\right)\).
Taking the first conjunct gives \(a_{1} \in s\).
Use the hypothesis \(\{a_{1}\} \cup \left(\{b\} \cup \{c\}\right) \subseteq s\).
Rewriting with \paperlink{premise:46696e7365742e696e736572745f7375627365745f696666}{Prerequisite~P\ref{premise:46696e7365742e696e736572745f7375627365745f696666}} and \paperlink{premise:46696e7365742e73696e676c65746f6e5f7375627365745f696666}{Prerequisite~P\ref{premise:46696e7365742e73696e676c65746f6e5f7375627365745f696666}} in the preceding result gives \(a_{1} \in s \land \left(b \in s \land c \in s\right)\).
Taking the second conjunct gives \(b \in s \land c \in s\).
Taking the first conjunct gives \(b \in s\).
Use the hypothesis \(\{a_{1}\} \cup \left(\{b\} \cup \{c\}\right) \subseteq s\).
Rewriting with \paperlink{premise:46696e7365742e696e736572745f7375627365745f696666}{Prerequisite~P\ref{premise:46696e7365742e696e736572745f7375627365745f696666}} and \paperlink{premise:46696e7365742e73696e676c65746f6e5f7375627365745f696666}{Prerequisite~P\ref{premise:46696e7365742e73696e676c65746f6e5f7375627365745f696666}} in the preceding result gives \(a_{1} \in s \land \left(b \in s \land c \in s\right)\).
Taking the second conjunct gives \(b \in s \land c \in s\).
Taking the second conjunct gives \(c \in s\).
Rewriting with \(a_{1} \in s\), \(b \in s\), \(c \in s\), \paperlink{premise:65715f66616c7365}{Prerequisite~P\ref{premise:65715f66616c7365}}, \(a_{1} \ne b\), \paperlink{premise:6e6f745f66616c73655f65715f74727565}{Prerequisite~P\ref{premise:6e6f745f66616c73655f65715f74727565}}, \(a_{1} \ne c\), \(b \ne c\) and \paperlink{premise:616e645f73656c66}{Prerequisite~P\ref{premise:616e645f73656c66}} gives \(a_{1} \in s \land (b \in s \land (c \in s \land (\neg (a_{1} = b) \land (\neg (a_{1} = c) \land \neg (b = c)))))\).
This witness proves existence. Hence \(\exists c_{1}\, :\, \alpha,\; a_{1} \in s \land (b \in s \land (c_{1} \in s \land (a_{1} \ne b \land (a_{1} \ne c_{1} \land b \ne c_{1}))))\).
This witness proves existence. Hence \(\exists b_{1}\, :\, \alpha,\; \exists c_{1}\, :\, \alpha,\; a_{1} \in s \land (b_{1} \in s \land (c_{1} \in s \land (a_{1} \ne b_{1} \land (a_{1} \ne c_{1} \land b_{1} \ne c_{1}))))\).
This witness proves existence. Hence \(\exists a_{2}\, :\, \alpha,\; \exists b_{1}\, :\, \alpha,\; \exists c_{1}\, :\, \alpha,\; a_{2} \in s \land (b_{1} \in s \land (c_{1} \in s \land (a_{2} \ne b_{1} \land (a_{2} \ne c_{1} \land b_{1} \ne c_{1}))))\).
\end{proofstep}
Use the hypothesis \(t = \{a_{1}\} \cup \left(\{b\} \cup \{c\}\right)\).
Reversing the equality, we obtain \(\{a_{1}\} \cup \left(\{b\} \cup \{c\}\right) = t\).
Use the hypothesis \(t \subseteq s\).
Substituting equals, we obtain \(\exists a_{2}\, :\, \alpha,\; \exists b_{1}\, :\, \alpha,\; \exists c_{1}\, :\, \alpha,\; a_{2} \in s \land (b_{1} \in s \land (c_{1} \in s \land (a_{2} \ne b_{1} \land (a_{2} \ne c_{1} \land b_{1} \ne c_{1}))))\).
\end{proofstep}
Using the two conjuncts, we obtain \(\exists a_{2}\, :\, \alpha,\; \exists b_{1}\, :\, \alpha,\; \exists c_{1}\, :\, \alpha,\; a_{2} \in s \land (b_{1} \in s \land (c_{1} \in s \land (a_{2} \ne b_{1} \land (a_{2} \ne c_{1} \land b_{1} \ne c_{1}))))\).
\end{proofstep}
Using the two conjuncts, we obtain \(\exists a_{2}\, :\, \alpha,\; \exists b_{1}\, :\, \alpha,\; \exists c_{1}\, :\, \alpha,\; a_{2} \in s \land (b_{1} \in s \land (c_{1} \in s \land (a_{2} \ne b_{1} \land (a_{2} \ne c_{1} \land b_{1} \ne c_{1}))))\).
\end{proofstep}
Using the two conjuncts, we obtain \(\exists a_{2}\, :\, \alpha,\; \exists b_{1}\, :\, \alpha,\; \exists c_{1}\, :\, \alpha,\; a_{2} \in s \land (b_{1} \in s \land (c_{1} \in s \land (a_{2} \ne b_{1} \land (a_{2} \ne c_{1} \land b_{1} \ne c_{1}))))\).
\end{proofstep}
Using the two conjuncts, we obtain \(\exists a_{2}\, :\, \alpha,\; \exists b_{1}\, :\, \alpha,\; \exists c_{1}\, :\, \alpha,\; a_{2} \in s \land (b_{1} \in s \land (c_{1} \in s \land (a_{2} \ne b_{1} \land (a_{2} \ne c_{1} \land b_{1} \ne c_{1}))))\).
\end{proofstep}
Using the existential hypothesis, we obtain \(\exists a_{2}\, :\, \alpha,\; \exists b_{1}\, :\, \alpha,\; \exists c_{1}\, :\, \alpha,\; a_{2} \in s \land (b_{1} \in s \land (c_{1} \in s \land (a_{2} \ne b_{1} \land (a_{2} \ne c_{1} \land b_{1} \ne c_{1}))))\).
\end{proofstep}
Using the existential hypothesis, we obtain \(\exists a_{2}\, :\, \alpha,\; \exists b_{1}\, :\, \alpha,\; \exists c\, :\, \alpha,\; a_{2} \in s \land (b_{1} \in s \land (c \in s \land (a_{2} \ne b_{1} \land (a_{2} \ne c \land b_{1} \ne c))))\).
\end{proofstep}
Using the existential hypothesis, we obtain \(\exists a_{2}\, :\, \alpha,\; \exists b\, :\, \alpha,\; \exists c\, :\, \alpha,\; a_{2} \in s \land (b \in s \land (c \in s \land (a_{2} \ne b \land (a_{2} \ne c \land b \ne c))))\).
\end{proofstep}
Using the existential hypothesis, we obtain \(\exists a_{1}\, :\, \alpha,\; \exists b\, :\, \alpha,\; \exists c\, :\, \alpha,\; a_{1} \in s \land (b \in s \land (c \in s \land (a_{1} \ne b \land (a_{1} \ne c \land b \ne c))))\).
\end{proofstep}
Consider the following case.
\begin{proofstep}
Assume \(\exists a_{1}\, :\, \alpha,\; \exists b\, :\, \alpha,\; \exists c\, :\, \alpha,\; a_{1} \in s \land (b \in s \land (c \in s \land (a_{1} \ne b \land (a_{1} \ne c \land b \ne c))))\).
Use the hypothesis \(\exists a_{1}\, :\, \alpha,\; \exists b\, :\, \alpha,\; \exists c\, :\, \alpha,\; a_{1} \in s \land (b \in s \land (c \in s \land (a_{1} \ne b \land (a_{1} \ne c \land b \ne c))))\).
Consider the following case.
\begin{proofstep}
Fix \(a_{1}\) in \(\alpha\).
Assume \(\exists b\, :\, \alpha,\; \exists c\, :\, \alpha,\; a_{1} \in s \land (b \in s \land (c \in s \land (a_{1} \ne b \land (a_{1} \ne c \land b \ne c))))\).
Use the hypothesis \(\exists b\, :\, \alpha,\; \exists c\, :\, \alpha,\; a_{1} \in s \land (b \in s \land (c \in s \land (a_{1} \ne b \land (a_{1} \ne c \land b \ne c))))\).
Consider the following case.
\begin{proofstep}
Fix \(b\) in \(\alpha\).
Assume \(\exists c\, :\, \alpha,\; a_{1} \in s \land (b \in s \land (c \in s \land (a_{1} \ne b \land (a_{1} \ne c \land b \ne c))))\).
Use the hypothesis \(\exists c\, :\, \alpha,\; a_{1} \in s \land (b \in s \land (c \in s \land (a_{1} \ne b \land (a_{1} \ne c \land b \ne c))))\).
Consider the following case.
\begin{proofstep}
Fix \(c\) in \(\alpha\).
Assume \(a_{1} \in s \land (b \in s \land (c \in s \land (a_{1} \ne b \land (a_{1} \ne c \land b \ne c))))\).
Use the hypothesis \(a_{1} \in s \land (b \in s \land (c \in s \land (a_{1} \ne b \land (a_{1} \ne c \land b \ne c))))\).
Consider the following case.
\begin{proofstep}
Assume \(a_{1} \in s\).
Assume \(b \in s \land \left(c \in s \land \left(a_{1} \ne b \land \left(a_{1} \ne c \land b \ne c\right)\right)\right)\).
Use the hypothesis \(b \in s \land \left(c \in s \land \left(a_{1} \ne b \land \left(a_{1} \ne c \land b \ne c\right)\right)\right)\).
Consider the following case.
\begin{proofstep}
Assume \(b \in s\).
Assume \(c \in s \land \left(a_{1} \ne b \land \left(a_{1} \ne c \land b \ne c\right)\right)\).
Use the hypothesis \(c \in s \land \left(a_{1} \ne b \land \left(a_{1} \ne c \land b \ne c\right)\right)\).
Consider the following case.
\begin{proofstep}
Assume \(c \in s\).
Assume \(a_{1} \ne b \land \left(a_{1} \ne c \land b \ne c\right)\).
Use the hypothesis \(a_{1} \ne b \land \left(a_{1} \ne c \land b \ne c\right)\).
Consider the following case.
\begin{proofstep}
Assume \(a_{1} \ne b\).
Assume \(a_{1} \ne c \land b \ne c\).
Use the hypothesis \(a_{1} \ne c \land b \ne c\).
Consider the following case.
\begin{proofstep}
Assume \(a_{1} \ne c\).
Assume \(b \ne c\).
Choose the witness \(a_{1}\).
Choose the witness \(b\).
Choose the witness \(c\).
Choose the witness \(\{a_{1}\} \cup \left(\{b\} \cup \{c\}\right)\).
Rewriting with \paperlink{premise:46696e7365742e696e736572745f7375627365745f696666}{Prerequisite~P\ref{premise:46696e7365742e696e736572745f7375627365745f696666}}, \(a_{1} \in s\), \(b \in s\), \paperlink{premise:46696e7365742e73696e676c65746f6e5f7375627365745f696666}{Prerequisite~P\ref{premise:46696e7365742e73696e676c65746f6e5f7375627365745f696666}}, \(c \in s\), \paperlink{premise:616e645f73656c66}{Prerequisite~P\ref{premise:616e645f73656c66}}, \paperlink{premise:65715f66616c7365}{Prerequisite~P\ref{premise:65715f66616c7365}}, \(a_{1} \ne b\), \paperlink{premise:6e6f745f66616c73655f65715f74727565}{Prerequisite~P\ref{premise:6e6f745f66616c73655f65715f74727565}}, \(a_{1} \ne c\) and \(b \ne c\) gives \(\{a_{1}\} \cup (\{b\} \cup \{c\}) \subseteq s \land (\neg (a_{1} = b) \land (\neg (a_{1} = c) \land (\neg (b = c) \land \{a_{1}\} \cup (\{b\} \cup \{c\}) = \{a_{1}\} \cup (\{b\} \cup \{c\}))))\).
This witness proves existence. Hence \(\exists a_{2}\, :\, \mathcal{P}_{\mathrm{fin}}(\alpha),\; a_{2} \subseteq s \land (a_{1} \ne b \land (a_{1} \ne c \land (b \ne c \land a_{2} = \{a_{1}\} \cup (\{b\} \cup \{c\}))))\).
This witness proves existence. Hence \(\exists b_{1}\, :\, \alpha,\; \exists a_{2}\, :\, \mathcal{P}_{\mathrm{fin}}(\alpha),\; a_{2} \subseteq s \land (a_{1} \ne b \land (a_{1} \ne b_{1} \land (b \ne b_{1} \land a_{2} = \{a_{1}\} \cup (\{b\} \cup \{b_{1}\}))))\).
This witness proves existence. Hence \(\exists b_{1}\, :\, \alpha,\; \exists b_{2}\, :\, \alpha,\; \exists a_{2}\, :\, \mathcal{P}_{\mathrm{fin}}(\alpha),\; a_{2} \subseteq s \land (a_{1} \ne b_{1} \land (a_{1} \ne b_{2} \land (b_{1} \ne b_{2} \land a_{2} = \{a_{1}\} \cup (\{b_{1}\} \cup \{b_{2}\}))))\).
This witness proves existence. Hence \(\exists b_{1}\, :\, \alpha,\; \exists b_{2}\, :\, \alpha,\; \exists b_{3}\, :\, \alpha,\; \exists a_{2}\, :\, \mathcal{P}_{\mathrm{fin}}(\alpha),\; a_{2} \subseteq s \land (b_{1} \ne b_{2} \land (b_{1} \ne b_{3} \land (b_{2} \ne b_{3} \land a_{2} = \{b_{1}\} \cup (\{b_{2}\} \cup \{b_{3}\}))))\).
\end{proofstep}
Using the two conjuncts, we obtain \(\exists b_{1}\, :\, \alpha,\; \exists b_{2}\, :\, \alpha,\; \exists b_{3}\, :\, \alpha,\; \exists a_{2}\, :\, \mathcal{P}_{\mathrm{fin}}(\alpha),\; a_{2} \subseteq s \land (b_{1} \ne b_{2} \land (b_{1} \ne b_{3} \land (b_{2} \ne b_{3} \land a_{2} = \{b_{1}\} \cup (\{b_{2}\} \cup \{b_{3}\}))))\).
\end{proofstep}
Using the two conjuncts, we obtain \(\exists b_{1}\, :\, \alpha,\; \exists b_{2}\, :\, \alpha,\; \exists b_{3}\, :\, \alpha,\; \exists a_{2}\, :\, \mathcal{P}_{\mathrm{fin}}(\alpha),\; a_{2} \subseteq s \land (b_{1} \ne b_{2} \land (b_{1} \ne b_{3} \land (b_{2} \ne b_{3} \land a_{2} = \{b_{1}\} \cup (\{b_{2}\} \cup \{b_{3}\}))))\).
\end{proofstep}
Using the two conjuncts, we obtain \(\exists b_{1}\, :\, \alpha,\; \exists b_{2}\, :\, \alpha,\; \exists b_{3}\, :\, \alpha,\; \exists a_{2}\, :\, \mathcal{P}_{\mathrm{fin}}(\alpha),\; a_{2} \subseteq s \land (b_{1} \ne b_{2} \land (b_{1} \ne b_{3} \land (b_{2} \ne b_{3} \land a_{2} = \{b_{1}\} \cup (\{b_{2}\} \cup \{b_{3}\}))))\).
\end{proofstep}
Using the two conjuncts, we obtain \(\exists b_{1}\, :\, \alpha,\; \exists b_{2}\, :\, \alpha,\; \exists b_{3}\, :\, \alpha,\; \exists a_{2}\, :\, \mathcal{P}_{\mathrm{fin}}(\alpha),\; a_{2} \subseteq s \land (b_{1} \ne b_{2} \land (b_{1} \ne b_{3} \land (b_{2} \ne b_{3} \land a_{2} = \{b_{1}\} \cup (\{b_{2}\} \cup \{b_{3}\}))))\).
\end{proofstep}
Using the two conjuncts, we obtain \(\exists b_{1}\, :\, \alpha,\; \exists b_{2}\, :\, \alpha,\; \exists b_{3}\, :\, \alpha,\; \exists a_{2}\, :\, \mathcal{P}_{\mathrm{fin}}(\alpha),\; a_{2} \subseteq s \land (b_{1} \ne b_{2} \land (b_{1} \ne b_{3} \land (b_{2} \ne b_{3} \land a_{2} = \{b_{1}\} \cup (\{b_{2}\} \cup \{b_{3}\}))))\).
\end{proofstep}
Using the existential hypothesis, we obtain \(\exists b_{1}\, :\, \alpha,\; \exists b_{2}\, :\, \alpha,\; \exists b_{3}\, :\, \alpha,\; \exists a_{2}\, :\, \mathcal{P}_{\mathrm{fin}}(\alpha),\; a_{2} \subseteq s \land (b_{1} \ne b_{2} \land (b_{1} \ne b_{3} \land (b_{2} \ne b_{3} \land a_{2} = \{b_{1}\} \cup (\{b_{2}\} \cup \{b_{3}\}))))\).
\end{proofstep}
Using the existential hypothesis, we obtain \(\exists b\, :\, \alpha,\; \exists b_{1}\, :\, \alpha,\; \exists b_{2}\, :\, \alpha,\; \exists a_{2}\, :\, \mathcal{P}_{\mathrm{fin}}(\alpha),\; a_{2} \subseteq s \land (b \ne b_{1} \land (b \ne b_{2} \land (b_{1} \ne b_{2} \land a_{2} = \{b\} \cup (\{b_{1}\} \cup \{b_{2}\}))))\).
\end{proofstep}
Using the existential hypothesis, we obtain \(\exists b\, :\, \alpha,\; \exists b_{1}\, :\, \alpha,\; \exists b_{2}\, :\, \alpha,\; \exists a_{1}\, :\, \mathcal{P}_{\mathrm{fin}}(\alpha),\; a_{1} \subseteq s \land (b \ne b_{1} \land (b \ne b_{2} \land (b_{1} \ne b_{2} \land a_{1} = \{b\} \cup (\{b_{1}\} \cup \{b_{2}\}))))\).
\end{proofstep}
Together the two implications give the equivalence. Hence \((\exists b\, :\, \alpha,\; \exists b_{1}\, :\, \alpha,\; \exists b_{2}\, :\, \alpha,\; \exists a\, :\, \mathcal{P}_{\mathrm{fin}}(\alpha),\; a \subseteq s \land (b \ne b_{1} \land (b \ne b_{2} \land (b_{1} \ne b_{2} \land a = \{b\} \cup (\{b_{1}\} \cup \{b_{2}\}))))) \iff (\exists a\, :\, \alpha,\; \exists b\, :\, \alpha,\; \exists c\, :\, \alpha,\; a \in s \land (b \in s \land (c \in s \land (a \ne b \land (a \ne c \land b \ne c)))))\).
Rewriting with \paperlink{premise:6578697374735f636f6d6d}{Prerequisite~P\ref{premise:6578697374735f636f6d6d}}, \paperlink{premise:66756e657874}{\emph{Function extensionality}} (\paperlink{premise:66756e657874}{Prerequisite~P\ref{premise:66756e657874}}) and the identity established above in the preceding result gives \((\exists t\, :\, \mathcal{P}_{\mathrm{fin}}(\alpha),\; \exists x\, :\, \alpha,\; \exists x_{1}\, :\, \alpha,\; \exists x_{2}\, :\, \alpha,\; t \subseteq s \land (x \ne x_{1} \land (x \ne x_{2} \land (x_{1} \ne x_{2} \land t = \{x\} \cup (\{x_{1}\} \cup \{x_{2}\}))))) \iff (\exists a\, :\, \alpha,\; \exists b\, :\, \alpha,\; \exists c\, :\, \alpha,\; a \in s \land (b \in s \land (c \in s \land (a \ne b \land (a \ne c \land b \ne c)))))\).
Rewriting with \paperlink{premise:66756e657874}{\emph{Function extensionality}} (\paperlink{premise:66756e657874}{Prerequisite~P\ref{premise:66756e657874}}) and the identity established above in the preceding result gives \((\exists t\, :\, \mathcal{P}_{\mathrm{fin}}(\alpha),\; t \subseteq s \land (\exists x\, :\, \alpha,\; \exists y\, :\, \alpha,\; \exists z\, :\, \alpha,\; x \ne y \land (x \ne z \land (y \ne z \land t = \{x\} \cup (\{y\} \cup \{z\}))))) \iff (\exists a\, :\, \alpha,\; \exists b\, :\, \alpha,\; \exists c\, :\, \alpha,\; a \in s \land (b \in s \land (c \in s \land (a \ne b \land (a \ne c \land b \ne c)))))\).
Rewriting with \paperlink{premise:66756e657874}{\emph{Function extensionality}} (\paperlink{premise:66756e657874}{Prerequisite~P\ref{premise:66756e657874}}) and the identity established above in the preceding result gives \((\exists t\, :\, \mathcal{P}_{\mathrm{fin}}(\alpha),\; t \subseteq s \land |t| = 2 + 1) \iff (\exists a\, :\, \alpha,\; \exists b\, :\, \alpha,\; \exists c\, :\, \alpha,\; a \in s \land (b \in s \land (c \in s \land (a \ne b \land (a \ne c \land b \ne c)))))\).
Rewriting with \paperlink{theorem:104}{Lemma~\ref{theorem:104}} in the preceding result gives \(2 + 1 \le |s| \iff (\exists a\, :\, \alpha,\; \exists b\, :\, \alpha,\; \exists c\, :\, \alpha,\; a \in s \land (b \in s \land (c \in s \land (a \ne b \land (a \ne c \land b \ne c)))))\).
Rewriting with \paperlink{premise:4e61742e6c745f6966665f6164645f6f6e655f6c65}{Prerequisite~P\ref{premise:4e61742e6c745f6966665f6164645f6f6e655f6c65}} in the preceding result gives \(2 < |s| \iff (\exists a\, :\, \alpha,\; \exists b\, :\, \alpha,\; \exists c\, :\, \alpha,\; a \in s \land (b \in s \land (c \in s \land (a \ne b \land (a \ne c \land b \ne c)))))\).
\end{proof}

\begin{theorem}\label{theorem:124}
Let \(\alpha\) be a type. Let \(s\) be a finite subset of \(\alpha\). Then
\begin{center}
\(\displaystyle 2 < |s| \iff (\exists a\, :\, \alpha,\; a \in s \land (\exists b\, :\, \alpha,\; b \in s \land (\exists c\, :\, \alpha,\; c \in s \land (a \ne b \land (a \ne c \land b \ne c))))).\)
\end{center}
\end{theorem}

\begin{proof}
The substitution uses the following argument.
\begin{proofstep}
Fix \(a\) in \(\alpha\).
For the required implication, argue as follows.
\begin{proofstep}
Fix \(b\) in \(\alpha\).
By \paperlink{premise:6578697374735f616e645f6c656674}{Prerequisite~P\ref{premise:6578697374735f616e645f6c656674}}, we obtain \((\exists x\, :\, \alpha,\; a \in s \land (x_{1} \mapsto b \in s \land (x_{1} \in s \land (a \ne b \land (a \ne x_{1} \land b \ne x_{1}))))(x)) \iff a \in s \land (\exists x\, :\, \alpha,\; (x_{1} \mapsto b \in s \land (x_{1} \in s \land (a \ne b \land (a \ne x_{1} \land b \ne x_{1}))))(x))\).
By propositional extensionality, we obtain \((\exists x\, :\, \alpha,\; a \in s \land (b \in s \land (x \in s \land (a \ne b \land (a \ne x \land b \ne x))))) = (a \in s \land (\exists x\, :\, \alpha,\; b \in s \land (x \in s \land (a \ne b \land (a \ne x \land b \ne x)))))\).
By \paperlink{premise:6578697374735f616e645f6c656674}{Prerequisite~P\ref{premise:6578697374735f616e645f6c656674}}, we obtain \((\exists x\, :\, \alpha,\; b \in s \land (x_{1} \mapsto x_{1} \in s \land (a \ne b \land (a \ne x_{1} \land b \ne x_{1})))(x)) \iff b \in s \land (\exists x\, :\, \alpha,\; (x_{1} \mapsto x_{1} \in s \land (a \ne b \land (a \ne x_{1} \land b \ne x_{1})))(x))\).
By propositional extensionality, we obtain \((\exists x\, :\, \alpha,\; b \in s \land (x \in s \land (a \ne b \land (a \ne x \land b \ne x)))) = (b \in s \land (\exists x\, :\, \alpha,\; x \in s \land (a \ne b \land (a \ne x \land b \ne x))))\).
Applying the same function to equal arguments, we obtain \((a \in s \land (\exists x\, :\, \alpha,\; b \in s \land (x \in s \land (a \ne b \land (a \ne x \land b \ne x))))) = (a \in s \land (b \in s \land (\exists x\, :\, \alpha,\; x \in s \land (a \ne b \land (a \ne x \land b \ne x)))))\).
Chaining the equalities, we obtain \((\exists x\, :\, \alpha,\; a \in s \land (b \in s \land (x \in s \land (a \ne b \land (a \ne x \land b \ne x))))) = (a \in s \land (b \in s \land (\exists x\, :\, \alpha,\; x \in s \land (a \ne b \land (a \ne x \land b \ne x)))))\).
\end{proofstep}
By function extensionality, we obtain \((x \mapsto \exists x_{1}\, :\, \alpha,\; a \in s \land (x \in s \land (x_{1} \in s \land (a \ne x \land (a \ne x_{1} \land x \ne x_{1}))))) = (x \mapsto a \in s \land (x \in s \land (\exists x_{1}\, :\, \alpha,\; x_{1} \in s \land (a \ne x \land (a \ne x_{1} \land x \ne x_{1})))))\).
Applying the same function to equal arguments, we obtain \((\exists x\, :\, \alpha,\; \exists x_{1}\, :\, \alpha,\; a \in s \land (x \in s \land (x_{1} \in s \land (a \ne x \land (a \ne x_{1} \land x \ne x_{1}))))) = (\exists x\, :\, \alpha,\; a \in s \land (x \in s \land (\exists x_{1}\, :\, \alpha,\; x_{1} \in s \land (a \ne x \land (a \ne x_{1} \land x \ne x_{1})))))\).
By \paperlink{premise:6578697374735f616e645f6c656674}{Prerequisite~P\ref{premise:6578697374735f616e645f6c656674}}, we obtain \((\exists x\, :\, \alpha,\; a \in s \land (x_{1} \mapsto x_{1} \in s \land (\exists x_{2}\, :\, \alpha,\; x_{2} \in s \land (a \ne x_{1} \land (a \ne x_{2} \land x_{1} \ne x_{2}))))(x)) \iff a \in s \land (\exists x\, :\, \alpha,\; (x_{1} \mapsto x_{1} \in s \land (\exists x_{2}\, :\, \alpha,\; x_{2} \in s \land (a \ne x_{1} \land (a \ne x_{2} \land x_{1} \ne x_{2}))))(x))\).
By propositional extensionality, we obtain \((\exists x\, :\, \alpha,\; a \in s \land (x \in s \land (\exists x_{1}\, :\, \alpha,\; x_{1} \in s \land (a \ne x \land (a \ne x_{1} \land x \ne x_{1}))))) = (a \in s \land (\exists x\, :\, \alpha,\; x \in s \land (\exists x_{1}\, :\, \alpha,\; x_{1} \in s \land (a \ne x \land (a \ne x_{1} \land x \ne x_{1})))))\).
Chaining the equalities, we obtain \((\exists x\, :\, \alpha,\; \exists x_{1}\, :\, \alpha,\; a \in s \land (x \in s \land (x_{1} \in s \land (a \ne x \land (a \ne x_{1} \land x \ne x_{1}))))) = (a \in s \land (\exists x\, :\, \alpha,\; x \in s \land (\exists x_{1}\, :\, \alpha,\; x_{1} \in s \land (a \ne x \land (a \ne x_{1} \land x \ne x_{1})))))\).
\end{proofstep}
Rewriting with \paperlink{premise:66756e657874}{\emph{Function extensionality}} (\paperlink{premise:66756e657874}{Prerequisite~P\ref{premise:66756e657874}}), \paperlink{premise:6966665f73656c66}{Prerequisite~P\ref{premise:6966665f73656c66}} and the identity established above gives \((\exists a\, :\, \alpha,\; \exists b\, :\, \alpha,\; \exists c\, :\, \alpha,\; a \in s \land (b \in s \land (c \in s \land (a \ne b \land (a \ne c \land b \ne c))))) \iff (\exists a\, :\, \alpha,\; a \in s \land (\exists b\, :\, \alpha,\; b \in s \land (\exists c\, :\, \alpha,\; c \in s \land (a \ne b \land (a \ne c \land b \ne c)))))\).
Rewriting with \paperlink{theorem:123}{Lemma~\ref{theorem:123}} in the preceding result gives \(2 < |s| \iff (\exists a\, :\, \alpha,\; a \in s \land (\exists b\, :\, \alpha,\; b \in s \land (\exists c\, :\, \alpha,\; c \in s \land (a \ne b \land (a \ne c \land b \ne c)))))\).
\end{proof}

\begin{theorem}\label{theorem:125}
Let \(\alpha\) be a type. Let \(p\) be a function of type \(\mathcal{P}_{\mathrm{fin}}(\alpha) \to \mathsf{Sort}\). Let \(H\) be a function of type \(\Pi_{s\, :\, \mathcal{P}_{\mathrm{fin}}(\alpha)}\, (\Pi_{t\, :\, \mathcal{P}_{\mathrm{fin}}(\alpha)}\, t \subsetneq s \to p(t)) \to p(s)\). Let \(s\) be a finite subset of \(\alpha\). Then
\begin{center}
\(\displaystyle \operatorname{strong\,Induction}(H,\allowbreak{} s) = H(s,\allowbreak{} t \mapsto x \mapsto \operatorname{strong\,Induction}(H,\allowbreak{} t)).\)
\end{center}
\end{theorem}

\begin{proof}
By definition, we obtain \(H(s,\allowbreak{} t \mapsto h \mapsto \operatorname{strong\,Induction}(H,\allowbreak{} t)) = H(s,\allowbreak{} t \mapsto x \mapsto \operatorname{strong\,Induction}(H,\allowbreak{} t))\).
Rewriting with \paperlink{defeq:2}{defining equation~D\ref{defeq:2}} in the preceding result gives \(\operatorname{strong\,Induction}(H,\allowbreak{} s) = H(s,\allowbreak{} t \mapsto x \mapsto \operatorname{strong\,Induction}(H,\allowbreak{} t))\).
\end{proof}

\begin{theorem}\label{theorem:126}
Let \(\alpha\) be a type. Let \(p\) be a function of type \(\mathcal{P}_{\mathrm{fin}}(\alpha) \to \mathsf{Sort}\). Let \(s\) be a finite subset of \(\alpha\). Let \(H\) be a function of type \(\Pi_{s_{1}\, :\, \mathcal{P}_{\mathrm{fin}}(\alpha)}\, (\Pi_{t\, :\, \mathcal{P}_{\mathrm{fin}}(\alpha)}\, t \subsetneq s_{1} \to p(t)) \to p(s_{1})\). Then
\begin{center}
\(\displaystyle \operatorname{strong\,Induction\,On}(s,\allowbreak{} H) = H(s,\allowbreak{} t \mapsto x \mapsto \operatorname{strong\,Induction\,On}(t,\allowbreak{} H)).\)
\end{center}
\end{theorem}

\begin{proof}
By definition, we obtain \(H(s,\allowbreak{} t \mapsto h \mapsto \operatorname{strong\,Induction}(H,\allowbreak{} t)) = H(s,\allowbreak{} t \mapsto x \mapsto \operatorname{strong\,Induction}(H,\allowbreak{} t))\).
Rewriting with \paperlink{defeq:2}{defining equation~D\ref{defeq:2}} in the preceding result gives \(\operatorname{strong\,Induction\,On}(s,\allowbreak{} H) = H(s,\allowbreak{} t \mapsto x \mapsto \operatorname{strong\,Induction\,On}(t,\allowbreak{} H))\).
\end{proof}

\begin{theorem}\label{theorem:127}
Let \(\alpha\) be a type. Assume equality on \(\alpha\) is decidable. Let \(p\) be a function of type \(\mathcal{P}_{\mathrm{fin}}(\alpha) \to \mathsf{Prop}\). Let \(s\) be a finite subset of \(\alpha\). Assume \(p\left(\varnothing\right)\). Assume \(\forall a\, :\, \alpha,\; \forall s_{1}\, :\, \mathcal{P}_{\mathrm{fin}}(\alpha),\; a \notin s_{1} \implies (\forall t\, :\, \mathcal{P}_{\mathrm{fin}}(\alpha),\; t \subseteq s_{1} \implies p(t)) \implies p(\{a\} \cup s_{1})\). Then \(p\left(s\right)\).
\end{theorem}

\begin{proof}
Use strong induction on \(s\), assuming the claim for every strict subset in the induction step.
Consider the following case.
\begin{proofstep}
Fix \(s_{1}\) in \(\mathcal{P}_{\mathrm{fin}}(\alpha)\).
For the required implication, argue as follows.
\begin{proofstep}
Assume \(\forall t\, :\, \mathcal{P}_{\mathrm{fin}}(\alpha),\; t \subsetneq \varnothing \implies p\left(t\right)\).
Use the hypothesis \(p\left(\varnothing\right)\).
\end{proofstep}
For the required implication, argue as follows.
\begin{proofstep}
Fix \(a\) in \(\alpha\).
Fix \(s_{2}\) in \(\mathcal{P}_{\mathrm{fin}}(\alpha)\).
Assume \(a \notin s_{2}\).
Assume \((\forall t\, :\, \mathcal{P}_{\mathrm{fin}}(\alpha),\; t \subsetneq s_{2} \implies p(t)) \implies p(s_{2})\).
Assume as induction hypothesis that \(\forall t\, :\, \mathcal{P}_{\mathrm{fin}}(\alpha),\; t \subsetneq \{a\} \cup s_{2} \implies p\left(t\right)\).
Use the hypothesis \(a \notin s_{2}\).
For the required implication, argue as follows.
\begin{proofstep}
Fix \(t\) in \(\mathcal{P}_{\mathrm{fin}}(\alpha)\).
Assume \(t \subseteq s_{2}\).
Use the hypothesis \(t \subseteq s_{2}\).
Use the hypothesis \(a \notin s_{2}\).
By \paperlink{premise:46696e7365742e737375627365745f696e73657274}{Prerequisite~P\ref{premise:46696e7365742e737375627365745f696e73657274}}, we obtain \(s_{2} \subsetneq \{a\} \cup s_{2}\).
By \paperlink{premise:6c745f6f665f6c655f6f665f6c74}{Prerequisite~P\ref{premise:6c745f6f665f6c655f6f665f6c74}} with \(\mathcal{P}_{\mathrm{fin}}(\alpha)\) as its type parameter \(\alpha\), we obtain \(t < \{a\} \cup s_{2}\).
Applying the induction hypothesis, we obtain \(p\left(t\right)\).
\end{proofstep}
Applying the hypothesis \(\forall a\, :\, \alpha,\; \forall s_{1}\, :\, \mathcal{P}_{\mathrm{fin}}(\alpha),\; a \notin s_{1} \implies (\forall t\, :\, \mathcal{P}_{\mathrm{fin}}(\alpha),\; t \subseteq s_{1} \implies p(t)) \implies p(\{a\} \cup s_{1})\), we obtain \(p\left(\{a\} \cup s_{2}\right)\).
\end{proofstep}
By \paperlink{premise:46696e7365742e696e64756374696f6e5f6f6e}{Prerequisite~P\ref{premise:46696e7365742e696e64756374696f6e5f6f6e}} with \(s_{1}\) in place of its parameter \(s\), we obtain \((\forall t\, :\, \mathcal{P}_{\mathrm{fin}}(\alpha),\; t \subsetneq s_{1} \implies p(t)) \implies p(s_{1})\).
\end{proofstep}
By strong induction on strict subsets, we obtain \(p\left(s\right)\).
\end{proof}

\begin{theorem}\label{theorem:128}
Let \(\alpha\) be a type. Let \(p\) be a function of type \(\Pi_{s\, :\, \mathcal{P}_{\mathrm{fin}}(\alpha)}\, s\ne\varnothing \to \mathsf{Prop}\). Assume \(\forall a\, :\, \alpha,\; p\left(\{a\}, \operatorname{singleton\,nonempty}\left(a\right)\right)\). Assume \(\forall s\, :\, \mathcal{P}_{\mathrm{fin}}(\alpha),\; \forall \mathit{hs}\, :\, s\text{ contains two distinct elements},\; (\forall t\, :\, \mathcal{P}_{\mathrm{fin}}(\alpha),\; \forall \mathit{ht}\, :\, t\ne\varnothing,\; t \subsetneq s \implies p(t,\allowbreak{} \mathit{ht})) \implies p(s,\allowbreak{} \operatorname{nonempty}(\mathit{hs}))\). Let \(s\) be a finite subset of \(\alpha\). Let \(\mathit{hs}\) be a proof of \(s\ne\varnothing\). Then \(p\left(s, \mathit{hs}\right)\).
\end{theorem}

\begin{proof}
A descending chain would make the natural-number measure \(x \mapsto \left|\pi_1(x)\right|\) strictly decrease at each step, so the relation is well-founded.
Proceed by well-founded induction to prove \(p\left(s, \mathit{hs}\right)\). For an arbitrary element we may assume the claim for each strictly smaller element.
\begin{proofstep}
Fix \(x\) in \(\{(s_{1}, h)\mid s_{1}\, :\, \mathcal{P}_{\mathrm{fin}}(\alpha),\; h\text{ proves }s_{1}\ne\varnothing\}\).
Assume \(\forall y\, :\, \{(s_{1}, h)\mid s_{1}\, :\, \mathcal{P}_{\mathrm{fin}}(\alpha),\; h\text{ proves }s_{1}\ne\varnothing\},\; |\pi_1(y)| < |\pi_1(x)| \implies p(\pi_1(y),\allowbreak{} \pi_2(y))\).
Use the hypothesis \(\forall y\, :\, \{(s_{1}, h)\mid s_{1}\, :\, \mathcal{P}_{\mathrm{fin}}(\alpha),\; h\text{ proves }s_{1}\ne\varnothing\},\; |\pi_1(y)| < |\pi_1(x)| \implies p(\pi_1(y),\allowbreak{} \pi_2(y))\).
Write the pair \(x\) in terms of its two components.
\begin{proofstep}
Fix \(s_{1}\) in \(\mathcal{P}_{\mathrm{fin}}(\alpha)\).
Assume \(s_{1}\ne\varnothing\).
Assume \(\forall y\, :\, \{(s_{2}, h)\mid s_{2}\, :\, \mathcal{P}_{\mathrm{fin}}(\alpha),\; h\text{ proves }s_{2}\ne\varnothing\},\; |\pi_1(y)| < |s_{1}| \implies p(\pi_1(y),\allowbreak{} \pi_2(y))\).
Use the hypothesis \(s_{1}\ne\varnothing\).
By \paperlink{premise:46696e7365742e4e6f6e656d7074792e6578697374735f65715f73696e676c65746f6e5f6f725f6e6f6e7472697669616c}{Prerequisite~P\ref{premise:46696e7365742e4e6f6e656d7074792e6578697374735f65715f73696e676c65746f6e5f6f725f6e6f6e7472697669616c}}, we obtain \(\left(\exists a_{2}\, :\, \alpha,\; s_{1} = \{a_{2}\}\right) \lor s_{1}\text{ contains two distinct elements}\).
Consider the following case.
\begin{proofstep}
Assume \(\exists a_{2}\, :\, \alpha,\; s_{1} = \{a_{2}\}\).
Use the hypothesis \(\exists a_{2}\, :\, \alpha,\; s_{1} = \{a_{2}\}\).
Consider the following case.
\begin{proofstep}
Fix \(a_{2}\) in \(\alpha\).
Assume \(s_{1} = \{a_{2}\}\).
For the required implication, argue as follows.
\begin{proofstep}
Assume \(\{a_{2}\}\ne\varnothing\).
Applying the hypothesis \(\forall a\, :\, \alpha,\; p\left(\{a\}, \operatorname{singleton\,nonempty}\left(a\right)\right)\), we obtain \(p\left(\{a_{2}\}, \operatorname{singleton\,nonempty}\left(a_{2}\right)\right)\).
\end{proofstep}
Use the hypothesis \(s_{1} = \{a_{2}\}\).
Reversing the equality, we obtain \(\{a_{2}\} = s_{1}\).
Use the hypothesis \(s_{1}\ne\varnothing\).
Substituting equals, we obtain \(p\left(s_{1}, \mathit{hs}_{1}\right)\).
\end{proofstep}
Using the existential hypothesis, we obtain \(p\left(s_{1}, \mathit{hs}_{1}\right)\).
\end{proofstep}
Consider the following case.
\begin{proofstep}
Assume \(s_{1}\text{ contains two distinct elements}\).
Use the hypothesis \(s_{1}\text{ contains two distinct elements}\).
For the required implication, argue as follows.
\begin{proofstep}
Fix \(t\) in \(\mathcal{P}_{\mathrm{fin}}(\alpha)\).
Assume \(t\ne\varnothing\).
Assume \(t \subsetneq s_{1}\).
By \paperlink{theorem:64}{Lemma~\ref{theorem:64}}, we obtain \(\left|t\right| < \left|s_{1}\right|\).
Use the previously proved fact \(\left|t\right| < \left|s_{1}\right|\).
Applying the hypothesis \(\forall y\, :\, \{(s_{2}, h)\mid s_{2}\, :\, \mathcal{P}_{\mathrm{fin}}(\alpha),\; h\text{ proves }s_{2}\ne\varnothing\},\; |\pi_1(y)| < |s_{1}| \implies p(\pi_1(y),\allowbreak{} \pi_2(y))\), we obtain \(p\left(t, \mathit{ht}\right)\).
\end{proofstep}
Applying the hypothesis \(\forall s\, :\, \mathcal{P}_{\mathrm{fin}}(\alpha),\; \forall \mathit{hs}\, :\, s\text{ contains two distinct elements},\; (\forall t\, :\, \mathcal{P}_{\mathrm{fin}}(\alpha),\; \forall \mathit{ht}\, :\, t\ne\varnothing,\; t \subsetneq s \implies p(t,\allowbreak{} \mathit{ht})) \implies p(s,\allowbreak{} \operatorname{nonempty}(\mathit{hs}))\), we obtain \(p\left(s_{1}, \operatorname{nonempty}\left(\mathit{hs}_{2}\right)\right)\).
\end{proofstep}
The two alternatives exhaust the cases. Hence \(p\left(s_{1}, \mathit{hs}_{1}\right)\).
\end{proofstep}
This proves \(p\left(\pi_1(x), \pi_2(x)\right)\).
\end{proofstep}
The induction gives \(p\left(s, \mathit{hs}\right)\).
\end{proof}

\begin{theorem}\label{theorem:129}
Let \(\alpha\) be a type. Let \(n\) be a natural number. Let \(p\) be a function of type \(\mathcal{P}_{\mathrm{fin}}(\alpha) \to \mathsf{Sort}\). Let \(H\) be a function of type \(\Pi_{t_{1}\, :\, \mathcal{P}_{\mathrm{fin}}(\alpha)}\, (\Pi_{t_{2}\, :\, \mathcal{P}_{\mathrm{fin}}(\alpha)}\, |t_{2}| \le n \to t_{1} \subsetneq t_{2} \to p(t_{2})) \to |t_{1}| \le n \to p(t_{1})\). Let \(s\) be a finite subset of \(\alpha\). Then
\begin{center}
\(\displaystyle \operatorname{strong\,Downward\,Induction}(H,\allowbreak{} s) = H(s,\allowbreak{} t \mapsto \mathit{ht} \mapsto x \mapsto \operatorname{strong\,Downward\,Induction}(H,\allowbreak{} t)).\)
\end{center}
\end{theorem}

\begin{proof}
By definition, we obtain \(H(s,\allowbreak{} t \mapsto \mathit{ht} \mapsto h \mapsto  \operatorname{strong\,Downward\,Induction}(H,\allowbreak{} t)|_{\mathit{this}=\operatorname{Finset\mathord{:}card\,lt\,card}(h)}) = H(s,\allowbreak{} t \mapsto \mathit{ht} \mapsto x \mapsto \operatorname{strong\,Downward\,Induction}(H,\allowbreak{} t))\).
Rewriting with \paperlink{defeq:4}{defining equation~D\ref{defeq:4}} in the preceding result gives \(\operatorname{strong\,Downward\,Induction}(H,\allowbreak{} s) = H(s,\allowbreak{} t \mapsto \mathit{ht} \mapsto x \mapsto \operatorname{strong\,Downward\,Induction}(H,\allowbreak{} t))\).
\end{proof}

\begin{theorem}\label{theorem:130}
Let \(\alpha\) be a type. Let \(n\) be a natural number. Let \(p\) be a function of type \(\mathcal{P}_{\mathrm{fin}}(\alpha) \to \mathsf{Sort}\). Let \(s\) be a finite subset of \(\alpha\). Let \(H\) be a function of type \(\Pi_{t_{1}\, :\, \mathcal{P}_{\mathrm{fin}}(\alpha)}\, (\Pi_{t_{2}\, :\, \mathcal{P}_{\mathrm{fin}}(\alpha)}\, |t_{2}| \le n \to t_{1} \subsetneq t_{2} \to p(t_{2})) \to |t_{1}| \le n \to p(t_{1})\). Then
\begin{center}
\(\displaystyle (a \mapsto \operatorname{strong\,Downward\,Induction\,On}(s,\allowbreak{} H)) = H(s,\allowbreak{} t \mapsto \mathit{ht} \mapsto x \mapsto \operatorname{strong\,Downward\,Induction\,On}(t,\allowbreak{} H)).\)
\end{center}
\end{theorem}

\begin{proof}
By definition, we obtain \((a \mapsto H(s,\allowbreak{} t \mapsto \mathit{ht} \mapsto h \mapsto  \operatorname{strong\,Downward\,Induction}(H,\allowbreak{} t)|_{\mathit{this}=\operatorname{Finset\mathord{:}card\,lt\,card}(h)},\allowbreak{} a)) = H(s,\allowbreak{} t \mapsto \mathit{ht} \mapsto x \mapsto \operatorname{strong\,Downward\,Induction}(H,\allowbreak{} t))\).
Rewriting with \paperlink{defeq:4}{defining equation~D\ref{defeq:4}} in the preceding result gives \((a \mapsto \operatorname{strong\,Downward\,Induction\,On}(s,\allowbreak{} H)) = H(s,\allowbreak{} t \mapsto \mathit{ht} \mapsto x \mapsto \operatorname{strong\,Downward\,Induction\,On}(t,\allowbreak{} H))\).
\end{proof}

\begin{theorem}\label{theorem:131}
Let \(\alpha\) be a type. Then \(\operatorname{Well\,Founded}\left(x \mapsto y \mapsto x < y\right)\).
\end{theorem}

\begin{proof}
Fix \(x\) in \(\mathcal{P}_{\mathrm{fin}}(\alpha)\).
Fix \(x_{1}\) in \(\mathcal{P}_{\mathrm{fin}}(\alpha)\).
Assume \(x < x_{1}\).
By \paperlink{theorem:64}{Lemma~\ref{theorem:64}}, we obtain \(\left|x\right| < \left|x_{1}\right|\).
For the required implication, argue as follows.
\begin{proofstep}
Fix \(x\) in \(\mathcal{P}_{\mathrm{fin}}(\alpha)\).
Fix \(y\) in \(\mathcal{P}_{\mathrm{fin}}(\alpha)\).
Applying the hypothesis \(\operatorname{Subrelation}(x \mapsto y \mapsto x < y,\allowbreak{} \operatorname{Inv\,Image}(\mathit{x1} \mapsto \mathit{x2} \mapsto \mathit{x1} < \mathit{x2},\allowbreak{} \operatorname{Finset\mathord{:}card}))\), we obtain \(x < y \implies \operatorname{Inv\,Image}(\mathit{x1} \mapsto \mathit{x2} \mapsto \mathit{x1} < \mathit{x2},\allowbreak{} \operatorname{Finset\mathord{:}card},\allowbreak{} x,\allowbreak{} y)\).
\end{proofstep}
By \paperlink{premise:57656c6c466f756e64656452656c6174696f6e2e7766}{Prerequisite~P\ref{premise:57656c6c466f756e64656452656c6174696f6e2e7766}} with \(\mathbb{N}\) as its type parameter \(\alpha\), we obtain \(\operatorname{Well\,Founded}\left(\operatorname{rel}\right)\).
By \paperlink{premise:496e76496d6167652e7766}{Prerequisite~P\ref{premise:496e76496d6167652e7766}} with \(\mathcal{P}_{\mathrm{fin}}(\alpha)\) as its type parameter \(\alpha\) and \(\beta = \mathbb{N}\) and \(f = \operatorname{Finset\mathord{:}card}\), we obtain \(\operatorname{Well\,Founded}(\operatorname{Inv\,Image}(\mathit{x1} \mapsto \mathit{x2} \mapsto \mathit{x1} < \mathit{x2},\allowbreak{} \operatorname{Finset\mathord{:}card}))\).
By \paperlink{premise:53756272656c6174696f6e2e7766}{Prerequisite~P\ref{premise:53756272656c6174696f6e2e7766}} with \(\mathcal{P}_{\mathrm{fin}}(\alpha)\) as its type parameter \(\alpha\), we obtain \(\operatorname{Well\,Founded}\left(x \mapsto y \mapsto x < y\right)\).
\end{proof}

\end{document}
